% Options for packages loaded elsewhere
% Options for packages loaded elsewhere
\PassOptionsToPackage{unicode}{hyperref}
\PassOptionsToPackage{hyphens}{url}
\PassOptionsToPackage{dvipsnames,svgnames,x11names}{xcolor}
%
\documentclass[
  letterpaper,
  DIV=11,
  numbers=noendperiod]{scrreprt}
\usepackage{xcolor}
\usepackage{amsmath,amssymb}
\setcounter{secnumdepth}{5}
\usepackage{iftex}
\ifPDFTeX
  \usepackage[T1]{fontenc}
  \usepackage[utf8]{inputenc}
  \usepackage{textcomp} % provide euro and other symbols
\else % if luatex or xetex
  \usepackage{unicode-math} % this also loads fontspec
  \defaultfontfeatures{Scale=MatchLowercase}
  \defaultfontfeatures[\rmfamily]{Ligatures=TeX,Scale=1}
\fi
\usepackage{lmodern}
\ifPDFTeX\else
  % xetex/luatex font selection
\fi
% Use upquote if available, for straight quotes in verbatim environments
\IfFileExists{upquote.sty}{\usepackage{upquote}}{}
\IfFileExists{microtype.sty}{% use microtype if available
  \usepackage[]{microtype}
  \UseMicrotypeSet[protrusion]{basicmath} % disable protrusion for tt fonts
}{}
\makeatletter
\@ifundefined{KOMAClassName}{% if non-KOMA class
  \IfFileExists{parskip.sty}{%
    \usepackage{parskip}
  }{% else
    \setlength{\parindent}{0pt}
    \setlength{\parskip}{6pt plus 2pt minus 1pt}}
}{% if KOMA class
  \KOMAoptions{parskip=half}}
\makeatother
% Make \paragraph and \subparagraph free-standing
\makeatletter
\ifx\paragraph\undefined\else
  \let\oldparagraph\paragraph
  \renewcommand{\paragraph}{
    \@ifstar
      \xxxParagraphStar
      \xxxParagraphNoStar
  }
  \newcommand{\xxxParagraphStar}[1]{\oldparagraph*{#1}\mbox{}}
  \newcommand{\xxxParagraphNoStar}[1]{\oldparagraph{#1}\mbox{}}
\fi
\ifx\subparagraph\undefined\else
  \let\oldsubparagraph\subparagraph
  \renewcommand{\subparagraph}{
    \@ifstar
      \xxxSubParagraphStar
      \xxxSubParagraphNoStar
  }
  \newcommand{\xxxSubParagraphStar}[1]{\oldsubparagraph*{#1}\mbox{}}
  \newcommand{\xxxSubParagraphNoStar}[1]{\oldsubparagraph{#1}\mbox{}}
\fi
\makeatother

\usepackage{color}
\usepackage{fancyvrb}
\newcommand{\VerbBar}{|}
\newcommand{\VERB}{\Verb[commandchars=\\\{\}]}
\DefineVerbatimEnvironment{Highlighting}{Verbatim}{commandchars=\\\{\}}
% Add ',fontsize=\small' for more characters per line
\usepackage{framed}
\definecolor{shadecolor}{RGB}{241,243,245}
\newenvironment{Shaded}{\begin{snugshade}}{\end{snugshade}}
\newcommand{\AlertTok}[1]{\textcolor[rgb]{0.68,0.00,0.00}{#1}}
\newcommand{\AnnotationTok}[1]{\textcolor[rgb]{0.37,0.37,0.37}{#1}}
\newcommand{\AttributeTok}[1]{\textcolor[rgb]{0.40,0.45,0.13}{#1}}
\newcommand{\BaseNTok}[1]{\textcolor[rgb]{0.68,0.00,0.00}{#1}}
\newcommand{\BuiltInTok}[1]{\textcolor[rgb]{0.00,0.23,0.31}{#1}}
\newcommand{\CharTok}[1]{\textcolor[rgb]{0.13,0.47,0.30}{#1}}
\newcommand{\CommentTok}[1]{\textcolor[rgb]{0.37,0.37,0.37}{#1}}
\newcommand{\CommentVarTok}[1]{\textcolor[rgb]{0.37,0.37,0.37}{\textit{#1}}}
\newcommand{\ConstantTok}[1]{\textcolor[rgb]{0.56,0.35,0.01}{#1}}
\newcommand{\ControlFlowTok}[1]{\textcolor[rgb]{0.00,0.23,0.31}{\textbf{#1}}}
\newcommand{\DataTypeTok}[1]{\textcolor[rgb]{0.68,0.00,0.00}{#1}}
\newcommand{\DecValTok}[1]{\textcolor[rgb]{0.68,0.00,0.00}{#1}}
\newcommand{\DocumentationTok}[1]{\textcolor[rgb]{0.37,0.37,0.37}{\textit{#1}}}
\newcommand{\ErrorTok}[1]{\textcolor[rgb]{0.68,0.00,0.00}{#1}}
\newcommand{\ExtensionTok}[1]{\textcolor[rgb]{0.00,0.23,0.31}{#1}}
\newcommand{\FloatTok}[1]{\textcolor[rgb]{0.68,0.00,0.00}{#1}}
\newcommand{\FunctionTok}[1]{\textcolor[rgb]{0.28,0.35,0.67}{#1}}
\newcommand{\ImportTok}[1]{\textcolor[rgb]{0.00,0.46,0.62}{#1}}
\newcommand{\InformationTok}[1]{\textcolor[rgb]{0.37,0.37,0.37}{#1}}
\newcommand{\KeywordTok}[1]{\textcolor[rgb]{0.00,0.23,0.31}{\textbf{#1}}}
\newcommand{\NormalTok}[1]{\textcolor[rgb]{0.00,0.23,0.31}{#1}}
\newcommand{\OperatorTok}[1]{\textcolor[rgb]{0.37,0.37,0.37}{#1}}
\newcommand{\OtherTok}[1]{\textcolor[rgb]{0.00,0.23,0.31}{#1}}
\newcommand{\PreprocessorTok}[1]{\textcolor[rgb]{0.68,0.00,0.00}{#1}}
\newcommand{\RegionMarkerTok}[1]{\textcolor[rgb]{0.00,0.23,0.31}{#1}}
\newcommand{\SpecialCharTok}[1]{\textcolor[rgb]{0.37,0.37,0.37}{#1}}
\newcommand{\SpecialStringTok}[1]{\textcolor[rgb]{0.13,0.47,0.30}{#1}}
\newcommand{\StringTok}[1]{\textcolor[rgb]{0.13,0.47,0.30}{#1}}
\newcommand{\VariableTok}[1]{\textcolor[rgb]{0.07,0.07,0.07}{#1}}
\newcommand{\VerbatimStringTok}[1]{\textcolor[rgb]{0.13,0.47,0.30}{#1}}
\newcommand{\WarningTok}[1]{\textcolor[rgb]{0.37,0.37,0.37}{\textit{#1}}}

\usepackage{longtable,booktabs,array}
\newcounter{none} % for unnumbered tables
\usepackage{calc} % for calculating minipage widths
% Correct order of tables after \paragraph or \subparagraph
\usepackage{etoolbox}
\makeatletter
\patchcmd\longtable{\par}{\if@noskipsec\mbox{}\fi\par}{}{}
\makeatother
% Allow footnotes in longtable head/foot
\IfFileExists{footnotehyper.sty}{\usepackage{footnotehyper}}{\usepackage{footnote}}
\makesavenoteenv{longtable}
\usepackage{graphicx}
\makeatletter
\newsavebox\pandoc@box
\newcommand*\pandocbounded[1]{% scales image to fit in text height/width
  \sbox\pandoc@box{#1}%
  \Gscale@div\@tempa{\textheight}{\dimexpr\ht\pandoc@box+\dp\pandoc@box\relax}%
  \Gscale@div\@tempb{\linewidth}{\wd\pandoc@box}%
  \ifdim\@tempb\p@<\@tempa\p@\let\@tempa\@tempb\fi% select the smaller of both
  \ifdim\@tempa\p@<\p@\scalebox{\@tempa}{\usebox\pandoc@box}%
  \else\usebox{\pandoc@box}%
  \fi%
}
% Set default figure placement to htbp
\def\fps@figure{htbp}
\makeatother





\setlength{\emergencystretch}{3em} % prevent overfull lines

\providecommand{\tightlist}{%
  \setlength{\itemsep}{0pt}\setlength{\parskip}{0pt}}



 


\usepackage{booktabs}
\usepackage{caption}
\usepackage{longtable}
\usepackage{colortbl}
\usepackage{array}
\usepackage{anyfontsize}
\usepackage{multirow}
\KOMAoption{captions}{tableheading}
\makeatletter
\@ifpackageloaded{tcolorbox}{}{\usepackage[skins,breakable]{tcolorbox}}
\@ifpackageloaded{fontawesome5}{}{\usepackage{fontawesome5}}
\definecolor{quarto-callout-color}{HTML}{909090}
\definecolor{quarto-callout-note-color}{HTML}{0758E5}
\definecolor{quarto-callout-important-color}{HTML}{CC1914}
\definecolor{quarto-callout-warning-color}{HTML}{EB9113}
\definecolor{quarto-callout-tip-color}{HTML}{00A047}
\definecolor{quarto-callout-caution-color}{HTML}{FC5300}
\definecolor{quarto-callout-color-frame}{HTML}{acacac}
\definecolor{quarto-callout-note-color-frame}{HTML}{4582ec}
\definecolor{quarto-callout-important-color-frame}{HTML}{d9534f}
\definecolor{quarto-callout-warning-color-frame}{HTML}{f0ad4e}
\definecolor{quarto-callout-tip-color-frame}{HTML}{02b875}
\definecolor{quarto-callout-caution-color-frame}{HTML}{fd7e14}
\makeatother
\makeatletter
\@ifpackageloaded{bookmark}{}{\usepackage{bookmark}}
\makeatother
\makeatletter
\@ifpackageloaded{caption}{}{\usepackage{caption}}
\AtBeginDocument{%
\ifdefined\contentsname
  \renewcommand*\contentsname{Table of contents}
\else
  \newcommand\contentsname{Table of contents}
\fi
\ifdefined\listfigurename
  \renewcommand*\listfigurename{List of Figures}
\else
  \newcommand\listfigurename{List of Figures}
\fi
\ifdefined\listtablename
  \renewcommand*\listtablename{List of Tables}
\else
  \newcommand\listtablename{List of Tables}
\fi
\ifdefined\figurename
  \renewcommand*\figurename{Figure}
\else
  \newcommand\figurename{Figure}
\fi
\ifdefined\tablename
  \renewcommand*\tablename{Table}
\else
  \newcommand\tablename{Table}
\fi
}
\@ifpackageloaded{float}{}{\usepackage{float}}
\floatstyle{ruled}
\@ifundefined{c@chapter}{\newfloat{codelisting}{h}{lop}}{\newfloat{codelisting}{h}{lop}[chapter]}
\floatname{codelisting}{Listing}
\newcommand*\listoflistings{\listof{codelisting}{List of Listings}}
\makeatother
\makeatletter
\makeatother
\makeatletter
\@ifpackageloaded{caption}{}{\usepackage{caption}}
\@ifpackageloaded{subcaption}{}{\usepackage{subcaption}}
\makeatother
\usepackage{bookmark}
\IfFileExists{xurl.sty}{\usepackage{xurl}}{} % add URL line breaks if available
\urlstyle{same}
\hypersetup{
  pdftitle={Elements of Computational Techniques in Statistics},
  pdfauthor={Longhai Li},
  colorlinks=true,
  linkcolor={blue},
  filecolor={Maroon},
  citecolor={Blue},
  urlcolor={Blue},
  pdfcreator={LaTeX via pandoc}}


\title{Elements of Computational Techniques in Statistics}
\author{Longhai Li}
\date{2026-09-13}
\begin{document}
\maketitle

\renewcommand*\contentsname{Table of contents}
{
\setcounter{tocdepth}{2}
\tableofcontents
}

\bookmarksetup{startatroot}

\chapter*{Preface}\label{preface}
\addcontentsline{toc}{chapter}{Preface}

\markboth{Preface}{Preface}

This is a concise course about statistical inference, which was
developed for the course
\href{https://longhaisk.github.io/teaching/stat812/}{STAT 420/812} at
University of Saskatchewan.

\bookmarksetup{startatroot}

\chapter{Introduction to R and
Quarto}\label{introduction-to-r-and-quarto}

\section{A toy example}\label{a-toy-example}

Throughout this unit we will use one running example to introduce the
basic building blocks of R: vectors, strings, special values, matrices,
data frames, lists, functions, file I/O, and graphics.

\textbf{The mission:} a class of 100 students has taken graded work in
three \emph{categories}: five quizzes, two tests, and one final exam,
each scored out of 100 points. A student's \textbf{course score} is a
weighted average of the category means -- quizzes count 30\%, tests
count 30\%, and the final counts 40\%. To make the simulated data
realistic, a student's test and final scores are generated from their
quiz performance by a simple linear regression, with a different
intercept for Program A and Program B students, and a few quizzes are
missing entirely (a student who did not sit that quiz). We want to

\begin{enumerate}
\def\labelenumi{\arabic{enumi}.}
\tightlist
\item
  work out the calculation step by step for one student,
\item
  package that calculation into a reusable function,
\item
  apply the function to every student to get a full set of results, and
\item
  present those results as a table and as plots.
\end{enumerate}

This mirrors how real data-analysis code is usually built: explore on
one case, generalize into a function, then apply and report.

\section{Generating an artificial score
dataset}\label{generating-an-artificial-score-dataset}

\subsection{Vectors: recording quiz
scores}\label{vectors-recording-quiz-scores}

Each quiz, across all students, is naturally stored as a vector. We
simulate a class of 100 students, five quizzes scored out of 100, using
\texttt{rnorm} and clipping the result to a valid
\texttt{0}-\texttt{100} range with \texttt{pmin}/\texttt{pmax}:

\begin{Shaded}
\begin{Highlighting}[]
\FunctionTok{set.seed}\NormalTok{(}\DecValTok{123}\NormalTok{)}
\NormalTok{n }\OtherTok{\textless{}{-}} \DecValTok{100}

\NormalTok{gen.quiz }\OtherTok{\textless{}{-}} \ControlFlowTok{function}\NormalTok{(mean, sd)}
\NormalTok{\{}
\NormalTok{    x }\OtherTok{\textless{}{-}} \FunctionTok{round}\NormalTok{(}\FunctionTok{rnorm}\NormalTok{(n, mean, sd))}
    \FunctionTok{pmin}\NormalTok{(}\DecValTok{100}\NormalTok{, }\FunctionTok{pmax}\NormalTok{(}\DecValTok{0}\NormalTok{, x))}
\NormalTok{\}}

\NormalTok{Quiz1 }\OtherTok{\textless{}{-}} \FunctionTok{gen.quiz}\NormalTok{(}\DecValTok{75}\NormalTok{, }\DecValTok{20}\NormalTok{)}
\NormalTok{Quiz2 }\OtherTok{\textless{}{-}} \FunctionTok{gen.quiz}\NormalTok{(}\DecValTok{78}\NormalTok{, }\DecValTok{20}\NormalTok{)}
\NormalTok{Quiz3 }\OtherTok{\textless{}{-}} \FunctionTok{gen.quiz}\NormalTok{(}\DecValTok{72}\NormalTok{, }\DecValTok{20}\NormalTok{)}
\NormalTok{Quiz4 }\OtherTok{\textless{}{-}} \FunctionTok{gen.quiz}\NormalTok{(}\DecValTok{80}\NormalTok{, }\DecValTok{20}\NormalTok{)}
\NormalTok{Quiz5 }\OtherTok{\textless{}{-}} \FunctionTok{gen.quiz}\NormalTok{(}\DecValTok{74}\NormalTok{, }\DecValTok{20}\NormalTok{)}

\NormalTok{Quiz1}
\end{Highlighting}
\end{Shaded}

\begin{verbatim}
  [1]  64  70 100  76  78 100  84  50  61  66  99  82  83  77  64 100  85  36
 [19]  89  66  54  71  54  60  62  41  92  78  52 100  84  69  93  93  91  89
 [37]  86  74  69  67  61  71  50 100  99  53  67  66  91  73  80  74  74 100
 [55]  70 100  44  87  77  79  83  65  68  55  54  81  84  76  93 100  65  29
 [73]  95  61  61  96  69  51  79  72  75  83  68  88  71  82  97  84  68  98
 [91]  95  86  80  62 100  63 100 100  70  54
\end{verbatim}

\subsection{Strings: naming the categories and
students}\label{strings-naming-the-categories-and-students}

We could type out \texttt{"Quiz1"}, \texttt{"Quiz2"}, \ldots{} by hand,
but \texttt{paste0} builds such labels programmatically -- a habit that
pays off the moment there are 40 quizzes instead of 5:

\begin{Shaded}
\begin{Highlighting}[]
\NormalTok{quiz.names }\OtherTok{\textless{}{-}} \FunctionTok{paste0}\NormalTok{(}\StringTok{"Quiz"}\NormalTok{, }\DecValTok{1}\SpecialCharTok{:}\DecValTok{5}\NormalTok{)}
\NormalTok{test.names }\OtherTok{\textless{}{-}} \FunctionTok{paste0}\NormalTok{(}\StringTok{"Test"}\NormalTok{, }\DecValTok{1}\SpecialCharTok{:}\DecValTok{2}\NormalTok{)}

\NormalTok{quiz.names}
\end{Highlighting}
\end{Shaded}

\begin{verbatim}
[1] "Quiz1" "Quiz2" "Quiz3" "Quiz4" "Quiz5"
\end{verbatim}

\begin{Shaded}
\begin{Highlighting}[]
\NormalTok{test.names}
\end{Highlighting}
\end{Shaded}

\begin{verbatim}
[1] "Test1" "Test2"
\end{verbatim}

With 100 students, typing out individual names is not realistic either.
\texttt{sprintf} generates zero-padded IDs like \texttt{"S001"},
\texttt{"S002"}, \ldots{} in one line -- exactly the kind of job it is
built for:

\begin{Shaded}
\begin{Highlighting}[]
\NormalTok{students }\OtherTok{\textless{}{-}} \FunctionTok{sprintf}\NormalTok{(}\StringTok{"S\%03d"}\NormalTok{, }\DecValTok{1}\SpecialCharTok{:}\NormalTok{n)}

\FunctionTok{head}\NormalTok{(students)}
\end{Highlighting}
\end{Shaded}

\begin{verbatim}
[1] "S001" "S002" "S003" "S004" "S005" "S006"
\end{verbatim}

\subsection{Simulating the tests and the final using vectorized
math}\label{simulating-the-tests-and-the-final-using-vectorized-math}

In a real class, a student's test and final scores are not independent
of how they did on the quizzes. We simulate that by treating each
student's \textbf{mean quiz score} as a predictor in a simple linear
regression.

This is a great place to demonstrate \textbf{vectorized calculations} in
R. Instead of writing a \texttt{for} loop to calculate each student's
score one by one, we can simply add and multiply whole vectors of length
100 together. R aligns the elements automatically:

\begin{Shaded}
\begin{Highlighting}[]
\NormalTok{program }\OtherTok{\textless{}{-}} \FunctionTok{factor}\NormalTok{(}\FunctionTok{sample}\NormalTok{(}\FunctionTok{c}\NormalTok{(}\StringTok{"Program A"}\NormalTok{, }\StringTok{"Program B"}\NormalTok{), n, }\AttributeTok{replace =} \ConstantTok{TRUE}\NormalTok{))}

\NormalTok{quiz.mat }\OtherTok{\textless{}{-}} \FunctionTok{cbind}\NormalTok{(Quiz1, Quiz2, Quiz3, Quiz4, Quiz5)}

\FunctionTok{head}\NormalTok{(quiz.mat)}
\end{Highlighting}
\end{Shaded}

\begin{verbatim}
     Quiz1 Quiz2 Quiz3 Quiz4 Quiz5
[1,]    64    64   100    66    73
[2,]    70    83    98    65    51
[3,]   100    73    67    61    61
[4,]    76    71    83    59    73
[5,]    78    59    64    71    87
[6,]   100    77    62    87    41
\end{verbatim}

\begin{Shaded}
\begin{Highlighting}[]
\NormalTok{quiz.mean.true }\OtherTok{\textless{}{-}} \FunctionTok{rowMeans}\NormalTok{(quiz.mat)}
\FunctionTok{data.frame}\NormalTok{(program, quiz.mean.true)}
\end{Highlighting}
\end{Shaded}

\begin{verbatim}
      program quiz.mean.true
1   Program A           73.4
2   Program B           73.4
3   Program B           72.4
4   Program A           72.4
5   Program B           71.8
6   Program A           73.4
7   Program A           61.8
8   Program B           65.6
9   Program B           78.8
10  Program A           82.4
11  Program A           84.6
12  Program B           84.6
13  Program B           71.6
14  Program A           71.2
15  Program B           68.6
16  Program A           79.8
17  Program A           75.6
18  Program A           59.4
19  Program B           76.4
20  Program B           66.2
21  Program A           76.0
22  Program A           78.8
23  Program B           69.2
24  Program A           82.2
25  Program B           71.4
26  Program B           67.6
27  Program B           82.4
28  Program A           76.6
29  Program A           80.2
30  Program A           70.8
31  Program B           90.8
32  Program B           72.8
33  Program B           78.8
34  Program A           74.8
35  Program B           70.2
36  Program B           65.2
37  Program A           68.8
38  Program A           83.8
39  Program A           74.8
40  Program A           69.8
41  Program A           70.0
42  Program A           77.6
43  Program A           76.8
44  Program B           69.6
45  Program B           70.6
46  Program B           68.4
47  Program B           70.2
48  Program B           72.8
49  Program A           68.8
50  Program A           64.8
51  Program B           87.6
52  Program B           72.6
53  Program B           75.6
54  Program A           80.4
55  Program B           80.0
56  Program A           70.0
57  Program B           82.2
58  Program B           77.2
59  Program B           71.0
60  Program B           71.8
61  Program B           82.2
62  Program A           68.8
63  Program A           72.6
64  Program A           83.8
65  Program B           74.4
66  Program B           87.8
67  Program B           84.2
68  Program A           68.4
69  Program A           83.4
70  Program A           85.4
71  Program B           85.6
72  Program A           62.8
73  Program B           75.4
74  Program A           74.2
75  Program A           74.6
76  Program B           72.0
77  Program B           76.6
78  Program A           71.6
79  Program A           81.6
80  Program B           65.2
81  Program B           65.2
82  Program A           86.8
83  Program A           80.4
84  Program A           70.2
85  Program B           79.2
86  Program A           84.2
87  Program A           90.8
88  Program A           76.8
89  Program A           80.0
90  Program B           78.8
91  Program A           82.2
92  Program B           77.6
93  Program A           90.6
94  Program B           59.4
95  Program A           75.6
96  Program B           76.6
97  Program B           86.6
98  Program A           74.4
99  Program B           68.8
100 Program B           72.6
\end{verbatim}

The intercept depends on \texttt{program}, so we compute it with
\texttt{ifelse}; the slope is the same for every student:

\begin{Shaded}
\begin{Highlighting}[]
\NormalTok{b0.test }\OtherTok{\textless{}{-}} \FunctionTok{ifelse}\NormalTok{(program }\SpecialCharTok{==} \StringTok{"Program A"}\NormalTok{, }\DecValTok{12}\NormalTok{, }\DecValTok{6}\NormalTok{)}
\NormalTok{b1.test }\OtherTok{\textless{}{-}} \FloatTok{0.80}

\FunctionTok{data.frame}\NormalTok{(program, b0.test, b1.test)}
\end{Highlighting}
\end{Shaded}

\begin{verbatim}
      program b0.test b1.test
1   Program A      12     0.8
2   Program B       6     0.8
3   Program B       6     0.8
4   Program A      12     0.8
5   Program B       6     0.8
6   Program A      12     0.8
7   Program A      12     0.8
8   Program B       6     0.8
9   Program B       6     0.8
10  Program A      12     0.8
11  Program A      12     0.8
12  Program B       6     0.8
13  Program B       6     0.8
14  Program A      12     0.8
15  Program B       6     0.8
16  Program A      12     0.8
17  Program A      12     0.8
18  Program A      12     0.8
19  Program B       6     0.8
20  Program B       6     0.8
21  Program A      12     0.8
22  Program A      12     0.8
23  Program B       6     0.8
24  Program A      12     0.8
25  Program B       6     0.8
26  Program B       6     0.8
27  Program B       6     0.8
28  Program A      12     0.8
29  Program A      12     0.8
30  Program A      12     0.8
31  Program B       6     0.8
32  Program B       6     0.8
33  Program B       6     0.8
34  Program A      12     0.8
35  Program B       6     0.8
36  Program B       6     0.8
37  Program A      12     0.8
38  Program A      12     0.8
39  Program A      12     0.8
40  Program A      12     0.8
41  Program A      12     0.8
42  Program A      12     0.8
43  Program A      12     0.8
44  Program B       6     0.8
45  Program B       6     0.8
46  Program B       6     0.8
47  Program B       6     0.8
48  Program B       6     0.8
49  Program A      12     0.8
50  Program A      12     0.8
51  Program B       6     0.8
52  Program B       6     0.8
53  Program B       6     0.8
54  Program A      12     0.8
55  Program B       6     0.8
56  Program A      12     0.8
57  Program B       6     0.8
58  Program B       6     0.8
59  Program B       6     0.8
60  Program B       6     0.8
61  Program B       6     0.8
62  Program A      12     0.8
63  Program A      12     0.8
64  Program A      12     0.8
65  Program B       6     0.8
66  Program B       6     0.8
67  Program B       6     0.8
68  Program A      12     0.8
69  Program A      12     0.8
70  Program A      12     0.8
71  Program B       6     0.8
72  Program A      12     0.8
73  Program B       6     0.8
74  Program A      12     0.8
75  Program A      12     0.8
76  Program B       6     0.8
77  Program B       6     0.8
78  Program A      12     0.8
79  Program A      12     0.8
80  Program B       6     0.8
81  Program B       6     0.8
82  Program A      12     0.8
83  Program A      12     0.8
84  Program A      12     0.8
85  Program B       6     0.8
86  Program A      12     0.8
87  Program A      12     0.8
88  Program A      12     0.8
89  Program A      12     0.8
90  Program B       6     0.8
91  Program A      12     0.8
92  Program B       6     0.8
93  Program A      12     0.8
94  Program B       6     0.8
95  Program A      12     0.8
96  Program B       6     0.8
97  Program B       6     0.8
98  Program A      12     0.8
99  Program B       6     0.8
100 Program B       6     0.8
\end{verbatim}

The final exam uses its own intercept and slope, computed the same way:

\begin{Shaded}
\begin{Highlighting}[]
\NormalTok{b0.final }\OtherTok{\textless{}{-}} \FunctionTok{ifelse}\NormalTok{(program }\SpecialCharTok{==} \StringTok{"Program A"}\NormalTok{, }\DecValTok{15}\NormalTok{, }\DecValTok{8}\NormalTok{)}
\NormalTok{b1.final }\OtherTok{\textless{}{-}} \FloatTok{0.85}

\FunctionTok{data.frame}\NormalTok{(program, b0.final, b1.final)}
\end{Highlighting}
\end{Shaded}

\begin{verbatim}
      program b0.final b1.final
1   Program A       15     0.85
2   Program B        8     0.85
3   Program B        8     0.85
4   Program A       15     0.85
5   Program B        8     0.85
6   Program A       15     0.85
7   Program A       15     0.85
8   Program B        8     0.85
9   Program B        8     0.85
10  Program A       15     0.85
11  Program A       15     0.85
12  Program B        8     0.85
13  Program B        8     0.85
14  Program A       15     0.85
15  Program B        8     0.85
16  Program A       15     0.85
17  Program A       15     0.85
18  Program A       15     0.85
19  Program B        8     0.85
20  Program B        8     0.85
21  Program A       15     0.85
22  Program A       15     0.85
23  Program B        8     0.85
24  Program A       15     0.85
25  Program B        8     0.85
26  Program B        8     0.85
27  Program B        8     0.85
28  Program A       15     0.85
29  Program A       15     0.85
30  Program A       15     0.85
31  Program B        8     0.85
32  Program B        8     0.85
33  Program B        8     0.85
34  Program A       15     0.85
35  Program B        8     0.85
36  Program B        8     0.85
37  Program A       15     0.85
38  Program A       15     0.85
39  Program A       15     0.85
40  Program A       15     0.85
41  Program A       15     0.85
42  Program A       15     0.85
43  Program A       15     0.85
44  Program B        8     0.85
45  Program B        8     0.85
46  Program B        8     0.85
47  Program B        8     0.85
48  Program B        8     0.85
49  Program A       15     0.85
50  Program A       15     0.85
51  Program B        8     0.85
52  Program B        8     0.85
53  Program B        8     0.85
54  Program A       15     0.85
55  Program B        8     0.85
56  Program A       15     0.85
57  Program B        8     0.85
58  Program B        8     0.85
59  Program B        8     0.85
60  Program B        8     0.85
61  Program B        8     0.85
62  Program A       15     0.85
63  Program A       15     0.85
64  Program A       15     0.85
65  Program B        8     0.85
66  Program B        8     0.85
67  Program B        8     0.85
68  Program A       15     0.85
69  Program A       15     0.85
70  Program A       15     0.85
71  Program B        8     0.85
72  Program A       15     0.85
73  Program B        8     0.85
74  Program A       15     0.85
75  Program A       15     0.85
76  Program B        8     0.85
77  Program B        8     0.85
78  Program A       15     0.85
79  Program A       15     0.85
80  Program B        8     0.85
81  Program B        8     0.85
82  Program A       15     0.85
83  Program A       15     0.85
84  Program A       15     0.85
85  Program B        8     0.85
86  Program A       15     0.85
87  Program A       15     0.85
88  Program A       15     0.85
89  Program A       15     0.85
90  Program B        8     0.85
91  Program A       15     0.85
92  Program B        8     0.85
93  Program A       15     0.85
94  Program B        8     0.85
95  Program A       15     0.85
96  Program B        8     0.85
97  Program B        8     0.85
98  Program A       15     0.85
99  Program B        8     0.85
100 Program B        8     0.85
\end{verbatim}

\texttt{clip.round} rounds a score and clips it to the valid
\texttt{0}-\texttt{100} range -- we'll reuse it for every test and the
final:

\begin{Shaded}
\begin{Highlighting}[]
\NormalTok{clip.round }\OtherTok{\textless{}{-}} \ControlFlowTok{function}\NormalTok{(x) }\FunctionTok{round}\NormalTok{(}\FunctionTok{pmin}\NormalTok{(}\DecValTok{100}\NormalTok{, }\FunctionTok{pmax}\NormalTok{(}\DecValTok{0}\NormalTok{, x)))}
\end{Highlighting}
\end{Shaded}

Now we add random error to the linear predictor to get each student's
raw \texttt{Test1} score, then clip and round it:

\begin{Shaded}
\begin{Highlighting}[]
\NormalTok{error }\OtherTok{\textless{}{-}} \FunctionTok{rnorm}\NormalTok{(n, }\DecValTok{0}\NormalTok{, }\DecValTok{8}\NormalTok{)}
\NormalTok{Test1.raw }\OtherTok{\textless{}{-}}\NormalTok{ b0.test }\SpecialCharTok{+}\NormalTok{ b1.test}\SpecialCharTok{*}\NormalTok{quiz.mean.true }\SpecialCharTok{+}\NormalTok{ error}

\FunctionTok{data.frame}\NormalTok{(program, b0.test, b1.test, quiz.mean.true, error, Test1.raw)}
\end{Highlighting}
\end{Shaded}

\begin{verbatim}
      program b0.test b1.test quiz.mean.true       error Test1.raw
1   Program A      12     0.8           73.4   6.7592339  77.47923
2   Program B       6     0.8           73.4  -6.2576148  58.46239
3   Program B       6     0.8           72.4   8.8856913  72.80569
4   Program A      12     0.8           72.4   1.9985978  71.91860
5   Program B       6     0.8           71.8  13.2153231  76.65532
6   Program A      12     0.8           73.4 -11.6717658  59.04823
7   Program A      12     0.8           61.8  -0.4103831  61.02962
8   Program B       6     0.8           65.6  -4.2154014  54.26460
9   Program B       6     0.8           78.8  -1.5781190  67.46188
10  Program A      12     0.8           82.4  -5.0366299  72.88337
11  Program A      12     0.8           84.6  -6.6707486  73.00925
12  Program B       6     0.8           84.6   4.6297790  78.30978
13  Program B       6     0.8           71.6  -8.7006457  54.57935
14  Program A      12     0.8           71.2  11.8722475  80.83225
15  Program B       6     0.8           68.6  -9.4896527  51.39035
16  Program A      12     0.8           79.8   0.8086332  76.64863
17  Program A      12     0.8           75.6   4.2639143  76.74391
18  Program A      12     0.8           59.4   4.6938827  64.21388
19  Program B       6     0.8           76.4  -2.4139733  64.70603
20  Program B       6     0.8           66.2   0.6360160  59.59602
21  Program A      12     0.8           76.0   7.6901132  80.49011
22  Program A      12     0.8           78.8 -11.6517273  63.38827
23  Program B       6     0.8           69.2  -6.2539177  55.10608
24  Program A      12     0.8           82.2   2.5632185  80.32322
25  Program B       6     0.8           71.4  -3.5582558  59.56174
26  Program B       6     0.8           67.6  10.9600320  71.04003
27  Program B       6     0.8           82.4   5.3860309  77.30603
28  Program A      12     0.8           76.6   0.5773340  73.85733
29  Program A      12     0.8           80.2 -12.0620585  64.09794
30  Program A      12     0.8           70.8   0.2088018  68.84880
31  Program B       6     0.8           90.8  -2.5313269  76.10867
32  Program B       6     0.8           72.8  -0.8187721  63.42123
33  Program B       6     0.8           78.8  -9.4524738  59.58753
34  Program A      12     0.8           74.8   3.9892644  75.82926
35  Program B       6     0.8           70.2  -8.3116515  53.84835
36  Program B       6     0.8           65.2  -1.8097759  56.35022
37  Program A      12     0.8           68.8   3.0514066  70.09141
38  Program A      12     0.8           83.8  -6.2681263  72.77187
39  Program A      12     0.8           74.8   4.6639312  76.50393
40  Program A      12     0.8           69.8 -10.5320832  57.30792
41  Program A      12     0.8           70.0 -22.4781974  45.52180
42  Program A      12     0.8           77.6   3.7197439  77.79974
43  Program A      12     0.8           76.8   6.7243186  80.16432
44  Program B       6     0.8           69.6  -2.2867634  59.39324
45  Program B       6     0.8           70.6   4.0330100  66.51301
46  Program B       6     0.8           68.4  -9.2473322  51.47267
47  Program B       6     0.8           70.2  -1.0171889  61.14281
48  Program B       6     0.8           72.8 -15.5321470  48.70785
49  Program A      12     0.8           68.8   9.4494471  76.48945
50  Program A      12     0.8           64.8  14.8792869  78.71929
51  Program B       6     0.8           87.6   8.5920981  84.67210
52  Program B       6     0.8           72.6  -0.2187757  63.86122
53  Program B       6     0.8           75.6  -0.2666427  66.21336
54  Program A      12     0.8           80.4 -12.1285410  64.19146
55  Program B       6     0.8           80.0   6.3230828  76.32308
56  Program A      12     0.8           70.0  -1.6858734  66.31413
57  Program B       6     0.8           82.2  -5.2539434  66.50606
58  Program B       6     0.8           77.2 -11.2962063  56.46379
59  Program B       6     0.8           71.0  -2.3981000  60.40190
60  Program B       6     0.8           71.8  -6.7924891  56.64751
61  Program B       6     0.8           82.2  -3.1762442  68.58376
62  Program A      12     0.8           68.8  -9.7407999  57.29920
63  Program A      12     0.8           72.6  13.5007159  83.58072
64  Program A      12     0.8           83.8  -0.1280202  78.91198
65  Program B       6     0.8           74.4   8.5995606  74.11956
66  Program B       6     0.8           87.8 -20.8135974  55.42640
67  Program B       6     0.8           84.2  -3.6255826  69.73442
68  Program A      12     0.8           68.4  -5.4038583  61.31614
69  Program A      12     0.8           83.4  -9.7834094  68.93659
70  Program A      12     0.8           85.4  12.3728732  92.69287
71  Program B       6     0.8           85.6 -11.3222554  63.15774
72  Program A      12     0.8           62.8   2.5471221  64.78712
73  Program B       6     0.8           75.4   6.7714903  73.09149
74  Program A      12     0.8           74.2   1.4255215  72.78552
75  Program A      12     0.8           74.6  -7.0020438  64.67796
76  Program B       6     0.8           72.0   7.5293265  71.12933
77  Program B       6     0.8           76.6   1.3647047  68.64470
78  Program A      12     0.8           71.6  -8.5079833  60.77202
79  Program A      12     0.8           81.6 -11.1043924  66.17561
80  Program B       6     0.8           65.2  16.6937395  74.85374
81  Program B       6     0.8           65.2  -5.4280252  52.73197
82  Program A      12     0.8           86.8 -14.8445732  66.59543
83  Program A      12     0.8           80.4   4.2660749  80.58607
84  Program A      12     0.8           70.2   2.4818420  70.64184
85  Program B       6     0.8           79.2 -10.8306747  58.52933
86  Program A      12     0.8           84.2 -15.5436513  63.81635
87  Program A      12     0.8           90.8  -0.9304201  83.70958
88  Program A      12     0.8           76.8   9.1151703  82.55517
89  Program A      12     0.8           80.0   5.0889923  81.08899
90  Program B       6     0.8           78.8  -3.9434994  65.09650
91  Program A      12     0.8           82.2  -6.6735059  71.08649
92  Program B       6     0.8           77.6   2.1685341  70.24853
93  Program A      12     0.8           90.6   1.2588268  85.73883
94  Program B       6     0.8           59.4   5.0376940  58.55769
95  Program A      12     0.8           75.6  -3.1663836  69.31362
96  Program B       6     0.8           76.6   7.1948324  74.47483
97  Program B       6     0.8           86.6  -6.6464922  68.63351
98  Program A      12     0.8           74.4  -2.6443576  68.87564
99  Program B       6     0.8           68.8   5.9265162  66.96652
100 Program B       6     0.8           72.6   7.9197729  71.99977
\end{verbatim}

\begin{Shaded}
\begin{Highlighting}[]
\NormalTok{Test1 }\OtherTok{\textless{}{-}} \FunctionTok{clip.round}\NormalTok{(Test1.raw)}

\FunctionTok{data.frame}\NormalTok{(program, Test1.raw, Test1)}
\end{Highlighting}
\end{Shaded}

\begin{verbatim}
      program Test1.raw Test1
1   Program A  77.47923    77
2   Program B  58.46239    58
3   Program B  72.80569    73
4   Program A  71.91860    72
5   Program B  76.65532    77
6   Program A  59.04823    59
7   Program A  61.02962    61
8   Program B  54.26460    54
9   Program B  67.46188    67
10  Program A  72.88337    73
11  Program A  73.00925    73
12  Program B  78.30978    78
13  Program B  54.57935    55
14  Program A  80.83225    81
15  Program B  51.39035    51
16  Program A  76.64863    77
17  Program A  76.74391    77
18  Program A  64.21388    64
19  Program B  64.70603    65
20  Program B  59.59602    60
21  Program A  80.49011    80
22  Program A  63.38827    63
23  Program B  55.10608    55
24  Program A  80.32322    80
25  Program B  59.56174    60
26  Program B  71.04003    71
27  Program B  77.30603    77
28  Program A  73.85733    74
29  Program A  64.09794    64
30  Program A  68.84880    69
31  Program B  76.10867    76
32  Program B  63.42123    63
33  Program B  59.58753    60
34  Program A  75.82926    76
35  Program B  53.84835    54
36  Program B  56.35022    56
37  Program A  70.09141    70
38  Program A  72.77187    73
39  Program A  76.50393    77
40  Program A  57.30792    57
41  Program A  45.52180    46
42  Program A  77.79974    78
43  Program A  80.16432    80
44  Program B  59.39324    59
45  Program B  66.51301    67
46  Program B  51.47267    51
47  Program B  61.14281    61
48  Program B  48.70785    49
49  Program A  76.48945    76
50  Program A  78.71929    79
51  Program B  84.67210    85
52  Program B  63.86122    64
53  Program B  66.21336    66
54  Program A  64.19146    64
55  Program B  76.32308    76
56  Program A  66.31413    66
57  Program B  66.50606    67
58  Program B  56.46379    56
59  Program B  60.40190    60
60  Program B  56.64751    57
61  Program B  68.58376    69
62  Program A  57.29920    57
63  Program A  83.58072    84
64  Program A  78.91198    79
65  Program B  74.11956    74
66  Program B  55.42640    55
67  Program B  69.73442    70
68  Program A  61.31614    61
69  Program A  68.93659    69
70  Program A  92.69287    93
71  Program B  63.15774    63
72  Program A  64.78712    65
73  Program B  73.09149    73
74  Program A  72.78552    73
75  Program A  64.67796    65
76  Program B  71.12933    71
77  Program B  68.64470    69
78  Program A  60.77202    61
79  Program A  66.17561    66
80  Program B  74.85374    75
81  Program B  52.73197    53
82  Program A  66.59543    67
83  Program A  80.58607    81
84  Program A  70.64184    71
85  Program B  58.52933    59
86  Program A  63.81635    64
87  Program A  83.70958    84
88  Program A  82.55517    83
89  Program A  81.08899    81
90  Program B  65.09650    65
91  Program A  71.08649    71
92  Program B  70.24853    70
93  Program A  85.73883    86
94  Program B  58.55769    59
95  Program A  69.31362    69
96  Program B  74.47483    74
97  Program B  68.63351    69
98  Program A  68.87564    69
99  Program B  66.96652    67
100 Program B  71.99977    72
\end{verbatim}

Using the above method, we generate two tests and a final:

\begin{Shaded}
\begin{Highlighting}[]
\CommentTok{\# Vectorized math: b0.test + b1.test * quiz.mean.true + rnorm(...)}
\NormalTok{Test1 }\OtherTok{\textless{}{-}} \FunctionTok{clip.round}\NormalTok{(b0.test  }\SpecialCharTok{+}\NormalTok{ b1.test}\SpecialCharTok{*}\NormalTok{quiz.mean.true  }\SpecialCharTok{+} \FunctionTok{rnorm}\NormalTok{(n, }\DecValTok{0}\NormalTok{, }\DecValTok{8}\NormalTok{))}
\NormalTok{Test2 }\OtherTok{\textless{}{-}} \FunctionTok{clip.round}\NormalTok{(b0.test  }\SpecialCharTok{+}\NormalTok{ b1.test}\SpecialCharTok{*}\NormalTok{quiz.mean.true  }\SpecialCharTok{+} \FunctionTok{rnorm}\NormalTok{(n, }\DecValTok{0}\NormalTok{, }\DecValTok{8}\NormalTok{))}
\NormalTok{Final }\OtherTok{\textless{}{-}} \FunctionTok{clip.round}\NormalTok{(b0.final }\SpecialCharTok{+}\NormalTok{ b1.final}\SpecialCharTok{*}\NormalTok{quiz.mean.true }\SpecialCharTok{+} \FunctionTok{rnorm}\NormalTok{(n, }\DecValTok{0}\NormalTok{, }\DecValTok{6}\NormalTok{))}
\end{Highlighting}
\end{Shaded}

\subsection{Matrices: combining all the scores (Complete
Data)}\label{matrices-combining-all-the-scores-complete-data}

Before we simulate absences (missing values), let's bind the quiz, test,
and final columns together into one complete matrix. It will have one
row per student and one column per graded item.

\begin{Shaded}
\begin{Highlighting}[]
\NormalTok{scores }\OtherTok{\textless{}{-}} \FunctionTok{cbind}\NormalTok{(quiz.mat, Test1, Test2, Final)}
\FunctionTok{colnames}\NormalTok{(scores) }\OtherTok{\textless{}{-}} \FunctionTok{c}\NormalTok{(quiz.names, test.names, }\StringTok{"Final"}\NormalTok{)}
\FunctionTok{rownames}\NormalTok{(scores) }\OtherTok{\textless{}{-}}\NormalTok{ students}

\CommentTok{\# Here is our complete dataset before any missing values}
\FunctionTok{head}\NormalTok{(scores)}
\end{Highlighting}
\end{Shaded}

\begin{verbatim}
     Quiz1 Quiz2 Quiz3 Quiz4 Quiz5 Test1 Test2 Final
S001    64    64   100    66    73    55    83    81
S002    70    83    98    65    51    66    64    70
S003   100    73    67    61    61    69    68    76
S004    76    71    83    59    73    58    72    70
S005    78    59    64    71    87    67    62    83
S006   100    77    62    87    41    64    70    74
\end{verbatim}

\subsection{Adding and handling missing
values}\label{adding-and-handling-missing-values}

In reality, students occasionally miss a quiz. Let's randomly drop about
5\% of the quiz scores by assigning them the special value \texttt{NA}
(Not Available).

\begin{Shaded}
\begin{Highlighting}[]
\NormalTok{miss.rate }\OtherTok{\textless{}{-}} \FloatTok{0.05}
\FunctionTok{sample}\NormalTok{(n, }\FunctionTok{round}\NormalTok{(n}\SpecialCharTok{*}\NormalTok{miss.rate))}
\end{Highlighting}
\end{Shaded}

\begin{verbatim}
[1] 31 35 73 49 38
\end{verbatim}

\begin{Shaded}
\begin{Highlighting}[]
\CommentTok{\# We only want to introduce NAs into the first 5 columns (the quizzes)}
\ControlFlowTok{for}\NormalTok{(j }\ControlFlowTok{in} \DecValTok{1}\SpecialCharTok{:}\FunctionTok{ncol}\NormalTok{(scores)) \{}
\NormalTok{    scores[}\FunctionTok{sample}\NormalTok{(n, }\FunctionTok{round}\NormalTok{(n}\SpecialCharTok{*}\NormalTok{miss.rate)), j] }\OtherTok{\textless{}{-}} \ConstantTok{NA}
\NormalTok{\}}

\FunctionTok{head}\NormalTok{(scores, }\DecValTok{10}\NormalTok{)}
\end{Highlighting}
\end{Shaded}

\begin{verbatim}
     Quiz1 Quiz2 Quiz3 Quiz4 Quiz5 Test1 Test2 Final
S001    64    64   100    66    73    55    83    81
S002    70    83    98    65    51    66    64    70
S003   100    73    67    61    61    69    68    76
S004    76    71    83    59    73    58    NA    70
S005    78    59    64    NA    87    67    62    83
S006   100    77    62    87    41    64    70    74
S007    84    62    56    40    67    70    70    59
S008    50    45    NA    84    89    63    57    62
S009    61    70   100    NA    63    75    53    76
S010    66    96    71   100    79    79    76    95
\end{verbatim}

\texttt{is.na} locates missing values anywhere in an object. Only the
quiz columns have any, since we generated \texttt{Test1},
\texttt{Test2}, and \texttt{Final} without missingness:

\begin{Shaded}
\begin{Highlighting}[]
\FunctionTok{colSums}\NormalTok{(}\FunctionTok{is.na}\NormalTok{(scores))}
\end{Highlighting}
\end{Shaded}

\begin{verbatim}
Quiz1 Quiz2 Quiz3 Quiz4 Quiz5 Test1 Test2 Final 
    5     5     5     5     5     5     5     5 
\end{verbatim}

We will rely on \texttt{na.rm\ =\ TRUE} later so that a missing quiz
simply drops out of a student's quiz average, rather than returning
\texttt{NA} for their whole course score.

\subsection{Data frames: adding names and
program}\label{data-frames-adding-names-and-program}

A data frame lets us combine columns of different types -- here, the
students' names and program alongside their (numeric) scores -- into a
single table:

\begin{Shaded}
\begin{Highlighting}[]
\NormalTok{roster }\OtherTok{\textless{}{-}} \FunctionTok{data.frame}\NormalTok{(}\AttributeTok{name =}\NormalTok{ students, }\AttributeTok{program =}\NormalTok{ program, scores, }\AttributeTok{check.names =} \ConstantTok{FALSE}\NormalTok{)}

\NormalTok{roster}
\end{Highlighting}
\end{Shaded}

\begin{verbatim}
     name   program Quiz1 Quiz2 Quiz3 Quiz4 Quiz5 Test1 Test2 Final
S001 S001 Program A    64    64   100    66    73    55    83    81
S002 S002 Program B    70    83    98    65    51    66    64    70
S003 S003 Program B   100    73    67    61    61    69    68    76
S004 S004 Program A    76    71    83    59    73    58    NA    70
S005 S005 Program B    78    59    64    NA    87    67    62    83
S006 S006 Program A   100    77    62    87    41    64    70    74
S007 S007 Program A    84    62    56    40    67    70    70    59
S008 S008 Program B    50    45    NA    84    89    63    57    62
S009 S009 Program B    61    70   100    NA    63    75    53    76
S010 S010 Program A    66    96    71   100    79    79    76    95
S011 S011 Program A    99    66    74   100    84    87    84    92
S012 S012 Program B    82    90    77    95    79    85    79    81
S013 S013 Program B    83    46    97    45    87    79    68    61
S014 S014 Program A    77    77    62    68    72    69    55    70
S015 S015 Program B    64    NA    52    73    66    43    64    65
S016 S016 Program A   100    NA   100    94    21    76    84    98
S017 S017 Program A    NA    80    63    78    72    74    83    68
S018 S018 Program A    36    65    58    55    83    58    58    72
S019 S019 Program B    89    61    47   100    85    72    65    70
S020 S020 Program B    66    58    46    98    63    67    71    74
S021 S021 Program A    54    80    61    85   100    69    59    73
S022 S022 Program A    71    59    NA   100    80    73    72    83
S023 S023 Program B    54    68    94    53    NA    NA    65    60
S024 S024 Program A    60    73    86    93    99    73    65    90
S025 S025 Program B    62   100    65    70    60    NA    65    78
S026 S026 Program B    41    65    73    94    65    44    75    NA
S027 S027 Program B    92    83    58    79   100    63    72    NA
S028 S028 Program A    78    80    58    93    74    63    71    91
S029 S029 Program A    52    59    90   100   100    69    80    91
S030 S030 Program A   100    77    52    80    45    63    NA    74
S031 S031 Program B    84   100   100    NA    70    82    85    88
S032 S032 Program B    69    87    70    56    NA    72    71    64
S033 S033 Program B    93    79    76    66    80    63    69    74
S034 S034 Program A    93    70    57   100    54    64    62    86
S035 S035 Program B    91    37    61    88    74    54    52    71
S036 S036 Program B    89   100    46    39    52    55    53    66
S037 S037 Program A    86    49    68    53    88    65    73    67
S038 S038 Program A    74    93    80    76    96    83    NA    83
S039 S039 Program A    69   100    78    97    30    69    57    NA
S040 S040 Program A    67    49    56    78    99    51    63    77
S041 S041 Program A    61    92    56    92    NA    67    71    69
S042 S042 Program A    71    73    62    99    83    84    67    84
S043 S043 Program A    50    47   100   100    87    83    81    74
S044 S044 Program B   100    48    49    81    70    55    64    78
S045 S045 Program B    99    46    68    79    61    50    61    71
S046 S046 Program B    53    67   100    45    77    80    67    62
S047 S047 Program B    NA    49    70    82    83    61    65    69
S048 S048 Program B    66    92    45    69    92    63    61    68
S049 S049 Program A    91   100    59    61    33    70    69    76
S050 S050 Program A    73    52    82    76    41    51    69    72
S051 S051 Program B    80    94    64   100   100    70    79    76
S052 S052 Program B    74    93    61    40    95    52    59    65
S053 S053 Program B    74    NA    65    71    NA    61    73    74
S054 S054 Program A   100    58    74    82    88    77    86    93
S055 S055 Program B    70    76   100    62    92    59    72    83
S056 S056 Program A   100    72    70    87    21    73    69    71
S057 S057 Program B    44    89    94    88    96    74    71    83
S058 S058 Program B    87    71    85    79    64    NA    85    73
S059 S059 Program B    77    98    70    31    79    65    65    70
S060 S060 Program B    79    71    41   100    68    69    62    71
S061 S061 Program B    83    99    62    76    91    80    52    93
S062 S062 Program A    65    57    62    93    67    65    55    66
S063 S063 Program A    68    53    73    85    84    69    69    77
S064 S064 Program A    55   100    98   100    66    78    81    76
S065 S065 Program B    54    70    NA    96    52    77    68    75
S066 S066 Program B    81    84   100    76    98    85    83    89
S067 S067 Program B    84    91    69    88    NA    80    72    88
S068 S068 Program A    76    68    37    61   100    64    76    62
S069 S069 Program A    93    88    64    97    75    82    71    88
S070 S070 Program A   100    85    NA    NA    97    84    71    NA
S071 S071 Program B    65    74    89   100   100    66    65    78
S072 S072 Program A    29    79    91    47    68    48    68    67
S073 S073 Program B    95    77    86    71    48    71    66    72
S074 S074 Program A    61   100    44    97    69    59    78    80
S075 S075 Program A    61    63    89    90    70    72    70    82
S076 S076 Program B    96    56    63    68    77    66    59    77
S077 S077 Program B    NA    79    75    60   100    72    64    73
S078 S078 Program A    51    84    73    83    67    71    54    77
S079 S079 Program A    79    87    81    80    81    71    68    99
S080 S080 Program B    72    69    72    44    69    66    47    66
S081 S081 Program B    75    57    39    81    74    72    NA    65
S082 S082 Program A    NA   100    87    84    80    NA    77    81
S083 S083 Program A    68    71    80    83   100    61    83    83
S084 S084 Program A    88    61    67    59    76    79    80    73
S085 S085 Program B    71    73    74    90    88    69    68    75
S086 S086 Program A    82    74    NA   100    90    NA    82    83
S087 S087 Program A    97   100    76    89    92    90    71    88
S088 S088 Program A    84    80   100    57    63    73    67    76
S089 S089 Program A    68    93    68    71   100    75    76    82
S090 S090 Program B    98    68    75    NA    66    77    71    83
S091 S091 Program A    95    NA    95    67    72    83    72    91
S092 S092 Program B    86    72    93    37   100    76    90    72
S093 S093 Program A    NA    80    95    98   100   100    81    82
S094 S094 Program B    62    60    60    63    52    59    54    53
S095 S095 Program A   100    52   100    69    57    74    78    77
S096 S096 Program B    63   100    73   100    47    50    67    NA
S097 S097 Program B   100    NA   100    65    78    97    70    78
S098 S098 Program A   100    53    45    97    77    68    80    90
S099 S099 Program B    70    66    72    55    81    80    74    65
S100 S100 Program B    54    54    97    73    85    67    NA    62
\end{verbatim}

\section{Calculating a Single Student's
Score}\label{calculating-a-single-students-score}

Before writing a function, it often helps to work through the logic step
by step on a single example. Let's find a student who missed a quiz so
we can see how missing values are handled, and extract their record:

\begin{Shaded}
\begin{Highlighting}[]
\CommentTok{\# Find the index of the first student who has at least one NA in columns 1 to 5}
\NormalTok{example.ix }\OtherTok{\textless{}{-}} \FunctionTok{which}\NormalTok{(}\FunctionTok{rowSums}\NormalTok{(}\FunctionTok{is.na}\NormalTok{(scores[, }\DecValTok{1}\SpecialCharTok{:}\DecValTok{5}\NormalTok{])) }\SpecialCharTok{\textgreater{}} \DecValTok{0}\NormalTok{)[}\DecValTok{1}\NormalTok{]}
\NormalTok{student\_record }\OtherTok{\textless{}{-}}\NormalTok{ roster[example.ix, ]}

\NormalTok{student\_record}
\end{Highlighting}
\end{Shaded}

\begin{verbatim}
     name   program Quiz1 Quiz2 Quiz3 Quiz4 Quiz5 Test1 Test2 Final
S005 S005 Program B    78    59    64    NA    87    67    62    83
\end{verbatim}

To avoid hardcoding column positions, we can extract the column names
from this record and use \texttt{grep} to dynamically search for columns
starting with ``Quiz'', ``Test'', or ``Final'':

\begin{Shaded}
\begin{Highlighting}[]
\CommentTok{\# Extract all the column names from the student record}
\NormalTok{col\_names }\OtherTok{\textless{}{-}} \FunctionTok{names}\NormalTok{(student\_record)}
\NormalTok{col\_names}
\end{Highlighting}
\end{Shaded}

\begin{verbatim}
 [1] "name"    "program" "Quiz1"   "Quiz2"   "Quiz3"   "Quiz4"   "Quiz5"  
 [8] "Test1"   "Test2"   "Final"  
\end{verbatim}

\begin{Shaded}
\begin{Highlighting}[]
\CommentTok{\# Use grep to find the numeric indices of specific categories}
\FunctionTok{grep}\NormalTok{(}\StringTok{"\^{}Quiz"}\NormalTok{, col\_names, }\AttributeTok{value =} \ConstantTok{TRUE}\NormalTok{)}
\end{Highlighting}
\end{Shaded}

\begin{verbatim}
[1] "Quiz1" "Quiz2" "Quiz3" "Quiz4" "Quiz5"
\end{verbatim}

\begin{Shaded}
\begin{Highlighting}[]
\NormalTok{quiz.cols }\OtherTok{\textless{}{-}} \FunctionTok{grep}\NormalTok{(}\StringTok{"\^{}Quiz"}\NormalTok{, col\_names)}
\NormalTok{quiz.cols}
\end{Highlighting}
\end{Shaded}

\begin{verbatim}
[1] 3 4 5 6 7
\end{verbatim}

\begin{Shaded}
\begin{Highlighting}[]
\NormalTok{test.cols }\OtherTok{\textless{}{-}} \FunctionTok{grep}\NormalTok{(}\StringTok{"\^{}Test"}\NormalTok{, col\_names)}
\NormalTok{final.col }\OtherTok{\textless{}{-}} \FunctionTok{grep}\NormalTok{(}\StringTok{"\^{}Final"}\NormalTok{, col\_names)}
\end{Highlighting}
\end{Shaded}

\begin{Shaded}
\begin{Highlighting}[]
\CommentTok{\# Check which columns grep found for the quizzes}
\NormalTok{quiz.cols}
\end{Highlighting}
\end{Shaded}

\begin{verbatim}
[1] 3 4 5 6 7
\end{verbatim}

Next, we extract the scores using those indices, convert them to numeric
values, and calculate the mean for each category. Using
\texttt{na.rm\ =\ TRUE} ensures that the missing quiz is simply ignored
rather than turning the whole average into \texttt{NA}:

\begin{Shaded}
\begin{Highlighting}[]
\CommentTok{\# Calculate the mean for each category, ignoring NAs}
\NormalTok{quiz.mean }\OtherTok{\textless{}{-}} \FunctionTok{mean}\NormalTok{(}\FunctionTok{as.numeric}\NormalTok{(student\_record[quiz.cols]), }\AttributeTok{na.rm =} \ConstantTok{TRUE}\NormalTok{)}
\NormalTok{test.mean }\OtherTok{\textless{}{-}} \FunctionTok{mean}\NormalTok{(}\FunctionTok{as.numeric}\NormalTok{(student\_record[test.cols]), }\AttributeTok{na.rm =} \ConstantTok{TRUE}\NormalTok{)}
\NormalTok{final.score }\OtherTok{\textless{}{-}} \FunctionTok{as.numeric}\NormalTok{(student\_record[final.col])}

\FunctionTok{c}\NormalTok{(}\AttributeTok{Quiz =}\NormalTok{ quiz.mean, }\AttributeTok{Test =}\NormalTok{ test.mean, }\AttributeTok{Final =}\NormalTok{ final.score)}
\end{Highlighting}
\end{Shaded}

\begin{verbatim}
 Quiz  Test Final 
 72.0  64.5  83.0 
\end{verbatim}

Finally, we apply the syllabus weights to calculate the overall course
score:

\begin{Shaded}
\begin{Highlighting}[]
\CommentTok{\# Calculate the weighted course score}
\NormalTok{course.score }\OtherTok{\textless{}{-}}\NormalTok{ (}\FloatTok{0.3} \SpecialCharTok{*}\NormalTok{ quiz.mean) }\SpecialCharTok{+}\NormalTok{ (}\FloatTok{0.3} \SpecialCharTok{*}\NormalTok{ test.mean) }\SpecialCharTok{+}\NormalTok{ (}\FloatTok{0.4} \SpecialCharTok{*}\NormalTok{ final.score)}

\NormalTok{course.score}
\end{Highlighting}
\end{Shaded}

\begin{verbatim}
[1] 74.15
\end{verbatim}

\subsection{Aggregating into a Reusable
Function}\label{aggregating-into-a-reusable-function}

Now that we know the logic works perfectly for one student, we can
package those exact steps into a reusable function. This function will
take any \texttt{student\_record} as its input and return a named vector
containing the category means and the final score.

To make the grading policy more forgiving, we are adding a conditional
rule: if a student's final exam score is strictly higher than their test
average, the final exam score will replace the test average. Because
\texttt{compute.course.score} processes exactly one student at a time,
we evaluate a single logical condition using a standard \texttt{if}
statement rather than a vectorized alternative like \texttt{ifelse}. We
also use \texttt{!is.na()} checks to ensure the \texttt{if} statement
doesn't crash by attempting to evaluate a missing value (which could
happen if a student had no test or final scores).

\begin{Shaded}
\begin{Highlighting}[]
\NormalTok{compute.course.score }\OtherTok{\textless{}{-}} \ControlFlowTok{function}\NormalTok{(student\_record, }\AttributeTok{w\_quiz =} \FloatTok{0.3}\NormalTok{, }\AttributeTok{w\_test =} \FloatTok{0.3}\NormalTok{, }\AttributeTok{w\_final =} \FloatTok{0.4}\NormalTok{) \{}
    
\NormalTok{    col\_names }\OtherTok{\textless{}{-}} \FunctionTok{names}\NormalTok{(student\_record)}
\NormalTok{    quiz.cols }\OtherTok{\textless{}{-}} \FunctionTok{grep}\NormalTok{(}\StringTok{"\^{}Quiz"}\NormalTok{, col\_names)}
\NormalTok{    test.cols }\OtherTok{\textless{}{-}} \FunctionTok{grep}\NormalTok{(}\StringTok{"\^{}Test"}\NormalTok{, col\_names)}
\NormalTok{    final.col }\OtherTok{\textless{}{-}} \FunctionTok{grep}\NormalTok{(}\StringTok{"\^{}Final"}\NormalTok{, col\_names)}
    
\NormalTok{    quiz.mean }\OtherTok{\textless{}{-}} \FunctionTok{mean}\NormalTok{(}\FunctionTok{as.numeric}\NormalTok{(student\_record[quiz.cols]), }\AttributeTok{na.rm =} \ConstantTok{TRUE}\NormalTok{)}
\NormalTok{    test.mean }\OtherTok{\textless{}{-}} \FunctionTok{mean}\NormalTok{(}\FunctionTok{as.numeric}\NormalTok{(student\_record[test.cols]), }\AttributeTok{na.rm =} \ConstantTok{TRUE}\NormalTok{)}
\NormalTok{    final.score }\OtherTok{\textless{}{-}} \FunctionTok{as.numeric}\NormalTok{(student\_record[final.col])}
    
    \CommentTok{\# Explicit conditional check: replace test mean with final score if it is higher}
    \ControlFlowTok{if}\NormalTok{ (}\SpecialCharTok{!}\FunctionTok{is.na}\NormalTok{(final.score) }\SpecialCharTok{\&\&} \SpecialCharTok{!}\FunctionTok{is.na}\NormalTok{(test.mean) }\SpecialCharTok{\&\&}\NormalTok{ final.score }\SpecialCharTok{\textgreater{}}\NormalTok{ test.mean) \{}
\NormalTok{        test.mean }\OtherTok{\textless{}{-}}\NormalTok{ final.score}
\NormalTok{    \}}
    
    \CommentTok{\# The weights are now dynamic variables rather than hardcoded numbers}
\NormalTok{    course.score }\OtherTok{\textless{}{-}}\NormalTok{ (w\_quiz }\SpecialCharTok{*}\NormalTok{ quiz.mean) }\SpecialCharTok{+}\NormalTok{ (w\_test }\SpecialCharTok{*}\NormalTok{ test.mean) }\SpecialCharTok{+}\NormalTok{ (w\_final }\SpecialCharTok{*}\NormalTok{ final.score)}
    
    \FunctionTok{c}\NormalTok{(}\AttributeTok{Quiz =}\NormalTok{ quiz.mean, }\AttributeTok{Test =}\NormalTok{ test.mean, }\AttributeTok{Final =}\NormalTok{ final.score, }\AttributeTok{course.score =}\NormalTok{ course.score)}
\NormalTok{\}}
\end{Highlighting}
\end{Shaded}

Let's test our updated function on two specific scenarios to verify the
conditional logic. We will format these inputs as \textbf{named vectors}
rather than data frames. This mimics exactly what the \texttt{apply()}
function will do later when it passes rows of our dataset to this
function.

\textbf{Student A} struggled on the tests (average: 70) but aced the
final (95). The function should detect that
\texttt{95\ \textgreater{}\ 70} and replace their test mean with 95
before calculating the final course score. \textbf{Student B} did great
on the tests (average: 90) but dropped on the final (75). The function
should keep their original test mean.

\begin{Shaded}
\begin{Highlighting}[]
\CommentTok{\# Create the mock record for Student A as a named vector}
\NormalTok{student\_A }\OtherTok{\textless{}{-}} \FunctionTok{c}\NormalTok{(}
    \AttributeTok{name =} \StringTok{"Student\_A"}\NormalTok{, }\AttributeTok{program =} \StringTok{"Program A"}\NormalTok{,}
    \AttributeTok{Quiz1 =} \StringTok{"80"}\NormalTok{, }\AttributeTok{Quiz2 =} \StringTok{"80"}\NormalTok{, }\AttributeTok{Quiz3 =} \StringTok{"80"}\NormalTok{, }\AttributeTok{Quiz4 =} \StringTok{"80"}\NormalTok{, }\AttributeTok{Quiz5 =} \StringTok{"80"}\NormalTok{, }\CommentTok{\# Quiz mean = 80}
    \AttributeTok{Test1 =} \StringTok{"65"}\NormalTok{, }\AttributeTok{Test2 =} \StringTok{"75"}\NormalTok{,                                           }\CommentTok{\# Test mean = 70}
    \AttributeTok{Final =} \StringTok{"95"}
\NormalTok{)}

\CommentTok{\# Create the mock record for Student B as a named vector}
\NormalTok{student\_B }\OtherTok{\textless{}{-}} \FunctionTok{c}\NormalTok{(}
    \AttributeTok{name =} \StringTok{"Student\_B"}\NormalTok{, }\AttributeTok{program =} \StringTok{"Program B"}\NormalTok{,}
    \AttributeTok{Quiz1 =} \StringTok{"80"}\NormalTok{, }\AttributeTok{Quiz2 =} \StringTok{"80"}\NormalTok{, }\AttributeTok{Quiz3 =} \StringTok{"80"}\NormalTok{, }\AttributeTok{Quiz4 =} \StringTok{"80"}\NormalTok{, }\AttributeTok{Quiz5 =} \StringTok{"80"}\NormalTok{, }\CommentTok{\# Quiz mean = 80}
    \AttributeTok{Test1 =} \StringTok{"90"}\NormalTok{, }\AttributeTok{Test2 =} \StringTok{"90"}\NormalTok{,                                           }\CommentTok{\# Test mean = 90}
    \AttributeTok{Final =} \StringTok{"75"}
\NormalTok{)}

\CommentTok{\# Test Student A: The Test mean should be updated to 95 in the output}
\FunctionTok{cat}\NormalTok{(}\StringTok{"Student A Results:}\SpecialCharTok{\textbackslash{}n}\StringTok{"}\NormalTok{)}
\end{Highlighting}
\end{Shaded}

\begin{verbatim}
Student A Results:
\end{verbatim}

\begin{Shaded}
\begin{Highlighting}[]
\FunctionTok{compute.course.score}\NormalTok{(student\_A)}
\end{Highlighting}
\end{Shaded}

\begin{verbatim}
        Quiz         Test        Final course.score 
        80.0         95.0         95.0         90.5 
\end{verbatim}

\begin{Shaded}
\begin{Highlighting}[]
\CommentTok{\# Math check: (0.3 * 80) + (0.3 * 95) + (0.4 * 95) = 90.5}

\CommentTok{\# Test Student B: The Test mean should remain 90 in the output}
\FunctionTok{cat}\NormalTok{(}\StringTok{"}\SpecialCharTok{\textbackslash{}n}\StringTok{Student B Results:}\SpecialCharTok{\textbackslash{}n}\StringTok{"}\NormalTok{)}
\end{Highlighting}
\end{Shaded}

\begin{verbatim}

Student B Results:
\end{verbatim}

\begin{Shaded}
\begin{Highlighting}[]
\FunctionTok{compute.course.score}\NormalTok{(student\_B)}
\end{Highlighting}
\end{Shaded}

\begin{verbatim}
        Quiz         Test        Final course.score 
          80           90           75           81 
\end{verbatim}

\begin{Shaded}
\begin{Highlighting}[]
\CommentTok{\# Math check: (0.3 * 80) + (0.3 * 90) + (0.4 * 75) = 81}
\end{Highlighting}
\end{Shaded}

\section{Applying the Function to Every
Student}\label{applying-the-function-to-every-student}

\subsection{\texorpdfstring{Using \texttt{apply}
function}{Using apply function}}\label{using-apply-function}

Now that we have a function working for one student, we can apply it
across the entire class. The \texttt{apply} function is perfect for
this: it takes our \texttt{roster} data frame, loops over either rows
(\texttt{1}) or columns (\texttt{2}), and executes our custom function.

\begin{Shaded}
\begin{Highlighting}[]
\CommentTok{\# apply over margin 1 (rows)}
\NormalTok{all.results }\OtherTok{\textless{}{-}} \FunctionTok{t}\NormalTok{(}\FunctionTok{apply}\NormalTok{(roster, }\DecValTok{1}\NormalTok{, compute.course.score))}

\CommentTok{\# Combine the original demographic info with the calculated scores}
\NormalTok{results }\OtherTok{\textless{}{-}} \FunctionTok{cbind}\NormalTok{(roster[, }\FunctionTok{c}\NormalTok{(}\StringTok{"name"}\NormalTok{, }\StringTok{"program"}\NormalTok{)], all.results)}

\NormalTok{results}
\end{Highlighting}
\end{Shaded}

\begin{verbatim}
     name   program  Quiz Test Final course.score
S001 S001 Program A 73.40 81.0    81       78.720
S002 S002 Program B 73.40 70.0    70       71.020
S003 S003 Program B 72.40 76.0    76       74.920
S004 S004 Program A 72.40 70.0    70       70.720
S005 S005 Program B 72.00 83.0    83       79.700
S006 S006 Program A 73.40 74.0    74       73.820
S007 S007 Program A 61.80 70.0    59       63.140
S008 S008 Program B 67.00 62.0    62       63.500
S009 S009 Program B 73.50 76.0    76       75.250
S010 S010 Program A 82.40 95.0    95       91.220
S011 S011 Program A 84.60 92.0    92       89.780
S012 S012 Program B 84.60 82.0    81       82.380
S013 S013 Program B 71.60 73.5    61       67.930
S014 S014 Program A 71.20 70.0    70       70.360
S015 S015 Program B 63.75 65.0    65       64.625
S016 S016 Program A 78.75 98.0    98       92.225
S017 S017 Program A 73.25 78.5    68       72.725
S018 S018 Program A 59.40 72.0    72       68.220
S019 S019 Program B 76.40 70.0    70       71.920
S020 S020 Program B 66.20 74.0    74       71.660
S021 S021 Program A 76.00 73.0    73       73.900
S022 S022 Program A 77.50 83.0    83       81.350
S023 S023 Program B 67.25 65.0    60       63.675
S024 S024 Program A 82.20 90.0    90       87.660
S025 S025 Program B 71.40 78.0    78       76.020
S026 S026 Program B 67.60 59.5    NA           NA
S027 S027 Program B 82.40 67.5    NA           NA
S028 S028 Program A 76.60 91.0    91       86.680
S029 S029 Program A 80.20 91.0    91       87.760
S030 S030 Program A 70.80 74.0    74       73.040
S031 S031 Program B 88.50 88.0    88       88.150
S032 S032 Program B 70.50 71.5    64       68.200
S033 S033 Program B 78.80 74.0    74       75.440
S034 S034 Program A 74.80 86.0    86       82.640
S035 S035 Program B 70.20 71.0    71       70.760
S036 S036 Program B 65.20 66.0    66       65.760
S037 S037 Program A 68.80 69.0    67       68.140
S038 S038 Program A 83.80 83.0    83       83.240
S039 S039 Program A 74.80 63.0    NA           NA
S040 S040 Program A 69.80 77.0    77       74.840
S041 S041 Program A 75.25 69.0    69       70.875
S042 S042 Program A 77.60 84.0    84       82.080
S043 S043 Program A 76.80 82.0    74       77.240
S044 S044 Program B 69.60 78.0    78       75.480
S045 S045 Program B 70.60 71.0    71       70.880
S046 S046 Program B 68.40 73.5    62       67.370
S047 S047 Program B 71.00 69.0    69       69.600
S048 S048 Program B 72.80 68.0    68       69.440
S049 S049 Program A 68.80 76.0    76       73.840
S050 S050 Program A 64.80 72.0    72       69.840
S051 S051 Program B 87.60 76.0    76       79.480
S052 S052 Program B 72.60 65.0    65       67.280
S053 S053 Program B 70.00 74.0    74       72.800
S054 S054 Program A 80.40 93.0    93       89.220
S055 S055 Program B 80.00 83.0    83       82.100
S056 S056 Program A 70.00 71.0    71       70.700
S057 S057 Program B 82.20 83.0    83       82.760
S058 S058 Program B 77.20 85.0    73       77.860
S059 S059 Program B 71.00 70.0    70       70.300
S060 S060 Program B 71.80 71.0    71       71.240
S061 S061 Program B 82.20 93.0    93       89.760
S062 S062 Program A 68.80 66.0    66       66.840
S063 S063 Program A 72.60 77.0    77       75.680
S064 S064 Program A 83.80 79.5    76       79.390
S065 S065 Program B 68.00 75.0    75       72.900
S066 S066 Program B 87.80 89.0    89       88.640
S067 S067 Program B 83.00 88.0    88       86.500
S068 S068 Program A 68.40 70.0    62       66.320
S069 S069 Program A 83.40 88.0    88       86.620
S070 S070 Program A 94.00 77.5    NA           NA
S071 S071 Program B 85.60 78.0    78       80.280
S072 S072 Program A 62.80 67.0    67       65.740
S073 S073 Program B 75.40 72.0    72       73.020
S074 S074 Program A 74.20 80.0    80       78.260
S075 S075 Program A 74.60 82.0    82       79.780
S076 S076 Program B 72.00 77.0    77       75.500
S077 S077 Program B 78.50 73.0    73       74.650
S078 S078 Program A 71.60 77.0    77       75.380
S079 S079 Program A 81.60 99.0    99       93.780
S080 S080 Program B 65.20 66.0    66       65.760
S081 S081 Program B 65.20 72.0    65       67.160
S082 S082 Program A 87.75 81.0    81       83.025
S083 S083 Program A 80.40 83.0    83       82.220
S084 S084 Program A 70.20 79.5    73       74.110
S085 S085 Program B 79.20 75.0    75       76.260
S086 S086 Program A 86.50 83.0    83       84.050
S087 S087 Program A 90.80 88.0    88       88.840
S088 S088 Program A 76.80 76.0    76       76.240
S089 S089 Program A 80.00 82.0    82       81.400
S090 S090 Program B 76.75 83.0    83       81.125
S091 S091 Program A 82.25 91.0    91       88.375
S092 S092 Program B 77.60 83.0    72       76.980
S093 S093 Program A 93.25 90.5    82       87.925
S094 S094 Program B 59.40 56.5    53       55.970
S095 S095 Program A 75.60 77.0    77       76.580
S096 S096 Program B 76.60 58.5    NA           NA
S097 S097 Program B 85.75 83.5    78       81.975
S098 S098 Program A 74.40 90.0    90       85.320
S099 S099 Program B 68.80 77.0    65       69.740
S100 S100 Program B 72.60 67.0    62       66.680
\end{verbatim}

\subsection{\texorpdfstring{Looping the rows with
\texttt{for}}{Looping the rows with for}}\label{looping-the-rows-with-for}

While \texttt{apply()} is idiomatic R and very concise, we can also use
a traditional \texttt{for} loop to process each student. To do this
efficiently, we first pre-allocate an empty matrix of the correct size
to store our output. Then, we iterate through the row indices one by
one.

Inside the loop, we use \texttt{unlist()} to convert the single row of
our data frame into a named vector. This guarantees we are feeding our
\texttt{compute.course.score} function the exact same data structure
that \texttt{apply()} uses under the hood.

\begin{Shaded}
\begin{Highlighting}[]
\CommentTok{\# 1. Pre{-}allocate an empty matrix to hold the results}
\NormalTok{num\_students }\OtherTok{\textless{}{-}} \FunctionTok{nrow}\NormalTok{(roster)}
\NormalTok{loop\_results }\OtherTok{\textless{}{-}} \FunctionTok{matrix}\NormalTok{(}\ConstantTok{NA}\NormalTok{, }\AttributeTok{nrow =}\NormalTok{ num\_students, }\AttributeTok{ncol =} \DecValTok{4}\NormalTok{)}

\CommentTok{\# Set the column and row names to match our expected output}
\FunctionTok{colnames}\NormalTok{(loop\_results) }\OtherTok{\textless{}{-}} \FunctionTok{c}\NormalTok{(}\StringTok{"Quiz"}\NormalTok{, }\StringTok{"Test"}\NormalTok{, }\StringTok{"Final"}\NormalTok{, }\StringTok{"course.score"}\NormalTok{)}
\FunctionTok{rownames}\NormalTok{(loop\_results) }\OtherTok{\textless{}{-}}\NormalTok{ roster}\SpecialCharTok{$}\NormalTok{name}

\CommentTok{\# 2. Iterate over each row index}
\ControlFlowTok{for}\NormalTok{ (i }\ControlFlowTok{in} \DecValTok{1}\SpecialCharTok{:}\NormalTok{num\_students) \{}
    \CommentTok{\# Extract the single row as a named character vector}
\NormalTok{    student\_record }\OtherTok{\textless{}{-}} \FunctionTok{unlist}\NormalTok{(roster[i, ]) }
    
    \CommentTok{\# Compute the score and store it in the corresponding row of our matrix}
\NormalTok{    loop\_results[i, ] }\OtherTok{\textless{}{-}} \FunctionTok{compute.course.score}\NormalTok{(student\_record)}
\NormalTok{\}}

\CommentTok{\# Peek at the first few rows of the loop results}
\FunctionTok{head}\NormalTok{(loop\_results)}
\end{Highlighting}
\end{Shaded}

\begin{verbatim}
     Quiz Test Final course.score
S001 73.4   81    81        78.72
S002 73.4   70    70        71.02
S003 72.4   76    76        74.92
S004 72.4   70    70        70.72
S005 72.0   83    83        79.70
S006 73.4   74    74        73.82
\end{verbatim}

Now, let's run the exact same calculation using \texttt{apply()} and use
the \texttt{all.equal()} function to verify that both methods yield the
exact same mathematical results.

\begin{Shaded}
\begin{Highlighting}[]
\CommentTok{\# 3. Check if they are mathematically identical}
\CommentTok{\# (We set check.attributes = FALSE to ignore minor metadata differences, }
\CommentTok{\# like how apply() inherits row names vs. how we manually assigned them)}
\NormalTok{is\_same }\OtherTok{\textless{}{-}} \FunctionTok{all.equal}\NormalTok{(loop\_results, results, }\AttributeTok{check.attributes =} \ConstantTok{FALSE}\NormalTok{)}

\FunctionTok{print}\NormalTok{(}\FunctionTok{paste}\NormalTok{(}\StringTok{"Do the for loop and apply() give the exact same results?"}\NormalTok{, is\_same))}
\end{Highlighting}
\end{Shaded}

\begin{verbatim}
[1] "Do the for loop and apply() give the exact same results? target is matrix, current is data.frame"
\end{verbatim}

\subsection{Saving and reloading the
results}\label{saving-and-reloading-the-results}

The \texttt{results} table is worth keeping. A data frame can be written
as a plain csv file for spreadsheets, or saved as an \texttt{.RDS} file
to preserve it exactly (including the \texttt{program} factor) for
future R sessions:

\begin{Shaded}
\begin{Highlighting}[]
\FunctionTok{write.csv}\NormalTok{(results, }\AttributeTok{file =} \StringTok{"results.csv"}\NormalTok{, }\AttributeTok{row.names =} \ConstantTok{FALSE}\NormalTok{)}
\FunctionTok{read.csv}\NormalTok{ (}\StringTok{"results.csv"}\NormalTok{)}
\end{Highlighting}
\end{Shaded}

\begin{verbatim}
    name   program  Quiz Test Final course.score
1   S001 Program A 73.40 81.0    81       78.720
2   S002 Program B 73.40 70.0    70       71.020
3   S003 Program B 72.40 76.0    76       74.920
4   S004 Program A 72.40 70.0    70       70.720
5   S005 Program B 72.00 83.0    83       79.700
6   S006 Program A 73.40 74.0    74       73.820
7   S007 Program A 61.80 70.0    59       63.140
8   S008 Program B 67.00 62.0    62       63.500
9   S009 Program B 73.50 76.0    76       75.250
10  S010 Program A 82.40 95.0    95       91.220
11  S011 Program A 84.60 92.0    92       89.780
12  S012 Program B 84.60 82.0    81       82.380
13  S013 Program B 71.60 73.5    61       67.930
14  S014 Program A 71.20 70.0    70       70.360
15  S015 Program B 63.75 65.0    65       64.625
16  S016 Program A 78.75 98.0    98       92.225
17  S017 Program A 73.25 78.5    68       72.725
18  S018 Program A 59.40 72.0    72       68.220
19  S019 Program B 76.40 70.0    70       71.920
20  S020 Program B 66.20 74.0    74       71.660
21  S021 Program A 76.00 73.0    73       73.900
22  S022 Program A 77.50 83.0    83       81.350
23  S023 Program B 67.25 65.0    60       63.675
24  S024 Program A 82.20 90.0    90       87.660
25  S025 Program B 71.40 78.0    78       76.020
26  S026 Program B 67.60 59.5    NA           NA
27  S027 Program B 82.40 67.5    NA           NA
28  S028 Program A 76.60 91.0    91       86.680
29  S029 Program A 80.20 91.0    91       87.760
30  S030 Program A 70.80 74.0    74       73.040
31  S031 Program B 88.50 88.0    88       88.150
32  S032 Program B 70.50 71.5    64       68.200
33  S033 Program B 78.80 74.0    74       75.440
34  S034 Program A 74.80 86.0    86       82.640
35  S035 Program B 70.20 71.0    71       70.760
36  S036 Program B 65.20 66.0    66       65.760
37  S037 Program A 68.80 69.0    67       68.140
38  S038 Program A 83.80 83.0    83       83.240
39  S039 Program A 74.80 63.0    NA           NA
40  S040 Program A 69.80 77.0    77       74.840
41  S041 Program A 75.25 69.0    69       70.875
42  S042 Program A 77.60 84.0    84       82.080
43  S043 Program A 76.80 82.0    74       77.240
44  S044 Program B 69.60 78.0    78       75.480
45  S045 Program B 70.60 71.0    71       70.880
46  S046 Program B 68.40 73.5    62       67.370
47  S047 Program B 71.00 69.0    69       69.600
48  S048 Program B 72.80 68.0    68       69.440
49  S049 Program A 68.80 76.0    76       73.840
50  S050 Program A 64.80 72.0    72       69.840
51  S051 Program B 87.60 76.0    76       79.480
52  S052 Program B 72.60 65.0    65       67.280
53  S053 Program B 70.00 74.0    74       72.800
54  S054 Program A 80.40 93.0    93       89.220
55  S055 Program B 80.00 83.0    83       82.100
56  S056 Program A 70.00 71.0    71       70.700
57  S057 Program B 82.20 83.0    83       82.760
58  S058 Program B 77.20 85.0    73       77.860
59  S059 Program B 71.00 70.0    70       70.300
60  S060 Program B 71.80 71.0    71       71.240
61  S061 Program B 82.20 93.0    93       89.760
62  S062 Program A 68.80 66.0    66       66.840
63  S063 Program A 72.60 77.0    77       75.680
64  S064 Program A 83.80 79.5    76       79.390
65  S065 Program B 68.00 75.0    75       72.900
66  S066 Program B 87.80 89.0    89       88.640
67  S067 Program B 83.00 88.0    88       86.500
68  S068 Program A 68.40 70.0    62       66.320
69  S069 Program A 83.40 88.0    88       86.620
70  S070 Program A 94.00 77.5    NA           NA
71  S071 Program B 85.60 78.0    78       80.280
72  S072 Program A 62.80 67.0    67       65.740
73  S073 Program B 75.40 72.0    72       73.020
74  S074 Program A 74.20 80.0    80       78.260
75  S075 Program A 74.60 82.0    82       79.780
76  S076 Program B 72.00 77.0    77       75.500
77  S077 Program B 78.50 73.0    73       74.650
78  S078 Program A 71.60 77.0    77       75.380
79  S079 Program A 81.60 99.0    99       93.780
80  S080 Program B 65.20 66.0    66       65.760
81  S081 Program B 65.20 72.0    65       67.160
82  S082 Program A 87.75 81.0    81       83.025
83  S083 Program A 80.40 83.0    83       82.220
84  S084 Program A 70.20 79.5    73       74.110
85  S085 Program B 79.20 75.0    75       76.260
86  S086 Program A 86.50 83.0    83       84.050
87  S087 Program A 90.80 88.0    88       88.840
88  S088 Program A 76.80 76.0    76       76.240
89  S089 Program A 80.00 82.0    82       81.400
90  S090 Program B 76.75 83.0    83       81.125
91  S091 Program A 82.25 91.0    91       88.375
92  S092 Program B 77.60 83.0    72       76.980
93  S093 Program A 93.25 90.5    82       87.925
94  S094 Program B 59.40 56.5    53       55.970
95  S095 Program A 75.60 77.0    77       76.580
96  S096 Program B 76.60 58.5    NA           NA
97  S097 Program B 85.75 83.5    78       81.975
98  S098 Program A 74.40 90.0    90       85.320
99  S099 Program B 68.80 77.0    65       69.740
100 S100 Program B 72.60 67.0    62       66.680
\end{verbatim}

\begin{Shaded}
\begin{Highlighting}[]
\FunctionTok{saveRDS}\NormalTok{(results, }\AttributeTok{file =} \StringTok{"results.RDS"}\NormalTok{)}
\NormalTok{results.reloaded }\OtherTok{\textless{}{-}} \FunctionTok{readRDS}\NormalTok{(}\StringTok{"results.RDS"}\NormalTok{)}
\FunctionTok{identical}\NormalTok{(results, results.reloaded)}
\end{Highlighting}
\end{Shaded}

\begin{verbatim}
[1] TRUE
\end{verbatim}

\section{Writing a Function for the Whole Process with List
Output}\label{writing-a-function-for-the-whole-process-with-list-output}

Sometimes a function needs to return multiple distinct pieces of
information that don't fit neatly into a single table. A \textbf{list}
in R is a flexible container that can bundle objects of completely
different types and sizes together.

Let's write a wrapper function that takes our raw \texttt{roster} data
and our weighting scheme. It will calculate the scores, but instead of
just returning the table, it will return a list containing four very
different items:

\begin{enumerate}
\def\labelenumi{\arabic{enumi}.}
\tightlist
\item
  The full processed data frame.
\item
  A numeric vector of just the final course scores.
\item
  A frequency table of scores binned into 10-point intervals.
\item
  A character vector of the names of students who failed (scored below
  60).
\end{enumerate}

\begin{Shaded}
\begin{Highlighting}[]
\NormalTok{analyze\_class\_performance }\OtherTok{\textless{}{-}} \ControlFlowTok{function}\NormalTok{(roster\_data, }\AttributeTok{w\_quiz =} \FloatTok{0.3}\NormalTok{, }\AttributeTok{w\_test =} \FloatTok{0.3}\NormalTok{, }\AttributeTok{w\_final =} \FloatTok{0.4}\NormalTok{) \{}
    
    \CommentTok{\# 1. Pass the weights down into compute.course.score}
\NormalTok{    all.results }\OtherTok{\textless{}{-}} \FunctionTok{t}\NormalTok{(}\FunctionTok{apply}\NormalTok{(roster\_data, }\DecValTok{1}\NormalTok{, compute.course.score, }
                           \AttributeTok{w\_quiz =}\NormalTok{ w\_quiz, }\AttributeTok{w\_test =}\NormalTok{ w\_test, }\AttributeTok{w\_final =}\NormalTok{ w\_final))}
                           
\NormalTok{    processed\_table }\OtherTok{\textless{}{-}} \FunctionTok{cbind}\NormalTok{(roster\_data[, }\FunctionTok{c}\NormalTok{(}\StringTok{"name"}\NormalTok{, }\StringTok{"program"}\NormalTok{)], all.results)}
    
    \CommentTok{\# 2. Extract final course scores }
\NormalTok{    final\_scores }\OtherTok{\textless{}{-}}\NormalTok{ processed\_table}\SpecialCharTok{$}\NormalTok{course.score}
    
    \CommentTok{\# 3. Create bins (handling NAs)}
\NormalTok{    bins }\OtherTok{\textless{}{-}} \FunctionTok{seq}\NormalTok{(}\DecValTok{0}\NormalTok{, }\DecValTok{100}\NormalTok{, }\AttributeTok{by =} \DecValTok{10}\NormalTok{)}
\NormalTok{    score\_bins }\OtherTok{\textless{}{-}} \FunctionTok{table}\NormalTok{(}\FunctionTok{cut}\NormalTok{(final\_scores, }\AttributeTok{breaks =}\NormalTok{ bins, }\AttributeTok{include.lowest =} \ConstantTok{TRUE}\NormalTok{), }\AttributeTok{useNA =} \StringTok{"ifany"}\NormalTok{)}
    
    \CommentTok{\# 4. Extract names of failing students (safely ignoring NAs using which)}
\NormalTok{    failing\_students }\OtherTok{\textless{}{-}}\NormalTok{ processed\_table}\SpecialCharTok{$}\NormalTok{name[}\FunctionTok{which}\NormalTok{(final\_scores }\SpecialCharTok{\textless{}} \DecValTok{60}\NormalTok{)]}
    
    \CommentTok{\# 5. Extract names of students with missing overall scores}
\NormalTok{    incomplete\_students }\OtherTok{\textless{}{-}}\NormalTok{ processed\_table}\SpecialCharTok{$}\NormalTok{name[}\FunctionTok{is.na}\NormalTok{(final\_scores)]}
    
    \CommentTok{\# Bundle everything into a named list}
    \FunctionTok{list}\NormalTok{(}
        \AttributeTok{ProcessedTable =}\NormalTok{ processed\_table,}
        \AttributeTok{CourseScores =}\NormalTok{ final\_scores,}
        \AttributeTok{ScoreDistribution =}\NormalTok{ score\_bins,}
        \AttributeTok{Failing\_Students =}\NormalTok{ failing\_students,}
        \AttributeTok{incomplete\_students =}\NormalTok{ incomplete\_students}
\NormalTok{    )}
\NormalTok{\}}
\end{Highlighting}
\end{Shaded}

When we run this function, it returns a single list object. We can use
the \texttt{str()} function to quickly inspect the structure of our new
list and see all the different data types nested inside it:

\begin{Shaded}
\begin{Highlighting}[]
\CommentTok{\# Run the analysis}
\NormalTok{class\_report }\OtherTok{\textless{}{-}} \FunctionTok{analyze\_class\_performance}\NormalTok{(roster)}

\CommentTok{\# Inspect the structure of the returned list}
\FunctionTok{str}\NormalTok{(class\_report, }\AttributeTok{max.level =} \DecValTok{1}\NormalTok{)}
\end{Highlighting}
\end{Shaded}

\begin{verbatim}
List of 5
 $ ProcessedTable     :'data.frame':    100 obs. of  6 variables:
 $ CourseScores       : num [1:100] 78.7 71 74.9 70.7 79.7 ...
 $ ScoreDistribution  : 'table' int [1:11(1d)] 0 0 0 0 0 1 21 42 28 3 ...
  ..- attr(*, "dimnames")=List of 1
 $ Failing_Students   : chr "S094"
 $ incomplete_students: chr [1:5] "S026" "S027" "S039" "S070" ...
\end{verbatim}

To access individual elements inside a list, you use the \texttt{\$}
operator (just like with data frames) or double brackets
\texttt{{[}{[}\ {]}{]}}. Let's look at the score distribution and the
failing students:

\begin{Shaded}
\begin{Highlighting}[]
\CommentTok{\# Access the 10{-}point bin distribution}
\NormalTok{class\_report}\SpecialCharTok{$}\NormalTok{ScoreDistribution}
\end{Highlighting}
\end{Shaded}

\begin{verbatim}

  [0,10]  (10,20]  (20,30]  (30,40]  (40,50]  (50,60]  (60,70]  (70,80] 
       0        0        0        0        0        1       21       42 
 (80,90] (90,100]     <NA> 
      28        3        5 
\end{verbatim}

\begin{Shaded}
\begin{Highlighting}[]
\CommentTok{\# Access the names of the failing students}
\NormalTok{class\_report}\SpecialCharTok{$}\NormalTok{Failing\_Students}
\end{Highlighting}
\end{Shaded}

\begin{verbatim}
[1] "S094"
\end{verbatim}

\begin{Shaded}
\begin{Highlighting}[]
\NormalTok{class\_report}\SpecialCharTok{$}\NormalTok{incomplete\_students }
\end{Highlighting}
\end{Shaded}

\begin{verbatim}
[1] "S026" "S027" "S039" "S070" "S096"
\end{verbatim}

\section{\texorpdfstring{The use of pipe
\texttt{\textbar{}\textgreater{}}}{The use of pipe \textbar\textgreater{}}}\label{the-use-of-pipe}

The pipe operator \texttt{x\ \textbar{}\textgreater{}\ f(...)} is just
shorthand for \texttt{f(x,\ ...)} -- it takes whatever is on its left
and inserts it as the \emph{first} argument of the function call on its
right. This lets you write a sequence of transformations in the order
you think of them, instead of nesting calls inside out:

\begin{Shaded}
\begin{Highlighting}[]
\CommentTok{\# Without the pipe: read from the inside out}
\FunctionTok{round}\NormalTok{(}\FunctionTok{sqrt}\NormalTok{(}\FunctionTok{mean}\NormalTok{(}\FunctionTok{c}\NormalTok{(}\DecValTok{4}\NormalTok{, }\DecValTok{9}\NormalTok{, }\DecValTok{16}\NormalTok{, }\DecValTok{25}\NormalTok{))), }\DecValTok{2}\NormalTok{)}
\end{Highlighting}
\end{Shaded}

\begin{verbatim}
[1] 3.67
\end{verbatim}

\begin{Shaded}
\begin{Highlighting}[]
\CommentTok{\# With the pipe: read top to bottom, left to right}
\FunctionTok{c}\NormalTok{(}\DecValTok{4}\NormalTok{, }\DecValTok{9}\NormalTok{, }\DecValTok{16}\NormalTok{, }\DecValTok{25}\NormalTok{) }\SpecialCharTok{|\textgreater{}}
    \FunctionTok{mean}\NormalTok{() }\SpecialCharTok{|\textgreater{}}
    \FunctionTok{sqrt}\NormalTok{() }\SpecialCharTok{|\textgreater{}}
    \FunctionTok{round}\NormalTok{(}\DecValTok{2}\NormalTok{)}
\end{Highlighting}
\end{Shaded}

\begin{verbatim}
[1] 3.67
\end{verbatim}

Each step's output becomes the next step's first input, so
\texttt{mean()} receives the vector, \texttt{sqrt()} receives the mean,
and \texttt{round(2)} receives the square root (with \texttt{2} filling
the \emph{second} argument, \texttt{digits}). Extra arguments after the
first are passed straight through exactly as they would be in a normal
call, which matters once a pipeline mixes several multi-argument
functions -- much like the quiz-generating recipe used throughout this
unit:

\begin{Shaded}
\begin{Highlighting}[]
\FunctionTok{set.seed}\NormalTok{(}\DecValTok{123}\NormalTok{)}

\CommentTok{\# Same idea as gen.quiz(), written as one pipeline: each step\textquotesingle{}s output}
\CommentTok{\# fills the *first* argument of the next call, while the remaining,}
\CommentTok{\# named arguments are supplied just as in an ordinary function call}
\FunctionTok{rnorm}\NormalTok{(}\DecValTok{20}\NormalTok{, }\AttributeTok{mean =} \DecValTok{75}\NormalTok{, }\AttributeTok{sd =} \DecValTok{10}\NormalTok{) }\SpecialCharTok{|\textgreater{}}
    \FunctionTok{round}\NormalTok{(}\AttributeTok{digits =} \DecValTok{0}\NormalTok{) }\SpecialCharTok{|\textgreater{}}       \CommentTok{\# round(x, digits = 0)}
    \FunctionTok{pmin}\NormalTok{(}\DecValTok{100}\NormalTok{) }\SpecialCharTok{|\textgreater{}}                \CommentTok{\# pmin(x, 100)}
    \FunctionTok{pmax}\NormalTok{(}\DecValTok{0}\NormalTok{) }\SpecialCharTok{|\textgreater{}}                  \CommentTok{\# pmax(x, 0)}
    \FunctionTok{sort}\NormalTok{(}\AttributeTok{decreasing =} \ConstantTok{TRUE}\NormalTok{) }\SpecialCharTok{|\textgreater{}}  \CommentTok{\# sort(x, decreasing = TRUE)}
    \FunctionTok{head}\NormalTok{(}\AttributeTok{n =} \DecValTok{5}\NormalTok{)                 }\CommentTok{\# head(x, n = 5)}
\end{Highlighting}
\end{Shaded}

\begin{verbatim}
[1] 93 92 91 87 82
\end{verbatim}

Sometimes the piped value needs to land in an argument \emph{other than}
the first. The placeholder \texttt{\_} marks exactly where it goes, as
long as that argument is named:

\begin{Shaded}
\begin{Highlighting}[]
\CommentTok{\# seq(1, to = 10, by = 2) {-}{-} the piped 10 fills the "to" argument}
\DecValTok{10} \SpecialCharTok{|\textgreater{}} \FunctionTok{seq}\NormalTok{(}\DecValTok{1}\NormalTok{, }\AttributeTok{to =}\NormalTok{ \_, }\AttributeTok{by =} \DecValTok{2}\NormalTok{)}
\end{Highlighting}
\end{Shaded}

\begin{verbatim}
[1] 1 3 5 7 9
\end{verbatim}

This same idea -- chaining verbs with \texttt{\textbar{}\textgreater{}}
-- is what makes the dplyr pipelines below read like a recipe.

\section{\texorpdfstring{Data Manipulation with
\texttt{dplyr}}{Data Manipulation with dplyr}}\label{data-manipulation-with-dplyr}

Base R's indexing, \texttt{apply()}, and \texttt{for} loops can do
everything we've done so far, but \textbf{dplyr} offers a grammar of a
handful of verbs -- \texttt{mutate()}, \texttt{filter()},
\texttt{select()}, \texttt{arrange()}, \texttt{group\_by()},
\texttt{summarize()} -- that are chained together with the pipe
\texttt{\textbar{}\textgreater{}}, so a whole data-wrangling recipe
reads top to bottom like a sentence instead of nesting function calls
inside one another.

\subsection{\texorpdfstring{Generating a roster with
\texttt{mutate()}}{Generating a roster with mutate()}}\label{generating-a-roster-with-mutate}

\texttt{mutate()} isn't limited to computing new columns from data that
already exists in a data frame -- it can build up a data frame from
nothing at all, one column at a time, chaining together exactly the same
generative steps we used in
\hyperref[generating-an-artificial-score-dataset]{Generating an
artificial score dataset}. Within a single \texttt{mutate()} call, each
new column is available to the expressions that follow it, so we can go
straight from raw quiz scores to \texttt{quiz.mean.true} to
\texttt{Test1}/\texttt{Test2}/\texttt{Final}, all in one pipeline,
reusing the \texttt{gen.quiz()} and \texttt{clip.round()} helpers we
already defined:

\begin{Shaded}
\begin{Highlighting}[]
\FunctionTok{library}\NormalTok{(dplyr)}

\FunctionTok{set.seed}\NormalTok{(}\DecValTok{2026}\NormalTok{)   }\CommentTok{\# a fresh cohort, just for this demonstration}

\NormalTok{n2 }\OtherTok{\textless{}{-}} \DecValTok{100}
\NormalTok{students2 }\OtherTok{\textless{}{-}} \FunctionTok{sprintf}\NormalTok{(}\StringTok{"S\%03d"}\NormalTok{, }\DecValTok{1}\SpecialCharTok{:}\NormalTok{n2)}
\NormalTok{program2  }\OtherTok{\textless{}{-}} \FunctionTok{factor}\NormalTok{(}\FunctionTok{sample}\NormalTok{(}\FunctionTok{c}\NormalTok{(}\StringTok{"Program A"}\NormalTok{, }\StringTok{"Program B"}\NormalTok{), n2, }\AttributeTok{replace =} \ConstantTok{TRUE}\NormalTok{))}

\FunctionTok{data.frame}\NormalTok{(}\AttributeTok{name =}\NormalTok{ students2, }\AttributeTok{program =}\NormalTok{ program2) }\SpecialCharTok{|\textgreater{}}
    \FunctionTok{mutate}\NormalTok{(}
        \AttributeTok{Quiz1 =} \FunctionTok{gen.quiz}\NormalTok{(}\DecValTok{75}\NormalTok{, }\DecValTok{20}\NormalTok{),}
        \AttributeTok{Quiz2 =} \FunctionTok{gen.quiz}\NormalTok{(}\DecValTok{78}\NormalTok{, }\DecValTok{20}\NormalTok{),}
        \AttributeTok{Quiz3 =} \FunctionTok{gen.quiz}\NormalTok{(}\DecValTok{72}\NormalTok{, }\DecValTok{20}\NormalTok{),}
        \AttributeTok{Quiz4 =} \FunctionTok{gen.quiz}\NormalTok{(}\DecValTok{80}\NormalTok{, }\DecValTok{20}\NormalTok{),}
        \AttributeTok{Quiz5 =} \FunctionTok{gen.quiz}\NormalTok{(}\DecValTok{74}\NormalTok{, }\DecValTok{20}\NormalTok{),}
        \AttributeTok{quiz.mean.true =} \FunctionTok{rowMeans}\NormalTok{(}\FunctionTok{cbind}\NormalTok{(Quiz1, Quiz2, Quiz3, Quiz4, Quiz5)),}
        \AttributeTok{b0.test  =} \FunctionTok{if\_else}\NormalTok{(program }\SpecialCharTok{==} \StringTok{"Program A"}\NormalTok{, }\DecValTok{12}\NormalTok{, }\DecValTok{6}\NormalTok{),}
        \AttributeTok{b0.final =} \FunctionTok{if\_else}\NormalTok{(program }\SpecialCharTok{==} \StringTok{"Program A"}\NormalTok{, }\DecValTok{15}\NormalTok{, }\DecValTok{8}\NormalTok{),}
        \AttributeTok{Test1 =} \FunctionTok{clip.round}\NormalTok{(b0.test  }\SpecialCharTok{+} \FloatTok{0.80} \SpecialCharTok{*}\NormalTok{ quiz.mean.true }\SpecialCharTok{+} \FunctionTok{rnorm}\NormalTok{(}\FunctionTok{n}\NormalTok{(), }\DecValTok{0}\NormalTok{, }\DecValTok{8}\NormalTok{)),}
        \AttributeTok{Test2 =} \FunctionTok{clip.round}\NormalTok{(b0.test  }\SpecialCharTok{+} \FloatTok{0.80} \SpecialCharTok{*}\NormalTok{ quiz.mean.true }\SpecialCharTok{+} \FunctionTok{rnorm}\NormalTok{(}\FunctionTok{n}\NormalTok{(), }\DecValTok{0}\NormalTok{, }\DecValTok{8}\NormalTok{)),}
        \AttributeTok{Final =} \FunctionTok{clip.round}\NormalTok{(b0.final }\SpecialCharTok{+} \FloatTok{0.85} \SpecialCharTok{*}\NormalTok{ quiz.mean.true }\SpecialCharTok{+} \FunctionTok{rnorm}\NormalTok{(}\FunctionTok{n}\NormalTok{(), }\DecValTok{0}\NormalTok{, }\DecValTok{6}\NormalTok{)),}
        \FunctionTok{across}\NormalTok{(}\FunctionTok{c}\NormalTok{(}\FunctionTok{starts\_with}\NormalTok{(}\StringTok{"Quiz"}\NormalTok{), }\FunctionTok{starts\_with}\NormalTok{(}\StringTok{"Test"}\NormalTok{), Final),}
               \SpecialCharTok{\textasciitilde{}} \FunctionTok{replace}\NormalTok{(.x, }\FunctionTok{sample}\NormalTok{(}\FunctionTok{n}\NormalTok{(), }\FunctionTok{round}\NormalTok{(}\FunctionTok{n}\NormalTok{() }\SpecialCharTok{*}\NormalTok{ miss.rate)), }\ConstantTok{NA}\NormalTok{))}
\NormalTok{    ) }\OtherTok{{-}\textgreater{}}\NormalTok{ roster\_worktable}

\NormalTok{roster\_worktable }\SpecialCharTok{|\textgreater{}} \FunctionTok{slice\_head}\NormalTok{ (}\AttributeTok{n =} \DecValTok{6}\NormalTok{)}
\end{Highlighting}
\end{Shaded}

\begin{verbatim}
  name   program Quiz1 Quiz2 Quiz3 Quiz4 Quiz5 quiz.mean.true b0.test b0.final
1 S001 Program A    52    78    66    NA    84           76.0      12       15
2 S002 Program A    71    NA    86    NA    97             NA      12       15
3 S003 Program A    56    98    64   100    80           79.6      12       15
4 S004 Program B    84    88    70    65    75           76.4       6        8
5 S005 Program A    62    90   100    47   100           79.8      12       15
6 S006 Program A    70   100    59    78    88           79.0      12       15
  Test1 Test2 Final
1    80    75    82
2    78    79    84
3    80    77    85
4    65    60    72
5    70    78    82
6    87    86    73
\end{verbatim}

\begin{Shaded}
\begin{Highlighting}[]
\NormalTok{roster\_mutate }\OtherTok{\textless{}{-}}\NormalTok{ roster\_worktable }\SpecialCharTok{|\textgreater{}} \FunctionTok{select}\NormalTok{(}\SpecialCharTok{{-}}\NormalTok{quiz.mean.true, }\SpecialCharTok{{-}}\NormalTok{b0.test, }\SpecialCharTok{{-}}\NormalTok{b0.final) }

\NormalTok{roster\_mutate }\SpecialCharTok{|\textgreater{}} \FunctionTok{slice\_head}\NormalTok{ (}\AttributeTok{n=} \DecValTok{6}\NormalTok{)}
\end{Highlighting}
\end{Shaded}

\begin{verbatim}
  name   program Quiz1 Quiz2 Quiz3 Quiz4 Quiz5 Test1 Test2 Final
1 S001 Program A    52    78    66    NA    84    80    75    82
2 S002 Program A    71    NA    86    NA    97    78    79    84
3 S003 Program A    56    98    64   100    80    80    77    85
4 S004 Program B    84    88    70    65    75    65    60    72
5 S005 Program A    62    90   100    47   100    70    78    82
6 S006 Program A    70   100    59    78    88    87    86    73
\end{verbatim}

A few things are worth noticing here. \texttt{b0.test} and
\texttt{b0.final} are computed with \texttt{if\_else()} -- the
vectorized, type-checked dplyr analogue of \texttt{ifelse()} -- directly
from the \texttt{program} column, and then used (and finally dropped
again with \texttt{select(-...)}) just like the helper vectors from the
original script. \texttt{n()}, dplyr's built-in count of the current
number of rows, replaces the standalone \texttt{n\ \textless{}-\ 100} we
relied on earlier, so this pipeline needs nothing from the outside
except \texttt{students2}/\texttt{program2} to start from. The
tidyselect helper \texttt{starts\_with()} -- the dplyr analogue of the
\texttt{grep("\^{}Quiz",\ ...)} trick from earlier -- picks out groups
of columns by name, and wrapping it in \texttt{across()} applies the
\emph{same} NA-injection rule to all eight score columns at once, in
place of the explicit \texttt{for} loop over \texttt{1:ncol(scores)}.

Because this pipeline draws its own random numbers with a fresh seed,
\texttt{roster\_mutate} is a different simulated class from
\texttt{roster} -- not a reproduction of it -- but it is built by
\emph{exactly} the same recipe, entirely inside one \texttt{mutate()}
chain. The rest of this section goes back to analyzing the original
\texttt{roster}.

\subsection{\texorpdfstring{Recomputing course scores with
\texttt{mutate()} and
\texttt{rowwise()}}{Recomputing course scores with mutate() and rowwise()}}\label{recomputing-course-scores-with-mutate-and-rowwise}

Our course-score calculation is inherently a \textbf{row-wise}
operation: for each student we need to look across several quiz columns
and average just that student's values. \texttt{rowwise()} tells dplyr
to treat each row as its own little group, and \texttt{c\_across()} lets
\texttt{mean()} reach across a \emph{selection} of columns (again using
\texttt{starts\_with()}) within that row.

Inside a single \texttt{mutate()} call, later expressions can refer to
columns just created earlier in the \emph{same} call -- so we can
reproduce the ``replace the test mean with the final score if it's
higher'' rule from \texttt{compute.course.score()} using
\texttt{if\_else()} in place of the base-R \texttt{if} statement:

\begin{Shaded}
\begin{Highlighting}[]
\NormalTok{dplyr\_results }\OtherTok{\textless{}{-}}\NormalTok{ roster }\SpecialCharTok{|\textgreater{}}
    \FunctionTok{rowwise}\NormalTok{() }\SpecialCharTok{|\textgreater{}}
    \FunctionTok{mutate}\NormalTok{(}
        \AttributeTok{Quiz  =} \FunctionTok{mean}\NormalTok{(}\FunctionTok{c\_across}\NormalTok{(}\FunctionTok{starts\_with}\NormalTok{(}\StringTok{"Quiz"}\NormalTok{)), }\AttributeTok{na.rm =} \ConstantTok{TRUE}\NormalTok{),}
        \AttributeTok{Test  =} \FunctionTok{mean}\NormalTok{(}\FunctionTok{c\_across}\NormalTok{(}\FunctionTok{starts\_with}\NormalTok{(}\StringTok{"Test"}\NormalTok{)), }\AttributeTok{na.rm =} \ConstantTok{TRUE}\NormalTok{),}
        \AttributeTok{Test  =} \FunctionTok{if\_else}\NormalTok{(}\SpecialCharTok{!}\FunctionTok{is.na}\NormalTok{(Final) }\SpecialCharTok{\&} \SpecialCharTok{!}\FunctionTok{is.na}\NormalTok{(Test) }\SpecialCharTok{\&}\NormalTok{ Final }\SpecialCharTok{\textgreater{}}\NormalTok{ Test, Final, Test),}
        \AttributeTok{course.score =} \FloatTok{0.3} \SpecialCharTok{*}\NormalTok{ Quiz }\SpecialCharTok{+} \FloatTok{0.3} \SpecialCharTok{*}\NormalTok{ Test }\SpecialCharTok{+} \FloatTok{0.4} \SpecialCharTok{*}\NormalTok{ Final}
\NormalTok{    ) }\SpecialCharTok{|\textgreater{}}
    \FunctionTok{ungroup}\NormalTok{() }\SpecialCharTok{|\textgreater{}}
    \FunctionTok{select}\NormalTok{(name, program, Quiz, Test, Final, course.score)}

\NormalTok{dplyr\_results}
\end{Highlighting}
\end{Shaded}

\begin{verbatim}
# A tibble: 100 x 6
   name  program    Quiz  Test Final course.score
   <chr> <fct>     <dbl> <dbl> <dbl>        <dbl>
 1 S001  Program A  73.4    81    81         78.7
 2 S002  Program B  73.4    70    70         71.0
 3 S003  Program B  72.4    76    76         74.9
 4 S004  Program A  72.4    70    70         70.7
 5 S005  Program B  72      83    83         79.7
 6 S006  Program A  73.4    74    74         73.8
 7 S007  Program A  61.8    70    59         63.1
 8 S008  Program B  67      62    62         63.5
 9 S009  Program B  73.5    76    76         75.2
10 S010  Program A  82.4    95    95         91.2
# i 90 more rows
\end{verbatim}

\texttt{rowwise()} groups by row, so we \texttt{ungroup()} again
immediately afterward -- otherwise every later verb would keep operating
one row at a time. As a sanity check, this dplyr pipeline should produce
\emph{exactly} the same numbers as \texttt{results}, our earlier
\texttt{apply()}-based table, just as the \texttt{for}-loop version did:

\begin{Shaded}
\begin{Highlighting}[]
\FunctionTok{all.equal}\NormalTok{(}\FunctionTok{as.data.frame}\NormalTok{(dplyr\_results), results, }\AttributeTok{check.attributes =} \ConstantTok{FALSE}\NormalTok{)}
\end{Highlighting}
\end{Shaded}

\begin{verbatim}
[1] TRUE
\end{verbatim}

\subsection{\texorpdfstring{Summarizing by group with
\texttt{group\_by()} and
\texttt{summarize()}}{Summarizing by group with group\_by() and summarize()}}\label{summarizing-by-group-with-group_by-and-summarize}

This is where dplyr's advantage over base R really shows. Producing a
by-program summary table with \texttt{tapply()} or nested
\texttt{sapply()} calls in base R takes some fiddling; with dplyr it's a
direct continuation of the same pipeline. \texttt{group\_by()} splits
the data by \texttt{program}, and every subsequent \texttt{summarize()}
argument is computed separately within each group:

\begin{Shaded}
\begin{Highlighting}[]
\NormalTok{dplyr\_results }\SpecialCharTok{|\textgreater{}}
    \FunctionTok{group\_by}\NormalTok{(program) }\SpecialCharTok{|\textgreater{}}
    \FunctionTok{summarize}\NormalTok{(}
        \AttributeTok{n                 =} \FunctionTok{n}\NormalTok{(),}
        \AttributeTok{mean\_quiz         =} \FunctionTok{mean}\NormalTok{(Quiz, }\AttributeTok{na.rm =} \ConstantTok{TRUE}\NormalTok{),}
        \AttributeTok{mean\_test         =} \FunctionTok{mean}\NormalTok{(Test, }\AttributeTok{na.rm =} \ConstantTok{TRUE}\NormalTok{),}
        \AttributeTok{mean\_final        =} \FunctionTok{mean}\NormalTok{(Final, }\AttributeTok{na.rm =} \ConstantTok{TRUE}\NormalTok{),}
        \AttributeTok{mean\_course\_score =} \FunctionTok{mean}\NormalTok{(course.score, }\AttributeTok{na.rm =} \ConstantTok{TRUE}\NormalTok{),}
        \AttributeTok{.groups =} \StringTok{"drop"}
\NormalTok{    )}
\end{Highlighting}
\end{Shaded}

\begin{verbatim}
# A tibble: 2 x 6
  program       n mean_quiz mean_test mean_final mean_course_score
  <fct>     <int>     <dbl>     <dbl>      <dbl>             <dbl>
1 Program A    50      76.3      80.2       79.4              78.7
2 Program B    50      74.1      74.1       72.9              73.8
\end{verbatim}

Program A's mean course score comes out noticeably higher -- exactly the
intercept difference (\texttt{b0.test}, \texttt{b0.final}) we built into
the simulation, and the same pattern we'll see again in the by-program
boxplot below.

\subsection{\texorpdfstring{Filtering and sorting with \texttt{filter()}
and
\texttt{arrange()}}{Filtering and sorting with filter() and arrange()}}\label{filtering-and-sorting-with-filter-and-arrange}

\texttt{filter()} keeps rows matching a logical condition (replacing the
\texttt{which()}-based indexing we used for \texttt{Failing\_Students}),
and \texttt{arrange()} sorts rows by one or more columns --
\texttt{desc()} reverses the order. Chaining them together gives the
failing students, worst first, and separately the top five performers in
the class:

\begin{Shaded}
\begin{Highlighting}[]
\CommentTok{\# Failing students, worst first}
\NormalTok{dplyr\_results }\SpecialCharTok{|\textgreater{}}
    \FunctionTok{filter}\NormalTok{(course.score }\SpecialCharTok{\textless{}} \DecValTok{60}\NormalTok{) }\SpecialCharTok{|\textgreater{}}
    \FunctionTok{arrange}\NormalTok{(course.score) }\SpecialCharTok{|\textgreater{}}
    \FunctionTok{select}\NormalTok{(name, program, course.score)}
\end{Highlighting}
\end{Shaded}

\begin{verbatim}
# A tibble: 1 x 3
  name  program   course.score
  <chr> <fct>            <dbl>
1 S094  Program B         56.0
\end{verbatim}

\begin{Shaded}
\begin{Highlighting}[]
\CommentTok{\# Top 5 students in the class}
\NormalTok{dplyr\_results }\SpecialCharTok{|\textgreater{}}
    \FunctionTok{arrange}\NormalTok{(}\FunctionTok{desc}\NormalTok{(course.score)) }\SpecialCharTok{|\textgreater{}}
    \FunctionTok{select}\NormalTok{(name, program, course.score) }\SpecialCharTok{|\textgreater{}}
    \FunctionTok{slice\_head}\NormalTok{(}\AttributeTok{n =} \DecValTok{5}\NormalTok{)}
\end{Highlighting}
\end{Shaded}

\begin{verbatim}
# A tibble: 5 x 3
  name  program   course.score
  <chr> <fct>            <dbl>
1 S079  Program A         93.8
2 S016  Program A         92.2
3 S010  Program A         91.2
4 S011  Program A         89.8
5 S061  Program B         89.8
\end{verbatim}

We'll keep using \texttt{class\_report\$ProcessedTable} -- our original
\texttt{apply()}-based results -- for the reports and plots below, but
from here on in the course, dplyr's pipe-and-verb style will be our
default tool for wrangling real datasets.

\section{Presenting the Results}\label{presenting-the-results}

\subsection{\texorpdfstring{Presenting the Processed Results with
\texttt{gt}}{Presenting the Processed Results with gt}}\label{presenting-the-processed-results-with-gt}

Finally, \textbf{gt} turns the raw data frame of results into a
polished, presentation-ready report. We can extract the
\texttt{ProcessedTable} directly from our \texttt{class\_report} list
and enhance it.

Here we use \texttt{tab\_spanner()} to group the category scores and the
final score under distinct visual headers, and \texttt{fmt\_number()} to
ensure all grades display exactly two decimal places. The column headers
below spell out the underlying formulas: the quiz average
\(\sum_{i}^5{q_i}/5\), the test score \(T\), the final exam \(F_1\), and
the overall course score \(F\).

\begin{Shaded}
\begin{Highlighting}[]
\FunctionTok{library}\NormalTok{(gt)}

\CommentTok{\# Extract the full processed table from our list output}
\CommentTok{\# (We use head() to display just the first 10 rows in the report)}
\NormalTok{full\_results }\OtherTok{\textless{}{-}} \FunctionTok{head}\NormalTok{(class\_report}\SpecialCharTok{$}\NormalTok{ProcessedTable, }\DecValTok{10}\NormalTok{)}

\FunctionTok{gt}\NormalTok{(full\_results) }\SpecialCharTok{|\textgreater{}}
    \FunctionTok{tab\_header}\NormalTok{(}
        \AttributeTok{title =} \StringTok{"Class Performance Report"}\NormalTok{,}
        \AttributeTok{subtitle =} \StringTok{"Individual category means and final calculated course scores"}
\NormalTok{    ) }\SpecialCharTok{|\textgreater{}}
    \CommentTok{\# Student Information}
    \FunctionTok{tab\_spanner}\NormalTok{(}
        \AttributeTok{label =} \StringTok{"Student Information"}\NormalTok{,}
        \AttributeTok{columns =} \FunctionTok{c}\NormalTok{(name, program)}
\NormalTok{    )}\SpecialCharTok{|\textgreater{}}
    \CommentTok{\# Group the individual graded categories}
    \FunctionTok{tab\_spanner}\NormalTok{(}
        \AttributeTok{label =} \StringTok{"Category Breakdown"}\NormalTok{,}
        \AttributeTok{columns =} \FunctionTok{c}\NormalTok{(Quiz, Test, Final)}
\NormalTok{    ) }\SpecialCharTok{|\textgreater{}}
    \CommentTok{\# Group the final processed output}
    \FunctionTok{tab\_spanner}\NormalTok{(}
        \AttributeTok{label =} \StringTok{"Submitted Score"}\NormalTok{,}
        \AttributeTok{columns =} \FunctionTok{c}\NormalTok{(course.score)}
\NormalTok{    ) }\SpecialCharTok{|\textgreater{}}
    \CommentTok{\# Enforce two decimal places for all numeric score columns}
    \FunctionTok{fmt\_number}\NormalTok{(}
        \AttributeTok{columns =} \FunctionTok{c}\NormalTok{(Quiz, Test, Final),}
        \AttributeTok{decimals =} \DecValTok{2}
\NormalTok{    ) }\SpecialCharTok{|\textgreater{}}
    \FunctionTok{fmt\_number}\NormalTok{(}
        \AttributeTok{columns =} \FunctionTok{c}\NormalTok{(course.score),}
        \AttributeTok{decimals =} \DecValTok{0}
\NormalTok{    ) }\SpecialCharTok{|\textgreater{}}
    \FunctionTok{cols\_label}\NormalTok{(}
        \AttributeTok{name =} \StringTok{"Name"}\NormalTok{,}
        \AttributeTok{program =} \StringTok{"Program"}\NormalTok{,}
        \AttributeTok{Quiz =} \FunctionTok{md}\NormalTok{(}\StringTok{"Quizzes ($}\SpecialCharTok{\textbackslash{}\textbackslash{}}\StringTok{sum\_\{i\}\^{}5\{q\_i\}/5$)"}\NormalTok{),}
        \AttributeTok{Test =} \FunctionTok{md}\NormalTok{(}\StringTok{"Tests ($T$)"}\NormalTok{),}
        \AttributeTok{Final =} \FunctionTok{md}\NormalTok{(}\StringTok{"Final Exam ($F\_1$)"}\NormalTok{),}
        \AttributeTok{course.score =} \FunctionTok{md}\NormalTok{(}\StringTok{"Course Score ($F$)"}\NormalTok{)}
\NormalTok{    ) }\SpecialCharTok{|\textgreater{}}
    \CommentTok{\# Shrink the name/program columns so the score columns have more room}
    \FunctionTok{cols\_width}\NormalTok{(}
\NormalTok{        name }\SpecialCharTok{\textasciitilde{}} \FunctionTok{px}\NormalTok{(}\DecValTok{100}\NormalTok{),}
\NormalTok{        program }\SpecialCharTok{\textasciitilde{}} \FunctionTok{px}\NormalTok{(}\DecValTok{90}\NormalTok{)}
\NormalTok{    ) }\SpecialCharTok{|\textgreater{}}
    \FunctionTok{tab\_stubhead}\NormalTok{(}\AttributeTok{label =} \StringTok{"Student ID"}\NormalTok{)}
\end{Highlighting}
\end{Shaded}

\begin{table}
\caption*{
{\fontsize{20}{25}\selectfont  Class Performance Report\fontsize{12}{15}\selectfont } \\ 
{\fontsize{14}{17}\selectfont  Individual category means and final calculated course scores\fontsize{12}{15}\selectfont }
} 
\fontsize{12.0pt}{14.0pt}\selectfont
\begin{tabular*}{\linewidth}{@{\extracolsep{\fill}}>{\raggedright\arraybackslash}p{\dimexpr 75.00pt -2\tabcolsep-1.5\arrayrulewidth}>{\centering\arraybackslash}p{\dimexpr 67.50pt -2\tabcolsep-1.5\arrayrulewidth}rrrr}
\toprule
\multicolumn{2}{>{\centering\arraybackslash}m{\dimexpr 142.50pt -2\tabcolsep-1.5\arrayrulewidth}}{\parbox{\linewidth}{\centering {Student Information}}} & \multicolumn{3}{c}{{Category Breakdown}} & \multicolumn{1}{c}{{Submitted Score}} \\ 
\cmidrule(lr){1-2} \cmidrule(lr){3-5} \cmidrule(lr){6-6}
Name & Program & Quizzes (\(\sum_{i}^5{q_i}/5\)) & Tests (\(T\)) & Final Exam (\(F_1\)) & Course Score (\(F\)) \\ 
\midrule\addlinespace[2.5pt]
S001 & Program A & 73.40 & 81.00 & 81.00 & 79 \\ 
S002 & Program B & 73.40 & 70.00 & 70.00 & 71 \\ 
S003 & Program B & 72.40 & 76.00 & 76.00 & 75 \\ 
S004 & Program A & 72.40 & 70.00 & 70.00 & 71 \\ 
S005 & Program B & 72.00 & 83.00 & 83.00 & 80 \\ 
S006 & Program A & 73.40 & 74.00 & 74.00 & 74 \\ 
S007 & Program A & 61.80 & 70.00 & 59.00 & 63 \\ 
S008 & Program B & 67.00 & 62.00 & 62.00 & 64 \\ 
S009 & Program B & 73.50 & 76.00 & 76.00 & 75 \\ 
S010 & Program A & 82.40 & 95.00 & 95.00 & 91 \\ 
\bottomrule
\end{tabular*}
\end{table}

\subsection{Base R graphics}\label{base-r-graphics}

A quick histogram shows the distribution of course scores, with a
vertical line marking the class average:

\begin{Shaded}
\begin{Highlighting}[]
\FunctionTok{hist}\NormalTok{(class\_report}\SpecialCharTok{$}\NormalTok{ProcessedTable}\SpecialCharTok{$}\NormalTok{course.score, }\AttributeTok{breaks =} \DecValTok{15}\NormalTok{, }\AttributeTok{col =} \StringTok{"lightblue"}\NormalTok{,}
     \AttributeTok{main =} \StringTok{"Distribution of Course Scores"}\NormalTok{, }\AttributeTok{xlab =} \StringTok{"Course Score"}\NormalTok{)}
\FunctionTok{abline}\NormalTok{(}\AttributeTok{v =} \FunctionTok{mean}\NormalTok{(class\_report}\SpecialCharTok{$}\NormalTok{ProcessedTable}\SpecialCharTok{$}\NormalTok{course.score, }\AttributeTok{na.rm =} \ConstantTok{TRUE}\NormalTok{), }\AttributeTok{col =} \StringTok{"red"}\NormalTok{, }\AttributeTok{lwd =} \DecValTok{2}\NormalTok{)}
\end{Highlighting}
\end{Shaded}

\pandocbounded{\includegraphics[keepaspectratio]{chapters/unit01_introduction_longer_files/figure-pdf/unnamed-chunk-39-1.pdf}}

\subsection{Plotting with ggplot2}\label{plotting-with-ggplot2}

\textbf{ggplot2} builds a comparison across programs declaratively, by
mapping columns of \texttt{class\_report\$ProcessedTable} to visual
properties and layering geometries with \texttt{+}. Here the boxplot is
enhanced with the individual course scores jittered on top:

\begin{Shaded}
\begin{Highlighting}[]
\FunctionTok{library}\NormalTok{(ggplot2)}

\FunctionTok{ggplot}\NormalTok{(class\_report}\SpecialCharTok{$}\NormalTok{ProcessedTable, }\FunctionTok{aes}\NormalTok{(}\AttributeTok{x =}\NormalTok{ program, }\AttributeTok{y =}\NormalTok{ course.score, }\AttributeTok{fill =}\NormalTok{ program)) }\SpecialCharTok{+}
    \FunctionTok{geom\_boxplot}\NormalTok{(}\AttributeTok{show.legend =} \ConstantTok{FALSE}\NormalTok{, }\AttributeTok{alpha =} \FloatTok{0.6}\NormalTok{) }\SpecialCharTok{+}
    \FunctionTok{geom\_jitter}\NormalTok{(}\AttributeTok{width =} \FloatTok{0.05}\NormalTok{, }\AttributeTok{size =} \FloatTok{1.5}\NormalTok{, }\AttributeTok{alpha =} \FloatTok{0.6}\NormalTok{) }\SpecialCharTok{+}
    \FunctionTok{labs}\NormalTok{(}
        \AttributeTok{title =} \StringTok{"Course Score by Academic Program"}\NormalTok{,}
        \AttributeTok{x =} \StringTok{"Program"}\NormalTok{, }\AttributeTok{y =} \StringTok{"Course Score"}
\NormalTok{    ) }\SpecialCharTok{+}
    \FunctionTok{theme\_minimal}\NormalTok{()}
\end{Highlighting}
\end{Shaded}

\pandocbounded{\includegraphics[keepaspectratio]{chapters/unit01_introduction_longer_files/figure-pdf/unnamed-chunk-40-1.pdf}}

We built a program-specific intercept into the simulation when
generating the final exam scores; a scatter plot of the final versus
each student's mean quiz score, with a fitted line per program, should
reveal it:

\begin{Shaded}
\begin{Highlighting}[]
\FunctionTok{ggplot}\NormalTok{(class\_report}\SpecialCharTok{$}\NormalTok{ProcessedTable, }\FunctionTok{aes}\NormalTok{(}\AttributeTok{x =}\NormalTok{ Quiz, }\AttributeTok{y =}\NormalTok{ Final, }\AttributeTok{color =}\NormalTok{ program)) }\SpecialCharTok{+}
    \FunctionTok{geom\_point}\NormalTok{(}\AttributeTok{size =} \DecValTok{2}\NormalTok{, }\AttributeTok{alpha =} \FloatTok{0.6}\NormalTok{) }\SpecialCharTok{+}
    \FunctionTok{geom\_smooth}\NormalTok{(}\AttributeTok{method =} \StringTok{"lm"}\NormalTok{, }\AttributeTok{se =} \ConstantTok{FALSE}\NormalTok{) }\SpecialCharTok{+}
    \FunctionTok{labs}\NormalTok{(}
        \AttributeTok{title =} \StringTok{"Final Exam vs. Mean Quiz Score, by Program"}\NormalTok{,}
        \AttributeTok{x =} \StringTok{"Mean Quiz Score"}\NormalTok{, }\AttributeTok{y =} \StringTok{"Final Exam Score"}\NormalTok{, }\AttributeTok{color =} \StringTok{"Program"}
\NormalTok{    ) }\SpecialCharTok{+}
    \FunctionTok{theme\_minimal}\NormalTok{()}
\end{Highlighting}
\end{Shaded}

\pandocbounded{\includegraphics[keepaspectratio]{chapters/unit01_introduction_longer_files/figure-pdf/unnamed-chunk-41-1.pdf}}

The two fitted lines are roughly parallel but vertically offset --
exactly the intercept difference we built into \texttt{b0.final} when
simulating the data.

\bookmarksetup{startatroot}

\chapter{Computer Arithmetic}\label{computer-arithmetic}

Statistical computation is arithmetic on a machine that stores every
number in a fixed number of binary digits. This chapter describes how
integers and real numbers are represented, what goes wrong when a
computation exceeds the capacity of that representation (overflow,
underflow, and roundoff error), and the standard techniques for writing
statistical code that avoids these failures: reordering sums, rewriting
formulas, and working on the logarithmic scale.

\section{Representations of Numbers in
Computer}\label{representations-of-numbers-in-computer}

\subsection{Place value}\label{place-value}

Every number is a string of digits in boxes, and each box has a
\textbf{place value}. In base 10,

10³

10²

10¹

10⁰

1

3

4

8

\(1348 = 1\times10^3 + 3\times10^2 + 4\times10^1 + 8\times10^0\)

A computer uses base 2, so each box, called a \textbf{bit}, holds 0 or
1:

2³

2²

2¹

2⁰

1

1

0

1

\(1101_2 = 1\times2^3 + 1\times2^2 + 0\times2^1 + 1\times2^0 = 13\)

Everything that follows is a consequence of one fact: the number of
boxes is fixed in advance.

\subsection{Integers}\label{integers}

An integer is stored in a fixed number of bits, typically 32 or 64. With
\(p\) bits the largest unsigned value is obtained by filling every box
with 1, \[
\sum_{i=0}^{p-1} 2^i = 2^p - 1 .
\] A 4-bit example: the largest storable integer is \(1111_2 = 15\), and
\(15 + 1 = 16 = 10000_2\) needs a fifth box that does not exist.

2⁴

2³

2²

2¹

2⁰

1

0

0

0

0

The leading 1 has nowhere to go: \textbf{integer overflow}

To represent negative values, one of the boxes is used as a \textbf{sign
bit}, leaving \(p-1\) boxes for the magnitude. With 32 bits, a signed
integer \(u\) is \[
u = (-1)^{s}\sum_{i=0}^{30} d_i\,2^i, \qquad s\in\{0,1\}.
\] The largest integer has the sign bit \(0\) and the other 31 boxes all
equal to \(1\):

sign

2³⁰

2²⁹

\ldots{}

2¹

2⁰

+

1

1

\ldots{}

1

1

\(u^{\max} = \sum_{i=0}^{30} 2^i = 2^{31}-1 = 2147483647\)

and the smallest is its negative, \(u^{\min} = -(2^{31}-1)\). When an
integer computation leaves this range, R does not wrap around silently,
as C does; it returns \texttt{NA} with a warning.

\begin{Shaded}
\begin{Highlighting}[]
\NormalTok{.Machine}\SpecialCharTok{$}\NormalTok{integer.max}
\end{Highlighting}
\end{Shaded}

\begin{verbatim}
[1] 2147483647
\end{verbatim}

\begin{Shaded}
\begin{Highlighting}[]
\DocumentationTok{\#\# Check upper bounds}
\FunctionTok{as.integer}\NormalTok{(}\DecValTok{2}\SpecialCharTok{\^{}}\DecValTok{31} \SpecialCharTok{{-}} \DecValTok{1}\NormalTok{)                    }\CommentTok{\# maximum valid integer}
\end{Highlighting}
\end{Shaded}

\begin{verbatim}
[1] 2147483647
\end{verbatim}

\begin{Shaded}
\begin{Highlighting}[]
\FunctionTok{as.integer}\NormalTok{(}\DecValTok{2}\SpecialCharTok{\^{}}\DecValTok{31}\NormalTok{)                        }\CommentTok{\# exceeds limit, returns NA}
\end{Highlighting}
\end{Shaded}

\begin{verbatim}
[1] NA
\end{verbatim}

\begin{Shaded}
\begin{Highlighting}[]
\FunctionTok{as.integer}\NormalTok{(}\DecValTok{2}\SpecialCharTok{\^{}}\DecValTok{31} \SpecialCharTok{{-}} \DecValTok{1}\NormalTok{) }\SpecialCharTok{+} \FunctionTok{as.integer}\NormalTok{(}\DecValTok{1}\NormalTok{)    }\CommentTok{\# overflows to NA}
\end{Highlighting}
\end{Shaded}

\begin{verbatim}
[1] NA
\end{verbatim}

\begin{Shaded}
\begin{Highlighting}[]
\DocumentationTok{\#\# Check lower bounds}
\FunctionTok{as.integer}\NormalTok{(}\SpecialCharTok{{-}}\NormalTok{(}\DecValTok{2}\SpecialCharTok{\^{}}\DecValTok{31} \SpecialCharTok{{-}} \DecValTok{1}\NormalTok{))                 }\CommentTok{\# minimum valid integer}
\end{Highlighting}
\end{Shaded}

\begin{verbatim}
[1] -2147483647
\end{verbatim}

\begin{Shaded}
\begin{Highlighting}[]
\FunctionTok{as.integer}\NormalTok{(}\SpecialCharTok{{-}}\DecValTok{2}\SpecialCharTok{\^{}}\DecValTok{31}\NormalTok{)                       }\CommentTok{\# exceeds lower limit, returns NA}
\end{Highlighting}
\end{Shaded}

\begin{verbatim}
[1] NA
\end{verbatim}

A practical consequence: products of moderately large integers overflow
easily. \texttt{50000L\ *\ 50000L} is \texttt{NA} in R, while
\texttt{50000\ *\ 50000} (double precision) is fine. R converts to
double automatically for most arithmetic, so integer overflow is rarely
a problem in R code; it is a common one in C, Java, and database
systems.

\begin{Shaded}
\begin{Highlighting}[]
\DecValTok{50000}\NormalTok{L }\SpecialCharTok{*} \DecValTok{50000}\NormalTok{L}
\end{Highlighting}
\end{Shaded}

\begin{verbatim}
[1] NA
\end{verbatim}

\begin{Shaded}
\begin{Highlighting}[]
\DecValTok{50000} \SpecialCharTok{*} \DecValTok{50000}
\end{Highlighting}
\end{Shaded}

\begin{verbatim}
[1] 2.5e+09
\end{verbatim}

\subsection{Floating point numbers}\label{floating-point-numbers}

A real number is stored in the form of scientific notation: a
\textbf{sign}, a \textbf{mantissa} holding the significant digits, and
an \textbf{exponent}: \[
x = (-1)^{s} \times 1.d_1 d_2 \cdots d_p \times 2^{\,e}.
\]

sign

mantissa (fixed number of digits)

exponent

+

1

.

d₁

d₂ \ldots{} dp

×2\^{}

±

e

More precisely, a floating point number \(x\) is represented by a vector
of bits \((s,d_0,\ldots,d_{t-1},e_1,\ldots, e_{k-1})\) as, \[
x = (-1)^s \sum_{i=0}^{t-1}d_i2^{-i}\times 2^{\sum_{i=0}^{k-1}e_i2^i-2^{k-1}}.
\] The leading mantissa bit is always 1 for a normalized number and need
not be stored, and the exponent is stored with an offset so that both
negative and positive exponents fit in unsigned bits. Double precision
(\texttt{double} in R, the IEEE 754 standard) uses 64 bits:

{\def\LTcaptype{none} % do not increment counter
\begin{longtable}[]{@{}
  >{\raggedright\arraybackslash}p{(\linewidth - 4\tabcolsep) * \real{0.1333}}
  >{\raggedleft\arraybackslash}p{(\linewidth - 4\tabcolsep) * \real{0.1333}}
  >{\raggedright\arraybackslash}p{(\linewidth - 4\tabcolsep) * \real{0.7333}}@{}}
\toprule\noalign{}
\begin{minipage}[b]{\linewidth}\raggedright
part
\end{minipage} & \begin{minipage}[b]{\linewidth}\raggedleft
bits
\end{minipage} & \begin{minipage}[b]{\linewidth}\raggedright
consequence
\end{minipage} \\
\midrule\noalign{}
\endhead
\bottomrule\noalign{}
\endlastfoot
sign & 1 & \\
mantissa & 52 & about 16 significant decimal digits \\
exponent & 11 & \(e \in [-1022,\ 1023]\), i.e.~magnitudes from
\(10^{-308}\) to \(10^{308}\) \\
\end{longtable}
}

\subsection{A toy floating point
system}\label{a-toy-floating-point-system}

To see what goes wrong, it is easier to work in a decimal toy computer
with a \textbf{4-digit mantissa} and a \textbf{1-digit exponent},
\(e\in[-9,9]\):

sign

1

10⁻¹

10⁻²

10⁻³

exp

+

3

.

1

4

2

×10\^{}

0

\(3.142\times10^{0}\), the closest storable value to \(\pi\)

Three things can go wrong, all caused by the boxes being finite:

\begin{enumerate}
\def\labelenumi{\arabic{enumi}.}
\item
  \textbf{Overflow}: the exponent box is too small for a large number.
\item
  \textbf{Underflow}: the exponent box is too small for a tiny number.
\item
  \textbf{Roundoff}: the mantissa boxes are too few for all the digits.
\end{enumerate}

\subsection{Significance of the Mantissa: A Zoom
Model}\label{significance-of-the-mantissa-a-zoom-model}

Think of a floating point number as a \textbf{ruler} and a \textbf{zoom
setting}. The mantissa is a fixed ruler with equally spaced marks; the
exponent is a zoom factor \(10^{e}\) that stretches the ruler to cover
the interval \([10^{e}, 10^{e+1})\). The number of marks never changes,
so the more the ruler is stretched, the farther apart the marks land on
the number line.

The simplest case has a \textbf{one-digit mantissa}: the ruler has nine
marks, \(1, 2, \ldots, 9\), and a number is stored as \(d\times10^{e}\)
with \(d\in\{1,\ldots,9\}\).

1

exp

d

×10\^{}

e

One mantissa digit: the ruler marks are
\(1\times10^{e}, 2\times10^{e}, \ldots, 9\times10^{e}\)

{\def\LTcaptype{none} % do not increment counter
\begin{longtable}[]{@{}cllc@{}}
\toprule\noalign{}
zoom \(e\) & interval \([10^e, 10^{e+1})\) & storable numbers & gap \\
\midrule\noalign{}
\endhead
\bottomrule\noalign{}
\endlastfoot
0 & \([1, 10)\) & \(1, 2, 3, \ldots, 9\) & \(1\) \\
1 & \([10, 100)\) & \(10, 20, 30, \ldots, 90\) & \(10\) \\
2 & \([100, 1000)\) & \(100, 200, \ldots, 900\) & \(100\) \\
\(-1\) & \([0.1, 1)\) & \(0.1, 0.2, \ldots, 0.9\) & \(0.1\) \\
\end{longtable}
}

Within each zoom level the storable numbers are evenly spaced, and the
spacing is multiplied by \(10\) each time the zoom goes up one level.
Every real number between two marks is rounded to the nearer mark:
\(24\to20\), \(27\to30\), \(2.4\to2\), and \(0.27\to0.3\).

Adding mantissa digits adds marks to the ruler without changing the zoom
mechanism. With a leading digit and \(t\) further digits there are
\(9\times10^{t}\) marks in \([10^e, 10^{e+1})\), and the gap at zoom
\(e\) is \(10^{e-t}\).

\begin{figure}

\centering{

\pandocbounded{\includegraphics[keepaspectratio]{chapters/unit02_comparith_files/figure-pdf/fig-zoom-model-1.pdf}}

}

\caption{\label{fig-zoom-model}Representable numbers in \([1, 100)\) for
a mantissa of 1 and 2 decimal digits. The two zoom levels \([1, 10)\)
and \([10, 100)\) are shaded; within each, the marks are evenly spaced,
and the spacing is ten times larger at the higher zoom.}

\end{figure}%

The picture explains the two facts that matter for computation:

\begin{itemize}
\tightlist
\item
  The \textbf{absolute} gap grows with the magnitude, because the same
  ruler is stretched more; the \textbf{relative} gap, gap divided by
  magnitude, stays between \(10^{-t}\) and \(10^{-t-1}\) at every zoom
  level.
\item
  The relative gap near \(1\) is the \textbf{machine epsilon}
  \(\epsilon\): the smallest \(\epsilon\) such that
  \(1+\epsilon \neq 1\) on the computer. It is the distance between the
  first two marks of the ruler at zoom \(e = 0\). In the one-digit
  system \(\epsilon = 1\), so \(1 + 0.4\) rounds back to \(1\); in the
  four-digit toy system \(\epsilon = 0.001\).
\end{itemize}

A binary computer works the same way with a ruler of \(2^{t}\) marks per
doubling instead of \(9\times10^{t}\) marks per decade. For doubles,
\(t = 52\) and \(\epsilon = 2^{-52}\approx 2.2\times10^{-16}\).

\begin{Shaded}
\begin{Highlighting}[]
\NormalTok{.Machine}\SpecialCharTok{$}\NormalTok{double.eps      }\CommentTok{\# 2\^{}{-}52}
\end{Highlighting}
\end{Shaded}

\begin{verbatim}
[1] 2.220446e-16
\end{verbatim}

\begin{Shaded}
\begin{Highlighting}[]
\DecValTok{1} \SpecialCharTok{+} \DecValTok{2}\SpecialCharTok{\^{}{-}}\DecValTok{52} \SpecialCharTok{==} \DecValTok{1}
\end{Highlighting}
\end{Shaded}

\begin{verbatim}
[1] FALSE
\end{verbatim}

\begin{Shaded}
\begin{Highlighting}[]
\DecValTok{1} \SpecialCharTok{+} \DecValTok{2}\SpecialCharTok{\^{}{-}}\DecValTok{53} \SpecialCharTok{==} \DecValTok{1}           \CommentTok{\# half an epsilon rounds back to 1}
\end{Highlighting}
\end{Shaded}

\begin{verbatim}
[1] TRUE
\end{verbatim}

For doubles the gap at zoom level \([2^e, 2^{e+1})\) is \(2^{e-52}\).
Near \(2^{52}\) the spacing is exactly \(1\), and in
\([2^{53}, 2^{54})\) it is \(2\), so odd integers can no longer be
represented.

\begin{Shaded}
\begin{Highlighting}[]
\DecValTok{2}\SpecialCharTok{\^{}}\DecValTok{52} \SpecialCharTok{+} \DecValTok{1} \SpecialCharTok{==} \DecValTok{2}\SpecialCharTok{\^{}}\DecValTok{52}             \CommentTok{\# spacing is 1: still distinguishable}
\end{Highlighting}
\end{Shaded}

\begin{verbatim}
[1] FALSE
\end{verbatim}

\begin{Shaded}
\begin{Highlighting}[]
\DecValTok{2}\SpecialCharTok{\^{}}\DecValTok{53} \SpecialCharTok{+} \DecValTok{1} \SpecialCharTok{==} \DecValTok{2}\SpecialCharTok{\^{}}\DecValTok{53}             \CommentTok{\# spacing is 2: the +1 is lost}
\end{Highlighting}
\end{Shaded}

\begin{verbatim}
[1] TRUE
\end{verbatim}

\begin{Shaded}
\begin{Highlighting}[]
\FunctionTok{format}\NormalTok{(}\FunctionTok{c}\NormalTok{(}\DecValTok{2}\SpecialCharTok{\^{}}\DecValTok{53}\NormalTok{, }\DecValTok{2}\SpecialCharTok{\^{}}\DecValTok{53} \SpecialCharTok{+} \DecValTok{1}\NormalTok{, }\DecValTok{2}\SpecialCharTok{\^{}}\DecValTok{53} \SpecialCharTok{+} \DecValTok{2}\NormalTok{), }\AttributeTok{digits =} \DecValTok{17}\NormalTok{)}
\end{Highlighting}
\end{Shaded}

\begin{verbatim}
[1] "9007199254740992" "9007199254740992" "9007199254740994"
\end{verbatim}

This is why a \texttt{double} can hold every integer only up to
\(2^{53}\approx 9\times10^{15}\), and why counters or identifiers above
that size must be kept as integers or character strings.

\subsection{Overflow and underflow}\label{overflow-and-underflow}

The largest toy number fills every box with 9, \(9.999\times10^{9}\).
The next number, \(10^{10}\), needs an exponent of two digits and cannot
be stored: this is \textbf{overflow}, and the result is recorded as
\texttt{Inf}. At the other end, \(10^{-9}\) is the smallest positive toy
number, and \(10^{-10}\) \textbf{underflows} to \(0\).

+

1

.

0

0

0

×10\^{}

10

\(10^{10}\): overflow, stored as \texttt{Inf}

+

1

.

0

0

0

×10\^{}

-10

\(10^{-10}\): underflow, stored as \texttt{0}

For doubles the largest value is about \(1.8\times10^{308}\) (just below
\(2^{1024}\)). The smallest positive value with full 52-bit precision is
\(2^{-1022}\approx 2.2\times10^{-308}\); below that, the IEEE standard
allows ``subnormal'' numbers with progressively fewer significant bits
down to \(2^{-1074}\approx 4.9\times10^{-324}\), after which the value
is \(0\).

\begin{figure}

\centering{

\pandocbounded{\includegraphics[keepaspectratio]{chapters/unit02_comparith_files/figure-pdf/fig-number-line-1.pdf}}

}

\caption{\label{fig-number-line}Overflow and underflow regions on the
number line for double precision}

\end{figure}%

\begin{Shaded}
\begin{Highlighting}[]
\DocumentationTok{\#\# largest double precision floating{-}point number}
\NormalTok{.Machine}\SpecialCharTok{$}\NormalTok{double.xmax}
\end{Highlighting}
\end{Shaded}

\begin{verbatim}
[1] 1.797693e+308
\end{verbatim}

\begin{Shaded}
\begin{Highlighting}[]
\DecValTok{2}\SpecialCharTok{\^{}}\DecValTok{1023}
\end{Highlighting}
\end{Shaded}

\begin{verbatim}
[1] 8.988466e+307
\end{verbatim}

\begin{Shaded}
\begin{Highlighting}[]
\DecValTok{2}\SpecialCharTok{\^{}}\DecValTok{1024}          \CommentTok{\# overflows to Inf}
\end{Highlighting}
\end{Shaded}

\begin{verbatim}
[1] Inf
\end{verbatim}

\begin{Shaded}
\begin{Highlighting}[]
\DocumentationTok{\#\# smallest double precision floating{-}point numbers}
\NormalTok{.Machine}\SpecialCharTok{$}\NormalTok{double.xmin   }\CommentTok{\# smallest with full precision, 2\^{}{-}1022}
\end{Highlighting}
\end{Shaded}

\begin{verbatim}
[1] 2.225074e-308
\end{verbatim}

\begin{Shaded}
\begin{Highlighting}[]
\DecValTok{2}\SpecialCharTok{\^{}}\NormalTok{(}\SpecialCharTok{{-}}\DecValTok{1074}\NormalTok{)              }\CommentTok{\# smallest subnormal}
\end{Highlighting}
\end{Shaded}

\begin{verbatim}
[1] 4.940656e-324
\end{verbatim}

\begin{Shaded}
\begin{Highlighting}[]
\DecValTok{2}\SpecialCharTok{\^{}}\NormalTok{(}\SpecialCharTok{{-}}\DecValTok{1075}\NormalTok{)              }\CommentTok{\# underflows to 0}
\end{Highlighting}
\end{Shaded}

\begin{verbatim}
[1] 0
\end{verbatim}

Overflow is loud: \texttt{Inf} propagates and usually produces
\texttt{NaN} or an error somewhere downstream. Underflow is quiet: a
probability that underflows to \(0\) simply produces
\texttt{log(0)\ =\ -Inf} or a likelihood ratio of \texttt{0/0}, and the
code may run to completion with a wrong answer.

\section{Rounding Error and Methods for Preventing
It}\label{rounding-error-and-methods-for-preventing-it}

\subsection{Roundoff error}\label{roundoff-error}

Compute \(1 + 0.0001\) in the toy system. The exact answer \(1.0001\)
needs a fifth mantissa box:

1

10⁻¹

10⁻²

10⁻³

10⁻⁴

exp

+

1

.

0

0

0

1

×10\^{}

0

The red digit does not fit and is rounded away: the stored result is
exactly \(1.000\)

In general the stored result of \(a+b\) (or \(a\times b\), \(a/b\),
\ldots) is the \textbf{nearest representable number} to the exact
result. Each single operation therefore has relative error at most about
\(\epsilon\), which is harmless on its own. The concern is when these
errors \textbf{accumulate} over many operations or are
\textbf{amplified} by the structure of a formula. Two patterns do the
amplifying:

\begin{itemize}
\item
  adding a \textbf{large + small} number, where the small number is
  partly or wholly rounded away;
\item
  subtracting \textbf{large − large} nearly equal numbers, where the
  leading digits cancel and only the roundoff errors survive
  (\textbf{catastrophic cancellation}).
\end{itemize}

\subsection{Computing sums: avoiding large +
small}\label{computing-sums-avoiding-large-small}

The mantissa has a finite capacity, so adding something extremely small,
like \(2^{-54}\), to \(1\) loses the small term entirely because its
bits fall outside the 52 stored digits.

\begin{figure}

\centering{

\pandocbounded{\includegraphics[keepaspectratio]{chapters/unit02_comparith_files/figure-pdf/fig-large-plus-small-1.pdf}}

}

\caption{\label{fig-large-plus-small}Roundoff error when adding a large
and a small number}

\end{figure}%

Consider \(1 + 0.0001 + 0.0001\) in the toy system. Evaluated left to
right, \((1 + 0.0001) + 0.0001\), the first addition rounds to
\(1.000\), and so does the second: both small terms are lost. Evaluated
as \(1 + (0.0001 + 0.0001) = 1 + 0.0002\), the intermediate sum
\(2.000\times10^{-4}\) is exact, and only the final addition rounds.
With four mantissa digits the final result is still \(1.000\), but the
point is general: adding the small terms among themselves first commits
only \textbf{one} roundoff error, whereas adding them one at a time to
the large number commits one per term, and these accumulate.

The following R demonstration makes the same point at machine precision.
\texttt{f2} attempts to add \(2^{-53}\) (half an epsilon) \(2^{13}\)
times, but because \(1 + 2^{-53} = 1\) at every single step, the final
sum is simply \(1\) instead of \(1 + 2^{-40}\). Grouping the small
numbers and adding them together before adding them to the large number
(\texttt{f6s}) recovers the correct answer.

\begin{Shaded}
\begin{Highlighting}[]
\DocumentationTok{\#\# the target value}
\NormalTok{f0 }\OtherTok{\textless{}{-}} \DecValTok{1} \SpecialCharTok{+} \DecValTok{2}\SpecialCharTok{\^{}}\NormalTok{(}\SpecialCharTok{{-}}\DecValTok{40}\NormalTok{)}
\NormalTok{f0}
\end{Highlighting}
\end{Shaded}

\begin{verbatim}
[1] 1
\end{verbatim}

\begin{Shaded}
\begin{Highlighting}[]
\DocumentationTok{\#\# adding 2\^{}{-}52 (one epsilon) 2\^{}12 times works: each addition changes the sum}
\NormalTok{f1 }\OtherTok{\textless{}{-}} \DecValTok{1}\NormalTok{; i }\OtherTok{\textless{}{-}} \DecValTok{1}
\ControlFlowTok{while}\NormalTok{ (i }\SpecialCharTok{\textless{}=} \DecValTok{2}\SpecialCharTok{\^{}}\DecValTok{12}\NormalTok{) \{ f1 }\OtherTok{\textless{}{-}}\NormalTok{ f1 }\SpecialCharTok{+} \DecValTok{2}\SpecialCharTok{\^{}}\NormalTok{(}\SpecialCharTok{{-}}\DecValTok{52}\NormalTok{); i }\OtherTok{\textless{}{-}}\NormalTok{ i }\SpecialCharTok{+} \DecValTok{1}\NormalTok{ \}}
\NormalTok{f1}
\end{Highlighting}
\end{Shaded}

\begin{verbatim}
[1] 1
\end{verbatim}

\begin{Shaded}
\begin{Highlighting}[]
\DocumentationTok{\#\# adding 2\^{}{-}53 (half an epsilon) 2\^{}13 times fails: 1 + 2\^{}{-}53 rounds to 1 every time}
\NormalTok{f2 }\OtherTok{\textless{}{-}} \DecValTok{1}\NormalTok{; i }\OtherTok{\textless{}{-}} \DecValTok{1}
\ControlFlowTok{while}\NormalTok{ (i }\SpecialCharTok{\textless{}=} \DecValTok{2}\SpecialCharTok{\^{}}\DecValTok{13}\NormalTok{) \{ f2 }\OtherTok{\textless{}{-}}\NormalTok{ f2 }\SpecialCharTok{+} \DecValTok{2}\SpecialCharTok{\^{}}\NormalTok{(}\SpecialCharTok{{-}}\DecValTok{53}\NormalTok{); i }\OtherTok{\textless{}{-}}\NormalTok{ i }\SpecialCharTok{+} \DecValTok{1}\NormalTok{ \}}
\NormalTok{f2}
\end{Highlighting}
\end{Shaded}

\begin{verbatim}
[1] 1
\end{verbatim}

\begin{Shaded}
\begin{Highlighting}[]
\DocumentationTok{\#\# this difference is not due to number display in R}
\NormalTok{f0 }\SpecialCharTok{==}\NormalTok{ f1}
\end{Highlighting}
\end{Shaded}

\begin{verbatim}
[1] TRUE
\end{verbatim}

\begin{Shaded}
\begin{Highlighting}[]
\NormalTok{f0 }\SpecialCharTok{==}\NormalTok{ f2}
\end{Highlighting}
\end{Shaded}

\begin{verbatim}
[1] FALSE
\end{verbatim}

\begin{Shaded}
\begin{Highlighting}[]
\DocumentationTok{\#\# R\textquotesingle{}s built{-}in sum() uses pairwise summation and extended precision}
\NormalTok{f3 }\OtherTok{\textless{}{-}} \FunctionTok{sum}\NormalTok{(}\FunctionTok{c}\NormalTok{(}\DecValTok{1}\NormalTok{, }\FunctionTok{rep}\NormalTok{(}\DecValTok{2}\SpecialCharTok{\^{}}\NormalTok{(}\SpecialCharTok{{-}}\DecValTok{53}\NormalTok{), }\DecValTok{2}\SpecialCharTok{\^{}}\DecValTok{13}\NormalTok{))); f3}
\end{Highlighting}
\end{Shaded}

\begin{verbatim}
[1] 1
\end{verbatim}

\begin{Shaded}
\begin{Highlighting}[]
\NormalTok{f1 }\SpecialCharTok{==}\NormalTok{ f3}
\end{Highlighting}
\end{Shaded}

\begin{verbatim}
[1] FALSE
\end{verbatim}

\begin{Shaded}
\begin{Highlighting}[]
\NormalTok{f4 }\OtherTok{\textless{}{-}} \FunctionTok{sum}\NormalTok{(}\FunctionTok{c}\NormalTok{(}\DecValTok{1}\NormalTok{, }\FunctionTok{rep}\NormalTok{(}\DecValTok{2}\SpecialCharTok{\^{}}\NormalTok{(}\SpecialCharTok{{-}}\DecValTok{54}\NormalTok{), }\DecValTok{2}\SpecialCharTok{\^{}}\DecValTok{14}\NormalTok{))); f4}
\end{Highlighting}
\end{Shaded}

\begin{verbatim}
[1] 1
\end{verbatim}

\begin{Shaded}
\begin{Highlighting}[]
\NormalTok{f0 }\SpecialCharTok{==}\NormalTok{ f4}
\end{Highlighting}
\end{Shaded}

\begin{verbatim}
[1] FALSE
\end{verbatim}

\begin{Shaded}
\begin{Highlighting}[]
\DocumentationTok{\#\# the fix: accumulate the small numbers first, then add to 1}
\NormalTok{f6s }\OtherTok{\textless{}{-}} \DecValTok{0}\NormalTok{; i }\OtherTok{\textless{}{-}} \DecValTok{1}
\ControlFlowTok{while}\NormalTok{ (i }\SpecialCharTok{\textless{}=} \DecValTok{2}\SpecialCharTok{\^{}}\DecValTok{14}\NormalTok{) \{ f6s }\OtherTok{\textless{}{-}}\NormalTok{ f6s }\SpecialCharTok{+} \DecValTok{2}\SpecialCharTok{\^{}}\NormalTok{(}\SpecialCharTok{{-}}\DecValTok{54}\NormalTok{); i }\OtherTok{\textless{}{-}}\NormalTok{ i }\SpecialCharTok{+} \DecValTok{1}\NormalTok{ \}}
\NormalTok{f6s}
\end{Highlighting}
\end{Shaded}

\begin{verbatim}
[1] 9.094947e-13
\end{verbatim}

\begin{Shaded}
\begin{Highlighting}[]
\NormalTok{f6 }\OtherTok{\textless{}{-}}\NormalTok{ f6s }\SpecialCharTok{+} \DecValTok{1}\NormalTok{; f6}
\end{Highlighting}
\end{Shaded}

\begin{verbatim}
[1] 1
\end{verbatim}

\begin{Shaded}
\begin{Highlighting}[]
\NormalTok{f0 }\SpecialCharTok{==}\NormalTok{ f6}
\end{Highlighting}
\end{Shaded}

\begin{verbatim}
[1] TRUE
\end{verbatim}

\subsubsection*{Rules for accurate
summation}\label{rules-for-accurate-summation}
\addcontentsline{toc}{subsubsection}{Rules for accurate summation}

\begin{itemize}
\item
  Add numbers of similar magnitude together; when summing terms of very
  different sizes, add the smallest first (sort ascending).
\item
  \textbf{Pairwise summation}: sum the two halves of the vector
  recursively, so that partial sums stay of similar size. R's
  \texttt{sum()} uses this together with extended-precision
  accumulation, which is why \texttt{sum()} beat the naive loop above.
\item
  \textbf{Kahan compensated summation}: keep track of the part rounded
  away at each step and add it back at the next step.
\end{itemize}

\begin{Shaded}
\begin{Highlighting}[]
\NormalTok{x }\OtherTok{\textless{}{-}} \FunctionTok{c}\NormalTok{(}\DecValTok{1}\NormalTok{, }\FunctionTok{rep}\NormalTok{(}\FloatTok{1e{-}16}\NormalTok{, }\FloatTok{1e4}\NormalTok{))                 }\CommentTok{\# exact sum is 1 + 1e{-}12}

\NormalTok{naive\_sum }\OtherTok{\textless{}{-}} \ControlFlowTok{function}\NormalTok{(x) \{ s }\OtherTok{\textless{}{-}} \DecValTok{0}\NormalTok{; }\ControlFlowTok{for}\NormalTok{ (v }\ControlFlowTok{in}\NormalTok{ x) s }\OtherTok{\textless{}{-}}\NormalTok{ s }\SpecialCharTok{+}\NormalTok{ v; s \}}

\NormalTok{kahan\_sum }\OtherTok{\textless{}{-}} \ControlFlowTok{function}\NormalTok{(x) \{}
\NormalTok{  s }\OtherTok{\textless{}{-}} \DecValTok{0}\NormalTok{; c }\OtherTok{\textless{}{-}} \DecValTok{0}                            \CommentTok{\# c holds the running roundoff error}
  \ControlFlowTok{for}\NormalTok{ (v }\ControlFlowTok{in}\NormalTok{ x) \{ y }\OtherTok{\textless{}{-}}\NormalTok{ v }\SpecialCharTok{{-}}\NormalTok{ c; t }\OtherTok{\textless{}{-}}\NormalTok{ s }\SpecialCharTok{+}\NormalTok{ y; c }\OtherTok{\textless{}{-}}\NormalTok{ (t }\SpecialCharTok{{-}}\NormalTok{ s) }\SpecialCharTok{{-}}\NormalTok{ y; s }\OtherTok{\textless{}{-}}\NormalTok{ t \}}
\NormalTok{  s \}}

\FunctionTok{c}\NormalTok{(}\AttributeTok{naive   =} \FunctionTok{naive\_sum}\NormalTok{(x) }\SpecialCharTok{{-}} \DecValTok{1}\NormalTok{,}
  \AttributeTok{sorted  =} \FunctionTok{naive\_sum}\NormalTok{(}\FunctionTok{sort}\NormalTok{(x)) }\SpecialCharTok{{-}} \DecValTok{1}\NormalTok{,}
  \AttributeTok{kahan   =} \FunctionTok{kahan\_sum}\NormalTok{(x) }\SpecialCharTok{{-}} \DecValTok{1}\NormalTok{,}
  \AttributeTok{builtin =} \FunctionTok{sum}\NormalTok{(x) }\SpecialCharTok{{-}} \DecValTok{1}\NormalTok{,}
  \AttributeTok{exact   =} \FloatTok{1e{-}12}\NormalTok{)}
\end{Highlighting}
\end{Shaded}

\begin{verbatim}
       naive       sorted        kahan      builtin        exact 
0.000000e+00 1.000089e-12 1.000089e-12 0.000000e+00 1.000000e-12 
\end{verbatim}

The naive loop loses every one of the \(10^4\) small terms; sorting
recovers most of the answer, and Kahan summation and R's \texttt{sum()}
recover it to full precision.

We can use \texttt{log()} to translate the 52-bit limit into other
bases, which tells us at roughly what decimal place precision is lost
(about 15 to 16 significant decimal digits).

\begin{Shaded}
\begin{Highlighting}[]
\FunctionTok{log}\NormalTok{(}\DecValTok{2}\SpecialCharTok{\^{}}\NormalTok{(}\SpecialCharTok{{-}}\DecValTok{52}\NormalTok{))               }\CommentTok{\# base e: about {-}36}
\end{Highlighting}
\end{Shaded}

\begin{verbatim}
[1] -36.04365
\end{verbatim}

\begin{Shaded}
\begin{Highlighting}[]
\DecValTok{1} \SpecialCharTok{+} \FunctionTok{exp}\NormalTok{(}\SpecialCharTok{{-}}\DecValTok{35}\NormalTok{)}
\end{Highlighting}
\end{Shaded}

\begin{verbatim}
[1] 1
\end{verbatim}

\begin{Shaded}
\begin{Highlighting}[]
\DecValTok{1} \SpecialCharTok{+} \FunctionTok{exp}\NormalTok{(}\SpecialCharTok{{-}}\DecValTok{37}\NormalTok{)               }\CommentTok{\# evaluates to exactly 1}
\end{Highlighting}
\end{Shaded}

\begin{verbatim}
[1] 1
\end{verbatim}

\begin{Shaded}
\begin{Highlighting}[]
\FunctionTok{log}\NormalTok{(}\DecValTok{2}\SpecialCharTok{\^{}}\NormalTok{(}\SpecialCharTok{{-}}\DecValTok{52}\NormalTok{), }\AttributeTok{base =} \DecValTok{10}\NormalTok{)    }\CommentTok{\# base 10: about {-}15.7}
\end{Highlighting}
\end{Shaded}

\begin{verbatim}
[1] -15.65356
\end{verbatim}

\begin{Shaded}
\begin{Highlighting}[]
\DecValTok{1} \SpecialCharTok{+} \FloatTok{1e{-}15} \SpecialCharTok{==} \DecValTok{1}
\end{Highlighting}
\end{Shaded}

\begin{verbatim}
[1] FALSE
\end{verbatim}

\begin{Shaded}
\begin{Highlighting}[]
\DecValTok{1} \SpecialCharTok{+} \FloatTok{1e{-}16} \SpecialCharTok{==} \DecValTok{1}
\end{Highlighting}
\end{Shaded}

\begin{verbatim}
[1] TRUE
\end{verbatim}

\subsection{\texorpdfstring{Computing \(e^x\): avoiding large −
large}{Computing e\^{}x: avoiding large − large}}\label{computing-ex-avoiding-large-large}

The Taylor series, \[
e^x \approx \sum_{k=0}^{K}\frac{x^k}{k!} = 1 + x + \frac{x^2}{2!} + \frac{x^3}{3!} + \cdots,
\] works well for \(x > 0\) but fails for large negative \(x\). For
\(x = -20\) the terms alternate in sign and are huge relative to the
answer \(e^{-20}\approx 2\times10^{-9}\):

{\def\LTcaptype{none} % do not increment counter
\begin{longtable}[]{@{}lcccccc@{}}
\toprule\noalign{}
\(k\) & 0 & 1 & 2 & \ldots{} & 20 & \ldots{} \\
\midrule\noalign{}
\endhead
\bottomrule\noalign{}
\endlastfoot
\(x^k/k!\) & \(1\) & \(-20\) & \(200\) & \ldots{} & \(4.3\times10^{7}\)
& \ldots{} \\
\end{longtable}
}

Each term carries a roundoff error of relative size
\(\epsilon \approx 10^{-16}\), i.e.~absolute size
\(\approx 4\times10^{-9}\) for the largest terms, which is \textbf{as
large as the answer itself}. Subtracting nearly equal large numbers
leaves only their errors: catastrophic cancellation.

The function \texttt{fexp} below implements the series with a stopping
rule based on the relative size of the current term. With
\texttt{debug\ =\ TRUE} it returns the internal steps rather than the
final answer.

\begin{Shaded}
\begin{Highlighting}[]
\FunctionTok{library}\NormalTok{(knitr)}

\NormalTok{fexp }\OtherTok{\textless{}{-}} \ControlFlowTok{function}\NormalTok{(x, }\AttributeTok{debug =} \ConstantTok{FALSE}\NormalTok{) \{}
\NormalTok{  i }\OtherTok{\textless{}{-}} \DecValTok{0}\NormalTok{; expx }\OtherTok{\textless{}{-}} \DecValTok{1}\NormalTok{; u }\OtherTok{\textless{}{-}} \DecValTok{1}
  \ControlFlowTok{if}\NormalTok{ (debug) debug\_list }\OtherTok{\textless{}{-}} \FunctionTok{list}\NormalTok{()}
  \ControlFlowTok{while}\NormalTok{ (}\FunctionTok{abs}\NormalTok{(u) }\SpecialCharTok{\textgreater{}} \FloatTok{1e{-}20} \SpecialCharTok{*} \FunctionTok{abs}\NormalTok{(expx)) \{}
\NormalTok{    i }\OtherTok{\textless{}{-}}\NormalTok{ i }\SpecialCharTok{+} \DecValTok{1}
\NormalTok{    u }\OtherTok{\textless{}{-}}\NormalTok{ u }\SpecialCharTok{*}\NormalTok{ x }\SpecialCharTok{/}\NormalTok{ i}
\NormalTok{    expx }\OtherTok{\textless{}{-}}\NormalTok{ expx }\SpecialCharTok{+}\NormalTok{ u}
    \ControlFlowTok{if}\NormalTok{ (debug) debug\_list[[i]] }\OtherTok{\textless{}{-}} \FunctionTok{c}\NormalTok{(}\AttributeTok{Step =}\NormalTok{ i, }\AttributeTok{Term\_u =}\NormalTok{ u, }\AttributeTok{exp\_x =}\NormalTok{ expx)}
\NormalTok{  \}}
  \ControlFlowTok{if}\NormalTok{ (debug) }\FunctionTok{return}\NormalTok{(}\FunctionTok{do.call}\NormalTok{(rbind, debug\_list))}
\NormalTok{  expx}
\NormalTok{\}}

\NormalTok{test\_vals }\OtherTok{\textless{}{-}} \FunctionTok{c}\NormalTok{(}\DecValTok{10}\NormalTok{, }\DecValTok{20}\NormalTok{, }\DecValTok{60}\NormalTok{, }\SpecialCharTok{{-}}\DecValTok{1}\NormalTok{, }\SpecialCharTok{{-}}\DecValTok{10}\NormalTok{, }\SpecialCharTok{{-}}\DecValTok{20}\NormalTok{, }\SpecialCharTok{{-}}\DecValTok{50}\NormalTok{)}
\NormalTok{res1 }\OtherTok{\textless{}{-}} \FunctionTok{data.frame}\NormalTok{(}\AttributeTok{x =}\NormalTok{ test\_vals, }\AttributeTok{Built\_in\_Exp =} \FunctionTok{exp}\NormalTok{(test\_vals), }\AttributeTok{Taylor\_fexp =} \FunctionTok{sapply}\NormalTok{(test\_vals, fexp))}
\FunctionTok{kable}\NormalTok{(res1, }\AttributeTok{caption =} \StringTok{"Built{-}in exp() versus the naive Taylor series fexp()"}\NormalTok{)}
\end{Highlighting}
\end{Shaded}

\begin{longtable}[]{@{}rrr@{}}
\caption{Built-in exp() versus the naive Taylor series
fexp()}\tabularnewline
\toprule\noalign{}
x & Built\_in\_Exp & Taylor\_fexp \\
\midrule\noalign{}
\endfirsthead
\toprule\noalign{}
x & Built\_in\_Exp & Taylor\_fexp \\
\midrule\noalign{}
\endhead
\bottomrule\noalign{}
\endlastfoot
10 & 2.202647e+04 & 2.202647e+04 \\
20 & 4.851652e+08 & 4.851652e+08 \\
60 & 1.142007e+26 & 1.142007e+26 \\
-1 & 3.678794e-01 & 3.678794e-01 \\
-10 & 4.540000e-05 & 4.540000e-05 \\
-20 & 0.000000e+00 & 0.000000e+00 \\
-50 & 0.000000e+00 & 1.107293e+04 \\
\end{longtable}

For positive \(x\) the series is accurate to all printed digits. For
\(x = -10\) the answer is already off in the fourth digit, for
\(x = -20\) it is wrong in the first digit, and \texttt{fexp(-50)}
yields a completely wrong, negative number. Running it in debug mode
shows the massive terms cancelling each other out.

\begin{Shaded}
\begin{Highlighting}[]
\NormalTok{debug\_matrix }\OtherTok{\textless{}{-}} \FunctionTok{fexp}\NormalTok{(}\SpecialCharTok{{-}}\DecValTok{50}\NormalTok{, }\AttributeTok{debug =} \ConstantTok{TRUE}\NormalTok{)}
\NormalTok{debug\_summary }\OtherTok{\textless{}{-}} \FunctionTok{rbind}\NormalTok{(}
  \FunctionTok{head}\NormalTok{(debug\_matrix, }\DecValTok{10}\NormalTok{),}
  \FunctionTok{matrix}\NormalTok{(}\ConstantTok{NA}\NormalTok{, }\AttributeTok{nrow =} \DecValTok{1}\NormalTok{, }\AttributeTok{ncol =} \DecValTok{3}\NormalTok{, }\AttributeTok{dimnames =} \FunctionTok{list}\NormalTok{(}\StringTok{"..."}\NormalTok{, }\FunctionTok{colnames}\NormalTok{(debug\_matrix))),}
  \FunctionTok{tail}\NormalTok{(debug\_matrix, }\DecValTok{10}\NormalTok{))}
\FunctionTok{kable}\NormalTok{(debug\_summary, }\AttributeTok{caption =} \StringTok{"Debug trace for fexp({-}50) showing catastrophic cancellation"}\NormalTok{)}
\end{Highlighting}
\end{Shaded}

\begin{longtable}[]{@{}lrrr@{}}
\caption{Debug trace for fexp(-50) showing catastrophic
cancellation}\tabularnewline
\toprule\noalign{}
& Step & Term\_u & exp\_x \\
\midrule\noalign{}
\endfirsthead
\toprule\noalign{}
& Step & Term\_u & exp\_x \\
\midrule\noalign{}
\endhead
\bottomrule\noalign{}
\endlastfoot
& 1 & -5.000000e+01 & -4.900000e+01 \\
& 2 & 1.250000e+03 & 1.201000e+03 \\
& 3 & -2.083333e+04 & -1.963233e+04 \\
& 4 & 2.604167e+05 & 2.407843e+05 \\
& 5 & -2.604167e+06 & -2.363382e+06 \\
& 6 & 2.170139e+07 & 1.933801e+07 \\
& 7 & -1.550099e+08 & -1.356719e+08 \\
& 8 & 9.688120e+08 & 8.331401e+08 \\
& 9 & -5.382289e+09 & -4.549149e+09 \\
& 10 & 2.691144e+10 & 2.236230e+10 \\
\ldots{} & NA & NA & NA \\
{[}158,{]} & 158 & 0.000000e+00 & 1.107293e+04 \\
{[}159,{]} & 159 & 0.000000e+00 & 1.107293e+04 \\
{[}160,{]} & 160 & 0.000000e+00 & 1.107293e+04 \\
{[}161,{]} & 161 & 0.000000e+00 & 1.107293e+04 \\
{[}162,{]} & 162 & 0.000000e+00 & 1.107293e+04 \\
{[}163,{]} & 163 & 0.000000e+00 & 1.107293e+04 \\
{[}164,{]} & 164 & 0.000000e+00 & 1.107293e+04 \\
{[}165,{]} & 165 & 0.000000e+00 & 1.107293e+04 \\
{[}166,{]} & 166 & 0.000000e+00 & 1.107293e+04 \\
{[}167,{]} & 167 & 0.000000e+00 & 1.107293e+04 \\
\end{longtable}

The largest terms reach about \(10^{20}\); their roundoff errors are of
order \(10^{4}\), and the true answer \(e^{-50}\approx 2\times10^{-22}\)
is invisible beneath them.

\textbf{Fix by reformulating.} For \(x>0\) all terms are positive, so no
cancellation occurs. Therefore compute \(e^{|x|}\) and invert the result
if \(x\) was negative: \[
e^{-x} = \frac{1}{e^{x}}, \qquad x > 0 .
\]

\begin{Shaded}
\begin{Highlighting}[]
\NormalTok{fexp2 }\OtherTok{\textless{}{-}} \ControlFlowTok{function}\NormalTok{(x) \{}
\NormalTok{  xa }\OtherTok{\textless{}{-}} \FunctionTok{abs}\NormalTok{(x); i }\OtherTok{\textless{}{-}} \DecValTok{0}\NormalTok{; expx }\OtherTok{\textless{}{-}} \DecValTok{1}\NormalTok{; u }\OtherTok{\textless{}{-}} \DecValTok{1}
  \ControlFlowTok{while}\NormalTok{ (u }\SpecialCharTok{\textgreater{}} \FloatTok{1e{-}20} \SpecialCharTok{*}\NormalTok{ expx) \{ i }\OtherTok{\textless{}{-}}\NormalTok{ i }\SpecialCharTok{+} \DecValTok{1}\NormalTok{; u }\OtherTok{\textless{}{-}}\NormalTok{ u }\SpecialCharTok{*}\NormalTok{ xa }\SpecialCharTok{/}\NormalTok{ i; expx }\OtherTok{\textless{}{-}}\NormalTok{ expx }\SpecialCharTok{+}\NormalTok{ u \}}
  \ControlFlowTok{if}\NormalTok{ (x }\SpecialCharTok{\textgreater{}=} \DecValTok{0}\NormalTok{) expx }\ControlFlowTok{else} \DecValTok{1} \SpecialCharTok{/}\NormalTok{ expx}
\NormalTok{\}}

\NormalTok{test\_vals\_neg }\OtherTok{\textless{}{-}} \FunctionTok{c}\NormalTok{(}\SpecialCharTok{{-}}\DecValTok{10}\NormalTok{, }\SpecialCharTok{{-}}\DecValTok{20}\NormalTok{, }\SpecialCharTok{{-}}\DecValTok{50}\NormalTok{)}
\NormalTok{res2 }\OtherTok{\textless{}{-}} \FunctionTok{data.frame}\NormalTok{(}\AttributeTok{x =}\NormalTok{ test\_vals\_neg, }\AttributeTok{Built\_in\_Exp =} \FunctionTok{exp}\NormalTok{(test\_vals\_neg), }\AttributeTok{Robust\_fexp2 =} \FunctionTok{sapply}\NormalTok{(test\_vals\_neg, fexp2))}
\FunctionTok{kable}\NormalTok{(res2, }\AttributeTok{caption =} \StringTok{"Robust fexp2() for negative arguments"}\NormalTok{)}
\end{Highlighting}
\end{Shaded}

\begin{longtable}[]{@{}rrr@{}}
\caption{Robust fexp2() for negative arguments}\tabularnewline
\toprule\noalign{}
x & Built\_in\_Exp & Robust\_fexp2 \\
\midrule\noalign{}
\endfirsthead
\toprule\noalign{}
x & Built\_in\_Exp & Robust\_fexp2 \\
\midrule\noalign{}
\endhead
\bottomrule\noalign{}
\endlastfoot
-10 & 4.54e-05 & 4.54e-05 \\
-20 & 0.00e+00 & 0.00e+00 \\
-50 & 0.00e+00 & 0.00e+00 \\
\end{longtable}

The general principle: rearrange a formula so that quantities of similar
magnitude and \textbf{the same sign} are combined, and avoid computing a
small result as the difference of two large numbers.

\textbf{A second example of the same principle.} The function
\(f(x) = e^x - 1\) is needed in interest-rate and survival calculations
for small \(x\). Computing \texttt{exp(x)\ -\ 1} directly subtracts two
numbers near 1, and for \(x = 10^{-10}\) loses six of the sixteen
digits; \texttt{expm1(x)} uses a series for small \(x\) and is exact to
machine precision. The same holds for \(\log(1+x)\) and
\texttt{log1p(x)}.

\begin{Shaded}
\begin{Highlighting}[]
\NormalTok{x }\OtherTok{\textless{}{-}} \FloatTok{1e{-}10}
\FunctionTok{c}\NormalTok{(}\AttributeTok{naive =} \FunctionTok{exp}\NormalTok{(x) }\SpecialCharTok{{-}} \DecValTok{1}\NormalTok{, }\AttributeTok{expm1 =} \FunctionTok{expm1}\NormalTok{(x), }\AttributeTok{series =}\NormalTok{ x }\SpecialCharTok{+}\NormalTok{ x}\SpecialCharTok{\^{}}\DecValTok{2} \SpecialCharTok{/} \DecValTok{2}\NormalTok{)}
\end{Highlighting}
\end{Shaded}

\begin{verbatim}
 naive  expm1 series 
 1e-10  1e-10  1e-10 
\end{verbatim}

\begin{Shaded}
\begin{Highlighting}[]
\NormalTok{x }\OtherTok{\textless{}{-}} \FloatTok{1e{-}20}
\FunctionTok{c}\NormalTok{(}\AttributeTok{naive =} \FunctionTok{log}\NormalTok{(}\DecValTok{1} \SpecialCharTok{+}\NormalTok{ x), }\AttributeTok{log1p =} \FunctionTok{log1p}\NormalTok{(x))}
\end{Highlighting}
\end{Shaded}

\begin{verbatim}
naive log1p 
0e+00 1e-20 
\end{verbatim}

\subsection{Computing variance: avoiding large −
large}\label{computing-variance-avoiding-large-large}

Two algebraically identical formulas for the sum of squared deviations:
\[
\underbrace{\sum_{i=1}^{n}(x_i - \overline{x})^2}_{\text{two-pass}}
\;=\;
\underbrace{\sum_{i=1}^{n} x_i^2 - \frac{\left(\sum_i x_i\right)^2}{n}}_{\text{one-pass}}.
\]

The one-pass formula needs only \(\sum x_i\) and \(\sum x_i^2\),
computed in a single sweep of the data, which made it attractive for
hand calculators and for streaming data. But when the \(x_i\) are large
and close together, \(\sum x_i^2\) and \((\sum x_i)^2/n\) are both huge
and nearly equal, and the subtraction is large − large.

Take \(x = (1, 2, 3)\): the deviations are \((-1, 0, 1)\), so
\(\sum(x_i-\overline x)^2 = 2\). Now shift by \(M = 10^{8}\), giving
\(x = (M+1, M+2, M+3)\). The deviations and the answer are unchanged,
but the one-pass formula computes \[
\sum x_i^2 \approx 3\times10^{16}, \qquad \frac{(\sum x_i)^2}{n} \approx 3\times10^{16}, \qquad \text{difference} = 2.
\] With 16 significant digits, \(3\times10^{16}\) is stored with a gap
of about \(4\) between neighbouring representable numbers, so the true
answer \(2\) lives entirely in the lost digits.

16 mantissa digits

lost

3

.

0

0

0

0

0

0

1

2

0

0

0

0

0

1

4

×10\^{}16

\begin{Shaded}
\begin{Highlighting}[]
\DocumentationTok{\#\# one{-}pass (computational) formula}
\NormalTok{var1 }\OtherTok{\textless{}{-}} \ControlFlowTok{function}\NormalTok{(x) \{ n }\OtherTok{\textless{}{-}} \FunctionTok{length}\NormalTok{(x); (}\FunctionTok{sum}\NormalTok{(x}\SpecialCharTok{\^{}}\DecValTok{2}\NormalTok{) }\SpecialCharTok{{-}} \FunctionTok{sum}\NormalTok{(x)}\SpecialCharTok{\^{}}\DecValTok{2} \SpecialCharTok{/}\NormalTok{ n) }\SpecialCharTok{/}\NormalTok{ (n }\SpecialCharTok{{-}} \DecValTok{1}\NormalTok{) \}}

\DocumentationTok{\#\# two{-}pass formula (centering on the mean)}
\NormalTok{var2 }\OtherTok{\textless{}{-}} \ControlFlowTok{function}\NormalTok{(x) \{ n }\OtherTok{\textless{}{-}} \FunctionTok{length}\NormalTok{(x); }\FunctionTok{sum}\NormalTok{((x }\SpecialCharTok{{-}} \FunctionTok{mean}\NormalTok{(x))}\SpecialCharTok{\^{}}\DecValTok{2}\NormalTok{) }\SpecialCharTok{/}\NormalTok{ (n }\SpecialCharTok{{-}} \DecValTok{1}\NormalTok{) \}}

\DocumentationTok{\#\# small data works fine for both}
\NormalTok{x }\OtherTok{\textless{}{-}} \FunctionTok{c}\NormalTok{(}\DecValTok{1}\NormalTok{, }\DecValTok{2}\NormalTok{, }\DecValTok{3}\NormalTok{)}
\FunctionTok{c}\NormalTok{(}\FunctionTok{var1}\NormalTok{(x), }\FunctionTok{var2}\NormalTok{(x), }\FunctionTok{var}\NormalTok{(x))}
\end{Highlighting}
\end{Shaded}

\begin{verbatim}
[1] 1 1 1
\end{verbatim}

\begin{Shaded}
\begin{Highlighting}[]
\DocumentationTok{\#\# a large shift exposes catastrophic cancellation in var1}
\NormalTok{x }\OtherTok{\textless{}{-}} \FunctionTok{c}\NormalTok{(}\DecValTok{1}\NormalTok{, }\DecValTok{2}\NormalTok{, }\DecValTok{3}\NormalTok{) }\SpecialCharTok{+} \FloatTok{1e8}
\FunctionTok{c}\NormalTok{(}\FunctionTok{var1}\NormalTok{(x), }\FunctionTok{var2}\NormalTok{(x), }\FunctionTok{var}\NormalTok{(x))}
\end{Highlighting}
\end{Shaded}

\begin{verbatim}
[1] 0 1 1
\end{verbatim}

\begin{Shaded}
\begin{Highlighting}[]
\NormalTok{x }\OtherTok{\textless{}{-}} \FunctionTok{c}\NormalTok{(}\DecValTok{1}\NormalTok{, }\DecValTok{2}\NormalTok{, }\DecValTok{3}\NormalTok{) }\SpecialCharTok{+} \FloatTok{1e10}
\FunctionTok{c}\NormalTok{(}\FunctionTok{var1}\NormalTok{(x), }\FunctionTok{var2}\NormalTok{(x), }\FunctionTok{var}\NormalTok{(x))}
\end{Highlighting}
\end{Shaded}

\begin{verbatim}
[1] -32768      1      1
\end{verbatim}

The failure sets in gradually. The error of the one-pass formula is of
order \(\epsilon\,M^2 n\), so it becomes visible once \(M^2\epsilon\) is
comparable to the true variance, i.e.~around \(M\approx 10^{8}\) for a
variance of 1. The following scan shows the one-pass result drifting
away from 1 and eventually becoming negative, which is impossible for a
variance.

\begin{Shaded}
\begin{Highlighting}[]
\NormalTok{M }\OtherTok{\textless{}{-}} \DecValTok{10}\SpecialCharTok{\^{}}\NormalTok{(}\DecValTok{0}\SpecialCharTok{:}\DecValTok{10}\NormalTok{)}
\NormalTok{res }\OtherTok{\textless{}{-}} \FunctionTok{t}\NormalTok{(}\FunctionTok{sapply}\NormalTok{(M, }\ControlFlowTok{function}\NormalTok{(m) \{ x }\OtherTok{\textless{}{-}} \FunctionTok{c}\NormalTok{(}\DecValTok{1}\NormalTok{, }\DecValTok{2}\NormalTok{, }\DecValTok{3}\NormalTok{) }\SpecialCharTok{+}\NormalTok{ m; }\FunctionTok{c}\NormalTok{(}\AttributeTok{M =}\NormalTok{ m, }\AttributeTok{one\_pass =} \FunctionTok{var1}\NormalTok{(x), }\AttributeTok{two\_pass =} \FunctionTok{var2}\NormalTok{(x)) \}))}
\FunctionTok{kable}\NormalTok{(res, }\AttributeTok{caption =} \StringTok{"One{-}pass versus two{-}pass variance as the data are shifted by M"}\NormalTok{, }\AttributeTok{digits =} \DecValTok{6}\NormalTok{)}
\end{Highlighting}
\end{Shaded}

\begin{longtable}[]{@{}rrr@{}}
\caption{One-pass versus two-pass variance as the data are shifted by
M}\tabularnewline
\toprule\noalign{}
M & one\_pass & two\_pass \\
\midrule\noalign{}
\endfirsthead
\toprule\noalign{}
M & one\_pass & two\_pass \\
\midrule\noalign{}
\endhead
\bottomrule\noalign{}
\endlastfoot
1e+00 & 1 & 1 \\
1e+01 & 1 & 1 \\
1e+02 & 1 & 1 \\
1e+03 & 1 & 1 \\
1e+04 & 1 & 1 \\
1e+05 & 1 & 1 \\
1e+06 & 1 & 1 \\
1e+07 & 1 & 1 \\
1e+08 & 0 & 1 \\
1e+09 & 0 & 1 \\
1e+10 & -32768 & 1 \\
\end{longtable}

\subsubsection*{Fixes}\label{fixes}
\addcontentsline{toc}{subsubsection}{Fixes}

\begin{itemize}
\item
  \textbf{Two-pass}: compute \(\overline x\) first, then
  \(\sum (x_i - \overline x)^2\). The deviations are small and of
  similar size, so no cancellation occurs. This is what \texttt{var()}
  does.
\item
  \textbf{Shifted one-pass}: subtract any convenient constant
  \(c \approx \overline x\), for instance \(c = x_1\), before applying
  the one-pass formula. The identity \[
  \sum (x_i-\overline x)^2 = \sum (x_i - c)^2 - \frac{\left(\sum (x_i - c)\right)^2}{n},
  \] holds for every \(c\), and with \(c\) near the mean the two terms
  on the right are no longer huge.
\item
  \textbf{Welford's online algorithm} updates the running mean
  \(\overline x_k\) and running sum of squares
  \(S_k = \sum_{i\le k}(x_i - \overline x_k)^2\) one observation at a
  time, \[
  \overline x_k = \overline x_{k-1} + \frac{x_k - \overline x_{k-1}}{k}, \qquad S_k = S_{k-1} + (x_k - \overline x_{k-1})(x_k - \overline x_k),
  \] and is the method of choice for streaming data, where a two-pass
  computation is impossible.
\end{itemize}

\begin{Shaded}
\begin{Highlighting}[]
\NormalTok{var\_shifted }\OtherTok{\textless{}{-}} \ControlFlowTok{function}\NormalTok{(x) \{ n }\OtherTok{\textless{}{-}} \FunctionTok{length}\NormalTok{(x); d }\OtherTok{\textless{}{-}}\NormalTok{ x }\SpecialCharTok{{-}}\NormalTok{ x[}\DecValTok{1}\NormalTok{]; (}\FunctionTok{sum}\NormalTok{(d}\SpecialCharTok{\^{}}\DecValTok{2}\NormalTok{) }\SpecialCharTok{{-}} \FunctionTok{sum}\NormalTok{(d)}\SpecialCharTok{\^{}}\DecValTok{2} \SpecialCharTok{/}\NormalTok{ n) }\SpecialCharTok{/}\NormalTok{ (n }\SpecialCharTok{{-}} \DecValTok{1}\NormalTok{) \}}

\NormalTok{var\_welford }\OtherTok{\textless{}{-}} \ControlFlowTok{function}\NormalTok{(x) \{}
\NormalTok{  m }\OtherTok{\textless{}{-}} \DecValTok{0}\NormalTok{; S }\OtherTok{\textless{}{-}} \DecValTok{0}
  \ControlFlowTok{for}\NormalTok{ (k }\ControlFlowTok{in} \FunctionTok{seq\_along}\NormalTok{(x)) \{ m\_old }\OtherTok{\textless{}{-}}\NormalTok{ m; m }\OtherTok{\textless{}{-}}\NormalTok{ m }\SpecialCharTok{+}\NormalTok{ (x[k] }\SpecialCharTok{{-}}\NormalTok{ m) }\SpecialCharTok{/}\NormalTok{ k; S }\OtherTok{\textless{}{-}}\NormalTok{ S }\SpecialCharTok{+}\NormalTok{ (x[k] }\SpecialCharTok{{-}}\NormalTok{ m\_old) }\SpecialCharTok{*}\NormalTok{ (x[k] }\SpecialCharTok{{-}}\NormalTok{ m) \}}
\NormalTok{  S }\SpecialCharTok{/}\NormalTok{ (}\FunctionTok{length}\NormalTok{(x) }\SpecialCharTok{{-}} \DecValTok{1}\NormalTok{) \}}

\NormalTok{x }\OtherTok{\textless{}{-}} \FunctionTok{c}\NormalTok{(}\DecValTok{1}\NormalTok{, }\DecValTok{2}\NormalTok{, }\DecValTok{3}\NormalTok{) }\SpecialCharTok{+} \FloatTok{1e10}
\FunctionTok{c}\NormalTok{(}\AttributeTok{one\_pass =} \FunctionTok{var1}\NormalTok{(x), }\AttributeTok{shifted =} \FunctionTok{var\_shifted}\NormalTok{(x), }\AttributeTok{welford =} \FunctionTok{var\_welford}\NormalTok{(x), }\AttributeTok{two\_pass =} \FunctionTok{var2}\NormalTok{(x))}
\end{Highlighting}
\end{Shaded}

\begin{verbatim}
one_pass  shifted  welford two_pass 
  -32768        1        1        1 
\end{verbatim}

\section{Methods for Mitigating Overflow and Underflow
Errors}\label{methods-for-mitigating-overflow-and-underflow-errors}

\subsection{Where overflow and underflow arise in
statistics}\label{where-overflow-and-underflow-arise-in-statistics}

Likelihoods are \textbf{products of many probabilities}, and normalizing
constants involve \textbf{factorials} or \textbf{exponentials}: \[
\prod_{i=1}^{2000} p(x_i) = \left(\tfrac12\right)^{2000} \approx 10^{-602} \quad\to\quad \text{underflow to } 0,
\] \[
1000! \approx 4\times10^{2567} \quad\to\quad \text{overflow to } \infty,
\] \[
\text{logistic: } \frac{e^{x}}{1+e^{x}} \text{ at } x = 1000 \quad\to\quad \frac{\infty}{\infty} = \texttt{NaN}.
\]

Two remedies are available: change the algebraic expression so that the
intermediate quantities stay in range, or work on the logarithmic scale.

\subsection{Changing the algebraic
expression}\label{changing-the-algebraic-expression}

Two equivalent forms of the logistic (sigmoid) function: \[
\sigma(x) = \frac{e^{x}}{1+e^{x}} = \frac{1}{1+e^{-x}}.
\]

{\def\LTcaptype{none} % do not increment counter
\begin{longtable}[]{@{}
  >{\centering\arraybackslash}p{(\linewidth - 4\tabcolsep) * \real{0.4000}}
  >{\raggedright\arraybackslash}p{(\linewidth - 4\tabcolsep) * \real{0.3000}}
  >{\raggedright\arraybackslash}p{(\linewidth - 4\tabcolsep) * \real{0.3000}}@{}}
\toprule\noalign{}
\begin{minipage}[b]{\linewidth}\centering
\(x\)
\end{minipage} & \begin{minipage}[b]{\linewidth}\raggedright
\(e^{x}/(1+e^{x})\)
\end{minipage} & \begin{minipage}[b]{\linewidth}\raggedright
\(1/(1+e^{-x})\)
\end{minipage} \\
\midrule\noalign{}
\endhead
\bottomrule\noalign{}
\endlastfoot
\(+1000\) & \(e^{x}=\infty\), \$\infty/\infty = \$ \texttt{NaN} ✗ &
\(e^{-x}=0\), \(1/1 = 1\) ✓ \\
\(-1000\) & \(e^{x}=0\), \(0/1 = 0\) ✓ & \(e^{-x}=\infty\),
\(1/\infty = 0\) ✓ \\
\end{longtable}
}

Evaluating the first form at a large positive value yields
\(\infty/\infty\); the second form lets \(e^{-x}\) underflow harmlessly
to \(0\) and computes \(1/(1+0) = 1\). In this case the second form is
safe at both ends, because a division by \(\infty\) gives \(0\)
correctly. In general one picks the form in which the exponential has a
\textbf{non-positive} argument: \[
\sigma(x) = \begin{cases} 1/(1+e^{-x}) & x \ge 0\\ e^{x}/(1+e^{x}) & x < 0.\end{cases}
\]

\begin{Shaded}
\begin{Highlighting}[]
\NormalTok{p1 }\OtherTok{\textless{}{-}} \ControlFlowTok{function}\NormalTok{(theta) }\FunctionTok{exp}\NormalTok{(theta) }\SpecialCharTok{/}\NormalTok{ (}\DecValTok{1} \SpecialCharTok{+} \FunctionTok{exp}\NormalTok{(theta))}
\NormalTok{p2 }\OtherTok{\textless{}{-}} \ControlFlowTok{function}\NormalTok{(theta) }\DecValTok{1} \SpecialCharTok{/}\NormalTok{ (}\DecValTok{1} \SpecialCharTok{+} \FunctionTok{exp}\NormalTok{(}\SpecialCharTok{{-}}\NormalTok{theta))}

\DocumentationTok{\#\# large positive theta breaks p1, works in p2}
\NormalTok{theta }\OtherTok{\textless{}{-}} \DecValTok{2000}
\FunctionTok{c}\NormalTok{(}\AttributeTok{p1 =} \FunctionTok{p1}\NormalTok{(theta), }\AttributeTok{p2 =} \FunctionTok{p2}\NormalTok{(theta))}
\end{Highlighting}
\end{Shaded}

\begin{verbatim}
 p1  p2 
NaN   1 
\end{verbatim}

\begin{Shaded}
\begin{Highlighting}[]
\DocumentationTok{\#\# large negative theta works for both}
\NormalTok{theta }\OtherTok{\textless{}{-}} \SpecialCharTok{{-}}\DecValTok{2000}
\FunctionTok{c}\NormalTok{(}\AttributeTok{p1 =} \FunctionTok{p1}\NormalTok{(theta), }\AttributeTok{p2 =} \FunctionTok{p2}\NormalTok{(theta))}
\end{Highlighting}
\end{Shaded}

\begin{verbatim}
p1 p2 
 0  0 
\end{verbatim}

\begin{Shaded}
\begin{Highlighting}[]
\DocumentationTok{\#\# the piecewise form is safe everywhere and keeps full precision in both tails}
\NormalTok{sigmoid }\OtherTok{\textless{}{-}} \ControlFlowTok{function}\NormalTok{(x) }\FunctionTok{ifelse}\NormalTok{(x }\SpecialCharTok{\textgreater{}=} \DecValTok{0}\NormalTok{, }\DecValTok{1} \SpecialCharTok{/}\NormalTok{ (}\DecValTok{1} \SpecialCharTok{+} \FunctionTok{exp}\NormalTok{(}\SpecialCharTok{{-}}\NormalTok{x)), }\FunctionTok{exp}\NormalTok{(x) }\SpecialCharTok{/}\NormalTok{ (}\DecValTok{1} \SpecialCharTok{+} \FunctionTok{exp}\NormalTok{(x)))}
\FunctionTok{sigmoid}\NormalTok{(}\FunctionTok{c}\NormalTok{(}\SpecialCharTok{{-}}\DecValTok{2000}\NormalTok{, }\SpecialCharTok{{-}}\DecValTok{20}\NormalTok{, }\DecValTok{0}\NormalTok{, }\DecValTok{20}\NormalTok{, }\DecValTok{2000}\NormalTok{))}
\end{Highlighting}
\end{Shaded}

\begin{verbatim}
[1] 0.000000e+00 2.061154e-09 5.000000e-01 1.000000e+00 1.000000e+00
\end{verbatim}

\textbf{Built-in functions.} Software provides functions written
specifically to avoid cancellation and overflow, and they should be used
in preference to the naive composition:

{\def\LTcaptype{none} % do not increment counter
\begin{longtable}[]{@{}
  >{\raggedright\arraybackslash}p{(\linewidth - 4\tabcolsep) * \real{0.3333}}
  >{\raggedright\arraybackslash}p{(\linewidth - 4\tabcolsep) * \real{0.3333}}
  >{\raggedright\arraybackslash}p{(\linewidth - 4\tabcolsep) * \real{0.3333}}@{}}
\toprule\noalign{}
\begin{minipage}[b]{\linewidth}\raggedright
instead of
\end{minipage} & \begin{minipage}[b]{\linewidth}\raggedright
use
\end{minipage} & \begin{minipage}[b]{\linewidth}\raggedright
reason
\end{minipage} \\
\midrule\noalign{}
\endhead
\bottomrule\noalign{}
\endlastfoot
\texttt{log(1\ +\ x)} & \texttt{log1p(x)} & for tiny \(x\),
\texttt{1\ +\ x} rounds to 1 \\
\texttt{exp(x)\ -\ 1} & \texttt{expm1(x)} & for tiny \(x\),
\texttt{exp(x)} rounds to 1 \\
\texttt{log(factorial(n))} & \texttt{lfactorial(n)},
\texttt{lgamma(n\ +\ 1)} & \texttt{factorial(n)} overflows \\
\texttt{log(choose(n,\ k))} & \texttt{lchoose(n,\ k)} & overflow \\
\texttt{log(dnorm(x))} & \texttt{dnorm(x,\ log\ =\ TRUE)} &
\texttt{dnorm(x)} underflows \\
\texttt{log(pnorm(x))} & \texttt{pnorm(x,\ log.p\ =\ TRUE)} &
underflow \\
\texttt{log(1\ -\ pnorm(x))} &
\texttt{pnorm(x,\ lower.tail\ =\ FALSE,\ log.p\ =\ TRUE)} & cancellation
and underflow \\
\end{longtable}
}

\begin{Shaded}
\begin{Highlighting}[]
\FunctionTok{c}\NormalTok{(}\AttributeTok{naive =} \FunctionTok{log}\NormalTok{(}\FunctionTok{factorial}\NormalTok{(}\DecValTok{1000}\NormalTok{)), }\AttributeTok{lfactorial =} \FunctionTok{lfactorial}\NormalTok{(}\DecValTok{1000}\NormalTok{))}
\end{Highlighting}
\end{Shaded}

\begin{verbatim}
     naive lfactorial 
       Inf   5912.128 
\end{verbatim}

\begin{Shaded}
\begin{Highlighting}[]
\FunctionTok{c}\NormalTok{(}\AttributeTok{naive =} \FunctionTok{log}\NormalTok{(}\FunctionTok{dnorm}\NormalTok{(}\DecValTok{40}\NormalTok{)), }\AttributeTok{dnorm\_log =} \FunctionTok{dnorm}\NormalTok{(}\DecValTok{40}\NormalTok{, }\AttributeTok{log =} \ConstantTok{TRUE}\NormalTok{))}
\end{Highlighting}
\end{Shaded}

\begin{verbatim}
    naive dnorm_log 
     -Inf -800.9189 
\end{verbatim}

\begin{Shaded}
\begin{Highlighting}[]
\FunctionTok{c}\NormalTok{(}\AttributeTok{naive =} \FunctionTok{log}\NormalTok{(}\DecValTok{1} \SpecialCharTok{{-}} \FunctionTok{pnorm}\NormalTok{(}\DecValTok{10}\NormalTok{)), }\AttributeTok{stable =} \FunctionTok{pnorm}\NormalTok{(}\DecValTok{10}\NormalTok{, }\AttributeTok{lower.tail =} \ConstantTok{FALSE}\NormalTok{, }\AttributeTok{log.p =} \ConstantTok{TRUE}\NormalTok{))}
\end{Highlighting}
\end{Shaded}

\begin{verbatim}
    naive    stable 
     -Inf -53.23129 
\end{verbatim}

The last line shows both failures at once: \texttt{1\ -\ pnorm(10)} is
\texttt{1\ -\ 1\ =\ 0} in double precision, so the naive answer is
\texttt{-Inf}, while the true value is about \(-53.2\).

\subsection{Representing numbers in logarithmic
form}\label{representing-numbers-in-logarithmic-form}

Represent a positive number \(x\) by \(\ell_x = \log x\): \[
x = e^{2000} \iff \ell_x = 2000, \qquad x = e^{-2000} \iff \ell_x = -2000 .
\] Numbers that overflow or underflow in ordinary form are ordinary
numbers on the log scale. The cost is that each operation \(f\) must be
\textbf{translated} to the log scale: \[
x_1,\dots,x_n \xrightarrow{\;f\;} y \qquad\Longleftrightarrow\qquad \ell_1,\dots,\ell_n \xrightarrow{\;\tilde f\;} \ell_y = \log f(e^{\ell_1},\dots,e^{\ell_n}).
\]

\textbf{Products and quotients} are the easy case: \[
\log(x_1 x_2) = \ell_1 + \ell_2, \qquad \log\left(\prod_{i=1}^{n} x_i\right) = \sum_{i=1}^{n} \ell_i, \qquad \log\frac{x_1}{x_2} = \ell_1 - \ell_2.
\] In particular \(\log \prod_i p(x_i) = \sum_i \log p(x_i)\): always
compute the log-likelihood, never the likelihood.

\begin{Shaded}
\begin{Highlighting}[]
\NormalTok{log\_x }\OtherTok{\textless{}{-}} \DecValTok{800}\NormalTok{; x }\OtherTok{\textless{}{-}} \FunctionTok{exp}\NormalTok{(log\_x); x          }\CommentTok{\# overflow}
\end{Highlighting}
\end{Shaded}

\begin{verbatim}
[1] Inf
\end{verbatim}

\begin{Shaded}
\begin{Highlighting}[]
\NormalTok{log\_y }\OtherTok{\textless{}{-}} \DecValTok{805}\NormalTok{; y }\OtherTok{\textless{}{-}} \FunctionTok{exp}\NormalTok{(log\_y); y          }\CommentTok{\# overflow}
\end{Highlighting}
\end{Shaded}

\begin{verbatim}
[1] Inf
\end{verbatim}

\begin{Shaded}
\begin{Highlighting}[]
\NormalTok{x }\SpecialCharTok{/}\NormalTok{ y                                     }\CommentTok{\# Inf/Inf = NaN}
\end{Highlighting}
\end{Shaded}

\begin{verbatim}
[1] NaN
\end{verbatim}

\begin{Shaded}
\begin{Highlighting}[]
\NormalTok{log\_xovery }\OtherTok{\textless{}{-}}\NormalTok{ log\_x }\SpecialCharTok{{-}}\NormalTok{ log\_y               }\CommentTok{\# subtracting logs is safe}
\NormalTok{log\_xovery}
\end{Highlighting}
\end{Shaded}

\begin{verbatim}
[1] -5
\end{verbatim}

\begin{Shaded}
\begin{Highlighting}[]
\FunctionTok{exp}\NormalTok{(log\_xovery)}
\end{Highlighting}
\end{Shaded}

\begin{verbatim}
[1] 0.006737947
\end{verbatim}

\begin{Shaded}
\begin{Highlighting}[]
\DocumentationTok{\#\# likelihood versus log{-}likelihood for 2000 fair{-}coin tosses}
\NormalTok{p }\OtherTok{\textless{}{-}} \FunctionTok{rep}\NormalTok{(}\FloatTok{0.5}\NormalTok{, }\DecValTok{2000}\NormalTok{)}
\FunctionTok{c}\NormalTok{(}\AttributeTok{likelihood =} \FunctionTok{prod}\NormalTok{(p), }\AttributeTok{log\_likelihood =} \FunctionTok{sum}\NormalTok{(}\FunctionTok{log}\NormalTok{(p)))}
\end{Highlighting}
\end{Shaded}

\begin{verbatim}
    likelihood log_likelihood 
         0.000      -1386.294 
\end{verbatim}

\subsection{The log-sum-exp trick}\label{the-log-sum-exp-trick}

Sums are the hard case: \(\log(x_1 + x_2) \neq \ell_1 + \ell_2\). We
need \[
\log\left(\sum_{i=1}^{n} x_i\right) = \log\left(\sum_{i=1}^{n} e^{\ell_i}\right),
\] but the \(e^{\ell_i}\) may overflow or underflow. Let
\(m = \max_i \ell_i\) and factor it out: \[
\log\left(\sum_{i=1}^{n} e^{\ell_i}\right)
= \log\left(e^{m}\sum_{i=1}^{n} e^{\ell_i - m}\right)
= m + \log\left(\sum_{i=1}^{n} e^{\ell_i - m}\right).
\]

Every \(\ell_i - m \le 0\), so \(e^{\ell_i - m} \in (0, 1]\) and nothing
overflows; the largest term is \(e^{0} = 1\), so the sum is at least 1
and cannot underflow. Terms far below the maximum underflow to \(0\)
individually, but they would have contributed nothing at 16-digit
precision anyway.

Illustration with \(x_i = e^{2000}, e^{2001}, e^{2002}\), each of which
overflows to \texttt{Inf}:

{\def\LTcaptype{none} % do not increment counter
\begin{longtable}[]{@{}lccc@{}}
\toprule\noalign{}
& \(i=1\) & \(i=2\) & \(i=3\) \\
\midrule\noalign{}
\endhead
\bottomrule\noalign{}
\endlastfoot
\(\ell_i\) & 2000 & 2001 & 2002 \\
\(\ell_i - m\) (\(m = 2002\)) & \(-2\) & \(-1\) & \(0\) \\
\(e^{\ell_i - m}\) & \(0.135\) & \(0.368\) & \(1\) \\
\end{longtable}
}

\[
\log\sum_i x_i = 2002 + \log(0.135 + 0.368 + 1) = 2002.408
\]

\begin{Shaded}
\begin{Highlighting}[]
\NormalTok{log\_sum\_exp }\OtherTok{\textless{}{-}} \ControlFlowTok{function}\NormalTok{(log\_x) \{}
\NormalTok{  max\_log\_x }\OtherTok{\textless{}{-}} \FunctionTok{max}\NormalTok{(log\_x)}
\NormalTok{  max\_log\_x }\SpecialCharTok{+} \FunctionTok{log}\NormalTok{(}\FunctionTok{sum}\NormalTok{(}\FunctionTok{exp}\NormalTok{(log\_x }\SpecialCharTok{{-}}\NormalTok{ max\_log\_x)))}
\NormalTok{\}}

\DocumentationTok{\#\# values that would overflow directly}
\NormalTok{log\_x }\OtherTok{\textless{}{-}} \FunctionTok{c}\NormalTok{(}\DecValTok{2000}\NormalTok{, }\DecValTok{2010}\NormalTok{, }\DecValTok{2030}\NormalTok{)}
\FunctionTok{exp}\NormalTok{(log\_x)                                   }\CommentTok{\# Inf}
\end{Highlighting}
\end{Shaded}

\begin{verbatim}
[1] Inf Inf Inf
\end{verbatim}

\begin{Shaded}
\begin{Highlighting}[]
\FunctionTok{c}\NormalTok{(}\AttributeTok{direct =} \FunctionTok{log}\NormalTok{(}\FunctionTok{sum}\NormalTok{(}\FunctionTok{exp}\NormalTok{(log\_x))), }\AttributeTok{lse =} \FunctionTok{log\_sum\_exp}\NormalTok{(log\_x))}
\end{Highlighting}
\end{Shaded}

\begin{verbatim}
direct    lse 
   Inf   2030 
\end{verbatim}

\begin{Shaded}
\begin{Highlighting}[]
\DocumentationTok{\#\# values that would underflow directly}
\NormalTok{log\_x }\OtherTok{\textless{}{-}} \SpecialCharTok{{-}}\FunctionTok{c}\NormalTok{(}\DecValTok{2000}\NormalTok{, }\DecValTok{2010}\NormalTok{, }\DecValTok{2030}\NormalTok{)}
\FunctionTok{exp}\NormalTok{(log\_x)                                   }\CommentTok{\# 0}
\end{Highlighting}
\end{Shaded}

\begin{verbatim}
[1] 0 0 0
\end{verbatim}

\begin{Shaded}
\begin{Highlighting}[]
\FunctionTok{c}\NormalTok{(}\AttributeTok{direct =} \FunctionTok{log}\NormalTok{(}\FunctionTok{sum}\NormalTok{(}\FunctionTok{exp}\NormalTok{(log\_x))), }\AttributeTok{lse =} \FunctionTok{log\_sum\_exp}\NormalTok{(log\_x))}
\end{Highlighting}
\end{Shaded}

\begin{verbatim}
direct    lse 
  -Inf  -2000 
\end{verbatim}

\textbf{Example: log-likelihood of a mixture.} The density of a
two-component normal mixture is
\(f(y) = \pi_1\varphi(y;\mu_1,1) + \pi_2\varphi(y;\mu_2,1)\), a sum
inside a product. For an observation far from both means, both component
densities underflow and the naive log-likelihood is \texttt{-Inf}; with
log-sum-exp applied observation by observation it is finite and correct.

\begin{Shaded}
\begin{Highlighting}[]
\NormalTok{y }\OtherTok{\textless{}{-}} \FunctionTok{c}\NormalTok{(}\SpecialCharTok{{-}}\DecValTok{1}\NormalTok{, }\FloatTok{0.5}\NormalTok{, }\DecValTok{2}\NormalTok{, }\DecValTok{45}\NormalTok{)                      }\CommentTok{\# the last value is far from both means}
\NormalTok{mu }\OtherTok{\textless{}{-}} \FunctionTok{c}\NormalTok{(}\DecValTok{0}\NormalTok{, }\DecValTok{3}\NormalTok{); w }\OtherTok{\textless{}{-}} \FunctionTok{c}\NormalTok{(}\FloatTok{0.6}\NormalTok{, }\FloatTok{0.4}\NormalTok{)}
\NormalTok{naive }\OtherTok{\textless{}{-}} \FunctionTok{sum}\NormalTok{(}\FunctionTok{log}\NormalTok{(w[}\DecValTok{1}\NormalTok{] }\SpecialCharTok{*} \FunctionTok{dnorm}\NormalTok{(y, mu[}\DecValTok{1}\NormalTok{]) }\SpecialCharTok{+}\NormalTok{ w[}\DecValTok{2}\NormalTok{] }\SpecialCharTok{*} \FunctionTok{dnorm}\NormalTok{(y, mu[}\DecValTok{2}\NormalTok{])))}
\NormalTok{stable }\OtherTok{\textless{}{-}} \FunctionTok{sum}\NormalTok{(}\FunctionTok{sapply}\NormalTok{(y, }\ControlFlowTok{function}\NormalTok{(yi) }\FunctionTok{log\_sum\_exp}\NormalTok{(}\FunctionTok{log}\NormalTok{(w) }\SpecialCharTok{+} \FunctionTok{dnorm}\NormalTok{(yi, mu, }\AttributeTok{log =} \ConstantTok{TRUE}\NormalTok{))))}
\FunctionTok{c}\NormalTok{(}\AttributeTok{naive =}\NormalTok{ naive, }\AttributeTok{stable =}\NormalTok{ stable)}
\end{Highlighting}
\end{Shaded}

\begin{verbatim}
    naive    stable 
     -Inf -889.3333 
\end{verbatim}

\textbf{Example: log-softmax.} The softmax probabilities
\(p_k = e^{x_k}/\sum_j e^{x_j}\) appear in multinomial logistic
regression and neural networks. Their logarithm is
\(\log p_k = x_k - \log\sum_j e^{x_j}\), and the second term is a
log-sum-exp.

\begin{Shaded}
\begin{Highlighting}[]
\NormalTok{x }\OtherTok{\textless{}{-}} \FunctionTok{c}\NormalTok{(}\DecValTok{1000}\NormalTok{, }\DecValTok{1001}\NormalTok{, }\DecValTok{1005}\NormalTok{)}
\NormalTok{naive }\OtherTok{\textless{}{-}} \FunctionTok{log}\NormalTok{(}\FunctionTok{exp}\NormalTok{(x) }\SpecialCharTok{/} \FunctionTok{sum}\NormalTok{(}\FunctionTok{exp}\NormalTok{(x)))}
\NormalTok{stable }\OtherTok{\textless{}{-}}\NormalTok{ x }\SpecialCharTok{{-}} \FunctionTok{log\_sum\_exp}\NormalTok{(x)}
\FunctionTok{rbind}\NormalTok{(naive, stable, }\AttributeTok{probabilities =} \FunctionTok{exp}\NormalTok{(stable))}
\end{Highlighting}
\end{Shaded}

\begin{verbatim}
                      [,1]        [,2]        [,3]
naive                  NaN         NaN         NaN
stable        -5.024744890 -4.02474489 -0.02474489
probabilities  0.006573263  0.01786798  0.97555875
\end{verbatim}

\subsection{Differences on the log
scale}\label{differences-on-the-log-scale}

For \(x > y > 0\) with \(\ell_x = \log x\), \(\ell_y = \log y\): \[
\log(x - y) = \log\left(e^{\ell_x}\left(1 - e^{\ell_y - \ell_x}\right)\right) = \ell_x + \log\left(1 - e^{\ell_y - \ell_x}\right).
\] Since \(\ell_y - \ell_x < 0\), \(e^{\ell_y-\ell_x}\in(0,1)\) and
there is no overflow. When \(x \approx y\), the factor
\(1 - e^{\ell_y-\ell_x}\) is itself a large − large subtraction;
\texttt{log1p(-exp(ly\ -\ lx))} avoids it.

\begin{Shaded}
\begin{Highlighting}[]
\NormalTok{log\_minus\_exp }\OtherTok{\textless{}{-}} \ControlFlowTok{function}\NormalTok{(log\_x, log\_y) \{}
  \ControlFlowTok{if}\NormalTok{ (log\_x }\SpecialCharTok{\textless{}}\NormalTok{ log\_y) }\FunctionTok{stop}\NormalTok{(}\StringTok{"the first argument must be the larger"}\NormalTok{)}
\NormalTok{  log\_x }\SpecialCharTok{+} \FunctionTok{log1p}\NormalTok{(}\SpecialCharTok{{-}}\FunctionTok{exp}\NormalTok{(log\_y }\SpecialCharTok{{-}}\NormalTok{ log\_x))}
\NormalTok{\}}

\FunctionTok{log\_minus\_exp}\NormalTok{(}\DecValTok{2020}\NormalTok{, }\DecValTok{2000}\NormalTok{)}
\end{Highlighting}
\end{Shaded}

\begin{verbatim}
[1] 2020
\end{verbatim}

\begin{Shaded}
\begin{Highlighting}[]
\FunctionTok{exp}\NormalTok{(}\DecValTok{2000}\NormalTok{) }\SpecialCharTok{{-}} \FunctionTok{exp}\NormalTok{(}\DecValTok{1999}\NormalTok{)                       }\CommentTok{\# direct subtraction: Inf {-} Inf = NaN}
\end{Highlighting}
\end{Shaded}

\begin{verbatim}
[1] NaN
\end{verbatim}

\begin{Shaded}
\begin{Highlighting}[]
\DocumentationTok{\#\# nearly equal arguments: log1p keeps the precision that log(1 {-} .) would lose}
\NormalTok{lx }\OtherTok{\textless{}{-}} \DecValTok{10}\NormalTok{; ly }\OtherTok{\textless{}{-}} \DecValTok{10} \SpecialCharTok{{-}} \FloatTok{1e{-}12}
\FunctionTok{c}\NormalTok{(}\AttributeTok{with\_log1p =} \FunctionTok{log\_minus\_exp}\NormalTok{(lx, ly), }\AttributeTok{with\_log =}\NormalTok{ lx }\SpecialCharTok{+} \FunctionTok{log}\NormalTok{(}\DecValTok{1} \SpecialCharTok{{-}} \FunctionTok{exp}\NormalTok{(ly }\SpecialCharTok{{-}}\NormalTok{ lx)), }\AttributeTok{exact =} \FunctionTok{log}\NormalTok{(}\FunctionTok{exp}\NormalTok{(lx) }\SpecialCharTok{{-}} \FunctionTok{exp}\NormalTok{(ly)))}
\end{Highlighting}
\end{Shaded}

\begin{verbatim}
with_log1p   with_log      exact 
 -17.63093  -17.63093  -17.63079 
\end{verbatim}

A statistical use of the difference formula is the survival function of
a discrete distribution,
\(\log P(X > k) = \log\big(1 - P(X\le k)\big)\), when \(P(X\le k)\) is
available only as a log-probability close to \(0\).

\section{Summary}\label{summary}

\begin{itemize}
\item
  Integers occupy a fixed number of bits; exceeding \(2^{31}-1\) gives
  \texttt{NA} in R.
\item
  Doubles store a sign, a 52-bit mantissa (about 16 decimal digits), and
  an 11-bit exponent (magnitudes \(10^{-308}\) to \(10^{308}\)).
  Neighbouring doubles are spaced at a relative distance of about
  \(\epsilon = 2^{-52}\); integers above \(2^{53}\) are not all
  representable.
\item
  Overflow produces \texttt{Inf}, underflow produces \texttt{0}; the
  latter is silent and therefore more dangerous.
\item
  Roundoff error is amplified by adding large + small (fix: add small
  terms first, pairwise or compensated summation) and by subtracting
  large − large (fix: reformulate so that same-sign quantities of
  similar size are combined, as in \(e^{-x} = 1/e^{x}\) and the two-pass
  variance).
\item
  Overflow and underflow are avoided by rewriting expressions so that
  exponentials have non-positive arguments, by using purpose-built
  functions (\texttt{log1p}, \texttt{expm1}, \texttt{lgamma},
  \texttt{dnorm(log\ =\ TRUE)}), and by working on the log scale, where
  products become sums and sums are handled by the log-sum-exp trick.
\end{itemize}

\section{Exercises}\label{exercises}

\textbf{Question 1.} How do you compute \[
\log\left(\frac{1}{1+e^{-x}}\right),
\] so as to avoid overflow and underflow of \(e^{-x}\), for both large
positive and large negative \(x\)? (Hint: write it as
\(-\log(1+e^{-x})\) and consider \texttt{log1p} and the two forms of the
sigmoid.)

\textbf{Question 2 (log-softmax).} How do you compute \[
\log\left(\frac{e^{x_k}}{\sum_{j=1}^{K} e^{x_j}}\right),
\] so as to avoid underflow and overflow of \(e^{x_j}\)? Verify your
formula against the \texttt{log-softmax} example above for
\(x = (1000, 1001, 1005)\) and for \(x = (-1000, -1001, -1005)\).

\textbf{Question 3.} The one-pass variance formula fails for
\(x = (M+1, M+2, M+3)\) with \(M = 10^{8}\). Predict the smallest \(M\)
(a power of 10) at which the one-pass result in R first differs from
\(2\), then verify using the scan in the variance section.

\textbf{Question 4.} Write a function \texttt{log\_dbinom(k,\ n,\ p)}
that returns \(\log\binom{n}{k}p^k(1-p)^{n-k}\) for \(n = 10^5\) without
overflow or underflow, using \texttt{lchoose} and \texttt{log1p}, and
compare it with \texttt{dbinom(k,\ n,\ p,\ log\ =\ TRUE)}.

\textbf{Question 5.} Implement Welford's algorithm for the sample
covariance of two variables and check it against \texttt{cov()} on data
shifted by \(10^{10}\).

\bookmarksetup{startatroot}

\chapter{Random Number Generator and Monte Carlo
Elements}\label{random-number-generator-and-monte-carlo-elements}

\section{Pseudo random numbers}\label{pseudo-random-numbers}

Computers generate ``random'' numbers using deterministic algorithms.
One of the most common methods is the \textbf{Linear Congruential
Generator (LCG)}. It produces a sequence of pseudo-random integers
\(X_n\) defined by the recurrence relation:

\[
X_n = (a X_{n-1} + c) \pmod M
\]

where \(X_0\) is the seed, \(a\) is the multiplier, \(c\) is the
increment, and \(M\) is the modulus. To get numbers in the interval
\([0,1]\), we divide by \(M-1\). In this example, we use the parameters
\(a = 7^5\), \(c = 0\), and \(M = 2^{31}-1\) (a minimal standard LCG).

\begin{Shaded}
\begin{Highlighting}[]
\NormalTok{A }\OtherTok{\textless{}{-}} \DecValTok{7}\SpecialCharTok{\^{}}\DecValTok{5}
\NormalTok{M }\OtherTok{\textless{}{-}}  \DecValTok{2}\SpecialCharTok{\^{}}\DecValTok{31{-}1}

\NormalTok{N }\OtherTok{\textless{}{-}} \DecValTok{500}
\NormalTok{rn }\OtherTok{\textless{}{-}} \FunctionTok{rep}\NormalTok{ (}\DecValTok{0}\NormalTok{, N)}
\NormalTok{rn[}\DecValTok{1}\NormalTok{] }\OtherTok{\textless{}{-}} \DecValTok{10} \CommentTok{\# Initial seed}
\ControlFlowTok{for}\NormalTok{ (i }\ControlFlowTok{in} \DecValTok{2}\SpecialCharTok{:}\FunctionTok{length}\NormalTok{ (rn))}
\NormalTok{\{}
\NormalTok{    rn[i] }\OtherTok{\textless{}{-}}\NormalTok{ (A }\SpecialCharTok{*}\NormalTok{ rn[i}\DecValTok{{-}1}\NormalTok{] ) }\SpecialCharTok{\%\%}\NormalTok{ M}
\NormalTok{\}}

\DocumentationTok{\#\# Scale the numbers to fall between 0 and 1}
\NormalTok{nrn }\OtherTok{\textless{}{-}}\NormalTok{ rn}\SpecialCharTok{/}\NormalTok{(M}\DecValTok{{-}1}\NormalTok{)}

\FunctionTok{par}\NormalTok{(}\AttributeTok{mfrow=}\FunctionTok{c}\NormalTok{(}\DecValTok{1}\NormalTok{,}\DecValTok{3}\NormalTok{),}\AttributeTok{mar=}\FunctionTok{c}\NormalTok{(}\DecValTok{4}\NormalTok{,}\DecValTok{4}\NormalTok{,}\DecValTok{3}\NormalTok{,}\DecValTok{1}\NormalTok{))}
\FunctionTok{plot}\NormalTok{(nrn[}\DecValTok{1}\SpecialCharTok{:}\DecValTok{100}\NormalTok{], }\AttributeTok{main=}\StringTok{"Trace Plot"}\NormalTok{, }\AttributeTok{ylab=}\StringTok{"Value"}\NormalTok{, }\AttributeTok{xlab=}\StringTok{"Index"}\NormalTok{)}
\FunctionTok{acf}\NormalTok{(nrn, }\AttributeTok{main=}\StringTok{"Autocorrelation"}\NormalTok{)}
\FunctionTok{hist}\NormalTok{(nrn, }\AttributeTok{main=}\StringTok{"Histogram"}\NormalTok{)}
\end{Highlighting}
\end{Shaded}

\pandocbounded{\includegraphics[keepaspectratio]{chapters/unit03_sampling_basics_files/figure-pdf/unnamed-chunk-1-1.pdf}}

\begin{Shaded}
\begin{Highlighting}[]
\DocumentationTok{\#\# Compare with R\textquotesingle{}s built{-}in runif()}
\NormalTok{n }\OtherTok{\textless{}{-}} \DecValTok{500}
\NormalTok{a }\OtherTok{\textless{}{-}} \FunctionTok{runif}\NormalTok{(n)}
\FunctionTok{par}\NormalTok{(}\AttributeTok{mfrow=}\FunctionTok{c}\NormalTok{(}\DecValTok{1}\NormalTok{,}\DecValTok{2}\NormalTok{),}\AttributeTok{mar=}\FunctionTok{c}\NormalTok{(}\DecValTok{4}\NormalTok{,}\DecValTok{4}\NormalTok{,}\DecValTok{3}\NormalTok{,}\DecValTok{1}\NormalTok{))}
\FunctionTok{hist}\NormalTok{(a,}\AttributeTok{xlab=}\StringTok{"Random Numbers"}\NormalTok{,}\AttributeTok{main=}\StringTok{"Built{-}in runif()"}\NormalTok{)}
\FunctionTok{acf}\NormalTok{(a,}\AttributeTok{main=}\StringTok{"Built{-}in acf"}\NormalTok{)}
\end{Highlighting}
\end{Shaded}

\pandocbounded{\includegraphics[keepaspectratio]{chapters/unit03_sampling_basics_files/figure-pdf/unnamed-chunk-1-2.pdf}}

\section{Inverting CDF}\label{inverting-cdf}

If we can generate uniform random variables \(U \sim \text{Unif}(0,1)\),
we can generate a random variable \(X\) from any continuous distribution
using the \textbf{Inverse Transform Method}. If \(F(x)\) is the
cumulative distribution function (CDF) of \(X\), then \(X = F^{-1}(U)\)
has the desired distribution.

For an Exponential distribution with rate \(\lambda = 1\), the CDF is
\(F(x) = 1 - e^{-x}\). Setting \(U = 1 - e^{-x}\) and solving for \(x\)
gives:

\[
x = -\ln(1-U)
\]

\begin{Shaded}
\begin{Highlighting}[]
\DocumentationTok{\#\# generate exponential random numbers}
\DocumentationTok{\#\# use method of inverse cdf to generate iid sample from exp(1)}
\NormalTok{gen\_exp }\OtherTok{\textless{}{-}} \ControlFlowTok{function}\NormalTok{(n)}
\NormalTok{\{}
    \CommentTok{\#generate unif(0,1) random numbers}
\NormalTok{    u }\OtherTok{\textless{}{-}} \FunctionTok{runif}\NormalTok{(n)}
    \CommentTok{\#transform the random numbers}
    \SpecialCharTok{{-}}\FunctionTok{log}\NormalTok{(}\DecValTok{1}\SpecialCharTok{{-}}\NormalTok{u)}
\NormalTok{\}}

\NormalTok{b }\OtherTok{\textless{}{-}} \FunctionTok{gen\_exp}\NormalTok{(}\DecValTok{10000}\NormalTok{)}
\FunctionTok{par}\NormalTok{(}\AttributeTok{mfrow=}\FunctionTok{c}\NormalTok{(}\DecValTok{1}\NormalTok{,}\DecValTok{3}\NormalTok{),}\AttributeTok{mar=}\FunctionTok{c}\NormalTok{(}\DecValTok{4}\NormalTok{,}\DecValTok{4}\NormalTok{,}\DecValTok{3}\NormalTok{,}\DecValTok{1}\NormalTok{))}
\FunctionTok{hist}\NormalTok{(b, }\AttributeTok{main=}\StringTok{"Custom gen\_exp"}\NormalTok{)}
\FunctionTok{acf}\NormalTok{(b, }\AttributeTok{main=}\StringTok{"Custom acf"}\NormalTok{)}

\DocumentationTok{\#\#\# r built{-}in generator for comparison}
\FunctionTok{hist}\NormalTok{(}\FunctionTok{rexp}\NormalTok{ (}\DecValTok{10000}\NormalTok{), }\AttributeTok{main=}\StringTok{"Built{-}in rexp"}\NormalTok{)}
\end{Highlighting}
\end{Shaded}

\pandocbounded{\includegraphics[keepaspectratio]{chapters/unit03_sampling_basics_files/figure-pdf/unnamed-chunk-2-1.pdf}}

\section{A Special Transformation for Generating Normal
Sample}\label{a-special-transformation-for-generating-normal-sample}

The \textbf{Box-Muller Transform} is an elegant method for generating
pairs of independent Standard Normal random variables from independent
Uniform random variables. It relies on converting cartesian coordinates
\((X,Y)\) into polar coordinates \((R, \theta)\).

If \(X, Y \sim N(0,1)\) are independent, their squared radius
\(R^2 = X^2 + Y^2\) follows an Exponential distribution with rate
\(\lambda = 1/2\), and their angle \(\theta\) is uniformly distributed
over \([0, 2\pi]\). By generating \(R\) and \(\theta\), we can map them
back to Normal variables:

\[
X = R\cos(\theta) \quad \text{and} \quad Y = R\sin(\theta)
\]

\begin{Shaded}
\begin{Highlighting}[]
\NormalTok{gen\_normal }\OtherTok{\textless{}{-}} \ControlFlowTok{function}\NormalTok{(n)}
\NormalTok{\{}
    \CommentTok{\# calculates size of random samples (we generate pairs, so we need half)}
\NormalTok{    size\_sample }\OtherTok{\textless{}{-}} \FunctionTok{ceiling}\NormalTok{(n}\SpecialCharTok{/}\DecValTok{2}\NormalTok{)}
    
    \CommentTok{\# R\^{}2 \textasciitilde{} Exp(1/2), so we take the square root to get R}
\NormalTok{    R }\OtherTok{\textless{}{-}} \FunctionTok{sqrt}\NormalTok{(}\DecValTok{2}\SpecialCharTok{*}\FunctionTok{rexp}\NormalTok{(size\_sample))}
\NormalTok{    theta }\OtherTok{\textless{}{-}} \FunctionTok{runif}\NormalTok{(size\_sample,}\DecValTok{0}\NormalTok{,}\DecValTok{2}\SpecialCharTok{*}\NormalTok{pi)}
    
\NormalTok{    X }\OtherTok{\textless{}{-}}\NormalTok{ R}\SpecialCharTok{*}\FunctionTok{cos}\NormalTok{(theta)  }
\NormalTok{    Y }\OtherTok{\textless{}{-}}\NormalTok{ R}\SpecialCharTok{*}\FunctionTok{sin}\NormalTok{(theta)}
    
    \FunctionTok{c}\NormalTok{(X,Y)[}\DecValTok{1}\SpecialCharTok{:}\NormalTok{n]  }
\NormalTok{\}}

\NormalTok{normal\_sample }\OtherTok{\textless{}{-}} \FunctionTok{gen\_normal}\NormalTok{(}\DecValTok{1000}\NormalTok{)}

\FunctionTok{par}\NormalTok{(}\AttributeTok{mfrow=}\FunctionTok{c}\NormalTok{(}\DecValTok{1}\NormalTok{,}\DecValTok{2}\NormalTok{),}\AttributeTok{mar=}\FunctionTok{c}\NormalTok{(}\DecValTok{4}\NormalTok{,}\DecValTok{4}\NormalTok{,}\DecValTok{1}\NormalTok{,}\DecValTok{1}\NormalTok{))}
\FunctionTok{hist}\NormalTok{(normal\_sample,}\AttributeTok{main=}\StringTok{"Custom Normal"}\NormalTok{)}
\FunctionTok{qqnorm}\NormalTok{(normal\_sample, }\AttributeTok{main=}\StringTok{"Custom QQ"}\NormalTok{)}
\FunctionTok{qqline}\NormalTok{(normal\_sample)}
\end{Highlighting}
\end{Shaded}

\pandocbounded{\includegraphics[keepaspectratio]{chapters/unit03_sampling_basics_files/figure-pdf/unnamed-chunk-3-1.pdf}}

\begin{Shaded}
\begin{Highlighting}[]
\DocumentationTok{\#\# R\textquotesingle{}s built{-}in generator for comparison}
\FunctionTok{par}\NormalTok{(}\AttributeTok{mfrow=}\FunctionTok{c}\NormalTok{(}\DecValTok{1}\NormalTok{,}\DecValTok{2}\NormalTok{),}\AttributeTok{mar=}\FunctionTok{c}\NormalTok{(}\DecValTok{4}\NormalTok{,}\DecValTok{4}\NormalTok{,}\DecValTok{1}\NormalTok{,}\DecValTok{1}\NormalTok{))}
\NormalTok{nsample2 }\OtherTok{\textless{}{-}} \FunctionTok{rnorm}\NormalTok{ (}\DecValTok{1000}\NormalTok{) }
\FunctionTok{hist}\NormalTok{(nsample2, }\AttributeTok{main=}\StringTok{"Built{-}in Normal"}\NormalTok{)}
\FunctionTok{qqnorm}\NormalTok{(nsample2, }\AttributeTok{main=}\StringTok{"Built{-}in QQ"}\NormalTok{)}
\FunctionTok{qqline}\NormalTok{(nsample2)}
\end{Highlighting}
\end{Shaded}

\pandocbounded{\includegraphics[keepaspectratio]{chapters/unit03_sampling_basics_files/figure-pdf/unnamed-chunk-3-2.pdf}}

\section{Demonstration of CLT and
LLN}\label{demonstration-of-clt-and-lln}

The \textbf{Law of Large Numbers (LLN)} states that as the sample size
\(n\) grows, the sample mean \(\bar{X}_n\) converges to the expected
value \(\mu\): \[
\bar{X}_n \to \mu \quad \text{as} \quad n \to \infty
\]

The \textbf{Central Limit Theorem (CLT)} states that the standardized
sample mean converges in distribution to a Standard Normal \(N(0,1)\):
\[
Z = \frac{\bar{X}_n - \mu}{\sigma/\sqrt{n}} \xrightarrow{d} N(0,1) \quad \text{as} \quad n \to \infty
\]

Here we draw from a Gamma distribution with shape parameter
\(\alpha = 2\) and rate \(\beta = 1\). The theoretical mean is
\(\mu = \alpha/\beta = 2\), and the variance is
\(\sigma^2 = \alpha/\beta^2 = 2\).

\begin{Shaded}
\begin{Highlighting}[]
\NormalTok{n }\OtherTok{\textless{}{-}} \DecValTok{100}
\NormalTok{rn }\OtherTok{\textless{}{-}} \FunctionTok{rgamma}\NormalTok{(n, }\AttributeTok{shape =} \DecValTok{2}\NormalTok{)}

\FunctionTok{par}\NormalTok{(}\AttributeTok{mfrow =} \FunctionTok{c}\NormalTok{(}\DecValTok{2}\NormalTok{,}\DecValTok{2}\NormalTok{))}
\FunctionTok{plot}\NormalTok{(rn[}\DecValTok{1}\SpecialCharTok{:}\DecValTok{100}\NormalTok{], }\AttributeTok{main=}\StringTok{"Raw Gamma Draws"}\NormalTok{, }\AttributeTok{ylab=}\StringTok{"Value"}\NormalTok{)}
\FunctionTok{hist}\NormalTok{(rn, }\AttributeTok{main=}\StringTok{"Histogram of Draws"}\NormalTok{)}

\DocumentationTok{\#\# Calculate running averages (LLN)}
\NormalTok{xbar.rep }\OtherTok{\textless{}{-}} \FunctionTok{replicate}\NormalTok{(}\DecValTok{5}\NormalTok{, }\FunctionTok{cumsum}\NormalTok{(}\FunctionTok{rgamma}\NormalTok{(n, }\AttributeTok{shape =} \DecValTok{2}\NormalTok{))}\SpecialCharTok{/}\NormalTok{(}\DecValTok{1}\SpecialCharTok{:}\NormalTok{n))}
\DocumentationTok{\#\# Calculate running standardized means (CLT)}
\NormalTok{se.rep }\OtherTok{\textless{}{-}} \FunctionTok{replicate}\NormalTok{(}\DecValTok{5}\NormalTok{, (}\FunctionTok{cumsum}\NormalTok{(}\FunctionTok{rgamma}\NormalTok{(n, }\AttributeTok{shape =} \DecValTok{2}\NormalTok{))}\SpecialCharTok{/}\NormalTok{(}\DecValTok{1}\SpecialCharTok{:}\NormalTok{n) }\SpecialCharTok{{-}} \DecValTok{2}\NormalTok{) }\SpecialCharTok{/} \FunctionTok{sqrt}\NormalTok{(}\DecValTok{2}\SpecialCharTok{/}\NormalTok{(}\DecValTok{1}\SpecialCharTok{:}\NormalTok{n)))}

\DocumentationTok{\#\# Plot LLN: should converge to mu = 2}
\FunctionTok{plot}\NormalTok{(xbar.rep[,}\DecValTok{1}\NormalTok{], }\AttributeTok{main =} \StringTok{"Demonstrating LLN"}\NormalTok{, }\AttributeTok{log =} \StringTok{"x"}\NormalTok{, }\AttributeTok{type =} \StringTok{"b"}\NormalTok{, }\AttributeTok{ylim =} \FunctionTok{c}\NormalTok{(}\DecValTok{0}\NormalTok{, }\DecValTok{6}\NormalTok{), }\AttributeTok{ylab=}\StringTok{"Sample Mean"}\NormalTok{)}
\FunctionTok{abline}\NormalTok{(}\AttributeTok{h =} \DecValTok{2}\NormalTok{, }\AttributeTok{lwd=}\DecValTok{2}\NormalTok{)}
\ControlFlowTok{for}\NormalTok{ (i }\ControlFlowTok{in} \DecValTok{2}\SpecialCharTok{:}\FunctionTok{ncol}\NormalTok{(xbar.rep)) }\FunctionTok{points}\NormalTok{(xbar.rep[,i], }\AttributeTok{col =}\NormalTok{ i, }\AttributeTok{pch =}\NormalTok{ i, }\AttributeTok{type =} \StringTok{"b"}\NormalTok{)}

\DocumentationTok{\#\# Plot CLT: should stabilize within roughly [{-}3, 3] standard deviations}
\FunctionTok{plot}\NormalTok{(se.rep[,}\DecValTok{1}\NormalTok{], }\AttributeTok{ylim =} \FunctionTok{c}\NormalTok{(}\SpecialCharTok{{-}}\DecValTok{3}\NormalTok{,}\DecValTok{3}\NormalTok{), }\AttributeTok{type=}\StringTok{"b"}\NormalTok{, }\AttributeTok{log =} \StringTok{"x"}\NormalTok{, }\AttributeTok{main=}\StringTok{"Demonstrating CLT"}\NormalTok{, }\AttributeTok{ylab=}\StringTok{"Z{-}Score"}\NormalTok{)}
\ControlFlowTok{for}\NormalTok{ (i }\ControlFlowTok{in} \DecValTok{2}\SpecialCharTok{:}\FunctionTok{ncol}\NormalTok{(se.rep)) }\FunctionTok{points}\NormalTok{(se.rep[,i], }\AttributeTok{col =}\NormalTok{ i, }\AttributeTok{pch =}\NormalTok{ i, }\AttributeTok{type =} \StringTok{"b"}\NormalTok{)}
\FunctionTok{abline}\NormalTok{(}\AttributeTok{h =} \FunctionTok{c}\NormalTok{(}\DecValTok{0}\NormalTok{, }\SpecialCharTok{{-}}\DecValTok{2}\NormalTok{, }\DecValTok{2}\NormalTok{), }\AttributeTok{lty =} \FunctionTok{c}\NormalTok{(}\DecValTok{1}\NormalTok{,}\DecValTok{2}\NormalTok{,}\DecValTok{2}\NormalTok{), }\AttributeTok{lwd=}\DecValTok{2}\NormalTok{)}
\end{Highlighting}
\end{Shaded}

\pandocbounded{\includegraphics[keepaspectratio]{chapters/unit03_sampling_basics_files/figure-pdf/unnamed-chunk-4-1.pdf}}

\section{\texorpdfstring{An Example of Monte Carlo for Estimating
\(\pi\)}{An Example of Monte Carlo for Estimating \textbackslash pi}}\label{an-example-of-monte-carlo-for-estimating-pi}

We can estimate the value of \(\pi\) by throwing random darts at a
square with side length 2 (from \(-1\) to \(1\)) and counting how many
land inside the inscribed circle of radius \(1\).

The area of the square is \((2)^2 = 4\), and the area of the circle is
\(\pi(1)^2 = \pi\). Therefore, the probability that a uniform random
point \((X,Y)\) falls inside the circle is: \[
P(X^2 + Y^2 \le 1) = \frac{\text{Area of Circle}}{\text{Area of Square}} = \frac{\pi}{4}
\]

\subsubsection{Visualizing the Method}\label{visualizing-the-method}

Before we estimate \(\pi\), let's visually simulate throwing 1,000
darts:

\begin{Shaded}
\begin{Highlighting}[]
\FunctionTok{set.seed}\NormalTok{(}\DecValTok{123}\NormalTok{)}
\NormalTok{n\_visual }\OtherTok{\textless{}{-}} \DecValTok{1000}
\NormalTok{x\_pts }\OtherTok{\textless{}{-}} \FunctionTok{runif}\NormalTok{(n\_visual, }\SpecialCharTok{{-}}\DecValTok{1}\NormalTok{, }\DecValTok{1}\NormalTok{)}
\NormalTok{y\_pts }\OtherTok{\textless{}{-}} \FunctionTok{runif}\NormalTok{(n\_visual, }\SpecialCharTok{{-}}\DecValTok{1}\NormalTok{, }\DecValTok{1}\NormalTok{)}
\NormalTok{inside\_circle }\OtherTok{\textless{}{-}}\NormalTok{ (x\_pts}\SpecialCharTok{\^{}}\DecValTok{2} \SpecialCharTok{+}\NormalTok{ y\_pts}\SpecialCharTok{\^{}}\DecValTok{2}\NormalTok{) }\SpecialCharTok{\textless{}=} \DecValTok{1}

\DocumentationTok{\#\# Plot points}
\FunctionTok{plot}\NormalTok{(x\_pts, y\_pts, }\AttributeTok{col =} \FunctionTok{ifelse}\NormalTok{(inside\_circle, }\StringTok{"steelblue"}\NormalTok{, }\StringTok{"tomato"}\NormalTok{), }
     \AttributeTok{pch =} \DecValTok{20}\NormalTok{, }\AttributeTok{asp =} \DecValTok{1}\NormalTok{, }\AttributeTok{xlab =} \StringTok{"X"}\NormalTok{, }\AttributeTok{ylab =} \StringTok{"Y"}\NormalTok{, }
     \AttributeTok{main =} \FunctionTok{paste}\NormalTok{(}\FunctionTok{sum}\NormalTok{(inside\_circle), }\StringTok{"out of"}\NormalTok{, n\_visual, }\StringTok{"points in circle"}\NormalTok{))}

\DocumentationTok{\#\# Draw the boundaries}
\FunctionTok{rect}\NormalTok{(}\SpecialCharTok{{-}}\DecValTok{1}\NormalTok{, }\SpecialCharTok{{-}}\DecValTok{1}\NormalTok{, }\DecValTok{1}\NormalTok{, }\DecValTok{1}\NormalTok{, }\AttributeTok{border =} \StringTok{"black"}\NormalTok{, }\AttributeTok{lwd =} \DecValTok{2}\NormalTok{)}
\NormalTok{theta\_seq }\OtherTok{\textless{}{-}} \FunctionTok{seq}\NormalTok{(}\DecValTok{0}\NormalTok{, }\DecValTok{2}\SpecialCharTok{*}\NormalTok{pi, }\AttributeTok{length.out =} \DecValTok{100}\NormalTok{)}
\FunctionTok{lines}\NormalTok{(}\FunctionTok{cos}\NormalTok{(theta\_seq), }\FunctionTok{sin}\NormalTok{(theta\_seq), }\AttributeTok{col =} \StringTok{"black"}\NormalTok{, }\AttributeTok{lwd =} \DecValTok{2}\NormalTok{)}
\end{Highlighting}
\end{Shaded}

\pandocbounded{\includegraphics[keepaspectratio]{chapters/unit03_sampling_basics_files/figure-pdf/unnamed-chunk-5-1.pdf}}

\subsubsection{The Monte Carlo Estimate}\label{the-monte-carlo-estimate}

If we define an indicator variable \(I_i\) that equals 1 if the \(i\)-th
dart lands in the circle and 0 otherwise, then \(E[I_i] = \pi/4\). We
define \(Z_i = 4 \times I_i\), which gives \(E[Z_i] = \pi\).

Our Monte Carlo estimate \(\hat{\pi}\) is the sample average of the
\(Z_i\) values: \[
\hat{\pi} = \bar{Z} = \frac{1}{n}\sum_{i=1}^n Z_i
\]

The Standard Error (SE) measures the uncertainty of our estimate. Using
the sample standard deviation \(S_Z\), the 95\% margin of error is
calculated as: \[
\text{Margin of Error} = 1.96 \times \text{SE} = 1.96 \times \frac{S_Z}{\sqrt{n}}
\]

\subsubsection{Results by Sample Size}\label{results-by-sample-size}

We encapsulate the logic in a function and track how precision improves
as \(n\) increases. The table below outlines our findings.

\begin{Shaded}
\begin{Highlighting}[]
\FunctionTok{library}\NormalTok{(knitr)}

\DocumentationTok{\#\# n is the number of samples drawn uniformly from the rectangle ({-}1,1)  * ({-}1,1)}
\DocumentationTok{\#\# an estimate of pi is returned}
\NormalTok{pi\_est\_mc }\OtherTok{\textless{}{-}} \ControlFlowTok{function}\NormalTok{(n)}
\NormalTok{\{}
\NormalTok{    X }\OtherTok{\textless{}{-}} \FunctionTok{runif}\NormalTok{(n,}\SpecialCharTok{{-}}\DecValTok{1}\NormalTok{,}\DecValTok{1}\NormalTok{)}
\NormalTok{    Y }\OtherTok{\textless{}{-}} \FunctionTok{runif}\NormalTok{(n,}\SpecialCharTok{{-}}\DecValTok{1}\NormalTok{,}\DecValTok{1}\NormalTok{)}
    
    \CommentTok{\# 4 times the indicator function}
\NormalTok{    Z }\OtherTok{\textless{}{-}} \DecValTok{4} \SpecialCharTok{*}\NormalTok{ (X}\SpecialCharTok{\^{}}\DecValTok{2} \SpecialCharTok{+}\NormalTok{ Y}\SpecialCharTok{\^{}}\DecValTok{2} \SpecialCharTok{\textless{}=} \DecValTok{1}\NormalTok{)}
    
\NormalTok{    mu }\OtherTok{\textless{}{-}} \FunctionTok{mean}\NormalTok{(Z)}
\NormalTok{    error }\OtherTok{\textless{}{-}} \FloatTok{1.96} \SpecialCharTok{*} \FunctionTok{sd}\NormalTok{(Z) }\SpecialCharTok{/} \FunctionTok{sqrt}\NormalTok{(n)}
    
    \FunctionTok{list}\NormalTok{(}\AttributeTok{n =}\NormalTok{ n, }\AttributeTok{pi.est =}\NormalTok{ mu, }\AttributeTok{error.95perc =}\NormalTok{ error, }\AttributeTok{ci.lower =}\NormalTok{ mu }\SpecialCharTok{{-}}\NormalTok{ error, }\AttributeTok{ci.upper =}\NormalTok{ mu }\SpecialCharTok{+}\NormalTok{ error)}
\NormalTok{\}  }

\DocumentationTok{\#\# Run the estimation for varying sample sizes}
\NormalTok{sample\_sizes }\OtherTok{\textless{}{-}} \FunctionTok{c}\NormalTok{(}\DecValTok{100}\NormalTok{, }\DecValTok{10000}\NormalTok{, }\DecValTok{100000}\NormalTok{, }\DecValTok{10000000}\NormalTok{)}
\NormalTok{results\_list }\OtherTok{\textless{}{-}} \FunctionTok{lapply}\NormalTok{(sample\_sizes, pi\_est\_mc)}

\DocumentationTok{\#\# Combine results into a dataframe for tabular display}
\NormalTok{results\_df }\OtherTok{\textless{}{-}} \FunctionTok{data.frame}\NormalTok{(}
  \AttributeTok{N =} \FunctionTok{sapply}\NormalTok{(results\_list, }\ControlFlowTok{function}\NormalTok{(x) x}\SpecialCharTok{$}\NormalTok{n),}
  \AttributeTok{Estimate =} \FunctionTok{sapply}\NormalTok{(results\_list, }\ControlFlowTok{function}\NormalTok{(x) x}\SpecialCharTok{$}\NormalTok{pi.est),}
  \AttributeTok{Error\_Margin =} \FunctionTok{sapply}\NormalTok{(results\_list, }\ControlFlowTok{function}\NormalTok{(x) x}\SpecialCharTok{$}\NormalTok{error}\FloatTok{.95}\NormalTok{perc),}
  \AttributeTok{CI\_Lower =} \FunctionTok{sapply}\NormalTok{(results\_list, }\ControlFlowTok{function}\NormalTok{(x) x}\SpecialCharTok{$}\NormalTok{ci.lower),}
  \AttributeTok{CI\_Upper =} \FunctionTok{sapply}\NormalTok{(results\_list, }\ControlFlowTok{function}\NormalTok{(x) x}\SpecialCharTok{$}\NormalTok{ci.upper)}
\NormalTok{)}

\DocumentationTok{\#\# Render the formatted table}
\FunctionTok{kable}\NormalTok{(results\_df, }
      \AttributeTok{col.names =} \FunctionTok{c}\NormalTok{(}\StringTok{"Sample Size (N)"}\NormalTok{, }\StringTok{"Estimated Pi"}\NormalTok{, }\StringTok{"95\% Margin of Error"}\NormalTok{, }\StringTok{"Lower CI"}\NormalTok{, }\StringTok{"Upper CI"}\NormalTok{), }
      \AttributeTok{digits =} \DecValTok{6}\NormalTok{, }
      \AttributeTok{caption =} \StringTok{"Monte Carlo Estimation of Pi across different sample sizes"}\NormalTok{)}
\end{Highlighting}
\end{Shaded}

\begin{longtable}[]{@{}
  >{\raggedleft\arraybackslash}p{(\linewidth - 8\tabcolsep) * \real{0.2388}}
  >{\raggedleft\arraybackslash}p{(\linewidth - 8\tabcolsep) * \real{0.1940}}
  >{\raggedleft\arraybackslash}p{(\linewidth - 8\tabcolsep) * \real{0.2985}}
  >{\raggedleft\arraybackslash}p{(\linewidth - 8\tabcolsep) * \real{0.1343}}
  >{\raggedleft\arraybackslash}p{(\linewidth - 8\tabcolsep) * \real{0.1343}}@{}}
\caption{Monte Carlo Estimation of Pi across different sample
sizes}\tabularnewline
\toprule\noalign{}
\begin{minipage}[b]{\linewidth}\raggedleft
Sample Size (N)
\end{minipage} & \begin{minipage}[b]{\linewidth}\raggedleft
Estimated Pi
\end{minipage} & \begin{minipage}[b]{\linewidth}\raggedleft
95\% Margin of Error
\end{minipage} & \begin{minipage}[b]{\linewidth}\raggedleft
Lower CI
\end{minipage} & \begin{minipage}[b]{\linewidth}\raggedleft
Upper CI
\end{minipage} \\
\midrule\noalign{}
\endfirsthead
\toprule\noalign{}
\begin{minipage}[b]{\linewidth}\raggedleft
Sample Size (N)
\end{minipage} & \begin{minipage}[b]{\linewidth}\raggedleft
Estimated Pi
\end{minipage} & \begin{minipage}[b]{\linewidth}\raggedleft
95\% Margin of Error
\end{minipage} & \begin{minipage}[b]{\linewidth}\raggedleft
Lower CI
\end{minipage} & \begin{minipage}[b]{\linewidth}\raggedleft
Upper CI
\end{minipage} \\
\midrule\noalign{}
\endhead
\bottomrule\noalign{}
\endlastfoot
1e+02 & 2.920000 & 0.349818 & 2.570182 & 3.269818 \\
1e+04 & 3.156400 & 0.031985 & 3.124415 & 3.188385 \\
1e+05 & 3.143600 & 0.010170 & 3.133430 & 3.153770 \\
1e+07 & 3.141769 & 0.001018 & 3.140751 & 3.142787 \\
\end{longtable}

\bookmarksetup{startatroot}

\chapter{Point Estimation and Hypothesis
Testing}\label{point-estimation-and-hypothesis-testing}

\section{Point Estimation}\label{point-estimation}

Assume data \(X = (X_1, \ldots, X_n)\) is drawn from a distribution
\(f(x; \theta)\) indexed by an unknown parameter \(\theta \in \Theta\).
An estimator \(\hat\theta = T(X)\) is a function derived exclusively
from the data. Because it relies on random samples, \(\hat\theta\) is
itself a random variable. Its behavior is defined by its sampling
distribution.

The two fundamental summaries of an estimator's performance are bias and
variance:

\[\mathrm{bias}(\hat\theta) = E_{\theta}(\hat\theta) - \theta, \qquad \mathrm{Var}_{\theta}(\hat\theta) = E_{\theta}\left[\left(\hat\theta - E_{\theta}(\hat\theta)\right)^{2}\right].\]

The most comprehensive single-number summary is the Mean Squared Error
(MSE), defined as
\(\mathrm{MSE}_{\theta}(\hat\theta) = E_{\theta}[(\hat\theta - \theta)^{2}]\).
By adding and subtracting \(E_{\theta}(\hat\theta)\) inside the square,
the cross term vanishes, yielding the MSE decomposition:

\[\mathrm{MSE}_{\theta}(\hat\theta) = \mathrm{Var}_{\theta}(\hat\theta) + \left[\mathrm{bias}_{\theta}(\hat\theta)\right]^{2}.\]

\textbf{Key Insight:} Unbiasedness is not strictly necessary. An
estimator that introduces a small bias in exchange for a massive
reduction in variance will frequently achieve a lower MSE than a
strictly unbiased competitor.

Because MSE is a \emph{function} of \(\theta\), it acts as the risk
function under squared-error loss. Two estimators rarely have one
absolute winner; their risk curves often cross. To choose a single
estimator globally, statisticians summarize the risk curve via two main
criteria:

\begin{itemize}
\tightlist
\item
  \textbf{Minimax:} Minimize the worst-case scenario:
  \(\sup_{\theta \in \Theta} \mathrm{MSE}_{\theta}(\hat\theta)\)
\item
  \textbf{Bayes Risk:} Average the error against a prior distribution
  \(\pi\):
  \(\int_{\Theta} \mathrm{MSE}_{\theta}(\hat\theta)\, \pi(\theta)\, d\theta\)
\end{itemize}

Asymptotically, \(\hat\theta_n\) is consistent if it converges to
\(\theta\) in probability. Under regularity conditions, Maximum
Likelihood Estimators (MLEs) are consistent and asymptotically normal,
achieving the Cramér--Rao lower variance bound \(1/I(\theta)\). However,
in small samples, the MLE can often be outperformed.

\subsection{Monte Carlo Evaluation of an
Estimator}\label{monte-carlo-evaluation-of-an-estimator}

Risk functions are rarely available in closed form, but because they are
expectations, we can estimate them efficiently using Monte Carlo
simulation.

\textbf{The Monte Carlo Recipe:}

\begin{enumerate}
\def\labelenumi{\arabic{enumi}.}
\tightlist
\item
  Define a grid of \(\theta\) values.
\item
  For each \(\theta\), generate \(M\) independent datasets.
\item
  Compute the estimate \(\hat\theta^{(m)}\) for each dataset.
\item
  Average the squared errors:
\end{enumerate}

\[\widehat{\mathrm{MSE}}_{\theta} = \frac{1}{M} \sum_{m=1}^{M} \left(\hat\theta^{(m)} - \theta\right)^{2}\]

Because this represents an average of independent, identically
distributed (i.i.d.) quantities, its standard error is \(s/\sqrt{M}\),
where \(s^{2}\) is the sample variance of the squared errors. Always
report the Monte Carlo error boundaries
(\(\widehat{\mathrm{MSE}}_{\theta} \pm 1.96\, s/\sqrt{M}\)) to
distinguish true differences in risk curves from mere simulation noise.

\subsection{Two Estimators of
Proportions}\label{two-estimators-of-proportions}

Consider estimating a binomial proportion \(p\) from \(n\) independent
trials.

\begin{itemize}
\tightlist
\item
  \textbf{Maximum Likelihood Estimator (\(\hat{p}_1\)):} The sample
  mean, \(\hat{p}_1 = \bar{X}\). This is unbiased but exhibits high
  variance in small samples, especially near boundaries (\(p \approx 0\)
  or \(1\)).
\item
  \textbf{Smoothed Estimator (\(\hat{p}_2\)):} Often called the Laplace
  or ``add-one-in'' estimator, \(\hat{p}_2 = (\sum_i X_i + 1)/(n + 2)\).
  This introduces bias by shrinking the estimate toward \(1/2\) (the
  posterior mean under a uniform prior), but reduces variance by a
  factor of \([n/(n+2)]^2\).
\end{itemize}

We can estimate the risk curves on a grid of \(p\) values using the
Monte Carlo method. Both estimators rely strictly on the sufficient
statistic \(S = \sum_i X_i \sim \mathrm{Binomial}(n, p)\), allowing us
to vectorize the simulation without looping over replicates.

\begin{Shaded}
\begin{Highlighting}[]
\DocumentationTok{\#\#\#Estimators, written in terms of the sufficient statistic S = sum(X\_i)}
\NormalTok{p\_est1 }\OtherTok{\textless{}{-}} \ControlFlowTok{function}\NormalTok{(s, n) s }\SpecialCharTok{/}\NormalTok{ n}
\NormalTok{p\_est2 }\OtherTok{\textless{}{-}} \ControlFlowTok{function}\NormalTok{(s, n) (s }\SpecialCharTok{+} \DecValTok{1}\NormalTok{) }\SpecialCharTok{/}\NormalTok{ (n }\SpecialCharTok{+} \DecValTok{2}\NormalTok{)}

\DocumentationTok{\#\#\# Monte Carlo estimate of the MSE and its standard error.  All no\_sim}
\DocumentationTok{\#\#\# replicates are generated in a single vectorized draw.}
\NormalTok{mse\_est\_mc }\OtherTok{\textless{}{-}} \ControlFlowTok{function}\NormalTok{(n, p, no\_sim, p\_est)}
\NormalTok{\{}
\NormalTok{    sq\_error }\OtherTok{\textless{}{-}}\NormalTok{ (}\FunctionTok{p\_est}\NormalTok{(}\FunctionTok{rbinom}\NormalTok{(no\_sim, }\AttributeTok{size =}\NormalTok{ n, }\AttributeTok{prob =}\NormalTok{ p), n) }\SpecialCharTok{{-}}\NormalTok{ p)}\SpecialCharTok{\^{}}\DecValTok{2}
    \FunctionTok{c}\NormalTok{(}\AttributeTok{mse =} \FunctionTok{mean}\NormalTok{(sq\_error), }\AttributeTok{sd =} \FunctionTok{sqrt}\NormalTok{(}\FunctionTok{var}\NormalTok{(sq\_error) }\SpecialCharTok{/}\NormalTok{ no\_sim))}
\NormalTok{\}}

\DocumentationTok{\#\#\# MSE over a grid of p values, returned as a 2 by length(p\_set) matrix}
\NormalTok{mse\_curve }\OtherTok{\textless{}{-}} \ControlFlowTok{function}\NormalTok{(n, p\_set, no\_sim, p\_est)}
    \FunctionTok{vapply}\NormalTok{(p\_set, }\ControlFlowTok{function}\NormalTok{(p) }\FunctionTok{mse\_est\_mc}\NormalTok{(n, p, no\_sim, p\_est), }\FunctionTok{c}\NormalTok{(}\AttributeTok{mse =} \DecValTok{0}\NormalTok{, }\AttributeTok{sd =} \DecValTok{0}\NormalTok{))}
\end{Highlighting}
\end{Shaded}

\begin{Shaded}
\begin{Highlighting}[]
\NormalTok{no\_sim }\OtherTok{\textless{}{-}} \DecValTok{10000}
\NormalTok{p\_set }\OtherTok{\textless{}{-}} \FunctionTok{seq}\NormalTok{(}\DecValTok{0}\NormalTok{, }\DecValTok{1}\NormalTok{, }\AttributeTok{by =} \FloatTok{0.05}\NormalTok{)}
\NormalTok{n\_set }\OtherTok{\textless{}{-}} \FunctionTok{c}\NormalTok{(}\DecValTok{2}\NormalTok{, }\DecValTok{10}\NormalTok{, }\DecValTok{50}\NormalTok{, }\DecValTok{100}\NormalTok{)}

\FunctionTok{par}\NormalTok{(}\AttributeTok{mfrow =} \FunctionTok{c}\NormalTok{(}\DecValTok{2}\NormalTok{, }\DecValTok{2}\NormalTok{), }\AttributeTok{mar =} \FunctionTok{c}\NormalTok{(}\DecValTok{4}\NormalTok{, }\FloatTok{4.5}\NormalTok{, }\FloatTok{2.5}\NormalTok{, }\DecValTok{1}\NormalTok{))}

\ControlFlowTok{for}\NormalTok{ (n }\ControlFlowTok{in}\NormalTok{ n\_set)}
\NormalTok{\{}
\NormalTok{    r1 }\OtherTok{\textless{}{-}} \FunctionTok{mse\_curve}\NormalTok{(n, p\_set, no\_sim, p\_est1)}
\NormalTok{    r2 }\OtherTok{\textless{}{-}} \FunctionTok{mse\_curve}\NormalTok{(n, p\_set, no\_sim, p\_est2)}

\NormalTok{    bands }\OtherTok{\textless{}{-}} \FunctionTok{cbind}\NormalTok{(r1[}\StringTok{"mse"}\NormalTok{, ], r1[}\StringTok{"mse"}\NormalTok{, ] }\SpecialCharTok{+} \FloatTok{1.96} \SpecialCharTok{*}\NormalTok{ r1[}\StringTok{"sd"}\NormalTok{, ], r1[}\StringTok{"mse"}\NormalTok{, ] }\SpecialCharTok{{-}} \FloatTok{1.96} \SpecialCharTok{*}\NormalTok{ r1[}\StringTok{"sd"}\NormalTok{, ],}
\NormalTok{                   r2[}\StringTok{"mse"}\NormalTok{, ], r2[}\StringTok{"mse"}\NormalTok{, ] }\SpecialCharTok{+} \FloatTok{1.96} \SpecialCharTok{*}\NormalTok{ r2[}\StringTok{"sd"}\NormalTok{, ], r2[}\StringTok{"mse"}\NormalTok{, ] }\SpecialCharTok{{-}} \FloatTok{1.96} \SpecialCharTok{*}\NormalTok{ r2[}\StringTok{"sd"}\NormalTok{, ])}

    \FunctionTok{matplot}\NormalTok{(p\_set, bands, }\AttributeTok{type =} \StringTok{"l"}\NormalTok{, }\AttributeTok{lty =} \FunctionTok{c}\NormalTok{(}\DecValTok{1}\NormalTok{, }\DecValTok{2}\NormalTok{, }\DecValTok{2}\NormalTok{, }\DecValTok{1}\NormalTok{, }\DecValTok{2}\NormalTok{, }\DecValTok{2}\NormalTok{), }\AttributeTok{col =} \FunctionTok{rep}\NormalTok{(}\FunctionTok{c}\NormalTok{(}\DecValTok{1}\NormalTok{, }\DecValTok{2}\NormalTok{), }\AttributeTok{each =} \DecValTok{3}\NormalTok{),}
            \AttributeTok{xlab =} \FunctionTok{TeX}\NormalTok{(r}\StringTok{"($p$)"}\NormalTok{),}
            \AttributeTok{ylab =} \FunctionTok{TeX}\NormalTok{(r}\StringTok{"($E[(\textbackslash{}hat\{p\} {-} p)\^{}2]$)"}\NormalTok{),}
            \AttributeTok{main =} \FunctionTok{TeX}\NormalTok{(}\FunctionTok{sprintf}\NormalTok{(r}\StringTok{"($n = \%d$)"}\NormalTok{, n)))}

    \FunctionTok{legend}\NormalTok{(}\StringTok{"top"}\NormalTok{, }\AttributeTok{bty =} \StringTok{"n"}\NormalTok{, }\AttributeTok{lty =} \DecValTok{1}\NormalTok{, }\AttributeTok{col =} \FunctionTok{c}\NormalTok{(}\DecValTok{1}\NormalTok{, }\DecValTok{2}\NormalTok{), }\AttributeTok{cex =} \FloatTok{0.9}\NormalTok{,}
           \AttributeTok{legend =} \FunctionTok{TeX}\NormalTok{(}\FunctionTok{c}\NormalTok{(r}\StringTok{"($\textbackslash{}hat\{p\}\_1 = S/n$)"}\NormalTok{, r}\StringTok{"($\textbackslash{}hat\{p\}\_2 = (S+1)/(n+2)$)"}\NormalTok{)))}
\NormalTok{\}}
\end{Highlighting}
\end{Shaded}

\pandocbounded{\includegraphics[keepaspectratio]{chapters/unit04_simulation_files/figure-pdf/prop-mse-plot-1.pdf}}

When \(n\) is small, the smoothed estimator significantly outperforms
the MLE across most of the parameter space. As \(n\) grows, the curves
converge, reflecting the asymptotic efficiency of the MLE. The exception
occurs strictly at the extremes (\(p \approx 0\) or \(1\)), where the
shrinkage penalty overrides the variance reduction, making \(\hat{p}_1\)
superior. The confidence bands confirm this crossing is a true
mathematical property, not simulation variance.

\section{Hypothesis Testing}\label{hypothesis-testing}

A hypothesis test maps data to a binary decision regarding a null
hypothesis (\(H_0\)) versus an alternative (\(H_1\)). It relies on a
test statistic \(T(X)\), where extreme values serve as evidence against
\(H_0\).

\begin{itemize}
\tightlist
\item
  \textbf{Type I Error:} Rejecting a true \(H_0\).
\item
  \textbf{Type II Error:} Failing to reject a false \(H_0\).
\item
  \textbf{Power Function \(\beta(\theta)\):} The probability of
  rejecting \(H_0\). It represents the Type I error rate when
  \(\theta \in H_0\), and the statistical power when \(\theta \in H_1\).
\item
  \textbf{Size \& Level:} The maximum Type I error probability across
  all parameters in \(H_0\). A test has \emph{level \(\alpha\)} if its
  size is at most \(\alpha\).
\end{itemize}

The \textbf{p-value} repackages the test rule without locking into a
fixed \(\alpha\) in advance. It represents the probability of observing
a test statistic at least as extreme as the one observed, assuming
\(H_0\) is true:

\[pv = \Pr_{H_0}\left(T(X^{*}) \ge T(x^{obs})\right)\]

A functional p-value must maintain \textbf{validity}:
\(\Pr_{H_0}(pv \le \alpha) \le \alpha\) for every \(\alpha \in (0, 1)\).
If \(T\) possesses a continuous null distribution and is applied
exactly, this inequality becomes an equality, meaning the p-values will
follow a uniform distribution on \([0,1]\) under \(H_0\).

\textbf{Diagnostic:} Simulate multiple datasets under \(H_0\), compute
the p-value for each, and plot the histogram.

\begin{itemize}
\tightlist
\item
  \textbf{Flat histogram:} The test is valid.
\item
  \textbf{Left-skewed (mass near 0):} The test is invalid and liberal;
  true size exceeds nominal level.
\item
  \textbf{Right-skewed (mass near 1):} The test is excessively
  conservative; sacrifices statistical power.
\end{itemize}

P-values generally fail uniformity for two reasons:

\begin{enumerate}
\def\labelenumi{\arabic{enumi}.}
\tightlist
\item
  \textbf{Distributional Misspecification:} A foundational assumption
  (e.g., normality) fails, skewing the null distribution.
\item
  \textbf{Optimization Bias (Data-dependent procedures):} The null
  distribution is computed assuming a fixed procedure, but the actual
  procedure was chosen by inspecting the data (e.g., feature selection).
\end{enumerate}

\subsection{Permutation Tests}\label{permutation-tests}

Permutation tests offer a non-parametric alternative when analytical
assumptions are dubious. If the null hypothesis states two variables are
independent, the exact null distribution can be mapped without assuming
normality.

Because the variables are independent under \(H_0\), their labels can be
scrambled without altering the joint distribution (exchangeability). The
true p-value is simply the rank of the observed statistic among all
permuted statistics:

\[\widehat{pv} = \frac{1 + \#\{j : T^{(j)} \ge T^{obs}\}}{N + 1}\]

The \(+1\) correction counts the observed dataset alongside the \(N\)
permutations, ensuring \(\widehat{pv}\) remains mathematically valid for
any finite \(N\).

\subsection{Example: The t-test on Non-Normal
Observations}\label{example-the-t-test-on-non-normal-observations}

The standard Student's t-test assumes normally distributed data. Here,
we test its robustness when exposed to heavy-tailed distributions. This
highlights the first failure mode: applying the wrong null distribution
to the test statistic.

We simulate data using a t-distribution across varying degrees of
freedom (df). Lower df creates heavier tails, with df=1 representing the
Cauchy distribution.

\begin{Shaded}
\begin{Highlighting}[]
\NormalTok{t.test.pvalue }\OtherTok{\textless{}{-}} \ControlFlowTok{function}\NormalTok{ (x, }\AttributeTok{mu =} \DecValTok{0}\NormalTok{)}
\NormalTok{\{}
    \ControlFlowTok{if}\NormalTok{ (}\SpecialCharTok{!}\FunctionTok{is.numeric}\NormalTok{(x) }\SpecialCharTok{||} \SpecialCharTok{!}\FunctionTok{is.numeric}\NormalTok{(mu) }\SpecialCharTok{||} \FunctionTok{length}\NormalTok{(mu) }\SpecialCharTok{!=} \DecValTok{1}\NormalTok{) \{ }\FunctionTok{stop}\NormalTok{(}\StringTok{"Invalid argument"}\NormalTok{) \}}
\NormalTok{    n }\OtherTok{\textless{}{-}} \FunctionTok{length}\NormalTok{(x)}
    \ControlFlowTok{if}\NormalTok{ (n }\SpecialCharTok{\textless{}} \DecValTok{2}\NormalTok{) \{ }\FunctionTok{stop}\NormalTok{(}\StringTok{"Can\textquotesingle{}t do a t test with less than two data points"}\NormalTok{) \}}
    
\NormalTok{    t }\OtherTok{\textless{}{-}}\NormalTok{ (}\FunctionTok{mean}\NormalTok{(x) }\SpecialCharTok{{-}}\NormalTok{ mu) }\SpecialCharTok{/} \FunctionTok{sqrt}\NormalTok{(}\FunctionTok{var}\NormalTok{(x) }\SpecialCharTok{/}\NormalTok{ n)}
    \DecValTok{2} \SpecialCharTok{*} \FunctionTok{pt}\NormalTok{ (}\SpecialCharTok{{-}}\FunctionTok{abs}\NormalTok{(t), n }\SpecialCharTok{{-}} \DecValTok{1}\NormalTok{)}
\NormalTok{\}}

\NormalTok{t.test.sim }\OtherTok{\textless{}{-}} \ControlFlowTok{function}\NormalTok{ (N, n, df, }\AttributeTok{mu.true =} \DecValTok{0}\NormalTok{, }\AttributeTok{mu =} \DecValTok{0}\NormalTok{)}
\NormalTok{\{}
\NormalTok{    X }\OtherTok{\textless{}{-}} \FunctionTok{matrix}\NormalTok{(}\FunctionTok{rt}\NormalTok{(n }\SpecialCharTok{*}\NormalTok{ N, df), n, N) }\SpecialCharTok{+}\NormalTok{ mu.true}
\NormalTok{    m }\OtherTok{\textless{}{-}} \FunctionTok{colMeans}\NormalTok{(X)}
\NormalTok{    v }\OtherTok{\textless{}{-}} \FunctionTok{colSums}\NormalTok{((X }\SpecialCharTok{{-}} \FunctionTok{rep}\NormalTok{(m, }\AttributeTok{each =}\NormalTok{ n))}\SpecialCharTok{\^{}}\DecValTok{2}\NormalTok{) }\SpecialCharTok{/}\NormalTok{ (n }\SpecialCharTok{{-}} \DecValTok{1}\NormalTok{)}
    \DecValTok{2} \SpecialCharTok{*} \FunctionTok{pt}\NormalTok{(}\SpecialCharTok{{-}}\FunctionTok{abs}\NormalTok{((m }\SpecialCharTok{{-}}\NormalTok{ mu) }\SpecialCharTok{/} \FunctionTok{sqrt}\NormalTok{(v }\SpecialCharTok{/}\NormalTok{ n)), n }\SpecialCharTok{{-}} \DecValTok{1}\NormalTok{)}
\NormalTok{\}}

\NormalTok{hist.pv }\OtherTok{\textless{}{-}} \ControlFlowTok{function}\NormalTok{ (pv, n, df, }\AttributeTok{bins =} \DecValTok{20}\NormalTok{)}
\NormalTok{\{}
    \FunctionTok{hist}\NormalTok{(pv, }\AttributeTok{nclass =}\NormalTok{ bins, }\AttributeTok{prob =} \ConstantTok{TRUE}\NormalTok{, }\AttributeTok{col =} \StringTok{"grey90"}\NormalTok{,}
         \AttributeTok{xlab =} \FunctionTok{TeX}\NormalTok{(}\FunctionTok{sprintf}\NormalTok{(r}\StringTok{"($\%d$ p{-}values, $n = \%d$, $}\SpecialCharTok{\textbackslash{}n}\StringTok{u = \%d$)"}\NormalTok{, }\FunctionTok{length}\NormalTok{(pv), n, df)),}
         \AttributeTok{main =} \FunctionTok{TeX}\NormalTok{(}\FunctionTok{sprintf}\NormalTok{(r}\StringTok{"(size at $}\SpecialCharTok{\textbackslash{}a}\StringTok{lpha = 0.05$: $\%.3f$)"}\NormalTok{, }\FunctionTok{mean}\NormalTok{(pv }\SpecialCharTok{\textless{}} \FloatTok{0.05}\NormalTok{))))}
    \FunctionTok{abline}\NormalTok{(}\AttributeTok{h =} \DecValTok{1}\NormalTok{, }\AttributeTok{col =} \DecValTok{4}\NormalTok{, }\AttributeTok{lwd =} \DecValTok{2}\NormalTok{)}
\NormalTok{\}}
\end{Highlighting}
\end{Shaded}

\subsubsection{Null Distribution
Diagnostics}\label{null-distribution-diagnostics}

We expect these histograms to sit perfectly flat against the blue
uniform density line.

\begin{Shaded}
\begin{Highlighting}[]
\NormalTok{N }\OtherTok{\textless{}{-}} \DecValTok{2000}
\NormalTok{bins }\OtherTok{\textless{}{-}} \DecValTok{20}
\end{Highlighting}
\end{Shaded}

\textbf{100 Degrees of Freedom (Nearly Normal)} With 100 df, the
distribution is practically Normal. The p-values form a perfect uniform
line, confirming validity.

\begin{Shaded}
\begin{Highlighting}[]
\FunctionTok{par}\NormalTok{ (}\AttributeTok{mfrow =} \FunctionTok{c}\NormalTok{(}\DecValTok{1}\NormalTok{, }\DecValTok{2}\NormalTok{), }\AttributeTok{mar =} \FunctionTok{c}\NormalTok{(}\DecValTok{4}\NormalTok{, }\DecValTok{4}\NormalTok{, }\FloatTok{2.5}\NormalTok{, }\DecValTok{1}\NormalTok{))}
\FunctionTok{hist.pv}\NormalTok{(}\FunctionTok{t.test.sim}\NormalTok{(N, }\DecValTok{2}\NormalTok{, }\DecValTok{100}\NormalTok{), }\AttributeTok{n =} \DecValTok{2}\NormalTok{, }\AttributeTok{df =} \DecValTok{100}\NormalTok{, bins)}
\FunctionTok{hist.pv}\NormalTok{(}\FunctionTok{t.test.sim}\NormalTok{(N, }\DecValTok{20}\NormalTok{, }\DecValTok{100}\NormalTok{), }\AttributeTok{n =} \DecValTok{20}\NormalTok{, }\AttributeTok{df =} \DecValTok{100}\NormalTok{, bins)}
\end{Highlighting}
\end{Shaded}

\pandocbounded{\includegraphics[keepaspectratio]{chapters/unit04_simulation_files/figure-pdf/t-test-df100-1.pdf}}

\textbf{3 Degrees of Freedom (Heavy Tails)} Heavier tails warp the
p-value distribution. The nominal Type I error rate is no longer
perfectly calibrated, leading to false confidence.

\begin{Shaded}
\begin{Highlighting}[]
\FunctionTok{par}\NormalTok{ (}\AttributeTok{mfrow =} \FunctionTok{c}\NormalTok{(}\DecValTok{1}\NormalTok{, }\DecValTok{2}\NormalTok{), }\AttributeTok{mar =} \FunctionTok{c}\NormalTok{(}\DecValTok{4}\NormalTok{, }\DecValTok{4}\NormalTok{, }\FloatTok{2.5}\NormalTok{, }\DecValTok{1}\NormalTok{))}
\FunctionTok{hist.pv}\NormalTok{(}\FunctionTok{t.test.sim}\NormalTok{(N, }\DecValTok{2}\NormalTok{, }\DecValTok{3}\NormalTok{), }\AttributeTok{n =} \DecValTok{2}\NormalTok{, }\AttributeTok{df =} \DecValTok{3}\NormalTok{, bins)}
\FunctionTok{hist.pv}\NormalTok{(}\FunctionTok{t.test.sim}\NormalTok{(N, }\DecValTok{20}\NormalTok{, }\DecValTok{3}\NormalTok{), }\AttributeTok{n =} \DecValTok{20}\NormalTok{, }\AttributeTok{df =} \DecValTok{3}\NormalTok{, bins)}
\end{Highlighting}
\end{Shaded}

\pandocbounded{\includegraphics[keepaspectratio]{chapters/unit04_simulation_files/figure-pdf/t-test-df3-1.pdf}}

\textbf{1 Degree of Freedom (Cauchy Distribution)} With an undefined
mean and variance, extreme outliers dominate. The t-test entirely breaks
down, clustering p-values near 1.0 and rendering the test virtually
useless.

\begin{Shaded}
\begin{Highlighting}[]
\FunctionTok{par}\NormalTok{ (}\AttributeTok{mfrow =} \FunctionTok{c}\NormalTok{(}\DecValTok{1}\NormalTok{, }\DecValTok{2}\NormalTok{), }\AttributeTok{mar =} \FunctionTok{c}\NormalTok{(}\DecValTok{4}\NormalTok{, }\DecValTok{4}\NormalTok{, }\FloatTok{2.5}\NormalTok{, }\DecValTok{1}\NormalTok{))}
\FunctionTok{hist.pv}\NormalTok{(}\FunctionTok{t.test.sim}\NormalTok{(N, }\DecValTok{2}\NormalTok{, }\DecValTok{1}\NormalTok{), }\AttributeTok{n =} \DecValTok{2}\NormalTok{, }\AttributeTok{df =} \DecValTok{1}\NormalTok{, bins)}
\FunctionTok{hist.pv}\NormalTok{(}\FunctionTok{t.test.sim}\NormalTok{(N, }\DecValTok{20}\NormalTok{, }\DecValTok{1}\NormalTok{), }\AttributeTok{n =} \DecValTok{20}\NormalTok{, }\AttributeTok{df =} \DecValTok{1}\NormalTok{, bins)}
\end{Highlighting}
\end{Shaded}

\pandocbounded{\includegraphics[keepaspectratio]{chapters/unit04_simulation_files/figure-pdf/t-test-df1-1.pdf}}

\textbf{Increasing Sample Size (\(n=30\))} As sample size scales, the
Central Limit Theorem rescues the underlying math. The p-value
distributions for df=5 and df=10 flatten back into uniformity.

\begin{Shaded}
\begin{Highlighting}[]
\FunctionTok{par}\NormalTok{ (}\AttributeTok{mfrow =} \FunctionTok{c}\NormalTok{(}\DecValTok{1}\NormalTok{, }\DecValTok{2}\NormalTok{), }\AttributeTok{mar =} \FunctionTok{c}\NormalTok{(}\DecValTok{4}\NormalTok{, }\DecValTok{4}\NormalTok{, }\FloatTok{2.5}\NormalTok{, }\DecValTok{1}\NormalTok{))}
\FunctionTok{hist.pv}\NormalTok{(}\FunctionTok{t.test.sim}\NormalTok{(N, }\DecValTok{30}\NormalTok{, }\DecValTok{10}\NormalTok{), }\AttributeTok{n =} \DecValTok{30}\NormalTok{, }\AttributeTok{df =} \DecValTok{10}\NormalTok{, bins)}
\FunctionTok{hist.pv}\NormalTok{(}\FunctionTok{t.test.sim}\NormalTok{(N, }\DecValTok{30}\NormalTok{, }\DecValTok{5}\NormalTok{), }\AttributeTok{n =} \DecValTok{30}\NormalTok{, }\AttributeTok{df =} \DecValTok{5}\NormalTok{, bins)}
\end{Highlighting}
\end{Shaded}

\pandocbounded{\includegraphics[keepaspectratio]{chapters/unit04_simulation_files/figure-pdf/t-test-n30-1.pdf}}

\subsubsection{Power of the t-test}\label{power-of-the-t-test}

Statistical power represents the ability to successfully detect a true
effect. Here, we inject a true mean of \(\mu = 0.5\) while testing
against \(H_0: \mu = 0\).

\begin{Shaded}
\begin{Highlighting}[]
\NormalTok{hist.power }\OtherTok{\textless{}{-}} \ControlFlowTok{function}\NormalTok{ (pv, n, df, mu.true)}
    \FunctionTok{hist}\NormalTok{(pv, }\AttributeTok{nclass =} \DecValTok{20}\NormalTok{, }\AttributeTok{col =} \StringTok{"grey90"}\NormalTok{,}
         \AttributeTok{xlab =} \FunctionTok{TeX}\NormalTok{(}\FunctionTok{sprintf}\NormalTok{(r}\StringTok{"($\%d$ p{-}values, $n = \%d$, $}\SpecialCharTok{\textbackslash{}n}\StringTok{u = \%d$, $\textbackslash{}mu = \%.1f$)"}\NormalTok{,}
                            \FunctionTok{length}\NormalTok{(pv), n, df, mu.true)),}
         \AttributeTok{main =} \FunctionTok{TeX}\NormalTok{(}\FunctionTok{sprintf}\NormalTok{(r}\StringTok{"(power at $}\SpecialCharTok{\textbackslash{}a}\StringTok{lpha = 0.05$: $\%.3f$)"}\NormalTok{, }\FunctionTok{mean}\NormalTok{(pv }\SpecialCharTok{\textless{}} \FloatTok{0.05}\NormalTok{))))}
\end{Highlighting}
\end{Shaded}

\begin{Shaded}
\begin{Highlighting}[]
\FunctionTok{par}\NormalTok{ (}\AttributeTok{mfrow =} \FunctionTok{c}\NormalTok{(}\DecValTok{1}\NormalTok{, }\DecValTok{2}\NormalTok{), }\AttributeTok{mar =} \FunctionTok{c}\NormalTok{(}\DecValTok{4}\NormalTok{, }\DecValTok{4}\NormalTok{, }\FloatTok{2.5}\NormalTok{, }\DecValTok{1}\NormalTok{))}
\FunctionTok{hist.power}\NormalTok{(}\FunctionTok{t.test.sim}\NormalTok{(N, }\DecValTok{30}\NormalTok{, }\DecValTok{5}\NormalTok{, }\AttributeTok{mu.true =} \FloatTok{0.5}\NormalTok{), }\AttributeTok{n =} \DecValTok{30}\NormalTok{, }\AttributeTok{df =} \DecValTok{5}\NormalTok{, }\AttributeTok{mu.true =} \FloatTok{0.5}\NormalTok{)}
\FunctionTok{hist.power}\NormalTok{(}\FunctionTok{t.test.sim}\NormalTok{(N, }\DecValTok{60}\NormalTok{, }\DecValTok{5}\NormalTok{, }\AttributeTok{mu.true =} \FloatTok{0.5}\NormalTok{), }\AttributeTok{n =} \DecValTok{60}\NormalTok{, }\AttributeTok{df =} \DecValTok{5}\NormalTok{, }\AttributeTok{mu.true =} \FloatTok{0.5}\NormalTok{)}
\end{Highlighting}
\end{Shaded}

\pandocbounded{\includegraphics[keepaspectratio]{chapters/unit04_simulation_files/figure-pdf/t-test-power-df5-1.pdf}}

\begin{Shaded}
\begin{Highlighting}[]
\FunctionTok{par}\NormalTok{ (}\AttributeTok{mfrow =} \FunctionTok{c}\NormalTok{(}\DecValTok{1}\NormalTok{, }\DecValTok{2}\NormalTok{), }\AttributeTok{mar =} \FunctionTok{c}\NormalTok{(}\DecValTok{4}\NormalTok{, }\DecValTok{4}\NormalTok{, }\FloatTok{2.5}\NormalTok{, }\DecValTok{1}\NormalTok{))}
\FunctionTok{hist.power}\NormalTok{(}\FunctionTok{t.test.sim}\NormalTok{(N, }\DecValTok{30}\NormalTok{, }\DecValTok{3}\NormalTok{, }\AttributeTok{mu.true =} \FloatTok{0.5}\NormalTok{), }\AttributeTok{n =} \DecValTok{30}\NormalTok{, }\AttributeTok{df =} \DecValTok{3}\NormalTok{, }\AttributeTok{mu.true =} \FloatTok{0.5}\NormalTok{)}
\FunctionTok{hist.power}\NormalTok{(}\FunctionTok{t.test.sim}\NormalTok{(N, }\DecValTok{60}\NormalTok{, }\DecValTok{3}\NormalTok{, }\AttributeTok{mu.true =} \FloatTok{0.5}\NormalTok{), }\AttributeTok{n =} \DecValTok{60}\NormalTok{, }\AttributeTok{df =} \DecValTok{3}\NormalTok{, }\AttributeTok{mu.true =} \FloatTok{0.5}\NormalTok{)}
\end{Highlighting}
\end{Shaded}

\pandocbounded{\includegraphics[keepaspectratio]{chapters/unit04_simulation_files/figure-pdf/t-test-power-df3-1.pdf}}

\subsection{Example: A Permutation Test for
Correlation}\label{example-a-permutation-test-for-correlation}

When evaluating the correlation between variables \(X\) and \(Y\),
dropping the analytical normality assumption in favor of a permutation
test avoids the first failure mode entirely.

\textbf{The Procedure:}

\begin{enumerate}
\def\labelenumi{\arabic{enumi}.}
\tightlist
\item
  Calculate the observed correlation \(\rho^{obs} = \hat\rho(X, Y)\).
\item
  For \(j = 1, \ldots, N\): randomly shuffle \(Y\) to create
  \(Y^{(j)}\), then calculate \(\rho^{(j)} = \hat\rho(X, Y^{(j)})\).
\item
  Calculate the permutation p-value as the proportion of shuffled
  correlations whose absolute value meets or exceeds the observed
  correlation.
\end{enumerate}

\begin{Shaded}
\begin{Highlighting}[]
\NormalTok{perm.cor.test }\OtherTok{\textless{}{-}} \ControlFlowTok{function}\NormalTok{ (x, y, }\AttributeTok{no.perm =} \DecValTok{100}\NormalTok{)}
\NormalTok{\{}
\NormalTok{    Yp }\OtherTok{\textless{}{-}} \FunctionTok{replicate}\NormalTok{(no.perm, }\FunctionTok{sample}\NormalTok{(y))          }\CommentTok{\# n by no.perm permuted responses}
\NormalTok{    permuted.abs.cor }\OtherTok{\textless{}{-}} \FunctionTok{abs}\NormalTok{(}\FunctionTok{as.vector}\NormalTok{(}\FunctionTok{cor}\NormalTok{(x, Yp)))}
\NormalTok{    real.cor }\OtherTok{\textless{}{-}} \FunctionTok{cor}\NormalTok{(x, y)}

    \FunctionTok{list}\NormalTok{ (}\AttributeTok{rho =}\NormalTok{ real.cor,}
          \AttributeTok{pv =}\NormalTok{ (}\DecValTok{1} \SpecialCharTok{+} \FunctionTok{sum}\NormalTok{(permuted.abs.cor }\SpecialCharTok{\textgreater{}=} \FunctionTok{abs}\NormalTok{(real.cor))) }\SpecialCharTok{/}\NormalTok{ (no.perm }\SpecialCharTok{+} \DecValTok{1}\NormalTok{),}
          \AttributeTok{perm.abs.cor =}\NormalTok{ permuted.abs.cor)}
\NormalTok{\}}
\end{Highlighting}
\end{Shaded}

\subsubsection{Performance on Non-Normal Dependent
Data}\label{performance-on-non-normal-dependent-data}

We generate non-normal exponential variables with a direct dependency.
The observed \(\vert{}\hat\rho\vert{}\) rests squarely in the far right
tail of the randomized permutation distribution.

\begin{Shaded}
\begin{Highlighting}[]
\FunctionTok{par}\NormalTok{ (}\AttributeTok{mfrow =} \FunctionTok{c}\NormalTok{(}\DecValTok{1}\NormalTok{, }\DecValTok{1}\NormalTok{), }\AttributeTok{mar =} \FunctionTok{c}\NormalTok{(}\DecValTok{4}\NormalTok{, }\DecValTok{4}\NormalTok{, }\DecValTok{3}\NormalTok{, }\DecValTok{1}\NormalTok{))}

\NormalTok{x }\OtherTok{\textless{}{-}} \FunctionTok{rexp}\NormalTok{(}\DecValTok{30}\NormalTok{); y }\OtherTok{\textless{}{-}} \FloatTok{0.4} \SpecialCharTok{*}\NormalTok{ x }\SpecialCharTok{+} \FunctionTok{rexp}\NormalTok{(}\DecValTok{30}\NormalTok{)     }\CommentTok{\# non{-}normal, dependent}
\NormalTok{out }\OtherTok{\textless{}{-}} \FunctionTok{perm.cor.test}\NormalTok{(x, y, }\AttributeTok{no.perm =} \DecValTok{2000}\NormalTok{)}

\FunctionTok{hist}\NormalTok{(out}\SpecialCharTok{$}\NormalTok{perm.abs.cor, }\AttributeTok{breaks =} \DecValTok{40}\NormalTok{, }\AttributeTok{col =} \StringTok{"grey90"}\NormalTok{,}
     \AttributeTok{xlab =} \FunctionTok{TeX}\NormalTok{(r}\StringTok{"($|\textbackslash{}hat\{}\SpecialCharTok{\textbackslash{}r}\StringTok{ho\}|$ from permuted data)"}\NormalTok{),}
     \AttributeTok{main =} \FunctionTok{TeX}\NormalTok{(}\FunctionTok{sprintf}\NormalTok{(r}\StringTok{"(Permutation distribution of $|\textbackslash{}hat\{}\SpecialCharTok{\textbackslash{}r}\StringTok{ho\}|$; $\textbackslash{}hat\{}\SpecialCharTok{\textbackslash{}r}\StringTok{ho\}\^{}\{obs\} = \%.3f$)"}\NormalTok{,}
\NormalTok{                        out}\SpecialCharTok{$}\NormalTok{rho)))}
\FunctionTok{abline}\NormalTok{(}\AttributeTok{v =} \FunctionTok{abs}\NormalTok{(out}\SpecialCharTok{$}\NormalTok{rho), }\AttributeTok{col =} \DecValTok{2}\NormalTok{, }\AttributeTok{lwd =} \DecValTok{2}\NormalTok{)}
\end{Highlighting}
\end{Shaded}

\pandocbounded{\includegraphics[keepaspectratio]{chapters/unit04_simulation_files/figure-pdf/perm-cor-illustration-1.pdf}}

\begin{Shaded}
\begin{Highlighting}[]
\FunctionTok{c}\NormalTok{(}\AttributeTok{rho\_obs =}\NormalTok{ out}\SpecialCharTok{$}\NormalTok{rho, }\AttributeTok{perm\_pv =}\NormalTok{ out}\SpecialCharTok{$}\NormalTok{pv, }\AttributeTok{pearson\_pv =} \FunctionTok{cor.test}\NormalTok{(x, y)}\SpecialCharTok{$}\NormalTok{p.value)}
\end{Highlighting}
\end{Shaded}

\begin{verbatim}
   rho_obs    perm_pv pearson_pv 
0.33310306 0.06396802 0.07206288 
\end{verbatim}

\subsubsection{Evaluating Null Validity}\label{evaluating-null-validity}

Under an identical, normal distribution, the permutation test perfectly
mirrors a standard analytical test, maintaining true uniformity under
\(H_0\).

\begin{Shaded}
\begin{Highlighting}[]
\NormalTok{test.perm.norm }\OtherTok{\textless{}{-}} \ControlFlowTok{function}\NormalTok{ (no.sim, no.perm, }\AttributeTok{n =} \DecValTok{20}\NormalTok{)}
    \FunctionTok{replicate}\NormalTok{(no.sim, }\FunctionTok{perm.cor.test}\NormalTok{(}\FunctionTok{rnorm}\NormalTok{(n), }\FunctionTok{rnorm}\NormalTok{(n), no.perm)}\SpecialCharTok{$}\NormalTok{pv)}

\NormalTok{pv }\OtherTok{\textless{}{-}} \FunctionTok{test.perm.norm}\NormalTok{(}\DecValTok{2000}\NormalTok{, }\DecValTok{100}\NormalTok{)}

\FunctionTok{par}\NormalTok{ (}\AttributeTok{mfrow =} \FunctionTok{c}\NormalTok{(}\DecValTok{1}\NormalTok{, }\DecValTok{1}\NormalTok{), }\AttributeTok{mar =} \FunctionTok{c}\NormalTok{(}\DecValTok{4}\NormalTok{, }\DecValTok{4}\NormalTok{, }\DecValTok{3}\NormalTok{, }\DecValTok{1}\NormalTok{))}
\FunctionTok{hist}\NormalTok{(pv, }\DecValTok{20}\NormalTok{, }\AttributeTok{prob =} \ConstantTok{TRUE}\NormalTok{, }\AttributeTok{col =} \StringTok{"grey90"}\NormalTok{,}
     \AttributeTok{xlab =} \FunctionTok{TeX}\NormalTok{(r}\StringTok{"($2000$ permutation p{-}values, normal data)"}\NormalTok{),}
     \AttributeTok{main =} \FunctionTok{TeX}\NormalTok{(}\FunctionTok{sprintf}\NormalTok{(r}\StringTok{"(size at $}\SpecialCharTok{\textbackslash{}a}\StringTok{lpha = 0.05$: $\%.3f$)"}\NormalTok{, }\FunctionTok{mean}\NormalTok{(pv }\SpecialCharTok{\textless{}} \FloatTok{0.05}\NormalTok{))))}
\FunctionTok{abline}\NormalTok{(}\AttributeTok{h =} \DecValTok{1}\NormalTok{, }\AttributeTok{col =} \DecValTok{4}\NormalTok{, }\AttributeTok{lwd =} \DecValTok{2}\NormalTok{)}
\end{Highlighting}
\end{Shaded}

\pandocbounded{\includegraphics[keepaspectratio]{chapters/unit04_simulation_files/figure-pdf/perm-cor-normal-h0-1.pdf}}

Crucially, when testing non-normal, entirely independent data, the
permutation test remains flawless. The standard Pearson test (relying on
a normal approximation) struggles, while the permutation test yields a
perfect uniform distribution without needing normality.

\begin{Shaded}
\begin{Highlighting}[]
\NormalTok{res }\OtherTok{\textless{}{-}} \FunctionTok{replicate}\NormalTok{(}\DecValTok{500}\NormalTok{, \{}
\NormalTok{  x }\OtherTok{\textless{}{-}} \FunctionTok{rexp}\NormalTok{(}\DecValTok{30}\NormalTok{); y }\OtherTok{\textless{}{-}} \FunctionTok{rexp}\NormalTok{(}\DecValTok{30}\NormalTok{)                    }\CommentTok{\# independent, non{-}normal}
  \FunctionTok{c}\NormalTok{(}\AttributeTok{perm    =} \FunctionTok{perm.cor.test}\NormalTok{(x, y, }\AttributeTok{no.perm =} \DecValTok{500}\NormalTok{)}\SpecialCharTok{$}\NormalTok{pv,}
    \AttributeTok{pearson =} \FunctionTok{cor.test}\NormalTok{(x, y)}\SpecialCharTok{$}\NormalTok{p.value)}
\NormalTok{\})}

\FunctionTok{par}\NormalTok{(}\AttributeTok{mfrow =} \FunctionTok{c}\NormalTok{(}\DecValTok{1}\NormalTok{, }\DecValTok{2}\NormalTok{), }\AttributeTok{mar =} \FunctionTok{c}\NormalTok{(}\DecValTok{4}\NormalTok{, }\DecValTok{4}\NormalTok{, }\DecValTok{3}\NormalTok{, }\DecValTok{1}\NormalTok{))}

\ControlFlowTok{for}\NormalTok{ (m }\ControlFlowTok{in} \FunctionTok{c}\NormalTok{(}\StringTok{"perm"}\NormalTok{, }\StringTok{"pearson"}\NormalTok{)) \{}
  \FunctionTok{hist}\NormalTok{(res[m, ], }\AttributeTok{breaks =} \DecValTok{20}\NormalTok{, }\AttributeTok{prob =} \ConstantTok{TRUE}\NormalTok{, }\AttributeTok{col =} \StringTok{"grey90"}\NormalTok{,}
       \AttributeTok{xlab =} \FunctionTok{TeX}\NormalTok{(}\FunctionTok{sprintf}\NormalTok{(r}\StringTok{"($500$ \%s p{-}values under $H\_0$)"}\NormalTok{, m)),}
       \AttributeTok{main =} \FunctionTok{TeX}\NormalTok{(}\FunctionTok{sprintf}\NormalTok{(r}\StringTok{"(size at $}\SpecialCharTok{\textbackslash{}a}\StringTok{lpha = 0.05$: $\%.3f$)"}\NormalTok{, }\FunctionTok{mean}\NormalTok{(res[m, ] }\SpecialCharTok{\textless{}} \FloatTok{0.05}\NormalTok{))))}
  \FunctionTok{abline}\NormalTok{(}\AttributeTok{h =} \DecValTok{1}\NormalTok{, }\AttributeTok{col =} \DecValTok{4}\NormalTok{, }\AttributeTok{lwd =} \DecValTok{2}\NormalTok{)}
\NormalTok{\}}
\end{Highlighting}
\end{Shaded}

\pandocbounded{\includegraphics[keepaspectratio]{chapters/unit04_simulation_files/figure-pdf/perm-cor-nonnormal-h0-1.pdf}}

\subsection{Example: Variable Selection and Optimization
Bias}\label{example-variable-selection-and-optimization-bias}

This section addresses the second failure mode: calculating a p-value as
if the model was fixed in advance, when the model was actually chosen by
actively querying the data (p-hacking).

\textbf{The Scenario:} You have a response \(Y\) and 20 candidate
predictors, \(X_1 \dots X_{20}\), with a sample size of \(n=50\). You
test the null hypothesis that \(Y\) has no relationship with any \(X\).

A deeply flawed but common practice is to:

\begin{enumerate}
\def\labelenumi{\arabic{enumi}.}
\tightlist
\item
  Scan for the single predictor \(X_{k^*}\) that happens to correlate
  highest with \(Y\).
\item
  Run a standard linear regression \(Y \sim X_{k^*}\).
\item
  Report the standard t-test p-value for the slope.
\end{enumerate}

\begin{Shaded}
\begin{Highlighting}[]
\NormalTok{test\_lm\_sel }\OtherTok{\textless{}{-}} \ControlFlowTok{function}\NormalTok{(y, X) \{}
\NormalTok{  k }\OtherTok{\textless{}{-}} \FunctionTok{which.max}\NormalTok{(}\FunctionTok{abs}\NormalTok{(}\FunctionTok{cor}\NormalTok{(X, y)))}
  \FunctionTok{summary}\NormalTok{(}\FunctionTok{lm}\NormalTok{(y }\SpecialCharTok{\textasciitilde{}}\NormalTok{ X[, k]))}\SpecialCharTok{$}\NormalTok{coefficients[}\DecValTok{2}\NormalTok{, }\DecValTok{4}\NormalTok{]    }\CommentTok{\# t{-}test p{-}value of selected slope}
\NormalTok{\}}
\end{Highlighting}
\end{Shaded}

\begin{Shaded}
\begin{Highlighting}[]
\DocumentationTok{\#\#\# two{-}sided p{-}value of the slope in a simple linear regression, vectorized in r}
\NormalTok{cor\_pvalue }\OtherTok{\textless{}{-}} \ControlFlowTok{function}\NormalTok{(r, n) }\DecValTok{2} \SpecialCharTok{*} \FunctionTok{pt}\NormalTok{(}\SpecialCharTok{{-}}\FunctionTok{abs}\NormalTok{(r }\SpecialCharTok{*} \FunctionTok{sqrt}\NormalTok{((n }\SpecialCharTok{{-}} \DecValTok{2}\NormalTok{) }\SpecialCharTok{/}\NormalTok{ (}\DecValTok{1} \SpecialCharTok{{-}}\NormalTok{ r}\SpecialCharTok{\^{}}\DecValTok{2}\NormalTok{))), n }\SpecialCharTok{{-}} \DecValTok{2}\NormalTok{)}
\NormalTok{test\_lm\_sel\_fast }\OtherTok{\textless{}{-}} \ControlFlowTok{function}\NormalTok{(y, X) }\FunctionTok{cor\_pvalue}\NormalTok{(}\FunctionTok{max}\NormalTok{(}\FunctionTok{abs}\NormalTok{(}\FunctionTok{cor}\NormalTok{(X, y))), }\FunctionTok{length}\NormalTok{(y))}
\NormalTok{sim\_null }\OtherTok{\textless{}{-}} \ControlFlowTok{function}\NormalTok{(}\AttributeTok{n =} \DecValTok{50}\NormalTok{, }\AttributeTok{p =} \DecValTok{20}\NormalTok{) }\FunctionTok{list}\NormalTok{(}\AttributeTok{y =} \FunctionTok{rnorm}\NormalTok{(n), }\AttributeTok{X =} \FunctionTok{matrix}\NormalTok{(}\FunctionTok{rnorm}\NormalTok{(n }\SpecialCharTok{*}\NormalTok{ p), n, p))}

\NormalTok{d }\OtherTok{\textless{}{-}} \FunctionTok{sim\_null}\NormalTok{()}
\FunctionTok{c}\NormalTok{(}\AttributeTok{lm =} \FunctionTok{test\_lm\_sel}\NormalTok{(d}\SpecialCharTok{$}\NormalTok{y, d}\SpecialCharTok{$}\NormalTok{X), }\AttributeTok{fast =} \FunctionTok{test\_lm\_sel\_fast}\NormalTok{(d}\SpecialCharTok{$}\NormalTok{y, d}\SpecialCharTok{$}\NormalTok{X))   }\CommentTok{\# identical}
\end{Highlighting}
\end{Shaded}

\begin{verbatim}
        lm       fast 
0.03412572 0.03412572 
\end{verbatim}

By intentionally fishing for the maximum correlation first, the naive
test is overwhelmingly likely to discover ``significant'' correlations
strictly out of random noise.

\begin{Shaded}
\begin{Highlighting}[]
\NormalTok{pv\_sel }\OtherTok{\textless{}{-}} \FunctionTok{replicate}\NormalTok{(}\DecValTok{2000}\NormalTok{, \{ d }\OtherTok{\textless{}{-}} \FunctionTok{sim\_null}\NormalTok{(); }\FunctionTok{test\_lm\_sel\_fast}\NormalTok{(d}\SpecialCharTok{$}\NormalTok{y, d}\SpecialCharTok{$}\NormalTok{X) \})}

\FunctionTok{par}\NormalTok{ (}\AttributeTok{mfrow =} \FunctionTok{c}\NormalTok{(}\DecValTok{1}\NormalTok{, }\DecValTok{1}\NormalTok{), }\AttributeTok{mar =} \FunctionTok{c}\NormalTok{(}\DecValTok{4}\NormalTok{, }\DecValTok{4}\NormalTok{, }\DecValTok{3}\NormalTok{, }\DecValTok{1}\NormalTok{))}
\FunctionTok{hist}\NormalTok{(pv\_sel, }\AttributeTok{breaks =} \DecValTok{20}\NormalTok{, }\AttributeTok{prob =} \ConstantTok{TRUE}\NormalTok{, }\AttributeTok{col =} \StringTok{"grey90"}\NormalTok{,}
     \AttributeTok{xlab =} \FunctionTok{TeX}\NormalTok{(r}\StringTok{"($2000$ naive p{-}values after selection, $H\_0$ true)"}\NormalTok{),}
     \AttributeTok{main =} \FunctionTok{TeX}\NormalTok{(}\FunctionTok{sprintf}\NormalTok{(r}\StringTok{"(size at $}\SpecialCharTok{\textbackslash{}a}\StringTok{lpha = 0.05$: $\%.3f$)"}\NormalTok{, }\FunctionTok{mean}\NormalTok{(pv\_sel }\SpecialCharTok{\textless{}} \FloatTok{0.05}\NormalTok{))))}
\FunctionTok{abline}\NormalTok{(}\AttributeTok{h =} \DecValTok{1}\NormalTok{, }\AttributeTok{col =} \DecValTok{4}\NormalTok{, }\AttributeTok{lwd =} \DecValTok{2}\NormalTok{)}
\end{Highlighting}
\end{Shaded}

\pandocbounded{\includegraphics[keepaspectratio]{chapters/unit04_simulation_files/figure-pdf/test-lm-sel-h0-1.pdf}}

As the histogram proves, the false-positive rate under this naive
approach explodes far above the accepted 0.05 threshold.

\subsubsection{The Permutation Fix}\label{the-permutation-fix}

To repair the p-value, we must embed the \emph{entire selection process}
inside the permutation loop.

\begin{enumerate}
\def\labelenumi{\arabic{enumi}.}
\tightlist
\item
  Record the naive p-value from the observed data.
\item
  Shuffle \(Y\), re-select the best predictor, and compute a new
  p-value. Repeat \(N\) times.
\item
  The true valid p-value is the proportion of permuted runs that beat
  the original observed run.
\end{enumerate}

\begin{Shaded}
\begin{Highlighting}[]
\NormalTok{perm\_lm\_sel }\OtherTok{\textless{}{-}} \ControlFlowTok{function}\NormalTok{(y, X, }\AttributeTok{N =} \DecValTok{500}\NormalTok{) \{}
\NormalTok{  n }\OtherTok{\textless{}{-}} \FunctionTok{length}\NormalTok{(y)}
\NormalTok{  pobs }\OtherTok{\textless{}{-}} \FunctionTok{test\_lm\_sel\_fast}\NormalTok{(y, X)}

\NormalTok{  Yp }\OtherTok{\textless{}{-}} \FunctionTok{replicate}\NormalTok{(N, }\FunctionTok{sample}\NormalTok{(y))            }\CommentTok{\# n by N permuted responses}
\NormalTok{  rmax }\OtherTok{\textless{}{-}} \FunctionTok{apply}\NormalTok{(}\FunctionTok{abs}\NormalTok{(}\FunctionTok{cor}\NormalTok{(X, Yp)), }\DecValTok{2}\NormalTok{, max)   }\CommentTok{\# selection redone in each column}
\NormalTok{  pperm }\OtherTok{\textless{}{-}} \FunctionTok{cor\_pvalue}\NormalTok{(rmax, n)}

\NormalTok{  (}\DecValTok{1} \SpecialCharTok{+} \FunctionTok{sum}\NormalTok{(pperm }\SpecialCharTok{\textless{}=}\NormalTok{ pobs)) }\SpecialCharTok{/}\NormalTok{ (N }\SpecialCharTok{+} \DecValTok{1}\NormalTok{)}
\NormalTok{\}}
\end{Highlighting}
\end{Shaded}

\begin{Shaded}
\begin{Highlighting}[]
\NormalTok{sim\_alt }\OtherTok{\textless{}{-}} \ControlFlowTok{function}\NormalTok{(}\AttributeTok{n =} \DecValTok{50}\NormalTok{, }\AttributeTok{p =} \DecValTok{20}\NormalTok{, }\AttributeTok{beta =} \FloatTok{0.5}\NormalTok{) \{}
\NormalTok{  X }\OtherTok{\textless{}{-}} \FunctionTok{matrix}\NormalTok{(}\FunctionTok{rnorm}\NormalTok{(n }\SpecialCharTok{*}\NormalTok{ p), n, p); }\FunctionTok{list}\NormalTok{(}\AttributeTok{y =}\NormalTok{ beta }\SpecialCharTok{*}\NormalTok{ X[, }\DecValTok{1}\NormalTok{] }\SpecialCharTok{+} \FunctionTok{rnorm}\NormalTok{(n), }\AttributeTok{X =}\NormalTok{ X)}
\NormalTok{\}}

\NormalTok{pv\_perm0 }\OtherTok{\textless{}{-}} \FunctionTok{replicate}\NormalTok{(}\DecValTok{300}\NormalTok{, \{ d }\OtherTok{\textless{}{-}} \FunctionTok{sim\_null}\NormalTok{(); }\FunctionTok{perm\_lm\_sel}\NormalTok{(d}\SpecialCharTok{$}\NormalTok{y, d}\SpecialCharTok{$}\NormalTok{X) \})}
\NormalTok{pv\_perm1 }\OtherTok{\textless{}{-}} \FunctionTok{replicate}\NormalTok{(}\DecValTok{300}\NormalTok{, \{ d }\OtherTok{\textless{}{-}} \FunctionTok{sim\_alt}\NormalTok{();  }\FunctionTok{perm\_lm\_sel}\NormalTok{(d}\SpecialCharTok{$}\NormalTok{y, d}\SpecialCharTok{$}\NormalTok{X) \})}
\NormalTok{pv\_sel1  }\OtherTok{\textless{}{-}} \FunctionTok{replicate}\NormalTok{(}\DecValTok{300}\NormalTok{, \{ d }\OtherTok{\textless{}{-}} \FunctionTok{sim\_alt}\NormalTok{();  }\FunctionTok{test\_lm\_sel\_fast}\NormalTok{(d}\SpecialCharTok{$}\NormalTok{y, d}\SpecialCharTok{$}\NormalTok{X) \})}

\FunctionTok{par}\NormalTok{ (}\AttributeTok{mfrow =} \FunctionTok{c}\NormalTok{(}\DecValTok{1}\NormalTok{, }\DecValTok{1}\NormalTok{), }\AttributeTok{mar =} \FunctionTok{c}\NormalTok{(}\DecValTok{4}\NormalTok{, }\DecValTok{4}\NormalTok{, }\DecValTok{3}\NormalTok{, }\DecValTok{1}\NormalTok{))}
\FunctionTok{hist}\NormalTok{(pv\_perm0, }\AttributeTok{breaks =} \DecValTok{20}\NormalTok{, }\AttributeTok{prob =} \ConstantTok{TRUE}\NormalTok{, }\AttributeTok{col =} \StringTok{"grey90"}\NormalTok{,}
     \AttributeTok{xlab =} \FunctionTok{TeX}\NormalTok{(r}\StringTok{"($300$ permutation p{-}values, $H\_0$ true)"}\NormalTok{),}
     \AttributeTok{main =} \FunctionTok{TeX}\NormalTok{(}\FunctionTok{sprintf}\NormalTok{(r}\StringTok{"(size at $}\SpecialCharTok{\textbackslash{}a}\StringTok{lpha = 0.05$: $\%.3f$)"}\NormalTok{, }\FunctionTok{mean}\NormalTok{(pv\_perm0 }\SpecialCharTok{\textless{}} \FloatTok{0.05}\NormalTok{))))}
\FunctionTok{abline}\NormalTok{(}\AttributeTok{h =} \DecValTok{1}\NormalTok{, }\AttributeTok{col =} \DecValTok{4}\NormalTok{, }\AttributeTok{lwd =} \DecValTok{2}\NormalTok{)}
\end{Highlighting}
\end{Shaded}

\pandocbounded{\includegraphics[keepaspectratio]{chapters/unit04_simulation_files/figure-pdf/perm-lm-sel-size-power-1.pdf}}

\begin{Shaded}
\begin{Highlighting}[]
\FunctionTok{rbind}\NormalTok{(}\AttributeTok{test\_lm\_sel =} \FunctionTok{c}\NormalTok{(}\AttributeTok{size =} \FunctionTok{mean}\NormalTok{(pv\_sel }\SpecialCharTok{\textless{}} \FloatTok{0.05}\NormalTok{),   }\AttributeTok{power =} \FunctionTok{mean}\NormalTok{(pv\_sel1 }\SpecialCharTok{\textless{}} \FloatTok{0.05}\NormalTok{)),}
      \AttributeTok{permutation =} \FunctionTok{c}\NormalTok{(}\AttributeTok{size =} \FunctionTok{mean}\NormalTok{(pv\_perm0 }\SpecialCharTok{\textless{}} \FloatTok{0.05}\NormalTok{), }\AttributeTok{power =} \FunctionTok{mean}\NormalTok{(pv\_perm1 }\SpecialCharTok{\textless{}} \FloatTok{0.05}\NormalTok{)))}
\end{Highlighting}
\end{Shaded}

\begin{verbatim}
                  size     power
test_lm_sel 0.65700000 0.9733333
permutation 0.06666667 0.6033333
\end{verbatim}

The permutation distribution restores the flat, uniform distribution.
Because power can only be honestly compared between tests of the same
size, the permutation test provides the mathematically sound baseline.

\subsection{Example: A Power Regression
Test}\label{example-a-power-regression-test}

The same optimization bias occurs when iterating through continuous
power transformations (e.g., \(X^{0.5}, X^2\)) to maximize correlation
between \(X^q\) and \(Y\). Because the selection rule adapts to the
data, applying standard analytical p-values directly to the final
``optimized'' model will yield artificially inflated significance.

We employ the same fix: encapsulate the entire search process inside the
permutation randomization.

\begin{Shaded}
\begin{Highlighting}[]
\NormalTok{pow.matrix }\OtherTok{\textless{}{-}} \ControlFlowTok{function}\NormalTok{ (x, }\AttributeTok{powers =} \FunctionTok{seq}\NormalTok{(}\FloatTok{0.1}\NormalTok{, }\FloatTok{4.0}\NormalTok{, }\AttributeTok{by =} \FloatTok{0.1}\NormalTok{))}
\NormalTok{\{}
    \ControlFlowTok{if}\NormalTok{ (}\SpecialCharTok{!}\FunctionTok{is.vector}\NormalTok{(x, }\StringTok{"numeric"}\NormalTok{)) }\FunctionTok{stop}\NormalTok{(}\StringTok{"Argument should be a numeric vector"}\NormalTok{)}
    \ControlFlowTok{if}\NormalTok{ (}\FunctionTok{any}\NormalTok{(x }\SpecialCharTok{\textless{}} \DecValTok{0}\NormalTok{)) }\FunctionTok{stop}\NormalTok{(}\StringTok{"Predictors must be non{-}negative"}\NormalTok{)}
    \FunctionTok{outer}\NormalTok{(x, powers, }\StringTok{"\^{}"}\NormalTok{)}
\NormalTok{\}}

\NormalTok{pvalue.lm.pow }\OtherTok{\textless{}{-}} \ControlFlowTok{function}\NormalTok{ (Y, Xp)}
\NormalTok{\{}
\NormalTok{    Y }\OtherTok{\textless{}{-}} \FunctionTok{as.matrix}\NormalTok{(Y)}
\NormalTok{    rmax }\OtherTok{\textless{}{-}} \FunctionTok{apply}\NormalTok{(}\FunctionTok{abs}\NormalTok{(}\FunctionTok{cor}\NormalTok{(Xp, Y)), }\DecValTok{2}\NormalTok{, max)}
    \FunctionTok{cor\_pvalue}\NormalTok{(rmax, }\FunctionTok{nrow}\NormalTok{(Y))}
\NormalTok{\}}

\NormalTok{test.lm.power }\OtherTok{\textless{}{-}} \ControlFlowTok{function}\NormalTok{ (Xp, N)}
    \FunctionTok{pvalue.lm.pow}\NormalTok{(}\FunctionTok{matrix}\NormalTok{(}\FunctionTok{rnorm}\NormalTok{(}\FunctionTok{nrow}\NormalTok{(Xp) }\SpecialCharTok{*}\NormalTok{ N), }\FunctionTok{nrow}\NormalTok{(Xp), N), Xp)}

\NormalTok{pvalue.perm.lm.pow }\OtherTok{\textless{}{-}} \ControlFlowTok{function}\NormalTok{ (y, Xp, }\AttributeTok{no.perm =} \DecValTok{100}\NormalTok{)}
\NormalTok{\{}
\NormalTok{    actual }\OtherTok{\textless{}{-}} \FunctionTok{pvalue.lm.pow}\NormalTok{(y, Xp)}
\NormalTok{    permuted }\OtherTok{\textless{}{-}} \FunctionTok{pvalue.lm.pow}\NormalTok{(}\FunctionTok{replicate}\NormalTok{(no.perm, }\FunctionTok{sample}\NormalTok{(y)), Xp)}
\NormalTok{    (}\DecValTok{1} \SpecialCharTok{+} \FunctionTok{sum}\NormalTok{(permuted }\SpecialCharTok{\textless{}=}\NormalTok{ actual)) }\SpecialCharTok{/}\NormalTok{ (no.perm }\SpecialCharTok{+} \DecValTok{1}\NormalTok{)}
\NormalTok{\}}

\NormalTok{test.perm.power }\OtherTok{\textless{}{-}} \ControlFlowTok{function}\NormalTok{ (x, Xp, power.sim, }\AttributeTok{N =} \DecValTok{100}\NormalTok{, }\AttributeTok{resid.std.dev =} \DecValTok{1}\NormalTok{, }\AttributeTok{no.perm =} \DecValTok{99}\NormalTok{)}
\NormalTok{\{}
\NormalTok{    n }\OtherTok{\textless{}{-}} \FunctionTok{length}\NormalTok{(x)}
    \FunctionTok{vapply}\NormalTok{(}\DecValTok{1}\SpecialCharTok{:}\NormalTok{N,}
           \ControlFlowTok{function}\NormalTok{(k) }\FunctionTok{pvalue.perm.lm.pow}\NormalTok{(x}\SpecialCharTok{\^{}}\NormalTok{power.sim }\SpecialCharTok{+} \FunctionTok{rnorm}\NormalTok{(n, }\DecValTok{0}\NormalTok{, resid.std.dev),}
\NormalTok{                                          Xp, no.perm),}
           \FunctionTok{numeric}\NormalTok{(}\DecValTok{1}\NormalTok{))}
\NormalTok{\}}
\end{Highlighting}
\end{Shaded}

\subsubsection{Diagnostic Uniformity and Statistical
Power}\label{diagnostic-uniformity-and-statistical-power}

As observed previously, the naive approach causes a massive
false-positive cascade (left). The permutation algorithm absorbs the
optimization phase, correcting the bias and yielding perfectly valid,
uniform p-values (right).

\begin{Shaded}
\begin{Highlighting}[]
\NormalTok{x }\OtherTok{\textless{}{-}} \FunctionTok{exp}\NormalTok{(}\FunctionTok{rnorm}\NormalTok{(}\DecValTok{20}\NormalTok{))}
\NormalTok{Xp }\OtherTok{\textless{}{-}} \FunctionTok{pow.matrix}\NormalTok{(x)}

\NormalTok{naive.pv }\OtherTok{\textless{}{-}} \FunctionTok{test.lm.power}\NormalTok{(Xp, }\DecValTok{1000}\NormalTok{)}
\NormalTok{perm.pv }\OtherTok{\textless{}{-}} \FunctionTok{test.perm.power}\NormalTok{(x, Xp, }\AttributeTok{power.sim =} \DecValTok{0}\NormalTok{, }\AttributeTok{N =} \DecValTok{1000}\NormalTok{, }\AttributeTok{no.perm =} \DecValTok{99}\NormalTok{)}

\FunctionTok{par}\NormalTok{(}\AttributeTok{mfrow =} \FunctionTok{c}\NormalTok{(}\DecValTok{1}\NormalTok{, }\DecValTok{2}\NormalTok{), }\AttributeTok{mar =} \FunctionTok{c}\NormalTok{(}\DecValTok{4}\NormalTok{, }\DecValTok{4}\NormalTok{, }\DecValTok{3}\NormalTok{, }\DecValTok{1}\NormalTok{))}

\FunctionTok{hist}\NormalTok{(naive.pv, }\AttributeTok{nclass =} \DecValTok{20}\NormalTok{, }\AttributeTok{prob =} \ConstantTok{TRUE}\NormalTok{, }\AttributeTok{col =} \StringTok{"grey90"}\NormalTok{,}
     \AttributeTok{xlab =} \FunctionTok{TeX}\NormalTok{(r}\StringTok{"($1000$ naive p{-}values, $H\_0$ true)"}\NormalTok{),}
     \AttributeTok{main =} \FunctionTok{TeX}\NormalTok{(}\FunctionTok{sprintf}\NormalTok{(r}\StringTok{"(naive test, size $= \%.3f$)"}\NormalTok{, }\FunctionTok{mean}\NormalTok{(naive.pv }\SpecialCharTok{\textless{}} \FloatTok{0.05}\NormalTok{))))}
\FunctionTok{abline}\NormalTok{(}\AttributeTok{h =} \DecValTok{1}\NormalTok{, }\AttributeTok{col =} \DecValTok{4}\NormalTok{, }\AttributeTok{lwd =} \DecValTok{2}\NormalTok{)}

\FunctionTok{hist}\NormalTok{(perm.pv, }\AttributeTok{nclass =} \DecValTok{20}\NormalTok{, }\AttributeTok{prob =} \ConstantTok{TRUE}\NormalTok{, }\AttributeTok{col =} \StringTok{"grey90"}\NormalTok{,}
     \AttributeTok{xlab =} \FunctionTok{TeX}\NormalTok{(r}\StringTok{"($1000$ permutation p{-}values, $H\_0$ true)"}\NormalTok{),}
     \AttributeTok{main =} \FunctionTok{TeX}\NormalTok{(}\FunctionTok{sprintf}\NormalTok{(r}\StringTok{"(permutation test, size $= \%.3f$)"}\NormalTok{, }\FunctionTok{mean}\NormalTok{(perm.pv }\SpecialCharTok{\textless{}} \FloatTok{0.05}\NormalTok{))))}
\FunctionTok{abline}\NormalTok{(}\AttributeTok{h =} \DecValTok{1}\NormalTok{, }\AttributeTok{col =} \DecValTok{4}\NormalTok{, }\AttributeTok{lwd =} \DecValTok{2}\NormalTok{)}
\end{Highlighting}
\end{Shaded}

\pandocbounded{\includegraphics[keepaspectratio]{chapters/unit04_simulation_files/figure-pdf/power-regression-h0-1.pdf}}

With the size strictly regulated, we inject mathematical signals
(\(Y = X^{1.5} + \epsilon\) and \(Y = X^{0.5} + \epsilon\)) to observe
statistical power dynamically capturing the non-linear true signals.

\begin{Shaded}
\begin{Highlighting}[]
\NormalTok{perm.pv.alt1 }\OtherTok{\textless{}{-}} \FunctionTok{test.perm.power}\NormalTok{(x, Xp, }\AttributeTok{power.sim =} \FloatTok{1.5}\NormalTok{, }\AttributeTok{N =} \DecValTok{1000}\NormalTok{, }\AttributeTok{no.perm =} \DecValTok{99}\NormalTok{)}
\NormalTok{perm.pv.alt2 }\OtherTok{\textless{}{-}} \FunctionTok{test.perm.power}\NormalTok{(x, Xp, }\AttributeTok{power.sim =} \FloatTok{0.5}\NormalTok{, }\AttributeTok{N =} \DecValTok{1000}\NormalTok{, }\AttributeTok{no.perm =} \DecValTok{99}\NormalTok{)}

\FunctionTok{par}\NormalTok{(}\AttributeTok{mfrow =} \FunctionTok{c}\NormalTok{(}\DecValTok{1}\NormalTok{, }\DecValTok{2}\NormalTok{), }\AttributeTok{mar =} \FunctionTok{c}\NormalTok{(}\DecValTok{4}\NormalTok{, }\DecValTok{4}\NormalTok{, }\DecValTok{3}\NormalTok{, }\DecValTok{1}\NormalTok{))}

\FunctionTok{hist}\NormalTok{(perm.pv.alt1, }\AttributeTok{nclass =} \DecValTok{20}\NormalTok{, }\AttributeTok{col =} \StringTok{"grey90"}\NormalTok{,}
     \AttributeTok{xlab =} \FunctionTok{TeX}\NormalTok{(r}\StringTok{"($1000$ permutation p{-}values)"}\NormalTok{),}
     \AttributeTok{main =} \FunctionTok{TeX}\NormalTok{(}\FunctionTok{sprintf}\NormalTok{(r}\StringTok{"($H\_1: y = x\^{}\{1.5\} + }\SpecialCharTok{\textbackslash{}e}\StringTok{psilon$, power $= \%.3f$)"}\NormalTok{,}
                        \FunctionTok{mean}\NormalTok{(perm.pv.alt1 }\SpecialCharTok{\textless{}} \FloatTok{0.05}\NormalTok{))))}

\FunctionTok{hist}\NormalTok{(perm.pv.alt2, }\AttributeTok{nclass =} \DecValTok{20}\NormalTok{, }\AttributeTok{col =} \StringTok{"grey90"}\NormalTok{,}
     \AttributeTok{xlab =} \FunctionTok{TeX}\NormalTok{(r}\StringTok{"($1000$ permutation p{-}values)"}\NormalTok{),}
     \AttributeTok{main =} \FunctionTok{TeX}\NormalTok{(}\FunctionTok{sprintf}\NormalTok{(r}\StringTok{"($H\_1: y = x\^{}\{0.5\} + }\SpecialCharTok{\textbackslash{}e}\StringTok{psilon$, power $= \%.3f$)"}\NormalTok{,}
                        \FunctionTok{mean}\NormalTok{(perm.pv.alt2 }\SpecialCharTok{\textless{}} \FloatTok{0.05}\NormalTok{))))}
\end{Highlighting}
\end{Shaded}

\pandocbounded{\includegraphics[keepaspectratio]{chapters/unit04_simulation_files/figure-pdf/power-regression-power-1.pdf}}

\bookmarksetup{startatroot}

\chapter{Univariate Optimization for
MLE}\label{univariate-optimization-for-mle}

\section{Maximum Likelihood: Theory}\label{maximum-likelihood-theory}

\subsection{Likelihood, Score, and
Hessian}\label{likelihood-score-and-hessian}

A model specifies the density or probability of the data \(D\) given a
parameter \(\theta\), written \(f(D \mid \theta)\). As a function of
\(D\) with \(\theta\) fixed it is the sampling distribution; as a
function of \(\theta\) with \(D\) observed it is the likelihood. For
independent observations \(x_1, \ldots, x_n\),
\[L(\theta) = \prod_{i=1}^{n} f(x_i \mid \theta), \qquad \ell(\theta) = \log L(\theta) = \sum_{i=1}^{n} \log f(x_i \mid \theta).\]

We always work with \(\ell\): the product underflows to zero for even
moderate \(n\), and sums are easier to differentiate. The maximum
likelihood estimator maximizes \(\ell\), and where \(\ell\) is smooth
with an interior maximum, \(\hat\theta\) solves the \textbf{score
equation}
\[S(\theta) = \ell'(\theta) = 0, \qquad \ell''(\hat\theta) < 0 .\]

Three quantities recur throughout the chapter:

\begin{itemize}
\tightlist
\item
  the \textbf{score} \(S(\theta) = \ell'(\theta)\), which locates the
  maximum;
\item
  the \textbf{observed information} \(J(\theta) = -\ell''(\theta)\), the
  curvature of the log-likelihood at \(\theta\);
\item
  the \textbf{expected information}
  \(n I(\theta) = -E_D\left[\ell''(\theta; D)\right]\), its expectation
  under the model.
\end{itemize}

The second derivative therefore plays a double role. It determines how a
Newton-type algorithm steps toward the maximum, and it determines how
precisely the data pin the parameter down. An optimizer that computes
the curvature has computed the standard error at the same time.

\subsection{Closed Form: the Normal
Mean}\label{closed-form-the-normal-mean}

Let \(X_1, \ldots, X_n\) be iid \(N(\mu, \sigma_0^{2})\) with
\(\sigma_0^{2}\) known. Using the identity
\(\sum_i (x_i - \mu)^2 = \sum_i (x_i - \bar x)^2 + n(\bar x - \mu)^2\),
\[\ell(\mu) = -\frac{n}{2}\log\left(2\pi\sigma_0^{2}\right) - \frac{1}{2\sigma_0^{2}}\sum_{i=1}^{n}(x_i - \mu)^{2} = C - \frac{n}{2\sigma_0^{2}}(\mu - \bar x)^{2} .\]

Differentiating,
\[S(\mu) = \frac{n}{\sigma_0^{2}}\left(\bar x - \mu\right), \qquad \ell''(\mu) = -\frac{n}{\sigma_0^{2}}, \qquad \hat\mu = \bar x .\]

Every feature of the general theory is visible here in exact form. The
log-likelihood is a downward parabola, so it has a unique maximum. The
score is linear, so the score equation is solved in one line. The
curvature \(J(\mu) = n/\sigma_0^{2}\) does not depend on \(\mu\) or on
the data, so observed and expected information coincide, \(J = nI\), and
both grow linearly with \(n\): doubling the sample doubles the sharpness
of the peak and halves the variance of \(\hat\mu\).

\begin{Shaded}
\begin{Highlighting}[]
\NormalTok{x20 }\OtherTok{\textless{}{-}} \FunctionTok{rnorm}\NormalTok{(}\DecValTok{20}\NormalTok{, }\AttributeTok{mean =} \DecValTok{1}\NormalTok{); x50 }\OtherTok{\textless{}{-}} \FunctionTok{rnorm}\NormalTok{(}\DecValTok{50}\NormalTok{, }\AttributeTok{mean =} \DecValTok{1}\NormalTok{)}
\NormalTok{mu }\OtherTok{\textless{}{-}} \FunctionTok{seq}\NormalTok{(}\DecValTok{0}\NormalTok{, }\DecValTok{2}\NormalTok{, }\AttributeTok{length.out =} \DecValTok{200}\NormalTok{)}

\NormalTok{ll\_norm }\OtherTok{\textless{}{-}} \ControlFlowTok{function}\NormalTok{(mu, x) }\FunctionTok{vapply}\NormalTok{(mu, }\ControlFlowTok{function}\NormalTok{(m) }\FunctionTok{sum}\NormalTok{(}\FunctionTok{dnorm}\NormalTok{(x, m, }\AttributeTok{log =} \ConstantTok{TRUE}\NormalTok{)), }\FunctionTok{numeric}\NormalTok{(}\DecValTok{1}\NormalTok{))}

\FunctionTok{par}\NormalTok{(}\AttributeTok{mfrow =} \FunctionTok{c}\NormalTok{(}\DecValTok{1}\NormalTok{, }\DecValTok{3}\NormalTok{), }\AttributeTok{mar =} \FunctionTok{c}\NormalTok{(}\DecValTok{4}\NormalTok{, }\FloatTok{4.5}\NormalTok{, }\FloatTok{2.5}\NormalTok{, }\DecValTok{1}\NormalTok{))}

\FunctionTok{matplot}\NormalTok{(mu, }\FunctionTok{cbind}\NormalTok{(}\FunctionTok{ll\_norm}\NormalTok{(mu, x20) }\SpecialCharTok{{-}} \FunctionTok{max}\NormalTok{(}\FunctionTok{ll\_norm}\NormalTok{(mu, x20)),}
                  \FunctionTok{ll\_norm}\NormalTok{(mu, x50) }\SpecialCharTok{{-}} \FunctionTok{max}\NormalTok{(}\FunctionTok{ll\_norm}\NormalTok{(mu, x50))),}
        \AttributeTok{type =} \StringTok{"l"}\NormalTok{, }\AttributeTok{lty =} \DecValTok{1}\NormalTok{, }\AttributeTok{lwd =} \DecValTok{2}\NormalTok{, }\AttributeTok{col =} \FunctionTok{c}\NormalTok{(}\DecValTok{1}\NormalTok{, }\DecValTok{2}\NormalTok{),}
        \AttributeTok{xlab =} \FunctionTok{TeX}\NormalTok{(r}\StringTok{"($\textbackslash{}mu$)"}\NormalTok{), }\AttributeTok{ylab =} \FunctionTok{TeX}\NormalTok{(r}\StringTok{"($}\SpecialCharTok{\textbackslash{}e}\StringTok{ll(\textbackslash{}mu) {-} }\SpecialCharTok{\textbackslash{}e}\StringTok{ll(\textbackslash{}hat\{\textbackslash{}mu\})$)"}\NormalTok{),}
        \AttributeTok{main =} \FunctionTok{TeX}\NormalTok{(r}\StringTok{"(log{-}likelihood)"}\NormalTok{))}
\FunctionTok{legend}\NormalTok{(}\StringTok{"bottom"}\NormalTok{, }\AttributeTok{bty =} \StringTok{"n"}\NormalTok{, }\AttributeTok{lty =} \DecValTok{1}\NormalTok{, }\AttributeTok{lwd =} \DecValTok{2}\NormalTok{, }\AttributeTok{col =} \FunctionTok{c}\NormalTok{(}\DecValTok{1}\NormalTok{, }\DecValTok{2}\NormalTok{),}
       \AttributeTok{legend =} \FunctionTok{TeX}\NormalTok{(}\FunctionTok{c}\NormalTok{(r}\StringTok{"($n = 20$)"}\NormalTok{, r}\StringTok{"($n = 50$)"}\NormalTok{)))}

\FunctionTok{matplot}\NormalTok{(mu, }\FunctionTok{cbind}\NormalTok{(}\DecValTok{20} \SpecialCharTok{*}\NormalTok{ (}\FunctionTok{mean}\NormalTok{(x20) }\SpecialCharTok{{-}}\NormalTok{ mu), }\DecValTok{50} \SpecialCharTok{*}\NormalTok{ (}\FunctionTok{mean}\NormalTok{(x50) }\SpecialCharTok{{-}}\NormalTok{ mu)),}
        \AttributeTok{type =} \StringTok{"l"}\NormalTok{, }\AttributeTok{lty =} \DecValTok{1}\NormalTok{, }\AttributeTok{lwd =} \DecValTok{2}\NormalTok{, }\AttributeTok{col =} \FunctionTok{c}\NormalTok{(}\DecValTok{1}\NormalTok{, }\DecValTok{2}\NormalTok{),}
        \AttributeTok{xlab =} \FunctionTok{TeX}\NormalTok{(r}\StringTok{"($\textbackslash{}mu$)"}\NormalTok{), }\AttributeTok{ylab =} \FunctionTok{TeX}\NormalTok{(r}\StringTok{"($S(\textbackslash{}mu)$)"}\NormalTok{), }\AttributeTok{main =} \FunctionTok{TeX}\NormalTok{(r}\StringTok{"(score)"}\NormalTok{))}
\FunctionTok{abline}\NormalTok{(}\AttributeTok{h =} \DecValTok{0}\NormalTok{, }\AttributeTok{lty =} \DecValTok{3}\NormalTok{)}

\FunctionTok{matplot}\NormalTok{(mu, }\FunctionTok{cbind}\NormalTok{(}\FunctionTok{rep}\NormalTok{(}\DecValTok{20}\NormalTok{, }\DecValTok{200}\NormalTok{), }\FunctionTok{rep}\NormalTok{(}\DecValTok{50}\NormalTok{, }\DecValTok{200}\NormalTok{)), }\AttributeTok{type =} \StringTok{"l"}\NormalTok{, }\AttributeTok{lty =} \DecValTok{1}\NormalTok{, }\AttributeTok{lwd =} \DecValTok{2}\NormalTok{, }\AttributeTok{col =} \FunctionTok{c}\NormalTok{(}\DecValTok{1}\NormalTok{, }\DecValTok{2}\NormalTok{),}
        \AttributeTok{ylim =} \FunctionTok{c}\NormalTok{(}\DecValTok{0}\NormalTok{, }\DecValTok{60}\NormalTok{), }\AttributeTok{xlab =} \FunctionTok{TeX}\NormalTok{(r}\StringTok{"($\textbackslash{}mu$)"}\NormalTok{), }\AttributeTok{ylab =} \FunctionTok{TeX}\NormalTok{(r}\StringTok{"(${-}}\SpecialCharTok{\textbackslash{}e}\StringTok{ll\textquotesingle{}\textquotesingle{}(\textbackslash{}mu)$)"}\NormalTok{),}
        \AttributeTok{main =} \FunctionTok{TeX}\NormalTok{(r}\StringTok{"(curvature)"}\NormalTok{))}
\end{Highlighting}
\end{Shaded}

\pandocbounded{\includegraphics[keepaspectratio]{chapters/unit05_MLE1_files/figure-pdf/normal-loglik-score-curvature-1.pdf}}

The score crosses zero exactly where \(\ell\) peaks, and the larger
sample gives a more peaked log-likelihood, a steeper score, and a larger
curvature.

\subsection{Consistency, Asymptotic Normality, and Fisher
Information}\label{consistency-asymptotic-normality-and-fisher-information}

Under regularity conditions --- the support of \(f\) not depending on
\(\theta\), the true value interior to the parameter space, and two
derivatives passing under the integral sign --- the MLE has three
properties as \(n \to \infty\), with \(\theta_*\) the true parameter.

\textbf{Consistency.} \(\hat\theta \to \theta_*\) in probability. The
reason is that \(\ell(\theta)/n\) converges to
\(E_{\theta_*}[\log f(X \mid \theta)]\), which by the information
inequality is maximized at \(\theta = \theta_*\): the log-likelihood
surface converges to one whose peak is at the truth.

\textbf{Asymptotic normality.}
\[\hat\theta \;\dot\sim\; N\left(\theta_*,\ \left[n I(\theta_*)\right]^{-1}\right), \qquad I(\theta) = -E_X\left[\frac{\partial^{2}}{\partial\theta^{2}}\log f(X \mid \theta)\right] = E_X\left[\left(\frac{\partial}{\partial\theta}\log f(X\mid\theta)\right)^{2}\right].\]

A sketch: expand \(S(\hat\theta) = 0\) about \(\theta_*\) to get
\(\hat\theta - \theta_* \approx S(\theta_*)/\{-\ell''(\theta_*)\}\). The
numerator is a sum of iid mean-zero terms, hence asymptotically normal
with variance \(nI(\theta_*)\); the denominator converges to
\(nI(\theta_*)\) by the law of large numbers. Dividing gives variance
\([nI(\theta_*)]^{-1}\).

\textbf{Asymptotic efficiency.} No consistent estimator has smaller
asymptotic variance --- the Cramér--Rao bound is attained.

Since \(\theta_*\) is unknown, the variance is estimated by substituting
\(\hat\theta\). Two versions are available: the expected information
\(nI(\hat\theta)\), which requires the expectation analytically, and the
observed information \(J(\hat\theta) = -\ell''(\hat\theta)\), which
requires only the curvature that a Newton step already computes. By the
law of large numbers \(J(\hat\theta)/n \approx I(\hat\theta)\), and both
give
\[\widehat{\mathrm{SE}}(\hat\theta) = J(\hat\theta)^{-1/2} \quad\text{or}\quad \left[nI(\hat\theta)\right]^{-1/2}, \qquad \hat\theta \pm 1.96\,\widehat{\mathrm{SE}}(\hat\theta).\]

These are asymptotic statements. A flat log-likelihood --- small
curvature --- gives a large standard error, and a log-likelihood that is
far from quadratic gives a standard error that should not be trusted at
all; the last section of this chapter shows a case of exactly that.

\subsection{Demonstration: the Cauchy Location
Model}\label{demonstration-the-cauchy-location-model}

The Cauchy density \(f(x \mid \theta) = 1/\{\pi[1 + (x - \theta)^2]\}\)
has no closed-form MLE:
\[\ell(\theta) = -n\log\pi - \sum_{i=1}^{n}\log\left(1 + (x_i - \theta)^{2}\right), \qquad
S(\theta) = \sum_{i=1}^{n}\frac{2(x_i - \theta)}{1 + (x_i - \theta)^{2}} .\]

The score equation is a polynomial equation of degree \(2n - 1\).
Plotting \(\ell\) for samples of increasing size shows the peak
sharpening around the truth, which is consistency and growing
information in one picture.

\begin{Shaded}
\begin{Highlighting}[]
\NormalTok{log\_like\_cauchy }\OtherTok{\textless{}{-}} \ControlFlowTok{function}\NormalTok{(theta, x) }\FunctionTok{sum}\NormalTok{(}\FunctionTok{dcauchy}\NormalTok{(x, theta, }\AttributeTok{log =} \ConstantTok{TRUE}\NormalTok{))}

\NormalTok{thetas }\OtherTok{\textless{}{-}} \FunctionTok{seq}\NormalTok{(}\DecValTok{15}\NormalTok{, }\DecValTok{25}\NormalTok{, }\AttributeTok{length.out =} \DecValTok{200}\NormalTok{)}
\NormalTok{n\_set }\OtherTok{\textless{}{-}} \FunctionTok{c}\NormalTok{(}\DecValTok{50}\NormalTok{, }\DecValTok{100}\NormalTok{, }\DecValTok{200}\NormalTok{)}

\NormalTok{ll\_cauchy }\OtherTok{\textless{}{-}} \FunctionTok{sapply}\NormalTok{(n\_set, }\ControlFlowTok{function}\NormalTok{(n) \{}
\NormalTok{  d }\OtherTok{\textless{}{-}} \FunctionTok{rcauchy}\NormalTok{(n, }\AttributeTok{location =} \DecValTok{20}\NormalTok{)}
  \FunctionTok{sapply}\NormalTok{(thetas, log\_like\_cauchy, }\AttributeTok{x =}\NormalTok{ d)}
\NormalTok{\})}
\NormalTok{ll\_cauchy }\OtherTok{\textless{}{-}} \FunctionTok{sweep}\NormalTok{(ll\_cauchy, }\DecValTok{2}\NormalTok{, }\FunctionTok{apply}\NormalTok{(ll\_cauchy, }\DecValTok{2}\NormalTok{, max))   }\CommentTok{\# align peaks at 0}

\FunctionTok{par}\NormalTok{(}\AttributeTok{mar =} \FunctionTok{c}\NormalTok{(}\DecValTok{4}\NormalTok{, }\FloatTok{4.5}\NormalTok{, }\FloatTok{2.5}\NormalTok{, }\DecValTok{1}\NormalTok{))}
\FunctionTok{matplot}\NormalTok{(thetas, ll\_cauchy, }\AttributeTok{type =} \StringTok{"l"}\NormalTok{, }\AttributeTok{lty =} \DecValTok{1}\NormalTok{, }\AttributeTok{lwd =} \DecValTok{2}\NormalTok{, }\AttributeTok{col =} \FunctionTok{c}\NormalTok{(}\DecValTok{1}\NormalTok{, }\DecValTok{2}\NormalTok{, }\DecValTok{4}\NormalTok{), }\AttributeTok{ylim =} \FunctionTok{c}\NormalTok{(}\SpecialCharTok{{-}}\DecValTok{20}\NormalTok{, }\DecValTok{0}\NormalTok{),}
        \AttributeTok{xlab =} \FunctionTok{TeX}\NormalTok{(r}\StringTok{"($}\SpecialCharTok{\textbackslash{}t}\StringTok{heta$)"}\NormalTok{), }\AttributeTok{ylab =} \FunctionTok{TeX}\NormalTok{(r}\StringTok{"($}\SpecialCharTok{\textbackslash{}e}\StringTok{ll(}\SpecialCharTok{\textbackslash{}t}\StringTok{heta) {-} }\SpecialCharTok{\textbackslash{}e}\StringTok{ll(\textbackslash{}hat\{}\SpecialCharTok{\textbackslash{}t}\StringTok{heta\})$)"}\NormalTok{),}
        \AttributeTok{main =} \FunctionTok{TeX}\NormalTok{(r}\StringTok{"(Cauchy log{-}likelihood, $}\SpecialCharTok{\textbackslash{}t}\StringTok{heta\_* = 20$)"}\NormalTok{))}
\FunctionTok{abline}\NormalTok{(}\AttributeTok{v =} \DecValTok{20}\NormalTok{, }\AttributeTok{lty =} \DecValTok{3}\NormalTok{)}
\FunctionTok{legend}\NormalTok{(}\StringTok{"bottomright"}\NormalTok{, }\AttributeTok{bty =} \StringTok{"n"}\NormalTok{, }\AttributeTok{lty =} \DecValTok{1}\NormalTok{, }\AttributeTok{lwd =} \DecValTok{2}\NormalTok{, }\AttributeTok{col =} \FunctionTok{c}\NormalTok{(}\DecValTok{1}\NormalTok{, }\DecValTok{2}\NormalTok{, }\DecValTok{4}\NormalTok{),}
       \AttributeTok{legend =} \FunctionTok{TeX}\NormalTok{(}\FunctionTok{sprintf}\NormalTok{(r}\StringTok{"($n = \%d$)"}\NormalTok{, n\_set)))}
\end{Highlighting}
\end{Shaded}

\pandocbounded{\includegraphics[keepaspectratio]{chapters/unit05_MLE1_files/figure-pdf/cauchy-loglik-by-n-1.pdf}}

The variability of the MLE across repeated samples is the sampling
distribution the theory describes. Repeating the experiment 50 times at
two sample sizes shows the curves clustering more tightly, and their
maxima concentrating, as \(n\) grows.

\begin{Shaded}
\begin{Highlighting}[]
\NormalTok{loglik\_curves }\OtherTok{\textless{}{-}} \ControlFlowTok{function}\NormalTok{(n, }\AttributeTok{nrep =} \DecValTok{50}\NormalTok{)}
  \FunctionTok{replicate}\NormalTok{(nrep, \{}
\NormalTok{    d }\OtherTok{\textless{}{-}} \FunctionTok{rcauchy}\NormalTok{(n, }\AttributeTok{location =} \DecValTok{20}\NormalTok{)}
\NormalTok{    lc }\OtherTok{\textless{}{-}} \FunctionTok{sapply}\NormalTok{(thetas, log\_like\_cauchy, }\AttributeTok{x =}\NormalTok{ d)}
\NormalTok{    lc }\SpecialCharTok{{-}} \FunctionTok{max}\NormalTok{(lc)}
\NormalTok{  \})}

\NormalTok{curves\_small }\OtherTok{\textless{}{-}} \FunctionTok{loglik\_curves}\NormalTok{(}\DecValTok{20}\NormalTok{)}
\NormalTok{curves\_large }\OtherTok{\textless{}{-}} \FunctionTok{loglik\_curves}\NormalTok{(}\DecValTok{200}\NormalTok{)}

\FunctionTok{par}\NormalTok{(}\AttributeTok{mfrow =} \FunctionTok{c}\NormalTok{(}\DecValTok{1}\NormalTok{, }\DecValTok{2}\NormalTok{), }\AttributeTok{mar =} \FunctionTok{c}\NormalTok{(}\DecValTok{4}\NormalTok{, }\FloatTok{4.5}\NormalTok{, }\FloatTok{2.5}\NormalTok{, }\DecValTok{1}\NormalTok{))}
\ControlFlowTok{for}\NormalTok{ (obj }\ControlFlowTok{in} \FunctionTok{list}\NormalTok{(}\FunctionTok{list}\NormalTok{(}\AttributeTok{m =}\NormalTok{ curves\_small, }\AttributeTok{n =} \DecValTok{20}\NormalTok{), }\FunctionTok{list}\NormalTok{(}\AttributeTok{m =}\NormalTok{ curves\_large, }\AttributeTok{n =} \DecValTok{200}\NormalTok{)))}
  \FunctionTok{matplot}\NormalTok{(thetas, obj}\SpecialCharTok{$}\NormalTok{m, }\AttributeTok{type =} \StringTok{"l"}\NormalTok{, }\AttributeTok{lty =} \DecValTok{1}\NormalTok{, }\AttributeTok{col =} \FunctionTok{adjustcolor}\NormalTok{(}\DecValTok{1}\NormalTok{, }\FloatTok{0.35}\NormalTok{), }\AttributeTok{ylim =} \FunctionTok{c}\NormalTok{(}\SpecialCharTok{{-}}\DecValTok{10}\NormalTok{, }\DecValTok{0}\NormalTok{),}
          \AttributeTok{xlab =} \FunctionTok{TeX}\NormalTok{(r}\StringTok{"($}\SpecialCharTok{\textbackslash{}t}\StringTok{heta$)"}\NormalTok{), }\AttributeTok{ylab =} \FunctionTok{TeX}\NormalTok{(r}\StringTok{"($}\SpecialCharTok{\textbackslash{}e}\StringTok{ll(}\SpecialCharTok{\textbackslash{}t}\StringTok{heta) {-} }\SpecialCharTok{\textbackslash{}e}\StringTok{ll(\textbackslash{}hat\{}\SpecialCharTok{\textbackslash{}t}\StringTok{heta\})$)"}\NormalTok{),}
          \AttributeTok{main =} \FunctionTok{TeX}\NormalTok{(}\FunctionTok{sprintf}\NormalTok{(r}\StringTok{"(sample size $n = \%d$)"}\NormalTok{, obj}\SpecialCharTok{$}\NormalTok{n)))}
\end{Highlighting}
\end{Shaded}

\pandocbounded{\includegraphics[keepaspectratio]{chapters/unit05_MLE1_files/figure-pdf/cauchy-loglik-variability-1.pdf}}

\begin{Shaded}
\begin{Highlighting}[]
\DocumentationTok{\#\#\# sampling distribution of the MLE at the two sample sizes}
\NormalTok{mle\_at }\OtherTok{\textless{}{-}} \ControlFlowTok{function}\NormalTok{(m) thetas[}\FunctionTok{apply}\NormalTok{(m, }\DecValTok{2}\NormalTok{, which.max)]}
\FunctionTok{rbind}\NormalTok{(}\StringTok{\textasciigrave{}}\AttributeTok{n = 20}\StringTok{\textasciigrave{}}  \OtherTok{=} \FunctionTok{c}\NormalTok{(}\AttributeTok{mean =} \FunctionTok{mean}\NormalTok{(}\FunctionTok{mle\_at}\NormalTok{(curves\_small)), }\AttributeTok{sd =} \FunctionTok{sd}\NormalTok{(}\FunctionTok{mle\_at}\NormalTok{(curves\_small))),}
      \StringTok{\textasciigrave{}}\AttributeTok{n = 200}\StringTok{\textasciigrave{}} \OtherTok{=} \FunctionTok{c}\NormalTok{(}\AttributeTok{mean =} \FunctionTok{mean}\NormalTok{(}\FunctionTok{mle\_at}\NormalTok{(curves\_large)), }\AttributeTok{sd =} \FunctionTok{sd}\NormalTok{(}\FunctionTok{mle\_at}\NormalTok{(curves\_large))))}
\end{Highlighting}
\end{Shaded}

\begin{verbatim}
            mean        sd
n = 20  20.05930 0.3843737
n = 200 20.01005 0.0920729
\end{verbatim}

The standard deviation falls by roughly the factor \(\sqrt{10}\)
predicted by \([nI(\theta_*)]^{-1/2}\). Two features of these curves
also warn us about computation: for small \(n\) the log-likelihood has
several local maxima, and its shape away from the peak is far from
quadratic. Numerical optimization is required, and it must be used with
care.

\section{Univariate Optimization
Algorithms}\label{univariate-optimization-algorithms}

\subsection{Three Equivalent
Formulations}\label{three-equivalent-formulations}

The same task can be posed three ways, and the algorithms below attack
different versions of it:

\begin{itemize}
\tightlist
\item
  maximize \(\ell(\theta)\);
\item
  minimize \(f(\theta) = -\ell(\theta)\);
\item
  find a root of \(S(\theta) = \ell'(\theta)\).
\end{itemize}

Root-finders are usually faster because they use derivative information,
but they locate any stationary point, including minima and saddle
points. Direct optimizers need only function values and always move
toward a maximum, but converge more slowly. Which to prefer depends on
whether derivatives are available and whether the function is smooth.

Each method below is illustrated with a schematic showing its first few
iterations. To keep the pictures comparable they all use the same small
problem: the per-observation log-likelihood of the zero-truncated
Poisson model of the next section, with \(\bar x\) fixed at 1.8. The
root of its score is marked in every figure.

\begin{Shaded}
\begin{Highlighting}[]
\NormalTok{xbar\_d }\OtherTok{\textless{}{-}} \FloatTok{1.8}                                             \CommentTok{\# fixed sample mean}

\NormalTok{ll\_d }\OtherTok{\textless{}{-}} \ControlFlowTok{function}\NormalTok{(l) xbar\_d }\SpecialCharTok{*} \FunctionTok{log}\NormalTok{(l) }\SpecialCharTok{{-}}\NormalTok{ l }\SpecialCharTok{{-}} \FunctionTok{log}\NormalTok{(}\DecValTok{1} \SpecialCharTok{{-}} \FunctionTok{exp}\NormalTok{(}\SpecialCharTok{{-}}\NormalTok{l))          }\CommentTok{\# log{-}likelihood}
\NormalTok{S\_d  }\OtherTok{\textless{}{-}} \ControlFlowTok{function}\NormalTok{(l) xbar\_d }\SpecialCharTok{/}\NormalTok{ l }\SpecialCharTok{{-}} \DecValTok{1} \SpecialCharTok{/}\NormalTok{ (}\DecValTok{1} \SpecialCharTok{{-}} \FunctionTok{exp}\NormalTok{(}\SpecialCharTok{{-}}\NormalTok{l))                  }\CommentTok{\# score}
\NormalTok{dS\_d }\OtherTok{\textless{}{-}} \ControlFlowTok{function}\NormalTok{(l) }\SpecialCharTok{{-}}\NormalTok{xbar\_d }\SpecialCharTok{/}\NormalTok{ l}\SpecialCharTok{\^{}}\DecValTok{2} \SpecialCharTok{+} \FunctionTok{exp}\NormalTok{(}\SpecialCharTok{{-}}\NormalTok{l) }\SpecialCharTok{/}\NormalTok{ (}\DecValTok{1} \SpecialCharTok{{-}} \FunctionTok{exp}\NormalTok{(}\SpecialCharTok{{-}}\NormalTok{l))}\SpecialCharTok{\^{}}\DecValTok{2}       \CommentTok{\# its derivative}
\NormalTok{I\_d  }\OtherTok{\textless{}{-}} \ControlFlowTok{function}\NormalTok{(l) }\DecValTok{1} \SpecialCharTok{/}\NormalTok{ (l }\SpecialCharTok{*}\NormalTok{ (}\DecValTok{1} \SpecialCharTok{{-}} \FunctionTok{exp}\NormalTok{(}\SpecialCharTok{{-}}\NormalTok{l))) }\SpecialCharTok{{-}} \FunctionTok{exp}\NormalTok{(}\SpecialCharTok{{-}}\NormalTok{l) }\SpecialCharTok{/}\NormalTok{ (}\DecValTok{1} \SpecialCharTok{{-}} \FunctionTok{exp}\NormalTok{(}\SpecialCharTok{{-}}\NormalTok{l))}\SpecialCharTok{\^{}}\DecValTok{2}   \CommentTok{\# information}
\NormalTok{g\_d  }\OtherTok{\textless{}{-}} \ControlFlowTok{function}\NormalTok{(l) xbar\_d }\SpecialCharTok{*}\NormalTok{ (}\DecValTok{1} \SpecialCharTok{{-}} \FunctionTok{exp}\NormalTok{(}\SpecialCharTok{{-}}\NormalTok{l))                          }\CommentTok{\# fixed{-}point map}

\NormalTok{lstar }\OtherTok{\textless{}{-}} \FunctionTok{uniroot}\NormalTok{(S\_d, }\FunctionTok{c}\NormalTok{(}\FloatTok{0.2}\NormalTok{, }\DecValTok{5}\NormalTok{))}\SpecialCharTok{$}\NormalTok{root}
\FunctionTok{c}\NormalTok{(}\AttributeTok{root =}\NormalTok{ lstar, }\AttributeTok{contraction\_factor =}\NormalTok{ xbar\_d }\SpecialCharTok{*} \FunctionTok{exp}\NormalTok{(}\SpecialCharTok{{-}}\NormalTok{lstar))}
\end{Highlighting}
\end{Shaded}

\begin{verbatim}
              root contraction_factor 
          1.318364           0.481631 
\end{verbatim}

\subsection{Fixed-Point (Simple)
Iteration}\label{fixed-point-simple-iteration}

Rearrange the score equation algebraically into the form
\(\theta = g(\theta)\) and iterate \(\theta_{k+1} = g(\theta_k)\). If
\(g\) maps an interval into itself and \(|g'(\theta^{*})| < 1\) at the
fixed point, the mean value theorem gives
\[\left|\theta_{k+1} - \theta^{*}\right| = \left|g(\theta_k) - g(\theta^{*})\right| \le \left|g'(\xi)\right| \left|\theta_k - \theta^{*}\right| ,\]
so the error contracts by roughly the constant factor
\(|g'(\theta^{*})|\) at each step. This is \textbf{linear convergence}:
a fixed number of digits is gained per iteration, and the closer
\(|g'|\) is to 1 the slower the progress.

\begin{Shaded}
\begin{Highlighting}[]
\NormalTok{fixed\_point }\OtherTok{\textless{}{-}} \ControlFlowTok{function}\NormalTok{(g, x0, }\AttributeTok{tol =} \FloatTok{1e{-}10}\NormalTok{, }\AttributeTok{maxit =} \DecValTok{500}\NormalTok{) \{}
  \ControlFlowTok{for}\NormalTok{ (k }\ControlFlowTok{in} \DecValTok{1}\SpecialCharTok{:}\NormalTok{maxit) \{}
\NormalTok{    x1 }\OtherTok{\textless{}{-}} \FunctionTok{g}\NormalTok{(x0)}
    \ControlFlowTok{if}\NormalTok{ (}\FunctionTok{abs}\NormalTok{(x1 }\SpecialCharTok{{-}}\NormalTok{ x0) }\SpecialCharTok{\textless{}}\NormalTok{ tol) }\FunctionTok{return}\NormalTok{(}\FunctionTok{c}\NormalTok{(}\AttributeTok{root =}\NormalTok{ x1, }\AttributeTok{iter =}\NormalTok{ k))}
\NormalTok{    x0 }\OtherTok{\textless{}{-}}\NormalTok{ x1}
\NormalTok{  \}}
  \FunctionTok{stop}\NormalTok{(}\StringTok{"not converged"}\NormalTok{)}
\NormalTok{\}}
\end{Highlighting}
\end{Shaded}

The mechanism is easiest to see as a \textbf{cobweb}: from \(\lambda_k\)
on the diagonal, move vertically to the curve \(g\) to obtain
\(\lambda_{k+1}\), then horizontally back to the diagonal to start the
next step. The iterates walk toward the intersection of \(g\) with the
\(45^{\circ}\) line, and the shallower \(g\) is at that crossing --- the
smaller \(|g'|\) --- the larger each stride toward it.

\begin{Shaded}
\begin{Highlighting}[]
\NormalTok{lam }\OtherTok{\textless{}{-}} \FunctionTok{seq}\NormalTok{(}\FloatTok{0.2}\NormalTok{, }\DecValTok{3}\NormalTok{, }\AttributeTok{length.out =} \DecValTok{300}\NormalTok{)}

\FunctionTok{par}\NormalTok{(}\AttributeTok{mar =} \FunctionTok{c}\NormalTok{(}\DecValTok{4}\NormalTok{, }\FloatTok{4.5}\NormalTok{, }\FloatTok{2.5}\NormalTok{, }\DecValTok{1}\NormalTok{))}
\FunctionTok{plot}\NormalTok{(lam, }\FunctionTok{g\_d}\NormalTok{(lam), }\AttributeTok{type =} \StringTok{"l"}\NormalTok{, }\AttributeTok{lwd =} \DecValTok{2}\NormalTok{, }\AttributeTok{xlim =} \FunctionTok{c}\NormalTok{(}\FloatTok{0.2}\NormalTok{, }\DecValTok{3}\NormalTok{), }\AttributeTok{ylim =} \FunctionTok{c}\NormalTok{(}\FloatTok{0.2}\NormalTok{, }\DecValTok{3}\NormalTok{),}
     \AttributeTok{xlab =} \FunctionTok{TeX}\NormalTok{(r}\StringTok{"($\textbackslash{}lambda\_k$)"}\NormalTok{), }\AttributeTok{ylab =} \FunctionTok{TeX}\NormalTok{(r}\StringTok{"($g(\textbackslash{}lambda\_k) = \textbackslash{}lambda\_\{k+1\}$)"}\NormalTok{),}
     \AttributeTok{main =} \FunctionTok{TeX}\NormalTok{(r}\StringTok{"(fixed{-}point iteration: the cobweb)"}\NormalTok{))}
\FunctionTok{abline}\NormalTok{(}\DecValTok{0}\NormalTok{, }\DecValTok{1}\NormalTok{, }\AttributeTok{lty =} \DecValTok{2}\NormalTok{)}

\NormalTok{x0 }\OtherTok{\textless{}{-}} \FloatTok{2.8}
\ControlFlowTok{for}\NormalTok{ (k }\ControlFlowTok{in} \DecValTok{1}\SpecialCharTok{:}\DecValTok{5}\NormalTok{) \{}
\NormalTok{  x1 }\OtherTok{\textless{}{-}} \FunctionTok{g\_d}\NormalTok{(x0)}
  \FunctionTok{segments}\NormalTok{(x0, x0, x0, x1, }\AttributeTok{col =} \DecValTok{2}\NormalTok{, }\AttributeTok{lwd =} \DecValTok{2}\NormalTok{)          }\CommentTok{\# up or down to the curve}
  \FunctionTok{segments}\NormalTok{(x0, x1, x1, x1, }\AttributeTok{col =} \DecValTok{2}\NormalTok{, }\AttributeTok{lwd =} \DecValTok{2}\NormalTok{)          }\CommentTok{\# across to the diagonal}
  \FunctionTok{points}\NormalTok{(x0, x1, }\AttributeTok{pch =} \DecValTok{19}\NormalTok{, }\AttributeTok{col =} \DecValTok{2}\NormalTok{, }\AttributeTok{cex =} \FloatTok{0.7}\NormalTok{)}
  \ControlFlowTok{if}\NormalTok{ (k }\SpecialCharTok{\textless{}=} \DecValTok{3}\NormalTok{) }\FunctionTok{text}\NormalTok{(x0, }\FloatTok{0.25}\NormalTok{, }\FunctionTok{TeX}\NormalTok{(}\FunctionTok{sprintf}\NormalTok{(r}\StringTok{"($\textbackslash{}lambda\_\%d$)"}\NormalTok{, k }\SpecialCharTok{{-}} \DecValTok{1}\NormalTok{)), }\AttributeTok{col =} \DecValTok{2}\NormalTok{, }\AttributeTok{cex =} \FloatTok{0.9}\NormalTok{)}
\NormalTok{  x0 }\OtherTok{\textless{}{-}}\NormalTok{ x1}
\NormalTok{\}}
\FunctionTok{points}\NormalTok{(lstar, lstar, }\AttributeTok{pch =} \DecValTok{19}\NormalTok{, }\AttributeTok{cex =} \FloatTok{1.2}\NormalTok{)}
\FunctionTok{text}\NormalTok{(lstar, lstar, }\FunctionTok{TeX}\NormalTok{(r}\StringTok{"($\textbackslash{}hat\{\textbackslash{}lambda\}$)"}\NormalTok{), }\AttributeTok{pos =} \DecValTok{4}\NormalTok{)}
\end{Highlighting}
\end{Shaded}

\pandocbounded{\includegraphics[keepaspectratio]{chapters/unit05_MLE1_files/figure-pdf/scheme-fixed-point-1.pdf}}

The method needs no derivatives, but the rearrangement must be found by
hand, and different rearrangements of the same equation can converge at
different speeds or diverge outright. Geometrically, a rearrangement
with \(|g'| > 1\) at the crossing produces a cobweb that spirals outward
instead of inward.

\subsection{Bisection}\label{bisection}

If \(S\) is continuous and \(S(a)\) and \(S(b)\) have opposite signs, a
root lies in \([a, b]\). Evaluate \(S\) at the midpoint and retain
whichever half still changes sign.

\begin{Shaded}
\begin{Highlighting}[]
\NormalTok{bisection }\OtherTok{\textless{}{-}} \ControlFlowTok{function}\NormalTok{(S, a, b, }\AttributeTok{tol =} \FloatTok{1e{-}10}\NormalTok{) \{}
  \ControlFlowTok{if}\NormalTok{ (}\FunctionTok{S}\NormalTok{(a) }\SpecialCharTok{*} \FunctionTok{S}\NormalTok{(b) }\SpecialCharTok{\textgreater{}} \DecValTok{0}\NormalTok{) }\FunctionTok{stop}\NormalTok{(}\StringTok{"endpoints do not bracket a root"}\NormalTok{)}
\NormalTok{  k }\OtherTok{\textless{}{-}} \DecValTok{0}
  \ControlFlowTok{while}\NormalTok{ (b }\SpecialCharTok{{-}}\NormalTok{ a }\SpecialCharTok{\textgreater{}}\NormalTok{ tol) \{}
\NormalTok{    m }\OtherTok{\textless{}{-}}\NormalTok{ (a }\SpecialCharTok{+}\NormalTok{ b) }\SpecialCharTok{/} \DecValTok{2}
    \ControlFlowTok{if}\NormalTok{ (}\FunctionTok{S}\NormalTok{(a) }\SpecialCharTok{*} \FunctionTok{S}\NormalTok{(m) }\SpecialCharTok{\textless{}=} \DecValTok{0}\NormalTok{) b }\OtherTok{\textless{}{-}}\NormalTok{ m }\ControlFlowTok{else}\NormalTok{ a }\OtherTok{\textless{}{-}}\NormalTok{ m}
\NormalTok{    k }\OtherTok{\textless{}{-}}\NormalTok{ k }\SpecialCharTok{+} \DecValTok{1}
\NormalTok{  \}}
  \FunctionTok{c}\NormalTok{(}\AttributeTok{root =}\NormalTok{ (a }\SpecialCharTok{+}\NormalTok{ b) }\SpecialCharTok{/} \DecValTok{2}\NormalTok{, }\AttributeTok{iter =}\NormalTok{ k)}
\NormalTok{\}}
\end{Highlighting}
\end{Shaded}

The schematic below draws the bracket at each iteration as a bar beneath
the score curve. The dot is the midpoint, and the dotted line carries it
up to the curve, where its sign decides which half survives.

\begin{Shaded}
\begin{Highlighting}[]
\FunctionTok{par}\NormalTok{(}\AttributeTok{mar =} \FunctionTok{c}\NormalTok{(}\DecValTok{4}\NormalTok{, }\FloatTok{4.5}\NormalTok{, }\FloatTok{2.5}\NormalTok{, }\DecValTok{1}\NormalTok{))}
\FunctionTok{plot}\NormalTok{(lam, }\FunctionTok{S\_d}\NormalTok{(lam), }\AttributeTok{type =} \StringTok{"l"}\NormalTok{, }\AttributeTok{lwd =} \DecValTok{2}\NormalTok{, }\AttributeTok{ylim =} \FunctionTok{c}\NormalTok{(}\SpecialCharTok{{-}}\FloatTok{1.6}\NormalTok{, }\FloatTok{3.5}\NormalTok{),}
     \AttributeTok{xlab =} \FunctionTok{TeX}\NormalTok{(r}\StringTok{"($\textbackslash{}lambda$)"}\NormalTok{), }\AttributeTok{ylab =} \FunctionTok{TeX}\NormalTok{(r}\StringTok{"($S(\textbackslash{}lambda)$)"}\NormalTok{),}
     \AttributeTok{main =} \FunctionTok{TeX}\NormalTok{(r}\StringTok{"(bisection: each step halves the bracket)"}\NormalTok{))}
\FunctionTok{abline}\NormalTok{(}\AttributeTok{h =} \DecValTok{0}\NormalTok{, }\AttributeTok{lty =} \DecValTok{3}\NormalTok{)}
\FunctionTok{points}\NormalTok{(lstar, }\DecValTok{0}\NormalTok{, }\AttributeTok{pch =} \DecValTok{19}\NormalTok{, }\AttributeTok{cex =} \FloatTok{1.2}\NormalTok{)}

\NormalTok{a }\OtherTok{\textless{}{-}} \FloatTok{0.3}\NormalTok{; b }\OtherTok{\textless{}{-}} \DecValTok{3}
\NormalTok{ylev }\OtherTok{\textless{}{-}} \FunctionTok{seq}\NormalTok{(}\SpecialCharTok{{-}}\FloatTok{0.5}\NormalTok{, }\SpecialCharTok{{-}}\FloatTok{1.4}\NormalTok{, }\AttributeTok{length.out =} \DecValTok{4}\NormalTok{)}
\ControlFlowTok{for}\NormalTok{ (k }\ControlFlowTok{in} \DecValTok{1}\SpecialCharTok{:}\DecValTok{4}\NormalTok{) \{}
\NormalTok{  m }\OtherTok{\textless{}{-}}\NormalTok{ (a }\SpecialCharTok{+}\NormalTok{ b) }\SpecialCharTok{/} \DecValTok{2}
  \FunctionTok{segments}\NormalTok{(a, ylev[k], b, ylev[k], }\AttributeTok{col =} \DecValTok{4}\NormalTok{, }\AttributeTok{lwd =} \DecValTok{3}\NormalTok{)}
  \FunctionTok{points}\NormalTok{(}\FunctionTok{c}\NormalTok{(a, b), }\FunctionTok{rep}\NormalTok{(ylev[k], }\DecValTok{2}\NormalTok{), }\AttributeTok{pch =} \StringTok{"|"}\NormalTok{, }\AttributeTok{col =} \DecValTok{4}\NormalTok{)}
  \FunctionTok{segments}\NormalTok{(m, ylev[k], m, }\FunctionTok{S\_d}\NormalTok{(m), }\AttributeTok{lty =} \DecValTok{3}\NormalTok{, }\AttributeTok{col =} \DecValTok{2}\NormalTok{)}
  \FunctionTok{points}\NormalTok{(m, ylev[k], }\AttributeTok{pch =} \DecValTok{19}\NormalTok{, }\AttributeTok{col =} \DecValTok{2}\NormalTok{, }\AttributeTok{cex =} \FloatTok{0.8}\NormalTok{)}
  \FunctionTok{text}\NormalTok{(}\FloatTok{0.25}\NormalTok{, ylev[k], k, }\AttributeTok{cex =} \FloatTok{0.8}\NormalTok{, }\AttributeTok{col =} \DecValTok{4}\NormalTok{)}
  \ControlFlowTok{if}\NormalTok{ (}\FunctionTok{S\_d}\NormalTok{(a) }\SpecialCharTok{*} \FunctionTok{S\_d}\NormalTok{(m) }\SpecialCharTok{\textless{}=} \DecValTok{0}\NormalTok{) b }\OtherTok{\textless{}{-}}\NormalTok{ m }\ControlFlowTok{else}\NormalTok{ a }\OtherTok{\textless{}{-}}\NormalTok{ m}
\NormalTok{\}}
\end{Highlighting}
\end{Shaded}

\pandocbounded{\includegraphics[keepaspectratio]{chapters/unit05_MLE1_files/figure-pdf/scheme-bisection-1.pdf}}

The bracket halves at every step, so after \(k\) iterations the error is
at most \((b - a)/2^{k}\) and the iteration count is known in advance:
about \(\log_2\{(b-a)/\text{tol}\}\), regardless of the shape of \(S\).
Bisection cannot diverge, which makes it the natural fallback when a
faster method misbehaves. R's \texttt{uniroot} uses Brent's method,
which retains the bracketing guarantee while interpolating for speed.

\subsection{Newton--Raphson}\label{newtonraphson}

\subsubsection{Derivation}\label{derivation}

Approximate the score by its tangent line at \(\theta_k\) and solve for
where the tangent crosses zero:
\[S(\theta) \approx S(\theta_k) + S'(\theta_k)(\theta - \theta_k) = 0
\implies \theta_{k+1} = \theta_k - \frac{S(\theta_k)}{S'(\theta_k)} = \theta_k - \frac{\ell'(\theta_k)}{\ell''(\theta_k)} .\]

Equivalently, replace \(\ell\) by its second-order Taylor expansion at
\(\theta_k\) and jump to the maximum of that parabola. For an exactly
quadratic \(\ell(\theta) = -a(\theta - u)^2 + C\) with \(a > 0\) we have
\(\ell' = -2a(\theta - u)\) and \(\ell'' = -2a\), so
\[\theta_{k+1} = \theta_k - \frac{-2a(\theta_k - u)}{-2a} = u\] in one
step from any starting value. The normal-mean example of the previous
section is exactly of this form: Newton converges immediately.

The left panel below traces three steps. At each iterate, the tangent to
the score is followed down to the horizontal axis, and where it lands is
the next iterate. As the curve straightens near the root, the tangent
becomes an excellent approximation and the steps land almost exactly on
the answer.

The right panel contrasts one Newton step with one step of the method of
scoring from the same point. Newton uses the actual slope
\(S'(\lambda_k)\); scoring uses \(-nI(\lambda_k)\), the slope averaged
over hypothetical datasets. The two lines differ away from the root and
coincide as \(\lambda_k\) approaches it.

\begin{Shaded}
\begin{Highlighting}[]
\FunctionTok{par}\NormalTok{(}\AttributeTok{mfrow =} \FunctionTok{c}\NormalTok{(}\DecValTok{1}\NormalTok{, }\DecValTok{2}\NormalTok{), }\AttributeTok{mar =} \FunctionTok{c}\NormalTok{(}\DecValTok{4}\NormalTok{, }\FloatTok{4.5}\NormalTok{, }\FloatTok{2.5}\NormalTok{, }\DecValTok{1}\NormalTok{))}

\DocumentationTok{\#\#\# three Newton steps}
\FunctionTok{plot}\NormalTok{(lam, }\FunctionTok{S\_d}\NormalTok{(lam), }\AttributeTok{type =} \StringTok{"l"}\NormalTok{, }\AttributeTok{lwd =} \DecValTok{2}\NormalTok{, }\AttributeTok{ylim =} \FunctionTok{c}\NormalTok{(}\SpecialCharTok{{-}}\DecValTok{1}\NormalTok{, }\FloatTok{3.5}\NormalTok{),}
     \AttributeTok{xlab =} \FunctionTok{TeX}\NormalTok{(r}\StringTok{"($\textbackslash{}lambda$)"}\NormalTok{), }\AttributeTok{ylab =} \FunctionTok{TeX}\NormalTok{(r}\StringTok{"($S(\textbackslash{}lambda)$)"}\NormalTok{),}
     \AttributeTok{main =} \FunctionTok{TeX}\NormalTok{(r}\StringTok{"(Newton{-}Raphson: follow the tangent)"}\NormalTok{))}
\FunctionTok{abline}\NormalTok{(}\AttributeTok{h =} \DecValTok{0}\NormalTok{, }\AttributeTok{lty =} \DecValTok{3}\NormalTok{)}
\FunctionTok{points}\NormalTok{(lstar, }\DecValTok{0}\NormalTok{, }\AttributeTok{pch =} \DecValTok{19}\NormalTok{, }\AttributeTok{cex =} \FloatTok{1.2}\NormalTok{)}

\NormalTok{x0 }\OtherTok{\textless{}{-}} \FloatTok{0.45}
\ControlFlowTok{for}\NormalTok{ (k }\ControlFlowTok{in} \DecValTok{1}\SpecialCharTok{:}\DecValTok{3}\NormalTok{) \{}
\NormalTok{  x1 }\OtherTok{\textless{}{-}}\NormalTok{ x0 }\SpecialCharTok{{-}} \FunctionTok{S\_d}\NormalTok{(x0) }\SpecialCharTok{/} \FunctionTok{dS\_d}\NormalTok{(x0)}
  \FunctionTok{segments}\NormalTok{(x0, }\FunctionTok{S\_d}\NormalTok{(x0), x1, }\DecValTok{0}\NormalTok{, }\AttributeTok{col =} \DecValTok{2}\NormalTok{, }\AttributeTok{lwd =} \DecValTok{2}\NormalTok{)}
  \FunctionTok{segments}\NormalTok{(x1, }\DecValTok{0}\NormalTok{, x1, }\FunctionTok{S\_d}\NormalTok{(x1), }\AttributeTok{lty =} \DecValTok{3}\NormalTok{, }\AttributeTok{col =} \StringTok{"grey40"}\NormalTok{)}
  \FunctionTok{points}\NormalTok{(}\FunctionTok{c}\NormalTok{(x0, x1), }\FunctionTok{c}\NormalTok{(}\FunctionTok{S\_d}\NormalTok{(x0), }\DecValTok{0}\NormalTok{), }\AttributeTok{pch =} \DecValTok{19}\NormalTok{, }\AttributeTok{col =} \DecValTok{2}\NormalTok{, }\AttributeTok{cex =} \FloatTok{0.7}\NormalTok{)}
  \FunctionTok{text}\NormalTok{(x0, }\FunctionTok{S\_d}\NormalTok{(x0), }\FunctionTok{TeX}\NormalTok{(}\FunctionTok{sprintf}\NormalTok{(r}\StringTok{"($\textbackslash{}lambda\_\%d$)"}\NormalTok{, k }\SpecialCharTok{{-}} \DecValTok{1}\NormalTok{)), }\AttributeTok{pos =} \DecValTok{3}\NormalTok{, }\AttributeTok{col =} \DecValTok{2}\NormalTok{, }\AttributeTok{cex =} \FloatTok{0.9}\NormalTok{)}
\NormalTok{  x0 }\OtherTok{\textless{}{-}}\NormalTok{ x1}
\NormalTok{\}}

\DocumentationTok{\#\# Newton step versus scoring step from the same point}
\FunctionTok{plot}\NormalTok{(lam, }\FunctionTok{S\_d}\NormalTok{(lam), }\AttributeTok{type =} \StringTok{"l"}\NormalTok{, }\AttributeTok{lwd =} \DecValTok{2}\NormalTok{, }\AttributeTok{ylim =} \FunctionTok{c}\NormalTok{(}\SpecialCharTok{{-}}\DecValTok{1}\NormalTok{, }\FloatTok{3.5}\NormalTok{),}
     \AttributeTok{xlab =} \FunctionTok{TeX}\NormalTok{(r}\StringTok{"($\textbackslash{}lambda$)"}\NormalTok{), }\AttributeTok{ylab =} \FunctionTok{TeX}\NormalTok{(r}\StringTok{"($S(\textbackslash{}lambda)$)"}\NormalTok{),}
     \AttributeTok{main =} \FunctionTok{TeX}\NormalTok{(r}\StringTok{"(observed versus expected curvature)"}\NormalTok{))}
\FunctionTok{abline}\NormalTok{(}\AttributeTok{h =} \DecValTok{0}\NormalTok{, }\AttributeTok{lty =} \DecValTok{3}\NormalTok{)}
\FunctionTok{points}\NormalTok{(lstar, }\DecValTok{0}\NormalTok{, }\AttributeTok{pch =} \DecValTok{19}\NormalTok{, }\AttributeTok{cex =} \FloatTok{1.2}\NormalTok{)}

\NormalTok{x0 }\OtherTok{\textless{}{-}} \FloatTok{0.45}
\NormalTok{x\_nr }\OtherTok{\textless{}{-}}\NormalTok{ x0 }\SpecialCharTok{{-}} \FunctionTok{S\_d}\NormalTok{(x0) }\SpecialCharTok{/} \FunctionTok{dS\_d}\NormalTok{(x0)          }\CommentTok{\# slope  S\textquotesingle{}(lambda)}
\NormalTok{x\_fs }\OtherTok{\textless{}{-}}\NormalTok{ x0 }\SpecialCharTok{+} \FunctionTok{S\_d}\NormalTok{(x0) }\SpecialCharTok{/} \FunctionTok{I\_d}\NormalTok{(x0)           }\CommentTok{\# slope {-}nI(lambda)}
\FunctionTok{segments}\NormalTok{(x0, }\FunctionTok{S\_d}\NormalTok{(x0), x\_nr, }\DecValTok{0}\NormalTok{, }\AttributeTok{col =} \DecValTok{2}\NormalTok{, }\AttributeTok{lwd =} \DecValTok{2}\NormalTok{)}
\FunctionTok{segments}\NormalTok{(x0, }\FunctionTok{S\_d}\NormalTok{(x0), x\_fs, }\DecValTok{0}\NormalTok{, }\AttributeTok{col =} \DecValTok{4}\NormalTok{, }\AttributeTok{lwd =} \DecValTok{2}\NormalTok{, }\AttributeTok{lty =} \DecValTok{5}\NormalTok{)}
\FunctionTok{points}\NormalTok{(}\FunctionTok{c}\NormalTok{(x0, x\_nr, x\_fs), }\FunctionTok{c}\NormalTok{(}\FunctionTok{S\_d}\NormalTok{(x0), }\DecValTok{0}\NormalTok{, }\DecValTok{0}\NormalTok{), }\AttributeTok{pch =} \DecValTok{19}\NormalTok{, }\AttributeTok{cex =} \FloatTok{0.7}\NormalTok{, }\AttributeTok{col =} \FunctionTok{c}\NormalTok{(}\DecValTok{1}\NormalTok{, }\DecValTok{2}\NormalTok{, }\DecValTok{4}\NormalTok{))}
\FunctionTok{legend}\NormalTok{(}\StringTok{"topright"}\NormalTok{, }\AttributeTok{bty =} \StringTok{"n"}\NormalTok{, }\AttributeTok{lwd =} \DecValTok{2}\NormalTok{, }\AttributeTok{lty =} \FunctionTok{c}\NormalTok{(}\DecValTok{1}\NormalTok{, }\DecValTok{5}\NormalTok{), }\AttributeTok{col =} \FunctionTok{c}\NormalTok{(}\DecValTok{2}\NormalTok{, }\DecValTok{4}\NormalTok{),}
       \AttributeTok{legend =} \FunctionTok{TeX}\NormalTok{(}\FunctionTok{c}\NormalTok{(r}\StringTok{"(Newton: slope $S\textquotesingle{}(\textbackslash{}lambda\_k)$)"}\NormalTok{, r}\StringTok{"(scoring: slope ${-}nI(\textbackslash{}lambda\_k)$)"}\NormalTok{)))}
\end{Highlighting}
\end{Shaded}

\pandocbounded{\includegraphics[keepaspectratio]{chapters/unit05_MLE1_files/figure-pdf/scheme-newton-1.pdf}}

\subsection{Advantage: quadratic
convergence}\label{advantage-quadratic-convergence}

Let \(e_k = \theta_k - \hat\theta\) and expand \(S(\hat\theta) = 0\)
about \(\theta_k\):
\[0 = S(\theta_k) - S'(\theta_k)e_k + \tfrac12 S''(\xi)e_k^{2}
\implies e_{k+1} = e_k - \frac{S(\theta_k)}{S'(\theta_k)} = \frac{S''(\xi)}{2S'(\theta_k)}\,e_k^{2} .\]

The error is squared at each step, so the number of correct digits
roughly doubles. Two conditions appear in the constant:
\(S'(\hat\theta) \ne 0\), meaning the maximum is a genuine peak and not
a flat spot, and \(S''\) bounded near the root. Asymptotic normality
says that log-likelihoods are approximately quadratic near the MLE,
which is why Newton is so effective there.

\begin{Shaded}
\begin{Highlighting}[]
\NormalTok{newton }\OtherTok{\textless{}{-}} \ControlFlowTok{function}\NormalTok{(S, dS, x0, }\AttributeTok{tol =} \FloatTok{1e{-}10}\NormalTok{, }\AttributeTok{maxit =} \DecValTok{100}\NormalTok{, }\AttributeTok{trace =} \ConstantTok{FALSE}\NormalTok{) \{}
  \ControlFlowTok{for}\NormalTok{ (k }\ControlFlowTok{in} \DecValTok{1}\SpecialCharTok{:}\NormalTok{maxit) \{}
\NormalTok{    x1 }\OtherTok{\textless{}{-}}\NormalTok{ x0 }\SpecialCharTok{{-}} \FunctionTok{S}\NormalTok{(x0) }\SpecialCharTok{/} \FunctionTok{dS}\NormalTok{(x0)}
    \ControlFlowTok{if}\NormalTok{ (trace) }\FunctionTok{cat}\NormalTok{(}\FunctionTok{sprintf}\NormalTok{(}\StringTok{"iter \%2d: \%.10f}\SpecialCharTok{\textbackslash{}n}\StringTok{"}\NormalTok{, k, x1))}
    \ControlFlowTok{if}\NormalTok{ (}\SpecialCharTok{!}\FunctionTok{is.finite}\NormalTok{(x1)) }\FunctionTok{stop}\NormalTok{(}\StringTok{"diverged"}\NormalTok{)}
    \ControlFlowTok{if}\NormalTok{ (}\FunctionTok{abs}\NormalTok{(x1 }\SpecialCharTok{{-}}\NormalTok{ x0) }\SpecialCharTok{\textless{}}\NormalTok{ tol) }\FunctionTok{return}\NormalTok{(}\FunctionTok{c}\NormalTok{(}\AttributeTok{root =}\NormalTok{ x1, }\AttributeTok{iter =}\NormalTok{ k))}
\NormalTok{    x0 }\OtherTok{\textless{}{-}}\NormalTok{ x1}
\NormalTok{  \}}
  \FunctionTok{stop}\NormalTok{(}\StringTok{"not converged"}\NormalTok{)}
\NormalTok{\}}
\end{Highlighting}
\end{Shaded}

\subsection{Problem: divergence}\label{problem-divergence}

Newton uses only local information, so the quadratic model can be badly
wrong away from the maximum. Three failures occur:

\begin{itemize}
\tightlist
\item
  \textbf{wrong curvature} --- where \(\ell'' > 0\) the model is a bowl,
  and the step heads toward a minimum or off to infinity;
\item
  \textbf{overshooting} --- where the score is nearly flat,
  \(S'(\theta_k) \approx 0\) produces an enormous step, possibly outside
  the parameter space;
\item
  \textbf{multiple roots} --- the iteration converges to whichever
  stationary point it happens to reach, which need not be the global
  maximum.
\end{itemize}

The single-observation Cauchy likelihood is the standard illustration.
With one observation at zero,
\(f(\theta) = -\ell(\theta) = \log(1 + \theta^{2})\) and
\(s(\theta) = f'(\theta) = 2\theta/(1 + \theta^{2})\), whose only root
is \(\theta = 0\). The score rises, turns over at \(\theta = 1\), and
decays to zero in the tails --- exactly the flat region that defeats
Newton.

\begin{Shaded}
\begin{Highlighting}[]
\NormalTok{lik }\OtherTok{\textless{}{-}} \ControlFlowTok{function}\NormalTok{(t) }\FunctionTok{log}\NormalTok{(}\DecValTok{1} \SpecialCharTok{+}\NormalTok{ t}\SpecialCharTok{\^{}}\DecValTok{2}\NormalTok{)                    }\CommentTok{\# negative log{-}likelihood}
\NormalTok{s   }\OtherTok{\textless{}{-}} \ControlFlowTok{function}\NormalTok{(t) }\DecValTok{2} \SpecialCharTok{*}\NormalTok{ t }\SpecialCharTok{/}\NormalTok{ (}\DecValTok{1} \SpecialCharTok{+}\NormalTok{ t}\SpecialCharTok{\^{}}\DecValTok{2}\NormalTok{)               }\CommentTok{\# its derivative}
\NormalTok{ds  }\OtherTok{\textless{}{-}} \ControlFlowTok{function}\NormalTok{(t) }\DecValTok{2} \SpecialCharTok{*}\NormalTok{ (}\DecValTok{1} \SpecialCharTok{{-}}\NormalTok{ t}\SpecialCharTok{\^{}}\DecValTok{2}\NormalTok{) }\SpecialCharTok{/}\NormalTok{ (}\DecValTok{1} \SpecialCharTok{+}\NormalTok{ t}\SpecialCharTok{\^{}}\DecValTok{2}\NormalTok{)}\SpecialCharTok{\^{}}\DecValTok{2}

\FunctionTok{par}\NormalTok{(}\AttributeTok{mfrow =} \FunctionTok{c}\NormalTok{(}\DecValTok{1}\NormalTok{, }\DecValTok{2}\NormalTok{), }\AttributeTok{mar =} \FunctionTok{c}\NormalTok{(}\DecValTok{4}\NormalTok{, }\FloatTok{4.5}\NormalTok{, }\FloatTok{2.5}\NormalTok{, }\DecValTok{1}\NormalTok{))}
\FunctionTok{curve}\NormalTok{(lik, }\SpecialCharTok{{-}}\DecValTok{5}\NormalTok{, }\DecValTok{5}\NormalTok{, }\AttributeTok{n =} \DecValTok{400}\NormalTok{, }\AttributeTok{lwd =} \DecValTok{2}\NormalTok{, }\AttributeTok{xlab =} \FunctionTok{TeX}\NormalTok{(r}\StringTok{"($}\SpecialCharTok{\textbackslash{}t}\StringTok{heta$)"}\NormalTok{), }\AttributeTok{ylab =} \FunctionTok{TeX}\NormalTok{(r}\StringTok{"(${-}}\SpecialCharTok{\textbackslash{}e}\StringTok{ll(}\SpecialCharTok{\textbackslash{}t}\StringTok{heta)$)"}\NormalTok{),}
      \AttributeTok{main =} \FunctionTok{TeX}\NormalTok{(r}\StringTok{"(negative log{-}likelihood)"}\NormalTok{))}
\FunctionTok{curve}\NormalTok{(s, }\SpecialCharTok{{-}}\DecValTok{5}\NormalTok{, }\DecValTok{5}\NormalTok{, }\AttributeTok{n =} \DecValTok{400}\NormalTok{, }\AttributeTok{lwd =} \DecValTok{2}\NormalTok{, }\AttributeTok{xlab =} \FunctionTok{TeX}\NormalTok{(r}\StringTok{"($}\SpecialCharTok{\textbackslash{}t}\StringTok{heta$)"}\NormalTok{), }\AttributeTok{ylab =} \FunctionTok{TeX}\NormalTok{(r}\StringTok{"($s(}\SpecialCharTok{\textbackslash{}t}\StringTok{heta)$)"}\NormalTok{),}
      \AttributeTok{main =} \FunctionTok{TeX}\NormalTok{(r}\StringTok{"(derivative, with the flat tails)"}\NormalTok{))}
\FunctionTok{abline}\NormalTok{(}\AttributeTok{h =} \DecValTok{0}\NormalTok{, }\AttributeTok{lty =} \DecValTok{3}\NormalTok{); }\FunctionTok{abline}\NormalTok{(}\AttributeTok{v =} \FunctionTok{c}\NormalTok{(}\SpecialCharTok{{-}}\DecValTok{1}\NormalTok{, }\DecValTok{1}\NormalTok{), }\AttributeTok{col =} \DecValTok{2}\NormalTok{, }\AttributeTok{lty =} \DecValTok{2}\NormalTok{)}
\end{Highlighting}
\end{Shaded}

\pandocbounded{\includegraphics[keepaspectratio]{chapters/unit05_MLE1_files/figure-pdf/cauchy-single-obs-1.pdf}}

Tracing the iterates from three starting values shows the boundary at
\(|\theta_0| = 1\), where the derivative \(s\) turns over:

\begin{Shaded}
\begin{Highlighting}[]
\NormalTok{newton\_trace }\OtherTok{\textless{}{-}} \ControlFlowTok{function}\NormalTok{(s, ds, x0, }\AttributeTok{iters =} \DecValTok{8}\NormalTok{) \{}
\NormalTok{  path }\OtherTok{\textless{}{-}} \FunctionTok{numeric}\NormalTok{(iters }\SpecialCharTok{+} \DecValTok{1}\NormalTok{); path[}\DecValTok{1}\NormalTok{] }\OtherTok{\textless{}{-}}\NormalTok{ x0}
  \ControlFlowTok{for}\NormalTok{ (k }\ControlFlowTok{in} \DecValTok{1}\SpecialCharTok{:}\NormalTok{iters) \{}
\NormalTok{    x0 }\OtherTok{\textless{}{-}}\NormalTok{ x0 }\SpecialCharTok{{-}} \FunctionTok{s}\NormalTok{(x0) }\SpecialCharTok{/} \FunctionTok{ds}\NormalTok{(x0)}
\NormalTok{    path[k }\SpecialCharTok{+} \DecValTok{1}\NormalTok{] }\OtherTok{\textless{}{-}} \ControlFlowTok{if}\NormalTok{ (}\FunctionTok{is.finite}\NormalTok{(x0)) x0 }\ControlFlowTok{else} \ConstantTok{NA}
\NormalTok{  \}}
\NormalTok{  path}
\NormalTok{\}}

\NormalTok{starts }\OtherTok{\textless{}{-}} \FunctionTok{c}\NormalTok{(}\FloatTok{0.5}\NormalTok{, }\FloatTok{0.8}\NormalTok{, }\FloatTok{1.2}\NormalTok{)}
\NormalTok{traces }\OtherTok{\textless{}{-}} \FunctionTok{vapply}\NormalTok{(starts, }\ControlFlowTok{function}\NormalTok{(v) }\FunctionTok{newton\_trace}\NormalTok{(s, ds, v), }\FunctionTok{numeric}\NormalTok{(}\DecValTok{9}\NormalTok{))}
\FunctionTok{colnames}\NormalTok{(traces) }\OtherTok{\textless{}{-}} \FunctionTok{sprintf}\NormalTok{(}\StringTok{"start = \%.1f"}\NormalTok{, starts)}
\FunctionTok{rownames}\NormalTok{(traces) }\OtherTok{\textless{}{-}} \FunctionTok{paste}\NormalTok{(}\StringTok{"iter"}\NormalTok{, }\DecValTok{0}\SpecialCharTok{:}\DecValTok{8}\NormalTok{)}
\FunctionTok{signif}\NormalTok{(traces, }\DecValTok{5}\NormalTok{)}
\end{Highlighting}
\end{Shaded}

\begin{verbatim}
       start = 0.5 start = 0.8 start = 1.2
iter 0  5.0000e-01      0.8000      1.2000
iter 1 -3.3333e-01     -2.8444      7.8545
iter 2  8.3333e-02     -6.4912     15.9680
iter 3 -1.1655e-03    -13.2980     32.0620
iter 4  3.1664e-09    -26.7470     64.1860
iter 5  0.0000e+00    -53.5690    128.4000
iter 6  0.0000e+00   -107.1800    256.8200
iter 7  0.0000e+00   -214.3700    513.6500
iter 8  0.0000e+00   -428.7500   1027.3000
\end{verbatim}

From \(\theta_0 = 0.5\), inside the region where the quadratic model is
adequate, Newton converges quadratically to zero. From
\(\theta_0 = 0.8\) the first step already overshoots into the tail, and
from \(\theta_0 = 1.2\) --- beyond the inflection point of \(s\) --- the
iterates run away, each step landing further out where the derivative is
flatter still.

\begin{Shaded}
\begin{Highlighting}[]
\DocumentationTok{\#\# Interactive versions of the same three trajectories (requires the animation}
\DocumentationTok{\#\# package and a graphics device that can write animations)}
\FunctionTok{library}\NormalTok{(animation)}
\FunctionTok{ani.options}\NormalTok{(}\AttributeTok{interval =} \DecValTok{1}\NormalTok{, }\AttributeTok{nmax =} \DecValTok{50}\NormalTok{)}
\FunctionTok{newton.method}\NormalTok{(}\AttributeTok{FUN =}\NormalTok{ s, }\AttributeTok{init =} \FloatTok{0.5}\NormalTok{, }\AttributeTok{rg =} \FunctionTok{c}\NormalTok{(}\SpecialCharTok{{-}}\FloatTok{1.5}\NormalTok{, }\FloatTok{1.5}\NormalTok{))}
\FunctionTok{newton.method}\NormalTok{(}\AttributeTok{FUN =}\NormalTok{ s, }\AttributeTok{init =} \FloatTok{0.8}\NormalTok{, }\AttributeTok{rg =} \FunctionTok{c}\NormalTok{(}\SpecialCharTok{{-}}\DecValTok{8}\NormalTok{, }\FloatTok{1.5}\NormalTok{))}
\FunctionTok{newton.method}\NormalTok{(}\AttributeTok{FUN =}\NormalTok{ s, }\AttributeTok{init =} \FloatTok{1.2}\NormalTok{, }\AttributeTok{rg =} \FunctionTok{c}\NormalTok{(}\SpecialCharTok{{-}}\DecValTok{2}\NormalTok{, }\DecValTok{20}\NormalTok{))}
\end{Highlighting}
\end{Shaded}

\subsection{Safeguards}\label{safeguards}

Three remedies cover most cases. Start from a robust estimate --- the
sample median for a location model. Apply \textbf{step halving}: accept
the full Newton step only if it improves the log-likelihood, otherwise
halve it until it does. Or replace the observed curvature by the
expected information, giving the \textbf{method of scoring},
\[\theta_{k+1} = \theta_k + \frac{S(\theta_k)}{n I(\theta_k)} ,\] in
which the divisor is positive by construction, so every step ascends.
Near the MLE \(nI(\hat\theta) \approx -\ell''(\hat\theta)\), so scoring
inherits Newton's speed while being far more stable away from the
optimum. Its cost is that the expectation must be available in closed
form.

\begin{Shaded}
\begin{Highlighting}[]
\NormalTok{safe\_newton }\OtherTok{\textless{}{-}} \ControlFlowTok{function}\NormalTok{(ll, S, dS, x0, }\AttributeTok{lower =} \SpecialCharTok{{-}}\ConstantTok{Inf}\NormalTok{, }\AttributeTok{upper =} \ConstantTok{Inf}\NormalTok{,}
                        \AttributeTok{tol =} \FloatTok{1e{-}10}\NormalTok{, }\AttributeTok{maxit =} \DecValTok{100}\NormalTok{, }\AttributeTok{max\_halve =} \DecValTok{40}\NormalTok{) \{}
\NormalTok{  f0 }\OtherTok{\textless{}{-}} \FunctionTok{ll}\NormalTok{(x0)}
  \ControlFlowTok{for}\NormalTok{ (k }\ControlFlowTok{in} \DecValTok{1}\SpecialCharTok{:}\NormalTok{maxit) \{}
\NormalTok{    step }\OtherTok{\textless{}{-}} \SpecialCharTok{{-}}\FunctionTok{S}\NormalTok{(x0) }\SpecialCharTok{/} \FunctionTok{dS}\NormalTok{(x0)}
    \ControlFlowTok{if}\NormalTok{ (}\FunctionTok{dS}\NormalTok{(x0) }\SpecialCharTok{\textgreater{}} \DecValTok{0}\NormalTok{) step }\OtherTok{\textless{}{-}} \FunctionTok{sign}\NormalTok{(}\FunctionTok{S}\NormalTok{(x0))          }\CommentTok{\# wrong curvature: just move uphill}
    \ControlFlowTok{for}\NormalTok{ (h }\ControlFlowTok{in} \DecValTok{0}\SpecialCharTok{:}\NormalTok{max\_halve) \{                     }\CommentTok{\# halve until feasible and improving}
\NormalTok{      x1 }\OtherTok{\textless{}{-}}\NormalTok{ x0 }\SpecialCharTok{+}\NormalTok{ step }\SpecialCharTok{/} \DecValTok{2}\SpecialCharTok{\^{}}\NormalTok{h}
      \ControlFlowTok{if}\NormalTok{ (x1 }\SpecialCharTok{\textgreater{}}\NormalTok{ lower }\SpecialCharTok{\&\&}\NormalTok{ x1 }\SpecialCharTok{\textless{}}\NormalTok{ upper }\SpecialCharTok{\&\&} \FunctionTok{is.finite}\NormalTok{(}\FunctionTok{ll}\NormalTok{(x1)) }\SpecialCharTok{\&\&} \FunctionTok{ll}\NormalTok{(x1) }\SpecialCharTok{\textgreater{}=}\NormalTok{ f0) }\ControlFlowTok{break}
\NormalTok{    \}}
    \ControlFlowTok{if}\NormalTok{ (}\FunctionTok{abs}\NormalTok{(x1 }\SpecialCharTok{{-}}\NormalTok{ x0) }\SpecialCharTok{\textless{}}\NormalTok{ tol) }\FunctionTok{return}\NormalTok{(}\FunctionTok{c}\NormalTok{(}\AttributeTok{root =}\NormalTok{ x1, }\AttributeTok{iter =}\NormalTok{ k, }\AttributeTok{halvings =}\NormalTok{ h))}
\NormalTok{    x0 }\OtherTok{\textless{}{-}}\NormalTok{ x1; f0 }\OtherTok{\textless{}{-}} \FunctionTok{ll}\NormalTok{(x1)}
\NormalTok{  \}}
  \FunctionTok{stop}\NormalTok{(}\StringTok{"not converged"}\NormalTok{)}
\NormalTok{\}}
\end{Highlighting}
\end{Shaded}

\subsection{Golden Section Search}\label{golden-section-search}

Golden section minimizes \(f = -\ell\) directly, using \textbf{no
derivative at all}, which makes it the method of choice when \(\ell\) is
not differentiable or its derivative is unreliable. Keep a bracket
\([x_l, x_u]\) containing the minimum of a unimodal \(f\), with two
interior points placed at the golden ratio
\(g = (3 - \sqrt5)/2 \approx 0.382\) of the way in from each end.
Comparing the two interior values discards one end, and the ratio is
chosen precisely so that the surviving interior point is correctly
positioned for the next iteration: one new evaluation per step, with the
bracket shrinking by the factor \(0.618\).

\begin{Shaded}
\begin{Highlighting}[]
\NormalTok{golden }\OtherTok{\textless{}{-}} \ControlFlowTok{function}\NormalTok{(f, brack.int, }\AttributeTok{eps =} \FloatTok{1e{-}4}\NormalTok{, ...) \{}
\NormalTok{  g }\OtherTok{\textless{}{-}}\NormalTok{ (}\DecValTok{3} \SpecialCharTok{{-}} \FunctionTok{sqrt}\NormalTok{(}\DecValTok{5}\NormalTok{)) }\SpecialCharTok{/} \DecValTok{2}
\NormalTok{  xl }\OtherTok{\textless{}{-}} \FunctionTok{min}\NormalTok{(brack.int); xu }\OtherTok{\textless{}{-}} \FunctionTok{max}\NormalTok{(brack.int)}
\NormalTok{  tmp }\OtherTok{\textless{}{-}}\NormalTok{ g }\SpecialCharTok{*}\NormalTok{ (xu }\SpecialCharTok{{-}}\NormalTok{ xl)}
\NormalTok{  xmu }\OtherTok{\textless{}{-}}\NormalTok{ xu }\SpecialCharTok{{-}}\NormalTok{ tmp; xml }\OtherTok{\textless{}{-}}\NormalTok{ xl }\SpecialCharTok{+}\NormalTok{ tmp}
\NormalTok{  fl }\OtherTok{\textless{}{-}} \FunctionTok{f}\NormalTok{(xml, ...); fu }\OtherTok{\textless{}{-}} \FunctionTok{f}\NormalTok{(xmu, ...)}
  \ControlFlowTok{while}\NormalTok{ (}\FunctionTok{abs}\NormalTok{(xu }\SpecialCharTok{{-}}\NormalTok{ xl) }\SpecialCharTok{\textgreater{}}\NormalTok{ (}\FloatTok{1e{-}5} \SpecialCharTok{+} \FunctionTok{abs}\NormalTok{(xl)) }\SpecialCharTok{*}\NormalTok{ eps) \{}
    \ControlFlowTok{if}\NormalTok{ (fl }\SpecialCharTok{\textless{}}\NormalTok{ fu) \{}
\NormalTok{      xu }\OtherTok{\textless{}{-}}\NormalTok{ xmu; xmu }\OtherTok{\textless{}{-}}\NormalTok{ xml; fu }\OtherTok{\textless{}{-}}\NormalTok{ fl}
\NormalTok{      xml }\OtherTok{\textless{}{-}}\NormalTok{ xl }\SpecialCharTok{+}\NormalTok{ g }\SpecialCharTok{*}\NormalTok{ (xu }\SpecialCharTok{{-}}\NormalTok{ xl); fl }\OtherTok{\textless{}{-}} \FunctionTok{f}\NormalTok{(xml, ...)}
\NormalTok{    \} }\ControlFlowTok{else}\NormalTok{ \{}
\NormalTok{      xl }\OtherTok{\textless{}{-}}\NormalTok{ xml; xml }\OtherTok{\textless{}{-}}\NormalTok{ xmu; fl }\OtherTok{\textless{}{-}}\NormalTok{ fu}
\NormalTok{      xmu }\OtherTok{\textless{}{-}}\NormalTok{ xu }\SpecialCharTok{{-}}\NormalTok{ g }\SpecialCharTok{*}\NormalTok{ (xu }\SpecialCharTok{{-}}\NormalTok{ xl); fu }\OtherTok{\textless{}{-}} \FunctionTok{f}\NormalTok{(xmu, ...)}
\NormalTok{    \}}
\NormalTok{  \}}
  \ControlFlowTok{if}\NormalTok{ (fl }\SpecialCharTok{\textless{}}\NormalTok{ fu) xml }\ControlFlowTok{else}\NormalTok{ xmu}
\NormalTok{\}}
\end{Highlighting}
\end{Shaded}

The schematic makes the reuse visible. Each bar is the current bracket
with its two interior points; the point that survives into the next
iteration is drawn in blue and sits at exactly the right position there,
so only the red point has to be evaluated afresh.

\begin{Shaded}
\begin{Highlighting}[]
\NormalTok{f\_d }\OtherTok{\textless{}{-}} \ControlFlowTok{function}\NormalTok{(l) }\SpecialCharTok{{-}}\FunctionTok{ll\_d}\NormalTok{(l)}
\NormalTok{gr }\OtherTok{\textless{}{-}}\NormalTok{ (}\DecValTok{3} \SpecialCharTok{{-}} \FunctionTok{sqrt}\NormalTok{(}\DecValTok{5}\NormalTok{)) }\SpecialCharTok{/} \DecValTok{2}

\FunctionTok{par}\NormalTok{(}\AttributeTok{mar =} \FunctionTok{c}\NormalTok{(}\DecValTok{4}\NormalTok{, }\FloatTok{4.5}\NormalTok{, }\FloatTok{2.5}\NormalTok{, }\DecValTok{1}\NormalTok{))}
\FunctionTok{plot}\NormalTok{(lam, }\FunctionTok{f\_d}\NormalTok{(lam), }\AttributeTok{type =} \StringTok{"l"}\NormalTok{, }\AttributeTok{lwd =} \DecValTok{2}\NormalTok{, }\AttributeTok{ylim =} \FunctionTok{c}\NormalTok{(}\FloatTok{0.2}\NormalTok{, }\FloatTok{1.25}\NormalTok{),}
     \AttributeTok{xlab =} \FunctionTok{TeX}\NormalTok{(r}\StringTok{"($\textbackslash{}lambda$)"}\NormalTok{), }\AttributeTok{ylab =} \FunctionTok{TeX}\NormalTok{(r}\StringTok{"(${-}}\SpecialCharTok{\textbackslash{}e}\StringTok{ll(\textbackslash{}lambda)$)"}\NormalTok{),}
     \AttributeTok{main =} \FunctionTok{TeX}\NormalTok{(r}\StringTok{"(golden section: one new evaluation per step)"}\NormalTok{))}
\FunctionTok{points}\NormalTok{(lstar, }\FunctionTok{f\_d}\NormalTok{(lstar), }\AttributeTok{pch =} \DecValTok{19}\NormalTok{, }\AttributeTok{cex =} \FloatTok{1.2}\NormalTok{)}

\NormalTok{xl }\OtherTok{\textless{}{-}} \FloatTok{0.3}\NormalTok{; xu }\OtherTok{\textless{}{-}} \DecValTok{3}
\NormalTok{ylev }\OtherTok{\textless{}{-}} \FunctionTok{seq}\NormalTok{(}\FloatTok{0.44}\NormalTok{, }\FloatTok{0.24}\NormalTok{, }\AttributeTok{length.out =} \DecValTok{3}\NormalTok{)}
\ControlFlowTok{for}\NormalTok{ (k }\ControlFlowTok{in} \DecValTok{1}\SpecialCharTok{:}\DecValTok{3}\NormalTok{) \{}
\NormalTok{  xml }\OtherTok{\textless{}{-}}\NormalTok{ xl }\SpecialCharTok{+}\NormalTok{ gr }\SpecialCharTok{*}\NormalTok{ (xu }\SpecialCharTok{{-}}\NormalTok{ xl); xmu }\OtherTok{\textless{}{-}}\NormalTok{ xu }\SpecialCharTok{{-}}\NormalTok{ gr }\SpecialCharTok{*}\NormalTok{ (xu }\SpecialCharTok{{-}}\NormalTok{ xl)}
  \FunctionTok{segments}\NormalTok{(xl, ylev[k], xu, ylev[k], }\AttributeTok{col =} \DecValTok{4}\NormalTok{, }\AttributeTok{lwd =} \DecValTok{3}\NormalTok{)}
  \FunctionTok{points}\NormalTok{(}\FunctionTok{c}\NormalTok{(xl, xu), }\FunctionTok{rep}\NormalTok{(ylev[k], }\DecValTok{2}\NormalTok{), }\AttributeTok{pch =} \StringTok{"|"}\NormalTok{, }\AttributeTok{col =} \DecValTok{4}\NormalTok{)}
\NormalTok{  keep\_lower }\OtherTok{\textless{}{-}} \FunctionTok{f\_d}\NormalTok{(xml) }\SpecialCharTok{\textless{}} \FunctionTok{f\_d}\NormalTok{(xmu)}
  \FunctionTok{points}\NormalTok{(xml, ylev[k], }\AttributeTok{pch =} \DecValTok{19}\NormalTok{, }\AttributeTok{cex =} \FloatTok{0.9}\NormalTok{, }\AttributeTok{col =} \ControlFlowTok{if}\NormalTok{ (keep\_lower) }\DecValTok{4} \ControlFlowTok{else} \DecValTok{2}\NormalTok{)}
  \FunctionTok{points}\NormalTok{(xmu, ylev[k], }\AttributeTok{pch =} \DecValTok{19}\NormalTok{, }\AttributeTok{cex =} \FloatTok{0.9}\NormalTok{, }\AttributeTok{col =} \ControlFlowTok{if}\NormalTok{ (keep\_lower) }\DecValTok{2} \ControlFlowTok{else} \DecValTok{4}\NormalTok{)}
  \FunctionTok{text}\NormalTok{(}\FloatTok{0.25}\NormalTok{, ylev[k], k, }\AttributeTok{cex =} \FloatTok{0.8}\NormalTok{, }\AttributeTok{col =} \DecValTok{4}\NormalTok{)}
  \ControlFlowTok{if}\NormalTok{ (keep\_lower) xu }\OtherTok{\textless{}{-}}\NormalTok{ xmu }\ControlFlowTok{else}\NormalTok{ xl }\OtherTok{\textless{}{-}}\NormalTok{ xml}
\NormalTok{\}}
\end{Highlighting}
\end{Shaded}

\pandocbounded{\includegraphics[keepaspectratio]{chapters/unit05_MLE1_files/figure-pdf/scheme-golden-1.pdf}}

The first test is the negative log-likelihood of the classical genetic
linkage model, with cell counts \((1997, 906, 904, 32)\) and
probabilities
\(\left((2+\theta)/4, (1-\theta)/4, (1-\theta)/4, \theta/4\right)\); it
is smooth, so any method would work. The second, \(f(x) = |x - 10|\), is
not differentiable at its minimum, so no root-finder applies --- golden
section handles it without modification.

\begin{Shaded}
\begin{Highlighting}[]
\NormalTok{f1 }\OtherTok{\textless{}{-}} \ControlFlowTok{function}\NormalTok{(theta) }\SpecialCharTok{{-}}\DecValTok{1997} \SpecialCharTok{*} \FunctionTok{log}\NormalTok{(}\DecValTok{2} \SpecialCharTok{+}\NormalTok{ theta) }\SpecialCharTok{{-}} \DecValTok{1810} \SpecialCharTok{*} \FunctionTok{log}\NormalTok{(}\DecValTok{1} \SpecialCharTok{{-}}\NormalTok{ theta) }\SpecialCharTok{{-}} \DecValTok{32} \SpecialCharTok{*} \FunctionTok{log}\NormalTok{(theta)}
\NormalTok{f2 }\OtherTok{\textless{}{-}} \ControlFlowTok{function}\NormalTok{(x) }\FunctionTok{abs}\NormalTok{(x }\SpecialCharTok{{-}} \DecValTok{10}\NormalTok{)}

\FunctionTok{par}\NormalTok{(}\AttributeTok{mfrow =} \FunctionTok{c}\NormalTok{(}\DecValTok{1}\NormalTok{, }\DecValTok{2}\NormalTok{), }\AttributeTok{mar =} \FunctionTok{c}\NormalTok{(}\DecValTok{4}\NormalTok{, }\FloatTok{4.5}\NormalTok{, }\FloatTok{2.5}\NormalTok{, }\DecValTok{1}\NormalTok{))}
\FunctionTok{curve}\NormalTok{(f1, }\FloatTok{0.001}\NormalTok{, }\FloatTok{0.999}\NormalTok{, }\AttributeTok{n =} \DecValTok{400}\NormalTok{, }\AttributeTok{lwd =} \DecValTok{2}\NormalTok{, }\AttributeTok{xlab =} \FunctionTok{TeX}\NormalTok{(r}\StringTok{"($}\SpecialCharTok{\textbackslash{}t}\StringTok{heta$)"}\NormalTok{), }\AttributeTok{ylab =} \FunctionTok{TeX}\NormalTok{(r}\StringTok{"(${-}}\SpecialCharTok{\textbackslash{}e}\StringTok{ll(}\SpecialCharTok{\textbackslash{}t}\StringTok{heta)$)"}\NormalTok{),}
      \AttributeTok{main =} \FunctionTok{TeX}\NormalTok{(r}\StringTok{"(genetic linkage model)"}\NormalTok{))}
\FunctionTok{abline}\NormalTok{(}\AttributeTok{v =} \FunctionTok{golden}\NormalTok{(f1, }\FunctionTok{c}\NormalTok{(}\FloatTok{0.001}\NormalTok{, }\FloatTok{0.999}\NormalTok{)), }\AttributeTok{col =} \DecValTok{2}\NormalTok{, }\AttributeTok{lty =} \DecValTok{2}\NormalTok{, }\AttributeTok{lwd =} \DecValTok{2}\NormalTok{)}

\FunctionTok{curve}\NormalTok{(f2, }\DecValTok{0}\NormalTok{, }\DecValTok{20}\NormalTok{, }\AttributeTok{n =} \DecValTok{400}\NormalTok{, }\AttributeTok{lwd =} \DecValTok{2}\NormalTok{, }\AttributeTok{xlab =} \FunctionTok{TeX}\NormalTok{(r}\StringTok{"($x$)"}\NormalTok{), }\AttributeTok{ylab =} \FunctionTok{TeX}\NormalTok{(r}\StringTok{"($|x {-} 10|$)"}\NormalTok{),}
      \AttributeTok{main =} \FunctionTok{TeX}\NormalTok{(r}\StringTok{"(not differentiable at the minimum)"}\NormalTok{))}
\FunctionTok{abline}\NormalTok{(}\AttributeTok{v =} \FunctionTok{golden}\NormalTok{(f2, }\FunctionTok{c}\NormalTok{(}\DecValTok{0}\NormalTok{, }\DecValTok{20}\NormalTok{)), }\AttributeTok{col =} \DecValTok{2}\NormalTok{, }\AttributeTok{lty =} \DecValTok{2}\NormalTok{, }\AttributeTok{lwd =} \DecValTok{2}\NormalTok{)}
\end{Highlighting}
\end{Shaded}

\pandocbounded{\includegraphics[keepaspectratio]{chapters/unit05_MLE1_files/figure-pdf/golden-tests-1.pdf}}

\begin{Shaded}
\begin{Highlighting}[]
\FunctionTok{c}\NormalTok{(}\AttributeTok{linkage =} \FunctionTok{golden}\NormalTok{(f1, }\FunctionTok{c}\NormalTok{(}\FloatTok{0.001}\NormalTok{, }\FloatTok{0.999}\NormalTok{)), }\AttributeTok{abs\_value =} \FunctionTok{golden}\NormalTok{(f2, }\FunctionTok{c}\NormalTok{(}\DecValTok{0}\NormalTok{, }\DecValTok{20}\NormalTok{)))}
\end{Highlighting}
\end{Shaded}

\begin{verbatim}
    linkage   abs_value 
 0.03571247 10.00009645 
\end{verbatim}

\subsection{Summary of the Methods}\label{summary-of-the-methods}

{\def\LTcaptype{none} % do not increment counter
\begin{longtable}[]{@{}
  >{\raggedright\arraybackslash}p{(\linewidth - 6\tabcolsep) * \real{0.2500}}
  >{\raggedright\arraybackslash}p{(\linewidth - 6\tabcolsep) * \real{0.2500}}
  >{\raggedright\arraybackslash}p{(\linewidth - 6\tabcolsep) * \real{0.2500}}
  >{\raggedright\arraybackslash}p{(\linewidth - 6\tabcolsep) * \real{0.2500}}@{}}
\toprule\noalign{}
\begin{minipage}[b]{\linewidth}\raggedright
method
\end{minipage} & \begin{minipage}[b]{\linewidth}\raggedright
needs
\end{minipage} & \begin{minipage}[b]{\linewidth}\raggedright
convergence
\end{minipage} & \begin{minipage}[b]{\linewidth}\raggedright
fails when
\end{minipage} \\
\midrule\noalign{}
\endhead
\bottomrule\noalign{}
\endlastfoot
fixed point & a rearrangement \(\theta = g(\theta)\) & linear, factor
\(\lvert g'\rvert\) & \(\lvert g'\rvert \ge 1\) \\
bisection & \(S\) and a sign-changing bracket & linear, factor \(1/2\) &
no bracket available \\
Newton--Raphson & \(S\) and \(S'\) & quadratic near the root & flat or
wrong curvature \\
method of scoring & \(S\) and \(I(\theta)\) & linear to quadratic &
\(E[\cdot]\) unavailable \\
golden section & \(\ell\) only, and a bracket & linear, factor \(0.618\)
& \(\ell\) not unimodal \\
\end{longtable}
}

\section{Example: Poisson Data With Zeros
Unobserved}\label{example-poisson-data-with-zeros-unobserved}

\subsection{The Model}\label{the-model}

Counts are often recorded only when positive: a study of hospital stays
sees admitted patients, a survey of litters sees litters that exist. If
\(Y \sim \mathrm{Poisson}(\lambda)\) and we observe
\(X = Y \mid Y > 0\),
\[P(x \mid \lambda) = \frac{e^{-\lambda}\lambda^{x}}{x!}\cdot\frac{1}{1 - e^{-\lambda}}, \qquad x = 1, 2, \ldots\]

Dropping the factorial terms, which do not involve \(\lambda\),
\[\ell(\lambda) = \left(\sum_i x_i\right)\log\lambda - n\left[\lambda + \log\left(1 - e^{-\lambda}\right)\right],\]
\[S(\lambda) = \frac{\sum_i x_i}{\lambda} - \frac{n}{1 - e^{-\lambda}}, \qquad
S'(\lambda) = -\frac{\sum_i x_i}{\lambda^{2}} + \frac{n e^{-\lambda}}{\left(1 - e^{-\lambda}\right)^{2}} .\]

Since \(E(X) = \lambda/(1 - e^{-\lambda})\), taking \(-E[S'(\lambda)]\)
gives the expected information
\[n I(\lambda) = \frac{n}{\lambda\left(1 - e^{-\lambda}\right)} - \frac{n e^{-\lambda}}{\left(1 - e^{-\lambda}\right)^{2}} .\]

The score equation has no closed-form root, but the likelihood is
unimodal --- an ideal problem on which to compare all the methods
against a single known answer.

\begin{Shaded}
\begin{Highlighting}[]
\DocumentationTok{\#\# log likelihood, omitting factorial terms; vectorized in lambda}
\NormalTok{nzp.log.likelihood }\OtherTok{\textless{}{-}} \ControlFlowTok{function}\NormalTok{(n, lambda)}
  \FunctionTok{sum}\NormalTok{(n) }\SpecialCharTok{*} \FunctionTok{log}\NormalTok{(lambda) }\SpecialCharTok{{-}} \FunctionTok{length}\NormalTok{(n) }\SpecialCharTok{*}\NormalTok{ (lambda }\SpecialCharTok{+} \FunctionTok{log}\NormalTok{(}\DecValTok{1} \SpecialCharTok{{-}} \FunctionTok{exp}\NormalTok{(}\SpecialCharTok{{-}}\NormalTok{lambda)))}

\NormalTok{nzp.scorefunction }\OtherTok{\textless{}{-}} \ControlFlowTok{function}\NormalTok{(n, lambda)}
  \FunctionTok{sum}\NormalTok{(n) }\SpecialCharTok{/}\NormalTok{ lambda }\SpecialCharTok{{-}} \FunctionTok{length}\NormalTok{(n) }\SpecialCharTok{/}\NormalTok{ (}\DecValTok{1} \SpecialCharTok{{-}} \FunctionTok{exp}\NormalTok{(}\SpecialCharTok{{-}}\NormalTok{lambda))}

\NormalTok{nzp.dscore }\OtherTok{\textless{}{-}} \ControlFlowTok{function}\NormalTok{(n, lambda)}
  \SpecialCharTok{{-}}\FunctionTok{sum}\NormalTok{(n) }\SpecialCharTok{/}\NormalTok{ lambda}\SpecialCharTok{\^{}}\DecValTok{2} \SpecialCharTok{+} \FunctionTok{length}\NormalTok{(n) }\SpecialCharTok{*} \FunctionTok{exp}\NormalTok{(}\SpecialCharTok{{-}}\NormalTok{lambda) }\SpecialCharTok{/}\NormalTok{ (}\DecValTok{1} \SpecialCharTok{{-}} \FunctionTok{exp}\NormalTok{(}\SpecialCharTok{{-}}\NormalTok{lambda))}\SpecialCharTok{\^{}}\DecValTok{2}

\NormalTok{nzp.information }\OtherTok{\textless{}{-}} \ControlFlowTok{function}\NormalTok{(n, lambda)}
  \FunctionTok{length}\NormalTok{(n) }\SpecialCharTok{/}\NormalTok{ (lambda }\SpecialCharTok{*}\NormalTok{ (}\DecValTok{1} \SpecialCharTok{{-}} \FunctionTok{exp}\NormalTok{(}\SpecialCharTok{{-}}\NormalTok{lambda))) }\SpecialCharTok{{-}} \FunctionTok{length}\NormalTok{(n) }\SpecialCharTok{*} \FunctionTok{exp}\NormalTok{(}\SpecialCharTok{{-}}\NormalTok{lambda) }\SpecialCharTok{/}\NormalTok{ (}\DecValTok{1} \SpecialCharTok{{-}} \FunctionTok{exp}\NormalTok{(}\SpecialCharTok{{-}}\NormalTok{lambda))}\SpecialCharTok{\^{}}\DecValTok{2}

\DocumentationTok{\#\# simulated data: Poisson counts with the zeros discarded}
\NormalTok{x }\OtherTok{\textless{}{-}} \FunctionTok{rpois}\NormalTok{(}\DecValTok{100}\NormalTok{, }\AttributeTok{lambda =} \FloatTok{1.5}\NormalTok{)}
\NormalTok{xp }\OtherTok{\textless{}{-}}\NormalTok{ x[x }\SpecialCharTok{\textgreater{}} \DecValTok{0}\NormalTok{]}

\FunctionTok{c}\NormalTok{(}\AttributeTok{n\_observed =} \FunctionTok{length}\NormalTok{(xp), }\AttributeTok{sample\_mean =} \FunctionTok{mean}\NormalTok{(xp), }\AttributeTok{true\_lambda =} \FloatTok{1.5}\NormalTok{)}
\end{Highlighting}
\end{Shaded}

\begin{verbatim}
 n_observed sample_mean true_lambda 
  83.000000    2.084337    1.500000 
\end{verbatim}

The sample mean overshoots 1.5, because the zeros that would have pulled
it down were never recorded.

\begin{Shaded}
\begin{Highlighting}[]
\NormalTok{lambdas }\OtherTok{\textless{}{-}} \FunctionTok{seq}\NormalTok{(}\FloatTok{0.1}\NormalTok{, }\DecValTok{8}\NormalTok{, }\AttributeTok{length.out =} \DecValTok{300}\NormalTok{)}

\FunctionTok{par}\NormalTok{(}\AttributeTok{mfrow =} \FunctionTok{c}\NormalTok{(}\DecValTok{1}\NormalTok{, }\DecValTok{2}\NormalTok{), }\AttributeTok{mar =} \FunctionTok{c}\NormalTok{(}\DecValTok{4}\NormalTok{, }\FloatTok{4.5}\NormalTok{, }\FloatTok{2.5}\NormalTok{, }\DecValTok{1}\NormalTok{))}
\FunctionTok{plot}\NormalTok{(lambdas, }\FunctionTok{nzp.log.likelihood}\NormalTok{(xp, lambdas), }\AttributeTok{type =} \StringTok{"l"}\NormalTok{, }\AttributeTok{lwd =} \DecValTok{2}\NormalTok{,}
     \AttributeTok{xlab =} \FunctionTok{TeX}\NormalTok{(r}\StringTok{"($\textbackslash{}lambda$)"}\NormalTok{), }\AttributeTok{ylab =} \FunctionTok{TeX}\NormalTok{(r}\StringTok{"($}\SpecialCharTok{\textbackslash{}e}\StringTok{ll(\textbackslash{}lambda)$)"}\NormalTok{), }\AttributeTok{main =} \FunctionTok{TeX}\NormalTok{(r}\StringTok{"(log{-}likelihood)"}\NormalTok{))}
\FunctionTok{plot}\NormalTok{(lambdas, }\FunctionTok{nzp.scorefunction}\NormalTok{(xp, lambdas), }\AttributeTok{type =} \StringTok{"l"}\NormalTok{, }\AttributeTok{lwd =} \DecValTok{2}\NormalTok{,}
     \AttributeTok{xlab =} \FunctionTok{TeX}\NormalTok{(r}\StringTok{"($\textbackslash{}lambda$)"}\NormalTok{), }\AttributeTok{ylab =} \FunctionTok{TeX}\NormalTok{(r}\StringTok{"($S(\textbackslash{}lambda)$)"}\NormalTok{), }\AttributeTok{main =} \FunctionTok{TeX}\NormalTok{(r}\StringTok{"(score)"}\NormalTok{))}
\FunctionTok{abline}\NormalTok{(}\AttributeTok{h =} \DecValTok{0}\NormalTok{, }\AttributeTok{lty =} \DecValTok{3}\NormalTok{)}
\end{Highlighting}
\end{Shaded}

\pandocbounded{\includegraphics[keepaspectratio]{chapters/unit05_MLE1_files/figure-pdf/nzp-loglik-score-1.pdf}}

\subsection{The Three Iterations}\label{the-three-iterations}

Each method below returns the whole sequence of iterates together with
the log-likelihood attained, so the trajectories can be compared
directly.

\textbf{Simple iteration} rearranges \(S(\lambda) = 0\) as
\[\frac{\bar x}{\lambda} = \frac{1}{1 - e^{-\lambda}} \iff \lambda = \bar x\left(1 - e^{-\lambda}\right) = g(\lambda),\]
with contraction factor \(g'(\lambda) = \bar x e^{-\lambda}\).

\begin{Shaded}
\begin{Highlighting}[]
\NormalTok{nzp.simple.iteration }\OtherTok{\textless{}{-}} \ControlFlowTok{function}\NormalTok{(n, r, }\AttributeTok{lambda0 =} \FunctionTok{mean}\NormalTok{(n)) \{}
\NormalTok{  mean.n }\OtherTok{\textless{}{-}} \FunctionTok{mean}\NormalTok{(n); lambda }\OtherTok{\textless{}{-}} \FunctionTok{numeric}\NormalTok{(r }\SpecialCharTok{+} \DecValTok{1}\NormalTok{); lambda[}\DecValTok{1}\NormalTok{] }\OtherTok{\textless{}{-}}\NormalTok{ lambda0}
  \ControlFlowTok{for}\NormalTok{ (i }\ControlFlowTok{in} \DecValTok{1}\SpecialCharTok{:}\NormalTok{r) lambda[i }\SpecialCharTok{+} \DecValTok{1}\NormalTok{] }\OtherTok{\textless{}{-}}\NormalTok{ mean.n }\SpecialCharTok{*}\NormalTok{ (}\DecValTok{1} \SpecialCharTok{{-}} \FunctionTok{exp}\NormalTok{(}\SpecialCharTok{{-}}\NormalTok{lambda[i]))}
  \FunctionTok{data.frame}\NormalTok{(}\AttributeTok{lambda =}\NormalTok{ lambda, }\AttributeTok{log.lik =} \FunctionTok{nzp.log.likelihood}\NormalTok{(n, lambda))}
\NormalTok{\}}

\NormalTok{nzp.newton.raphson }\OtherTok{\textless{}{-}} \ControlFlowTok{function}\NormalTok{(n, r, }\AttributeTok{lambda0 =} \FunctionTok{mean}\NormalTok{(n)) \{}
\NormalTok{  mean.n }\OtherTok{\textless{}{-}} \FunctionTok{mean}\NormalTok{(n); lambda }\OtherTok{\textless{}{-}} \FunctionTok{numeric}\NormalTok{(r }\SpecialCharTok{+} \DecValTok{1}\NormalTok{); lambda[}\DecValTok{1}\NormalTok{] }\OtherTok{\textless{}{-}}\NormalTok{ lambda0}
  \ControlFlowTok{for}\NormalTok{ (i }\ControlFlowTok{in} \DecValTok{1}\SpecialCharTok{:}\NormalTok{r) \{}
\NormalTok{    e }\OtherTok{\textless{}{-}} \FunctionTok{exp}\NormalTok{(}\SpecialCharTok{{-}}\NormalTok{lambda[i])}
\NormalTok{    lambda[i }\SpecialCharTok{+} \DecValTok{1}\NormalTok{] }\OtherTok{\textless{}{-}}\NormalTok{ lambda[i] }\SpecialCharTok{{-}}
\NormalTok{      (mean.n }\SpecialCharTok{/}\NormalTok{ lambda[i] }\SpecialCharTok{{-}} \DecValTok{1} \SpecialCharTok{/}\NormalTok{ (}\DecValTok{1} \SpecialCharTok{{-}}\NormalTok{ e)) }\SpecialCharTok{/}\NormalTok{ (e }\SpecialCharTok{/}\NormalTok{ (}\DecValTok{1} \SpecialCharTok{{-}}\NormalTok{ e)}\SpecialCharTok{\^{}}\DecValTok{2} \SpecialCharTok{{-}}\NormalTok{ mean.n }\SpecialCharTok{/}\NormalTok{ lambda[i]}\SpecialCharTok{\^{}}\DecValTok{2}\NormalTok{)}
\NormalTok{  \}}
  \FunctionTok{data.frame}\NormalTok{(}\AttributeTok{lambda =}\NormalTok{ lambda, }\AttributeTok{log.lik =} \FunctionTok{nzp.log.likelihood}\NormalTok{(n, lambda))}
\NormalTok{\}}

\NormalTok{nzp.method.of.scoring }\OtherTok{\textless{}{-}} \ControlFlowTok{function}\NormalTok{(n, r, }\AttributeTok{lambda0 =} \FunctionTok{mean}\NormalTok{(n)) \{}
\NormalTok{  mean.n }\OtherTok{\textless{}{-}} \FunctionTok{mean}\NormalTok{(n); lambda }\OtherTok{\textless{}{-}} \FunctionTok{numeric}\NormalTok{(r }\SpecialCharTok{+} \DecValTok{1}\NormalTok{); lambda[}\DecValTok{1}\NormalTok{] }\OtherTok{\textless{}{-}}\NormalTok{ lambda0}
  \ControlFlowTok{for}\NormalTok{ (i }\ControlFlowTok{in} \DecValTok{1}\SpecialCharTok{:}\NormalTok{r) \{}
\NormalTok{    e }\OtherTok{\textless{}{-}} \FunctionTok{exp}\NormalTok{(}\SpecialCharTok{{-}}\NormalTok{lambda[i])}
\NormalTok{    lambda[i }\SpecialCharTok{+} \DecValTok{1}\NormalTok{] }\OtherTok{\textless{}{-}}\NormalTok{ lambda[i] }\SpecialCharTok{{-}}
\NormalTok{      (mean.n }\SpecialCharTok{/}\NormalTok{ lambda[i] }\SpecialCharTok{{-}} \DecValTok{1} \SpecialCharTok{/}\NormalTok{ (}\DecValTok{1} \SpecialCharTok{{-}}\NormalTok{ e)) }\SpecialCharTok{/}\NormalTok{ ((e }\SpecialCharTok{/}\NormalTok{ (}\DecValTok{1} \SpecialCharTok{{-}}\NormalTok{ e) }\SpecialCharTok{{-}} \DecValTok{1} \SpecialCharTok{/}\NormalTok{ lambda[i]) }\SpecialCharTok{/}\NormalTok{ (}\DecValTok{1} \SpecialCharTok{{-}}\NormalTok{ e))}
\NormalTok{  \}}
  \FunctionTok{data.frame}\NormalTok{(}\AttributeTok{lambda =}\NormalTok{ lambda, }\AttributeTok{log.lik =} \FunctionTok{nzp.log.likelihood}\NormalTok{(n, lambda))}
\NormalTok{\}}
\end{Highlighting}
\end{Shaded}

\subsection{Comparison From Several Starting
Values}\label{comparison-from-several-starting-values}

The three methods are run for 15 iterations from a sensible start, a
start far below, and a start far above.

\begin{Shaded}
\begin{Highlighting}[]
\NormalTok{starts }\OtherTok{\textless{}{-}} \FunctionTok{c}\NormalTok{(}\FunctionTok{mean}\NormalTok{(xp), }\FloatTok{0.1}\NormalTok{, }\DecValTok{100}\NormalTok{)}

\NormalTok{trace\_tab }\OtherTok{\textless{}{-}} \ControlFlowTok{function}\NormalTok{(fun, }\AttributeTok{r =} \DecValTok{15}\NormalTok{)}
  \FunctionTok{vapply}\NormalTok{(starts, }\ControlFlowTok{function}\NormalTok{(s) }\FunctionTok{fun}\NormalTok{(xp, r, s)}\SpecialCharTok{$}\NormalTok{lambda, }\FunctionTok{numeric}\NormalTok{(r }\SpecialCharTok{+} \DecValTok{1}\NormalTok{))}

\NormalTok{paths }\OtherTok{\textless{}{-}} \FunctionTok{list}\NormalTok{(}\AttributeTok{simple =} \FunctionTok{trace\_tab}\NormalTok{(nzp.simple.iteration),}
              \AttributeTok{newton =} \FunctionTok{trace\_tab}\NormalTok{(nzp.newton.raphson),}
              \AttributeTok{scoring =} \FunctionTok{trace\_tab}\NormalTok{(nzp.method.of.scoring))}
\end{Highlighting}
\end{Shaded}

\begin{verbatim}
Warning in log(lambda): NaNs produced
\end{verbatim}

\begin{verbatim}
Warning in log(1 - exp(-lambda)): NaNs produced
\end{verbatim}

\begin{Shaded}
\begin{Highlighting}[]
\DocumentationTok{\#\#\# first six iterates from each starting value}
\ControlFlowTok{for}\NormalTok{ (m }\ControlFlowTok{in} \FunctionTok{names}\NormalTok{(paths)) \{}
  \FunctionTok{cat}\NormalTok{(}\StringTok{"}\SpecialCharTok{\textbackslash{}n}\StringTok{"}\NormalTok{, m, }\StringTok{"}\SpecialCharTok{\textbackslash{}n}\StringTok{"}\NormalTok{)}
\NormalTok{  tab }\OtherTok{\textless{}{-}}\NormalTok{ paths[[m]][}\DecValTok{1}\SpecialCharTok{:}\DecValTok{6}\NormalTok{, ]}
  \FunctionTok{dimnames}\NormalTok{(tab) }\OtherTok{\textless{}{-}} \FunctionTok{list}\NormalTok{(}\FunctionTok{paste}\NormalTok{(}\StringTok{"iter"}\NormalTok{, }\DecValTok{0}\SpecialCharTok{:}\DecValTok{5}\NormalTok{), }\FunctionTok{sprintf}\NormalTok{(}\StringTok{"start = \%.4g"}\NormalTok{, starts))}
  \FunctionTok{print}\NormalTok{(}\FunctionTok{signif}\NormalTok{(tab, }\DecValTok{6}\NormalTok{))}
\NormalTok{\}}
\end{Highlighting}
\end{Shaded}

\begin{verbatim}

 simple 
       start = 2.084 start = 0.1 start = 100
iter 0       2.08434    0.100000   100.00000
iter 1       1.82507    0.198351     2.08434
iter 2       1.74833    0.375010     1.82507
iter 3       1.72153    0.651809     1.74833
iter 4       1.71167    0.998184     1.72153
iter 5       1.70798    1.316160     1.71167

 newton 
       start = 2.084 start = 0.1 start = 100
iter 0       2.08434    0.100000      100.00
iter 1       1.63694    0.195239    -4597.69
iter 2       1.70334    0.371813         NaN
iter 3       1.70574    0.672763         NaN
iter 4       1.70574    1.099100         NaN
iter 5       1.70574    1.505560         NaN

 scoring 
       start = 2.084 start = 0.1 start = 100
iter 0       2.08434     0.10000   100.00000
iter 1       1.71600     2.10035     2.08434
iter 2       1.70575     1.71682     1.71600
iter 3       1.70574     1.70575     1.70575
iter 4       1.70574     1.70574     1.70574
iter 5       1.70574     1.70574     1.70574
\end{verbatim}

The pattern is the one the theory predicts. Newton--Raphson takes the
fewest iterations from a good start, and its digits lock in from the
left as the error is squared each step. Simple iteration converges from
every start, including the extreme ones, but slowly and at a steady
rate. The method of scoring is intermediate: always ascending, and
nearly as fast as Newton once close.

\begin{Shaded}
\begin{Highlighting}[]
\NormalTok{lhat }\OtherTok{\textless{}{-}} \FunctionTok{nzp.newton.raphson}\NormalTok{(xp, }\DecValTok{30}\NormalTok{, }\FunctionTok{mean}\NormalTok{(xp))}\SpecialCharTok{$}\NormalTok{lambda[}\DecValTok{31}\NormalTok{]}

\NormalTok{err }\OtherTok{\textless{}{-}} \FunctionTok{sapply}\NormalTok{(paths, }\ControlFlowTok{function}\NormalTok{(p) }\FunctionTok{abs}\NormalTok{(p[, }\DecValTok{1}\NormalTok{] }\SpecialCharTok{{-}}\NormalTok{ lhat))   }\CommentTok{\# from the mean(xp) start}

\FunctionTok{par}\NormalTok{(}\AttributeTok{mar =} \FunctionTok{c}\NormalTok{(}\DecValTok{4}\NormalTok{, }\FloatTok{4.5}\NormalTok{, }\FloatTok{2.5}\NormalTok{, }\DecValTok{1}\NormalTok{))}
\FunctionTok{matplot}\NormalTok{(}\DecValTok{0}\SpecialCharTok{:}\DecValTok{15}\NormalTok{, }\FunctionTok{pmax}\NormalTok{(err, }\FloatTok{1e{-}17}\NormalTok{), }\AttributeTok{log =} \StringTok{"y"}\NormalTok{, }\AttributeTok{type =} \StringTok{"b"}\NormalTok{, }\AttributeTok{pch =} \DecValTok{19}\NormalTok{, }\AttributeTok{lty =} \DecValTok{1}\NormalTok{, }\AttributeTok{col =} \FunctionTok{c}\NormalTok{(}\DecValTok{4}\NormalTok{, }\DecValTok{1}\NormalTok{, }\DecValTok{2}\NormalTok{),}
        \AttributeTok{xlab =} \FunctionTok{TeX}\NormalTok{(r}\StringTok{"(iteration $k$)"}\NormalTok{), }\AttributeTok{ylab =} \FunctionTok{TeX}\NormalTok{(r}\StringTok{"($|\textbackslash{}lambda\_k {-} \textbackslash{}hat\{\textbackslash{}lambda\}|$)"}\NormalTok{),}
        \AttributeTok{main =} \FunctionTok{TeX}\NormalTok{(r}\StringTok{"(error decay from $\textbackslash{}lambda\_0 = }\SpecialCharTok{\textbackslash{}b}\StringTok{ar\{x\}$)"}\NormalTok{))}
\FunctionTok{legend}\NormalTok{(}\StringTok{"bottomleft"}\NormalTok{, }\AttributeTok{bty =} \StringTok{"n"}\NormalTok{, }\AttributeTok{pch =} \DecValTok{19}\NormalTok{, }\AttributeTok{lty =} \DecValTok{1}\NormalTok{, }\AttributeTok{col =} \FunctionTok{c}\NormalTok{(}\DecValTok{4}\NormalTok{, }\DecValTok{1}\NormalTok{, }\DecValTok{2}\NormalTok{),}
       \AttributeTok{legend =} \FunctionTok{c}\NormalTok{(}\StringTok{"simple iteration"}\NormalTok{, }\StringTok{"Newton{-}Raphson"}\NormalTok{, }\StringTok{"method of scoring"}\NormalTok{))}
\end{Highlighting}
\end{Shaded}

\pandocbounded{\includegraphics[keepaspectratio]{chapters/unit05_MLE1_files/figure-pdf/nzp-convergence-plot-1.pdf}}

On a log scale a linear method is a straight line; Newton's curve bends
downward until it reaches the floor set by machine precision.

\begin{Shaded}
\begin{Highlighting}[]
\NormalTok{S\_xp  }\OtherTok{\textless{}{-}} \ControlFlowTok{function}\NormalTok{(l) }\FunctionTok{nzp.scorefunction}\NormalTok{(xp, l)}
\NormalTok{dS\_xp }\OtherTok{\textless{}{-}} \ControlFlowTok{function}\NormalTok{(l) }\FunctionTok{nzp.dscore}\NormalTok{(xp, l)}
\NormalTok{ll\_xp }\OtherTok{\textless{}{-}} \ControlFlowTok{function}\NormalTok{(l) }\FunctionTok{nzp.log.likelihood}\NormalTok{(xp, l)}

\FunctionTok{c}\NormalTok{(}\AttributeTok{newton      =} \FunctionTok{unname}\NormalTok{(}\FunctionTok{newton}\NormalTok{(S\_xp, dS\_xp, }\FunctionTok{mean}\NormalTok{(xp))[}\StringTok{"root"}\NormalTok{]),}
  \AttributeTok{fixed\_point =} \FunctionTok{unname}\NormalTok{(}\FunctionTok{fixed\_point}\NormalTok{(}\ControlFlowTok{function}\NormalTok{(l) }\FunctionTok{mean}\NormalTok{(xp) }\SpecialCharTok{*}\NormalTok{ (}\DecValTok{1} \SpecialCharTok{{-}} \FunctionTok{exp}\NormalTok{(}\SpecialCharTok{{-}}\NormalTok{l)), }\FunctionTok{mean}\NormalTok{(xp))[}\StringTok{"root"}\NormalTok{]),}
  \AttributeTok{bisection   =} \FunctionTok{unname}\NormalTok{(}\FunctionTok{bisection}\NormalTok{(S\_xp, }\FloatTok{0.1}\NormalTok{, }\DecValTok{20}\NormalTok{)[}\StringTok{"root"}\NormalTok{]),}
  \AttributeTok{golden      =} \FunctionTok{golden}\NormalTok{(}\ControlFlowTok{function}\NormalTok{(l) }\SpecialCharTok{{-}}\FunctionTok{ll\_xp}\NormalTok{(l), }\FunctionTok{c}\NormalTok{(}\FloatTok{0.1}\NormalTok{, }\DecValTok{20}\NormalTok{), }\AttributeTok{eps =} \FloatTok{1e{-}8}\NormalTok{),}
  \AttributeTok{uniroot     =} \FunctionTok{uniroot}\NormalTok{(S\_xp, }\FunctionTok{c}\NormalTok{(}\FloatTok{0.1}\NormalTok{, }\DecValTok{20}\NormalTok{), }\AttributeTok{tol =} \FloatTok{1e{-}10}\NormalTok{)}\SpecialCharTok{$}\NormalTok{root,}
  \AttributeTok{optimize    =} \FunctionTok{optimize}\NormalTok{(}\ControlFlowTok{function}\NormalTok{(l) }\SpecialCharTok{{-}}\FunctionTok{ll\_xp}\NormalTok{(l), }\FunctionTok{c}\NormalTok{(}\FloatTok{0.1}\NormalTok{, }\DecValTok{20}\NormalTok{), }\AttributeTok{tol =} \FloatTok{1e{-}10}\NormalTok{)}\SpecialCharTok{$}\NormalTok{minimum)}
\end{Highlighting}
\end{Shaded}

\begin{verbatim}
     newton fixed_point   bisection      golden     uniroot    optimize 
   1.705744    1.705744    1.705744    1.705744    1.705744    1.705744 
\end{verbatim}

All six agree, which is the point of running them together: the methods
differ in cost and robustness, not in the answer.

\subsection{Standard Error From the
Curvature}\label{standard-error-from-the-curvature}

The optimizer has already produced everything inference needs.

\begin{Shaded}
\begin{Highlighting}[]
\NormalTok{se\_obs }\OtherTok{\textless{}{-}} \DecValTok{1} \SpecialCharTok{/} \FunctionTok{sqrt}\NormalTok{(}\SpecialCharTok{{-}}\FunctionTok{nzp.dscore}\NormalTok{(xp, lhat))}
\NormalTok{se\_exp }\OtherTok{\textless{}{-}} \DecValTok{1} \SpecialCharTok{/} \FunctionTok{sqrt}\NormalTok{(}\FunctionTok{nzp.information}\NormalTok{(xp, lhat))}

\FunctionTok{c}\NormalTok{(}\AttributeTok{mle =}\NormalTok{ lhat, }\AttributeTok{se\_observed =}\NormalTok{ se\_obs, }\AttributeTok{se\_expected =}\NormalTok{ se\_exp)}
\end{Highlighting}
\end{Shaded}

\begin{verbatim}
        mle se_observed se_expected 
   1.705744    0.164514    0.164514 
\end{verbatim}

\begin{Shaded}
\begin{Highlighting}[]
\NormalTok{lhat }\SpecialCharTok{+} \FunctionTok{c}\NormalTok{(}\SpecialCharTok{{-}}\FloatTok{1.96}\NormalTok{, }\FloatTok{1.96}\NormalTok{) }\SpecialCharTok{*}\NormalTok{ se\_obs                     }\CommentTok{\# 95\% Wald interval}
\end{Highlighting}
\end{Shaded}

\begin{verbatim}
[1] 1.383297 2.028192
\end{verbatim}

\begin{Shaded}
\begin{Highlighting}[]
\NormalTok{lam }\OtherTok{\textless{}{-}} \FunctionTok{seq}\NormalTok{(lhat }\SpecialCharTok{{-}} \DecValTok{4} \SpecialCharTok{*}\NormalTok{ se\_obs, lhat }\SpecialCharTok{+} \DecValTok{4} \SpecialCharTok{*}\NormalTok{ se\_obs, }\AttributeTok{length.out =} \DecValTok{300}\NormalTok{)}
\NormalTok{quad }\OtherTok{\textless{}{-}} \SpecialCharTok{{-}}\FloatTok{0.5} \SpecialCharTok{*}\NormalTok{ (}\SpecialCharTok{{-}}\FunctionTok{nzp.dscore}\NormalTok{(xp, lhat)) }\SpecialCharTok{*}\NormalTok{ (lam }\SpecialCharTok{{-}}\NormalTok{ lhat)}\SpecialCharTok{\^{}}\DecValTok{2}

\FunctionTok{par}\NormalTok{(}\AttributeTok{mar =} \FunctionTok{c}\NormalTok{(}\DecValTok{4}\NormalTok{, }\FloatTok{4.5}\NormalTok{, }\FloatTok{2.5}\NormalTok{, }\DecValTok{1}\NormalTok{))}
\FunctionTok{plot}\NormalTok{(lam, }\FunctionTok{ll\_xp}\NormalTok{(lam) }\SpecialCharTok{{-}} \FunctionTok{ll\_xp}\NormalTok{(lhat), }\AttributeTok{type =} \StringTok{"l"}\NormalTok{, }\AttributeTok{lwd =} \DecValTok{2}\NormalTok{,}
     \AttributeTok{xlab =} \FunctionTok{TeX}\NormalTok{(r}\StringTok{"($\textbackslash{}lambda$)"}\NormalTok{), }\AttributeTok{ylab =} \FunctionTok{TeX}\NormalTok{(r}\StringTok{"($}\SpecialCharTok{\textbackslash{}e}\StringTok{ll(\textbackslash{}lambda) {-} }\SpecialCharTok{\textbackslash{}e}\StringTok{ll(\textbackslash{}hat\{\textbackslash{}lambda\})$)"}\NormalTok{),}
     \AttributeTok{main =} \FunctionTok{TeX}\NormalTok{(r}\StringTok{"(log{-}likelihood and its quadratic approximation at $\textbackslash{}hat\{\textbackslash{}lambda\}$)"}\NormalTok{))}
\FunctionTok{lines}\NormalTok{(lam, quad, }\AttributeTok{col =} \DecValTok{2}\NormalTok{, }\AttributeTok{lwd =} \DecValTok{2}\NormalTok{, }\AttributeTok{lty =} \DecValTok{2}\NormalTok{)}
\FunctionTok{abline}\NormalTok{(}\AttributeTok{h =} \SpecialCharTok{{-}}\FloatTok{1.92}\NormalTok{, }\AttributeTok{col =} \DecValTok{4}\NormalTok{, }\AttributeTok{lty =} \DecValTok{3}\NormalTok{)}
\FunctionTok{abline}\NormalTok{(}\AttributeTok{v =}\NormalTok{ lhat }\SpecialCharTok{+} \FunctionTok{c}\NormalTok{(}\SpecialCharTok{{-}}\FloatTok{1.96}\NormalTok{, }\FloatTok{1.96}\NormalTok{) }\SpecialCharTok{*}\NormalTok{ se\_obs, }\AttributeTok{col =} \DecValTok{4}\NormalTok{, }\AttributeTok{lty =} \DecValTok{3}\NormalTok{)}
\FunctionTok{legend}\NormalTok{(}\StringTok{"bottom"}\NormalTok{, }\AttributeTok{bty =} \StringTok{"n"}\NormalTok{, }\AttributeTok{lwd =} \DecValTok{2}\NormalTok{, }\AttributeTok{lty =} \FunctionTok{c}\NormalTok{(}\DecValTok{1}\NormalTok{, }\DecValTok{2}\NormalTok{), }\AttributeTok{col =} \FunctionTok{c}\NormalTok{(}\DecValTok{1}\NormalTok{, }\DecValTok{2}\NormalTok{),}
       \AttributeTok{legend =} \FunctionTok{TeX}\NormalTok{(}\FunctionTok{c}\NormalTok{(r}\StringTok{"($}\SpecialCharTok{\textbackslash{}e}\StringTok{ll(\textbackslash{}lambda)$)"}\NormalTok{, r}\StringTok{"(quadratic from $J(\textbackslash{}hat\{\textbackslash{}lambda\})$)"}\NormalTok{)))}
\end{Highlighting}
\end{Shaded}

\pandocbounded{\includegraphics[keepaspectratio]{chapters/unit05_MLE1_files/figure-pdf/nzp-quadratic-approx-1.pdf}}

The parabola implied by the observed information tracks the
log-likelihood closely over the width of the interval, so the Wald
interval is trustworthy here.

\subsection{Evaluating the Naive
Estimator}\label{evaluating-the-naive-estimator}

All of this machinery is worth the effort only if the MLE beats the easy
alternative. The naive estimator ignores the truncation and reports
\(\bar x\). Simulating many truncated samples gives the sampling
distributions of both estimators.

\begin{Shaded}
\begin{Highlighting}[]
\NormalTok{mle\_nzp }\OtherTok{\textless{}{-}} \ControlFlowTok{function}\NormalTok{(n)}
  \FunctionTok{uniroot}\NormalTok{(}\ControlFlowTok{function}\NormalTok{(l) }\FunctionTok{nzp.scorefunction}\NormalTok{(n, l), }\FunctionTok{c}\NormalTok{(}\FloatTok{1e{-}6}\NormalTok{, }\DecValTok{50}\NormalTok{), }\AttributeTok{tol =} \FloatTok{1e{-}10}\NormalTok{)}\SpecialCharTok{$}\NormalTok{root}

\NormalTok{sim\_nzp }\OtherTok{\textless{}{-}} \ControlFlowTok{function}\NormalTok{(lambda, }\AttributeTok{n\_gen =} \DecValTok{100}\NormalTok{, }\AttributeTok{B =} \DecValTok{500}\NormalTok{)}
  \FunctionTok{vapply}\NormalTok{(}\DecValTok{1}\SpecialCharTok{:}\NormalTok{B, }\ControlFlowTok{function}\NormalTok{(b) \{}
\NormalTok{    y }\OtherTok{\textless{}{-}} \FunctionTok{rpois}\NormalTok{(n\_gen, lambda); n }\OtherTok{\textless{}{-}}\NormalTok{ y[y }\SpecialCharTok{\textgreater{}} \DecValTok{0}\NormalTok{]}
    \FunctionTok{c}\NormalTok{(}\AttributeTok{naive =} \FunctionTok{mean}\NormalTok{(n), }\AttributeTok{mle =} \FunctionTok{mle\_nzp}\NormalTok{(n))}
\NormalTok{  \}, }\FunctionTok{c}\NormalTok{(}\AttributeTok{naive =} \DecValTok{0}\NormalTok{, }\AttributeTok{mle =} \DecValTok{0}\NormalTok{))}

\ControlFlowTok{for}\NormalTok{ (lambda\_true }\ControlFlowTok{in} \FunctionTok{c}\NormalTok{(}\FloatTok{0.5}\NormalTok{, }\FloatTok{1.5}\NormalTok{)) \{}
\NormalTok{  est }\OtherTok{\textless{}{-}} \FunctionTok{sim\_nzp}\NormalTok{(lambda\_true)}
\NormalTok{  tab }\OtherTok{\textless{}{-}} \FunctionTok{t}\NormalTok{(}\FunctionTok{apply}\NormalTok{(est, }\DecValTok{1}\NormalTok{, }\ControlFlowTok{function}\NormalTok{(e)}
    \FunctionTok{c}\NormalTok{(}\AttributeTok{bias =} \FunctionTok{mean}\NormalTok{(e) }\SpecialCharTok{{-}}\NormalTok{ lambda\_true, }\AttributeTok{sd =} \FunctionTok{sd}\NormalTok{(e), }\AttributeTok{rmse =} \FunctionTok{sqrt}\NormalTok{(}\FunctionTok{mean}\NormalTok{((e }\SpecialCharTok{{-}}\NormalTok{ lambda\_true)}\SpecialCharTok{\^{}}\DecValTok{2}\NormalTok{)))))}
  \FunctionTok{cat}\NormalTok{(}\StringTok{"}\SpecialCharTok{\textbackslash{}n}\StringTok{lambda ="}\NormalTok{, lambda\_true, }\StringTok{"}\SpecialCharTok{\textbackslash{}n}\StringTok{"}\NormalTok{); }\FunctionTok{print}\NormalTok{(}\FunctionTok{round}\NormalTok{(tab, }\DecValTok{4}\NormalTok{))}
\NormalTok{\}}
\end{Highlighting}
\end{Shaded}

\begin{verbatim}

lambda = 0.5 
         bias     sd   rmse
naive  0.7666 0.0845 0.7712
mle   -0.0101 0.1448 0.1450

lambda = 1.5 
         bias     sd   rmse
naive  0.4201 0.1162 0.4359
mle   -0.0169 0.1590 0.1597
\end{verbatim}

\begin{Shaded}
\begin{Highlighting}[]
\NormalTok{est }\OtherTok{\textless{}{-}} \FunctionTok{sim\_nzp}\NormalTok{(}\FloatTok{1.5}\NormalTok{)}
\FunctionTok{par}\NormalTok{(}\AttributeTok{mfrow =} \FunctionTok{c}\NormalTok{(}\DecValTok{1}\NormalTok{, }\DecValTok{2}\NormalTok{), }\AttributeTok{mar =} \FunctionTok{c}\NormalTok{(}\DecValTok{4}\NormalTok{, }\FloatTok{4.5}\NormalTok{, }\FloatTok{2.5}\NormalTok{, }\DecValTok{1}\NormalTok{))}
\ControlFlowTok{for}\NormalTok{ (m }\ControlFlowTok{in} \FunctionTok{c}\NormalTok{(}\StringTok{"naive"}\NormalTok{, }\StringTok{"mle"}\NormalTok{)) \{}
  \FunctionTok{hist}\NormalTok{(est[m, ], }\AttributeTok{breaks =} \DecValTok{25}\NormalTok{, }\AttributeTok{col =} \StringTok{"grey90"}\NormalTok{,}
       \AttributeTok{xlim =} \FunctionTok{range}\NormalTok{(est), }\AttributeTok{xlab =} \FunctionTok{TeX}\NormalTok{(r}\StringTok{"($\textbackslash{}hat\{\textbackslash{}lambda\}$)"}\NormalTok{),}
       \AttributeTok{main =} \FunctionTok{TeX}\NormalTok{(}\FunctionTok{sprintf}\NormalTok{(r}\StringTok{"(\%s, bias $= \%.3f$)"}\NormalTok{, m, }\FunctionTok{mean}\NormalTok{(est[m, ]) }\SpecialCharTok{{-}} \FloatTok{1.5}\NormalTok{)))}
  \FunctionTok{abline}\NormalTok{(}\AttributeTok{v =} \FloatTok{1.5}\NormalTok{, }\AttributeTok{col =} \DecValTok{2}\NormalTok{, }\AttributeTok{lwd =} \DecValTok{2}\NormalTok{, }\AttributeTok{lty =} \DecValTok{2}\NormalTok{)}
\NormalTok{\}}
\end{Highlighting}
\end{Shaded}

\pandocbounded{\includegraphics[keepaspectratio]{chapters/unit05_MLE1_files/figure-pdf/nzp-mle-vs-naive-1.pdf}}

The naive estimator is biased upward, and severely so when \(\lambda\)
is small, since that is when most of the probability mass sits on the
discarded zero. Its bias does not shrink with sample size: it reflects
the wrong model, not too little data. The MLE is centred on the truth
with a standard deviation matching the information-based standard error
of the previous section. The optimization was the means; the correct
estimator was the end.

\section{Postscript: Interval Data and a Flat
Likelihood}\label{postscript-interval-data-and-a-flat-likelihood}

Not every likelihood is well approximated by a parabola, and when it is
not, the information-based standard error misleads. Suppose each Poisson
observation is known only to lie in an interval \([l_i, h_i]\). Then
\[\ell(\lambda) = \sum_{i=1}^{n} \log \Pr\left(l_i \le X_i \le h_i \mid \lambda\right)
= \sum_{i=1}^{n} \log \sum_{x = l_i}^{h_i} \frac{e^{-\lambda}\lambda^{x}}{x!} .\]

The inner sum is computed on the log scale with the log-sum-exp trick to
avoid underflow, and the maximization is delegated to \texttt{nlm},
which returns a numerical Hessian and hence a standard error.

\begin{Shaded}
\begin{Highlighting}[]
\NormalTok{log\_sum\_exp }\OtherTok{\textless{}{-}} \ControlFlowTok{function}\NormalTok{(lx) \{ mlx }\OtherTok{\textless{}{-}} \FunctionTok{max}\NormalTok{(lx); mlx }\SpecialCharTok{+} \FunctionTok{log}\NormalTok{(}\FunctionTok{sum}\NormalTok{(}\FunctionTok{exp}\NormalTok{(lx }\SpecialCharTok{{-}}\NormalTok{ mlx))) \}}

\NormalTok{poisson\_log\_likelihood }\OtherTok{\textless{}{-}} \ControlFlowTok{function}\NormalTok{(lambda, low, high)}
  \FunctionTok{sum}\NormalTok{(}\FunctionTok{vapply}\NormalTok{(}\FunctionTok{seq\_along}\NormalTok{(low),}
             \ControlFlowTok{function}\NormalTok{(i) }\FunctionTok{log\_sum\_exp}\NormalTok{(}\FunctionTok{dpois}\NormalTok{(low[i]}\SpecialCharTok{:}\NormalTok{high[i], lambda, }\AttributeTok{log =} \ConstantTok{TRUE}\NormalTok{)),}
             \FunctionTok{numeric}\NormalTok{(}\DecValTok{1}\NormalTok{)))}
\DocumentationTok{\#\#\# equivalently, and faster: sum(log(ppois(high, lambda) {-} ppois(low {-} 1, lambda)))}

\NormalTok{poisson\_mle }\OtherTok{\textless{}{-}} \ControlFlowTok{function}\NormalTok{(low, high) \{}
  \ControlFlowTok{if}\NormalTok{ (}\FunctionTok{length}\NormalTok{(low) }\SpecialCharTok{!=} \FunctionTok{length}\NormalTok{(high)) }\FunctionTok{stop}\NormalTok{(}\StringTok{"low and high have different lengths"}\NormalTok{)}
  \ControlFlowTok{if}\NormalTok{ (}\FunctionTok{any}\NormalTok{(}\FunctionTok{floor}\NormalTok{(low) }\SpecialCharTok{!=}\NormalTok{ low) }\SpecialCharTok{||} \FunctionTok{any}\NormalTok{(}\FunctionTok{floor}\NormalTok{(high) }\SpecialCharTok{!=}\NormalTok{ high)) }\FunctionTok{stop}\NormalTok{(}\StringTok{"interval ends are not integers"}\NormalTok{)}

\NormalTok{  fit }\OtherTok{\textless{}{-}} \FunctionTok{nlm}\NormalTok{(}\ControlFlowTok{function}\NormalTok{(l) }\SpecialCharTok{{-}}\FunctionTok{poisson\_log\_likelihood}\NormalTok{(l, low, high),}
             \AttributeTok{p =}\NormalTok{ (}\FunctionTok{mean}\NormalTok{(low) }\SpecialCharTok{+} \FunctionTok{mean}\NormalTok{(high)) }\SpecialCharTok{/} \DecValTok{2}\NormalTok{, }\AttributeTok{hessian =} \ConstantTok{TRUE}\NormalTok{)}
\NormalTok{  est }\OtherTok{\textless{}{-}}\NormalTok{ fit}\SpecialCharTok{$}\NormalTok{estimate; sd }\OtherTok{\textless{}{-}} \FunctionTok{sqrt}\NormalTok{(}\DecValTok{1} \SpecialCharTok{/}\NormalTok{ fit}\SpecialCharTok{$}\NormalTok{hessian[}\DecValTok{1}\NormalTok{, }\DecValTok{1}\NormalTok{])}
  \FunctionTok{list}\NormalTok{(}\AttributeTok{est =}\NormalTok{ est, }\AttributeTok{sd =}\NormalTok{ sd, }\AttributeTok{CI =}\NormalTok{ est }\SpecialCharTok{+} \FunctionTok{c}\NormalTok{(}\SpecialCharTok{{-}}\FloatTok{1.96}\NormalTok{, }\FloatTok{1.96}\NormalTok{) }\SpecialCharTok{*}\NormalTok{ sd)}
\NormalTok{\}}
\end{Highlighting}
\end{Shaded}

\begin{Shaded}
\begin{Highlighting}[]
\NormalTok{pn }\OtherTok{\textless{}{-}} \FunctionTok{rpois}\NormalTok{(}\DecValTok{1000}\NormalTok{, }\AttributeTok{lambda =} \DecValTok{20}\NormalTok{)}

\NormalTok{tight }\OtherTok{\textless{}{-}} \FunctionTok{poisson\_mle}\NormalTok{(}\FunctionTok{pmax}\NormalTok{(pn }\SpecialCharTok{{-}} \DecValTok{1}\NormalTok{, }\DecValTok{0}\NormalTok{), pn }\SpecialCharTok{+} \DecValTok{4}\NormalTok{)      }\CommentTok{\# narrow intervals: much information}
\NormalTok{wide  }\OtherTok{\textless{}{-}} \FunctionTok{poisson\_mle}\NormalTok{(}\FunctionTok{pmax}\NormalTok{(pn }\SpecialCharTok{{-}} \DecValTok{5}\NormalTok{, }\DecValTok{0}\NormalTok{), pn }\SpecialCharTok{+} \DecValTok{5}\NormalTok{)     }\CommentTok{\# wide intervals: little information}

\FunctionTok{rbind}\NormalTok{(}\AttributeTok{tight =} \FunctionTok{c}\NormalTok{(}\AttributeTok{est =}\NormalTok{ tight}\SpecialCharTok{$}\NormalTok{est, }\AttributeTok{sd =}\NormalTok{ tight}\SpecialCharTok{$}\NormalTok{sd, tight}\SpecialCharTok{$}\NormalTok{CI),}
      \AttributeTok{wide  =} \FunctionTok{c}\NormalTok{(}\AttributeTok{est =}\NormalTok{ wide}\SpecialCharTok{$}\NormalTok{est,  }\AttributeTok{sd =}\NormalTok{ wide}\SpecialCharTok{$}\NormalTok{sd,  wide}\SpecialCharTok{$}\NormalTok{CI))}
\end{Highlighting}
\end{Shaded}

\begin{verbatim}
           est        sd                  
tight 21.53996 0.1565054 21.23321 21.84671
wide  20.00541 0.1752922 19.66183 20.34898
\end{verbatim}

\begin{Shaded}
\begin{Highlighting}[]
\NormalTok{lam\_grid }\OtherTok{\textless{}{-}} \FunctionTok{seq}\NormalTok{(}\DecValTok{1}\NormalTok{, }\DecValTok{50}\NormalTok{, }\AttributeTok{by =} \FloatTok{0.5}\NormalTok{)}

\FunctionTok{par}\NormalTok{(}\AttributeTok{mfrow =} \FunctionTok{c}\NormalTok{(}\DecValTok{1}\NormalTok{, }\DecValTok{2}\NormalTok{), }\AttributeTok{mar =} \FunctionTok{c}\NormalTok{(}\DecValTok{4}\NormalTok{, }\FloatTok{4.5}\NormalTok{, }\FloatTok{2.5}\NormalTok{, }\DecValTok{1}\NormalTok{))}
\ControlFlowTok{for}\NormalTok{ (iv }\ControlFlowTok{in} \FunctionTok{list}\NormalTok{(}\FunctionTok{list}\NormalTok{(}\AttributeTok{lo =} \FunctionTok{pmax}\NormalTok{(pn }\SpecialCharTok{{-}} \DecValTok{1}\NormalTok{, }\DecValTok{0}\NormalTok{), }\AttributeTok{hi =}\NormalTok{ pn }\SpecialCharTok{+} \DecValTok{4}\NormalTok{, }\AttributeTok{lab =} \StringTok{"narrow intervals"}\NormalTok{),}
                \FunctionTok{list}\NormalTok{(}\AttributeTok{lo =} \FunctionTok{pmax}\NormalTok{(pn }\SpecialCharTok{{-}} \DecValTok{50}\NormalTok{, }\DecValTok{0}\NormalTok{), }\AttributeTok{hi =}\NormalTok{ pn }\SpecialCharTok{+} \DecValTok{5}\NormalTok{, }\AttributeTok{lab =} \StringTok{"wide intervals"}\NormalTok{))) \{}
\NormalTok{  lp }\OtherTok{\textless{}{-}} \FunctionTok{vapply}\NormalTok{(lam\_grid, poisson\_log\_likelihood, }\FunctionTok{numeric}\NormalTok{(}\DecValTok{1}\NormalTok{), }\AttributeTok{low =}\NormalTok{ iv}\SpecialCharTok{$}\NormalTok{lo, }\AttributeTok{high =}\NormalTok{ iv}\SpecialCharTok{$}\NormalTok{hi)}
  \FunctionTok{plot}\NormalTok{(lam\_grid, lp }\SpecialCharTok{{-}} \FunctionTok{max}\NormalTok{(lp), }\AttributeTok{type =} \StringTok{"l"}\NormalTok{, }\AttributeTok{lwd =} \DecValTok{2}\NormalTok{, }\AttributeTok{ylim =} \FunctionTok{c}\NormalTok{(}\SpecialCharTok{{-}}\DecValTok{50}\NormalTok{, }\DecValTok{0}\NormalTok{),}
       \AttributeTok{xlab =} \FunctionTok{TeX}\NormalTok{(r}\StringTok{"($\textbackslash{}lambda$)"}\NormalTok{), }\AttributeTok{ylab =} \FunctionTok{TeX}\NormalTok{(r}\StringTok{"($}\SpecialCharTok{\textbackslash{}e}\StringTok{ll(\textbackslash{}lambda) {-} }\SpecialCharTok{\textbackslash{}e}\StringTok{ll(\textbackslash{}hat\{\textbackslash{}lambda\})$)"}\NormalTok{),}
       \AttributeTok{main =} \FunctionTok{TeX}\NormalTok{(}\FunctionTok{sprintf}\NormalTok{(r}\StringTok{"(\%s)"}\NormalTok{, iv}\SpecialCharTok{$}\NormalTok{lab)))}
  \FunctionTok{abline}\NormalTok{(}\AttributeTok{h =} \SpecialCharTok{{-}}\FloatTok{1.92}\NormalTok{, }\AttributeTok{col =} \DecValTok{4}\NormalTok{, }\AttributeTok{lty =} \DecValTok{3}\NormalTok{)}
\NormalTok{\}}
\end{Highlighting}
\end{Shaded}

\pandocbounded{\includegraphics[keepaspectratio]{chapters/unit05_MLE1_files/figure-pdf/interval-poisson-loglik-1.pdf}}

With narrow intervals the log-likelihood is sharply peaked and the
standard error is small and reliable. With wide intervals, most of which
include the lower boundary at zero, the log-likelihood is nearly flat
over a broad range: the curvature at the maximum is tiny, the reported
standard error is large, and --- because the curve is visibly asymmetric
and non-quadratic --- even that large value understates the true
uncertainty. The information-based standard error is an asymptotic
approximation resting on local quadratic behaviour, and this is what its
failure looks like. When the log-likelihood is this shape, report a
likelihood-based interval, read off where \(\ell\) drops by 1.92,
instead of a Wald interval.

\section{\texorpdfstring{A Practical Newton-Raphson Algorithm for
Finding
\(\sqrt{a}\)}{A Practical Newton-Raphson Algorithm for Finding \textbackslash sqrt\{a\}}}\label{a-practical-newton-raphson-algorithm-for-finding-sqrta}

\textbf{Derivation.} To find \(\sqrt{a}\), take \(f(x) = x^2 - a\),
whose positive root is \(\sqrt a\). Then \(f'(x) = 2x\), and the Newton
step is

\[x_{k+1} = x_k - \frac{f(x_k)}{f'(x_k)} = x_k - \frac{x_k^{2} - a}{2x_k} = \frac{1}{2}\left(x_k + \frac{a}{x_k}\right).\]

The final form is the Babylonian rule: if \(x_k\) is too small then
\(a/x_k\) is too large, and the root lies between them, so average the
pair. Newton's tangent line lands exactly on the midpoint.

The error follows the general quadratic law. Writing
\(e_k = x_k - \sqrt a\),

\[e_{k+1} = \frac{1}{2}\left(x_k + \frac{a}{x_k}\right) - \sqrt a = \frac{\left(x_k - \sqrt a\right)^{2}}{2x_k} = \frac{e_k^{2}}{2x_k},\]

so the number of correct digits roughly doubles each step.

\begin{Shaded}
\begin{Highlighting}[]
\NormalTok{babylonian }\OtherTok{\textless{}{-}} \ControlFlowTok{function}\NormalTok{(a, }\AttributeTok{x0 =} \DecValTok{1}\NormalTok{, }\AttributeTok{iters =} \DecValTok{6}\NormalTok{)}
  \FunctionTok{Reduce}\NormalTok{(}\ControlFlowTok{function}\NormalTok{(x, k) (x }\SpecialCharTok{+}\NormalTok{ a }\SpecialCharTok{/}\NormalTok{ x) }\SpecialCharTok{/} \DecValTok{2}\NormalTok{, }\DecValTok{1}\SpecialCharTok{:}\NormalTok{iters, }\AttributeTok{init =}\NormalTok{ x0, }\AttributeTok{accumulate =} \ConstantTok{TRUE}\NormalTok{)}

\NormalTok{traces }\OtherTok{\textless{}{-}} \FunctionTok{sapply}\NormalTok{(}\FunctionTok{c}\NormalTok{(}\DecValTok{10}\NormalTok{, }\DecValTok{20}\NormalTok{, }\DecValTok{50}\NormalTok{), babylonian)}
\FunctionTok{dimnames}\NormalTok{(traces) }\OtherTok{\textless{}{-}} \FunctionTok{list}\NormalTok{(}\FunctionTok{paste}\NormalTok{(}\StringTok{"iter"}\NormalTok{, }\DecValTok{0}\SpecialCharTok{:}\DecValTok{6}\NormalTok{), }\FunctionTok{sprintf}\NormalTok{(}\StringTok{"sqrt(\%d)"}\NormalTok{, }\FunctionTok{c}\NormalTok{(}\DecValTok{10}\NormalTok{, }\DecValTok{20}\NormalTok{, }\DecValTok{50}\NormalTok{)))}

\FunctionTok{print}\NormalTok{(traces, }\AttributeTok{digits =} \DecValTok{12}\NormalTok{)}
\end{Highlighting}
\end{Shaded}

\begin{verbatim}
            sqrt(10)       sqrt(20)       sqrt(50)
iter 0 1.00000000000  1.00000000000  1.00000000000
iter 1 5.50000000000 10.50000000000 25.50000000000
iter 2 3.65909090909  6.20238095238 13.73039215686
iter 3 3.19600508187  4.71347454529  8.68597437190
iter 4 3.16245562280  4.47831444547  7.22119047433
iter 5 3.16227766518  4.47214021707  7.07262827574
iter 6 3.16227766017  4.47213595500  7.07106798401
\end{verbatim}

\begin{Shaded}
\begin{Highlighting}[]
\FunctionTok{sqrt}\NormalTok{(}\FunctionTok{c}\NormalTok{(}\DecValTok{10}\NormalTok{, }\DecValTok{20}\NormalTok{, }\DecValTok{50}\NormalTok{))                       }\CommentTok{\# for comparison}
\end{Highlighting}
\end{Shaded}

\begin{verbatim}
[1] 3.162278 4.472136 7.071068
\end{verbatim}

\begin{Shaded}
\begin{Highlighting}[]
\FunctionTok{signif}\NormalTok{(}\FunctionTok{abs}\NormalTok{(traces }\SpecialCharTok{{-}} \FunctionTok{rep}\NormalTok{(}\FunctionTok{sqrt}\NormalTok{(}\FunctionTok{c}\NormalTok{(}\DecValTok{10}\NormalTok{, }\DecValTok{20}\NormalTok{, }\DecValTok{50}\NormalTok{)), }\AttributeTok{each =} \DecValTok{7}\NormalTok{)), }\DecValTok{3}\NormalTok{)    }\CommentTok{\# errors}
\end{Highlighting}
\end{Shaded}

\begin{verbatim}
       sqrt(10) sqrt(20) sqrt(50)
iter 0 2.16e+00 3.47e+00 6.07e+00
iter 1 2.34e+00 6.03e+00 1.84e+01
iter 2 4.97e-01 1.73e+00 6.66e+00
iter 3 3.37e-02 2.41e-01 1.61e+00
iter 4 1.78e-04 6.18e-03 1.50e-01
iter 5 5.01e-09 4.26e-06 1.56e-03
iter 6 4.44e-16 2.03e-12 1.72e-07
\end{verbatim}

Starting from \(x_0 = 1\) in every case, the first two steps do the
coarse work --- the larger \(a\) is, the further the opening leap ---
and from the third step on the error squares: roughly \(10^{-1}\), then
\(10^{-3}\), then \(10^{-6}\), then machine precision. Six iterations
suffice for all three, though \(\sqrt{50}\) spends one more step than
\(\sqrt{10}\) recovering from the poor start.

\bookmarksetup{startatroot}

\chapter{Multivariate Optimization Methods for Finding
MLE}\label{multivariate-optimization-methods-for-finding-mle}

In the previous chapter the parameter was a scalar, and the
log-likelihood could be inspected by plotting it. Once
\(\theta \in \mathbb{R}^p\) neither is true: we cannot see the surface,
and the notion of ``bracketing a root'' that underlies bisection and
golden section has no direct multivariate analogue. What survives is the
local quadratic approximation, and the Newton--Raphson algorithm is its
direct consequence.

This chapter develops the multivariate Newton--Raphson method and its
statistical variant, Fisher scoring, applies both to logistic
regression, examines interactively how and why the iteration can
diverge, and closes with a survey of the derivative-free and
gradient-only alternatives.

\section{The Newton--Raphson Algorithm in Several
Dimensions}\label{the-newtonraphson-algorithm-in-several-dimensions}

\subsection{Derivation from a quadratic
approximation}\label{derivation-from-a-quadratic-approximation}

Write the log-likelihood as \(\ell(\theta)\) for
\(\theta \in \mathbb{R}^p\), and define

\[
S(\theta) = \frac{\partial \ell(\theta)}{\partial \theta}
\quad (p \times 1, \text{ the score}),
\qquad
H(\theta) = \frac{\partial^2 \ell(\theta)}{\partial \theta \, \partial \theta^{\top}}
\quad (p \times p, \text{ the Hessian}).
\]

The negative Hessian \(J(\theta) = -H(\theta)\) is the \textbf{observed
information matrix}.

Expand \(\ell\) to second order about the current iterate
\(\theta^{(k)}\):

\[
\ell(\theta) \;\approx\; \ell(\theta^{(k)})
 + S(\theta^{(k)})^{\top}(\theta - \theta^{(k)})
 - \tfrac{1}{2}\,(\theta - \theta^{(k)})^{\top} J(\theta^{(k)}) \,(\theta - \theta^{(k)}).
\]

This is a quadratic in \(\theta\). If \(J(\theta^{(k)})\) is positive
definite the quadratic has a unique maximum, found by setting its
derivative to zero:

\[
S(\theta^{(k)}) - J(\theta^{(k)})(\theta - \theta^{(k)}) = 0 .
\]

Solving for \(\theta\) gives the \textbf{Newton--Raphson update}:

\[
\boxed{\;\theta^{(k+1)} = \theta^{(k)} + J(\theta^{(k)})^{-1} S(\theta^{(k)}) \;}
\]

The algorithm therefore replaces the log-likelihood by the quadratic
that matches its value, slope and curvature at the current point, and
jumps to the maximum of that quadratic. Everything that follows --- the
fast convergence, and the modes of failure --- comes from how good or
bad that replacement is.

Three remarks are worth making explicit.

\emph{The inverse is never formed.} Write the update as
\(\theta^{(k+1)} = \theta^{(k)} + \delta^{(k)}\) where \(\delta^{(k)}\)
solves the linear system \(J(\theta^{(k)})\,\delta = S(\theta^{(k)})\).
In R this is \texttt{solve(J,\ S)}, not \texttt{solve(J)\ \%*\%\ S}. The
former is faster and numerically more stable; the latter computes
\(p^2\) numbers in order to use \(p\) of them.

\emph{The step is not the gradient direction.}
\(\delta^{(k)} = J^{-1}S\) points along the gradient only when \(J\) is
a multiple of the identity. The information matrix rescales and rotates
the ascent direction to account for differing curvature and correlation
among the parameters, which is exactly why Newton--Raphson is
insensitive to the scaling of the parameters while gradient descent is
not.

\emph{Positive definiteness is an assumption, not a guarantee.} Away
from the maximum \(J(\theta^{(k)})\) can be indefinite or singular, in
which case the ``maximum of the quadratic'' is a saddle point or does
not exist, and \(\delta^{(k)}\) may point downhill.
Section~\ref{sec-divergence} examines this.

\subsubsection{Fisher scoring}\label{fisher-scoring}

The observed information \(J(\theta)\) depends on the data. Its
expectation under the model,

\[
I(\theta) = E_{\theta}\{J(\theta)\} = -E_{\theta}\left\{ \frac{\partial^2 \ell(\theta)}{\partial \theta \, \partial \theta^{\top}} \right\},
\]

is the \textbf{expected (Fisher) information matrix}. Replacing \(J\) by
\(I\) in the update gives the \textbf{method of scoring}:

\[
\theta^{(k+1)} = \theta^{(k)} + I(\theta^{(k)})^{-1} S(\theta^{(k)}).
\]

The two algorithms have different iterates but the same fixed point,
since both stop where \(S(\theta) = 0\). The practical differences:

\begin{longtable}[]{@{}
  >{\raggedright\arraybackslash}p{(\linewidth - 4\tabcolsep) * \real{0.3333}}
  >{\raggedright\arraybackslash}p{(\linewidth - 4\tabcolsep) * \real{0.3333}}
  >{\raggedright\arraybackslash}p{(\linewidth - 4\tabcolsep) * \real{0.3333}}@{}}
\caption{Newton--Raphson compared with Fisher
scoring}\label{tbl-nr-scoring}\tabularnewline
\toprule\noalign{}
\begin{minipage}[b]{\linewidth}\raggedright
\end{minipage} & \begin{minipage}[b]{\linewidth}\raggedright
Newton--Raphson
\end{minipage} & \begin{minipage}[b]{\linewidth}\raggedright
Fisher scoring
\end{minipage} \\
\midrule\noalign{}
\endfirsthead
\toprule\noalign{}
\begin{minipage}[b]{\linewidth}\raggedright
\end{minipage} & \begin{minipage}[b]{\linewidth}\raggedright
Newton--Raphson
\end{minipage} & \begin{minipage}[b]{\linewidth}\raggedright
Fisher scoring
\end{minipage} \\
\midrule\noalign{}
\endhead
\bottomrule\noalign{}
\endlastfoot
Curvature matrix & Observed information \(J(\theta)\) & Expected
information \(I(\theta)\) \\
Derivation & Second derivative of \(\ell\) & Expectation of that second
derivative \\
Positive definite? & Only guaranteed near the maximum & Always, wherever
the model is identified \\
Algebra & Can be messy & Often simpler, as data-dependent terms average
out \\
Local convergence & Quadratic & Linear, but typically with a very small
rate constant \\
Stability from poor starts & Less stable & More stable \\
\end{longtable}

The guaranteed positive definiteness of \(I(\theta)\) is the main
practical argument for scoring: the update direction is always an ascent
direction, so the iteration cannot be sent uphill-seeking in a downhill
direction by an ill-behaved Hessian.

For generalised linear models with the \textbf{canonical link} the two
coincide exactly, because the second derivative of the log-likelihood
does not involve the response. Logistic regression uses the canonical
logit link, so the example below is simultaneously Newton--Raphson,
Fisher scoring, and --- as we note at the end of the section ---
iteratively reweighted least squares.

\subsection{Convergence properties}\label{convergence-properties}

Near a maximum \(\hat\theta\) at which \(J(\hat\theta)\) is positive
definite, Newton--Raphson converges \textbf{quadratically}: the error
satisfies

\[
\|\theta^{(k+1)} - \hat\theta\| \;\le\; C \,\|\theta^{(k)} - \hat\theta\|^{2}
\]

for some constant \(C\) depending on the third derivatives. Informally,
the number of correct digits doubles at each iteration, which is why
five or six iterations are usually enough and why a well-started run
appears to converge instantly.

The qualification ``near'' is essential. The theorem is local: it
guarantees nothing about a start far from \(\hat\theta\), and nothing at
all if \(J\) fails to be positive definite along the way.

\subsection{Standard errors as a
by-product}\label{standard-errors-as-a-by-product}

At convergence the same matrix that drove the iteration supplies the
asymptotic variance. Under the usual regularity conditions,

\[
\hat\theta \;\dot\sim\; N\!\left(\theta, \; J(\hat\theta)^{-1}\right),
\]

so
\(\operatorname{SE}(\hat\theta_j) = \sqrt{[J(\hat\theta)^{-1}]_{jj}}\).
This is a genuine practical advantage of Newton-type methods over
derivative-free ones: the curvature information needed for inference has
already been computed. Methods that never form a Hessian require a
separate numerical differentiation step, which is what
\texttt{hessian\ =\ TRUE} requests from \texttt{nlm} and \texttt{optim}.

\section{Worked Example: Logistic
Regression}\label{worked-example-logistic-regression}

\subsection{The model and its
derivatives}\label{the-model-and-its-derivatives}

With a single covariate \(z_i\) and binary response \(r_i\), let
\(u_i = \beta_0 + \beta_1 z_i\) and \(p_i = e^{u_i}/(1 + e^{u_i})\). The
log-likelihood is

\[
\ell(\beta) = \sum_{i=1}^{n}\left\{ r_i u_i - \log\left(1 + e^{u_i}\right) \right\}.
\]

Writing \(X\) for the \(n \times 2\) design matrix with rows
\((1, z_i)\), the score and information are

\[
S(\beta) = X^{\top}(r - p), \qquad
J(\beta) = X^{\top} W X, \qquad W = \operatorname{diag}\{p_i(1 - p_i)\}.
\]

Note that \(J\) does not involve \(r\). This is the canonical-link
property mentioned above, and it means \(J(\beta) = I(\beta)\)
identically.

Two numerical points, both carried over from the chapter on computer
arithmetic. First, \(\log(1 + e^{u})\) overflows for \(u \gtrsim 710\);
the stable form is \(\log(1+e^u) = \max(u,0) + \log(1 + e^{-|u|})\).
Second, \(p_i(1-p_i) \to 0\) when \(|u_i|\) is large, so \(W\) --- and
hence \(J\) --- degenerates toward singularity when the fitted
probabilities are extreme. That observation is the whole explanation of
the divergence we are about to see.

\subsection{Implementation}\label{implementation}

\begin{Shaded}
\begin{Highlighting}[]
\DocumentationTok{\#\#\# numerically stable log(1 + exp(u))}
\NormalTok{log1pexp }\OtherTok{\textless{}{-}} \ControlFlowTok{function}\NormalTok{(u) }\FunctionTok{pmax}\NormalTok{(u, }\DecValTok{0}\NormalTok{) }\SpecialCharTok{+} \FunctionTok{log1p}\NormalTok{(}\FunctionTok{exp}\NormalTok{(}\SpecialCharTok{{-}}\FunctionTok{abs}\NormalTok{(u)))}

\DocumentationTok{\#\#\# negative log{-}likelihood, negative score, and information matrix}
\NormalTok{lik\_score\_info }\OtherTok{\textless{}{-}} \ControlFlowTok{function}\NormalTok{(b, r, z) \{}
\NormalTok{    u }\OtherTok{\textless{}{-}}\NormalTok{ b[}\DecValTok{1}\NormalTok{] }\SpecialCharTok{+}\NormalTok{ b[}\DecValTok{2}\NormalTok{] }\SpecialCharTok{*}\NormalTok{ z}
\NormalTok{    p }\OtherTok{\textless{}{-}} \FunctionTok{plogis}\NormalTok{(u)}
\NormalTok{    w }\OtherTok{\textless{}{-}}\NormalTok{ p }\SpecialCharTok{*}\NormalTok{ (}\DecValTok{1} \SpecialCharTok{{-}}\NormalTok{ p)}
    \FunctionTok{list}\NormalTok{(}
        \AttributeTok{neg\_loglik =} \SpecialCharTok{{-}}\FunctionTok{sum}\NormalTok{(u }\SpecialCharTok{*}\NormalTok{ r }\SpecialCharTok{{-}} \FunctionTok{log1pexp}\NormalTok{(u)),}
        \AttributeTok{neg\_score  =} \SpecialCharTok{{-}}\FunctionTok{c}\NormalTok{(}\FunctionTok{sum}\NormalTok{(r }\SpecialCharTok{{-}}\NormalTok{ p), }\FunctionTok{sum}\NormalTok{(z }\SpecialCharTok{*}\NormalTok{ (r }\SpecialCharTok{{-}}\NormalTok{ p))),}
        \AttributeTok{inf        =} \FunctionTok{matrix}\NormalTok{(}\FunctionTok{c}\NormalTok{(}\FunctionTok{sum}\NormalTok{(w), }\FunctionTok{sum}\NormalTok{(z }\SpecialCharTok{*}\NormalTok{ w),}
                              \FunctionTok{sum}\NormalTok{(z }\SpecialCharTok{*}\NormalTok{ w), }\FunctionTok{sum}\NormalTok{(z }\SpecialCharTok{*}\NormalTok{ z }\SpecialCharTok{*}\NormalTok{ w)), }\DecValTok{2}\NormalTok{, }\DecValTok{2}\NormalTok{)}
\NormalTok{    )}
\NormalTok{\}}

\NormalTok{neg\_logp\_logistic }\OtherTok{\textless{}{-}} \ControlFlowTok{function}\NormalTok{(b, z, r) \{}
\NormalTok{    u }\OtherTok{\textless{}{-}}\NormalTok{ b[}\DecValTok{1}\NormalTok{] }\SpecialCharTok{+}\NormalTok{ b[}\DecValTok{2}\NormalTok{] }\SpecialCharTok{*}\NormalTok{ z}
    \SpecialCharTok{{-}}\FunctionTok{sum}\NormalTok{(u }\SpecialCharTok{*}\NormalTok{ r }\SpecialCharTok{{-}} \FunctionTok{log1pexp}\NormalTok{(u))}
\NormalTok{\}}
\end{Highlighting}
\end{Shaded}

Because we minimise the \emph{negative} log-likelihood, the update
\(\beta^{(k+1)} = \beta^{(k)} + J^{-1}S\) becomes
\(\beta^{(k+1)} = \beta^{(k)} - J^{-1} \times (\text{negative score})\),
which is the form coded below.

\begin{Shaded}
\begin{Highlighting}[]
\NormalTok{mle\_logistic\_nr }\OtherTok{\textless{}{-}} \ControlFlowTok{function}\NormalTok{(b0, no\_iter, r, z, }\AttributeTok{debug =} \ConstantTok{FALSE}\NormalTok{) \{}
\NormalTok{    out }\OtherTok{\textless{}{-}} \FunctionTok{matrix}\NormalTok{(}\ConstantTok{NA\_real\_}\NormalTok{, no\_iter }\SpecialCharTok{+} \DecValTok{1}\NormalTok{, }\DecValTok{3}\NormalTok{,}
                  \AttributeTok{dimnames =} \FunctionTok{list}\NormalTok{(}\ConstantTok{NULL}\NormalTok{, }\FunctionTok{c}\NormalTok{(}\StringTok{"beta0"}\NormalTok{, }\StringTok{"beta1"}\NormalTok{, }\StringTok{"neg\_loglike"}\NormalTok{)))}
\NormalTok{    b }\OtherTok{\textless{}{-}}\NormalTok{ b0}
    \ControlFlowTok{for}\NormalTok{ (i }\ControlFlowTok{in} \FunctionTok{seq\_len}\NormalTok{(no\_iter }\SpecialCharTok{+} \DecValTok{1}\NormalTok{)) \{}
\NormalTok{        q }\OtherTok{\textless{}{-}} \FunctionTok{lik\_score\_info}\NormalTok{(b, r, z)}
\NormalTok{        out[i, ] }\OtherTok{\textless{}{-}} \FunctionTok{c}\NormalTok{(b, q}\SpecialCharTok{$}\NormalTok{neg\_loglik)}
        \ControlFlowTok{if}\NormalTok{ (debug) }\FunctionTok{print}\NormalTok{(out[i, ])}
        \ControlFlowTok{if}\NormalTok{ (i }\SpecialCharTok{\textgreater{}}\NormalTok{ no\_iter) }\ControlFlowTok{break}
\NormalTok{        step }\OtherTok{\textless{}{-}} \FunctionTok{tryCatch}\NormalTok{(}\FunctionTok{solve}\NormalTok{(q}\SpecialCharTok{$}\NormalTok{inf, q}\SpecialCharTok{$}\NormalTok{neg\_score),}
                         \AttributeTok{error =} \ControlFlowTok{function}\NormalTok{(e) }\FunctionTok{rep}\NormalTok{(}\ConstantTok{NA\_real\_}\NormalTok{, }\DecValTok{2}\NormalTok{))}
        \ControlFlowTok{if}\NormalTok{ (}\FunctionTok{any}\NormalTok{(}\SpecialCharTok{!}\FunctionTok{is.finite}\NormalTok{(step))) }\ControlFlowTok{break}
\NormalTok{        b }\OtherTok{\textless{}{-}}\NormalTok{ b }\SpecialCharTok{{-}}\NormalTok{ step}
\NormalTok{    \}}
\NormalTok{    out}
\NormalTok{\}}
\end{Highlighting}
\end{Shaded}

\begin{Shaded}
\begin{Highlighting}[]
\NormalTok{gen\_logistic\_data }\OtherTok{\textless{}{-}} \ControlFlowTok{function}\NormalTok{(b, n) \{}
\NormalTok{    z }\OtherTok{\textless{}{-}} \FunctionTok{sort}\NormalTok{(}\FunctionTok{runif}\NormalTok{(n, }\SpecialCharTok{{-}}\DecValTok{2}\NormalTok{, }\DecValTok{2}\NormalTok{))}
\NormalTok{    p }\OtherTok{\textless{}{-}} \FunctionTok{plogis}\NormalTok{(b[}\DecValTok{1}\NormalTok{] }\SpecialCharTok{+}\NormalTok{ b[}\DecValTok{2}\NormalTok{] }\SpecialCharTok{*}\NormalTok{ z)}
\NormalTok{    r }\OtherTok{\textless{}{-}} \FunctionTok{rbinom}\NormalTok{(n, }\DecValTok{1}\NormalTok{, p)}
    \FunctionTok{plot}\NormalTok{(z, p, }\AttributeTok{type =} \StringTok{"l"}\NormalTok{, }\AttributeTok{ylim =} \FunctionTok{c}\NormalTok{(}\DecValTok{0}\NormalTok{, }\DecValTok{1}\NormalTok{), }\AttributeTok{ylab =} \StringTok{"P(r = 1)"}\NormalTok{)}
    \FunctionTok{points}\NormalTok{(z, r, }\AttributeTok{col =}\NormalTok{ r }\SpecialCharTok{+} \DecValTok{1}\NormalTok{)}
    \FunctionTok{list}\NormalTok{(}\AttributeTok{z =}\NormalTok{ z, }\AttributeTok{r =}\NormalTok{ r)}
\NormalTok{\}}

\FunctionTok{set.seed}\NormalTok{(}\DecValTok{1}\NormalTok{)}
\NormalTok{data }\OtherTok{\textless{}{-}} \FunctionTok{gen\_logistic\_data}\NormalTok{(}\FunctionTok{c}\NormalTok{(}\DecValTok{0}\NormalTok{, }\FloatTok{1.5}\NormalTok{), }\DecValTok{200}\NormalTok{)}
\end{Highlighting}
\end{Shaded}

\begin{figure}[H]

{\centering \pandocbounded{\includegraphics[keepaspectratio]{chapters/unit06_MLEm_files/figure-pdf/logistic-data-1.pdf}}

}

\caption{Simulated logistic data: true success probability and observed
responses.}

\end{figure}%

A start close to the truth converges in a handful of iterations:

\begin{Shaded}
\begin{Highlighting}[]
\FunctionTok{round}\NormalTok{(}\FunctionTok{mle\_logistic\_nr}\NormalTok{(}\FunctionTok{c}\NormalTok{(}\DecValTok{0}\NormalTok{, }\DecValTok{3}\NormalTok{), }\DecValTok{8}\NormalTok{, data}\SpecialCharTok{$}\NormalTok{r, data}\SpecialCharTok{$}\NormalTok{z), }\DecValTok{6}\NormalTok{)}
\end{Highlighting}
\end{Shaded}

\begin{verbatim}
          beta0     beta1 neg_loglike
 [1,]  0.000000  3.000000   117.03598
 [2,] -0.505055 -0.306054   164.86061
 [3,] -0.040097  1.281039    99.03765
 [4,] -0.251124  1.470987    97.74624
 [5,] -0.272205  1.502838    97.72964
 [6,] -0.272532  1.503424    97.72964
 [7,] -0.272532  1.503424    97.72964
 [8,] -0.272532  1.503424    97.72964
 [9,] -0.272532  1.503424    97.72964
\end{verbatim}

The doubling of correct digits predicted by the quadratic-convergence
result is visible in the final rows.

\subsection{\texorpdfstring{Comparison with \texttt{nlm} and
\texttt{glm}}{Comparison with nlm and glm}}\label{comparison-with-nlm-and-glm}

\begin{Shaded}
\begin{Highlighting}[]
\NormalTok{logit\_nlm }\OtherTok{\textless{}{-}} \FunctionTok{nlm}\NormalTok{(neg\_logp\_logistic, }\FunctionTok{c}\NormalTok{(}\DecValTok{0}\NormalTok{, }\DecValTok{0}\NormalTok{),}
                 \AttributeTok{z =}\NormalTok{ data}\SpecialCharTok{$}\NormalTok{z, }\AttributeTok{r =}\NormalTok{ data}\SpecialCharTok{$}\NormalTok{r, }\AttributeTok{hessian =} \ConstantTok{TRUE}\NormalTok{)}
\NormalTok{logit\_nlm}\SpecialCharTok{$}\NormalTok{estimate}
\end{Highlighting}
\end{Shaded}

\begin{verbatim}
[1] -0.2725325  1.5034237
\end{verbatim}

\begin{Shaded}
\begin{Highlighting}[]
\FunctionTok{sqrt}\NormalTok{(}\FunctionTok{diag}\NormalTok{(}\FunctionTok{solve}\NormalTok{(logit\_nlm}\SpecialCharTok{$}\NormalTok{hessian)))   }\CommentTok{\# asymptotic standard errors}
\end{Highlighting}
\end{Shaded}

\begin{verbatim}
[1] 0.1797391 0.2089673
\end{verbatim}

\begin{Shaded}
\begin{Highlighting}[]
\FunctionTok{summary}\NormalTok{(}\FunctionTok{glm}\NormalTok{(r }\SpecialCharTok{\textasciitilde{}}\NormalTok{ z, }\AttributeTok{family =} \FunctionTok{binomial}\NormalTok{(), }\AttributeTok{data =}\NormalTok{ data))}\SpecialCharTok{$}\NormalTok{coefficients}
\end{Highlighting}
\end{Shaded}

\begin{verbatim}
              Estimate Std. Error   z value     Pr(>|z|)
(Intercept) -0.2725322  0.1797371 -1.516283 1.294479e-01
z            1.5034244  0.2089473  7.195234 6.235374e-13
\end{verbatim}

The three routes agree to printing accuracy. \texttt{glm} is doing the
same arithmetic as \texttt{mle\_logistic\_nr}: for a GLM the
Newton--Raphson step can be rearranged as a weighted least-squares fit
of the working response \(u_i + (r_i - p_i)/\{p_i(1-p_i)\}\) on \(X\)
with weights \(p_i(1-p_i)\), giving \textbf{iteratively reweighted least
squares (IRLS)}. This is not a different algorithm but a different
arrangement of the same one, and it explains why \texttt{glm} reports
``Fisher Scoring iterations'' in its output.

\section{Divergence}\label{sec-divergence}

\subsection{What goes wrong}\label{what-goes-wrong}

The quadratic approximation is trustworthy only in a neighbourhood of
the current iterate. Starting from a large \(|\beta_1|\), nearly all
fitted probabilities are pushed to \(0\) or \(1\), so the weights
\(p_i(1-p_i)\) are essentially zero, \(J\) is nearly singular, and
\(J^{-1}\) is enormous. The resulting step is huge, the new iterate
lands far outside the region where the quadratic was a reasonable fit,
and the process can oscillate or run off to infinity --- even though the
logistic log-likelihood is globally concave and has a unique maximum.

The interactive figure below lets you choose a starting value and step
through the iterations. Press the play button on the iteration slider to
animate the path over the contours of the negative log-likelihood. This
app runs live only in the HTML edition of this book; readers of other
formats can follow the link below it.

\protect\phantomsection\label{sec-nr-divergence-app}{}

This section contains a live, in-browser Shiny app that only runs in the
HTML edition of this book. View it at
\url{https://longhaisk.github.io/compstat/chapters/unit06_MLEm.html\#sec-nr-divergence-app}.

\begin{tcolorbox}[enhanced jigsaw, arc=.35mm, bottomrule=.15mm, bottomtitle=1mm, breakable, colback=white, colbacktitle=quarto-callout-note-color!10!white, colframe=quarto-callout-note-color-frame, coltitle=black, left=2mm, leftrule=.75mm, opacityback=0, opacitybacktitle=0.6, rightrule=.15mm, title=\textcolor{quarto-callout-note-color}{\faInfo}\hspace{0.5em}{Reproducing this figure}, titlerule=0mm, toprule=.15mm, toptitle=1mm]

The chunk above uses the
\href{https://github.com/quarto-ext/shinylive}{shinylive} extension,
which compiles the app to WebAssembly so that it runs in the reader's
browser with no Shiny server. Install it once per project with

\begin{Shaded}
\begin{Highlighting}[]
\ExtensionTok{quarto}\NormalTok{ add quarto{-}ext/shinylive}
\end{Highlighting}
\end{Shaded}

and add \texttt{filters:\ {[}shinylive{]}} to the document or project
YAML. Only packages available in webR may be used inside the app, so the
code above is restricted to base R and \texttt{shiny}.

\end{tcolorbox}

\subsection{Safeguards}\label{safeguards-1}

The divergence is not a defect of the Newton direction but of the Newton
\emph{step length}. The standard remedies all keep the direction and
control the distance.

\textbf{Step halving (backtracking).} Take \(\delta^{(k)}\) as usual but
move \(\theta^{(k+1)} = \theta^{(k)} + s\,\delta^{(k)}\), starting from
\(s = 1\) and halving until the log-likelihood actually increases. This
costs a few extra function evaluations, never damages quadratic
convergence near the maximum (where \(s = 1\) is accepted on the first
try), and rescues most divergent runs. Toggle it in the app above to see
the effect.

\textbf{Ridging (Levenberg--Marquardt).} Replace \(J\) by
\(J + \lambda \mathbb{I}\) with \(\lambda > 0\). Large \(\lambda\)
shrinks the step toward a small gradient step; \(\lambda \to 0\)
recovers Newton. This also repairs indefiniteness, and is the approach
\texttt{nlm} takes internally.

\textbf{Trust regions.} Maximise the quadratic approximation subject to
\(\|\theta - \theta^{(k)}\| \le \Delta_k\), adapting \(\Delta_k\)
according to how well the quadratic predicted the realised change. This
is \texttt{nlminb} and the \texttt{PORT} routines.

\textbf{Reparameterisation and good starts.} Often the cheapest fix.
Centring and scaling covariates improves conditioning of \(J\)
substantially, and for GLMs the null-model fit is a reliable starting
point --- which is what \texttt{glm} uses when \texttt{start} is not
supplied.

\section{Other Multivariate Optimization
Methods}\label{other-multivariate-optimization-methods}

Newton--Raphson is the method of choice when the second derivatives are
available and \(p\) is modest. Outside that regime, three considerations
push toward alternatives: derivatives may be unavailable or painful to
code; the \(O(p^3)\) cost of solving the Newton system and the
\(O(p^2)\) storage of \(J\) become prohibitive for large \(p\); and the
objective may be non-smooth, noisy or multimodal.

\subsection{Classification by derivative
information}\label{classification-by-derivative-information}

Methods are most usefully organised by what they require of the
objective.

\begin{longtable}[]{@{}
  >{\raggedleft\arraybackslash}p{(\linewidth - 6\tabcolsep) * \real{0.2500}}
  >{\raggedright\arraybackslash}p{(\linewidth - 6\tabcolsep) * \real{0.2500}}
  >{\raggedright\arraybackslash}p{(\linewidth - 6\tabcolsep) * \real{0.2500}}
  >{\raggedright\arraybackslash}p{(\linewidth - 6\tabcolsep) * \real{0.2500}}@{}}
\caption{Optimization methods classified by the derivative information
required}\label{tbl-orders}\tabularnewline
\toprule\noalign{}
\begin{minipage}[b]{\linewidth}\raggedleft
Order
\end{minipage} & \begin{minipage}[b]{\linewidth}\raggedright
Requires
\end{minipage} & \begin{minipage}[b]{\linewidth}\raggedright
Methods
\end{minipage} & \begin{minipage}[b]{\linewidth}\raggedright
Typical use
\end{minipage} \\
\midrule\noalign{}
\endfirsthead
\toprule\noalign{}
\begin{minipage}[b]{\linewidth}\raggedleft
Order
\end{minipage} & \begin{minipage}[b]{\linewidth}\raggedright
Requires
\end{minipage} & \begin{minipage}[b]{\linewidth}\raggedright
Methods
\end{minipage} & \begin{minipage}[b]{\linewidth}\raggedright
Typical use
\end{minipage} \\
\midrule\noalign{}
\endhead
\bottomrule\noalign{}
\endlastfoot
0 & Function values only & Nelder--Mead, simulated annealing, genetic
algorithms & Non-smooth, discontinuous, or awkward-to-differentiate
objectives; small \(p\) \\
1 & Gradient & Steepest descent, conjugate gradient, quasi-Newton (BFGS,
L-BFGS), stochastic gradient descent & The workhorse class; large \(p\);
gradients cheap by automatic differentiation \\
2 & Gradient and Hessian & Newton--Raphson, Fisher scoring, IRLS,
trust-region Newton & Statistical models of modest dimension where the
Hessian is also wanted for inference \\
\end{longtable}

\subsection{The main alternatives}\label{the-main-alternatives}

\textbf{Nelder--Mead (simplex).} Maintains \(p+1\) points forming a
simplex and moves it downhill by reflecting, expanding, contracting and
shrinking. It needs no derivatives and tolerates mild non-smoothness,
but converges slowly, degrades badly as \(p\) grows past roughly 10--20,
and can stall at non-stationary points. It is \texttt{optim}'s default,
which makes it by far the most-used and most-overused method in R.
Reasonable as a first look at an unfamiliar objective; a poor choice for
a likelihood you will fit thousands of times.

\textbf{Steepest descent with line search.} Move along \(-\nabla f\),
choosing the distance by a one-dimensional search. Conceptually the
simplest gradient method and the one that best illustrates the role of a
line search, but the convergence is only linear, and the rate
deteriorates with the condition number of the Hessian: on an elongated
valley the iterates zig-zag across it rather than travelling along it.
Rarely used unmodified.

\textbf{Conjugate gradient.} Corrects the zig-zag by choosing each
direction to be conjugate to the previous ones with respect to the local
curvature, without ever forming or storing a Hessian. Storage is
\(O(p)\), which is its defining advantage. Fletcher--Reeves and
Polak--Ribière differ only in the formula for the correction
coefficient. The method is sensitive to the accuracy of the line search
and usually needs periodic restarting.

\textbf{Quasi-Newton (BFGS and L-BFGS).} Builds an approximation to
\(J^{-1}\) from successive gradient differences, so it attains
superlinear convergence with only first derivatives. BFGS stores a full
\(p \times p\) approximation; L-BFGS keeps only the last few update
vectors, reducing storage to \(O(mp)\) and making it the default choice
for high-dimensional smooth problems. For most likelihood problems where
the Hessian is awkward to derive, BFGS with an analytic gradient is the
best practical compromise.

\textbf{Stochastic gradient descent.} For an objective that is a sum
over \(n\) observations, estimate the gradient from a small random
subset (``minibatch'') at each step. The per-iteration cost becomes
independent of \(n\), at the price of a noisy trajectory requiring a
decreasing step-size sequence (\(\sum_k \alpha_k = \infty\),
\(\sum_k \alpha_k^2 < \infty\)). Modern variants --- momentum, AdaGrad,
RMSProp, Adam --- adapt the step size per coordinate. SGD is the
foundation of neural-network fitting and is treated in its own right in
a later chapter; for classical likelihood problems with \(n\) in the
thousands it offers nothing over BFGS.

\subsection{Comparison}\label{comparison}

\begin{longtable}[]{@{}
  >{\raggedright\arraybackslash}p{(\linewidth - 12\tabcolsep) * \real{0.1429}}
  >{\raggedright\arraybackslash}p{(\linewidth - 12\tabcolsep) * \real{0.1429}}
  >{\raggedright\arraybackslash}p{(\linewidth - 12\tabcolsep) * \real{0.1429}}
  >{\raggedright\arraybackslash}p{(\linewidth - 12\tabcolsep) * \real{0.1429}}
  >{\raggedright\arraybackslash}p{(\linewidth - 12\tabcolsep) * \real{0.1429}}
  >{\raggedright\arraybackslash}p{(\linewidth - 12\tabcolsep) * \real{0.1429}}
  >{\raggedright\arraybackslash}p{(\linewidth - 12\tabcolsep) * \real{0.1429}}@{}}
\caption{Comparison of multivariate optimization methods. \(p\) is the
number of parameters and \(m\) the L-BFGS memory length (typically
5--20).}\label{tbl-comparison}\tabularnewline
\toprule\noalign{}
\begin{minipage}[b]{\linewidth}\raggedright
Method
\end{minipage} & \begin{minipage}[b]{\linewidth}\raggedright
Derivatives
\end{minipage} & \begin{minipage}[b]{\linewidth}\raggedright
Storage
\end{minipage} & \begin{minipage}[b]{\linewidth}\raggedright
Cost per iteration
\end{minipage} & \begin{minipage}[b]{\linewidth}\raggedright
Local rate
\end{minipage} & \begin{minipage}[b]{\linewidth}\raggedright
Robust start
\end{minipage} & \begin{minipage}[b]{\linewidth}\raggedright
Hessian for SEs
\end{minipage} \\
\midrule\noalign{}
\endfirsthead
\toprule\noalign{}
\begin{minipage}[b]{\linewidth}\raggedright
Method
\end{minipage} & \begin{minipage}[b]{\linewidth}\raggedright
Derivatives
\end{minipage} & \begin{minipage}[b]{\linewidth}\raggedright
Storage
\end{minipage} & \begin{minipage}[b]{\linewidth}\raggedright
Cost per iteration
\end{minipage} & \begin{minipage}[b]{\linewidth}\raggedright
Local rate
\end{minipage} & \begin{minipage}[b]{\linewidth}\raggedright
Robust start
\end{minipage} & \begin{minipage}[b]{\linewidth}\raggedright
Hessian for SEs
\end{minipage} \\
\midrule\noalign{}
\endhead
\bottomrule\noalign{}
\endlastfoot
Nelder--Mead & none & \(O(p^2)\) & \(O(p)\) evaluations & sublinear &
good & no \\
Steepest descent & gradient & \(O(p)\) & 1 gradient + line search &
linear (slow) & good & no \\
Conjugate gradient & gradient & \(O(p)\) & 1 gradient + line search &
linear (fast) & moderate & no \\
BFGS & gradient & \(O(p^2)\) & 1 gradient + \(O(p^2)\) & superlinear &
moderate & approximate \\
L-BFGS & gradient & \(O(mp)\) & 1 gradient + \(O(mp)\) & superlinear &
moderate & no \\
Newton--Raphson & gradient, Hessian & \(O(p^2)\) & 1 Hessian +
\(O(p^3)\) & quadratic & poor & exact \\
Fisher scoring / IRLS & score, expected information & \(O(p^2)\) & 1
weighted LS fit & near-quadratic & fair & exact \\
Stochastic gradient & minibatch gradient & \(O(p)\) &
\(O(\text{batch})\) & sublinear & good & no \\
\end{longtable}

\subsection{Implementations in R}\label{implementations-in-r}

\begin{longtable}[]{@{}
  >{\raggedright\arraybackslash}p{(\linewidth - 6\tabcolsep) * \real{0.2500}}
  >{\raggedright\arraybackslash}p{(\linewidth - 6\tabcolsep) * \real{0.2500}}
  >{\raggedright\arraybackslash}p{(\linewidth - 6\tabcolsep) * \real{0.2500}}
  >{\raggedright\arraybackslash}p{(\linewidth - 6\tabcolsep) * \real{0.2500}}@{}}
\caption{R interfaces to multivariate
optimizers}\label{tbl-r-functions}\tabularnewline
\toprule\noalign{}
\begin{minipage}[b]{\linewidth}\raggedright
Function
\end{minipage} & \begin{minipage}[b]{\linewidth}\raggedright
Package
\end{minipage} & \begin{minipage}[b]{\linewidth}\raggedright
Methods available
\end{minipage} & \begin{minipage}[b]{\linewidth}\raggedright
Notes
\end{minipage} \\
\midrule\noalign{}
\endfirsthead
\toprule\noalign{}
\begin{minipage}[b]{\linewidth}\raggedright
Function
\end{minipage} & \begin{minipage}[b]{\linewidth}\raggedright
Package
\end{minipage} & \begin{minipage}[b]{\linewidth}\raggedright
Methods available
\end{minipage} & \begin{minipage}[b]{\linewidth}\raggedright
Notes
\end{minipage} \\
\midrule\noalign{}
\endhead
\bottomrule\noalign{}
\endlastfoot
\texttt{optim} & stats & Nelder--Mead, BFGS, CG, L-BFGS-B, SANN &
General workhorse; supply \texttt{gr} or the gradient is differenced;
\texttt{hessian\ =\ TRUE} for inference; L-BFGS-B allows box
constraints \\
\texttt{nlm} & stats & Newton-type with ridging & Uses analytic
gradient/Hessian if attached as attributes, otherwise numerical; returns
\texttt{\$hessian} \\
\texttt{nlminb} & stats & Trust-region (PORT) & Reliable on badly scaled
problems; supports box constraints \\
\texttt{optimize}, \texttt{uniroot} & stats & Golden section, Brent &
One dimension only \\
\texttt{glm} & stats & IRLS (Fisher scoring) & The specialised case; use
\texttt{start} when separation is suspected \\
\texttt{nls} & stats & Gauss--Newton, Golub--Pereyra & Nonlinear least
squares; \texttt{port} algorithm adds constraints \\
\texttt{optimx} & optimx & Unified front end to a dozen optimizers &
Useful for comparing methods on the same objective \\
\texttt{nloptr} & nloptr & Large NLopt collection, with constraints &
Global and derivative-free options \\
\texttt{Rsolnp}, \texttt{alabama} & --- & Nonlinear equality and
inequality constraints & For genuinely constrained problems \\
\end{longtable}

\subsection{Choosing a method}\label{choosing-a-method}

\begin{longtable}[]{@{}
  >{\raggedright\arraybackslash}p{(\linewidth - 2\tabcolsep) * \real{0.5000}}
  >{\raggedright\arraybackslash}p{(\linewidth - 2\tabcolsep) * \real{0.5000}}@{}}
\caption{Practical guidance on method
selection}\label{tbl-choosing}\tabularnewline
\toprule\noalign{}
\begin{minipage}[b]{\linewidth}\raggedright
Situation
\end{minipage} & \begin{minipage}[b]{\linewidth}\raggedright
Suggested approach
\end{minipage} \\
\midrule\noalign{}
\endfirsthead
\toprule\noalign{}
\begin{minipage}[b]{\linewidth}\raggedright
Situation
\end{minipage} & \begin{minipage}[b]{\linewidth}\raggedright
Suggested approach
\end{minipage} \\
\midrule\noalign{}
\endhead
\bottomrule\noalign{}
\endlastfoot
Hessian derivable, \(p\) small, standard errors wanted & Newton--Raphson
or Fisher scoring \\
GLM or a model expressible as one & \texttt{glm} (IRLS) \\
Gradient available, \(p\) moderate & BFGS via
\texttt{optim(method\ =\ "BFGS")} \\
Gradient available, \(p\) large & L-BFGS \\
Derivatives unavailable, \(p < 10\) & Nelder--Mead as a first pass, then
refine \\
Objective noisy or non-smooth & Nelder--Mead or simulated annealing \\
Divergence or poor conditioning suspected & \texttt{nlm} or
\texttt{nlminb}; rescale the parameters \\
Objective a sum over very many observations & Stochastic gradient
descent and variants \\
Multiple local maxima suspected & Multiple random starts with any of the
above \\
\end{longtable}

The last row deserves emphasis. Every method in this chapter is a
\emph{local} method: each converges to a stationary point near where it
started, and none can certify that the point found is the global
maximum. Restarting from a dispersed set of initial values and comparing
the attained log-likelihoods is the only general safeguard, and it
should be routine for any likelihood that is not known to be concave.

\bookmarksetup{startatroot}

\chapter{The EM Algorithm}\label{the-em-algorithm}

\section{Overview of EM}\label{overview-of-em}

\subsection{Setup and motivation}\label{setup-and-motivation}

Let \(Y\) denote the observed data and \(\theta\) the parameter, with
likelihood \(L(\theta) = P(Y \mid \theta)\). In many models this
quantity is itself an integral (or sum) over latent or missing data
\(Z\):

\[
P(Y \mid \theta) = \int P(Y, Z \mid \theta) \, dZ .
\]

The integral is what makes direct maximization awkward: it destroys the
product structure that would otherwise make the log-likelihood a sum of
simple terms. The EM algorithm avoids it by working with the
complete-data likelihood, which retains that structure, and correcting
for the fact that \(Z\) is unobserved.

\begin{longtable}[]{@{}
  >{\raggedright\arraybackslash}p{(\linewidth - 4\tabcolsep) * \real{0.3333}}
  >{\raggedright\arraybackslash}p{(\linewidth - 4\tabcolsep) * \real{0.3333}}
  >{\raggedright\arraybackslash}p{(\linewidth - 4\tabcolsep) * \real{0.3333}}@{}}
\caption{Observed and complete
likelihoods}\label{tbl-terminology}\tabularnewline
\toprule\noalign{}
\begin{minipage}[b]{\linewidth}\raggedright
\end{minipage} & \begin{minipage}[b]{\linewidth}\raggedright
Likelihood
\end{minipage} & \begin{minipage}[b]{\linewidth}\raggedright
Log-likelihood
\end{minipage} \\
\midrule\noalign{}
\endfirsthead
\toprule\noalign{}
\begin{minipage}[b]{\linewidth}\raggedright
\end{minipage} & \begin{minipage}[b]{\linewidth}\raggedright
Likelihood
\end{minipage} & \begin{minipage}[b]{\linewidth}\raggedright
Log-likelihood
\end{minipage} \\
\midrule\noalign{}
\endhead
\bottomrule\noalign{}
\endlastfoot
Observed & \(L_{\text{obs}}(\theta) = P(Y \mid \theta)\) &
\(\ell_{\text{obs}}(\theta) = \log P(Y \mid \theta)\) \\
Complete & \(L_{\text{comp}}(\theta) = P(Y, Z \mid \theta)\) &
\(\ell_{\text{comp}}(\theta; Y, Z) = \log P(Y, Z \mid \theta)\) \\
\end{longtable}

\subsection{The algorithm}\label{the-algorithm}

Start from \(\hat\theta^{(0)}\). For \(t = 0, 1, 2, \ldots\):

\textbf{E-step.} Replace \(\ell_{\text{obs}}(\theta)\) by the expected
complete log-likelihood,

\[
Q\!\left(\theta \mid \hat\theta^{(t)}\right)
= E_{Z \mid Y, \hat\theta^{(t)}}\!\left[\ell_{\text{comp}}(\theta; Y, Z)\right].
\]

\textbf{M-step.} Maximize it,

\[
\hat\theta^{(t+1)} = \arg\max_{\theta} \; Q\!\left(\theta \mid \hat\theta^{(t)}\right).
\]

Iterate until \(\hat\theta^{(t)}\) converges.

\begin{tcolorbox}[enhanced jigsaw, arc=.35mm, bottomrule=.15mm, bottomtitle=1mm, breakable, colback=white, colbacktitle=quarto-callout-warning-color!10!white, colframe=quarto-callout-warning-color-frame, coltitle=black, left=2mm, leftrule=.75mm, opacityback=0, opacitybacktitle=0.6, rightrule=.15mm, title=\textcolor{quarto-callout-warning-color}{\faExclamationTriangle}\hspace{0.5em}{A common misreading}, titlerule=0mm, toprule=.15mm, toptitle=1mm]

\(Q\) is an \emph{average of} \(\ell_{\text{comp}}(\theta; Y, Z)\) over
\(P(Z \mid Y, \hat\theta^{(t)})\). It is \textbf{not}
\(\ell_{\text{comp}}\!\left(\theta; Y, E(Z \mid Y, \hat\theta^{(t)})\right)\)
--- that is, EM is not ``fill in the missing values with their
expectations and maximize''. The two coincide only when
\(\ell_{\text{comp}}\) happens to be linear in \(Z\), which is true in
both examples of this chapter and false in general.

\end{tcolorbox}

\subsection{\texorpdfstring{\(Q\) as an average of complete
log-likelihoods}{Q as an average of complete log-likelihoods}}\label{q-as-an-average-of-complete-log-likelihoods}

Fix \(\hat\theta^{(t)}\). Each draw
\(Z^{(s)} \sim P(Z \mid Y, \hat\theta^{(t)})\) gives one curve
\(\ell_{\text{comp}}(\theta; Y, Z^{(s)})\), and
\(Q(\theta \mid \hat\theta^{(t)})\) is their expectation. The figure
uses the censored Poisson model of Section~\ref{sec-poisson} with
\(\lambda^{(t)} = 1.5\).

\begin{Shaded}
\begin{Highlighting}[]
\NormalTok{lt }\OtherTok{\textless{}{-}} \FloatTok{1.5}\NormalTok{; S }\OtherTok{\textless{}{-}} \DecValTok{10}
\NormalTok{Z }\OtherTok{\textless{}{-}} \FunctionTok{replicate}\NormalTok{(S, }\FunctionTok{rbinom}\NormalTok{(ncen\_d, }\DecValTok{1}\NormalTok{, lt }\SpecialCharTok{/}\NormalTok{ (}\DecValTok{1} \SpecialCharTok{+}\NormalTok{ lt)))}
\NormalTok{curves }\OtherTok{\textless{}{-}} \FunctionTok{sapply}\NormalTok{(}\DecValTok{1}\SpecialCharTok{:}\NormalTok{S, }\ControlFlowTok{function}\NormalTok{(s) }\FunctionTok{sapply}\NormalTok{(lam, lcomp, }\AttributeTok{z =}\NormalTok{ Z[, s]))}
\NormalTok{Qexact }\OtherTok{\textless{}{-}} \FunctionTok{sapply}\NormalTok{(lam, Qfun, }\AttributeTok{lt =}\NormalTok{ lt)}
\FunctionTok{par}\NormalTok{(}\AttributeTok{mar =} \FunctionTok{c}\NormalTok{(}\DecValTok{4}\NormalTok{, }\DecValTok{4}\NormalTok{, }\DecValTok{1}\NormalTok{, }\DecValTok{1}\NormalTok{))}
\FunctionTok{matplot}\NormalTok{(lam, curves, }\AttributeTok{type =} \StringTok{"l"}\NormalTok{, }\AttributeTok{lty =} \DecValTok{1}\NormalTok{, }\AttributeTok{lwd =} \DecValTok{1}\NormalTok{, }\AttributeTok{col =} \StringTok{"gray60"}\NormalTok{,}
        \AttributeTok{xlab =} \FunctionTok{expression}\NormalTok{(theta), }\AttributeTok{ylab =} \StringTok{"log{-}likelihood"}\NormalTok{,}
        \AttributeTok{ylim =} \FunctionTok{range}\NormalTok{(curves) }\SpecialCharTok{+} \FunctionTok{c}\NormalTok{(}\DecValTok{0}\NormalTok{, }\DecValTok{2}\NormalTok{))}
\FunctionTok{lines}\NormalTok{(lam, }\FunctionTok{rowMeans}\NormalTok{(curves), }\AttributeTok{lwd =} \DecValTok{3}\NormalTok{, }\AttributeTok{col =} \StringTok{"firebrick"}\NormalTok{)}
\FunctionTok{lines}\NormalTok{(lam, Qexact, }\AttributeTok{lwd =} \DecValTok{3}\NormalTok{, }\AttributeTok{col =} \StringTok{"black"}\NormalTok{, }\AttributeTok{lty =} \DecValTok{2}\NormalTok{)}
\FunctionTok{abline}\NormalTok{(}\AttributeTok{v =}\NormalTok{ lt, }\AttributeTok{lty =} \DecValTok{3}\NormalTok{)}
\FunctionTok{legend}\NormalTok{(}\StringTok{"bottomright"}\NormalTok{, }\AttributeTok{bty =} \StringTok{"n"}\NormalTok{, }\AttributeTok{lwd =} \FunctionTok{c}\NormalTok{(}\DecValTok{1}\NormalTok{, }\DecValTok{3}\NormalTok{, }\DecValTok{3}\NormalTok{), }\AttributeTok{lty =} \FunctionTok{c}\NormalTok{(}\DecValTok{1}\NormalTok{, }\DecValTok{1}\NormalTok{, }\DecValTok{2}\NormalTok{),}
       \AttributeTok{col =} \FunctionTok{c}\NormalTok{(}\StringTok{"gray60"}\NormalTok{, }\StringTok{"firebrick"}\NormalTok{, }\StringTok{"black"}\NormalTok{),}
       \AttributeTok{legend =} \FunctionTok{c}\NormalTok{(}\FunctionTok{expression}\NormalTok{(l[comp](theta}\SpecialCharTok{*}\StringTok{";"}\SpecialCharTok{*}\NormalTok{Y}\SpecialCharTok{*}\StringTok{","}\SpecialCharTok{*}\NormalTok{Z}\SpecialCharTok{\^{}}\NormalTok{(s)) }\SpecialCharTok{\textasciitilde{}} \StringTok{"for 10 draws"}\NormalTok{),}
                  \StringTok{"average of the 10 curves"}\NormalTok{,}
                  \FunctionTok{expression}\NormalTok{(}\FunctionTok{Q}\NormalTok{(theta}\SpecialCharTok{*}\StringTok{"|"}\SpecialCharTok{*}\FunctionTok{hat}\NormalTok{(theta)}\SpecialCharTok{\^{}}\NormalTok{(t)) }\SpecialCharTok{\textasciitilde{}} \StringTok{"(exact expectation)"}\NormalTok{)))}
\end{Highlighting}
\end{Shaded}

\begin{figure}[H]

\centering{

\pandocbounded{\includegraphics[keepaspectratio]{chapters/unit07_em_files/figure-pdf/fig-q-average-1.pdf}}

}

\caption{\label{fig-q-average}Ten complete-data log-likelihood curves,
their average, and the exact \(Q\) function.}

\end{figure}%

\subsection{\texorpdfstring{\(Q\) functions as tangent lower
bounds}{Q functions as tangent lower bounds}}\label{q-functions-as-tangent-lower-bounds}

The geometric picture that explains why EM works: each
\(Q(\cdot \mid \hat\theta^{(t)})\), shifted by a constant, lies
\textbf{below} \(\ell_{\text{obs}}\) and \textbf{touches} it at
\(\hat\theta^{(t)}\). Maximizing the bound therefore moves uphill on
\(\ell_{\text{obs}}\) itself.

\begin{Shaded}
\begin{Highlighting}[]
\NormalTok{lo }\OtherTok{\textless{}{-}} \FunctionTok{sapply}\NormalTok{(lam, lobs); top }\OtherTok{\textless{}{-}} \FunctionTok{max}\NormalTok{(lo)}
\FunctionTok{par}\NormalTok{(}\AttributeTok{mar =} \FunctionTok{c}\NormalTok{(}\DecValTok{4}\NormalTok{, }\DecValTok{4}\NormalTok{, }\DecValTok{1}\NormalTok{, }\DecValTok{1}\NormalTok{))}
\FunctionTok{plot}\NormalTok{(lam, lo, }\AttributeTok{type =} \StringTok{"l"}\NormalTok{, }\AttributeTok{lwd =} \DecValTok{3}\NormalTok{, }\AttributeTok{xlab =} \FunctionTok{expression}\NormalTok{(theta),}
     \AttributeTok{ylab =} \StringTok{"log{-}likelihood"}\NormalTok{, }\AttributeTok{ylim =} \FunctionTok{c}\NormalTok{(top }\SpecialCharTok{{-}} \DecValTok{16}\NormalTok{, top }\SpecialCharTok{+} \FloatTok{2.5}\NormalTok{))}
\NormalTok{lt }\OtherTok{\textless{}{-}} \FloatTok{0.5}\NormalTok{; cols }\OtherTok{\textless{}{-}} \FunctionTok{c}\NormalTok{(}\StringTok{"firebrick"}\NormalTok{, }\StringTok{"darkorange"}\NormalTok{, }\StringTok{"forestgreen"}\NormalTok{)}
\ControlFlowTok{for}\NormalTok{ (k }\ControlFlowTok{in} \DecValTok{1}\SpecialCharTok{:}\DecValTok{3}\NormalTok{) \{}
\NormalTok{  shift }\OtherTok{\textless{}{-}} \FunctionTok{lobs}\NormalTok{(lt) }\SpecialCharTok{{-}} \FunctionTok{Qfun}\NormalTok{(lt, lt)            }\CommentTok{\# make Q touch l\_obs at theta\^{}(t)}
  \FunctionTok{lines}\NormalTok{(lam, }\FunctionTok{sapply}\NormalTok{(lam, Qfun, }\AttributeTok{lt =}\NormalTok{ lt) }\SpecialCharTok{+}\NormalTok{ shift, }\AttributeTok{col =}\NormalTok{ cols[k], }\AttributeTok{lwd =} \DecValTok{2}\NormalTok{, }\AttributeTok{lty =} \DecValTok{2}\NormalTok{)}
  \FunctionTok{points}\NormalTok{(lt, }\FunctionTok{lobs}\NormalTok{(lt), }\AttributeTok{pch =} \DecValTok{19}\NormalTok{, }\AttributeTok{col =}\NormalTok{ cols[k])}
\NormalTok{  lt\_new }\OtherTok{\textless{}{-}}\NormalTok{ (}\FunctionTok{sum}\NormalTok{(yobs\_d) }\SpecialCharTok{+}\NormalTok{ ncen\_d }\SpecialCharTok{*}\NormalTok{ lt }\SpecialCharTok{/}\NormalTok{ (}\DecValTok{1} \SpecialCharTok{+}\NormalTok{ lt)) }\SpecialCharTok{/}\NormalTok{ n\_d     }\CommentTok{\# M{-}step}
  \FunctionTok{segments}\NormalTok{(lt\_new, top }\SpecialCharTok{{-}} \DecValTok{16}\NormalTok{, lt\_new, }\FunctionTok{lobs}\NormalTok{(lt\_new), }\AttributeTok{col =}\NormalTok{ cols[k], }\AttributeTok{lty =} \DecValTok{3}\NormalTok{)}
  \FunctionTok{text}\NormalTok{(lt, }\FunctionTok{lobs}\NormalTok{(lt), }\FunctionTok{bquote}\NormalTok{(}\FunctionTok{hat}\NormalTok{(theta)}\SpecialCharTok{\^{}}\NormalTok{(.(k }\SpecialCharTok{{-}} \DecValTok{1}\NormalTok{))), }\AttributeTok{pos =} \DecValTok{3}\NormalTok{, }\AttributeTok{col =}\NormalTok{ cols[k], }\AttributeTok{xpd =} \ConstantTok{NA}\NormalTok{)}
\NormalTok{  lt }\OtherTok{\textless{}{-}}\NormalTok{ lt\_new}
\NormalTok{\}}
\FunctionTok{legend}\NormalTok{(}\StringTok{"bottomright"}\NormalTok{, }\AttributeTok{bty =} \StringTok{"n"}\NormalTok{, }\AttributeTok{lwd =} \FunctionTok{c}\NormalTok{(}\DecValTok{3}\NormalTok{, }\DecValTok{2}\NormalTok{), }\AttributeTok{lty =} \FunctionTok{c}\NormalTok{(}\DecValTok{1}\NormalTok{, }\DecValTok{2}\NormalTok{),}
       \AttributeTok{legend =} \FunctionTok{c}\NormalTok{(}\FunctionTok{expression}\NormalTok{(l[obs](theta)),}
                  \FunctionTok{expression}\NormalTok{(}\FunctionTok{Q}\NormalTok{(theta}\SpecialCharTok{*}\StringTok{"|"}\SpecialCharTok{*}\FunctionTok{hat}\NormalTok{(theta)}\SpecialCharTok{\^{}}\NormalTok{(t)) }\SpecialCharTok{+} \StringTok{"const"}\NormalTok{)))}
\end{Highlighting}
\end{Shaded}

\begin{figure}[H]

\centering{

\pandocbounded{\includegraphics[keepaspectratio]{chapters/unit07_em_files/figure-pdf/fig-tangent-bound-1.pdf}}

}

\caption{\label{fig-tangent-bound}Three EM iterations viewed as
successive tangent minorants of the observed log-likelihood.}

\end{figure}%

\subsection{The ascending property}\label{the-ascending-property}

Write \(h(Y, Z \mid \theta) = P(Y, Z \mid \theta)\),
\(g(Y \mid \theta) = \int h(Y, Z \mid \theta) \, dZ\), and
\(K(Z \mid Y, \theta) = h / g\).

\begin{tcolorbox}[enhanced jigsaw, arc=.35mm, bottomrule=.15mm, bottomtitle=1mm, breakable, colback=white, colbacktitle=quarto-callout-note-color!10!white, colframe=quarto-callout-note-color-frame, coltitle=black, left=2mm, leftrule=.75mm, opacityback=0, opacitybacktitle=0.6, rightrule=.15mm, title=\textcolor{quarto-callout-note-color}{\faInfo}\hspace{0.5em}{Theorem}, titlerule=0mm, toprule=.15mm, toptitle=1mm]

\(\ell_{\text{obs}}\!\left(\hat\theta^{(t+1)}\right) \ \ge\ \ell_{\text{obs}}\!\left(\hat\theta^{(t)}\right)\)
at every iteration.

\end{tcolorbox}

\textbf{Key identity.} Since
\(\log g(Y \mid \theta) = \log h(Y, Z \mid \theta) - \log K(Z \mid Y, \theta)\)
holds for every \(Z\), taking \(E_{Z \mid Y, \theta_0}\) of both sides
gives, for any \(\theta_0\),

\[
\ell_{\text{obs}}(\theta)
= \underbrace{E_{Z \mid Y, \theta_0}\!\left[\log h(Y, Z \mid \theta)\right]}_{Q(\theta \mid \theta_0)}
- \underbrace{E_{Z \mid Y, \theta_0}\!\left[\log K(Z \mid Y, \theta)\right]}_{H(\theta \mid \theta_0)} .
\]

The left side does not depend on \(Z\), so the expectation leaves it
unchanged; this is the whole content of the step.

\textbf{Proof.} Set \(\theta_0 = \hat\theta^{(t)}\) and subtract the
identity evaluated at \(\theta = \hat\theta^{(t)}\) from the one
evaluated at \(\theta = \hat\theta^{(t+1)}\):

\[
\begin{aligned}
\ell_{\text{obs}}\!\left(\hat\theta^{(t+1)}\right) - \ell_{\text{obs}}\!\left(\hat\theta^{(t)}\right)
&= \underbrace{\left[Q\!\left(\hat\theta^{(t+1)} \mid \hat\theta^{(t)}\right) - Q\!\left(\hat\theta^{(t)} \mid \hat\theta^{(t)}\right)\right]}_{\ge\, 0 \text{ by the M-step}} \\
&\quad - \underbrace{\left[H\!\left(\hat\theta^{(t+1)} \mid \hat\theta^{(t)}\right) - H\!\left(\hat\theta^{(t)} \mid \hat\theta^{(t)}\right)\right]}_{\le\, 0 \text{ by Jensen, below}} .
\end{aligned}
\]

\textbf{The \(H\) term.} We need
\(H(\theta \mid \hat\theta^{(t)}) \le H(\hat\theta^{(t)} \mid \hat\theta^{(t)})\)
for every \(\theta\), that is

\[
E_{Z \mid Y, \hat\theta^{(t)}}\!\left[\log \frac{K(Z \mid Y, \theta)}{K(Z \mid Y, \hat\theta^{(t)})}\right] \le 0 .
\]

Because \(\log\) is concave, Jensen's inequality gives
\(E[\log W] \le \log E[W]\), so

\[
\begin{aligned}
E_{Z \mid Y, \hat\theta^{(t)}}\!\left[\log \frac{K(Z \mid Y, \theta)}{K(Z \mid Y, \hat\theta^{(t)})}\right]
&\le \log \int \frac{K(Z \mid Y, \theta)}{K(Z \mid Y, \hat\theta^{(t)})} \, K(Z \mid Y, \hat\theta^{(t)}) \, dZ \\
&= \log \int K(Z \mid Y, \theta) \, dZ = \log 1 = 0 . \qquad \blacksquare
\end{aligned}
\]

The quantity on the left is
\(-\mathrm{KL}\!\left(K(\cdot \mid Y, \hat\theta^{(t)}) \,\|\, K(\cdot \mid Y, \theta)\right)\).
Note carefully what the theorem does and does not say: the observed
likelihood never decreases, but that guarantees convergence only to a
stationary point, which may be a local maximum or even a saddle point.

\subsection{EM as alternating
optimization}\label{em-as-alternating-optimization}

There is a second formulation that makes the algorithm's structure
clearer and generalizes better. Let \(\tilde P(\cdot)\) be any
distribution on the hidden data \(Z\), and define

\[
F\!\left(\tilde P, \theta\right)
= \int \tilde P(z) \log P(Y, z \mid \theta) \, dz
- \int \tilde P(z) \log \tilde P(z) \, dz .
\]

Using \(P(Y, z \mid \theta) = P(Y \mid \theta) \, P(z \mid Y, \theta)\),

\[
F\!\left(\tilde P, \theta\right)
= \ell_{\text{obs}}(\theta) - \mathrm{KL}\!\left(\tilde P \,\|\, P(\cdot \mid Y, \theta)\right),
\qquad
\mathrm{KL}(P \| q) = \int P(z) \log \frac{P(z)}{q(z)} \, dz \; \ge\; 0 .
\]

Both steps of EM are now coordinate ascent on the single function \(F\):

\begin{itemize}
\tightlist
\item
  \textbf{E-step} maximizes \(F\) over \(\tilde P\) with
  \(\theta = \hat\theta^{(t)}\) fixed. Since \(\mathrm{KL} \ge 0\) with
  equality if and only if
  \(\tilde P = P(\cdot \mid Y, \hat\theta^{(t)})\), the E-step sets
  \(\tilde P\) to the conditional distribution of \(Z\).
\item
  \textbf{M-step} maximizes \(F\) over \(\theta\) with \(\tilde P\)
  fixed. The entropy term is free of \(\theta\), so this is
  \(\arg\max_\theta Q(\theta \mid \hat\theta^{(t)})\).
\end{itemize}

The ascending property is then immediate: alternating maximization of a
single objective cannot decrease it. This view also explains the
variational methods used when the E-step is intractable --- restrict
\(\tilde P\) to a manageable family and accept a bound rather than an
equality.

For completeness, \(\mathrm{KL} \ge 0\) follows from
\(\log x \le x - 1\):

\[
\mathrm{KL}(P \| q) = -\int P(z) \log \frac{q(z)}{P(z)} \, dz
\;\ge\; -\int P(z) \left(\frac{q(z)}{P(z)} - 1\right) dz
= -\int \left[q(z) - P(z)\right] dz = 0 .
\]

\section{Example 1: Two-Component Normal Mixture}\label{sec-mixture}

\subsection{Model and likelihoods}\label{model-and-likelihoods}

\[
Z_i \sim \text{Bern}(p), \qquad
Y_i \mid Z_i = 1 \sim N(\mu_1, 1), \qquad
Y_i \mid Z_i = 0 \sim N(\mu_0, 1) .
\]

Observed: \(y_1, \ldots, y_n\). Hidden: \(z_1, \ldots, z_n\). Parameter
\(\theta = (\mu_0, \mu_1, p)\).

\textbf{Complete likelihood}

\[
P(y, z \mid \theta) = \prod_{i=1}^{n} P(y_i \mid z_i, \theta) \, P(z_i \mid \theta)
= \prod_{i=1}^{n} \varphi(y_i \mid \mu_{z_i}, 1) \, p^{z_i} (1-p)^{1-z_i} .
\]

\textbf{Observed likelihood}, summing out each \(z_i\),

\[
P(y \mid \theta) = \prod_{i=1}^{n}\left[\varphi(y_i \mid \mu_1, 1)\, p + \varphi(y_i \mid \mu_0, 1)(1 - p)\right] .
\]

The sum inside the product is what makes direct maximization awkward.

\begin{Shaded}
\begin{Highlighting}[]
\FunctionTok{par}\NormalTok{(}\AttributeTok{mar =} \FunctionTok{c}\NormalTok{(}\DecValTok{4}\NormalTok{, }\DecValTok{4}\NormalTok{, }\DecValTok{1}\NormalTok{, }\DecValTok{1}\NormalTok{)); mu0 }\OtherTok{\textless{}{-}} \DecValTok{0}\NormalTok{; mu1 }\OtherTok{\textless{}{-}} \DecValTok{4}\NormalTok{; p }\OtherTok{\textless{}{-}} \FloatTok{0.3}
\FunctionTok{curve}\NormalTok{((}\DecValTok{1} \SpecialCharTok{{-}}\NormalTok{ p) }\SpecialCharTok{*} \FunctionTok{dnorm}\NormalTok{(x, mu0), }\SpecialCharTok{{-}}\FloatTok{3.5}\NormalTok{, }\FloatTok{7.5}\NormalTok{, }\AttributeTok{lwd =} \DecValTok{2}\NormalTok{, }\AttributeTok{col =} \StringTok{"steelblue"}\NormalTok{,}
      \AttributeTok{ylab =} \StringTok{"density"}\NormalTok{, }\AttributeTok{xlab =} \StringTok{"y"}\NormalTok{, }\AttributeTok{ylim =} \FunctionTok{c}\NormalTok{(}\DecValTok{0}\NormalTok{, }\FloatTok{0.3}\NormalTok{))}
\FunctionTok{curve}\NormalTok{(p }\SpecialCharTok{*} \FunctionTok{dnorm}\NormalTok{(x, mu1), }\AttributeTok{add =} \ConstantTok{TRUE}\NormalTok{, }\AttributeTok{lwd =} \DecValTok{2}\NormalTok{, }\AttributeTok{col =} \StringTok{"firebrick"}\NormalTok{)}
\FunctionTok{curve}\NormalTok{((}\DecValTok{1} \SpecialCharTok{{-}}\NormalTok{ p) }\SpecialCharTok{*} \FunctionTok{dnorm}\NormalTok{(x, mu0) }\SpecialCharTok{+}\NormalTok{ p }\SpecialCharTok{*} \FunctionTok{dnorm}\NormalTok{(x, mu1), }\AttributeTok{add =} \ConstantTok{TRUE}\NormalTok{, }\AttributeTok{lwd =} \DecValTok{3}\NormalTok{)}
\FunctionTok{legend}\NormalTok{(}\StringTok{"topright"}\NormalTok{, }\AttributeTok{bty =} \StringTok{"n"}\NormalTok{, }\AttributeTok{lwd =} \FunctionTok{c}\NormalTok{(}\DecValTok{2}\NormalTok{, }\DecValTok{2}\NormalTok{, }\DecValTok{3}\NormalTok{),}
       \AttributeTok{col =} \FunctionTok{c}\NormalTok{(}\StringTok{"steelblue"}\NormalTok{, }\StringTok{"firebrick"}\NormalTok{, }\StringTok{"black"}\NormalTok{),}
       \AttributeTok{legend =} \FunctionTok{c}\NormalTok{(}\FunctionTok{expression}\NormalTok{((}\DecValTok{1}\SpecialCharTok{{-}}\NormalTok{p)}\SpecialCharTok{*}\FunctionTok{phi}\NormalTok{(y}\SpecialCharTok{*}\StringTok{"|"}\SpecialCharTok{*}\NormalTok{mu[}\DecValTok{0}\NormalTok{])),}
                  \FunctionTok{expression}\NormalTok{(p}\SpecialCharTok{*}\FunctionTok{phi}\NormalTok{(y}\SpecialCharTok{*}\StringTok{"|"}\SpecialCharTok{*}\NormalTok{mu[}\DecValTok{1}\NormalTok{])),}
                  \FunctionTok{expression}\NormalTok{(}\FunctionTok{f}\NormalTok{(y}\SpecialCharTok{*}\StringTok{"|"}\SpecialCharTok{*}\NormalTok{theta) }\SpecialCharTok{==} \StringTok{"sum"}\NormalTok{)))}
\end{Highlighting}
\end{Shaded}

\begin{figure}[H]

\centering{

\pandocbounded{\includegraphics[keepaspectratio]{chapters/unit07_em_files/figure-pdf/fig-mixture-components-1.pdf}}

}

\caption{\label{fig-mixture-components}The observed density is a sum of
two weighted normal components.}

\end{figure}%

The marginal density may be bimodal or unimodal depending on \(p\) and
\(|\mu_1 - \mu_0|\) --- a point worth keeping in mind, since the
unimodal case is where EM struggles most.

\begin{Shaded}
\begin{Highlighting}[]
\FunctionTok{par}\NormalTok{(}\AttributeTok{mfrow =} \FunctionTok{c}\NormalTok{(}\DecValTok{1}\NormalTok{, }\DecValTok{3}\NormalTok{), }\AttributeTok{mar =} \FunctionTok{c}\NormalTok{(}\DecValTok{4}\NormalTok{, }\DecValTok{4}\NormalTok{, }\DecValTok{2}\NormalTok{, }\DecValTok{1}\NormalTok{))}
\NormalTok{mix }\OtherTok{\textless{}{-}} \ControlFlowTok{function}\NormalTok{(x, mu0, mu1, p) (}\DecValTok{1} \SpecialCharTok{{-}}\NormalTok{ p) }\SpecialCharTok{*} \FunctionTok{dnorm}\NormalTok{(x, mu0) }\SpecialCharTok{+}\NormalTok{ p }\SpecialCharTok{*} \FunctionTok{dnorm}\NormalTok{(x, mu1)}
\ControlFlowTok{for}\NormalTok{ (s }\ControlFlowTok{in} \FunctionTok{list}\NormalTok{(}\FunctionTok{c}\NormalTok{(}\DecValTok{0}\NormalTok{, }\DecValTok{4}\NormalTok{, }\FloatTok{0.5}\NormalTok{), }\FunctionTok{c}\NormalTok{(}\DecValTok{0}\NormalTok{, }\DecValTok{4}\NormalTok{, }\FloatTok{0.15}\NormalTok{), }\FunctionTok{c}\NormalTok{(}\DecValTok{0}\NormalTok{, }\FloatTok{1.5}\NormalTok{, }\FloatTok{0.5}\NormalTok{))) \{}
  \FunctionTok{curve}\NormalTok{(}\FunctionTok{mix}\NormalTok{(x, s[}\DecValTok{1}\NormalTok{], s[}\DecValTok{2}\NormalTok{], s[}\DecValTok{3}\NormalTok{]), }\SpecialCharTok{{-}}\DecValTok{4}\NormalTok{, }\DecValTok{7}\NormalTok{, }\AttributeTok{lwd =} \DecValTok{3}\NormalTok{, }\AttributeTok{xlab =} \StringTok{"y"}\NormalTok{, }\AttributeTok{ylab =} \StringTok{"density"}\NormalTok{,}
        \AttributeTok{main =} \FunctionTok{bquote}\NormalTok{(mu[}\DecValTok{0}\NormalTok{] }\SpecialCharTok{==}\NormalTok{ .(s[}\DecValTok{1}\NormalTok{]) }\SpecialCharTok{\textasciitilde{}} \StringTok{","} \SpecialCharTok{\textasciitilde{}}\NormalTok{ mu[}\DecValTok{1}\NormalTok{] }\SpecialCharTok{==}\NormalTok{ .(s[}\DecValTok{2}\NormalTok{]) }\SpecialCharTok{\textasciitilde{}} \StringTok{","} \SpecialCharTok{\textasciitilde{}}\NormalTok{ p }\SpecialCharTok{==}\NormalTok{ .(s[}\DecValTok{3}\NormalTok{])))}
\NormalTok{  y }\OtherTok{\textless{}{-}} \FunctionTok{ifelse}\NormalTok{(}\FunctionTok{rbinom}\NormalTok{(}\DecValTok{200}\NormalTok{, }\DecValTok{1}\NormalTok{, s[}\DecValTok{3}\NormalTok{]) }\SpecialCharTok{==} \DecValTok{1}\NormalTok{, }\FunctionTok{rnorm}\NormalTok{(}\DecValTok{200}\NormalTok{, s[}\DecValTok{2}\NormalTok{]), }\FunctionTok{rnorm}\NormalTok{(}\DecValTok{200}\NormalTok{, s[}\DecValTok{1}\NormalTok{]))}
  \FunctionTok{rug}\NormalTok{(y, }\AttributeTok{col =} \StringTok{"gray40"}\NormalTok{)}
\NormalTok{\}}
\end{Highlighting}
\end{Shaded}

\begin{figure}[H]

\centering{

\pandocbounded{\includegraphics[keepaspectratio]{chapters/unit07_em_files/figure-pdf/fig-mixture-shapes-1.pdf}}

}

\caption{\label{fig-mixture-shapes}Separation and mixing proportion
determine whether the mixture is visibly bimodal.}

\end{figure}%

\subsection{Complete log-likelihood}\label{complete-log-likelihood}

With
\(\log \varphi(y \mid \mu, 1) = -\tfrac{1}{2}(y - \mu)^2 - \log\sqrt{2\pi}\),

\[
\begin{aligned}
\ell_{\text{comp}}(\theta; y, z)
&= \sum_{i=1}^{n}\left[z_i \log p + (1 - z_i)\log(1-p) - \tfrac{1}{2}(y_i - \mu_{z_i})^2\right] + C \\
&= \left(\sum_{i=1}^{n} z_i\right) \log \frac{p}{1-p} + n \log(1-p) \\
&\quad - \frac{1}{2}\sum_{i=1}^{n}\left[z_i (y_i - \mu_1)^2 + (1 - z_i)(y_i - \mu_0)^2\right] + C .
\end{aligned}
\]

Two features drive everything that follows. The expression is
\textbf{linear in each \(z_i\)}, so the E-step needs only
\(E(Z_i \mid y_i, \hat\theta^{(t)})\) rather than the full conditional
distribution. And given \(z\), the parameters \textbf{separate}: \(p\)
is a Bernoulli proportion while \(\mu_1\) and \(\mu_0\) are group means.

\subsection{E-step: soft
classification}\label{e-step-soft-classification}

Given
\(\hat\theta^{(t)} = (\hat\mu_0^{(t)}, \hat\mu_1^{(t)}, \hat p^{(t)})\),
Bayes' rule gives

\[
\hat p_i^{(t)} \equiv P\!\left(Z_i = 1 \mid y_i, \hat\theta^{(t)}\right)
= \frac{\varphi\!\left(y_i \mid \hat\mu_1^{(t)}, 1\right) \hat p^{(t)}}
       {\varphi\!\left(y_i \mid \hat\mu_0^{(t)}, 1\right)\!\left(1 - \hat p^{(t)}\right)
        + \varphi\!\left(y_i \mid \hat\mu_1^{(t)}, 1\right) \hat p^{(t)}} ,
\]

and therefore

\[
\begin{aligned}
Q\!\left(\theta \mid \hat\theta^{(t)}\right)
&= \left(\sum_i \hat p_i^{(t)}\right) \log \frac{p}{1-p} + n \log(1-p) \\
&\quad - \frac{1}{2}\sum_i \left[\hat p_i^{(t)} (y_i - \mu_1)^2 + \left(1 - \hat p_i^{(t)}\right)(y_i - \mu_0)^2\right] .
\end{aligned}
\]

Computationally, \(\hat p_i^{(t)}\) must be formed on the log scale.
Writing
\(a_i = \log \varphi(y_i \mid \hat\mu_1^{(t)}, 1) + \log \hat p^{(t)}\)
and
\(b_i = \log \varphi(y_i \mid \hat\mu_0^{(t)}, 1) + \log(1 - \hat p^{(t)})\),
we have \(\hat p_i^{(t)} = \exp\{a_i - \operatorname{lse}(a_i, b_i)\}\)
where \(\operatorname{lse}\) is the log-sum-exp function of the chapter
on computer arithmetic. Forming the ratio directly underflows as soon as
an observation sits far from both means.

\begin{Shaded}
\begin{Highlighting}[]
\FunctionTok{par}\NormalTok{(}\AttributeTok{mfrow =} \FunctionTok{c}\NormalTok{(}\DecValTok{1}\NormalTok{, }\DecValTok{2}\NormalTok{), }\AttributeTok{mar =} \FunctionTok{c}\NormalTok{(}\DecValTok{4}\NormalTok{, }\DecValTok{4}\NormalTok{, }\DecValTok{1}\NormalTok{, }\DecValTok{1}\NormalTok{)); mu0 }\OtherTok{\textless{}{-}} \DecValTok{0}\NormalTok{; mu1 }\OtherTok{\textless{}{-}} \DecValTok{4}\NormalTok{; p }\OtherTok{\textless{}{-}} \FloatTok{0.3}
\NormalTok{f0 }\OtherTok{\textless{}{-}} \ControlFlowTok{function}\NormalTok{(x) (}\DecValTok{1} \SpecialCharTok{{-}}\NormalTok{ p) }\SpecialCharTok{*} \FunctionTok{dnorm}\NormalTok{(x, mu0)}
\NormalTok{f1 }\OtherTok{\textless{}{-}} \ControlFlowTok{function}\NormalTok{(x) p }\SpecialCharTok{*} \FunctionTok{dnorm}\NormalTok{(x, mu1)}
\FunctionTok{curve}\NormalTok{(f0, }\SpecialCharTok{{-}}\FloatTok{3.5}\NormalTok{, }\FloatTok{7.5}\NormalTok{, }\AttributeTok{lwd =} \DecValTok{2}\NormalTok{, }\AttributeTok{col =} \StringTok{"steelblue"}\NormalTok{, }\AttributeTok{ylab =} \StringTok{"weighted density"}\NormalTok{,}
      \AttributeTok{xlab =} \StringTok{"y"}\NormalTok{, }\AttributeTok{ylim =} \FunctionTok{c}\NormalTok{(}\DecValTok{0}\NormalTok{, }\FloatTok{0.3}\NormalTok{))}
\FunctionTok{curve}\NormalTok{(f1, }\AttributeTok{add =} \ConstantTok{TRUE}\NormalTok{, }\AttributeTok{lwd =} \DecValTok{2}\NormalTok{, }\AttributeTok{col =} \StringTok{"firebrick"}\NormalTok{)}
\NormalTok{yi }\OtherTok{\textless{}{-}} \FunctionTok{c}\NormalTok{(}\SpecialCharTok{{-}}\FloatTok{0.5}\NormalTok{, }\FloatTok{1.8}\NormalTok{, }\FloatTok{3.2}\NormalTok{)}
\ControlFlowTok{for}\NormalTok{ (y }\ControlFlowTok{in}\NormalTok{ yi) \{}
  \FunctionTok{segments}\NormalTok{(y, }\DecValTok{0}\NormalTok{, y, }\FunctionTok{f0}\NormalTok{(y), }\AttributeTok{col =} \StringTok{"steelblue"}\NormalTok{, }\AttributeTok{lwd =} \DecValTok{4}\NormalTok{)}
  \FunctionTok{segments}\NormalTok{(y, }\FunctionTok{f0}\NormalTok{(y), y, }\FunctionTok{f0}\NormalTok{(y) }\SpecialCharTok{+} \FunctionTok{f1}\NormalTok{(y), }\AttributeTok{col =} \StringTok{"firebrick"}\NormalTok{, }\AttributeTok{lwd =} \DecValTok{4}\NormalTok{)}
  \FunctionTok{text}\NormalTok{(y, }\FunctionTok{f0}\NormalTok{(y) }\SpecialCharTok{+} \FunctionTok{f1}\NormalTok{(y), }\FunctionTok{sprintf}\NormalTok{(}\StringTok{"\%.2f"}\NormalTok{, }\FunctionTok{f1}\NormalTok{(y) }\SpecialCharTok{/}\NormalTok{ (}\FunctionTok{f0}\NormalTok{(y) }\SpecialCharTok{+} \FunctionTok{f1}\NormalTok{(y))), }\AttributeTok{pos =} \DecValTok{3}\NormalTok{, }\AttributeTok{cex =} \FloatTok{0.9}\NormalTok{)}
\NormalTok{\}}
\FunctionTok{legend}\NormalTok{(}\StringTok{"topright"}\NormalTok{, }\AttributeTok{bty =} \StringTok{"n"}\NormalTok{, }\AttributeTok{lwd =} \DecValTok{2}\NormalTok{, }\AttributeTok{col =} \FunctionTok{c}\NormalTok{(}\StringTok{"steelblue"}\NormalTok{, }\StringTok{"firebrick"}\NormalTok{),}
       \AttributeTok{legend =} \FunctionTok{c}\NormalTok{(}\FunctionTok{expression}\NormalTok{((}\DecValTok{1}\SpecialCharTok{{-}}\FunctionTok{hat}\NormalTok{(p))}\SpecialCharTok{*}\FunctionTok{phi}\NormalTok{(y}\SpecialCharTok{*}\StringTok{"|"}\SpecialCharTok{*}\FunctionTok{hat}\NormalTok{(mu)[}\DecValTok{0}\NormalTok{])),}
                  \FunctionTok{expression}\NormalTok{(}\FunctionTok{hat}\NormalTok{(p)}\SpecialCharTok{*}\FunctionTok{phi}\NormalTok{(y}\SpecialCharTok{*}\StringTok{"|"}\SpecialCharTok{*}\FunctionTok{hat}\NormalTok{(mu)[}\DecValTok{1}\NormalTok{]))))}
\FunctionTok{curve}\NormalTok{(}\FunctionTok{f1}\NormalTok{(x) }\SpecialCharTok{/}\NormalTok{ (}\FunctionTok{f0}\NormalTok{(x) }\SpecialCharTok{+} \FunctionTok{f1}\NormalTok{(x)), }\SpecialCharTok{{-}}\FloatTok{3.5}\NormalTok{, }\FloatTok{7.5}\NormalTok{, }\AttributeTok{lwd =} \DecValTok{3}\NormalTok{, }\AttributeTok{xlab =} \StringTok{"y"}\NormalTok{,}
      \AttributeTok{ylab =} \FunctionTok{expression}\NormalTok{(}\FunctionTok{P}\NormalTok{(Z[i] }\SpecialCharTok{==} \DecValTok{1}\SpecialCharTok{*}\StringTok{"|"}\SpecialCharTok{*}\NormalTok{y[i]}\SpecialCharTok{*}\StringTok{","}\SpecialCharTok{*}\FunctionTok{hat}\NormalTok{(theta)}\SpecialCharTok{\^{}}\NormalTok{(t))))}
\FunctionTok{abline}\NormalTok{(}\AttributeTok{h =} \FunctionTok{c}\NormalTok{(}\DecValTok{0}\NormalTok{, }\DecValTok{1}\NormalTok{), }\AttributeTok{lty =} \DecValTok{3}\NormalTok{); }\FunctionTok{points}\NormalTok{(yi, }\FunctionTok{f1}\NormalTok{(yi) }\SpecialCharTok{/}\NormalTok{ (}\FunctionTok{f0}\NormalTok{(yi) }\SpecialCharTok{+} \FunctionTok{f1}\NormalTok{(yi)), }\AttributeTok{pch =} \DecValTok{19}\NormalTok{)}
\end{Highlighting}
\end{Shaded}

\begin{figure}[H]

\centering{

\pandocbounded{\includegraphics[keepaspectratio]{chapters/unit07_em_files/figure-pdf/fig-estep-1.pdf}}

}

\caption{\label{fig-estep}Left: the responsibility at each \(y_i\) is
the red share of the total height. Right: responsibility as a function
of \(y\).}

\end{figure}%

\subsection{M-step: weighted proportions and
means}\label{m-step-weighted-proportions-and-means}

\(Q\) separates into one term in \(p\) and two in \(\mu_1\), \(\mu_0\),
each maximized in closed form.

\begin{longtable}[]{@{}
  >{\raggedright\arraybackslash}p{(\linewidth - 2\tabcolsep) * \real{0.5000}}
  >{\raggedright\arraybackslash}p{(\linewidth - 2\tabcolsep) * \real{0.5000}}@{}}
\caption{M-step updates for the two-component
mixture}\label{tbl-mstep-mixture}\tabularnewline
\toprule\noalign{}
\begin{minipage}[b]{\linewidth}\raggedright
Update
\end{minipage} & \begin{minipage}[b]{\linewidth}\raggedright
Stationary equation \(\Longrightarrow\) solution
\end{minipage} \\
\midrule\noalign{}
\endfirsthead
\toprule\noalign{}
\begin{minipage}[b]{\linewidth}\raggedright
Update
\end{minipage} & \begin{minipage}[b]{\linewidth}\raggedright
Stationary equation \(\Longrightarrow\) solution
\end{minipage} \\
\midrule\noalign{}
\endhead
\bottomrule\noalign{}
\endlastfoot
\(p\) &
\(\dfrac{\partial Q}{\partial p} = \dfrac{\sum_i \hat p_i^{(t)}}{p(1-p)} - \dfrac{n}{1-p} = 0 \;\Longrightarrow\; \hat p^{(t+1)} = \dfrac{1}{n}\sum_{i=1}^{n} \hat p_i^{(t)}\) \\
\(\mu_1\) &
\(\dfrac{\partial Q}{\partial \mu_1} = \sum_i \hat p_i^{(t)}(y_i - \mu_1) = 0 \;\Longrightarrow\; \hat\mu_1^{(t+1)} = \dfrac{\sum_i \hat p_i^{(t)} y_i}{\sum_i \hat p_i^{(t)}}\) \\
\(\mu_0\) & weights \(1 - \hat p_i^{(t)}\):
\(\;\hat\mu_0^{(t+1)} = \dfrac{\sum_i \left(1 - \hat p_i^{(t)}\right) y_i}{\sum_i \left(1 - \hat p_i^{(t)}\right)}\) \\
\end{longtable}

In words: the M-step is the complete-data MLE with each observation
split between the two groups in proportion to its soft classification.
Return to the E-step with \(\hat\theta^{(t+1)}\) and recompute the
responsibilities.

\subsection{Implementation}\label{implementation-1}

\begin{Shaded}
\begin{Highlighting}[]
\NormalTok{log\_sum\_exp }\OtherTok{\textless{}{-}} \ControlFlowTok{function}\NormalTok{(log\_x) \{}
\NormalTok{    max\_log\_x }\OtherTok{\textless{}{-}} \FunctionTok{max}\NormalTok{(log\_x)}
\NormalTok{    max\_log\_x }\SpecialCharTok{+} \FunctionTok{log}\NormalTok{(}\FunctionTok{sum}\NormalTok{(}\FunctionTok{exp}\NormalTok{(log\_x }\SpecialCharTok{{-}}\NormalTok{ max\_log\_x)))}
\NormalTok{\}}

\DocumentationTok{\#\#\# observed log{-}likelihood, theta = (p, mu1, mu0)}
\NormalTok{log\_like\_obs }\OtherTok{\textless{}{-}} \ControlFlowTok{function}\NormalTok{(theta, Y) \{}
\NormalTok{    log\_joint }\OtherTok{\textless{}{-}} \FunctionTok{cbind}\NormalTok{(}\FunctionTok{log}\NormalTok{(theta[}\DecValTok{1}\NormalTok{])     }\SpecialCharTok{+} \FunctionTok{dnorm}\NormalTok{(Y, theta[}\DecValTok{2}\NormalTok{], }\DecValTok{1}\NormalTok{, }\AttributeTok{log =} \ConstantTok{TRUE}\NormalTok{),}
                       \FunctionTok{log}\NormalTok{(}\DecValTok{1} \SpecialCharTok{{-}}\NormalTok{ theta[}\DecValTok{1}\NormalTok{]) }\SpecialCharTok{+} \FunctionTok{dnorm}\NormalTok{(Y, theta[}\DecValTok{3}\NormalTok{], }\DecValTok{1}\NormalTok{, }\AttributeTok{log =} \ConstantTok{TRUE}\NormalTok{))}
    \FunctionTok{sum}\NormalTok{(}\FunctionTok{apply}\NormalTok{(log\_joint, }\DecValTok{1}\NormalTok{, log\_sum\_exp))}
\NormalTok{\}}

\DocumentationTok{\#\#\# E{-}step: responsibilities P(Z\_i = 1 | y\_i, theta)}
\NormalTok{e\_step }\OtherTok{\textless{}{-}} \ControlFlowTok{function}\NormalTok{(theta, Y) \{}
\NormalTok{    l1 }\OtherTok{\textless{}{-}} \FunctionTok{dnorm}\NormalTok{(Y, theta[}\DecValTok{2}\NormalTok{], }\DecValTok{1}\NormalTok{, }\AttributeTok{log =} \ConstantTok{TRUE}\NormalTok{) }\SpecialCharTok{+} \FunctionTok{log}\NormalTok{(theta[}\DecValTok{1}\NormalTok{])}
\NormalTok{    l0 }\OtherTok{\textless{}{-}} \FunctionTok{dnorm}\NormalTok{(Y, theta[}\DecValTok{3}\NormalTok{], }\DecValTok{1}\NormalTok{, }\AttributeTok{log =} \ConstantTok{TRUE}\NormalTok{) }\SpecialCharTok{+} \FunctionTok{log}\NormalTok{(}\DecValTok{1} \SpecialCharTok{{-}}\NormalTok{ theta[}\DecValTok{1}\NormalTok{])}
    \FunctionTok{exp}\NormalTok{(l1 }\SpecialCharTok{{-}} \FunctionTok{apply}\NormalTok{(}\FunctionTok{cbind}\NormalTok{(l1, l0), }\DecValTok{1}\NormalTok{, log\_sum\_exp))}
\NormalTok{\}}

\DocumentationTok{\#\#\# M{-}step: closed{-}form weighted updates}
\NormalTok{m\_step }\OtherTok{\textless{}{-}} \ControlFlowTok{function}\NormalTok{(w, Y) \{}
    \FunctionTok{c}\NormalTok{(}\FunctionTok{mean}\NormalTok{(w),}
      \FunctionTok{sum}\NormalTok{(Y }\SpecialCharTok{*}\NormalTok{ w) }\SpecialCharTok{/} \FunctionTok{sum}\NormalTok{(w),}
      \FunctionTok{sum}\NormalTok{(Y }\SpecialCharTok{*}\NormalTok{ (}\DecValTok{1} \SpecialCharTok{{-}}\NormalTok{ w)) }\SpecialCharTok{/} \FunctionTok{sum}\NormalTok{(}\DecValTok{1} \SpecialCharTok{{-}}\NormalTok{ w))}
\NormalTok{\}}

\NormalTok{em\_mixnorm }\OtherTok{\textless{}{-}} \ControlFlowTok{function}\NormalTok{(theta0, Y, no\_iters) \{}
\NormalTok{    out }\OtherTok{\textless{}{-}} \FunctionTok{matrix}\NormalTok{(}\ConstantTok{NA\_real\_}\NormalTok{, no\_iters }\SpecialCharTok{+} \DecValTok{1}\NormalTok{, }\DecValTok{4}\NormalTok{,}
                  \AttributeTok{dimnames =} \FunctionTok{list}\NormalTok{(}\ConstantTok{NULL}\NormalTok{, }\FunctionTok{c}\NormalTok{(}\StringTok{"p"}\NormalTok{, }\StringTok{"mu1"}\NormalTok{, }\StringTok{"mu0"}\NormalTok{, }\StringTok{"log\_lik"}\NormalTok{)))}
\NormalTok{    theta }\OtherTok{\textless{}{-}}\NormalTok{ theta0}
\NormalTok{    out[}\DecValTok{1}\NormalTok{, ] }\OtherTok{\textless{}{-}} \FunctionTok{c}\NormalTok{(theta, }\FunctionTok{log\_like\_obs}\NormalTok{(theta, Y))}
    \ControlFlowTok{for}\NormalTok{ (i }\ControlFlowTok{in} \FunctionTok{seq\_len}\NormalTok{(no\_iters)) \{}
\NormalTok{        theta }\OtherTok{\textless{}{-}} \FunctionTok{m\_step}\NormalTok{(}\FunctionTok{e\_step}\NormalTok{(theta, Y), Y)}
\NormalTok{        out[i }\SpecialCharTok{+} \DecValTok{1}\NormalTok{, ] }\OtherTok{\textless{}{-}} \FunctionTok{c}\NormalTok{(theta, }\FunctionTok{log\_like\_obs}\NormalTok{(theta, Y))}
\NormalTok{    \}}
\NormalTok{    out}
\NormalTok{\}}
\end{Highlighting}
\end{Shaded}

\begin{Shaded}
\begin{Highlighting}[]
\NormalTok{gen\_mixnorm }\OtherTok{\textless{}{-}} \ControlFlowTok{function}\NormalTok{(theta, n) \{}
\NormalTok{    Z }\OtherTok{\textless{}{-}} \FunctionTok{rbinom}\NormalTok{(n, }\DecValTok{1}\NormalTok{, theta[}\DecValTok{1}\NormalTok{])}
\NormalTok{    Y }\OtherTok{\textless{}{-}} \FunctionTok{rnorm}\NormalTok{(n, }\AttributeTok{mean =} \FunctionTok{ifelse}\NormalTok{(Z }\SpecialCharTok{==} \DecValTok{1}\NormalTok{, theta[}\DecValTok{2}\NormalTok{], theta[}\DecValTok{3}\NormalTok{]), }\AttributeTok{sd =} \DecValTok{1}\NormalTok{)}
    \FunctionTok{plot}\NormalTok{(Y, Z, }\AttributeTok{col =} \FunctionTok{cm.colors}\NormalTok{(}\DecValTok{2}\NormalTok{)[Z }\SpecialCharTok{+} \DecValTok{1}\NormalTok{], }\AttributeTok{pch =} \DecValTok{19}\NormalTok{,}
         \AttributeTok{ylab =} \StringTok{"component label Z"}\NormalTok{)}
\NormalTok{    Y}
\NormalTok{\}}

\FunctionTok{set.seed}\NormalTok{(}\DecValTok{812}\NormalTok{)}
\NormalTok{data }\OtherTok{\textless{}{-}} \FunctionTok{gen\_mixnorm}\NormalTok{(}\FunctionTok{c}\NormalTok{(}\FloatTok{0.3}\NormalTok{, }\DecValTok{0}\NormalTok{, }\DecValTok{3}\NormalTok{), }\DecValTok{200}\NormalTok{)}
\end{Highlighting}
\end{Shaded}

\begin{figure}[H]

{\centering \pandocbounded{\includegraphics[keepaspectratio]{chapters/unit07_em_files/figure-pdf/mix-data-1.pdf}}

}

\caption{Simulated mixture data, coloured by the (normally unobserved)
component label.}

\end{figure}%

A well-chosen start converges quickly:

\begin{Shaded}
\begin{Highlighting}[]
\FunctionTok{round}\NormalTok{(}\FunctionTok{em\_mixnorm}\NormalTok{(}\FunctionTok{c}\NormalTok{(}\FloatTok{0.5}\NormalTok{, }\DecValTok{0}\NormalTok{, }\DecValTok{3}\NormalTok{), data, }\DecValTok{10}\NormalTok{), }\DecValTok{4}\NormalTok{)}
\end{Highlighting}
\end{Shaded}

\begin{verbatim}
           p    mu1    mu0   log_lik
 [1,] 0.5000 0.0000 3.0000 -375.3490
 [2,] 0.3377 0.2695 3.0137 -361.6709
 [3,] 0.3242 0.2319 2.9768 -361.3077
 [4,] 0.3154 0.1956 2.9584 -361.1388
 [5,] 0.3094 0.1694 2.9463 -361.0581
 [6,] 0.3054 0.1513 2.9380 -361.0202
 [7,] 0.3026 0.1389 2.9323 -361.0026
 [8,] 0.3007 0.1306 2.9285 -360.9945
 [9,] 0.2995 0.1249 2.9258 -360.9908
[10,] 0.2986 0.1211 2.9240 -360.9891
[11,] 0.2980 0.1185 2.9228 -360.9884
\end{verbatim}

The monotone increase of the last column is the ascending property in
action; it is the single most useful diagnostic when debugging an EM
implementation, since any decrease indicates an error in the E- or
M-step.

\subsection{Interactive demonstration}\label{sec-mixture-app}

The figure below replaces the frame-by-frame animation of earlier
versions of this chapter. Choose a data set and a starting value ---
either from the presets in the left panel or by typing your own --- then
step through the fit with \textbf{Prev} and \textbf{Next}, or press
\textbf{Play} to run it automatically. \textbf{Regenerate data} draws a
fresh sample from the same model, which is the quickest way to see how
much of the algorithm's behaviour is a property of the model and how
much is an accident of one data set. Changing the data, the sample size
or any starting value resets the display to iteration 0.

Upper panel: the two weighted component densities, with each observation
coloured by its current responsibility \(\hat p_i^{(t)}\) and arrows
showing where the M-step moves the two means. Lower panel: the observed
log-likelihood, which must never decrease.

\begin{Shaded}
\begin{Highlighting}[]
\NormalTok{\#| label: em{-}mixture{-}app}
\NormalTok{\#| standalone: true}
\NormalTok{\#| viewerHeight: 760}

\NormalTok{library(shiny)}

\NormalTok{\#\#\# {-}{-}{-}{-}{-}{-}{-}{-}{-}{-} model functions {-}{-}{-}{-}{-}{-}{-}{-}{-}{-}{-}{-}{-}{-}{-}{-}{-}{-}{-}{-}{-}{-}{-}{-}{-}{-}{-}{-}{-}{-}{-}{-}{-}{-}{-}{-}{-}{-}{-}{-}{-}{-}{-}{-}{-}{-}{-}{-}{-}}
\NormalTok{log\_sum\_exp \textless{}{-} function(lx) \{ mx \textless{}{-} max(lx); mx + log(sum(exp(lx {-} mx))) \}}

\NormalTok{log\_like\_obs \textless{}{-} function(theta, Y) \{}
\NormalTok{    A \textless{}{-} cbind(log(theta[1])     + dnorm(Y, theta[2], 1, log = TRUE),}
\NormalTok{               log(1 {-} theta[1]) + dnorm(Y, theta[3], 1, log = TRUE))}
\NormalTok{    sum(apply(A, 1, log\_sum\_exp))}
\NormalTok{\}}

\NormalTok{e\_step \textless{}{-} function(theta, Y) \{}
\NormalTok{    l1 \textless{}{-} dnorm(Y, theta[2], 1, log = TRUE) + log(theta[1])}
\NormalTok{    l0 \textless{}{-} dnorm(Y, theta[3], 1, log = TRUE) + log(1 {-} theta[1])}
\NormalTok{    exp(l1 {-} apply(cbind(l1, l0), 1, log\_sum\_exp))}
\NormalTok{\}}

\NormalTok{m\_step \textless{}{-} function(w, Y) \{}
\NormalTok{    c(mean(w), sum(Y * w) / sum(w), sum(Y * (1 {-} w)) / sum(1 {-} w))}
\NormalTok{\}}

\NormalTok{em\_path \textless{}{-} function(theta0, Y, no\_iters) \{}
\NormalTok{    out \textless{}{-} matrix(NA\_real\_, no\_iters + 1, 4,}
\NormalTok{                  dimnames = list(NULL, c("p", "mu1", "mu0", "log\_lik")))}
\NormalTok{    th \textless{}{-} theta0}
\NormalTok{    out[1, ] \textless{}{-} c(th, log\_like\_obs(th, Y))}
\NormalTok{    for (i in seq\_len(no\_iters)) \{}
\NormalTok{        th \textless{}{-} m\_step(e\_step(th, Y), Y)}
\NormalTok{        if (any(!is.finite(th))) break}
\NormalTok{        out[i + 1, ] \textless{}{-} c(th, log\_like\_obs(th, Y))}
\NormalTok{    \}}
\NormalTok{    out}
\NormalTok{\}}

\NormalTok{\#\#\# {-}{-}{-}{-}{-}{-}{-}{-}{-}{-} preset starting values {-}{-}{-}{-}{-}{-}{-}{-}{-}{-}{-}{-}{-}{-}{-}{-}{-}{-}{-}{-}{-}{-}{-}{-}{-}{-}{-}{-}{-}{-}{-}{-}{-}{-}{-}{-}{-}{-}{-}{-}{-}}
\NormalTok{presets \textless{}{-} list(}
\NormalTok{    "near the truth: (0.5, 0, 3)"                 = c(0.5,  0.0,  3.0),}
\NormalTok{    "far outside the data: (0.5, {-}10, 5)"         = c(0.5, {-}10.0, 5.0),}
\NormalTok{    "both means too high: (0.5, 5, 10)"           = c(0.5,  5.0, 10.0),}
\NormalTok{    "means nearly equal: (0.5, 0.1, 0.2)"         = c(0.5,  0.1,  0.2),}
\NormalTok{    "label{-}switched: (0.5, 3, 0)"                 = c(0.5,  3.0,  0.0),}
\NormalTok{    "random labels"                               = NA,}
\NormalTok{    "manual entry"                                = NULL}
\NormalTok{)}

\NormalTok{datasets \textless{}{-} list(}
\NormalTok{    "well separated: p = 0.3, mu = (3, 0)"  = c(0.3, 3, 0),}
\NormalTok{    "widely separated: p = 0.3, mu = (5, 0)" = c(0.3, 5, 0),}
\NormalTok{    "overlapping: p = 0.3, mu = (1.5, 0)"    = c(0.3, 1.5, 0),}
\NormalTok{    "balanced and overlapping: p = 0.5, mu = (1.5, 0)" = c(0.5, 1.5, 0)}
\NormalTok{)}

\NormalTok{\#\#\# {-}{-}{-}{-}{-}{-}{-}{-}{-}{-} UI {-}{-}{-}{-}{-}{-}{-}{-}{-}{-}{-}{-}{-}{-}{-}{-}{-}{-}{-}{-}{-}{-}{-}{-}{-}{-}{-}{-}{-}{-}{-}{-}{-}{-}{-}{-}{-}{-}{-}{-}{-}{-}{-}{-}{-}{-}{-}{-}{-}{-}{-}{-}{-}{-}{-}{-}{-}{-}{-}{-}{-}{-}{-}}
\NormalTok{ui \textless{}{-} fluidPage(}
\NormalTok{    sidebarLayout(}
\NormalTok{        sidebarPanel(}
\NormalTok{            width = 4,}
\NormalTok{            h5("Data"),}
\NormalTok{            selectInput("dataset", NULL, choices = names(datasets)),}
\NormalTok{            sliderInput("n", "sample size", 50, 400, 200, step = 50),}
\NormalTok{            actionButton("newdata", "Regenerate data", width = "100\%"),}
\NormalTok{            tags$hr(),}
\NormalTok{            h5("Starting value"),}
\NormalTok{            selectInput("preset", NULL, choices = names(presets),}
\NormalTok{                        selected = "far outside the data: (0.5, {-}10, 5)"),}
\NormalTok{            fluidRow(}
\NormalTok{                column(4, numericInput("p0",   "p",       0.5, min = 0.01,}
\NormalTok{                                       max = 0.99, step = 0.05)),}
\NormalTok{                column(4, numericInput("mu10", "mu1",   {-}10,   step = 0.5)),}
\NormalTok{                column(4, numericInput("mu00", "mu0",     5,   step = 0.5))}
\NormalTok{            ),}
\NormalTok{            helpText("Selecting a preset fills the three boxes. The boxes are ",}
\NormalTok{                     "always live: edit them to use your own values."),}
\NormalTok{            tags$hr(),}
\NormalTok{            sliderInput("maxit", "iterations to run", 5, 60, 30, step = 5),}
\NormalTok{            h5(textOutput("stepLabel", inline = TRUE)),}
\NormalTok{            fluidRow(}
\NormalTok{                column(4, actionButton("prv",  "\textbackslash{}u25c0 Prev", width = "100\%")),}
\NormalTok{                column(4, actionButton("nxt",  "Next \textbackslash{}u25b6", width = "100\%")),}
\NormalTok{                column(4, actionButton("play", "Play",         width = "100\%"))}
\NormalTok{            ),}
\NormalTok{            tags$br(),}
\NormalTok{            tableOutput("tab")}
\NormalTok{        ),}
\NormalTok{        mainPanel(}
\NormalTok{            width = 8,}
\NormalTok{            plotOutput("fit",   height = "400px"),}
\NormalTok{            plotOutput("trace", height = "260px")}
\NormalTok{        )}
\NormalTok{    )}
\NormalTok{)}

\NormalTok{\#\#\# {-}{-}{-}{-}{-}{-}{-}{-}{-}{-} server {-}{-}{-}{-}{-}{-}{-}{-}{-}{-}{-}{-}{-}{-}{-}{-}{-}{-}{-}{-}{-}{-}{-}{-}{-}{-}{-}{-}{-}{-}{-}{-}{-}{-}{-}{-}{-}{-}{-}{-}{-}{-}{-}{-}{-}{-}{-}{-}{-}{-}{-}{-}{-}{-}{-}{-}{-}{-}{-}}
\NormalTok{server \textless{}{-} function(input, output, session) \{}

\NormalTok{    seed \textless{}{-} reactiveVal(812)}
\NormalTok{    observeEvent(input$newdata, seed(seed() + 1))}

\NormalTok{    Y \textless{}{-} reactive(\{}
\NormalTok{        s \textless{}{-} datasets[[input$dataset]]}
\NormalTok{        set.seed(seed())}
\NormalTok{        Z \textless{}{-} rbinom(input$n, 1, s[1])}
\NormalTok{        rnorm(input$n, mean = ifelse(Z == 1, s[2], s[3]), sd = 1)}
\NormalTok{    \})}

\NormalTok{    \#\# preset {-}\textgreater{} fill the three boxes}
\NormalTok{    observeEvent(list(input$preset, input$dataset, input$n, seed()), \{}
\NormalTok{        pr \textless{}{-} presets[[input$preset]]}
\NormalTok{        if (is.null(pr)) return(invisible(NULL))          \# manual entry}
\NormalTok{        if (length(pr) == 1 \&\& is.na(pr)) \{               \# random labels}
\NormalTok{            y \textless{}{-} Y()}
\NormalTok{            lab \textless{}{-} rbinom(length(y), 1, 0.5)}
\NormalTok{            pr \textless{}{-} c(0.5, mean(y[lab == 1]), mean(y[lab == 0]))}
\NormalTok{        \}}
\NormalTok{        updateNumericInput(session, "p0",   value = round(pr[1], 3))}
\NormalTok{        updateNumericInput(session, "mu10", value = round(pr[2], 3))}
\NormalTok{        updateNumericInput(session, "mu00", value = round(pr[3], 3))}
\NormalTok{    \}, ignoreInit = FALSE)}

\NormalTok{    \#\# {-}{-}{-}{-} iteration state driven by the three buttons {-}{-}{-}{-}{-}{-}{-}{-}{-}{-}{-}{-}{-}{-}{-}{-}{-}{-}{-}{-}{-}{-}}
\NormalTok{    t\_cur   \textless{}{-} reactiveVal(0)      \# iteration currently displayed}
\NormalTok{    playing \textless{}{-} reactiveVal(FALSE)}

\NormalTok{    stop\_play \textless{}{-} function() \{}
\NormalTok{        playing(FALSE)}
\NormalTok{        updateActionButton(session, "play", label = "Play")}
\NormalTok{    \}}

\NormalTok{    \#\# any change to the data or the starting value restarts the fit}
\NormalTok{    observeEvent(list(input$dataset, input$n, seed(), input$maxit,}
\NormalTok{                      input$p0, input$mu10, input$mu00), \{}
\NormalTok{        t\_cur(0)}
\NormalTok{        stop\_play()}
\NormalTok{    \})}

\NormalTok{    observeEvent(input$nxt, \{}
\NormalTok{        stop\_play()}
\NormalTok{        t\_cur(min(t\_cur() + 1, input$maxit))}
\NormalTok{    \})}

\NormalTok{    observeEvent(input$prv, \{}
\NormalTok{        stop\_play()}
\NormalTok{        t\_cur(max(t\_cur() {-} 1, 0))}
\NormalTok{    \})}

\NormalTok{    observeEvent(input$play, \{}
\NormalTok{        if (playing()) \{}
\NormalTok{            stop\_play()}
\NormalTok{        \} else \{}
\NormalTok{            if (t\_cur() \textgreater{}= input$maxit) t\_cur(0)   \# replay from the start}
\NormalTok{            playing(TRUE)}
\NormalTok{            updateActionButton(session, "play", label = "Pause")}
\NormalTok{        \}}
\NormalTok{    \})}

\NormalTok{    \#\# the timer: advances one iteration per tick while playing}
\NormalTok{    observe(\{}
\NormalTok{        if (!playing()) return(NULL)}
\NormalTok{        invalidateLater(700, session)}
\NormalTok{        isolate(\{}
\NormalTok{            if (t\_cur() \textgreater{}= input$maxit) stop\_play() else t\_cur(t\_cur() + 1)}
\NormalTok{        \})}
\NormalTok{    \})}

\NormalTok{    output$stepLabel \textless{}{-} renderText(}
\NormalTok{        sprintf("Iteration \%d of \%d", t\_cur(), input$maxit)}
\NormalTok{    )}

\NormalTok{    theta0 \textless{}{-} reactive(\{}
\NormalTok{        p \textless{}{-} input$p0}
\NormalTok{        validate(need(is.finite(p) \&\& p \textgreater{} 0 \&\& p \textless{} 1, "p must lie strictly between 0 and 1"),}
\NormalTok{                 need(is.finite(input$mu10) \&\& is.finite(input$mu00),}
\NormalTok{                      "both means must be numeric"))}
\NormalTok{        c(p, input$mu10, input$mu00)}
\NormalTok{    \})}

\NormalTok{    path \textless{}{-} reactive(em\_path(theta0(), Y(), input$maxit))}

\NormalTok{    output$fit \textless{}{-} renderPlot(\{}
\NormalTok{        y \textless{}{-} Y(); P \textless{}{-} path()}
\NormalTok{        k  \textless{}{-} min(t\_cur() + 1, nrow(P))}
\NormalTok{        th \textless{}{-} P[k, 1:3]}
\NormalTok{        validate(need(all(is.finite(th)), "the iteration broke down numerically"))}
\NormalTok{        w  \textless{}{-} e\_step(th, y)}
\NormalTok{        cw \textless{}{-} cm.colors(100)}

\NormalTok{        xr \textless{}{-} range(y, th[2:3]) + c({-}2, 2)}
\NormalTok{        xg \textless{}{-} seq(xr[1], xr[2], length.out = 400)}
\NormalTok{        d1 \textless{}{-} dnorm(xg, th[2]) * th[1]}
\NormalTok{        d0 \textless{}{-} dnorm(xg, th[3]) * (1 {-} th[1])}

\NormalTok{        par(mar = c(4, 4, 3, 1))}
\NormalTok{        plot(xg, d1, type = "l", lwd = 2, col = cw[100],}
\NormalTok{             ylim = c(0, max(0.42, d1, d0)), xlab = "y", ylab = "weighted density",}
\NormalTok{             main = sprintf("step \%d:  p = \%.3f,  mu1 = \%.3f,  mu0 = \%.3f,  loglik = \%.2f",}
\NormalTok{                            k {-} 1, th[1], th[2], th[3], P[k, 4]))}
\NormalTok{        lines(xg, d0, lwd = 2, col = cw[1])}
\NormalTok{        points(y, rep(0, length(y)), col = cw[pmax(1, ceiling(w * 100))], pch = 19,}
\NormalTok{               cex = 0.7)}
\NormalTok{        if (k \textless{} nrow(P) \&\& all(is.finite(P[k + 1, 1:3]))) \{}
\NormalTok{            arrows(th[2], 0.02, P[k + 1, 2], 0.02, length = 0.08, lwd = 2, col = cw[100])}
\NormalTok{            arrows(th[3], 0.02, P[k + 1, 3], 0.02, length = 0.08, lwd = 2, col = cw[1])}
\NormalTok{        \}}
\NormalTok{        legend("topright", bty = "n", lwd = 2, col = c(cw[100], cw[1]),}
\NormalTok{               legend = c(expression(hat(p) * phi(y * "|" * hat(mu)[1])),}
\NormalTok{                          expression((1 {-} hat(p)) * phi(y * "|" * hat(mu)[0]))))}
\NormalTok{    \})}

\NormalTok{    output$trace \textless{}{-} renderPlot(\{}
\NormalTok{        P \textless{}{-} path(); k \textless{}{-} min(t\_cur() + 1, nrow(P))}
\NormalTok{        par(mfrow = c(1, 2), mar = c(4, 4, 2, 1))}
\NormalTok{        it \textless{}{-} 0:(nrow(P) {-} 1)}
\NormalTok{        plot(it, P[, 4], type = "n", xlab = "iteration",}
\NormalTok{             ylab = "observed log{-}likelihood", main = "must never decrease")}
\NormalTok{        lines(it[1:k], P[1:k, 4], type = "b", pch = 19, col = "firebrick")}
\NormalTok{        yl \textless{}{-} range(P[, 2:3], na.rm = TRUE)}
\NormalTok{        plot(it, P[, 2], type = "n", ylim = yl, xlab = "iteration",}
\NormalTok{             ylab = "component means", main = "parameter paths")}
\NormalTok{        lines(it[1:k], P[1:k, 2], type = "b", pch = 19, col = "firebrick")}
\NormalTok{        lines(it[1:k], P[1:k, 3], type = "b", pch = 19, col = "steelblue")}
\NormalTok{        legend("right", bty = "n", lwd = 2, col = c("firebrick", "steelblue"),}
\NormalTok{               legend = c(expression(hat(mu)[1]), expression(hat(mu)[0])))}
\NormalTok{    \})}

\NormalTok{    output$tab \textless{}{-} renderTable(\{}
\NormalTok{        P \textless{}{-} path(); k \textless{}{-} min(t\_cur() + 1, nrow(P))}
\NormalTok{        d \textless{}{-} as.data.frame(P[seq\_len(k), , drop = FALSE])}
\NormalTok{        d$iter \textless{}{-} seq\_len(k) {-} 1}
\NormalTok{        tail(d[, c("iter", "p", "mu1", "mu0", "log\_lik")], 5)}
\NormalTok{    \}, digits = 4, rownames = FALSE)}
\NormalTok{\}}

\NormalTok{shinyApp(ui, server)}
\end{Highlighting}
\end{Shaded}

Four behaviours are worth producing deliberately.

\textbf{Slow drift from a distant start.} With \((0.5, -10, 5)\) both
means begin far outside the data. The early iterations move them very
little, because every observation is assigned almost the same
responsibility; progress accelerates only once one mean reaches the
data.

\textbf{Label switching.} The starts \((0.5, 0, 3)\) and \((0.5, 3, 0)\)
converge to the same fitted mixture with the roles of the two components
exchanged. The likelihood is invariant to this relabelling, so the
fitted density is identical while the parameter estimates differ --- a
genuine identifiability problem for the model, not a defect of the
algorithm.

\textbf{Stalling at a symmetric point.} With \(\mu_1^{(0)} = 0.1\) and
\(\mu_0^{(0)} = 0.2\) the two components are nearly indistinguishable,
every responsibility is close to \(\hat p^{(0)}\), and the algorithm
creeps away from the one-component solution very slowly. That
configuration is close to a saddle point of \(\ell_{\text{obs}}\).

\textbf{Separation and speed.} Switch the data set from
\(\mu = (1.5, 0)\) to \(\mu = (5, 0)\). Well-separated components make
the responsibilities nearly \(0\) or \(1\), the E-step becomes almost a
hard classification, and convergence takes a handful of iterations
rather than dozens. This is the practical face of the general result
that EM's linear convergence rate is governed by the fraction of missing
information.

\subsection{Comparison with general-purpose
optimizers}\label{comparison-with-general-purpose-optimizers}

EM is not the only way to maximize the observed likelihood. Applying a
general-purpose optimizer directly requires attention to the parameter
space: \(p\) must lie in \((0,1)\), so we reparameterize with the logit.

\begin{Shaded}
\begin{Highlighting}[]
\NormalTok{neg\_loglike\_obs }\OtherTok{\textless{}{-}} \ControlFlowTok{function}\NormalTok{(theta, Y) }\SpecialCharTok{{-}}\FunctionTok{log\_like\_obs}\NormalTok{(theta, Y)}

\DocumentationTok{\#\#\# transformed parameterization: ttheta[1] = logit(p)}
\NormalTok{neg\_loglike\_obs\_transf }\OtherTok{\textless{}{-}} \ControlFlowTok{function}\NormalTok{(ttheta, Y) \{}
\NormalTok{    theta }\OtherTok{\textless{}{-}}\NormalTok{ ttheta}
\NormalTok{    theta[}\DecValTok{1}\NormalTok{] }\OtherTok{\textless{}{-}} \DecValTok{1} \SpecialCharTok{/}\NormalTok{ (}\DecValTok{1} \SpecialCharTok{+} \FunctionTok{exp}\NormalTok{(}\SpecialCharTok{{-}}\NormalTok{ttheta[}\DecValTok{1}\NormalTok{]))}
    \FunctionTok{neg\_loglike\_obs}\NormalTok{(theta, Y)}
\NormalTok{\}}

\NormalTok{lp0 }\OtherTok{\textless{}{-}} \FunctionTok{log}\NormalTok{(}\FloatTok{0.3} \SpecialCharTok{/} \FloatTok{0.7}\NormalTok{)}

\DocumentationTok{\#\#\# Newton{-}type, untransformed: p can leave (0, 1)}
\FunctionTok{try}\NormalTok{(}\FunctionTok{nlm}\NormalTok{(neg\_loglike\_obs, }\AttributeTok{p =} \FunctionTok{c}\NormalTok{(}\FloatTok{0.3}\NormalTok{, }\DecValTok{0}\NormalTok{, }\DecValTok{3}\NormalTok{), }\AttributeTok{Y =}\NormalTok{ data)}\SpecialCharTok{$}\NormalTok{estimate)}
\end{Highlighting}
\end{Shaded}

\begin{verbatim}
[1] 0.2968473 0.1131871 2.9203484
\end{verbatim}

\begin{Shaded}
\begin{Highlighting}[]
\FunctionTok{try}\NormalTok{(}\FunctionTok{nlm}\NormalTok{(neg\_loglike\_obs, }\AttributeTok{p =} \FunctionTok{c}\NormalTok{(}\FloatTok{0.5}\NormalTok{, }\SpecialCharTok{{-}}\DecValTok{10}\NormalTok{, }\DecValTok{5}\NormalTok{), }\AttributeTok{Y =}\NormalTok{ data)}\SpecialCharTok{$}\NormalTok{estimate)}
\end{Highlighting}
\end{Shaded}

\begin{verbatim}
[1]  4.872163e-06 -1.000004e+01  4.038012e+00
\end{verbatim}

\begin{Shaded}
\begin{Highlighting}[]
\DocumentationTok{\#\#\# Newton{-}type, transformed}
\FunctionTok{try}\NormalTok{(}\FunctionTok{nlm}\NormalTok{(neg\_loglike\_obs\_transf, }\AttributeTok{p =} \FunctionTok{c}\NormalTok{(lp0, }\DecValTok{0}\NormalTok{, }\DecValTok{3}\NormalTok{), }\AttributeTok{Y =}\NormalTok{ data)}\SpecialCharTok{$}\NormalTok{estimate)}
\end{Highlighting}
\end{Shaded}

\begin{verbatim}
[1] -0.8623555  0.1131868  2.9203483
\end{verbatim}

\begin{Shaded}
\begin{Highlighting}[]
\FunctionTok{try}\NormalTok{(}\FunctionTok{nlm}\NormalTok{(neg\_loglike\_obs\_transf, }\AttributeTok{p =} \FunctionTok{c}\NormalTok{(}\DecValTok{0}\NormalTok{, }\SpecialCharTok{{-}}\DecValTok{10}\NormalTok{, }\DecValTok{5}\NormalTok{), }\AttributeTok{Y =}\NormalTok{ data)}\SpecialCharTok{$}\NormalTok{estimate)}
\end{Highlighting}
\end{Shaded}

\begin{verbatim}
[1] -16.605894  -9.996458   2.087049
\end{verbatim}

\begin{Shaded}
\begin{Highlighting}[]
\DocumentationTok{\#\#\# Nelder{-}Mead}
\FunctionTok{try}\NormalTok{(}\FunctionTok{optim}\NormalTok{(}\FunctionTok{c}\NormalTok{(lp0, }\DecValTok{0}\NormalTok{, }\DecValTok{3}\NormalTok{), neg\_loglike\_obs\_transf, }\AttributeTok{Y =}\NormalTok{ data, }\AttributeTok{method =} \StringTok{"Nelder{-}Mead"}\NormalTok{)}\SpecialCharTok{$}\NormalTok{par)}
\end{Highlighting}
\end{Shaded}

\begin{verbatim}
[1] -0.8623197  0.1132163  2.9203325
\end{verbatim}

\begin{Shaded}
\begin{Highlighting}[]
\FunctionTok{try}\NormalTok{(}\FunctionTok{optim}\NormalTok{(}\FunctionTok{c}\NormalTok{(}\DecValTok{0}\NormalTok{, }\SpecialCharTok{{-}}\DecValTok{10}\NormalTok{, }\DecValTok{5}\NormalTok{), neg\_loglike\_obs\_transf, }\AttributeTok{Y =}\NormalTok{ data, }\AttributeTok{method =} \StringTok{"Nelder{-}Mead"}\NormalTok{)}\SpecialCharTok{$}\NormalTok{par)}
\end{Highlighting}
\end{Shaded}

\begin{verbatim}
[1] -0.8622804  0.1131101  2.9206253
\end{verbatim}

\begin{Shaded}
\begin{Highlighting}[]
\DocumentationTok{\#\#\# Conjugate gradient}
\FunctionTok{try}\NormalTok{(}\FunctionTok{optim}\NormalTok{(}\FunctionTok{c}\NormalTok{(}\DecValTok{0}\NormalTok{, }\DecValTok{0}\NormalTok{, }\DecValTok{3}\NormalTok{), neg\_loglike\_obs\_transf, }\AttributeTok{Y =}\NormalTok{ data, }\AttributeTok{method =} \StringTok{"CG"}\NormalTok{)}\SpecialCharTok{$}\NormalTok{par)}
\end{Highlighting}
\end{Shaded}

\begin{verbatim}
[1] -0.8623543  0.1131880  2.9203502
\end{verbatim}

\begin{Shaded}
\begin{Highlighting}[]
\FunctionTok{try}\NormalTok{(}\FunctionTok{optim}\NormalTok{(}\FunctionTok{c}\NormalTok{(}\DecValTok{0}\NormalTok{, }\SpecialCharTok{{-}}\DecValTok{10}\NormalTok{, }\DecValTok{5}\NormalTok{), neg\_loglike\_obs\_transf, }\AttributeTok{Y =}\NormalTok{ data, }\AttributeTok{method =} \StringTok{"CG"}\NormalTok{)}\SpecialCharTok{$}\NormalTok{par)}
\end{Highlighting}
\end{Shaded}

\begin{verbatim}
[1] -6.090903 -9.998146  2.080878
\end{verbatim}

Applied to the observed log-likelihood, these algorithms are more
fragile than EM when the initial values are poorly chosen: the
untransformed runs can step outside \(p \in (0,1)\) and produce
\texttt{NaN}, and the gradient-based methods can stall in the flat
region that a distant start creates. Nelder--Mead is the most stable of
the three here, at the cost of many more function evaluations.

\begin{longtable}[]{@{}
  >{\raggedright\arraybackslash}p{(\linewidth - 4\tabcolsep) * \real{0.3333}}
  >{\raggedright\arraybackslash}p{(\linewidth - 4\tabcolsep) * \real{0.3333}}
  >{\raggedright\arraybackslash}p{(\linewidth - 4\tabcolsep) * \real{0.3333}}@{}}
\caption{EM compared with direct numerical
maximization}\label{tbl-em-vs-optim}\tabularnewline
\toprule\noalign{}
\begin{minipage}[b]{\linewidth}\raggedright
\end{minipage} & \begin{minipage}[b]{\linewidth}\raggedright
EM
\end{minipage} & \begin{minipage}[b]{\linewidth}\raggedright
General optimizer on \(\ell_{\text{obs}}\)
\end{minipage} \\
\midrule\noalign{}
\endfirsthead
\toprule\noalign{}
\begin{minipage}[b]{\linewidth}\raggedright
\end{minipage} & \begin{minipage}[b]{\linewidth}\raggedright
EM
\end{minipage} & \begin{minipage}[b]{\linewidth}\raggedright
General optimizer on \(\ell_{\text{obs}}\)
\end{minipage} \\
\midrule\noalign{}
\endhead
\bottomrule\noalign{}
\endlastfoot
Constraints & Respected automatically by the M-step & Require
reparameterization \\
Monotonicity & Guaranteed; a decrease signals a coding error & None; a
bad step can be accepted \\
Derivatives & None needed & Gradient helpful, Hessian better \\
Convergence rate & Linear, rate set by the missing information &
Superlinear or quadratic near the optimum \\
Standard errors & Not produced; need Louis' method or a bootstrap &
Hessian available directly \\
Behaviour from poor starts & Slow but stable & Can fail outright \\
\end{longtable}

The last row is the practical trade. EM rarely fails and rarely needs
tuning, which is why it dominates in mixture and latent-variable
modelling; it is also slow near the optimum and gives no standard
errors, which is why a hybrid --- EM to get close, then a Newton step to
polish and to obtain the Hessian --- is common in production code.

\section{Example 2: Censored Poisson Data}\label{sec-poisson}

\subsection{Truncated distributions in
brief}\label{truncated-distributions-in-brief}

Let \(X \sim f(x)\) and \(Y = I(X > c)\). Then
\(Y \sim \text{Bern}\{P(X > c)\}\), and the conditional distribution of
\(X\) given \(Y\) is \(f\) restricted to one side of \(c\) and
renormalized:

\[
f(x \mid Y = 1) = \frac{f(x)\, I(x > c)}{\int_c^{\infty} f(t)\,dt},
\qquad
f(x \mid Y = 0) = \frac{f(x)\, I(x \le c)}{\int_{-\infty}^{c} f(t)\,dt} .
\]

For discrete \(X\) the integrals are sums. This is exactly the
distribution the E-step needs when an observation is known only to lie
on one side of \(c\).

\begin{Shaded}
\begin{Highlighting}[]
\FunctionTok{par}\NormalTok{(}\AttributeTok{mfrow =} \FunctionTok{c}\NormalTok{(}\DecValTok{1}\NormalTok{, }\DecValTok{2}\NormalTok{), }\AttributeTok{mar =} \FunctionTok{c}\NormalTok{(}\DecValTok{4}\NormalTok{, }\DecValTok{4}\NormalTok{, }\DecValTok{2}\NormalTok{, }\DecValTok{1}\NormalTok{)); cc }\OtherTok{\textless{}{-}} \DecValTok{1}
\FunctionTok{curve}\NormalTok{(dnorm, }\SpecialCharTok{{-}}\FloatTok{3.5}\NormalTok{, }\FloatTok{3.5}\NormalTok{, }\AttributeTok{lwd =} \DecValTok{3}\NormalTok{, }\AttributeTok{xlab =} \StringTok{"x"}\NormalTok{, }\AttributeTok{ylab =} \StringTok{"density"}\NormalTok{,}
      \AttributeTok{main =} \StringTok{"f(x) with cut point c"}\NormalTok{)}
\NormalTok{xs }\OtherTok{\textless{}{-}} \FunctionTok{seq}\NormalTok{(cc, }\FloatTok{3.5}\NormalTok{, }\AttributeTok{length.out =} \DecValTok{100}\NormalTok{)}
\FunctionTok{polygon}\NormalTok{(}\FunctionTok{c}\NormalTok{(xs, }\FunctionTok{rev}\NormalTok{(xs)), }\FunctionTok{c}\NormalTok{(}\FunctionTok{dnorm}\NormalTok{(xs), }\FunctionTok{rep}\NormalTok{(}\DecValTok{0}\NormalTok{, }\DecValTok{100}\NormalTok{)),}
        \AttributeTok{col =} \FunctionTok{adjustcolor}\NormalTok{(}\StringTok{"firebrick"}\NormalTok{, }\FloatTok{0.3}\NormalTok{), }\AttributeTok{border =} \ConstantTok{NA}\NormalTok{)}
\NormalTok{xs }\OtherTok{\textless{}{-}} \FunctionTok{seq}\NormalTok{(}\SpecialCharTok{{-}}\FloatTok{3.5}\NormalTok{, cc, }\AttributeTok{length.out =} \DecValTok{100}\NormalTok{)}
\FunctionTok{polygon}\NormalTok{(}\FunctionTok{c}\NormalTok{(xs, }\FunctionTok{rev}\NormalTok{(xs)), }\FunctionTok{c}\NormalTok{(}\FunctionTok{dnorm}\NormalTok{(xs), }\FunctionTok{rep}\NormalTok{(}\DecValTok{0}\NormalTok{, }\DecValTok{100}\NormalTok{)),}
        \AttributeTok{col =} \FunctionTok{adjustcolor}\NormalTok{(}\StringTok{"steelblue"}\NormalTok{, }\FloatTok{0.3}\NormalTok{), }\AttributeTok{border =} \ConstantTok{NA}\NormalTok{)}
\FunctionTok{abline}\NormalTok{(}\AttributeTok{v =}\NormalTok{ cc, }\AttributeTok{lty =} \DecValTok{2}\NormalTok{); }\FunctionTok{text}\NormalTok{(cc, }\FloatTok{0.38}\NormalTok{, }\StringTok{"c"}\NormalTok{, }\AttributeTok{pos =} \DecValTok{4}\NormalTok{)}
\FunctionTok{text}\NormalTok{(}\FloatTok{1.8}\NormalTok{, }\FloatTok{0.05}\NormalTok{, }\StringTok{"P(X \textgreater{} c)"}\NormalTok{, }\AttributeTok{col =} \StringTok{"firebrick"}\NormalTok{)}
\FunctionTok{text}\NormalTok{(}\SpecialCharTok{{-}}\DecValTok{1}\NormalTok{, }\FloatTok{0.1}\NormalTok{, }\FunctionTok{expression}\NormalTok{(}\FunctionTok{P}\NormalTok{(X }\SpecialCharTok{\textless{}=}\NormalTok{ c)), }\AttributeTok{col =} \StringTok{"steelblue"}\NormalTok{)}
\FunctionTok{curve}\NormalTok{(}\FunctionTok{dnorm}\NormalTok{(x) }\SpecialCharTok{*}\NormalTok{ (x }\SpecialCharTok{\textgreater{}}\NormalTok{ cc) }\SpecialCharTok{/} \FunctionTok{pnorm}\NormalTok{(cc, }\AttributeTok{lower.tail =} \ConstantTok{FALSE}\NormalTok{), }\SpecialCharTok{{-}}\FloatTok{3.5}\NormalTok{, }\FloatTok{3.5}\NormalTok{, }\AttributeTok{lwd =} \DecValTok{3}\NormalTok{,}
      \AttributeTok{col =} \StringTok{"firebrick"}\NormalTok{, }\AttributeTok{n =} \DecValTok{1001}\NormalTok{, }\AttributeTok{xlab =} \StringTok{"x"}\NormalTok{, }\AttributeTok{ylab =} \StringTok{"density"}\NormalTok{,}
      \AttributeTok{main =} \StringTok{"conditional densities"}\NormalTok{, }\AttributeTok{ylim =} \FunctionTok{c}\NormalTok{(}\DecValTok{0}\NormalTok{, }\FloatTok{1.6}\NormalTok{))}
\FunctionTok{curve}\NormalTok{(}\FunctionTok{dnorm}\NormalTok{(x) }\SpecialCharTok{*}\NormalTok{ (x }\SpecialCharTok{\textless{}=}\NormalTok{ cc) }\SpecialCharTok{/} \FunctionTok{pnorm}\NormalTok{(cc), }\AttributeTok{add =} \ConstantTok{TRUE}\NormalTok{, }\AttributeTok{lwd =} \DecValTok{3}\NormalTok{, }\AttributeTok{col =} \StringTok{"steelblue"}\NormalTok{, }\AttributeTok{n =} \DecValTok{1001}\NormalTok{)}
\FunctionTok{abline}\NormalTok{(}\AttributeTok{v =}\NormalTok{ cc, }\AttributeTok{lty =} \DecValTok{2}\NormalTok{)}
\FunctionTok{legend}\NormalTok{(}\StringTok{"topleft"}\NormalTok{, }\AttributeTok{bty =} \StringTok{"n"}\NormalTok{, }\AttributeTok{lwd =} \DecValTok{3}\NormalTok{, }\AttributeTok{col =} \FunctionTok{c}\NormalTok{(}\StringTok{"steelblue"}\NormalTok{, }\StringTok{"firebrick"}\NormalTok{),}
       \AttributeTok{legend =} \FunctionTok{c}\NormalTok{(}\StringTok{"f(x | Y = 0)"}\NormalTok{, }\StringTok{"f(x | Y = 1)"}\NormalTok{))}
\end{Highlighting}
\end{Shaded}

\begin{figure}[H]

\centering{

\pandocbounded{\includegraphics[keepaspectratio]{chapters/unit07_em_files/figure-pdf/fig-truncation-1.pdf}}

}

\caption{\label{fig-truncation}A density split at a cut point, and the
two renormalized conditional densities.}

\end{figure}%

\subsection{Data and likelihoods}\label{data-and-likelihoods}

Let \(y_i \overset{iid}{\sim} \text{Pois}(\lambda)\) for
\(i = 1, \ldots, n\). The first \(m\) values are observed exactly; for
\(i > m\) we know only that \(y_i < 2\), recorded as
\(x_i = I(y_i < 2) = 1\).

\textbf{Observed data}: \(y_1, \ldots, y_m\) together with
\(x_{m+1} = \cdots = x_n = 1\). \textbf{Missing}:
\(y_{m+1}, \ldots, y_n\).

\textbf{Complete likelihood}

\[
\prod_{i=1}^{n} f^{\text{Pois}}(y_i \mid \lambda) \cdot \prod_{i=m+1}^{n} I(y_i < 2) .
\]

\textbf{Observed likelihood}

\[
\prod_{i=1}^{m} f^{\text{Pois}}(y_i \mid \lambda) \cdot \prod_{i=m+1}^{n} P(y_i < 2 \mid \lambda),
\qquad P(y_i < 2 \mid \lambda) = e^{-\lambda}(1 + \lambda) .
\]

Taking logarithms and dropping constants,

\[
\ell_{\text{obs}}(\lambda) = m \bar y_{1:m} \log \lambda - n \lambda + (n - m)\log(1 + \lambda) .
\]

The second factor makes the MLE nonlinear in \(\lambda\); EM sidesteps
it.

\subsection{E-step}\label{e-step}

\[
\ell_{\text{comp}}(\lambda) = -n\lambda + (\log \lambda) \sum_{i=1}^{m} y_i
+ (\log \lambda) \underbrace{\sum_{i=m+1}^{n} y_i}_{\text{missing}}
- \sum_{i=1}^{n} \log(y_i!) .
\]

For \(i > m\) we have \(y_i \in \{0, 1\}\), and the Poisson pmf
truncated to that set given \(\lambda^{(t)}\) is

\[
P\!\left(y_i = 0 \mid x_i = 1\right) = \frac{1}{1 + \lambda^{(t)}},
\qquad
P\!\left(y_i = 1 \mid x_i = 1\right) = \frac{\lambda^{(t)}}{1 + \lambda^{(t)}} ,
\]

so that

\[
\hat y^{(t)} \equiv E\!\left(y_i \mid x_i = 1, \lambda^{(t)}\right) = \frac{\lambda^{(t)}}{1 + \lambda^{(t)}} .
\]

Since \(y_i! = 1\) for \(y_i \in \{0, 1\}\), the \(\log(y_i!)\) term is
a known constant and no further expectation is needed.

\begin{Shaded}
\begin{Highlighting}[]
\FunctionTok{par}\NormalTok{(}\AttributeTok{mfrow =} \FunctionTok{c}\NormalTok{(}\DecValTok{1}\NormalTok{, }\DecValTok{2}\NormalTok{), }\AttributeTok{mar =} \FunctionTok{c}\NormalTok{(}\DecValTok{4}\NormalTok{, }\DecValTok{4}\NormalTok{, }\DecValTok{2}\NormalTok{, }\DecValTok{1}\NormalTok{)); lt }\OtherTok{\textless{}{-}} \FloatTok{1.5}\NormalTok{; k }\OtherTok{\textless{}{-}} \DecValTok{0}\SpecialCharTok{:}\DecValTok{6}
\FunctionTok{barplot}\NormalTok{(}\FunctionTok{dpois}\NormalTok{(k, lt), }\AttributeTok{names.arg =}\NormalTok{ k,}
        \AttributeTok{col =} \FunctionTok{c}\NormalTok{(}\StringTok{"firebrick"}\NormalTok{, }\StringTok{"firebrick"}\NormalTok{, }\FunctionTok{rep}\NormalTok{(}\StringTok{"gray80"}\NormalTok{, }\DecValTok{5}\NormalTok{)),}
        \AttributeTok{xlab =} \StringTok{"y"}\NormalTok{, }\AttributeTok{ylab =} \StringTok{"probability"}\NormalTok{,}
        \AttributeTok{main =} \FunctionTok{bquote}\NormalTok{(}\StringTok{"Pois("} \SpecialCharTok{*}\NormalTok{ lambda}\SpecialCharTok{\^{}}\NormalTok{(t) }\SpecialCharTok{==}\NormalTok{ .(lt) }\SpecialCharTok{*} \StringTok{")"}\NormalTok{))}
\FunctionTok{barplot}\NormalTok{(}\FunctionTok{c}\NormalTok{(}\DecValTok{1}\NormalTok{, lt) }\SpecialCharTok{/}\NormalTok{ (}\DecValTok{1} \SpecialCharTok{+}\NormalTok{ lt), }\AttributeTok{names.arg =} \DecValTok{0}\SpecialCharTok{:}\DecValTok{1}\NormalTok{, }\AttributeTok{col =} \StringTok{"firebrick"}\NormalTok{,}
        \AttributeTok{xlab =} \StringTok{"y"}\NormalTok{, }\AttributeTok{ylab =} \StringTok{"probability"}\NormalTok{,}
        \AttributeTok{main =} \FunctionTok{expression}\NormalTok{(}\StringTok{"truncated to \{0, 1\}: "} \SpecialCharTok{*} \FunctionTok{P}\NormalTok{(y[i]}\SpecialCharTok{*}\StringTok{"|"}\SpecialCharTok{*}\NormalTok{x[i] }\SpecialCharTok{==} \DecValTok{1}\NormalTok{)))}
\end{Highlighting}
\end{Shaded}

\begin{figure}[H]

\centering{

\pandocbounded{\includegraphics[keepaspectratio]{chapters/unit07_em_files/figure-pdf/fig-poisson-estep-1.pdf}}

}

\caption{\label{fig-poisson-estep}The conditional distribution of a
censored count is the Poisson pmf restricted to \(\{0,1\}\) and
renormalized.}

\end{figure}%

\subsection{M-step and the EM update}\label{m-step-and-the-em-update}

\[
Q\!\left(\lambda \mid \lambda^{(t)}\right)
= -n\lambda + (\log \lambda)\left[\sum_{i=1}^{m} y_i + (n - m)\,\hat y^{(t)}\right] + C^{(t)} ,
\]

\[
\frac{\partial Q}{\partial \lambda}
= -n + \frac{1}{\lambda}\left[\sum_{i=1}^{m} y_i + (n - m)\,\hat y^{(t)}\right] = 0 ,
\]

\[
\Longrightarrow \quad
\lambda^{(t+1)} = \frac{m\, \bar y_{1:m} + (n - m)\dfrac{\lambda^{(t)}}{1 + \lambda^{(t)}}}{n} .
\]

The update is the ordinary Poisson MLE --- the sample mean --- with each
censored value replaced by its conditional expectation. As warned in the
overview, this works \emph{only} because \(\ell_{\text{comp}}\) is
linear in the missing \(y_i\).

\subsection{Implementation}\label{implementation-2}

\begin{Shaded}
\begin{Highlighting}[]
\DocumentationTok{\#\#\# observed log{-}likelihood, up to an additive constant}
\DocumentationTok{\#\#   m    : number of exactly observed counts}
\DocumentationTok{\#\#   ybar : their mean}
\DocumentationTok{\#\#   ncen : number of counts known only to be \textless{} 2}
\NormalTok{loglik\_obs\_pois }\OtherTok{\textless{}{-}} \ControlFlowTok{function}\NormalTok{(m, ybar, ncen, lambda) \{}
\NormalTok{    m }\SpecialCharTok{*}\NormalTok{ ybar }\SpecialCharTok{*} \FunctionTok{log}\NormalTok{(lambda) }\SpecialCharTok{{-}}\NormalTok{ (m }\SpecialCharTok{+}\NormalTok{ ncen) }\SpecialCharTok{*}\NormalTok{ lambda }\SpecialCharTok{+}\NormalTok{ ncen }\SpecialCharTok{*} \FunctionTok{log}\NormalTok{(}\DecValTok{1} \SpecialCharTok{+}\NormalTok{ lambda)}
\NormalTok{\}}

\NormalTok{em\_censored\_poisson }\OtherTok{\textless{}{-}} \ControlFlowTok{function}\NormalTok{(m, ybar, ncen,}
                                \AttributeTok{lambda0 =}\NormalTok{ (m }\SpecialCharTok{*}\NormalTok{ ybar }\SpecialCharTok{+}\NormalTok{ ncen) }\SpecialCharTok{/}\NormalTok{ (m }\SpecialCharTok{+}\NormalTok{ ncen),}
                                \AttributeTok{iterations =} \DecValTok{20}\NormalTok{) \{}
\NormalTok{    out }\OtherTok{\textless{}{-}} \FunctionTok{matrix}\NormalTok{(}\ConstantTok{NA\_real\_}\NormalTok{, iterations }\SpecialCharTok{+} \DecValTok{1}\NormalTok{, }\DecValTok{2}\NormalTok{,}
                  \AttributeTok{dimnames =} \FunctionTok{list}\NormalTok{(}\ConstantTok{NULL}\NormalTok{, }\FunctionTok{c}\NormalTok{(}\StringTok{"lambda"}\NormalTok{, }\StringTok{"log\_lik"}\NormalTok{)))}
\NormalTok{    lambda }\OtherTok{\textless{}{-}}\NormalTok{ lambda0}
\NormalTok{    out[}\DecValTok{1}\NormalTok{, ] }\OtherTok{\textless{}{-}} \FunctionTok{c}\NormalTok{(lambda, }\FunctionTok{loglik\_obs\_pois}\NormalTok{(m, ybar, ncen, lambda))}
    \ControlFlowTok{for}\NormalTok{ (i }\ControlFlowTok{in} \FunctionTok{seq\_len}\NormalTok{(iterations)) \{}
\NormalTok{        y\_mis  }\OtherTok{\textless{}{-}}\NormalTok{ lambda }\SpecialCharTok{/}\NormalTok{ (}\DecValTok{1} \SpecialCharTok{+}\NormalTok{ lambda)                       }\CommentTok{\# E{-}step}
\NormalTok{        lambda }\OtherTok{\textless{}{-}}\NormalTok{ (m }\SpecialCharTok{*}\NormalTok{ ybar }\SpecialCharTok{+}\NormalTok{ ncen }\SpecialCharTok{*}\NormalTok{ y\_mis) }\SpecialCharTok{/}\NormalTok{ (m }\SpecialCharTok{+}\NormalTok{ ncen)       }\CommentTok{\# M{-}step}
\NormalTok{        out[i }\SpecialCharTok{+} \DecValTok{1}\NormalTok{, ] }\OtherTok{\textless{}{-}} \FunctionTok{c}\NormalTok{(lambda, }\FunctionTok{loglik\_obs\_pois}\NormalTok{(m, ybar, ncen, lambda))}
\NormalTok{    \}}
\NormalTok{    out}
\NormalTok{\}}
\end{Highlighting}
\end{Shaded}

\begin{Shaded}
\begin{Highlighting}[]
\FunctionTok{set.seed}\NormalTok{(}\DecValTok{812}\NormalTok{)}
\NormalTok{y    }\OtherTok{\textless{}{-}} \FunctionTok{rpois}\NormalTok{(}\DecValTok{200}\NormalTok{, }\DecValTok{3}\NormalTok{)}
\NormalTok{ncen }\OtherTok{\textless{}{-}} \FunctionTok{sum}\NormalTok{(y }\SpecialCharTok{\textless{}} \DecValTok{2}\NormalTok{)}
\NormalTok{m    }\OtherTok{\textless{}{-}} \FunctionTok{length}\NormalTok{(y) }\SpecialCharTok{{-}}\NormalTok{ ncen}
\NormalTok{ybar }\OtherTok{\textless{}{-}} \FunctionTok{mean}\NormalTok{(y[y }\SpecialCharTok{\textgreater{}=} \DecValTok{2}\NormalTok{])}
\FunctionTok{c}\NormalTok{(}\AttributeTok{m =}\NormalTok{ m, }\AttributeTok{ncen =}\NormalTok{ ncen, }\AttributeTok{ybar\_observed =}\NormalTok{ ybar)}
\end{Highlighting}
\end{Shaded}

\begin{verbatim}
            m          ncen ybar_observed 
   157.000000     43.000000      3.528662 
\end{verbatim}

Discarding the censored observations and averaging the rest gives an
estimate clearly above the true mean of \(3\): conditioning on
\(y \ge 2\) removes the small counts and biases the mean upward. EM
recovers them.

\begin{Shaded}
\begin{Highlighting}[]
\FunctionTok{round}\NormalTok{(}\FunctionTok{em\_censored\_poisson}\NormalTok{(m, ybar, ncen, }\AttributeTok{lambda0 =} \DecValTok{100}\NormalTok{, }\AttributeTok{iterations =} \DecValTok{10}\NormalTok{), }\DecValTok{5}\NormalTok{)}
\end{Highlighting}
\end{Shaded}

\begin{verbatim}
         lambda      log_lik
 [1,] 100.00000 -17250.28553
 [2,]   2.98287     68.31091
 [3,]   2.93102     68.40282
 [4,]   2.93031     68.40284
 [5,]   2.93030     68.40284
 [6,]   2.93030     68.40284
 [7,]   2.93030     68.40284
 [8,]   2.93030     68.40284
 [9,]   2.93030     68.40284
[10,]   2.93030     68.40284
[11,]   2.93030     68.40284
\end{verbatim}

\begin{Shaded}
\begin{Highlighting}[]
\FunctionTok{round}\NormalTok{(}\FunctionTok{em\_censored\_poisson}\NormalTok{(m, ybar, ncen, }\AttributeTok{lambda0 =}\NormalTok{ ybar, }\AttributeTok{iterations =} \DecValTok{10}\NormalTok{), }\DecValTok{5}\NormalTok{)}
\end{Highlighting}
\end{Shaded}

\begin{verbatim}
       lambda  log_lik
 [1,] 3.52866 57.76492
 [2,] 2.93752 68.40109
 [3,] 2.93040 68.40284
 [4,] 2.93030 68.40284
 [5,] 2.93030 68.40284
 [6,] 2.93030 68.40284
 [7,] 2.93030 68.40284
 [8,] 2.93030 68.40284
 [9,] 2.93030 68.40284
[10,] 2.93030 68.40284
[11,] 2.93030 68.40284
\end{verbatim}

Both starting values reach the same limit, and the log-likelihood
increases monotonically from each. Note how slowly the run from
\(\lambda^{(0)} = 100\) descends at first --- a single-parameter
illustration of the same slow-drift behaviour seen in the mixture app.

\subsection{Comparison with general-purpose
optimizers}\label{comparison-with-general-purpose-optimizers-1}

\begin{Shaded}
\begin{Highlighting}[]
\NormalTok{neg\_loglik\_pois }\OtherTok{\textless{}{-}} \ControlFlowTok{function}\NormalTok{(lambda, m, ybar, ncen) \{}
    \SpecialCharTok{{-}}\FunctionTok{loglik\_obs\_pois}\NormalTok{(m, ybar, ncen, lambda)}
\NormalTok{\}}

\FunctionTok{try}\NormalTok{(}\FunctionTok{nlm}\NormalTok{(neg\_loglik\_pois, }\AttributeTok{p =}\NormalTok{ ybar, }\AttributeTok{m =}\NormalTok{ m, }\AttributeTok{ybar =}\NormalTok{ ybar, }\AttributeTok{ncen =}\NormalTok{ ncen)}\SpecialCharTok{$}\NormalTok{estimate)}
\end{Highlighting}
\end{Shaded}

\begin{verbatim}
[1] 2.930295
\end{verbatim}

\begin{Shaded}
\begin{Highlighting}[]
\FunctionTok{try}\NormalTok{(}\FunctionTok{nlm}\NormalTok{(neg\_loglik\_pois, }\AttributeTok{p =} \DecValTok{100}\NormalTok{,  }\AttributeTok{m =}\NormalTok{ m, }\AttributeTok{ybar =}\NormalTok{ ybar, }\AttributeTok{ncen =}\NormalTok{ ncen)}\SpecialCharTok{$}\NormalTok{estimate)}
\end{Highlighting}
\end{Shaded}

\begin{verbatim}
[1] 2.930295
\end{verbatim}

\begin{Shaded}
\begin{Highlighting}[]
\FunctionTok{try}\NormalTok{(}\FunctionTok{optim}\NormalTok{(ybar, neg\_loglik\_pois, }\AttributeTok{method =} \StringTok{"CG"}\NormalTok{,}
          \AttributeTok{m =}\NormalTok{ m, }\AttributeTok{ybar =}\NormalTok{ ybar, }\AttributeTok{ncen =}\NormalTok{ ncen)}\SpecialCharTok{$}\NormalTok{par)}
\end{Highlighting}
\end{Shaded}

\begin{verbatim}
[1] 2.930297
\end{verbatim}

\begin{Shaded}
\begin{Highlighting}[]
\FunctionTok{try}\NormalTok{(}\FunctionTok{optim}\NormalTok{(}\DecValTok{100}\NormalTok{,  neg\_loglik\_pois, }\AttributeTok{method =} \StringTok{"CG"}\NormalTok{,}
          \AttributeTok{m =}\NormalTok{ m, }\AttributeTok{ybar =}\NormalTok{ ybar, }\AttributeTok{ncen =}\NormalTok{ ncen)}\SpecialCharTok{$}\NormalTok{par)}
\end{Highlighting}
\end{Shaded}

\begin{verbatim}
[1] 2.930297
\end{verbatim}

With one parameter and a smooth concave objective this is an easy
problem for any method; the contrast with the mixture example is
instructive, since there the multimodality and the constraint on \(p\)
were what made the general optimizers fragile.

\section{More Complex Applications}\label{more-complex-applications}

\subsection{Model-based clustering}\label{model-based-clustering}

\[
Z_i \sim \text{Categorical}(\pi_1, \ldots, \pi_K),
\qquad Y_i \mid Z_i = k \sim N_d(\mu_k, \Sigma_k) .
\]

This is the multivariate, \(K\)-component generalization of
Section~\ref{sec-mixture}. The E-step computes soft cluster memberships
\(P(Z_i = k \mid y_i, \hat\theta^{(t)})\); the M-step returns weighted
means, covariances and proportions. The \texttt{mclust} package
implements this with a family of constrained covariance structures.

\begin{Shaded}
\begin{Highlighting}[]
\FunctionTok{par}\NormalTok{(}\AttributeTok{mar =} \FunctionTok{c}\NormalTok{(}\DecValTok{4}\NormalTok{, }\DecValTok{4}\NormalTok{, }\DecValTok{1}\NormalTok{, }\DecValTok{1}\NormalTok{))}
\NormalTok{ctr }\OtherTok{\textless{}{-}} \FunctionTok{rbind}\NormalTok{(}\FunctionTok{c}\NormalTok{(}\DecValTok{0}\NormalTok{, }\DecValTok{0}\NormalTok{), }\FunctionTok{c}\NormalTok{(}\DecValTok{4}\NormalTok{, }\DecValTok{3}\NormalTok{), }\FunctionTok{c}\NormalTok{(}\DecValTok{1}\NormalTok{, }\DecValTok{5}\NormalTok{))}
\NormalTok{Sig }\OtherTok{\textless{}{-}} \FunctionTok{list}\NormalTok{(}\FunctionTok{diag}\NormalTok{(}\DecValTok{2}\NormalTok{) }\SpecialCharTok{*} \FloatTok{0.8}\NormalTok{, }\FunctionTok{matrix}\NormalTok{(}\FunctionTok{c}\NormalTok{(}\DecValTok{1}\NormalTok{, }\FloatTok{0.7}\NormalTok{, }\FloatTok{0.7}\NormalTok{, }\DecValTok{1}\NormalTok{), }\DecValTok{2}\NormalTok{), }\FunctionTok{diag}\NormalTok{(}\FunctionTok{c}\NormalTok{(}\FloatTok{1.2}\NormalTok{, }\FloatTok{0.4}\NormalTok{)))}
\NormalTok{cols }\OtherTok{\textless{}{-}} \FunctionTok{c}\NormalTok{(}\StringTok{"steelblue"}\NormalTok{, }\StringTok{"firebrick"}\NormalTok{, }\StringTok{"forestgreen"}\NormalTok{)}
\NormalTok{ang }\OtherTok{\textless{}{-}} \FunctionTok{seq}\NormalTok{(}\DecValTok{0}\NormalTok{, }\DecValTok{2} \SpecialCharTok{*}\NormalTok{ pi, }\AttributeTok{length.out =} \DecValTok{100}\NormalTok{)}
\FunctionTok{plot}\NormalTok{(}\ConstantTok{NA}\NormalTok{, }\AttributeTok{xlim =} \FunctionTok{c}\NormalTok{(}\SpecialCharTok{{-}}\DecValTok{3}\NormalTok{, }\DecValTok{7}\NormalTok{), }\AttributeTok{ylim =} \FunctionTok{c}\NormalTok{(}\SpecialCharTok{{-}}\DecValTok{3}\NormalTok{, }\DecValTok{8}\NormalTok{),}
     \AttributeTok{xlab =} \FunctionTok{expression}\NormalTok{(y[}\DecValTok{1}\NormalTok{]), }\AttributeTok{ylab =} \FunctionTok{expression}\NormalTok{(y[}\DecValTok{2}\NormalTok{]))}
\ControlFlowTok{for}\NormalTok{ (k }\ControlFlowTok{in} \DecValTok{1}\SpecialCharTok{:}\DecValTok{3}\NormalTok{) \{}
\NormalTok{  R }\OtherTok{\textless{}{-}} \FunctionTok{chol}\NormalTok{(Sig[[k]])}
\NormalTok{  pts }\OtherTok{\textless{}{-}} \FunctionTok{matrix}\NormalTok{(}\FunctionTok{rnorm}\NormalTok{(}\DecValTok{120}\NormalTok{), }\DecValTok{60}\NormalTok{) }\SpecialCharTok{\%*\%}\NormalTok{ R }\SpecialCharTok{+} \FunctionTok{rep}\NormalTok{(ctr[k, ], }\AttributeTok{each =} \DecValTok{60}\NormalTok{)}
  \FunctionTok{points}\NormalTok{(pts, }\AttributeTok{col =} \FunctionTok{adjustcolor}\NormalTok{(cols[k], }\FloatTok{0.6}\NormalTok{), }\AttributeTok{pch =} \DecValTok{19}\NormalTok{, }\AttributeTok{cex =} \FloatTok{0.8}\NormalTok{)}
\NormalTok{  ell }\OtherTok{\textless{}{-}}\NormalTok{ (}\FunctionTok{cbind}\NormalTok{(}\FunctionTok{cos}\NormalTok{(ang), }\FunctionTok{sin}\NormalTok{(ang)) }\SpecialCharTok{*} \FloatTok{2.1}\NormalTok{) }\SpecialCharTok{\%*\%}\NormalTok{ R }\SpecialCharTok{+} \FunctionTok{rep}\NormalTok{(ctr[k, ], }\AttributeTok{each =} \DecValTok{100}\NormalTok{)}
  \FunctionTok{lines}\NormalTok{(ell, }\AttributeTok{col =}\NormalTok{ cols[k], }\AttributeTok{lwd =} \DecValTok{2}\NormalTok{)}
  \FunctionTok{points}\NormalTok{(ctr[k, }\DecValTok{1}\NormalTok{], ctr[k, }\DecValTok{2}\NormalTok{], }\AttributeTok{pch =} \DecValTok{4}\NormalTok{, }\AttributeTok{lwd =} \DecValTok{3}\NormalTok{, }\AttributeTok{col =}\NormalTok{ cols[k])}
\NormalTok{\}}
\end{Highlighting}
\end{Shaded}

\begin{figure}[H]

\centering{

\pandocbounded{\includegraphics[keepaspectratio]{chapters/unit07_em_files/figure-pdf/fig-clustering-1.pdf}}

}

\caption{\label{fig-clustering}Three-component multivariate normal
mixture with unequal covariances.}

\end{figure}%

\subsection{Hidden Markov models}\label{hidden-markov-models}

Hidden states \(h_1 \to h_2 \to \cdots \to h_n\) form a Markov chain,
and each observation depends only on its own state:

\[
P(y_{1:n} \mid h_{1:n}, \theta) = \prod_{i=1}^{n} P(y_i \mid h_i, \theta) .
\]

The E-step is the forward--backward algorithm, which computes
\(P(h_i \mid y_{1:n}, \hat\theta^{(t)})\) and
\(P(h_i, h_{i+1} \mid y_{1:n}, \hat\theta^{(t)})\) without enumerating
all \(K^n\) state sequences; the M-step produces weighted transition and
emission estimates. The combination is known as the Baum--Welch
algorithm, and it is the clearest example of an E-step that is itself a
nontrivial computation.

\begin{Shaded}
\begin{Highlighting}[]
\FunctionTok{par}\NormalTok{(}\AttributeTok{mar =} \FunctionTok{c}\NormalTok{(}\DecValTok{4}\NormalTok{, }\DecValTok{4}\NormalTok{, }\DecValTok{1}\NormalTok{, }\DecValTok{1}\NormalTok{)); N }\OtherTok{\textless{}{-}} \DecValTok{120}
\NormalTok{Pm }\OtherTok{\textless{}{-}} \FunctionTok{matrix}\NormalTok{(}\FunctionTok{c}\NormalTok{(}\FloatTok{0.95}\NormalTok{, }\FloatTok{0.05}\NormalTok{, }\FloatTok{0.1}\NormalTok{, }\FloatTok{0.9}\NormalTok{), }\DecValTok{2}\NormalTok{, }\AttributeTok{byrow =} \ConstantTok{TRUE}\NormalTok{)}
\NormalTok{h }\OtherTok{\textless{}{-}} \FunctionTok{integer}\NormalTok{(N); h[}\DecValTok{1}\NormalTok{] }\OtherTok{\textless{}{-}} \DecValTok{1}
\ControlFlowTok{for}\NormalTok{ (i }\ControlFlowTok{in} \DecValTok{2}\SpecialCharTok{:}\NormalTok{N) h[i] }\OtherTok{\textless{}{-}} \FunctionTok{sample}\NormalTok{(}\DecValTok{1}\SpecialCharTok{:}\DecValTok{2}\NormalTok{, }\DecValTok{1}\NormalTok{, }\AttributeTok{prob =}\NormalTok{ Pm[h[i }\SpecialCharTok{{-}} \DecValTok{1}\NormalTok{], ])}
\NormalTok{yy }\OtherTok{\textless{}{-}} \FunctionTok{rnorm}\NormalTok{(N, }\AttributeTok{mean =} \FunctionTok{c}\NormalTok{(}\DecValTok{0}\NormalTok{, }\DecValTok{3}\NormalTok{)[h])}
\FunctionTok{plot}\NormalTok{(yy, }\AttributeTok{type =} \StringTok{"n"}\NormalTok{, }\AttributeTok{xlab =} \StringTok{"i"}\NormalTok{, }\AttributeTok{ylab =} \FunctionTok{expression}\NormalTok{(y[i]))}
\NormalTok{r }\OtherTok{\textless{}{-}} \FunctionTok{rle}\NormalTok{(h); e }\OtherTok{\textless{}{-}} \FunctionTok{cumsum}\NormalTok{(r}\SpecialCharTok{$}\NormalTok{lengths); s }\OtherTok{\textless{}{-}} \FunctionTok{c}\NormalTok{(}\DecValTok{1}\NormalTok{, }\FunctionTok{head}\NormalTok{(e, }\SpecialCharTok{{-}}\DecValTok{1}\NormalTok{) }\SpecialCharTok{+} \DecValTok{1}\NormalTok{)}
\FunctionTok{rect}\NormalTok{(s }\SpecialCharTok{{-}} \FloatTok{0.5}\NormalTok{, }\FunctionTok{par}\NormalTok{(}\StringTok{"usr"}\NormalTok{)[}\DecValTok{3}\NormalTok{], e }\SpecialCharTok{+} \FloatTok{0.5}\NormalTok{, }\FunctionTok{par}\NormalTok{(}\StringTok{"usr"}\NormalTok{)[}\DecValTok{4}\NormalTok{],}
     \AttributeTok{col =} \FunctionTok{adjustcolor}\NormalTok{(}\FunctionTok{c}\NormalTok{(}\StringTok{"steelblue"}\NormalTok{, }\StringTok{"firebrick"}\NormalTok{)[r}\SpecialCharTok{$}\NormalTok{values], }\FloatTok{0.15}\NormalTok{), }\AttributeTok{border =} \ConstantTok{NA}\NormalTok{)}
\FunctionTok{lines}\NormalTok{(yy, }\AttributeTok{col =} \StringTok{"gray30"}\NormalTok{)}
\FunctionTok{points}\NormalTok{(yy, }\AttributeTok{pch =} \DecValTok{19}\NormalTok{, }\AttributeTok{col =} \FunctionTok{c}\NormalTok{(}\StringTok{"steelblue"}\NormalTok{, }\StringTok{"firebrick"}\NormalTok{)[h], }\AttributeTok{cex =} \FloatTok{0.7}\NormalTok{)}
\FunctionTok{legend}\NormalTok{(}\StringTok{"topright"}\NormalTok{, }\AttributeTok{bty =} \StringTok{"n"}\NormalTok{, }\AttributeTok{pch =} \DecValTok{19}\NormalTok{, }\AttributeTok{col =} \FunctionTok{c}\NormalTok{(}\StringTok{"steelblue"}\NormalTok{, }\StringTok{"firebrick"}\NormalTok{),}
       \AttributeTok{legend =} \FunctionTok{c}\NormalTok{(}\FunctionTok{expression}\NormalTok{(h[i] }\SpecialCharTok{==} \DecValTok{1}\NormalTok{), }\FunctionTok{expression}\NormalTok{(h[i] }\SpecialCharTok{==} \DecValTok{2}\NormalTok{)), }\AttributeTok{horiz =} \ConstantTok{TRUE}\NormalTok{)}
\end{Highlighting}
\end{Shaded}

\begin{figure}[H]

\centering{

\pandocbounded{\includegraphics[keepaspectratio]{chapters/unit07_em_files/figure-pdf/fig-hmm-1.pdf}}

}

\caption{\label{fig-hmm}A two-state hidden Markov chain with normal
emissions; background shading shows the hidden state.}

\end{figure}%

\subsection{Factor analysis and the common
pattern}\label{factor-analysis-and-the-common-pattern}

\[
y_i = \Gamma x_i + \varepsilon_i,
\qquad x_i \sim N_q(0, I) \text{ hidden},
\qquad \varepsilon_i \sim N_d(0, \Psi) .
\]

The E-step uses the fact that \(x_i \mid y_i, \hat\theta^{(t)}\) is
normal with closed-form mean and covariance; the M-step is a regression
of \(y\) on the expected factors.

Across all these applications the recipe is the same:

\begin{enumerate}
\def\labelenumi{\arabic{enumi}.}
\tightlist
\item
  Introduce a hidden variable that makes the complete-data model
  tractable.
\item
  \textbf{E-step}: compute the conditional distribution of the hidden
  data, or just the expected sufficient statistics if
  \(\ell_{\text{comp}}\) is linear in them.
\item
  \textbf{M-step}: compute the complete-data MLE with those expected
  sufficient statistics in place of the unobserved ones.
\end{enumerate}

The art is entirely in step 1. The same observed-data model often admits
several choices of \(Z\), and they lead to EM algorithms with quite
different convergence rates --- faster when the augmentation adds little
information, slower when the missing information is a large fraction of
the total.

\section{Summary}\label{summary-1}

\begin{itemize}
\tightlist
\item
  EM maximizes \(\ell_{\text{obs}}\) by alternating between forming
  \(Q(\theta \mid \hat\theta^{(t)}) = E_{Z \mid Y, \hat\theta^{(t)}}[\ell_{\text{comp}}]\)
  and maximizing it.
\item
  Every iteration weakly increases \(\ell_{\text{obs}}\). Equivalently,
  EM is coordinate ascent on
  \(F(\tilde P, \theta) = \ell_{\text{obs}}(\theta) - \mathrm{KL}(\tilde P \,\|\, P(Z \mid Y, \theta))\).
\item
  When \(\ell_{\text{comp}}\) is linear in sufficient statistics of
  \(Z\), the E-step reduces to computing their conditional expectations.
  Otherwise the full conditional distribution is required, and ``plug in
  the expectation'' is wrong.
\item
  Convergence is linear, often slow, and only to a stationary point. Run
  from several starting values and compare the attained log-likelihoods.
\item
  Standard errors are not a by-product. Use Louis' identity, the SEM
  algorithm, a bootstrap, or a final Newton step.
\end{itemize}

\bookmarksetup{startatroot}

\chapter{Numerical Quadrature}\label{numerical-quadrature}

\begin{Shaded}
\begin{Highlighting}[]
\FunctionTok{library}\NormalTok{(metRology)}
\end{Highlighting}
\end{Shaded}

The previous two chapters found the mode of a log-likelihood or
log-posterior by optimization. Many Bayesian calculations instead need
an \emph{area}: a tail probability, a posterior expectation, or the
normalizing constant of a posterior density. This chapter develops
\textbf{numerical quadrature} --- approximating an integral by a
weighted sum of function values at a finite set of points --- and
applies it to the central example of the rest of the book, the marginal
likelihood of a Bayesian model. The next chapter (Laplace's method)
revisits the same example with a method based on the mode instead of a
grid.

\section{The Problem and
Transformations}\label{the-problem-and-transformations}

\subsection{The general setup}\label{the-general-setup}

The goal is to approximate \[
A = \int_a^b f(x)\,dx
\] for finite \(a,b\) by evaluating \(f\) at \(n\) points
\(x_1,\ldots,x_n\) and forming a weighted sum \[
\tilde A_n = \sum_{i=1}^n w_i\, f(x_i).
\] The quadrature rules below differ only in how the points and weights
are chosen. Typical statistical uses are probabilities \(P(a < X < b)\),
expectations \(E[h(X)]\), and normalizing constants such as the marginal
likelihood \(P(y) = \int L(\theta)\pi(\theta)\,d\theta\) in Bayesian
inference.

\subsection{Infinite ranges: the logistic
transformation}\label{infinite-ranges-the-logistic-transformation}

A grid needs a finite interval, so when the natural range of integration
is \((-\infty,\infty)\) we first change variables. The logistic function
maps \(x \in (-\infty,\infty)\) to \(y \in (0,1)\), \[
y = \frac{1}{1+e^{-x}}, \qquad x = t(y) = \log\frac{y}{1-y}, \qquad
t'(y) = \frac{1}{y(1-y)},
\] and the standard change-of-variables identity gives \[
\int_{-\infty}^{\infty} f(x)\,dx
 = \int_0^1 f\big(t(y)\big)\,t'(y)\,dy
 = \int_0^1 g(y)\,dy, \qquad
g(y) = \frac{f\big(\log\frac{y}{1-y}\big)}{y(1-y)}.
\] Any quadrature rule for a finite interval now applies to \(g\) on
\((0,1)\). Two consequences matter later in the chapter:

\begin{itemize}
\tightlist
\item
  Points equally spaced in \(y\) correspond to points in \(x\) that are
  dense near \(x=0\) and increasingly sparse in the tails --- the grid
  automatically concentrates effort where the logistic curve is steep.
\item
  The factor \(1/\{y(1-y)\}\) blows up at \(y=0\) and \(y=1\), so \(g\)
  is typically undefined at the endpoints even when \(f\) is perfectly
  well behaved. This is the reason the midpoint rule, which never
  evaluates the integrand at the endpoints of a panel, is the natural
  choice here.
\end{itemize}

For a half-line \((0,\infty)\) the same idea applies with
\(x = -\log(1-y)\) or \(x = y/(1-y)\) in place of the logistic map.

\begin{Shaded}
\begin{Highlighting}[]
\FunctionTok{par}\NormalTok{(}\AttributeTok{mfrow =} \FunctionTok{c}\NormalTok{(}\DecValTok{1}\NormalTok{, }\DecValTok{3}\NormalTok{), }\AttributeTok{mar =} \FunctionTok{c}\NormalTok{(}\DecValTok{4}\NormalTok{, }\DecValTok{4}\NormalTok{, }\DecValTok{2}\NormalTok{, }\DecValTok{1}\NormalTok{))}

\FunctionTok{curve}\NormalTok{(}\DecValTok{1} \SpecialCharTok{/}\NormalTok{ (}\DecValTok{1} \SpecialCharTok{+} \FunctionTok{exp}\NormalTok{(}\SpecialCharTok{{-}}\NormalTok{x)), }\SpecialCharTok{{-}}\DecValTok{6}\NormalTok{, }\DecValTok{6}\NormalTok{, }\AttributeTok{lwd =} \DecValTok{3}\NormalTok{, }\AttributeTok{xlab =} \StringTok{"x"}\NormalTok{, }\AttributeTok{ylab =} \StringTok{"y"}\NormalTok{,}
      \AttributeTok{main =} \FunctionTok{expression}\NormalTok{(y }\SpecialCharTok{==} \DecValTok{1} \SpecialCharTok{/}\NormalTok{ (}\DecValTok{1} \SpecialCharTok{+}\NormalTok{ e}\SpecialCharTok{\^{}{-}}\NormalTok{x)))}
\NormalTok{ys }\OtherTok{\textless{}{-}} \FunctionTok{seq}\NormalTok{(}\FloatTok{0.05}\NormalTok{, }\FloatTok{0.95}\NormalTok{, }\AttributeTok{by =} \FloatTok{0.1}\NormalTok{)}
\NormalTok{xs }\OtherTok{\textless{}{-}} \FunctionTok{log}\NormalTok{(ys }\SpecialCharTok{/}\NormalTok{ (}\DecValTok{1} \SpecialCharTok{{-}}\NormalTok{ ys))}
\FunctionTok{segments}\NormalTok{(xs, }\DecValTok{0}\NormalTok{, xs, ys, }\AttributeTok{lty =} \DecValTok{3}\NormalTok{, }\AttributeTok{col =} \StringTok{"gray50"}\NormalTok{)}
\FunctionTok{segments}\NormalTok{(}\SpecialCharTok{{-}}\DecValTok{6}\NormalTok{, ys, xs, ys, }\AttributeTok{lty =} \DecValTok{3}\NormalTok{, }\AttributeTok{col =} \StringTok{"gray50"}\NormalTok{)}
\FunctionTok{points}\NormalTok{(xs, ys, }\AttributeTok{pch =} \DecValTok{19}\NormalTok{, }\AttributeTok{col =} \StringTok{"firebrick"}\NormalTok{)}

\NormalTok{f }\OtherTok{\textless{}{-}} \ControlFlowTok{function}\NormalTok{(x) }\FloatTok{0.6} \SpecialCharTok{*} \FunctionTok{dnorm}\NormalTok{(x, }\SpecialCharTok{{-}}\DecValTok{1}\NormalTok{, }\DecValTok{1}\NormalTok{) }\SpecialCharTok{+} \FloatTok{0.4} \SpecialCharTok{*} \FunctionTok{dnorm}\NormalTok{(x, }\DecValTok{2}\NormalTok{, }\FloatTok{1.5}\NormalTok{)}
\FunctionTok{curve}\NormalTok{(f, }\SpecialCharTok{{-}}\DecValTok{6}\NormalTok{, }\DecValTok{6}\NormalTok{, }\AttributeTok{lwd =} \DecValTok{3}\NormalTok{, }\AttributeTok{xlab =} \StringTok{"x"}\NormalTok{, }\AttributeTok{ylab =} \StringTok{"f(x)"}\NormalTok{, }\AttributeTok{main =} \StringTok{"original integrand"}\NormalTok{)}
\NormalTok{xs2 }\OtherTok{\textless{}{-}} \FunctionTok{seq}\NormalTok{(}\SpecialCharTok{{-}}\DecValTok{6}\NormalTok{, }\DecValTok{6}\NormalTok{, }\AttributeTok{length.out =} \DecValTok{400}\NormalTok{)}
\FunctionTok{polygon}\NormalTok{(}\FunctionTok{c}\NormalTok{(xs2, }\FunctionTok{rev}\NormalTok{(xs2)), }\FunctionTok{c}\NormalTok{(}\FunctionTok{f}\NormalTok{(xs2), }\FunctionTok{rep}\NormalTok{(}\DecValTok{0}\NormalTok{, }\DecValTok{400}\NormalTok{)), }\AttributeTok{col =} \FunctionTok{adjustcolor}\NormalTok{(}\StringTok{"firebrick"}\NormalTok{, }\FloatTok{0.3}\NormalTok{), }\AttributeTok{border =} \ConstantTok{NA}\NormalTok{)}
\FunctionTok{rug}\NormalTok{(xs, }\AttributeTok{col =} \StringTok{"firebrick"}\NormalTok{, }\AttributeTok{lwd =} \DecValTok{2}\NormalTok{)}

\NormalTok{g }\OtherTok{\textless{}{-}} \ControlFlowTok{function}\NormalTok{(y) }\FunctionTok{f}\NormalTok{(}\FunctionTok{log}\NormalTok{(y }\SpecialCharTok{/}\NormalTok{ (}\DecValTok{1} \SpecialCharTok{{-}}\NormalTok{ y))) }\SpecialCharTok{/}\NormalTok{ (y }\SpecialCharTok{*}\NormalTok{ (}\DecValTok{1} \SpecialCharTok{{-}}\NormalTok{ y))}
\FunctionTok{curve}\NormalTok{(g, }\FloatTok{0.001}\NormalTok{, }\FloatTok{0.999}\NormalTok{, }\AttributeTok{lwd =} \DecValTok{3}\NormalTok{, }\AttributeTok{n =} \DecValTok{501}\NormalTok{, }\AttributeTok{xlab =} \StringTok{"y"}\NormalTok{, }\AttributeTok{ylab =} \StringTok{"g(y)"}\NormalTok{,}
      \AttributeTok{main =} \StringTok{"transformed integrand on (0, 1)"}\NormalTok{)}
\NormalTok{ys2 }\OtherTok{\textless{}{-}} \FunctionTok{seq}\NormalTok{(}\FloatTok{0.001}\NormalTok{, }\FloatTok{0.999}\NormalTok{, }\AttributeTok{length.out =} \DecValTok{400}\NormalTok{)}
\FunctionTok{polygon}\NormalTok{(}\FunctionTok{c}\NormalTok{(ys2, }\FunctionTok{rev}\NormalTok{(ys2)), }\FunctionTok{c}\NormalTok{(}\FunctionTok{g}\NormalTok{(ys2), }\FunctionTok{rep}\NormalTok{(}\DecValTok{0}\NormalTok{, }\DecValTok{400}\NormalTok{)), }\AttributeTok{col =} \FunctionTok{adjustcolor}\NormalTok{(}\StringTok{"steelblue"}\NormalTok{, }\FloatTok{0.3}\NormalTok{), }\AttributeTok{border =} \ConstantTok{NA}\NormalTok{)}
\FunctionTok{rug}\NormalTok{(ys, }\AttributeTok{col =} \StringTok{"firebrick"}\NormalTok{, }\AttributeTok{lwd =} \DecValTok{2}\NormalTok{)}
\end{Highlighting}
\end{Shaded}

\begin{figure}[H]

\centering{

\pandocbounded{\includegraphics[keepaspectratio]{chapters/unit08_integral_files/figure-pdf/fig-logistic-transform-1.pdf}}

}

\caption{\label{fig-logistic-transform}Equally spaced points in \(y\)
(red ticks, right panel) map to unequally spaced points in \(x\) (red
ticks, middle panel). The two shaded areas are equal: the transform
changes where we look, not the value of the integral.}

\end{figure}%

\subsection{Probabilities as a special
case}\label{probabilities-as-a-special-case}

For \(X\) with density \(\varphi\),
e.g.~\(\varphi(x) = \frac{1}{\sqrt{2\pi}}e^{-x^2/2}\), \[
P(a < X < b) = \int_a^b \varphi(x)\,dx = \int_{-\infty}^{\infty} \varphi^{(c)}(x)\,dx,
\qquad \varphi^{(c)}(x) = \varphi(x)\,I(a < x < b).
\] Either integrate \(\varphi\) directly over the finite interval
\([a,b]\), or integrate the \emph{truncated} function \(\varphi^{(c)}\)
over the whole line, letting the indicator do the cutting. The second
view is the more useful one: the same machinery --- transform to
\((0,1)\), apply a quadrature rule --- computes \(P(X \in S)\) for an
arbitrary set \(S\) and, more generally, \(E[h(X)] = \int
h(x)\varphi(x)\,dx\) for an arbitrary \(h\).

\section{Quadrature Rules on a Finite
Interval}\label{quadrature-rules-on-a-finite-interval}

All three rules below partition \([a,b]\) into \(n\) panels of width
\(h_n =
(b-a)/n\) and differ only in how \(f\) is approximated \emph{within}
each panel.

\subsection{Midpoint rule}\label{midpoint-rule}

Approximate \(f\) on each panel by its value at the panel's midpoint (a
constant, i.e.~a rectangle): \[
\tilde A_n = h_n \sum_{i=1}^n f\!\Big(a + \big(i - \tfrac12\big)h_n\Big).
\] For smooth \(f\) the error is
\(A - \tilde A_n = \dfrac{(b-a)h_n^2}{24}f''(\xi) =
O(h_n^2)\) for some \(\xi \in (a,b)\). The rule never evaluates \(f\) at
the panel edges, so in particular it never touches \(a\) or \(b\) ---
convenient when \(f\) is singular or undefined at the endpoints, exactly
the situation created by the logistic transform above.

\subsection{Trapezoidal rule}\label{trapezoidal-rule}

Approximate \(f\) on \([x_{i-1},x_i]\) by the straight line through the
two endpoint values, so each panel contributes a trapezoid,
\(\tfrac12[f(x_{i-1})+f(x_i)](x_i-x_{i-1})\). Summing over panels,
interior points are shared by two neighbouring panels and each
contributes once (not twice) to the total: \[
\tilde A_n = \tfrac12 h_n\big[f(a) + f(b)\big] + h_n \sum_{i=1}^{n-1} f(a+ih_n).
\] The error is
\(A - \tilde A_n = -\dfrac{(b-a)h_n^2}{12}f''(\xi) = O(h_n^2)\) --- the
same order as the midpoint rule, but with twice the constant and the
opposite sign. It is exact for linear \(f\).

\subsection{Simpson's rule}\label{simpsons-rule}

Take \(n = 2m\) even. On each \emph{pair} of panels \([x_0,x_1]\) with
\(x_1-x_0=2h_n\), replace \(f\) by the quadratic through the two
endpoints and the midpoint; integrating that quadratic exactly gives \[
\int_{x_0}^{x_1} f(x)\,dx \approx \frac{x_1-x_0}{6}
\Big[f(x_0) + 4f\big(\tfrac{x_0+x_1}{2}\big) + f(x_1)\Big].
\] Summing over the \(m\) panel-pairs, with each even interior point
shared by two pairs, \[
\tilde A_{2m} = \frac{h_n}{3}\left[f(a) + f(b)
 + 4\sum_{i=1}^{m} f\big(a+(2i-1)h_n\big)
 + 2\sum_{i=1}^{m-1} f\big(a+2ih_n\big)\right].
\] The error is
\(A - \tilde A_{2m} = -\dfrac{(b-a)h_n^4}{180}f^{(4)}(\xi) =
O(h_n^4)\), two orders faster than the other two rules, and Simpson's
rule is exact for cubics. It can also be written as a weighted average
of the other two rules on the same panel pair:
\(\text{Simpson} = \tfrac23\,\text{midpoint} +
\tfrac13\,\text{trapezoid}\).

Table~\ref{tbl-rules-compared} summarizes the three rules.

\begin{longtable}[]{@{}
  >{\raggedright\arraybackslash}p{(\linewidth - 8\tabcolsep) * \real{0.2000}}
  >{\raggedright\arraybackslash}p{(\linewidth - 8\tabcolsep) * \real{0.2000}}
  >{\raggedright\arraybackslash}p{(\linewidth - 8\tabcolsep) * \real{0.2000}}
  >{\raggedright\arraybackslash}p{(\linewidth - 8\tabcolsep) * \real{0.2000}}
  >{\raggedright\arraybackslash}p{(\linewidth - 8\tabcolsep) * \real{0.2000}}@{}}
\caption{The three quadrature rules
compared}\label{tbl-rules-compared}\tabularnewline
\toprule\noalign{}
\begin{minipage}[b]{\linewidth}\raggedright
Rule
\end{minipage} & \begin{minipage}[b]{\linewidth}\raggedright
Approximates \(f\) on a panel by
\end{minipage} & \begin{minipage}[b]{\linewidth}\raggedright
Uses endpoints?
\end{minipage} & \begin{minipage}[b]{\linewidth}\raggedright
Error
\end{minipage} & \begin{minipage}[b]{\linewidth}\raggedright
Exact for
\end{minipage} \\
\midrule\noalign{}
\endfirsthead
\toprule\noalign{}
\begin{minipage}[b]{\linewidth}\raggedright
Rule
\end{minipage} & \begin{minipage}[b]{\linewidth}\raggedright
Approximates \(f\) on a panel by
\end{minipage} & \begin{minipage}[b]{\linewidth}\raggedright
Uses endpoints?
\end{minipage} & \begin{minipage}[b]{\linewidth}\raggedright
Error
\end{minipage} & \begin{minipage}[b]{\linewidth}\raggedright
Exact for
\end{minipage} \\
\midrule\noalign{}
\endhead
\bottomrule\noalign{}
\endlastfoot
Midpoint & constant at panel midpoint & no & \(O(h_n^2)\) & linear \\
Trapezoidal & line through panel endpoints & yes & \(O(h_n^2)\) &
linear \\
Simpson & quadratic through 3 points (2 panels) & yes & \(O(h_n^4)\) &
cubic \\
\end{longtable}

\subsection{Implementation and
accuracy}\label{implementation-and-accuracy}

All three rules are one line of vectorized R each --- no explicit loop
over panels is needed:

\begin{Shaded}
\begin{Highlighting}[]
\NormalTok{midpoint  }\OtherTok{\textless{}{-}} \ControlFlowTok{function}\NormalTok{(f, a, b, n) \{}
\NormalTok{    h }\OtherTok{\textless{}{-}}\NormalTok{ (b }\SpecialCharTok{{-}}\NormalTok{ a) }\SpecialCharTok{/}\NormalTok{ n}
\NormalTok{    h }\SpecialCharTok{*} \FunctionTok{sum}\NormalTok{(}\FunctionTok{f}\NormalTok{(a }\SpecialCharTok{+}\NormalTok{ (}\FunctionTok{seq\_len}\NormalTok{(n) }\SpecialCharTok{{-}} \FloatTok{0.5}\NormalTok{) }\SpecialCharTok{*}\NormalTok{ h))}
\NormalTok{\}}
\NormalTok{trapezoid }\OtherTok{\textless{}{-}} \ControlFlowTok{function}\NormalTok{(f, a, b, n) \{}
\NormalTok{    h }\OtherTok{\textless{}{-}}\NormalTok{ (b }\SpecialCharTok{{-}}\NormalTok{ a) }\SpecialCharTok{/}\NormalTok{ n}
\NormalTok{    x }\OtherTok{\textless{}{-}}\NormalTok{ a }\SpecialCharTok{+}\NormalTok{ (}\DecValTok{0}\SpecialCharTok{:}\NormalTok{n) }\SpecialCharTok{*}\NormalTok{ h}
\NormalTok{    h }\SpecialCharTok{*}\NormalTok{ (}\FunctionTok{sum}\NormalTok{(}\FunctionTok{f}\NormalTok{(x)) }\SpecialCharTok{{-}} \FloatTok{0.5} \SpecialCharTok{*}\NormalTok{ (}\FunctionTok{f}\NormalTok{(a) }\SpecialCharTok{+} \FunctionTok{f}\NormalTok{(b)))}
\NormalTok{\}}
\NormalTok{simpson   }\OtherTok{\textless{}{-}} \ControlFlowTok{function}\NormalTok{(f, a, b, n) \{}
    \FunctionTok{stopifnot}\NormalTok{(n }\SpecialCharTok{\%\%} \DecValTok{2} \SpecialCharTok{==} \DecValTok{0}\NormalTok{)}
\NormalTok{    h }\OtherTok{\textless{}{-}}\NormalTok{ (b }\SpecialCharTok{{-}}\NormalTok{ a) }\SpecialCharTok{/}\NormalTok{ n}
\NormalTok{    x }\OtherTok{\textless{}{-}}\NormalTok{ a }\SpecialCharTok{+}\NormalTok{ (}\DecValTok{0}\SpecialCharTok{:}\NormalTok{n) }\SpecialCharTok{*}\NormalTok{ h}
\NormalTok{    h }\SpecialCharTok{/} \DecValTok{3} \SpecialCharTok{*} \FunctionTok{sum}\NormalTok{(}\FunctionTok{f}\NormalTok{(x) }\SpecialCharTok{*} \FunctionTok{c}\NormalTok{(}\DecValTok{1}\NormalTok{, }\FunctionTok{rep}\NormalTok{(}\FunctionTok{c}\NormalTok{(}\DecValTok{4}\NormalTok{, }\DecValTok{2}\NormalTok{), n }\SpecialCharTok{/} \DecValTok{2} \SpecialCharTok{{-}} \DecValTok{1}\NormalTok{), }\DecValTok{4}\NormalTok{, }\DecValTok{1}\NormalTok{))}
\NormalTok{\}}
\end{Highlighting}
\end{Shaded}

We check all three against a case with a known answer,
\(P(0.5 < Z < 2)\) for \(Z \sim N(0,1)\):

\begin{Shaded}
\begin{Highlighting}[]
\NormalTok{exact }\OtherTok{\textless{}{-}} \FunctionTok{pnorm}\NormalTok{(}\DecValTok{2}\NormalTok{) }\SpecialCharTok{{-}} \FunctionTok{pnorm}\NormalTok{(}\FloatTok{0.5}\NormalTok{)}
\NormalTok{errs }\OtherTok{\textless{}{-}} \FunctionTok{sapply}\NormalTok{(}\FunctionTok{list}\NormalTok{(}\AttributeTok{midpoint =}\NormalTok{ midpoint, }\AttributeTok{trapezoid =}\NormalTok{ trapezoid, }\AttributeTok{simpson =}\NormalTok{ simpson),}
               \ControlFlowTok{function}\NormalTok{(rule) }\FunctionTok{rule}\NormalTok{(dnorm, }\FloatTok{0.5}\NormalTok{, }\DecValTok{2}\NormalTok{, }\DecValTok{10}\NormalTok{) }\SpecialCharTok{{-}}\NormalTok{ exact)}
\NormalTok{knitr}\SpecialCharTok{::}\FunctionTok{kable}\NormalTok{(}\FunctionTok{t}\NormalTok{(errs), }\AttributeTok{digits =} \DecValTok{10}\NormalTok{, }\AttributeTok{row.names =} \ConstantTok{FALSE}\NormalTok{)}
\end{Highlighting}
\end{Shaded}

{

\begin{longtable}[]{@{}rrr@{}}

\caption{\label{tbl-rules-accuracy}Error of each rule with only \(n=10\)
panels, relative to \texttt{pnorm}.}

\tabularnewline

\toprule\noalign{}
midpoint & trapezoid & simpson \\
\midrule\noalign{}
\endhead
\bottomrule\noalign{}
\endlastfoot
-6.4163e-05 & 0.0001280126 & -1.6895e-06 \\

\end{longtable}

}

With only 10 panels, Simpson's rule is already accurate to about
\(10^{-7}\), several orders of magnitude better than the other two.

\begin{Shaded}
\begin{Highlighting}[]
\NormalTok{ns  }\OtherTok{\textless{}{-}} \DecValTok{2}\SpecialCharTok{\^{}}\NormalTok{(}\DecValTok{1}\SpecialCharTok{:}\DecValTok{9}\NormalTok{)}
\NormalTok{err }\OtherTok{\textless{}{-}} \FunctionTok{sapply}\NormalTok{(}\FunctionTok{list}\NormalTok{(midpoint, trapezoid, simpson),}
              \ControlFlowTok{function}\NormalTok{(rule) }\FunctionTok{sapply}\NormalTok{(ns, }\ControlFlowTok{function}\NormalTok{(n) }\FunctionTok{abs}\NormalTok{(}\FunctionTok{rule}\NormalTok{(dnorm, }\FloatTok{0.5}\NormalTok{, }\DecValTok{2}\NormalTok{, n) }\SpecialCharTok{{-}}\NormalTok{ exact)))}
\FunctionTok{matplot}\NormalTok{(ns, err, }\AttributeTok{type =} \StringTok{"b"}\NormalTok{, }\AttributeTok{log =} \StringTok{"xy"}\NormalTok{, }\AttributeTok{pch =} \DecValTok{19}\NormalTok{, }\AttributeTok{lty =} \DecValTok{1}\NormalTok{, }\AttributeTok{lwd =} \DecValTok{2}\NormalTok{,}
        \AttributeTok{col =} \FunctionTok{c}\NormalTok{(}\StringTok{"steelblue"}\NormalTok{, }\StringTok{"firebrick"}\NormalTok{, }\StringTok{"forestgreen"}\NormalTok{),}
        \AttributeTok{xlab =} \StringTok{"number of panels n"}\NormalTok{, }\AttributeTok{ylab =} \StringTok{"absolute error"}\NormalTok{)}
\FunctionTok{legend}\NormalTok{(}\StringTok{"bottomleft"}\NormalTok{, }\AttributeTok{bty =} \StringTok{"n"}\NormalTok{, }\AttributeTok{lwd =} \DecValTok{2}\NormalTok{, }\AttributeTok{pch =} \DecValTok{19}\NormalTok{,}
       \AttributeTok{col =} \FunctionTok{c}\NormalTok{(}\StringTok{"steelblue"}\NormalTok{, }\StringTok{"firebrick"}\NormalTok{, }\StringTok{"forestgreen"}\NormalTok{),}
       \AttributeTok{legend =} \FunctionTok{c}\NormalTok{(}\StringTok{"midpoint, slope {-}2"}\NormalTok{, }\StringTok{"trapezoid, slope {-}2"}\NormalTok{, }\StringTok{"Simpson, slope {-}4"}\NormalTok{))}
\end{Highlighting}
\end{Shaded}

\begin{figure}[H]

\centering{

\pandocbounded{\includegraphics[keepaspectratio]{chapters/unit08_integral_files/figure-pdf/fig-convergence-1.pdf}}

}

\caption{\label{fig-convergence}Absolute error versus number of panels,
log-log scale. The slopes are the exponents of \(1/n\): \(-2\) for
midpoint and trapezoid, \(-4\) for Simpson, until floating-point
rounding error dominates near \(10^{-15}\).}

\end{figure}%

For a smooth, unimodal integrand like a normal density, Simpson's rule
dominates. The rest of this chapter nonetheless uses the
\textbf{midpoint rule}, because the application of interest --- Bayesian
marginal likelihood after a logistic transform --- produces an integrand
that is undefined exactly at the endpoints \(0\) and \(1\), which only
the midpoint rule avoids evaluating.

\section{Application: The Marginal
Likelihood}\label{application-the-marginal-likelihood}

\subsection{Setup}\label{setup}

With likelihood \(L(\theta) = \prod_{i=1}^n P(y_i \mid \theta)\) and
prior \(\pi(\theta)\), Bayes' rule is \[
P(\theta \mid y) = \frac{L(\theta)\pi(\theta)}{P(y)}, \qquad
P(y) = \int L(\theta)\pi(\theta)\,d\theta = E_\pi[L(\theta)].
\] \(P(y)\), the \textbf{marginal likelihood} (the likelihood averaged
over the prior), is the normalizing constant of the posterior and the
basis of Bayes factors for comparing models. It has no closed form
outside conjugate settings, so we approximate it by quadrature. Because
it depends on the \emph{whole} prior --- not just where the likelihood
happens to be large --- its value is far more sensitive to prior choice
than a posterior mean is, as the next figure shows.

\begin{Shaded}
\begin{Highlighting}[]
\FunctionTok{par}\NormalTok{(}\AttributeTok{mfrow =} \FunctionTok{c}\NormalTok{(}\DecValTok{1}\NormalTok{, }\DecValTok{3}\NormalTok{), }\AttributeTok{mar =} \FunctionTok{c}\NormalTok{(}\DecValTok{4}\NormalTok{, }\DecValTok{4}\NormalTok{, }\DecValTok{2}\NormalTok{, }\DecValTok{1}\NormalTok{))}
\NormalTok{L }\OtherTok{\textless{}{-}} \ControlFlowTok{function}\NormalTok{(t) }\FunctionTok{dnorm}\NormalTok{(t, }\DecValTok{2}\NormalTok{, }\DecValTok{1}\NormalTok{)}
\NormalTok{priors }\OtherTok{\textless{}{-}} \FunctionTok{list}\NormalTok{(}\FunctionTok{c}\NormalTok{(}\DecValTok{2}\NormalTok{, }\FloatTok{1.5}\NormalTok{, }\StringTok{"prior near likelihood"}\NormalTok{),}
               \FunctionTok{c}\NormalTok{(}\SpecialCharTok{{-}}\DecValTok{5}\NormalTok{, }\DecValTok{1}\NormalTok{, }\StringTok{"prior far away"}\NormalTok{),}
               \FunctionTok{c}\NormalTok{(}\DecValTok{0}\NormalTok{, }\DecValTok{10}\NormalTok{, }\StringTok{"vague prior N(0, 100)"}\NormalTok{))}
\ControlFlowTok{for}\NormalTok{ (pr }\ControlFlowTok{in}\NormalTok{ priors) \{}
\NormalTok{    m }\OtherTok{\textless{}{-}} \FunctionTok{as.numeric}\NormalTok{(pr[}\DecValTok{1}\NormalTok{]); s }\OtherTok{\textless{}{-}} \FunctionTok{as.numeric}\NormalTok{(pr[}\DecValTok{2}\NormalTok{])}
\NormalTok{    pi\_fn }\OtherTok{\textless{}{-}} \ControlFlowTok{function}\NormalTok{(t) }\FunctionTok{dnorm}\NormalTok{(t, m, s)}
\NormalTok{    ml }\OtherTok{\textless{}{-}} \FunctionTok{integrate}\NormalTok{(}\ControlFlowTok{function}\NormalTok{(t) }\FunctionTok{L}\NormalTok{(t) }\SpecialCharTok{*} \FunctionTok{pi\_fn}\NormalTok{(t), }\SpecialCharTok{{-}}\ConstantTok{Inf}\NormalTok{, }\ConstantTok{Inf}\NormalTok{)}\SpecialCharTok{$}\NormalTok{value}
    \FunctionTok{curve}\NormalTok{(L, }\SpecialCharTok{{-}}\DecValTok{9}\NormalTok{, }\DecValTok{8}\NormalTok{, }\AttributeTok{lwd =} \DecValTok{3}\NormalTok{, }\AttributeTok{col =} \StringTok{"firebrick"}\NormalTok{, }\AttributeTok{ylim =} \FunctionTok{c}\NormalTok{(}\DecValTok{0}\NormalTok{, }\FloatTok{0.42}\NormalTok{),}
          \AttributeTok{xlab =} \FunctionTok{expression}\NormalTok{(theta), }\AttributeTok{ylab =} \StringTok{""}\NormalTok{,}
          \AttributeTok{main =} \FunctionTok{sprintf}\NormalTok{(}\StringTok{"\%s}\SpecialCharTok{\textbackslash{}n}\StringTok{P(y) = \%.4f"}\NormalTok{, pr[}\DecValTok{3}\NormalTok{], ml))}
    \FunctionTok{curve}\NormalTok{(pi\_fn, }\AttributeTok{add =} \ConstantTok{TRUE}\NormalTok{, }\AttributeTok{lwd =} \DecValTok{3}\NormalTok{, }\AttributeTok{col =} \StringTok{"steelblue"}\NormalTok{)}
\NormalTok{    ts }\OtherTok{\textless{}{-}} \FunctionTok{seq}\NormalTok{(}\SpecialCharTok{{-}}\DecValTok{9}\NormalTok{, }\DecValTok{8}\NormalTok{, }\AttributeTok{length.out =} \DecValTok{300}\NormalTok{)}
    \FunctionTok{polygon}\NormalTok{(}\FunctionTok{c}\NormalTok{(ts, }\FunctionTok{rev}\NormalTok{(ts)), }\FunctionTok{c}\NormalTok{(}\FunctionTok{L}\NormalTok{(ts) }\SpecialCharTok{*} \FunctionTok{pi\_fn}\NormalTok{(ts), }\FunctionTok{rep}\NormalTok{(}\DecValTok{0}\NormalTok{, }\DecValTok{300}\NormalTok{)),}
            \AttributeTok{col =} \FunctionTok{adjustcolor}\NormalTok{(}\StringTok{"gray30"}\NormalTok{, }\FloatTok{0.4}\NormalTok{), }\AttributeTok{border =} \ConstantTok{NA}\NormalTok{)}
\NormalTok{\}}
\end{Highlighting}
\end{Shaded}

\begin{figure}[H]

\centering{

\pandocbounded{\includegraphics[keepaspectratio]{chapters/unit08_integral_files/figure-pdf/fig-prior-sensitivity-1.pdf}}

}

\caption{\label{fig-prior-sensitivity}The same likelihood (red) combined
with three priors (blue). \(P(y)\) is the area under
\(L(\theta)\pi(\theta)\) (shaded), computed here by \texttt{integrate}
for a single parameter as a sanity check.}

\end{figure}%

A prior centred where the likelihood is large gives a much bigger
marginal likelihood than one that is far away or merely diffuse ---
which is exactly why \(P(y)\) is useful for comparing models, and why we
must be able to compute it even when the parameter is multivariate and
there is no \texttt{integrate}.

\subsection{A two-parameter example: mean and variance
unknown}\label{a-two-parameter-example-mean-and-variance-unknown}

Let
\(y_1,\ldots,y_n \mid \mu, \sigma \overset{iid}{\sim} t_\nu(\mu,\sigma)\),
a location-scale \(t\) distribution with \(\nu\) degrees of freedom (the
special case \(\nu=\infty\) is the normal). Reparametrize the scale as
\(w = \log \sigma\) so that both parameters range over
\((-\infty,\infty)\), and put independent normal priors on them: \[
\mu \sim N(\mu_0, \sigma_\mu^2), \qquad w \sim N(w_0, \sigma_w^2).
\] The marginal likelihood is the two-dimensional integral \[
P(y) = \int_{-\infty}^{\infty}\!\int_{-\infty}^{\infty}
 L(\mu, w)\,\pi(\mu, w)\, d\mu\, dw,
\qquad L(\mu,w) = \prod_{i=1}^n t_\nu\!\big(y_i;\ \mu,\ e^w\big),
\] which has no closed form because the normal prior on \(w\) is not
conjugate to the likelihood. We handle it with a \textbf{product
midpoint rule}: apply the logistic transform to each of \(\mu\) and
\(w\) separately, \[
\mu = \log\frac{u}{1-u}, \qquad w = \log\frac{v}{1-v}, \qquad
d\mu\,dw = \frac{du\,dv}{u(1-u)\,v(1-v)},
\] giving \[
P(y) = \int_0^1\!\int_0^1 g(u,v)\,du\,dv, \qquad
g(u,v) = \frac{L(\mu(u),w(v))\,\pi(\mu(u),w(v))}{u(1-u)\,v(1-v)},
\] and evaluate \(g\) on an \(n \times n\) grid \(u_j = (j-\tfrac12)h\),
\(v_k =
(k-\tfrac12)h\) with \(h=1/n\): \[
P(y) \approx h^2 \sum_{j=1}^n \sum_{k=1}^n g(u_j, v_k).
\] The midpoint rule again earns its place here:
\(1/\{u(1-u)v(1-v)\} \to \infty\) at the boundary of the square, and the
midpoint grid never touches it.

\subsection{Computing on the log
scale}\label{computing-on-the-log-scale}

\(L(\theta)\) is a product of \(n\) densities and underflows to exactly
\(0\) in floating point for even moderate \(n\), so every calculation
below is carried out on the log scale, and sums of exponentials are
combined with \textbf{log-sum-exp}: \[
\log \sum_i e^{\ell_i} = M + \log \sum_i e^{\ell_i - M}, \qquad M = \max_i \ell_i.
\] Subtracting the maximum before exponentiating makes the largest term
exactly \(e^0=1\), so nothing overflows; terms far below \(M\) underflow
harmlessly to zero instead of corrupting the sum.

\begin{Shaded}
\begin{Highlighting}[]
\NormalTok{log\_sum\_exp }\OtherTok{\textless{}{-}} \ControlFlowTok{function}\NormalTok{(lx) \{}
\NormalTok{    m }\OtherTok{\textless{}{-}} \FunctionTok{max}\NormalTok{(lx)}
\NormalTok{    m }\SpecialCharTok{+} \FunctionTok{log}\NormalTok{(}\FunctionTok{sum}\NormalTok{(}\FunctionTok{exp}\NormalTok{(lx }\SpecialCharTok{{-}}\NormalTok{ m)))}
\NormalTok{\}}
\end{Highlighting}
\end{Shaded}

\subsection{An efficient
implementation}\label{an-efficient-implementation}

The functions below compute \(\log g(u,v)\) on the whole grid at once
with \texttt{outer()}, rather than looping in R or nesting two calls to
a generic 1-D integrator (the nested approach used in earlier drafts of
this material is correct but re-evaluates the logistic transform and
re-dispatches a function call \(n^2\) times through two layers of
\texttt{sapply}; building the grid directly is both clearer and several
times faster for the grid sizes used here).

\begin{Shaded}
\begin{Highlighting}[]
\DocumentationTok{\#\#\# logistic link and its log{-}derivative, both vectorized}
\NormalTok{logit\_inv     }\OtherTok{\textless{}{-}} \ControlFlowTok{function}\NormalTok{(u) }\FunctionTok{log}\NormalTok{(u }\SpecialCharTok{/}\NormalTok{ (}\DecValTok{1} \SpecialCharTok{{-}}\NormalTok{ u))}
\NormalTok{log\_der\_logit }\OtherTok{\textless{}{-}} \ControlFlowTok{function}\NormalTok{(u) }\SpecialCharTok{{-}}\FunctionTok{log}\NormalTok{(u) }\SpecialCharTok{{-}} \FunctionTok{log}\NormalTok{(}\DecValTok{1} \SpecialCharTok{{-}}\NormalTok{ u)}

\DocumentationTok{\#\#\# log likelihood of a location{-}scale t sample (df = Inf is the normal)}
\NormalTok{log\_lik }\OtherTok{\textless{}{-}} \ControlFlowTok{function}\NormalTok{(x, mu, w, }\AttributeTok{df =} \ConstantTok{Inf}\NormalTok{) }\FunctionTok{sum}\NormalTok{(}\FunctionTok{dt.scaled}\NormalTok{(x, df, mu, }\FunctionTok{exp}\NormalTok{(w), }\AttributeTok{log =} \ConstantTok{TRUE}\NormalTok{))}

\DocumentationTok{\#\#\# log joint prior on (mu, w)}
\NormalTok{log\_prior }\OtherTok{\textless{}{-}} \ControlFlowTok{function}\NormalTok{(mu, w, mu\_0, sigma\_mu, w\_0, sigma\_w) \{}
    \FunctionTok{dnorm}\NormalTok{(mu, mu\_0, sigma\_mu, }\AttributeTok{log =} \ConstantTok{TRUE}\NormalTok{) }\SpecialCharTok{+} \FunctionTok{dnorm}\NormalTok{(w, w\_0, sigma\_w, }\AttributeTok{log =} \ConstantTok{TRUE}\NormalTok{)}
\NormalTok{\}}

\DocumentationTok{\#\#\# log of the transformed integrand g(u, v) at one grid point}
\NormalTok{log\_g }\OtherTok{\textless{}{-}} \ControlFlowTok{function}\NormalTok{(u, v, x, mu\_0, sigma\_mu, w\_0, sigma\_w, }\AttributeTok{df =} \ConstantTok{Inf}\NormalTok{) \{}
\NormalTok{    mu }\OtherTok{\textless{}{-}} \FunctionTok{logit\_inv}\NormalTok{(u); w }\OtherTok{\textless{}{-}} \FunctionTok{logit\_inv}\NormalTok{(v)}
    \FunctionTok{log\_lik}\NormalTok{(x, mu, w, df) }\SpecialCharTok{+} \FunctionTok{log\_prior}\NormalTok{(mu, w, mu\_0, sigma\_mu, w\_0, sigma\_w) }\SpecialCharTok{+}
        \FunctionTok{log\_der\_logit}\NormalTok{(u) }\SpecialCharTok{+} \FunctionTok{log\_der\_logit}\NormalTok{(v)}
\NormalTok{\}}

\DocumentationTok{\#\#\# log marginal likelihood by the product midpoint rule on an n x n grid}
\NormalTok{log\_marlik\_mid }\OtherTok{\textless{}{-}} \ControlFlowTok{function}\NormalTok{(x, mu\_0, sigma\_mu, w\_0, sigma\_w, n, }\AttributeTok{df =} \ConstantTok{Inf}\NormalTok{) \{}
\NormalTok{    h }\OtherTok{\textless{}{-}} \DecValTok{1} \SpecialCharTok{/}\NormalTok{ n}
\NormalTok{    grid }\OtherTok{\textless{}{-}}\NormalTok{ (}\FunctionTok{seq\_len}\NormalTok{(n) }\SpecialCharTok{{-}} \FloatTok{0.5}\NormalTok{) }\SpecialCharTok{*}\NormalTok{ h}
\NormalTok{    G }\OtherTok{\textless{}{-}} \FunctionTok{outer}\NormalTok{(grid, grid, }\FunctionTok{Vectorize}\NormalTok{(}\ControlFlowTok{function}\NormalTok{(u, v)}
        \FunctionTok{log\_g}\NormalTok{(u, v, x, mu\_0, sigma\_mu, w\_0, sigma\_w, df)))}
    \FunctionTok{log\_sum\_exp}\NormalTok{(G) }\SpecialCharTok{+} \DecValTok{2} \SpecialCharTok{*} \FunctionTok{log}\NormalTok{(h)}
\NormalTok{\}}

\DocumentationTok{\#\#\# log marginal likelihood by naive Monte Carlo (for validating the above)}
\NormalTok{log\_marlik\_mc }\OtherTok{\textless{}{-}} \ControlFlowTok{function}\NormalTok{(x, mu\_0, sigma\_mu, w\_0, sigma\_w, iters\_mc, }\AttributeTok{df =} \ConstantTok{Inf}\NormalTok{) \{}
\NormalTok{    mus }\OtherTok{\textless{}{-}} \FunctionTok{rnorm}\NormalTok{(iters\_mc, mu\_0, sigma\_mu)}
\NormalTok{    ws  }\OtherTok{\textless{}{-}} \FunctionTok{rnorm}\NormalTok{(iters\_mc, w\_0, sigma\_w)}
\NormalTok{    v\_log\_lik }\OtherTok{\textless{}{-}} \FunctionTok{vapply}\NormalTok{(}\FunctionTok{seq\_len}\NormalTok{(iters\_mc), }\ControlFlowTok{function}\NormalTok{(i) }\FunctionTok{log\_lik}\NormalTok{(x, mus[i], ws[i], df), }\FunctionTok{numeric}\NormalTok{(}\DecValTok{1}\NormalTok{))}
    \FunctionTok{log\_sum\_exp}\NormalTok{(v\_log\_lik) }\SpecialCharTok{{-}} \FunctionTok{log}\NormalTok{(iters\_mc)}
\NormalTok{\}}
\end{Highlighting}
\end{Shaded}

\texttt{log\_marlik\_mc} draws \(\theta\) from the prior and averages
the likelihood --- a direct Monte Carlo estimate of \(E_\pi[L(\theta)]\)
--- and serves purely as an independent check on the quadrature answer,
since it makes no use of a grid or a transformation and so cannot share
their failure modes.

\subsection{Visualizing the grid and the
integrand}\label{visualizing-the-grid-and-the-integrand}

\begin{Shaded}
\begin{Highlighting}[]
\FunctionTok{set.seed}\NormalTok{(}\DecValTok{1}\NormalTok{)}
\NormalTok{y }\OtherTok{\textless{}{-}} \FunctionTok{rnorm}\NormalTok{(}\DecValTok{20}\NormalTok{, }\DecValTok{1}\NormalTok{, }\DecValTok{2}\NormalTok{)}
\NormalTok{n }\OtherTok{\textless{}{-}} \DecValTok{20}\NormalTok{; h }\OtherTok{\textless{}{-}} \DecValTok{1} \SpecialCharTok{/}\NormalTok{ n; grid }\OtherTok{\textless{}{-}}\NormalTok{ (}\FunctionTok{seq\_len}\NormalTok{(n) }\SpecialCharTok{{-}} \FloatTok{0.5}\NormalTok{) }\SpecialCharTok{*}\NormalTok{ h}
\NormalTok{G  }\OtherTok{\textless{}{-}} \FunctionTok{outer}\NormalTok{(grid, grid, }\FunctionTok{Vectorize}\NormalTok{(}\ControlFlowTok{function}\NormalTok{(u, v) }\FunctionTok{log\_g}\NormalTok{(u, v, y, }\DecValTok{0}\NormalTok{, }\DecValTok{5}\NormalTok{, }\DecValTok{0}\NormalTok{, }\DecValTok{2}\NormalTok{)))}

\NormalTok{mg }\OtherTok{\textless{}{-}} \FunctionTok{seq}\NormalTok{(}\SpecialCharTok{{-}}\DecValTok{6}\NormalTok{, }\DecValTok{6}\NormalTok{, }\AttributeTok{length.out =} \DecValTok{80}\NormalTok{); wg }\OtherTok{\textless{}{-}} \FunctionTok{seq}\NormalTok{(}\SpecialCharTok{{-}}\DecValTok{3}\NormalTok{, }\DecValTok{4}\NormalTok{, }\AttributeTok{length.out =} \DecValTok{80}\NormalTok{)}
\NormalTok{Gm }\OtherTok{\textless{}{-}} \FunctionTok{outer}\NormalTok{(mg, wg, }\FunctionTok{Vectorize}\NormalTok{(}\ControlFlowTok{function}\NormalTok{(m, w) }\FunctionTok{log\_lik}\NormalTok{(y, m, w) }\SpecialCharTok{+} \FunctionTok{log\_prior}\NormalTok{(m, w, }\DecValTok{0}\NormalTok{, }\DecValTok{5}\NormalTok{, }\DecValTok{0}\NormalTok{, }\DecValTok{2}\NormalTok{)))}

\FunctionTok{par}\NormalTok{(}\AttributeTok{mfrow =} \FunctionTok{c}\NormalTok{(}\DecValTok{1}\NormalTok{, }\DecValTok{2}\NormalTok{), }\AttributeTok{mar =} \FunctionTok{c}\NormalTok{(}\DecValTok{4}\NormalTok{, }\DecValTok{4}\NormalTok{, }\DecValTok{2}\NormalTok{, }\DecValTok{1}\NormalTok{))}
\FunctionTok{contour}\NormalTok{(mg, wg, }\FunctionTok{exp}\NormalTok{(Gm }\SpecialCharTok{{-}} \FunctionTok{max}\NormalTok{(Gm)), }\AttributeTok{nlevels =} \DecValTok{8}\NormalTok{,}
        \AttributeTok{xlab =} \FunctionTok{expression}\NormalTok{(mu), }\AttributeTok{ylab =} \FunctionTok{expression}\NormalTok{(w }\SpecialCharTok{==}\NormalTok{ log}\SpecialCharTok{\textasciitilde{}}\NormalTok{sigma}\SpecialCharTok{\^{}}\DecValTok{2}\NormalTok{),}
        \AttributeTok{main =} \FunctionTok{expression}\NormalTok{(}\FunctionTok{L}\NormalTok{(mu, w)}\SpecialCharTok{\textasciitilde{}}\FunctionTok{pi}\NormalTok{(mu, w)))}
\FunctionTok{image}\NormalTok{(grid, grid, }\FunctionTok{exp}\NormalTok{(G }\SpecialCharTok{{-}} \FunctionTok{max}\NormalTok{(G)), }\AttributeTok{col =} \FunctionTok{hcl.colors}\NormalTok{(}\DecValTok{30}\NormalTok{, }\StringTok{"YlOrRd"}\NormalTok{, }\AttributeTok{rev =} \ConstantTok{TRUE}\NormalTok{),}
      \AttributeTok{xlab =} \StringTok{"u"}\NormalTok{, }\AttributeTok{ylab =} \StringTok{"v"}\NormalTok{, }\AttributeTok{main =} \StringTok{"g(u, v) with midpoint grid"}\NormalTok{)}
\FunctionTok{abline}\NormalTok{(}\AttributeTok{v =} \FunctionTok{seq}\NormalTok{(}\DecValTok{0}\NormalTok{, }\DecValTok{1}\NormalTok{, }\AttributeTok{by =}\NormalTok{ h), }\AttributeTok{h =} \FunctionTok{seq}\NormalTok{(}\DecValTok{0}\NormalTok{, }\DecValTok{1}\NormalTok{, }\AttributeTok{by =}\NormalTok{ h), }\AttributeTok{col =} \FunctionTok{adjustcolor}\NormalTok{(}\StringTok{"gray30"}\NormalTok{, }\FloatTok{0.4}\NormalTok{))}
\FunctionTok{points}\NormalTok{(}\FunctionTok{expand.grid}\NormalTok{(grid, grid), }\AttributeTok{pch =} \StringTok{"."}\NormalTok{, }\AttributeTok{cex =} \FloatTok{1.5}\NormalTok{)}
\end{Highlighting}
\end{Shaded}

\begin{figure}[H]

\centering{

\pandocbounded{\includegraphics[keepaspectratio]{chapters/unit08_integral_files/figure-pdf/fig-grid-1.pdf}}

}

\caption{\label{fig-grid}Left: the posterior \(L(\mu,w)\pi(\mu,w)\) on
its natural scale. Right: the transformed integrand \(g(u,v)\) on the
unit square, with the midpoint grid at \(n=20\) overlaid. Equal
probability mass under the two surfaces corresponds to equal volume,
even though the shapes differ.}

\end{figure}%

\subsection{Checking the quadrature against Monte
Carlo}\label{checking-the-quadrature-against-monte-carlo}

\begin{Shaded}
\begin{Highlighting}[]
\FunctionTok{set.seed}\NormalTok{(}\DecValTok{2}\NormalTok{)}
\NormalTok{x1 }\OtherTok{\textless{}{-}} \FunctionTok{rt.scaled}\NormalTok{(}\DecValTok{50}\NormalTok{, }\AttributeTok{mean =} \DecValTok{2}\NormalTok{, }\AttributeTok{sd =} \DecValTok{2}\NormalTok{, }\AttributeTok{df =} \DecValTok{2}\NormalTok{)}
\NormalTok{x2 }\OtherTok{\textless{}{-}} \FunctionTok{rt.scaled}\NormalTok{(}\DecValTok{100}\NormalTok{, }\AttributeTok{mean =} \DecValTok{2}\NormalTok{, }\AttributeTok{sd =} \DecValTok{2}\NormalTok{, }\AttributeTok{df =} \ConstantTok{Inf}\NormalTok{)}

\NormalTok{comparison }\OtherTok{\textless{}{-}} \FunctionTok{data.frame}\NormalTok{(}
    \AttributeTok{dataset =} \FunctionTok{c}\NormalTok{(}\StringTok{"n=50, t (df=2) data"}\NormalTok{, }\StringTok{"n=100, normal data"}\NormalTok{),}
    \AttributeTok{midpoint\_n100 =} \FunctionTok{c}\NormalTok{(}\FunctionTok{log\_marlik\_mid}\NormalTok{(x1, }\DecValTok{0}\NormalTok{, }\DecValTok{10}\NormalTok{, }\DecValTok{0}\NormalTok{, }\DecValTok{10}\NormalTok{, }\DecValTok{100}\NormalTok{),}
                       \FunctionTok{log\_marlik\_mid}\NormalTok{(x2, }\DecValTok{0}\NormalTok{, }\DecValTok{10}\NormalTok{, }\DecValTok{0}\NormalTok{, }\DecValTok{10}\NormalTok{, }\DecValTok{100}\NormalTok{)),}
    \AttributeTok{monte\_carlo   =} \FunctionTok{c}\NormalTok{(}\FunctionTok{log\_marlik\_mc}\NormalTok{(x1, }\DecValTok{0}\NormalTok{, }\DecValTok{10}\NormalTok{, }\DecValTok{0}\NormalTok{, }\DecValTok{10}\NormalTok{, }\FloatTok{1e5}\NormalTok{),}
                       \FunctionTok{log\_marlik\_mc}\NormalTok{(x2, }\DecValTok{0}\NormalTok{, }\DecValTok{10}\NormalTok{, }\DecValTok{0}\NormalTok{, }\DecValTok{10}\NormalTok{, }\FloatTok{1e5}\NormalTok{))}
\NormalTok{)}
\NormalTok{knitr}\SpecialCharTok{::}\FunctionTok{kable}\NormalTok{(comparison, }\AttributeTok{digits =} \DecValTok{4}\NormalTok{)}
\end{Highlighting}
\end{Shaded}

{

\begin{longtable}[]{@{}lrr@{}}

\caption{\label{tbl-quad-vs-mc}Estimated log marginal likelihood:
product midpoint rule versus naive Monte Carlo, for two simulated data
sets.}

\tabularnewline

\toprule\noalign{}
dataset & midpoint\_n100 & monte\_carlo \\
\midrule\noalign{}
\endhead
\bottomrule\noalign{}
\endlastfoot
n=50, t (df=2) data & -173.6227 & -173.5375 \\
n=100, normal data & -221.7068 & -221.6300 \\

\end{longtable}

}

The two independent methods agree to well within a tenth of a log-unit,
which is reassuring: they share no code path, so agreement is evidence
that neither the grid nor the Monte Carlo sample is too coarse for this
posterior.

\subsection{Convergence as the grid is
refined}\label{convergence-as-the-grid-is-refined}

\begin{Shaded}
\begin{Highlighting}[]
\NormalTok{ns\_grid }\OtherTok{\textless{}{-}} \FunctionTok{seq}\NormalTok{(}\DecValTok{10}\NormalTok{, }\DecValTok{90}\NormalTok{, }\AttributeTok{by =} \DecValTok{10}\NormalTok{)}
\NormalTok{est }\OtherTok{\textless{}{-}} \FunctionTok{vapply}\NormalTok{(ns\_grid, }\ControlFlowTok{function}\NormalTok{(k) }\FunctionTok{log\_marlik\_mid}\NormalTok{(x2, }\DecValTok{0}\NormalTok{, }\DecValTok{10}\NormalTok{, }\DecValTok{0}\NormalTok{, }\DecValTok{10}\NormalTok{, k), }\FunctionTok{numeric}\NormalTok{(}\DecValTok{1}\NormalTok{))}
\NormalTok{knitr}\SpecialCharTok{::}\FunctionTok{kable}\NormalTok{(}\FunctionTok{data.frame}\NormalTok{(}\StringTok{\textasciigrave{}}\AttributeTok{grid points per axis}\StringTok{\textasciigrave{}} \OtherTok{=}\NormalTok{ ns\_grid, }\StringTok{\textasciigrave{}}\AttributeTok{log P(y)}\StringTok{\textasciigrave{}} \OtherTok{=}\NormalTok{ est,}
                         \AttributeTok{check.names =} \ConstantTok{FALSE}\NormalTok{), }\AttributeTok{digits =} \DecValTok{4}\NormalTok{)}
\end{Highlighting}
\end{Shaded}

{

\begin{longtable}[]{@{}rr@{}}

\caption{\label{tbl-quad-convergence}Estimated log marginal likelihood
as the number of grid points per axis increases, for the \(n=100\)
normal-data example. By \(n=40\) per axis the estimate has stabilized to
the precision shown.}

\tabularnewline

\toprule\noalign{}
grid points per axis & log P(y) \\
\midrule\noalign{}
\endhead
\bottomrule\noalign{}
\endlastfoot
10 & -222.7297 \\
20 & -221.5545 \\
30 & -221.7318 \\
40 & -221.7069 \\
50 & -221.7067 \\
60 & -221.7068 \\
70 & -221.7068 \\
80 & -221.7068 \\
90 & -221.7068 \\

\end{longtable}

}

\subsection{Comparing priors}\label{comparing-priors}

Holding the data fixed, the marginal likelihood falls as the prior on
\(\mu\) either narrows around the wrong value or spreads out
indefinitely:

\begin{Shaded}
\begin{Highlighting}[]
\FunctionTok{set.seed}\NormalTok{(}\DecValTok{3}\NormalTok{)}
\NormalTok{x }\OtherTok{\textless{}{-}} \FunctionTok{rnorm}\NormalTok{(}\DecValTok{100}\NormalTok{)}
\NormalTok{prior\_grid }\OtherTok{\textless{}{-}} \FunctionTok{expand.grid}\NormalTok{(}\AttributeTok{mu\_0 =} \FunctionTok{c}\NormalTok{(}\DecValTok{0}\NormalTok{, }\SpecialCharTok{{-}}\DecValTok{5}\NormalTok{), }\AttributeTok{sigma\_mu =} \FunctionTok{c}\NormalTok{(}\FloatTok{0.01}\NormalTok{, }\FloatTok{0.1}\NormalTok{, }\DecValTok{1}\NormalTok{, }\DecValTok{10}\NormalTok{, }\DecValTok{100}\NormalTok{, }\DecValTok{1000}\NormalTok{))}
\NormalTok{prior\_grid}\SpecialCharTok{$}\NormalTok{log\_P\_y }\OtherTok{\textless{}{-}} \FunctionTok{mapply}\NormalTok{(}\ControlFlowTok{function}\NormalTok{(m, s) }\FunctionTok{log\_marlik\_mid}\NormalTok{(x, m, s, }\DecValTok{0}\NormalTok{, }\DecValTok{1}\NormalTok{, }\DecValTok{100}\NormalTok{),}
\NormalTok{                              prior\_grid}\SpecialCharTok{$}\NormalTok{mu\_0, prior\_grid}\SpecialCharTok{$}\NormalTok{sigma\_mu)}
\NormalTok{knitr}\SpecialCharTok{::}\FunctionTok{kable}\NormalTok{(prior\_grid, }\AttributeTok{digits =} \DecValTok{2}\NormalTok{,}
             \AttributeTok{col.names =} \FunctionTok{c}\NormalTok{(}\StringTok{"prior mean $}\SpecialCharTok{\textbackslash{}\textbackslash{}}\StringTok{mu\_0$"}\NormalTok{, }\StringTok{"prior sd $}\SpecialCharTok{\textbackslash{}\textbackslash{}}\StringTok{sigma\_}\SpecialCharTok{\textbackslash{}\textbackslash{}}\StringTok{mu$"}\NormalTok{, }\StringTok{"log P(y)"}\NormalTok{))}
\end{Highlighting}
\end{Shaded}

{

\begin{longtable}[]{@{}rrr@{}}

\caption{\label{tbl-compare-priors}Log marginal likelihood for data
centred at 0, under priors on mu with the stated mean and standard
deviation (w-prior fixed at N(0,1)).}

\tabularnewline

\toprule\noalign{}
prior mean \(\mu_0\) & prior sd \(\sigma_\mu\) & log P(y) \\
\midrule\noalign{}
\endhead
\bottomrule\noalign{}
\endlastfoot
0 & 1e-02 & -129.39 \\
-5 & 1e-02 & -739.84 \\
0 & 1e-01 & -128.95 \\
-5 & 1e-01 & -316.30 \\
0 & 1e+00 & -130.97 \\
-5 & 1e+00 & -143.43 \\
0 & 1e+01 & -133.27 \\
-5 & 1e+01 & -133.40 \\
0 & 1e+02 & -135.57 \\
-5 & 1e+02 & -135.57 \\
0 & 1e+03 & -137.88 \\
-5 & 1e+03 & -137.88 \\

\end{longtable}

}

With \(\mu_0 = 0\) (the right neighbourhood for data centred at \(0\)),
the marginal likelihood is highest for a moderately informative prior
and degrades slowly as the prior widens. With \(\mu_0 = -5\) (the wrong
neighbourhood), a \emph{narrow} prior is actively harmful --- it insists
on a region the data reject --- while a wide prior recovers because it
eventually puts enough mass back near the truth. This is the general
shape of Bayesian model comparison: strong, correct priors win; strong,
wrong priors lose; vague priors are safe but not optimal.

\subsection{Comparing models for the
data}\label{comparing-models-for-the-data}

The marginal likelihood also compares different likelihoods for the same
data --- here, whether \(\nu\) (the degrees of freedom of the \(t\)
model) matches how the data were actually generated:

\begin{Shaded}
\begin{Highlighting}[]
\NormalTok{dfs }\OtherTok{\textless{}{-}} \FunctionTok{c}\NormalTok{(}\ConstantTok{Inf}\NormalTok{, }\DecValTok{2}\NormalTok{, }\DecValTok{1}\NormalTok{, }\FloatTok{0.5}\NormalTok{)}
\FunctionTok{set.seed}\NormalTok{(}\DecValTok{4}\NormalTok{)}
\NormalTok{x\_normal }\OtherTok{\textless{}{-}} \FunctionTok{rt.scaled}\NormalTok{(}\DecValTok{100}\NormalTok{, }\AttributeTok{mean =} \DecValTok{2}\NormalTok{, }\AttributeTok{sd =} \DecValTok{2}\NormalTok{, }\AttributeTok{df =} \ConstantTok{Inf}\NormalTok{)}
\NormalTok{x\_heavy  }\OtherTok{\textless{}{-}} \FunctionTok{rt.scaled}\NormalTok{(}\DecValTok{100}\NormalTok{, }\AttributeTok{mean =} \DecValTok{2}\NormalTok{, }\AttributeTok{sd =} \DecValTok{2}\NormalTok{, }\AttributeTok{df =} \DecValTok{2}\NormalTok{)}
\NormalTok{model\_compare }\OtherTok{\textless{}{-}} \FunctionTok{data.frame}\NormalTok{(}
    \AttributeTok{df\_fitted        =}\NormalTok{ dfs,}
    \AttributeTok{normal\_data       =} \FunctionTok{vapply}\NormalTok{(dfs, }\ControlFlowTok{function}\NormalTok{(d) }\FunctionTok{log\_marlik\_mid}\NormalTok{(x\_normal, }\DecValTok{0}\NormalTok{, }\DecValTok{10}\NormalTok{, }\DecValTok{0}\NormalTok{, }\DecValTok{10}\NormalTok{, }\DecValTok{100}\NormalTok{, }\AttributeTok{df =}\NormalTok{ d), }\FunctionTok{numeric}\NormalTok{(}\DecValTok{1}\NormalTok{)),}
    \AttributeTok{heavy\_tailed\_data =} \FunctionTok{vapply}\NormalTok{(dfs, }\ControlFlowTok{function}\NormalTok{(d) }\FunctionTok{log\_marlik\_mid}\NormalTok{(x\_heavy,  }\DecValTok{0}\NormalTok{, }\DecValTok{10}\NormalTok{, }\DecValTok{0}\NormalTok{, }\DecValTok{10}\NormalTok{, }\DecValTok{100}\NormalTok{, }\AttributeTok{df =}\NormalTok{ d), }\FunctionTok{numeric}\NormalTok{(}\DecValTok{1}\NormalTok{))}
\NormalTok{)}
\NormalTok{knitr}\SpecialCharTok{::}\FunctionTok{kable}\NormalTok{(model\_compare, }\AttributeTok{digits =} \DecValTok{3}\NormalTok{,}
             \AttributeTok{col.names =} \FunctionTok{c}\NormalTok{(}\StringTok{"fitted df"}\NormalTok{, }\StringTok{"data simulated as normal"}\NormalTok{, }\StringTok{"data simulated as t(df=2)"}\NormalTok{))}
\end{Highlighting}
\end{Shaded}

{

\begin{longtable}[]{@{}rrr@{}}

\caption{\label{tbl-compare-models}Log marginal likelihood under four
values of the t-model degrees of freedom nu, fit to data simulated from
a Gaussian and from a heavy-tailed t(df=2).}

\tabularnewline

\toprule\noalign{}
fitted df & data simulated as normal & data simulated as t(df=2) \\
\midrule\noalign{}
\endhead
\bottomrule\noalign{}
\endlastfoot
Inf & -210.675 & -300.951 \\
2.0 & -220.056 & -261.663 \\
1.0 & -231.060 & -265.646 \\
0.5 & -251.042 & -281.555 \\

\end{longtable}

}

Each column favours the \(\nu\) that actually generated it: normal data
are best explained by \(\nu=\infty\), and \(t(\text{df}=2)\) data are
best explained by \(\nu=2\), with the wrong tail weight costing several
log-likelihood units either way. As the earlier prior comparison showed,
the prior also shifts these numbers --- but the \emph{ranking} of the
correct model over the incorrect ones survives even under a badly
misspecified prior on \((\mu,w)\), which is why marginal likelihoods
remain useful for model comparison despite their prior-sensitivity.

\section{Difficulties of Numerical
Quadrature}\label{difficulties-of-numerical-quadrature}

\subsection{Peaked integrands can be missed by a fixed
grid}\label{peaked-integrands-can-be-missed-by-a-fixed-grid}

A grid of fixed width \(h_n\) can straddle a narrow spike without ever
landing a point near its peak, badly under- or over-estimating the area
it represents.

\begin{Shaded}
\begin{Highlighting}[]
\FunctionTok{par}\NormalTok{(}\AttributeTok{mfrow =} \FunctionTok{c}\NormalTok{(}\DecValTok{1}\NormalTok{, }\DecValTok{2}\NormalTok{), }\AttributeTok{mar =} \FunctionTok{c}\NormalTok{(}\DecValTok{4}\NormalTok{, }\DecValTok{4}\NormalTok{, }\DecValTok{2}\NormalTok{, }\DecValTok{1}\NormalTok{))}
\NormalTok{a }\OtherTok{\textless{}{-}} \SpecialCharTok{{-}}\DecValTok{3}\NormalTok{; b }\OtherTok{\textless{}{-}} \DecValTok{3}\NormalTok{; n\_panels }\OtherTok{\textless{}{-}} \DecValTok{8}\NormalTok{; h }\OtherTok{\textless{}{-}}\NormalTok{ (b }\SpecialCharTok{{-}}\NormalTok{ a) }\SpecialCharTok{/}\NormalTok{ n\_panels}
\NormalTok{xm }\OtherTok{\textless{}{-}}\NormalTok{ a }\SpecialCharTok{+}\NormalTok{ (}\FunctionTok{seq\_len}\NormalTok{(n\_panels) }\SpecialCharTok{{-}} \FloatTok{0.5}\NormalTok{) }\SpecialCharTok{*}\NormalTok{ h}
\ControlFlowTok{for}\NormalTok{ (s }\ControlFlowTok{in} \FunctionTok{c}\NormalTok{(}\DecValTok{1}\NormalTok{, }\FloatTok{0.03}\NormalTok{)) \{}
\NormalTok{    f }\OtherTok{\textless{}{-}} \ControlFlowTok{function}\NormalTok{(x) }\FunctionTok{dnorm}\NormalTok{(x, }\FloatTok{0.2}\NormalTok{, s)}
    \FunctionTok{curve}\NormalTok{(f, a, b, }\AttributeTok{n =} \DecValTok{2001}\NormalTok{, }\AttributeTok{lwd =} \DecValTok{3}\NormalTok{, }\AttributeTok{xlab =} \StringTok{"x"}\NormalTok{, }\AttributeTok{ylab =} \StringTok{"f(x)"}\NormalTok{,}
          \AttributeTok{main =} \FunctionTok{sprintf}\NormalTok{(}\StringTok{"sd = \%g: midpoint = \%.3f (exact = 1)"}\NormalTok{, s, h }\SpecialCharTok{*} \FunctionTok{sum}\NormalTok{(}\FunctionTok{f}\NormalTok{(xm))))}
    \FunctionTok{rect}\NormalTok{(xm }\SpecialCharTok{{-}}\NormalTok{ h }\SpecialCharTok{/} \DecValTok{2}\NormalTok{, }\DecValTok{0}\NormalTok{, xm }\SpecialCharTok{+}\NormalTok{ h }\SpecialCharTok{/} \DecValTok{2}\NormalTok{, }\FunctionTok{f}\NormalTok{(xm), }\AttributeTok{col =} \FunctionTok{adjustcolor}\NormalTok{(}\StringTok{"steelblue"}\NormalTok{, }\FloatTok{0.3}\NormalTok{), }\AttributeTok{border =} \StringTok{"steelblue"}\NormalTok{)}
    \FunctionTok{points}\NormalTok{(xm, }\FunctionTok{f}\NormalTok{(xm), }\AttributeTok{pch =} \DecValTok{19}\NormalTok{, }\AttributeTok{col =} \StringTok{"steelblue"}\NormalTok{)}
\NormalTok{\}}
\end{Highlighting}
\end{Shaded}

\begin{figure}[H]

\centering{

\pandocbounded{\includegraphics[keepaspectratio]{chapters/unit08_integral_files/figure-pdf/fig-peaked-1.pdf}}

}

\caption{\label{fig-peaked}Same rule, same 8-panel grid: the wide
density (left) integrates almost exactly, while the narrow spike (right)
is nearly invisible to the grid because no panel midpoint falls near its
mode.}

\end{figure}%

This is exactly the situation a Bayesian posterior creates as \(n\)
grows: its width shrinks like \(O(1/\sqrt n)\), so a grid tuned for
\(n=10\) observations can become useless for \(n=10{,}000\). The usual
remedy is to find the mode first (e.g.~by the optimization methods of
the previous two chapters), then centre and scale the grid around it ---
which is also the starting point for the Laplace approximation in the
next chapter.

\subsection{Transformed integrands can be squeezed against the
boundary}\label{transformed-integrands-can-be-squeezed-against-the-boundary}

The logistic transform has its own failure mode: if \(\theta\)'s
posterior sits at a large \(|\theta|\), it gets squeezed against \(y=0\)
or \(y=1\). For \(\theta \approx 50\),
\(y = 1/(1+e^{-50}) \approx 1 - 2\times 10^{-22}\) --- the entire
posterior mass occupies a sliver of \((0,1)\) far too thin for any
practical uniform grid to resolve.

\begin{Shaded}
\begin{Highlighting}[]
\FunctionTok{par}\NormalTok{(}\AttributeTok{mfrow =} \FunctionTok{c}\NormalTok{(}\DecValTok{1}\NormalTok{, }\DecValTok{2}\NormalTok{), }\AttributeTok{mar =} \FunctionTok{c}\NormalTok{(}\DecValTok{4}\NormalTok{, }\DecValTok{4}\NormalTok{, }\DecValTok{2}\NormalTok{, }\DecValTok{1}\NormalTok{))}
\NormalTok{f }\OtherTok{\textless{}{-}} \ControlFlowTok{function}\NormalTok{(x) }\FunctionTok{dnorm}\NormalTok{(x, }\DecValTok{8}\NormalTok{, }\DecValTok{1}\NormalTok{)}
\FunctionTok{curve}\NormalTok{(f, }\SpecialCharTok{{-}}\DecValTok{2}\NormalTok{, }\DecValTok{12}\NormalTok{, }\AttributeTok{lwd =} \DecValTok{3}\NormalTok{, }\AttributeTok{xlab =} \FunctionTok{expression}\NormalTok{(theta), }\AttributeTok{ylab =} \StringTok{""}\NormalTok{,}
      \AttributeTok{main =} \FunctionTok{expression}\NormalTok{(}\StringTok{"posterior at"} \SpecialCharTok{\textasciitilde{}}\NormalTok{ theta }\SpecialCharTok{==} \DecValTok{8}\NormalTok{))}
\NormalTok{g }\OtherTok{\textless{}{-}} \ControlFlowTok{function}\NormalTok{(y) }\FunctionTok{f}\NormalTok{(}\FunctionTok{log}\NormalTok{(y }\SpecialCharTok{/}\NormalTok{ (}\DecValTok{1} \SpecialCharTok{{-}}\NormalTok{ y))) }\SpecialCharTok{/}\NormalTok{ (y }\SpecialCharTok{*}\NormalTok{ (}\DecValTok{1} \SpecialCharTok{{-}}\NormalTok{ y))}
\FunctionTok{curve}\NormalTok{(g, }\FloatTok{0.001}\NormalTok{, }\FloatTok{0.9999}\NormalTok{, }\AttributeTok{n =} \DecValTok{5001}\NormalTok{, }\AttributeTok{lwd =} \DecValTok{3}\NormalTok{, }\AttributeTok{xlab =} \StringTok{"y"}\NormalTok{, }\AttributeTok{ylab =} \StringTok{""}\NormalTok{,}
      \AttributeTok{main =} \StringTok{"after logistic transform"}\NormalTok{)}
\FunctionTok{abline}\NormalTok{(}\AttributeTok{v =} \FunctionTok{seq}\NormalTok{(}\FloatTok{0.025}\NormalTok{, }\FloatTok{0.975}\NormalTok{, }\AttributeTok{by =} \FloatTok{0.05}\NormalTok{), }\AttributeTok{col =} \FunctionTok{adjustcolor}\NormalTok{(}\StringTok{"steelblue"}\NormalTok{, }\FloatTok{0.6}\NormalTok{), }\AttributeTok{lty =} \DecValTok{3}\NormalTok{)}
\end{Highlighting}
\end{Shaded}

\begin{figure}[H]

\centering{

\pandocbounded{\includegraphics[keepaspectratio]{chapters/unit08_integral_files/figure-pdf/fig-boundary-squeeze-1.pdf}}

}

\caption{\label{fig-boundary-squeeze}A posterior located at theta=8
already occupies a visibly asymmetric region after the logistic
transform (right); the grid lines (n=20) are far more informative on the
left than the right.}

\end{figure}%

We already met the extreme version of this problem at the start of this
chapter's development: with data simulated far from where the
untransformed grid was designed to look, the midpoint rule and Monte
Carlo disagree badly.

\begin{Shaded}
\begin{Highlighting}[]
\FunctionTok{set.seed}\NormalTok{(}\DecValTok{5}\NormalTok{)}
\NormalTok{x\_far }\OtherTok{\textless{}{-}} \FunctionTok{rt.scaled}\NormalTok{(}\DecValTok{100}\NormalTok{, }\AttributeTok{mean =} \DecValTok{50}\NormalTok{, }\AttributeTok{sd =} \DecValTok{2}\NormalTok{, }\AttributeTok{df =} \ConstantTok{Inf}\NormalTok{)}
\NormalTok{knitr}\SpecialCharTok{::}\FunctionTok{kable}\NormalTok{(}\FunctionTok{data.frame}\NormalTok{(}
    \AttributeTok{method =} \FunctionTok{c}\NormalTok{(}\StringTok{"midpoint (n = 100)"}\NormalTok{, }\StringTok{"Monte Carlo (n = 10,000)"}\NormalTok{),}
    \AttributeTok{log\_P\_y =} \FunctionTok{c}\NormalTok{(}\FunctionTok{log\_marlik\_mid}\NormalTok{(x\_far, }\DecValTok{0}\NormalTok{, }\DecValTok{100}\NormalTok{, }\DecValTok{0}\NormalTok{, }\DecValTok{100}\NormalTok{, }\DecValTok{100}\NormalTok{),}
                \FunctionTok{log\_marlik\_mc}\NormalTok{(x\_far, }\DecValTok{0}\NormalTok{, }\DecValTok{100}\NormalTok{, }\DecValTok{0}\NormalTok{, }\DecValTok{100}\NormalTok{, }\FloatTok{1e4}\NormalTok{))}
\NormalTok{), }\AttributeTok{digits =} \DecValTok{3}\NormalTok{)}
\end{Highlighting}
\end{Shaded}

{

\begin{longtable}[]{@{}lr@{}}

\caption{\label{tbl-quad-failure}Quadrature failure when the posterior
is squeezed against the boundary: data centred at 50, but the logistic
transform (with no recentring) puts almost the whole posterior mass
within about 1e-20 of y = 1.}

\tabularnewline

\toprule\noalign{}
method & log\_P\_y \\
\midrule\noalign{}
\endhead
\bottomrule\noalign{}
\endlastfoot
midpoint (n = 100) & -535.484 \\
Monte Carlo (n = 10,000) & -251.749 \\

\end{longtable}

}

The two estimates disagree by many orders of magnitude on the log scale
--- the grid is simply not looking where the posterior mass is. The
general remedy is the same as for the peaked-integrand problem:
\textbf{recentre the transform} on a preliminary location estimate
\(\hat\theta\) (e.g.~the MLE), using
\(y = 1/(1 + e^{-(\theta-\hat\theta)/s})\) for a suitable scale \(s\),
so that the bulk of the posterior maps back to the middle of \((0,1)\)
where the grid is dense.

\subsection{Explore it interactively}\label{explore-it-interactively}

The app below lets you move the true mean of the simulated data and
watch the midpoint-rule estimate diverge from the Monte Carlo baseline
as the posterior is pushed toward the boundary of the transformed
square, and shows the \(g(u,v)\) heat map so you can see the mass
leaving the grid.

\begin{Shaded}
\begin{Highlighting}[]
\NormalTok{\#| label: quadrature{-}failure{-}app}
\NormalTok{\#| standalone: true}
\NormalTok{\#| viewerHeight: 620}

\NormalTok{library(shiny)}

\NormalTok{logit\_inv     \textless{}{-} function(u) log(u / (1 {-} u))}
\NormalTok{log\_der\_logit \textless{}{-} function(u) {-}log(u) {-} log(1 {-} u)}
\NormalTok{log\_lik   \textless{}{-} function(x, mu, w) sum(dnorm(x, mu, exp(w), log = TRUE))}
\NormalTok{log\_prior \textless{}{-} function(mu, w, mu\_0, sigma\_mu, w\_0, sigma\_w) \{}
\NormalTok{    dnorm(mu, mu\_0, sigma\_mu, log = TRUE) + dnorm(w, w\_0, sigma\_w, log = TRUE)}
\NormalTok{\}}
\NormalTok{log\_g \textless{}{-} function(u, v, x, mu\_0, sigma\_mu, w\_0, sigma\_w) \{}
\NormalTok{    mu \textless{}{-} logit\_inv(u); w \textless{}{-} logit\_inv(v)}
\NormalTok{    log\_lik(x, mu, w) + log\_prior(mu, w, mu\_0, sigma\_mu, w\_0, sigma\_w) +}
\NormalTok{        log\_der\_logit(u) + log\_der\_logit(v)}
\NormalTok{\}}
\NormalTok{log\_sum\_exp \textless{}{-} function(lx) \{ m \textless{}{-} max(lx); m + log(sum(exp(lx {-} m))) \}}

\NormalTok{log\_marlik\_mid \textless{}{-} function(x, mu\_0, sigma\_mu, w\_0, sigma\_w, n) \{}
\NormalTok{    h \textless{}{-} 1 / n}
\NormalTok{    grid \textless{}{-} (seq\_len(n) {-} 0.5) * h}
\NormalTok{    G \textless{}{-} outer(grid, grid, Vectorize(function(u, v) log\_g(u, v, x, mu\_0, sigma\_mu, w\_0, sigma\_w)))}
\NormalTok{    list(log\_ml = log\_sum\_exp(G) + 2 * log(h), grid = grid, G = G)}
\NormalTok{\}}
\NormalTok{log\_marlik\_mc \textless{}{-} function(x, mu\_0, sigma\_mu, w\_0, sigma\_w, iters\_mc) \{}
\NormalTok{    mus \textless{}{-} rnorm(iters\_mc, mu\_0, sigma\_mu)}
\NormalTok{    ws  \textless{}{-} rnorm(iters\_mc, w\_0, sigma\_w)}
\NormalTok{    v\_log\_lik \textless{}{-} vapply(seq\_len(iters\_mc), function(i) log\_lik(x, mus[i], ws[i]), numeric(1))}
\NormalTok{    log\_sum\_exp(v\_log\_lik) {-} log(iters\_mc)}
\NormalTok{\}}

\NormalTok{ui \textless{}{-} fluidPage(}
\NormalTok{    titlePanel(NULL),}
\NormalTok{    sidebarLayout(}
\NormalTok{        sidebarPanel(}
\NormalTok{            width = 4,}
\NormalTok{            sliderInput("true\_mean", "true mean of the simulated data",}
\NormalTok{                        min = 0, max = 60, value = 0, step = 2),}
\NormalTok{            sliderInput("n\_grid", "grid points per axis (midpoint rule)",}
\NormalTok{                        min = 10, max = 60, value = 20, step = 5),}
\NormalTok{            helpText("The prior on mu is N(0, 100), i.e. it is not recentred ",}
\NormalTok{                     "on the data. As the true mean moves past about 15{-}20, the ",}
\NormalTok{                     "posterior for mu is squeezed against y = 1 after the ",}
\NormalTok{                     "logistic transform, and the midpoint estimate breaks away ",}
\NormalTok{                     "from the Monte Carlo baseline."),}
\NormalTok{            tableOutput("tab")}
\NormalTok{        ),}
\NormalTok{        mainPanel(}
\NormalTok{            width = 8,}
\NormalTok{            plotOutput("heat", height = "320px"),}
\NormalTok{            plotOutput("compare", height = "260px")}
\NormalTok{        )}
\NormalTok{    )}
\NormalTok{)}

\NormalTok{server \textless{}{-} function(input, output, session) \{}

\NormalTok{    dat \textless{}{-} reactive(\{}
\NormalTok{        set.seed(1)}
\NormalTok{        rnorm(60, input$true\_mean, 2)}
\NormalTok{    \})}

\NormalTok{    fit \textless{}{-} reactive(\{}
\NormalTok{        log\_marlik\_mid(dat(), 0, 100, 0, 100, input$n\_grid)}
\NormalTok{    \})}

\NormalTok{    mc \textless{}{-} reactive(\{}
\NormalTok{        set.seed(2)}
\NormalTok{        log\_marlik\_mc(dat(), 0, 100, 0, 100, 20000)}
\NormalTok{    \})}

\NormalTok{    output$heat \textless{}{-} renderPlot(\{}
\NormalTok{        f \textless{}{-} fit()}
\NormalTok{        par(mar = c(4, 4, 2, 1))}
\NormalTok{        image(f$grid, f$grid, exp(f$G {-} max(f$G)),}
\NormalTok{              col = hcl.colors(30, "YlOrRd", rev = TRUE),}
\NormalTok{              xlab = "u", ylab = "v",}
\NormalTok{              main = "g(u, v) on the unit square: where is the mass?")}
\NormalTok{    \})}

\NormalTok{    output$compare \textless{}{-} renderPlot(\{}
\NormalTok{        par(mar = c(4, 4, 1, 1))}
\NormalTok{        vals \textless{}{-} c(midpoint = fit()$log\_ml, monte\_carlo = mc())}
\NormalTok{        bp \textless{}{-} barplot(vals, col = c("steelblue", "firebrick"),}
\NormalTok{                      ylab = "estimated log P(y)",}
\NormalTok{                      ylim = range(vals) + c({-}1, 1) * 0.15 * diff(range(vals) + 1e{-}6))}
\NormalTok{        text(bp, vals, sprintf("\%.2f", vals), pos = 3)}
\NormalTok{    \})}

\NormalTok{    output$tab \textless{}{-} renderTable(\{}
\NormalTok{        data.frame(method = c("midpoint rule", "Monte Carlo"),}
\NormalTok{                   \textasciigrave{}log P(y)\textasciigrave{} = c(fit()$log\_ml, mc()), check.names = FALSE)}
\NormalTok{    \}, digits = 3, rownames = FALSE)}
\NormalTok{\}}

\NormalTok{shinyApp(ui, server)}
\end{Highlighting}
\end{Shaded}

\begin{tcolorbox}[enhanced jigsaw, arc=.35mm, bottomrule=.15mm, bottomtitle=1mm, breakable, colback=white, colbacktitle=quarto-callout-note-color!10!white, colframe=quarto-callout-note-color-frame, coltitle=black, left=2mm, leftrule=.75mm, opacityback=0, opacitybacktitle=0.6, rightrule=.15mm, title=\textcolor{quarto-callout-note-color}{\faInfo}\hspace{0.5em}{Reproducing this figure}, titlerule=0mm, toprule=.15mm, toptitle=1mm]

The chunk above uses the
\href{https://github.com/quarto-ext/shinylive}{shinylive} extension,
which compiles the app to WebAssembly so it runs in the reader's browser
with no Shiny server. Install it once per project with

\begin{Shaded}
\begin{Highlighting}[]
\ExtensionTok{quarto}\NormalTok{ add quarto{-}ext/shinylive}
\end{Highlighting}
\end{Shaded}

and add \texttt{filters:\ {[}shinylive{]}} to the document or project
YAML. Only packages available in webR may be used inside the app, so the
code above is restricted to base R and \texttt{shiny}.

\end{tcolorbox}

\subsection{The curse of
dimensionality}\label{the-curse-of-dimensionality}

A product grid with \(n\) points per axis needs \(n^d\) evaluations in
\(d\) dimensions:

\begin{longtable}[]{@{}rrr@{}}
\caption{Grid size grows exponentially with
dimension}\label{tbl-curse}\tabularnewline
\toprule\noalign{}
\(d\) & \(n=20\) & \(n=50\) \\
\midrule\noalign{}
\endfirsthead
\toprule\noalign{}
\(d\) & \(n=20\) & \(n=50\) \\
\midrule\noalign{}
\endhead
\bottomrule\noalign{}
\endlastfoot
1 & 20 & 50 \\
2 & 400 & 2,500 \\
5 & \(3.2\times10^6\) & \(3.1\times10^8\) \\
10 & \(1.0\times10^{13}\) & \(9.8\times10^{16}\) \\
\end{longtable}

Meanwhile the posterior occupies only a tiny, often tilted, region
inside that grid --- most evaluations are wasted on points where
\(L(\theta)\pi(\theta)
\approx 0\) --- and the error of a product rule, expressed in terms of
the total number of evaluations \(N=n^d\), degrades to \(O(N^{-2/d})\).

\begin{Shaded}
\begin{Highlighting}[]
\FunctionTok{par}\NormalTok{(}\AttributeTok{mfrow =} \FunctionTok{c}\NormalTok{(}\DecValTok{1}\NormalTok{, }\DecValTok{2}\NormalTok{), }\AttributeTok{mar =} \FunctionTok{c}\NormalTok{(}\DecValTok{4}\NormalTok{, }\DecValTok{4}\NormalTok{, }\DecValTok{2}\NormalTok{, }\DecValTok{1}\NormalTok{))}
\NormalTok{S }\OtherTok{\textless{}{-}} \FunctionTok{matrix}\NormalTok{(}\FunctionTok{c}\NormalTok{(}\DecValTok{1}\NormalTok{, }\FloatTok{0.85}\NormalTok{, }\FloatTok{0.85}\NormalTok{, }\DecValTok{1}\NormalTok{), }\DecValTok{2}\NormalTok{) }\SpecialCharTok{*} \FloatTok{0.15}
\NormalTok{Si }\OtherTok{\textless{}{-}} \FunctionTok{solve}\NormalTok{(S)}
\NormalTok{grd }\OtherTok{\textless{}{-}} \FunctionTok{seq}\NormalTok{(}\SpecialCharTok{{-}}\DecValTok{3}\NormalTok{, }\DecValTok{3}\NormalTok{, }\AttributeTok{length.out =} \DecValTok{150}\NormalTok{)}
\NormalTok{Z }\OtherTok{\textless{}{-}} \FunctionTok{outer}\NormalTok{(grd, grd, }\FunctionTok{Vectorize}\NormalTok{(}\ControlFlowTok{function}\NormalTok{(a, b) }\FunctionTok{exp}\NormalTok{(}\SpecialCharTok{{-}}\FloatTok{0.5} \SpecialCharTok{*} \FunctionTok{c}\NormalTok{(a, b) }\SpecialCharTok{\%*\%}\NormalTok{ Si }\SpecialCharTok{\%*\%} \FunctionTok{c}\NormalTok{(a, b))))}
\FunctionTok{contour}\NormalTok{(grd, grd, Z, }\AttributeTok{nlevels =} \DecValTok{6}\NormalTok{, }\AttributeTok{xlab =} \FunctionTok{expression}\NormalTok{(theta[}\DecValTok{1}\NormalTok{]), }\AttributeTok{ylab =} \FunctionTok{expression}\NormalTok{(theta[}\DecValTok{2}\NormalTok{]),}
        \AttributeTok{main =} \StringTok{"posterior inside a 12 x 12 grid"}\NormalTok{, }\AttributeTok{drawlabels =} \ConstantTok{FALSE}\NormalTok{, }\AttributeTok{col =} \StringTok{"firebrick"}\NormalTok{, }\AttributeTok{lwd =} \DecValTok{2}\NormalTok{)}
\NormalTok{gr }\OtherTok{\textless{}{-}} \FunctionTok{seq}\NormalTok{(}\SpecialCharTok{{-}}\DecValTok{3}\NormalTok{, }\DecValTok{3}\NormalTok{, }\AttributeTok{length.out =} \DecValTok{13}\NormalTok{)}
\FunctionTok{abline}\NormalTok{(}\AttributeTok{v =}\NormalTok{ gr, }\AttributeTok{h =}\NormalTok{ gr, }\AttributeTok{col =} \FunctionTok{adjustcolor}\NormalTok{(}\StringTok{"gray30"}\NormalTok{, }\FloatTok{0.5}\NormalTok{))}
\NormalTok{pts }\OtherTok{\textless{}{-}} \FunctionTok{expand.grid}\NormalTok{(gr, gr)}
\FunctionTok{points}\NormalTok{(pts, }\AttributeTok{pch =} \DecValTok{19}\NormalTok{, }\AttributeTok{cex =} \FloatTok{0.5}\NormalTok{)}
\NormalTok{inside }\OtherTok{\textless{}{-}} \FunctionTok{apply}\NormalTok{(pts, }\DecValTok{1}\NormalTok{, }\ControlFlowTok{function}\NormalTok{(p) }\FunctionTok{exp}\NormalTok{(}\SpecialCharTok{{-}}\FloatTok{0.5} \SpecialCharTok{*}\NormalTok{ p }\SpecialCharTok{\%*\%}\NormalTok{ Si }\SpecialCharTok{\%*\%}\NormalTok{ p)) }\SpecialCharTok{\textgreater{}} \FloatTok{0.01}
\FunctionTok{points}\NormalTok{(pts[inside, ], }\AttributeTok{pch =} \DecValTok{19}\NormalTok{, }\AttributeTok{col =} \StringTok{"firebrick"}\NormalTok{, }\AttributeTok{cex =} \FloatTok{0.9}\NormalTok{)}
\FunctionTok{mtext}\NormalTok{(}\FunctionTok{sprintf}\NormalTok{(}\StringTok{"\%d of \%d grid points carry non{-}negligible mass"}\NormalTok{, }\FunctionTok{sum}\NormalTok{(inside), }\FunctionTok{nrow}\NormalTok{(pts)),}
      \AttributeTok{side =} \DecValTok{1}\NormalTok{, }\AttributeTok{line =} \FloatTok{2.6}\NormalTok{, }\AttributeTok{cex =} \FloatTok{0.85}\NormalTok{)}

\FunctionTok{plot}\NormalTok{(}\DecValTok{1}\SpecialCharTok{:}\DecValTok{20}\NormalTok{, }\FunctionTok{sapply}\NormalTok{(}\DecValTok{1}\SpecialCharTok{:}\DecValTok{20}\NormalTok{, }\ControlFlowTok{function}\NormalTok{(d) }\DecValTok{20}\SpecialCharTok{\^{}}\NormalTok{d), }\AttributeTok{log =} \StringTok{"y"}\NormalTok{, }\AttributeTok{type =} \StringTok{"b"}\NormalTok{, }\AttributeTok{pch =} \DecValTok{19}\NormalTok{,}
     \AttributeTok{xlab =} \StringTok{"dimension d"}\NormalTok{, }\AttributeTok{ylab =} \StringTok{"evaluations with 20 points per axis"}\NormalTok{)}
\FunctionTok{abline}\NormalTok{(}\AttributeTok{h =} \FloatTok{1e12}\NormalTok{, }\AttributeTok{lty =} \DecValTok{2}\NormalTok{)}
\FunctionTok{text}\NormalTok{(}\DecValTok{3}\NormalTok{, }\FloatTok{3e12}\NormalTok{, }\StringTok{"one trillion"}\NormalTok{, }\AttributeTok{adj =} \DecValTok{0}\NormalTok{, }\AttributeTok{cex =} \FloatTok{0.9}\NormalTok{)}
\end{Highlighting}
\end{Shaded}

\begin{figure}[H]

\centering{

\pandocbounded{\includegraphics[keepaspectratio]{chapters/unit08_integral_files/figure-pdf/fig-curse-1.pdf}}

}

\caption{\label{fig-curse}Left: a correlated 2-D posterior inside a
12x12 grid; only the highlighted points carry non-negligible mass.
Right: evaluations needed for a 20-point-per-axis grid as dimension
increases.}

\end{figure}%

\textbf{Conclusion.} Quadrature is a reliable, cheap tool for
\(d \lesssim 3\), and it is what the rest of this book uses whenever a
low-dimensional integral needs checking against another method. Beyond
that, Monte Carlo methods --- whose error is \(O(N^{-1/2})\)
\emph{regardless of dimension} --- take over, at the cost of losing the
very fast convergence rates of Simpson's rule; this trade-off motivates
the sampling methods developed from the next few chapters onward.

\section{Summary}\label{summary-2}

\begin{itemize}
\tightlist
\item
  Quadrature approximates \(\int_a^b f\,dx\) by a weighted sum of \(f\)
  at grid points; infinite ranges are handled by a change of variable
  such as the logistic transform, always carrying along the Jacobian
  \(t'(y)\).
\item
  The midpoint and trapezoidal rules have \(O(h_n^2)\) error; Simpson's
  rule has \(O(h_n^4)\). The midpoint rule's extra virtue --- never
  evaluating the integrand at a panel edge --- matters whenever the
  transformed integrand is singular at the boundary, as it is here.
\item
  The marginal likelihood \(P(y) = \int L(\theta)\pi(\theta)\,d\theta\)
  is the key statistical use case in this book: always compute it on the
  log scale with log-sum-exp, and validate a quadrature answer against
  an independent Monte Carlo estimate.
\item
  Quadrature fails when the integrand is peaked relative to the grid, or
  when a transform squeezes the mass of interest into a region the grid
  cannot resolve; both problems are cured by locating the mass first
  (optimization) and centring the grid or transform on it --- the idea
  taken up systematically by the Laplace approximation in the next
  chapter.
\item
  Product grids cost \(n^d\) evaluations, so quadrature is practical
  only for small \(d\); higher-dimensional problems require Monte Carlo
  methods.
\end{itemize}

\bookmarksetup{startatroot}

\chapter{Laplace Approximation}\label{laplace-approximation}

The previous chapter approximated an integral by a weighted sum of
function values on a grid. A grid gets expensive fast: an
\(n\)-point-per-axis product grid costs \(n^d\) evaluations, so
quadrature is only practical for \(d \lesssim
3\). The \textbf{Laplace approximation} takes a completely different
route: instead of evaluating the integrand everywhere, it looks only at
the single most important point --- the mode --- and replaces the whole
integrand there by the Gaussian that matches its height and curvature.
No grid is needed at any dimension, at the price of assuming the
integrand is roughly bell-shaped.

\section{The Goal and the Idea}\label{the-goal-and-the-idea}

Many Bayesian calculations are posterior expectations \[
E_f\big[a(\theta)\big] = \frac{\int a(\theta)\,f(\theta)\,d\theta}{\int f(\theta)\,d\theta},
\] where \(\theta \in \mathbb{R}^p\), \(a(\theta)\) is some quantity of
interest (a parameter, a prediction, an indicator), and \(f(\theta)\) is
an \textbf{unnormalized} posterior, typically
\(f(\theta) = L(\theta)\pi(\theta)\). Both numerator and denominator are
\(p\)-dimensional integrals.

Write \(h(\theta) = -\log f(\theta)\), so
\(f(\theta) = e^{-h(\theta)}\). For independent data, \(h\) is (up to an
additive constant) minus the log-likelihood plus minus the log-prior ---
a sum of \(n\) terms --- so as \(n\) grows it becomes large and sharply
curved at its minimum. This is the same large-sample regularity that
makes the MLE asymptotically normal (seen in the maximum-likelihood
chapters); here it makes the \emph{posterior} asymptotically normal too.
The Laplace approximation makes this precise by expanding \(h\) to
second order at its minimum and integrating the resulting Gaussian
exactly --- no grid, no sampling, just one optimization and one Hessian.

\section{Derivation}\label{derivation-1}

Let \(\theta_0 = \arg\min_\theta h(\theta)\), the posterior mode, found
numerically (\texttt{optim(...,\ hessian\ =\ TRUE)} or
\texttt{nlm(...,\ hessian\ =\ TRUE)} in R --- exactly the optimizers
used to find the MLE in earlier chapters). Expand \(h\) around
\(\theta_0\): \[
h(\theta) = h(\theta_0) + \underbrace{\nabla h(\theta_0)^\top}_{=\,0\text{ at the minimum}}(\theta - \theta_0)
 + \tfrac12(\theta-\theta_0)^\top H (\theta - \theta_0) + O\big(\|\theta-\theta_0\|^3\big),
\qquad H = \nabla^2 h(\theta_0).
\] \(H\), the Hessian of \(-\log f\) at the mode, is positive definite
at a genuine minimum; for \(f = L\pi\) it is the observed information
plus the curvature of the log-prior. Dropping the cubic remainder and
exponentiating, \[
f(\theta) \approx \tilde f(\theta) = f(\theta_0)\exp\Big\{-\tfrac12(\theta-\theta_0)^\top H(\theta-\theta_0)\Big\}
 = \underbrace{f(\theta_0)(2\pi)^{p/2}|H|^{-1/2}}_{\text{constant}}
   \cdot \underbrace{(2\pi)^{-p/2}|H|^{1/2}e^{-\frac12(\theta-\theta_0)^\top H(\theta-\theta_0)}}_{N_p(\theta_0,\,H^{-1})\text{ density}}.
\] Two results follow immediately, since the bracketed factor is a
properly normalized \(N_p(\theta_0, H^{-1})\) density:

\[
\boxed{\ \theta \mid \text{data} \ \approx\ N_p\big(\theta_0,\ H^{-1}\big) \ }
\qquad\qquad
\boxed{\ \int f(\theta)\,d\theta \ \approx\ f(\theta_0)\,(2\pi)^{p/2}\,|H|^{-1/2} \ }
\]

Only the mode and the Hessian of \(-\log f\) at the mode are needed ---
both come for free from any optimizer that returns a Hessian. The
relative error of the second result is \(O(1/n)\) for regular models;
the error of the first is governed by how skewed the true posterior
actually is.

\section{Worked Example: A Poisson Rate with a Gamma
Prior}\label{worked-example-a-poisson-rate-with-a-gamma-prior}

To see the approximation against ground truth, use a model with a known
exact answer. For
\(y_1,\ldots,y_n \mid \lambda \sim \text{Poisson}(\lambda)\) and
\(\lambda \sim \text{Gamma}(\alpha,\beta)\), the posterior is exactly
\(\text{Gamma}(\alpha + \sum y_i,\ \beta+n)\), so both the exact log
marginal likelihood and the exact posterior are available in closed form
and the Laplace error can be measured directly.

\begin{Shaded}
\begin{Highlighting}[]
\NormalTok{alpha }\OtherTok{\textless{}{-}} \DecValTok{2}\NormalTok{; beta }\OtherTok{\textless{}{-}} \DecValTok{1}

\DocumentationTok{\#\#\# Laplace approximation: mode, Hessian, and the two boxed results above}
\NormalTok{laplace\_poisson }\OtherTok{\textless{}{-}} \ControlFlowTok{function}\NormalTok{(y) \{}
\NormalTok{  h  }\OtherTok{\textless{}{-}} \ControlFlowTok{function}\NormalTok{(l) }\SpecialCharTok{{-}}\NormalTok{(}\FunctionTok{sum}\NormalTok{(}\FunctionTok{dpois}\NormalTok{(y, l, }\AttributeTok{log =} \ConstantTok{TRUE}\NormalTok{)) }\SpecialCharTok{+} \FunctionTok{dgamma}\NormalTok{(l, alpha, beta, }\AttributeTok{log =} \ConstantTok{TRUE}\NormalTok{))}
\NormalTok{  op }\OtherTok{\textless{}{-}} \FunctionTok{optim}\NormalTok{(}\FunctionTok{mean}\NormalTok{(y) }\SpecialCharTok{+} \FloatTok{0.5}\NormalTok{, h, }\AttributeTok{method =} \StringTok{"L{-}BFGS{-}B"}\NormalTok{, }\AttributeTok{lower =} \FloatTok{1e{-}6}\NormalTok{, }\AttributeTok{hessian =} \ConstantTok{TRUE}\NormalTok{)}
  \FunctionTok{c}\NormalTok{(}\AttributeTok{mode =}\NormalTok{ op}\SpecialCharTok{$}\NormalTok{par, }\AttributeTok{se =} \FunctionTok{sqrt}\NormalTok{(}\DecValTok{1} \SpecialCharTok{/}\NormalTok{ op}\SpecialCharTok{$}\NormalTok{hessian[}\DecValTok{1}\NormalTok{, }\DecValTok{1}\NormalTok{]),}
    \AttributeTok{log\_ml =} \SpecialCharTok{{-}}\NormalTok{op}\SpecialCharTok{$}\NormalTok{value }\SpecialCharTok{+} \FloatTok{0.5} \SpecialCharTok{*} \FunctionTok{log}\NormalTok{(}\DecValTok{2} \SpecialCharTok{*}\NormalTok{ pi) }\SpecialCharTok{{-}} \FloatTok{0.5} \SpecialCharTok{*} \FunctionTok{log}\NormalTok{(op}\SpecialCharTok{$}\NormalTok{hessian[}\DecValTok{1}\NormalTok{, }\DecValTok{1}\NormalTok{]))}
\NormalTok{\}}

\DocumentationTok{\#\#\# exact log marginal likelihood, from the Gamma{-}Poisson conjugate identity}
\NormalTok{exact\_log\_ml }\OtherTok{\textless{}{-}} \ControlFlowTok{function}\NormalTok{(y) \{}
\NormalTok{  n }\OtherTok{\textless{}{-}} \FunctionTok{length}\NormalTok{(y); s }\OtherTok{\textless{}{-}} \FunctionTok{sum}\NormalTok{(y)}
\NormalTok{  alpha }\SpecialCharTok{*} \FunctionTok{log}\NormalTok{(beta) }\SpecialCharTok{{-}} \FunctionTok{lgamma}\NormalTok{(alpha) }\SpecialCharTok{+} \FunctionTok{lgamma}\NormalTok{(alpha }\SpecialCharTok{+}\NormalTok{ s) }\SpecialCharTok{{-}}
\NormalTok{    (alpha }\SpecialCharTok{+}\NormalTok{ s) }\SpecialCharTok{*} \FunctionTok{log}\NormalTok{(beta }\SpecialCharTok{+}\NormalTok{ n) }\SpecialCharTok{{-}} \FunctionTok{sum}\NormalTok{(}\FunctionTok{lfactorial}\NormalTok{(y))}
\NormalTok{\}}
\end{Highlighting}
\end{Shaded}

\begin{Shaded}
\begin{Highlighting}[]
\FunctionTok{set.seed}\NormalTok{(}\DecValTok{42}\NormalTok{)}
\NormalTok{res }\OtherTok{\textless{}{-}} \FunctionTok{t}\NormalTok{(}\FunctionTok{sapply}\NormalTok{(}\FunctionTok{c}\NormalTok{(}\DecValTok{5}\NormalTok{, }\DecValTok{20}\NormalTok{, }\DecValTok{100}\NormalTok{, }\DecValTok{500}\NormalTok{), }\ControlFlowTok{function}\NormalTok{(n) \{}
\NormalTok{  y }\OtherTok{\textless{}{-}} \FunctionTok{rpois}\NormalTok{(n, }\DecValTok{3}\NormalTok{)}
\NormalTok{  lap }\OtherTok{\textless{}{-}} \FunctionTok{laplace\_poisson}\NormalTok{(y)}
  \FunctionTok{c}\NormalTok{(}\AttributeTok{n =}\NormalTok{ n, }\AttributeTok{laplace\_log\_ml =}\NormalTok{ lap[}\StringTok{"log\_ml"}\NormalTok{], }\AttributeTok{exact\_log\_ml =} \FunctionTok{exact\_log\_ml}\NormalTok{(y),}
    \AttributeTok{error =}\NormalTok{ lap[}\StringTok{"log\_ml"}\NormalTok{] }\SpecialCharTok{{-}} \FunctionTok{exact\_log\_ml}\NormalTok{(y),}
    \AttributeTok{laplace\_mode =}\NormalTok{ lap[}\StringTok{"mode"}\NormalTok{], }\AttributeTok{posterior\_mean =}\NormalTok{ (alpha }\SpecialCharTok{+} \FunctionTok{sum}\NormalTok{(y)) }\SpecialCharTok{/}\NormalTok{ (beta }\SpecialCharTok{+}\NormalTok{ n))}
\NormalTok{\}))}
\NormalTok{knitr}\SpecialCharTok{::}\FunctionTok{kable}\NormalTok{(res, }\AttributeTok{digits =} \DecValTok{4}\NormalTok{, }\AttributeTok{row.names =} \ConstantTok{FALSE}\NormalTok{,}
             \AttributeTok{col.names =} \FunctionTok{c}\NormalTok{(}\StringTok{"n"}\NormalTok{, }\StringTok{"Laplace log P(y)"}\NormalTok{, }\StringTok{"exact log P(y)"}\NormalTok{, }\StringTok{"error"}\NormalTok{,}
                           \StringTok{"Laplace mode"}\NormalTok{, }\StringTok{"posterior mean"}\NormalTok{))}
\end{Highlighting}
\end{Shaded}

{

\begin{longtable}[]{@{}
  >{\raggedleft\arraybackslash}p{(\linewidth - 10\tabcolsep) * \real{0.0556}}
  >{\raggedleft\arraybackslash}p{(\linewidth - 10\tabcolsep) * \real{0.2361}}
  >{\raggedleft\arraybackslash}p{(\linewidth - 10\tabcolsep) * \real{0.2083}}
  >{\raggedleft\arraybackslash}p{(\linewidth - 10\tabcolsep) * \real{0.1111}}
  >{\raggedleft\arraybackslash}p{(\linewidth - 10\tabcolsep) * \real{0.1806}}
  >{\raggedleft\arraybackslash}p{(\linewidth - 10\tabcolsep) * \real{0.2083}}@{}}

\caption{\label{tbl-poisson-gamma-error}Laplace log marginal likelihood
versus the exact value, and the posterior mean versus the Laplace mode,
as n grows. The error shrinks roughly like 1/n: it falls by about a
factor of 4-5 each time n increases by that factor.}

\tabularnewline

\toprule\noalign{}
\begin{minipage}[b]{\linewidth}\raggedleft
n
\end{minipage} & \begin{minipage}[b]{\linewidth}\raggedleft
Laplace log P(y)
\end{minipage} & \begin{minipage}[b]{\linewidth}\raggedleft
exact log P(y)
\end{minipage} & \begin{minipage}[b]{\linewidth}\raggedleft
error
\end{minipage} & \begin{minipage}[b]{\linewidth}\raggedleft
Laplace mode
\end{minipage} & \begin{minipage}[b]{\linewidth}\raggedleft
posterior mean
\end{minipage} \\
\midrule\noalign{}
\endhead
\bottomrule\noalign{}
\endlastfoot
5 & -11.3822 & -11.3784 & -0.0038 & 3.6667 & 3.8333 \\
20 & -44.2274 & -44.2263 & -0.0011 & 3.6190 & 3.6667 \\
100 & -197.0022 & -197.0019 & -0.0003 & 2.9406 & 2.9505 \\
500 & -958.5983 & -958.5983 & -0.0001 & 2.8842 & 2.8862 \\

\end{longtable}

}

The error is a few thousandths at \(n=5\) and shrinks to a few
ten-thousandths by \(n=500\); the Laplace mode also converges to the
posterior mean as the Gamma posterior becomes more symmetric. Plotting
the exact posterior against its Laplace approximation makes this
concrete:

\begin{Shaded}
\begin{Highlighting}[]
\FunctionTok{par}\NormalTok{(}\AttributeTok{mfrow =} \FunctionTok{c}\NormalTok{(}\DecValTok{1}\NormalTok{, }\DecValTok{2}\NormalTok{), }\AttributeTok{mar =} \FunctionTok{c}\NormalTok{(}\DecValTok{4}\NormalTok{, }\DecValTok{4}\NormalTok{, }\DecValTok{2}\NormalTok{, }\DecValTok{1}\NormalTok{))}
\FunctionTok{set.seed}\NormalTok{(}\DecValTok{42}\NormalTok{)}
\ControlFlowTok{for}\NormalTok{ (n }\ControlFlowTok{in} \FunctionTok{c}\NormalTok{(}\DecValTok{5}\NormalTok{, }\DecValTok{100}\NormalTok{)) \{}
\NormalTok{  y }\OtherTok{\textless{}{-}} \FunctionTok{rpois}\NormalTok{(n, }\DecValTok{3}\NormalTok{); a1 }\OtherTok{\textless{}{-}}\NormalTok{ alpha }\SpecialCharTok{+} \FunctionTok{sum}\NormalTok{(y); b1 }\OtherTok{\textless{}{-}}\NormalTok{ beta }\SpecialCharTok{+}\NormalTok{ n}
\NormalTok{  lap }\OtherTok{\textless{}{-}} \FunctionTok{laplace\_poisson}\NormalTok{(y)}
  \FunctionTok{curve}\NormalTok{(}\FunctionTok{dgamma}\NormalTok{(x, a1, b1), }\FunctionTok{max}\NormalTok{(}\DecValTok{0}\NormalTok{, lap[}\StringTok{"mode"}\NormalTok{] }\SpecialCharTok{{-}} \DecValTok{4} \SpecialCharTok{*}\NormalTok{ lap[}\StringTok{"se"}\NormalTok{]), lap[}\StringTok{"mode"}\NormalTok{] }\SpecialCharTok{+} \DecValTok{4} \SpecialCharTok{*}\NormalTok{ lap[}\StringTok{"se"}\NormalTok{],}
        \AttributeTok{lwd =} \DecValTok{3}\NormalTok{, }\AttributeTok{col =} \StringTok{"firebrick"}\NormalTok{, }\AttributeTok{xlab =} \FunctionTok{expression}\NormalTok{(lambda), }\AttributeTok{ylab =} \StringTok{"density"}\NormalTok{,}
        \AttributeTok{main =} \FunctionTok{sprintf}\NormalTok{(}\StringTok{"n = \%d"}\NormalTok{, n))}
  \FunctionTok{curve}\NormalTok{(}\FunctionTok{dnorm}\NormalTok{(x, lap[}\StringTok{"mode"}\NormalTok{], lap[}\StringTok{"se"}\NormalTok{]), }\AttributeTok{add =} \ConstantTok{TRUE}\NormalTok{, }\AttributeTok{lwd =} \DecValTok{3}\NormalTok{, }\AttributeTok{col =} \StringTok{"forestgreen"}\NormalTok{)}
  \ControlFlowTok{if}\NormalTok{ (n }\SpecialCharTok{==} \DecValTok{5}\NormalTok{) }\FunctionTok{legend}\NormalTok{(}\StringTok{"topright"}\NormalTok{, }\AttributeTok{bty =} \StringTok{"n"}\NormalTok{, }\AttributeTok{lwd =} \DecValTok{3}\NormalTok{, }\AttributeTok{col =} \FunctionTok{c}\NormalTok{(}\StringTok{"firebrick"}\NormalTok{, }\StringTok{"forestgreen"}\NormalTok{),}
                      \AttributeTok{legend =} \FunctionTok{c}\NormalTok{(}\StringTok{"exact posterior"}\NormalTok{, }\StringTok{"Laplace approximation"}\NormalTok{))}
\NormalTok{\}}
\end{Highlighting}
\end{Shaded}

\begin{figure}[H]

\centering{

\pandocbounded{\includegraphics[keepaspectratio]{chapters/unit09_laplace_files/figure-pdf/fig-poisson-gamma-density-1.pdf}}

}

\caption{\label{fig-poisson-gamma-density}Exact Gamma posterior (red)
and the Laplace normal approximation (green) for n=5 and n=100. At n=5
the exact posterior is visibly right-skewed and the Gaussian is a rough
match; by n=100 the two are nearly indistinguishable, illustrating the
O(1/n) convergence in the table above.}

\end{figure}%

\subsection{Linear functionals: posterior mean and credible
intervals}\label{linear-functionals-posterior-mean-and-credible-intervals}

If \(a(\theta) = A^\top\theta + a_0\) is linear, replacing \(f\) by
\(\tilde f\) gives \(E_f[a(\theta)] \approx A^\top\theta_0 + a_0\): the
posterior mean is approximated by the posterior mode, and
\(\text{Var}(A^\top\theta) \approx
A^\top H^{-1}A\) gives Wald-style credible intervals \(\theta_{0j} \pm
z_{\alpha/2}\sqrt{(H^{-1})_{jj}}\) --- exactly the Bayesian counterpart
of an MLE plus its Wald standard error from the maximum-likelihood
chapters.

\begin{Shaded}
\begin{Highlighting}[]
\FunctionTok{set.seed}\NormalTok{(}\DecValTok{42}\NormalTok{)}
\NormalTok{n }\OtherTok{\textless{}{-}} \DecValTok{20}\NormalTok{; y }\OtherTok{\textless{}{-}} \FunctionTok{rpois}\NormalTok{(n, }\DecValTok{3}\NormalTok{); lap }\OtherTok{\textless{}{-}} \FunctionTok{laplace\_poisson}\NormalTok{(y)}
\NormalTok{a1 }\OtherTok{\textless{}{-}}\NormalTok{ alpha }\SpecialCharTok{+} \FunctionTok{sum}\NormalTok{(y); b1 }\OtherTok{\textless{}{-}}\NormalTok{ beta }\SpecialCharTok{+}\NormalTok{ n}
\NormalTok{comp }\OtherTok{\textless{}{-}} \FunctionTok{data.frame}\NormalTok{(}
  \AttributeTok{quantity =} \FunctionTok{c}\NormalTok{(}\StringTok{"Laplace (mode / se)"}\NormalTok{, }\StringTok{"exact posterior (mean / sd)"}\NormalTok{),}
  \AttributeTok{center =} \FunctionTok{c}\NormalTok{(lap[}\StringTok{"mode"}\NormalTok{], a1 }\SpecialCharTok{/}\NormalTok{ b1),}
  \AttributeTok{spread =} \FunctionTok{c}\NormalTok{(lap[}\StringTok{"se"}\NormalTok{], }\FunctionTok{sqrt}\NormalTok{(a1) }\SpecialCharTok{/}\NormalTok{ b1),}
  \AttributeTok{lower95 =} \FunctionTok{c}\NormalTok{(lap[}\StringTok{"mode"}\NormalTok{] }\SpecialCharTok{{-}} \FloatTok{1.96} \SpecialCharTok{*}\NormalTok{ lap[}\StringTok{"se"}\NormalTok{], }\FunctionTok{qgamma}\NormalTok{(}\FloatTok{0.025}\NormalTok{, a1, b1)),}
  \AttributeTok{upper95 =} \FunctionTok{c}\NormalTok{(lap[}\StringTok{"mode"}\NormalTok{] }\SpecialCharTok{+} \FloatTok{1.96} \SpecialCharTok{*}\NormalTok{ lap[}\StringTok{"se"}\NormalTok{], }\FunctionTok{qgamma}\NormalTok{(}\FloatTok{0.975}\NormalTok{, a1, b1))}
\NormalTok{)}
\NormalTok{knitr}\SpecialCharTok{::}\FunctionTok{kable}\NormalTok{(comp, }\AttributeTok{digits =} \DecValTok{4}\NormalTok{)}
\end{Highlighting}
\end{Shaded}

{

\begin{longtable}[]{@{}llrrrr@{}}

\caption{\label{tbl-poisson-gamma-ci}Laplace mode/SE-based 95\% interval
versus the exact Gamma posterior mean/sd-based interval, n = 20.}

\tabularnewline

\toprule\noalign{}
& quantity & center & spread & lower95 & upper95 \\
\midrule\noalign{}
\endhead
\bottomrule\noalign{}
\endlastfoot
mode & Laplace (mode / se) & 3.6190 & 0.4151 & 2.8054 & 4.4327 \\
& exact posterior (mean / sd) & 3.6667 & 0.4179 & 2.8937 & 4.5298 \\

\end{longtable}

}

The two intervals agree closely even at \(n=20\): the mild right-skew of
the Gamma posterior visible in Figure~\ref{fig-poisson-gamma-density}
has not yet noticeably distorted a Wald-style interval built from a
symmetric Gaussian.

\subsection{\texorpdfstring{BIC as a large-\(n\)
simplification}{BIC as a large-n simplification}}\label{bic-as-a-large-n-simplification}

With \(f(\theta) = L(\theta;D)\pi(\theta)\), the Laplace approximation
to the log marginal likelihood is \[
\log P(D) \approx \log\big[L(\theta_0;D)\,\pi(\theta_0)\big] + \frac p2\log(2\pi) - \frac12\log|H|.
\] For large \(n\) the Hessian is dominated by the likelihood,
\(H \approx
n\,I(\theta_0)\) with \(I\) the per-observation Fisher information, so
\(|H|
\approx n^p|I(\theta_0)|\) and \[
\log P(D) \approx \log L(\theta_0;D) - \frac p2\log n
 + \underbrace{\log\pi(\theta_0) + \frac p2\log(2\pi) - \frac12\log|I(\theta_0)|}_{O(1),\text{ dropped}}.
\] Multiplying by \(-2\) and replacing the posterior mode by the MLE
(the two coincide to \(O(1/n)\)) gives the \textbf{Bayesian information
criterion}, \(\text{BIC} = -2\log L(\hat\theta_{\text{MLE}}) + p\log n\)
--- a model-comparison score that needs no prior at all, because the
dropped \(O(1)\) terms do not grow with \(n\) and so wash out of
\emph{differences} between models fit to the same data. For the \(n=20\)
Poisson example from the credible-interval check above,
\(-2\times(\text{Laplace log }P(y))\) and the MLE-based BIC differ by
about 1.68 --- a constant offset consistent with the dropped \(O(1)\)
terms, not something that grows with \(n\).

\section{When Laplace Works Well, and When It Does
Not}\label{when-laplace-works-well-and-when-it-does-not}

\textbf{Works well} when \(n\) is large relative to \(p\) and the
posterior is unimodal and roughly symmetric, as in the \(n=100\) panel
above. Parameters should live on an unbounded scale --- transform first
(\(\log\sigma\), \(\text{logit}\,p\)) so the matching Gaussian is not
truncated at a boundary, remembering to include the Jacobian of the
transform in \(f\) (exactly the \(w=\log\sigma^2\) reparametrization
used throughout this book since the multivariate optimization chapter).

\textbf{Struggles} when the posterior is skewed (small \(n\), a
parameter near a boundary) --- the mode is then a poor stand-in for the
mean and the Gaussian tails are wrong, as the \(n=5\) panel above
already shows mildly; when the posterior is multimodal, since the
approximation only sees the one mode an optimizer happened to find; and
in hierarchical models where \(p\) grows with \(n\), so the \(O(1/n)\)
error does not vanish (nested Laplace approximation, INLA, handles this
by integrating out the high-dimensional part analytically).

\section{Comparing Laplace, Quadrature, and Monte Carlo on the Same
Problem}\label{comparing-laplace-quadrature-and-monte-carlo-on-the-same-problem}

The previous chapter (Numerical Quadrature) computed the marginal
likelihood of a normal sample with unknown mean \(\mu\) and log-variance
\(w=\log\sigma^2\), \[
P(y) = \int_{-\infty}^{\infty}\!\int_{-\infty}^{\infty} L(\mu,w)\,\pi(\mu,w)\,d\mu\,dw,
\qquad \mu\sim N(\mu_0,\sigma_\mu^2),\ w\sim N(w_0,\sigma_w^2),
\] by a product midpoint rule on a logistic-transformed \((0,1)^2\)
grid, checked against naive Monte Carlo. The Laplace approximation
attacks the same integral by optimization instead of a grid: find the
joint mode of \((\mu,w)\) and its \(2\times2\) Hessian, then apply the
boxed formula with \(p=2\).

\begin{Shaded}
\begin{Highlighting}[]
\DocumentationTok{\#\#\# log likelihood and log prior for the normal/log{-}variance model of Unit08}
\NormalTok{log\_lik\_gauss  }\OtherTok{\textless{}{-}} \ControlFlowTok{function}\NormalTok{(x, mu, w) }\FunctionTok{sum}\NormalTok{(}\FunctionTok{dnorm}\NormalTok{(x, mu, }\FunctionTok{exp}\NormalTok{(w), }\AttributeTok{log =} \ConstantTok{TRUE}\NormalTok{))}
\NormalTok{log\_prior\_gauss }\OtherTok{\textless{}{-}} \ControlFlowTok{function}\NormalTok{(mu, w, mu0, sigma\_mu, w0, sigma\_w) \{}
  \FunctionTok{dnorm}\NormalTok{(mu, mu0, sigma\_mu, }\AttributeTok{log =} \ConstantTok{TRUE}\NormalTok{) }\SpecialCharTok{+} \FunctionTok{dnorm}\NormalTok{(w, w0, sigma\_w, }\AttributeTok{log =} \ConstantTok{TRUE}\NormalTok{)}
\NormalTok{\}}
\NormalTok{neg\_log\_post }\OtherTok{\textless{}{-}} \ControlFlowTok{function}\NormalTok{(theta, x, mu0, sigma\_mu, w0, sigma\_w) \{}
  \SpecialCharTok{{-}}\FunctionTok{log\_lik\_gauss}\NormalTok{(x, theta[}\DecValTok{1}\NormalTok{], theta[}\DecValTok{2}\NormalTok{]) }\SpecialCharTok{{-}} \FunctionTok{log\_prior\_gauss}\NormalTok{(theta[}\DecValTok{1}\NormalTok{], theta[}\DecValTok{2}\NormalTok{], mu0, sigma\_mu, w0, sigma\_w)}
\NormalTok{\}}

\DocumentationTok{\#\#\# Laplace approximation for this two{-}parameter posterior}
\NormalTok{laplace\_gaussian }\OtherTok{\textless{}{-}} \ControlFlowTok{function}\NormalTok{(x, mu0, sigma\_mu, w0, sigma\_w) \{}
\NormalTok{  op }\OtherTok{\textless{}{-}} \FunctionTok{nlm}\NormalTok{(neg\_log\_post, }\AttributeTok{p =} \FunctionTok{c}\NormalTok{(}\FunctionTok{mean}\NormalTok{(x), }\FunctionTok{log}\NormalTok{(}\FunctionTok{sd}\NormalTok{(x))), }\AttributeTok{hessian =} \ConstantTok{TRUE}\NormalTok{,}
            \AttributeTok{x =}\NormalTok{ x, }\AttributeTok{mu0 =}\NormalTok{ mu0, }\AttributeTok{sigma\_mu =}\NormalTok{ sigma\_mu, }\AttributeTok{w0 =}\NormalTok{ w0, }\AttributeTok{sigma\_w =}\NormalTok{ sigma\_w)}
\NormalTok{  H }\OtherTok{\textless{}{-}}\NormalTok{ op}\SpecialCharTok{$}\NormalTok{hessian; p }\OtherTok{\textless{}{-}} \FunctionTok{length}\NormalTok{(op}\SpecialCharTok{$}\NormalTok{estimate)}
  \FunctionTok{list}\NormalTok{(}\AttributeTok{mode =}\NormalTok{ op}\SpecialCharTok{$}\NormalTok{estimate, }\AttributeTok{se =} \FunctionTok{sqrt}\NormalTok{(}\FunctionTok{diag}\NormalTok{(}\FunctionTok{solve}\NormalTok{(H))),}
       \AttributeTok{log\_ml =} \SpecialCharTok{{-}}\NormalTok{op}\SpecialCharTok{$}\NormalTok{minimum }\SpecialCharTok{+}\NormalTok{ p }\SpecialCharTok{/} \DecValTok{2} \SpecialCharTok{*} \FunctionTok{log}\NormalTok{(}\DecValTok{2} \SpecialCharTok{*}\NormalTok{ pi) }\SpecialCharTok{{-}} \FloatTok{0.5} \SpecialCharTok{*} \FunctionTok{log}\NormalTok{(}\FunctionTok{det}\NormalTok{(H)))}
\NormalTok{\}}

\DocumentationTok{\#\#\# naive Monte Carlo (samples from the prior; see Unit08 for the log{-}sum{-}exp helper)}
\NormalTok{log\_sum\_exp }\OtherTok{\textless{}{-}} \ControlFlowTok{function}\NormalTok{(lx) \{ m }\OtherTok{\textless{}{-}} \FunctionTok{max}\NormalTok{(lx); m }\SpecialCharTok{+} \FunctionTok{log}\NormalTok{(}\FunctionTok{sum}\NormalTok{(}\FunctionTok{exp}\NormalTok{(lx }\SpecialCharTok{{-}}\NormalTok{ m))) \}}
\NormalTok{log\_mar\_mc }\OtherTok{\textless{}{-}} \ControlFlowTok{function}\NormalTok{(x, mu0, sigma\_mu, w0, sigma\_w, iters) \{}
\NormalTok{  mus }\OtherTok{\textless{}{-}} \FunctionTok{rnorm}\NormalTok{(iters, mu0, sigma\_mu); ws }\OtherTok{\textless{}{-}} \FunctionTok{rnorm}\NormalTok{(iters, w0, sigma\_w)}
\NormalTok{  ll }\OtherTok{\textless{}{-}} \FunctionTok{vapply}\NormalTok{(}\FunctionTok{seq\_len}\NormalTok{(iters), }\ControlFlowTok{function}\NormalTok{(i) }\FunctionTok{log\_lik\_gauss}\NormalTok{(x, mus[i], ws[i]), }\FunctionTok{numeric}\NormalTok{(}\DecValTok{1}\NormalTok{))}
  \FunctionTok{log\_sum\_exp}\NormalTok{(ll) }\SpecialCharTok{{-}} \FunctionTok{log}\NormalTok{(iters)}
\NormalTok{\}}

\DocumentationTok{\#\#\# product midpoint quadrature on the logistic{-}transformed square (Unit08\textquotesingle{}s method)}
\NormalTok{logit\_inv     }\OtherTok{\textless{}{-}} \ControlFlowTok{function}\NormalTok{(u) }\FunctionTok{log}\NormalTok{(u }\SpecialCharTok{/}\NormalTok{ (}\DecValTok{1} \SpecialCharTok{{-}}\NormalTok{ u))}
\NormalTok{log\_der\_logit }\OtherTok{\textless{}{-}} \ControlFlowTok{function}\NormalTok{(u) }\SpecialCharTok{{-}}\FunctionTok{log}\NormalTok{(u) }\SpecialCharTok{{-}} \FunctionTok{log}\NormalTok{(}\DecValTok{1} \SpecialCharTok{{-}}\NormalTok{ u)}
\NormalTok{log\_g }\OtherTok{\textless{}{-}} \ControlFlowTok{function}\NormalTok{(u, v, x, mu0, sigma\_mu, w0, sigma\_w) \{}
\NormalTok{  mu }\OtherTok{\textless{}{-}} \FunctionTok{logit\_inv}\NormalTok{(u); w }\OtherTok{\textless{}{-}} \FunctionTok{logit\_inv}\NormalTok{(v)}
  \FunctionTok{log\_lik\_gauss}\NormalTok{(x, mu, w) }\SpecialCharTok{+} \FunctionTok{log\_prior\_gauss}\NormalTok{(mu, w, mu0, sigma\_mu, w0, sigma\_w) }\SpecialCharTok{+}
    \FunctionTok{log\_der\_logit}\NormalTok{(u) }\SpecialCharTok{+} \FunctionTok{log\_der\_logit}\NormalTok{(v)}
\NormalTok{\}}
\NormalTok{log\_mar\_mid }\OtherTok{\textless{}{-}} \ControlFlowTok{function}\NormalTok{(x, mu0, sigma\_mu, w0, sigma\_w, n) \{}
\NormalTok{  h }\OtherTok{\textless{}{-}} \DecValTok{1} \SpecialCharTok{/}\NormalTok{ n; grid }\OtherTok{\textless{}{-}}\NormalTok{ (}\FunctionTok{seq\_len}\NormalTok{(n) }\SpecialCharTok{{-}} \FloatTok{0.5}\NormalTok{) }\SpecialCharTok{*}\NormalTok{ h}
\NormalTok{  G }\OtherTok{\textless{}{-}} \FunctionTok{outer}\NormalTok{(grid, grid, }\FunctionTok{Vectorize}\NormalTok{(}\ControlFlowTok{function}\NormalTok{(u, v) }\FunctionTok{log\_g}\NormalTok{(u, v, x, mu0, sigma\_mu, w0, sigma\_w)))}
  \FunctionTok{log\_sum\_exp}\NormalTok{(G) }\SpecialCharTok{+} \DecValTok{2} \SpecialCharTok{*} \FunctionTok{log}\NormalTok{(h)}
\NormalTok{\}}
\end{Highlighting}
\end{Shaded}

\begin{Shaded}
\begin{Highlighting}[]
\FunctionTok{set.seed}\NormalTok{(}\DecValTok{7}\NormalTok{)}
\NormalTok{scenarios }\OtherTok{\textless{}{-}} \FunctionTok{list}\NormalTok{(}
  \FunctionTok{c}\NormalTok{(}\AttributeTok{desc =} \StringTok{"well{-}centred: mean 5"}\NormalTok{,    }\AttributeTok{mean =} \DecValTok{5}\NormalTok{,   }\AttributeTok{sd =} \DecValTok{1}\NormalTok{),}
  \FunctionTok{c}\NormalTok{(}\AttributeTok{desc =} \StringTok{"well{-}centred: mean {-}5"}\NormalTok{,   }\AttributeTok{mean =} \SpecialCharTok{{-}}\DecValTok{5}\NormalTok{,  }\AttributeTok{sd =} \DecValTok{1}\NormalTok{),}
  \FunctionTok{c}\NormalTok{(}\AttributeTok{desc =} \StringTok{"far from origin: mean {-}50, sd 4"}\NormalTok{,  }\AttributeTok{mean =} \SpecialCharTok{{-}}\DecValTok{50}\NormalTok{, }\AttributeTok{sd =} \DecValTok{4}\NormalTok{),}
  \FunctionTok{c}\NormalTok{(}\AttributeTok{desc =} \StringTok{"far from origin: mean {-}50, sd 10"}\NormalTok{, }\AttributeTok{mean =} \SpecialCharTok{{-}}\DecValTok{50}\NormalTok{, }\AttributeTok{sd =} \DecValTok{10}\NormalTok{)}
\NormalTok{)}
\NormalTok{three\_way }\OtherTok{\textless{}{-}} \FunctionTok{t}\NormalTok{(}\FunctionTok{sapply}\NormalTok{(scenarios, }\ControlFlowTok{function}\NormalTok{(s) \{}
\NormalTok{  x }\OtherTok{\textless{}{-}} \FunctionTok{rnorm}\NormalTok{(}\DecValTok{50}\NormalTok{, }\AttributeTok{mean =} \FunctionTok{as.numeric}\NormalTok{(s[}\StringTok{"mean"}\NormalTok{]), }\AttributeTok{sd =} \FunctionTok{as.numeric}\NormalTok{(s[}\StringTok{"sd"}\NormalTok{]))}
\NormalTok{  lap }\OtherTok{\textless{}{-}} \FunctionTok{laplace\_gaussian}\NormalTok{(x, }\DecValTok{0}\NormalTok{, }\DecValTok{100}\NormalTok{, }\DecValTok{0}\NormalTok{, }\DecValTok{5}\NormalTok{)}
  \FunctionTok{c}\NormalTok{(}\AttributeTok{Laplace =}\NormalTok{ lap}\SpecialCharTok{$}\NormalTok{log\_ml,}
    \AttributeTok{Monte\_Carlo =} \FunctionTok{log\_mar\_mc}\NormalTok{(x, }\DecValTok{0}\NormalTok{, }\DecValTok{100}\NormalTok{, }\DecValTok{0}\NormalTok{, }\DecValTok{5}\NormalTok{, }\FloatTok{1e6}\NormalTok{),}
    \AttributeTok{Quadrature =} \FunctionTok{log\_mar\_mid}\NormalTok{(x, }\DecValTok{0}\NormalTok{, }\DecValTok{100}\NormalTok{, }\DecValTok{0}\NormalTok{, }\DecValTok{5}\NormalTok{, }\DecValTok{100}\NormalTok{))}
\NormalTok{\}))}
\NormalTok{knitr}\SpecialCharTok{::}\FunctionTok{kable}\NormalTok{(}\FunctionTok{cbind}\NormalTok{(}\AttributeTok{scenario =} \FunctionTok{sapply}\NormalTok{(scenarios, }\StringTok{\textasciigrave{}}\AttributeTok{[}\StringTok{\textasciigrave{}}\NormalTok{, }\StringTok{"desc"}\NormalTok{), }\FunctionTok{round}\NormalTok{(three\_way, }\DecValTok{3}\NormalTok{)),}
             \AttributeTok{row.names =} \ConstantTok{FALSE}\NormalTok{, }\AttributeTok{col.names =} \FunctionTok{c}\NormalTok{(}\StringTok{"scenario"}\NormalTok{, }\StringTok{"Laplace"}\NormalTok{, }\StringTok{"Monte Carlo"}\NormalTok{, }\StringTok{"Quadrature (n=100 grid)"}\NormalTok{))}
\end{Highlighting}
\end{Shaded}

{

\begin{longtable}[]{@{}
  >{\raggedright\arraybackslash}p{(\linewidth - 6\tabcolsep) * \real{0.4231}}
  >{\raggedright\arraybackslash}p{(\linewidth - 6\tabcolsep) * \real{0.1154}}
  >{\raggedright\arraybackslash}p{(\linewidth - 6\tabcolsep) * \real{0.1538}}
  >{\raggedright\arraybackslash}p{(\linewidth - 6\tabcolsep) * \real{0.3077}}@{}}

\caption{\label{tbl-three-way}Log marginal likelihood by three
independent methods, for a vague prior N(0, 100) on mu and N(0, 5) on w,
n = 50 observations per scenario. Quadrature uses a 100x100 grid
(unchanged from Unit08); Monte Carlo uses 1e6 prior draws.}

\tabularnewline

\toprule\noalign{}
\begin{minipage}[b]{\linewidth}\raggedright
scenario
\end{minipage} & \begin{minipage}[b]{\linewidth}\raggedright
Laplace
\end{minipage} & \begin{minipage}[b]{\linewidth}\raggedright
Monte Carlo
\end{minipage} & \begin{minipage}[b]{\linewidth}\raggedright
Quadrature (n=100 grid)
\end{minipage} \\
\midrule\noalign{}
\endhead
\bottomrule\noalign{}
\endlastfoot
well-centred: mean 5 & -81.378 & -81.502 & -79.713 \\
well-centred: mean -5 & -84.421 & -84.26 & -85.081 \\
far from origin: mean -50, sd 4 & -141.175 & -141.065 & -271.232 \\
far from origin: mean -50, sd 10 & -191.395 & -191.377 & -271.311 \\

\end{longtable}

}

The pattern is striking. For data centred near the origin, all three
methods agree to within a unit or two on the log scale. But once the
data are far from the origin (mean \(-50\)), \textbf{Laplace and Monte
Carlo still agree closely with each other, while the fixed-grid
quadrature is off by well over 100 log-units} --- the same
boundary-squeeze failure diagnosed at the end of Unit08: the logistic
transform used there is centred at \(0\), so a posterior sitting at
\(\mu\approx-50\) gets squeezed into a sliver of the unit square that a
\(100\times100\) grid cannot resolve. Laplace approximation has no such
failure mode here, and the reason is structural, not a matter of tuning:
\textbf{it has no fixed grid to be squeezed against}. It optimizes to
find \(\theta_0\) wherever the mode actually is, then builds its
approximation \emph{around that point} --- so it automatically
``recentres'' itself for every dataset, which is exactly the manual fix
(recentre the transform on a location estimate) that Unit08 proposed for
quadrature. This is Laplace's core practical advantage over quadrature:
robustness to \emph{where} the posterior happens to sit, at the cost of
assuming it looks Gaussian once you get there.

\subsection{Explore it interactively}\label{explore-it-interactively-1}

The app below lets you move the true mean of the simulated data and
compare the three log marginal likelihood estimates live, so you can
watch quadrature's error grow while Laplace and Monte Carlo continue to
track each other.

\begin{Shaded}
\begin{Highlighting}[]
\NormalTok{\#| label: laplace{-}vs{-}quadrature{-}app}
\NormalTok{\#| standalone: true}
\NormalTok{\#| viewerHeight: 560}

\NormalTok{library(shiny)}

\NormalTok{log\_lik\_gauss  \textless{}{-} function(x, mu, w) sum(dnorm(x, mu, exp(w), log = TRUE))}
\NormalTok{log\_prior\_gauss \textless{}{-} function(mu, w, mu0, sigma\_mu, w0, sigma\_w) \{}
\NormalTok{  dnorm(mu, mu0, sigma\_mu, log = TRUE) + dnorm(w, w0, sigma\_w, log = TRUE)}
\NormalTok{\}}
\NormalTok{neg\_log\_post \textless{}{-} function(theta, x, mu0, sigma\_mu, w0, sigma\_w) \{}
\NormalTok{  {-}log\_lik\_gauss(x, theta[1], theta[2]) {-} log\_prior\_gauss(theta[1], theta[2], mu0, sigma\_mu, w0, sigma\_w)}
\NormalTok{\}}
\NormalTok{laplace\_gaussian \textless{}{-} function(x, mu0, sigma\_mu, w0, sigma\_w) \{}
\NormalTok{  op \textless{}{-} nlm(neg\_log\_post, p = c(mean(x), log(sd(x))), hessian = TRUE,}
\NormalTok{            x = x, mu0 = mu0, sigma\_mu = sigma\_mu, w0 = w0, sigma\_w = sigma\_w)}
\NormalTok{  H \textless{}{-} op$hessian; p \textless{}{-} length(op$estimate)}
\NormalTok{  list(mode = op$estimate, log\_ml = {-}op$minimum + p / 2 * log(2 * pi) {-} 0.5 * log(det(H)))}
\NormalTok{\}}
\NormalTok{log\_sum\_exp \textless{}{-} function(lx) \{ m \textless{}{-} max(lx); m + log(sum(exp(lx {-} m))) \}}
\NormalTok{log\_mar\_mc \textless{}{-} function(x, mu0, sigma\_mu, w0, sigma\_w, iters) \{}
\NormalTok{  mus \textless{}{-} rnorm(iters, mu0, sigma\_mu); ws \textless{}{-} rnorm(iters, w0, sigma\_w)}
\NormalTok{  ll \textless{}{-} vapply(seq\_len(iters), function(i) log\_lik\_gauss(x, mus[i], ws[i]), numeric(1))}
\NormalTok{  log\_sum\_exp(ll) {-} log(iters)}
\NormalTok{\}}
\NormalTok{logit\_inv     \textless{}{-} function(u) log(u / (1 {-} u))}
\NormalTok{log\_der\_logit \textless{}{-} function(u) {-}log(u) {-} log(1 {-} u)}
\NormalTok{log\_g \textless{}{-} function(u, v, x, mu0, sigma\_mu, w0, sigma\_w) \{}
\NormalTok{  mu \textless{}{-} logit\_inv(u); w \textless{}{-} logit\_inv(v)}
\NormalTok{  log\_lik\_gauss(x, mu, w) + log\_prior\_gauss(mu, w, mu0, sigma\_mu, w0, sigma\_w) +}
\NormalTok{    log\_der\_logit(u) + log\_der\_logit(v)}
\NormalTok{\}}
\NormalTok{log\_mar\_mid \textless{}{-} function(x, mu0, sigma\_mu, w0, sigma\_w, n) \{}
\NormalTok{  h \textless{}{-} 1 / n; grid \textless{}{-} (seq\_len(n) {-} 0.5) * h}
\NormalTok{  G \textless{}{-} outer(grid, grid, Vectorize(function(u, v) log\_g(u, v, x, mu0, sigma\_mu, w0, sigma\_w)))}
\NormalTok{  log\_sum\_exp(G) + 2 * log(h)}
\NormalTok{\}}

\NormalTok{ui \textless{}{-} fluidPage(}
\NormalTok{  titlePanel(NULL),}
\NormalTok{  sidebarLayout(}
\NormalTok{    sidebarPanel(}
\NormalTok{      width = 4,}
\NormalTok{      sliderInput("true\_mean", "true mean of the simulated data", min = {-}60, max = 10, value = 5, step = 2),}
\NormalTok{      sliderInput("n\_grid", "quadrature grid points per axis", min = 10, max = 100, value = 100, step = 10),}
\NormalTok{      helpText("Prior: N(0, 100) on mu, N(0, 5) on w = log(sd). As the true mean ",}
\NormalTok{               "moves away from 0, the fixed logistic{-}transform grid squeezes the ",}
\NormalTok{               "posterior into a corner and the quadrature estimate breaks away, ",}
\NormalTok{               "while Laplace (which re{-}optimizes for every dataset) tracks Monte Carlo."),}
\NormalTok{      tableOutput("tab")}
\NormalTok{    ),}
\NormalTok{    mainPanel(width = 8, plotOutput("bars", height = "420px"))}
\NormalTok{  )}
\NormalTok{)}

\NormalTok{server \textless{}{-} function(input, output, session) \{}
\NormalTok{  dat \textless{}{-} reactive(\{ set.seed(1); rnorm(50, input$true\_mean, 2) \})}
\NormalTok{  vals \textless{}{-} reactive(\{}
\NormalTok{    x \textless{}{-} dat()}
\NormalTok{    c(Laplace = laplace\_gaussian(x, 0, 100, 0, 5)$log\_ml,}
\NormalTok{      Monte\_Carlo = \{ set.seed(2); log\_mar\_mc(x, 0, 100, 0, 5, 2e5) \},}
\NormalTok{      Quadrature = log\_mar\_mid(x, 0, 100, 0, 5, input$n\_grid))}
\NormalTok{  \})}
\NormalTok{  output$bars \textless{}{-} renderPlot(\{}
\NormalTok{    v \textless{}{-} vals()}
\NormalTok{    par(mar = c(4, 4, 2, 1))}
\NormalTok{    bp \textless{}{-} barplot(v, col = c("forestgreen", "steelblue", "firebrick"),}
\NormalTok{                  ylab = "estimated log P(y)", main = "three methods, same data")}
\NormalTok{    text(bp, v, sprintf("\%.1f", v), pos = ifelse(v \textgreater{}= 0, 3, 1))}
\NormalTok{  \})}
\NormalTok{  output$tab \textless{}{-} renderTable(\{}
\NormalTok{    v \textless{}{-} vals()}
\NormalTok{    data.frame(method = names(v), \textasciigrave{}log P(y)\textasciigrave{} = as.numeric(v), check.names = FALSE)}
\NormalTok{  \}, digits = 2, rownames = FALSE)}
\NormalTok{\}}

\NormalTok{shinyApp(ui, server)}
\end{Highlighting}
\end{Shaded}

\begin{tcolorbox}[enhanced jigsaw, arc=.35mm, bottomrule=.15mm, bottomtitle=1mm, breakable, colback=white, colbacktitle=quarto-callout-note-color!10!white, colframe=quarto-callout-note-color-frame, coltitle=black, left=2mm, leftrule=.75mm, opacityback=0, opacitybacktitle=0.6, rightrule=.15mm, title=\textcolor{quarto-callout-note-color}{\faInfo}\hspace{0.5em}{Reproducing this figure}, titlerule=0mm, toprule=.15mm, toptitle=1mm]

As in earlier chapters, this app uses the
\href{https://github.com/quarto-ext/shinylive}{shinylive} extension
(\texttt{quarto\ add\ quarto-ext/shinylive}, with
\texttt{filters:\ {[}shinylive{]}} in the document YAML), compiled to
WebAssembly so it runs in the browser with no Shiny server. Only base R
and \texttt{shiny} are available inside the app.

\end{tcolorbox}

\section{Summary}\label{summary-3}

\begin{itemize}
\tightlist
\item
  The Laplace approximation expands \(h=-\log f\) to second order at its
  minimum \(\theta_0\) and integrates the resulting Gaussian exactly:
  \(\theta\mid\text{data} \approx N_p(\theta_0, H^{-1})\) and
  \(\int f(\theta)\,d\theta \approx f(\theta_0)(2\pi)^{p/2}|H|^{-1/2}\).
\item
  It needs only an optimizer and a Hessian --- the same machinery
  already used to find the MLE --- and its cost does not grow with
  dimension the way a quadrature grid's does.
\item
  BIC is the large-\(n\) simplification of the Laplace marginal
  likelihood, dropping \(O(1)\) terms that do not depend on the sample
  size.
\item
  On the running Gaussian marginal-likelihood example, Laplace matched
  Monte Carlo everywhere quadrature was tested, and remained accurate
  exactly where quadrature's fixed grid broke down --- because Laplace
  re-optimizes for every dataset instead of committing to one transform
  in advance.
\item
  It is deterministic, fast, and accurate to \(O(1/n)\), but blind to
  skewness and multimodality; the Monte Carlo methods of the next
  chapters remove these limitations at higher computational cost.
\end{itemize}

\bookmarksetup{startatroot}

\chapter{Rejection Sampling}\label{rejection-sampling}

\begin{Shaded}
\begin{Highlighting}[]
\FunctionTok{library}\NormalTok{(ars)}
\end{Highlighting}
\end{Shaded}

Monte Carlo integration estimates
\(E_f[a(\theta)] = \int a(\theta) f(\theta)\,d\theta\) from iid draws
\(\theta_1,\ldots,\theta_n \sim f\). The catch is the sampling step
itself: R has built-in samplers for standard families, but a posterior
\(f(\theta) \propto L(\theta)\pi(\theta)\) is typically known only up to
a constant, and there is no generic way to draw from it directly.
\textbf{Rejection sampling} solves this by producing \emph{exact} iid
draws from \(f\) using draws from an easier distribution \(g\),
following the recipe \emph{reject, reweight, repeat}.

\section{The Envelope and the
Algorithm}\label{the-envelope-and-the-algorithm}

\subsection{The envelope condition}\label{the-envelope-condition}

Suppose we can find a function \(g(x)\) such that

\begin{itemize}
\tightlist
\item
  \(g(x) \ge f(x)\) for every \(x\) --- an \textbf{envelope} of \(f\),
  and
\item
  we can sample from the normalized density \(g(x)/\!\int g(x)\,dx\).
\end{itemize}

Neither \(f\) nor \(g\) needs to be normalized: \(f\) may be an
unnormalized posterior and \(g\) a scaled density \(M\,g_0(x)\) for a
sampler \(g_0\) we already have. The overall acceptance rate of the
algorithm below turns out to be the ratio of the two areas, \[
P(\text{accept}) = \frac{\int f(x)\,dx}{\int g(x)\,dx},
\] so a tight envelope wastes few draws and a loose one wastes many ---
the entire design problem is finding a \(g\) that is both easy to sample
and close to \(f\).

\subsection{The algorithm}\label{the-algorithm-1}

\textbf{Input}: target \(f\), envelope \(g \ge f\) with a sampler for
\(g(x)/\!\int g(x)\,dx\). For each draw needed, repeat until acceptance:

\begin{enumerate}
\def\labelenumi{\arabic{enumi}.}
\tightlist
\item
  Draw \(x \sim g(x)/\!\int g(x)\,dx\).
\item
  Compute \(r = f(x)/g(x)\).
\item
  Draw \(U \sim \text{Unif}(0,1)\).
\item
  If \(U < r\), accept \(x\); otherwise return to step 1.
\end{enumerate}

Only the ratio \(f/g\) is ever evaluated, so unknown normalizing
constants in \(f\) (and in \(g\)) cancel automatically as long as the
\emph{versions actually used} satisfy \(g \ge f\). The number of
proposals needed per accepted draw is geometric with mean
\(1/P(\text{accept})\).

\subsection{An equivalent geometric
picture}\label{an-equivalent-geometric-picture}

An illuminating restatement: draw \(X \sim g\), then draw a height
\(V \sim \text{Unif}(0, g(X))\) (equivalently \(V = U\,g(X)\)), and
accept \(X\) if \(V < f(X)\). This is the same algorithm, since
\(V < f(X) \iff U < f(X)/g(X)\), but it exposes \emph{why} rejection
sampling works:

\begin{itemize}
\tightlist
\item
  Steps ``draw \(X\sim g\), then \(V\sim\text{Unif}(0,g(X))\)'' place
  \((X,V)\) \textbf{uniformly} in the region under \(g\),
  \(G = \{(x,v): 0<v<g(x)\}\).
\item
  Keeping only points with \(v < f(x)\) restricts to the sub-region
  under \(f\), \(F = \{(x,v): 0<v<f(x)\} \subset G\).
\item
  A uniform point on \(G\), conditioned on landing in \(F\), is uniform
  on \(F\). The \(x\)-coordinate of a uniform point on \(F\) has density
  proportional to the \emph{height} of \(F\) at \(x\), which is
  \(f(x)\). That is exactly a draw from \(f/\!\int f\).
\end{itemize}

\begin{Shaded}
\begin{Highlighting}[]
\NormalTok{f\_demo }\OtherTok{\textless{}{-}} \ControlFlowTok{function}\NormalTok{(x) }\FunctionTok{exp}\NormalTok{(}\SpecialCharTok{{-}}\NormalTok{x}\SpecialCharTok{\^{}}\DecValTok{2} \SpecialCharTok{/} \DecValTok{2}\NormalTok{)}
\NormalTok{g\_demo }\OtherTok{\textless{}{-}} \ControlFlowTok{function}\NormalTok{(x) }\FunctionTok{exp}\NormalTok{(}\FloatTok{0.5} \SpecialCharTok{{-}} \FunctionTok{abs}\NormalTok{(x))}
\NormalTok{xs }\OtherTok{\textless{}{-}} \FunctionTok{seq}\NormalTok{(}\SpecialCharTok{{-}}\DecValTok{6}\NormalTok{, }\DecValTok{6}\NormalTok{, }\AttributeTok{length.out =} \DecValTok{500}\NormalTok{)}

\FunctionTok{par}\NormalTok{(}\AttributeTok{mfrow =} \FunctionTok{c}\NormalTok{(}\DecValTok{1}\NormalTok{, }\DecValTok{2}\NormalTok{), }\AttributeTok{mar =} \FunctionTok{c}\NormalTok{(}\DecValTok{4}\NormalTok{, }\DecValTok{4}\NormalTok{, }\DecValTok{2}\NormalTok{, }\DecValTok{1}\NormalTok{))}
\FunctionTok{plot}\NormalTok{(xs, }\FunctionTok{g\_demo}\NormalTok{(xs), }\AttributeTok{type =} \StringTok{"n"}\NormalTok{, }\AttributeTok{ylim =} \FunctionTok{c}\NormalTok{(}\DecValTok{0}\NormalTok{, }\FloatTok{1.75}\NormalTok{), }\AttributeTok{xlab =} \StringTok{"x"}\NormalTok{, }\AttributeTok{ylab =} \StringTok{""}\NormalTok{)}
\FunctionTok{polygon}\NormalTok{(}\FunctionTok{c}\NormalTok{(xs, }\FunctionTok{rev}\NormalTok{(xs)), }\FunctionTok{c}\NormalTok{(}\FunctionTok{g\_demo}\NormalTok{(xs), }\FunctionTok{rep}\NormalTok{(}\DecValTok{0}\NormalTok{, }\DecValTok{500}\NormalTok{)), }\AttributeTok{col =} \FunctionTok{adjustcolor}\NormalTok{(}\StringTok{"steelblue"}\NormalTok{, }\FloatTok{0.3}\NormalTok{), }\AttributeTok{border =} \ConstantTok{NA}\NormalTok{)}
\FunctionTok{polygon}\NormalTok{(}\FunctionTok{c}\NormalTok{(xs, }\FunctionTok{rev}\NormalTok{(xs)), }\FunctionTok{c}\NormalTok{(}\FunctionTok{f\_demo}\NormalTok{(xs), }\FunctionTok{rep}\NormalTok{(}\DecValTok{0}\NormalTok{, }\DecValTok{500}\NormalTok{)), }\AttributeTok{col =} \FunctionTok{adjustcolor}\NormalTok{(}\StringTok{"firebrick"}\NormalTok{, }\FloatTok{0.4}\NormalTok{), }\AttributeTok{border =} \ConstantTok{NA}\NormalTok{)}
\FunctionTok{lines}\NormalTok{(xs, }\FunctionTok{g\_demo}\NormalTok{(xs), }\AttributeTok{col =} \StringTok{"steelblue"}\NormalTok{, }\AttributeTok{lwd =} \DecValTok{3}\NormalTok{); }\FunctionTok{lines}\NormalTok{(xs, }\FunctionTok{f\_demo}\NormalTok{(xs), }\AttributeTok{col =} \StringTok{"firebrick"}\NormalTok{, }\AttributeTok{lwd =} \DecValTok{3}\NormalTok{)}
\FunctionTok{legend}\NormalTok{(}\StringTok{"topright"}\NormalTok{, }\AttributeTok{bty =} \StringTok{"n"}\NormalTok{, }\AttributeTok{fill =} \FunctionTok{adjustcolor}\NormalTok{(}\FunctionTok{c}\NormalTok{(}\StringTok{"steelblue"}\NormalTok{, }\StringTok{"firebrick"}\NormalTok{), }\FloatTok{0.35}\NormalTok{), }\AttributeTok{border =} \ConstantTok{NA}\NormalTok{,}
       \AttributeTok{legend =} \FunctionTok{c}\NormalTok{(}\StringTok{"area under envelope g"}\NormalTok{, }\StringTok{"area under target f"}\NormalTok{))}

\FunctionTok{set.seed}\NormalTok{(}\DecValTok{1}\NormalTok{)}
\NormalTok{N }\OtherTok{\textless{}{-}} \DecValTok{600}
\NormalTok{rlaplace }\OtherTok{\textless{}{-}} \ControlFlowTok{function}\NormalTok{(n) \{ u }\OtherTok{\textless{}{-}} \FunctionTok{runif}\NormalTok{(n); }\FunctionTok{ifelse}\NormalTok{(u }\SpecialCharTok{\textless{}} \FloatTok{0.5}\NormalTok{, }\FunctionTok{log}\NormalTok{(}\DecValTok{2} \SpecialCharTok{*}\NormalTok{ u), }\SpecialCharTok{{-}}\FunctionTok{log}\NormalTok{(}\DecValTok{2} \SpecialCharTok{*}\NormalTok{ (}\DecValTok{1} \SpecialCharTok{{-}}\NormalTok{ u))) \}}
\NormalTok{xp }\OtherTok{\textless{}{-}} \FunctionTok{rlaplace}\NormalTok{(N); vp }\OtherTok{\textless{}{-}} \FunctionTok{runif}\NormalTok{(N) }\SpecialCharTok{*} \FunctionTok{g\_demo}\NormalTok{(xp); acc }\OtherTok{\textless{}{-}}\NormalTok{ vp }\SpecialCharTok{\textless{}} \FunctionTok{f\_demo}\NormalTok{(xp)}
\FunctionTok{plot}\NormalTok{(xs, }\FunctionTok{g\_demo}\NormalTok{(xs), }\AttributeTok{type =} \StringTok{"l"}\NormalTok{, }\AttributeTok{lwd =} \DecValTok{3}\NormalTok{, }\AttributeTok{col =} \StringTok{"steelblue"}\NormalTok{, }\AttributeTok{ylim =} \FunctionTok{c}\NormalTok{(}\DecValTok{0}\NormalTok{, }\FloatTok{1.75}\NormalTok{), }\AttributeTok{xlab =} \StringTok{"x"}\NormalTok{, }\AttributeTok{ylab =} \StringTok{"v"}\NormalTok{)}
\FunctionTok{lines}\NormalTok{(xs, }\FunctionTok{f\_demo}\NormalTok{(xs), }\AttributeTok{lwd =} \DecValTok{3}\NormalTok{, }\AttributeTok{col =} \StringTok{"firebrick"}\NormalTok{)}
\FunctionTok{points}\NormalTok{(xp[}\SpecialCharTok{!}\NormalTok{acc], vp[}\SpecialCharTok{!}\NormalTok{acc], }\AttributeTok{pch =} \DecValTok{4}\NormalTok{, }\AttributeTok{col =} \FunctionTok{adjustcolor}\NormalTok{(}\StringTok{"gray40"}\NormalTok{, }\FloatTok{0.7}\NormalTok{), }\AttributeTok{cex =} \FloatTok{0.7}\NormalTok{)}
\FunctionTok{points}\NormalTok{(xp[acc], vp[acc], }\AttributeTok{pch =} \DecValTok{19}\NormalTok{, }\AttributeTok{col =} \FunctionTok{adjustcolor}\NormalTok{(}\StringTok{"firebrick"}\NormalTok{, }\FloatTok{0.7}\NormalTok{), }\AttributeTok{cex =} \FloatTok{0.7}\NormalTok{)}
\FunctionTok{legend}\NormalTok{(}\StringTok{"topright"}\NormalTok{, }\AttributeTok{bty =} \StringTok{"n"}\NormalTok{, }\AttributeTok{pch =} \FunctionTok{c}\NormalTok{(}\DecValTok{19}\NormalTok{, }\DecValTok{4}\NormalTok{), }\AttributeTok{col =} \FunctionTok{c}\NormalTok{(}\StringTok{"firebrick"}\NormalTok{, }\StringTok{"gray40"}\NormalTok{),}
       \AttributeTok{legend =} \FunctionTok{c}\NormalTok{(}\FunctionTok{sprintf}\NormalTok{(}\StringTok{"accepted (\%d)"}\NormalTok{, }\FunctionTok{sum}\NormalTok{(acc)), }\FunctionTok{sprintf}\NormalTok{(}\StringTok{"rejected (\%d)"}\NormalTok{, }\FunctionTok{sum}\NormalTok{(}\SpecialCharTok{!}\NormalTok{acc))))}
\end{Highlighting}
\end{Shaded}

\begin{figure}[H]

\centering{

\pandocbounded{\includegraphics[keepaspectratio]{chapters/unit10_rejection_sampling_files/figure-pdf/fig-envelope-scatter-1.pdf}}

}

\caption{\label{fig-envelope-scatter}Left: target \(f(x)=e^{-x^2/2}\)
(unnormalized standard normal) under its Laplace envelope
\(g(x) = e^{1/2-|x|}\), with the two areas shaded. Right: 600 candidate
points drawn uniformly under \(g\); the accepted ones (red, under \(f\)
too) are uniformly distributed under \(f\), so their \(x\)-coordinates
are draws from the normalized \(f\).}

\end{figure}%

\section{Example: Sampling a Normal from a Laplace
Envelope}\label{example-sampling-a-normal-from-a-laplace-envelope}

The figure above already uses a concrete, closed-form example worth
stating in full, because it is simple enough to check every claim
exactly. The target is the unnormalized standard normal density
\(f(x) = e^{-x^2/2}\). Take \(g_0(x) = \tfrac12 e^{-|x|}\), the standard
Laplace (double-exponential) density, which is trivial to sample by
inversion. Since \(e^{-x^2/2} \le e^{1/2-|x|}\) for all \(x\) (equality
exactly at \(|x| = 1\), where the two curves in the figure touch), the
envelope \[
g(x) = 2e^{1/2}\,g_0(x) = e^{1/2 - |x|} \ \ge\ f(x)
\] is valid, and the exact acceptance rate is \[
P(\text{accept}) = \frac{\int f(x)\,dx}{\int g(x)\,dx}
 = \frac{\sqrt{2\pi}}{2e^{1/2}} \approx 0.760.
\]

\begin{Shaded}
\begin{Highlighting}[]
\NormalTok{rnorm\_reject }\OtherTok{\textless{}{-}} \ControlFlowTok{function}\NormalTok{(n) \{}
\NormalTok{    out }\OtherTok{\textless{}{-}} \FunctionTok{numeric}\NormalTok{(}\DecValTok{0}\NormalTok{); tries }\OtherTok{\textless{}{-}} \DecValTok{0}
    \ControlFlowTok{while}\NormalTok{ (}\FunctionTok{length}\NormalTok{(out) }\SpecialCharTok{\textless{}}\NormalTok{ n) \{}
\NormalTok{        m }\OtherTok{\textless{}{-}}\NormalTok{ n }\SpecialCharTok{{-}} \FunctionTok{length}\NormalTok{(out)}
\NormalTok{        x }\OtherTok{\textless{}{-}} \FunctionTok{rlaplace}\NormalTok{(m); tries }\OtherTok{\textless{}{-}}\NormalTok{ tries }\SpecialCharTok{+}\NormalTok{ m}
\NormalTok{        out }\OtherTok{\textless{}{-}} \FunctionTok{c}\NormalTok{(out, x[}\FunctionTok{runif}\NormalTok{(m) }\SpecialCharTok{\textless{}} \FunctionTok{f\_demo}\NormalTok{(x) }\SpecialCharTok{/} \FunctionTok{g\_demo}\NormalTok{(x)])}
\NormalTok{    \}}
    \FunctionTok{list}\NormalTok{(}\AttributeTok{x =}\NormalTok{ out, }\AttributeTok{acceptance =}\NormalTok{ n }\SpecialCharTok{/}\NormalTok{ tries)}
\NormalTok{\}}
\FunctionTok{set.seed}\NormalTok{(}\DecValTok{2}\NormalTok{)}
\NormalTok{res }\OtherTok{\textless{}{-}} \FunctionTok{rnorm\_reject}\NormalTok{(}\FloatTok{2e5}\NormalTok{)}
\NormalTok{knitr}\SpecialCharTok{::}\FunctionTok{kable}\NormalTok{(}\FunctionTok{data.frame}\NormalTok{(}
    \AttributeTok{quantity =} \FunctionTok{c}\NormalTok{(}\StringTok{"acceptance rate"}\NormalTok{, }\StringTok{"sample mean"}\NormalTok{, }\StringTok{"sample sd"}\NormalTok{),}
    \AttributeTok{theoretical =} \FunctionTok{c}\NormalTok{(}\FunctionTok{sqrt}\NormalTok{(}\DecValTok{2} \SpecialCharTok{*}\NormalTok{ pi) }\SpecialCharTok{/}\NormalTok{ (}\DecValTok{2} \SpecialCharTok{*} \FunctionTok{sqrt}\NormalTok{(}\FunctionTok{exp}\NormalTok{(}\DecValTok{1}\NormalTok{))), }\DecValTok{0}\NormalTok{, }\DecValTok{1}\NormalTok{),}
    \AttributeTok{simulated =} \FunctionTok{c}\NormalTok{(res}\SpecialCharTok{$}\NormalTok{acceptance, }\FunctionTok{mean}\NormalTok{(res}\SpecialCharTok{$}\NormalTok{x), }\FunctionTok{sd}\NormalTok{(res}\SpecialCharTok{$}\NormalTok{x))}
\NormalTok{), }\AttributeTok{digits =} \DecValTok{4}\NormalTok{)}
\end{Highlighting}
\end{Shaded}

\begin{longtable}[]{@{}lrr@{}}
\caption{Theoretical vs.~simulated acceptance rate and sample moments
for the Laplace-envelope normal sampler.}\tabularnewline
\toprule\noalign{}
quantity & theoretical & simulated \\
\midrule\noalign{}
\endfirsthead
\toprule\noalign{}
quantity & theoretical & simulated \\
\midrule\noalign{}
\endhead
\bottomrule\noalign{}
\endlastfoot
acceptance rate & 0.7602 & 0.7595 \\
sample mean & 0.0000 & 0.0031 \\
sample sd & 1.0000 & 0.9991 \\
\end{longtable}

Roughly 3 in 4 proposals are kept --- a tight, efficient envelope, and
the benchmark against which the next example's Cauchy envelope should be
judged.

\subsection{Trying it interactively: tight vs.~loose
envelopes}\label{trying-it-interactively-tight-vs.-loose-envelopes}

The app below scales the envelope by a factor \(c \ge 1\),
\(g_c(x) = c\, e^{1/2-|x|}\): larger \(c\) keeps the envelope valid but
wastes more draws (acceptance rate \(= 0.760/c\)). Dragging \(c\) below
\(1\) shows what an \emph{invalid} envelope looks like --- the target
pokes through it, and rejection sampling would silently under-represent
that region.

\begin{Shaded}
\begin{Highlighting}[]
\NormalTok{\#| label: envelope{-}tightness{-}app}
\NormalTok{\#| standalone: true}
\NormalTok{\#| viewerHeight: 560}

\NormalTok{library(shiny)}

\NormalTok{f\_demo \textless{}{-} function(x) exp({-}x\^{}2 / 2)}
\NormalTok{g\_demo \textless{}{-} function(x, c) c * exp(0.5 {-} abs(x))}
\NormalTok{rlaplace \textless{}{-} function(n) \{ u \textless{}{-} runif(n); ifelse(u \textless{} 0.5, log(2 * u), {-}log(2 * (1 {-} u))) \}}

\NormalTok{ui \textless{}{-} fluidPage(}
\NormalTok{    titlePanel(NULL),}
\NormalTok{    sidebarLayout(}
\NormalTok{        sidebarPanel(}
\NormalTok{            width = 4,}
\NormalTok{            sliderInput("c", "envelope scale c (g = c * e\^{}(1/2 {-} |x|))",}
\NormalTok{                        min = 0.6, max = 3, value = 1.2, step = 0.05),}
\NormalTok{            helpText("c \textgreater{}= 1 keeps g above f everywhere: a valid, if looser, ",}
\NormalTok{                     "envelope. c \textless{} 1 breaks the envelope condition near ",}
\NormalTok{                     "|x| = 1, where the target actually exceeds it."),}
\NormalTok{            actionButton("resample", "draw new points"),}
\NormalTok{            tableOutput("tab")}
\NormalTok{        ),}
\NormalTok{        mainPanel(}
\NormalTok{            width = 8,}
\NormalTok{            plotOutput("scatter", height = "440px")}
\NormalTok{        )}
\NormalTok{    )}
\NormalTok{)}

\NormalTok{server \textless{}{-} function(input, output, session) \{}

\NormalTok{    pts \textless{}{-} eventReactive(list(input$c, input$resample), \{}
\NormalTok{        set.seed(input$resample + 1)}
\NormalTok{        N \textless{}{-} 500}
\NormalTok{        x \textless{}{-} rlaplace(N)}
\NormalTok{        v \textless{}{-} runif(N) * g\_demo(x, input$c)}
\NormalTok{        list(x = x, v = v, acc = v \textless{} f\_demo(x))}
\NormalTok{    \})}

\NormalTok{    output$scatter \textless{}{-} renderPlot(\{}
\NormalTok{        p \textless{}{-} pts()}
\NormalTok{        xs \textless{}{-} seq({-}6, 6, length.out = 500)}
\NormalTok{        par(mar = c(4, 4, 2, 1))}
\NormalTok{        plot(xs, g\_demo(xs, input$c), type = "l", lwd = 3, col = "steelblue",}
\NormalTok{             ylim = c(0, max(2.2, 1.05 * input$c * exp(0.5))),}
\NormalTok{             xlab = "x", ylab = "v", main = sprintf("c = \%.2f", input$c))}
\NormalTok{        lines(xs, f\_demo(xs), lwd = 3, col = "firebrick")}
\NormalTok{        bad \textless{}{-} f\_demo(xs) \textgreater{} g\_demo(xs, input$c)}
\NormalTok{        if (any(bad)) lines(xs[bad], f\_demo(xs)[bad], col = "black", lwd = 5)}
\NormalTok{        points(p$x[!p$acc], p$v[!p$acc], pch = 4, col = adjustcolor("gray40", 0.6), cex = 0.7)}
\NormalTok{        points(p$x[p$acc], p$v[p$acc], pch = 19, col = adjustcolor("firebrick", 0.7), cex = 0.7)}
\NormalTok{        legend("topright", bty = "n", lwd = c(3, 3, 5), col = c("firebrick", "steelblue", "black"),}
\NormalTok{               legend = c("target f", "envelope g", "f exceeds g here: envelope invalid"))}
\NormalTok{    \})}

\NormalTok{    output$tab \textless{}{-} renderTable(\{}
\NormalTok{        p \textless{}{-} pts()}
\NormalTok{        data.frame(quantity = c("theoretical acceptance", "observed acceptance", "envelope valid?"),}
\NormalTok{                   value = c(sprintf("\%.3f", 0.7602 / input$c), sprintf("\%.3f", mean(p$acc)),}
\NormalTok{                             ifelse(input$c \textgreater{}= 1, "yes", "NO (c \textless{} 1)")))}
\NormalTok{    \}, colnames = FALSE)}
\NormalTok{\}}

\NormalTok{shinyApp(ui, server)}
\end{Highlighting}
\end{Shaded}

\begin{tcolorbox}[enhanced jigsaw, arc=.35mm, bottomrule=.15mm, bottomtitle=1mm, breakable, colback=white, colbacktitle=quarto-callout-note-color!10!white, colframe=quarto-callout-note-color-frame, coltitle=black, left=2mm, leftrule=.75mm, opacityback=0, opacitybacktitle=0.6, rightrule=.15mm, title=\textcolor{quarto-callout-note-color}{\faInfo}\hspace{0.5em}{Reproducing this figure}, titlerule=0mm, toprule=.15mm, toptitle=1mm]

The chunk above uses the
\href{https://github.com/quarto-ext/shinylive}{shinylive} extension,
compiled to WebAssembly so it runs in the reader's browser with no Shiny
server (\texttt{quarto\ add\ quarto-ext/shinylive}, then
\texttt{filters:\ {[}shinylive{]}} in the document or project YAML).
Only base R and \texttt{shiny} are available inside the app.

\end{tcolorbox}

\section{Example: Sampling a Gamma
Distribution}\label{example-sampling-a-gamma-distribution}

The Laplace envelope above was designed to match a normal target. A
different target needs a different envelope. Here is one built for
\(\text{Gamma}(\alpha, 1)\), \(\alpha > 1\), whose unnormalized
log-density is \((\alpha-1)\log x - x\).

\begin{Shaded}
\begin{Highlighting}[]
\DocumentationTok{\#\#\# log of a function that is (for alpha \textgreater{} 2) always above the Gamma(alpha, 1) log{-}density}
\NormalTok{log\_g\_gamma }\OtherTok{\textless{}{-}} \ControlFlowTok{function}\NormalTok{(x, alpha) \{}
\NormalTok{    (alpha }\SpecialCharTok{{-}} \DecValTok{1}\NormalTok{) }\SpecialCharTok{*}\NormalTok{ (}\FunctionTok{log}\NormalTok{(alpha }\SpecialCharTok{{-}} \DecValTok{1}\NormalTok{) }\SpecialCharTok{{-}} \DecValTok{1}\NormalTok{) }\SpecialCharTok{{-}} \FunctionTok{log}\NormalTok{(}\DecValTok{1} \SpecialCharTok{+}\NormalTok{ (x }\SpecialCharTok{{-}}\NormalTok{ (alpha }\SpecialCharTok{{-}} \DecValTok{1}\NormalTok{))}\SpecialCharTok{\^{}}\DecValTok{2} \SpecialCharTok{/}\NormalTok{ (}\DecValTok{2} \SpecialCharTok{*}\NormalTok{ alpha }\SpecialCharTok{{-}} \DecValTok{1}\NormalTok{))}
\NormalTok{\}}
\end{Highlighting}
\end{Shaded}

\subsection{Where the envelope comes
from}\label{where-the-envelope-comes-from}

\texttt{log\_g\_gamma} is not an arbitrary formula. Up to an additive
constant that depends only on \(\alpha\), it is the log-density of a
\textbf{rescaled, shifted Cauchy} distribution: writing
\(Y = (\alpha-1) + \sqrt{2\alpha-1}\,C\) for a standard Cauchy \(C\), a
change of variables gives
\(\log f_Y(y) = -\log\pi - \tfrac12\log(2\alpha-1) - \log\!\big(1 + (y-(\alpha-1))^2/(2\alpha-1)\big)\),
which differs from \texttt{log\_g\_gamma(y,\ alpha)} only by a term that
does not depend on \(y\). We can check this numerically rather than take
it on faith:

\begin{Shaded}
\begin{Highlighting}[]
\NormalTok{x\_check }\OtherTok{\textless{}{-}} \FunctionTok{seq}\NormalTok{(}\FloatTok{0.1}\NormalTok{, }\DecValTok{10}\NormalTok{, }\AttributeTok{by =} \FloatTok{0.37}\NormalTok{)}
\NormalTok{alpha\_check }\OtherTok{\textless{}{-}} \FloatTok{3.3}
\NormalTok{log\_f\_Y }\OtherTok{\textless{}{-}} \FunctionTok{dcauchy}\NormalTok{((x\_check }\SpecialCharTok{{-}}\NormalTok{ (alpha\_check }\SpecialCharTok{{-}} \DecValTok{1}\NormalTok{)) }\SpecialCharTok{/} \FunctionTok{sqrt}\NormalTok{(}\DecValTok{2} \SpecialCharTok{*}\NormalTok{ alpha\_check }\SpecialCharTok{{-}} \DecValTok{1}\NormalTok{), }\AttributeTok{log =} \ConstantTok{TRUE}\NormalTok{) }\SpecialCharTok{{-}}
    \FunctionTok{log}\NormalTok{(}\FunctionTok{sqrt}\NormalTok{(}\DecValTok{2} \SpecialCharTok{*}\NormalTok{ alpha\_check }\SpecialCharTok{{-}} \DecValTok{1}\NormalTok{))}
\FunctionTok{range}\NormalTok{(}\FunctionTok{log\_g\_gamma}\NormalTok{(x\_check, alpha\_check) }\SpecialCharTok{{-}}\NormalTok{ log\_f\_Y)   }\CommentTok{\# constant across x confirms the claim}
\end{Highlighting}
\end{Shaded}

\begin{verbatim}
[1] 1.621804 1.621804
\end{verbatim}

So sampling the envelope is just sampling a (rescaled) Cauchy, and the
constant in \texttt{log\_g\_gamma} was chosen so that \(g\) touches the
Gamma density exactly at the shared mode \(x = \alpha - 1\). Whether
\(g\) stays \emph{above} the Gamma density everywhere else is a separate
question, and the docstring's claim ``alpha must be \textgreater{} 2''
is a claim about exactly that:

\begin{Shaded}
\begin{Highlighting}[]
\NormalTok{alphas\_check }\OtherTok{\textless{}{-}} \FunctionTok{c}\NormalTok{(}\FloatTok{1.1}\NormalTok{, }\FloatTok{1.5}\NormalTok{, }\FloatTok{1.9}\NormalTok{, }\FloatTok{1.99}\NormalTok{, }\FloatTok{2.0}\NormalTok{, }\FloatTok{2.01}\NormalTok{, }\FloatTok{2.1}\NormalTok{, }\FloatTok{2.5}\NormalTok{, }\DecValTok{3}\NormalTok{, }\FloatTok{4.5}\NormalTok{)}
\NormalTok{xv }\OtherTok{\textless{}{-}} \FunctionTok{seq}\NormalTok{(}\FloatTok{1e{-}6}\NormalTok{, }\DecValTok{60}\NormalTok{, }\AttributeTok{length.out =} \FloatTok{2e4}\NormalTok{)}
\NormalTok{gaps }\OtherTok{\textless{}{-}} \FunctionTok{vapply}\NormalTok{(alphas\_check, }\ControlFlowTok{function}\NormalTok{(a) }\FunctionTok{min}\NormalTok{(}\FunctionTok{log\_g\_gamma}\NormalTok{(xv, a) }\SpecialCharTok{{-}} \FunctionTok{dgamma}\NormalTok{(xv, a, }\AttributeTok{log =} \ConstantTok{TRUE}\NormalTok{)), }\FunctionTok{numeric}\NormalTok{(}\DecValTok{1}\NormalTok{))}
\NormalTok{knitr}\SpecialCharTok{::}\FunctionTok{kable}\NormalTok{(}\FunctionTok{data.frame}\NormalTok{(}\AttributeTok{alpha =}\NormalTok{ alphas\_check, }\AttributeTok{min\_log\_gap =}\NormalTok{ gaps, }\AttributeTok{envelope\_valid =}\NormalTok{ gaps }\SpecialCharTok{\textgreater{}=} \DecValTok{0}\NormalTok{), }\AttributeTok{digits =} \DecValTok{5}\NormalTok{)}
\end{Highlighting}
\end{Shaded}

{

\begin{longtable}[]{@{}rrl@{}}

\caption{\label{tbl-envelope-threshold}Minimum of log g(x) - log f(x)
over x \textgreater{} 0, for several alpha. Negative means the envelope
is violated somewhere.}

\tabularnewline

\toprule\noalign{}
alpha & min\_log\_gap & envelope\_valid \\
\midrule\noalign{}
\endhead
\bottomrule\noalign{}
\endlastfoot
1.10 & -0.04987 & FALSE \\
1.50 & -0.12078 & FALSE \\
1.90 & -0.03898 & FALSE \\
1.99 & -0.00420 & FALSE \\
2.00 & 0.00000 & TRUE \\
2.01 & 0.00426 & TRUE \\
2.10 & 0.04544 & TRUE \\
2.50 & 0.28468 & TRUE \\
3.00 & 0.69315 & TRUE \\
4.50 & 2.45374 & TRUE \\

\end{longtable}

}

The gap is exactly \(0\) at \(\alpha = 2\) and negative below it: the
Cauchy-based envelope is tangent to the Gamma density at \(\alpha = 2\)
and crosses below it for smaller \(\alpha\), so \(\alpha > 2\) is
precisely the validity condition, not a conservative rule of thumb.

\subsection{Sampling function}\label{sampling-function}

A batch-and-refill implementation is both clearer and much faster than
drawing one candidate at a time in an R loop: generate a vector of
candidates, vectorize the accept/reject test, keep the survivors, and
top up only the shortfall.

\begin{Shaded}
\begin{Highlighting}[]
\NormalTok{sample\_gamma\_rej }\OtherTok{\textless{}{-}} \ControlFlowTok{function}\NormalTok{(n, alpha, }\AttributeTok{max\_tries =} \FloatTok{2e6}\NormalTok{) \{}
\NormalTok{    out }\OtherTok{\textless{}{-}} \FunctionTok{numeric}\NormalTok{(}\DecValTok{0}\NormalTok{); tries }\OtherTok{\textless{}{-}} \DecValTok{0}
    \ControlFlowTok{while}\NormalTok{ (}\FunctionTok{length}\NormalTok{(out) }\SpecialCharTok{\textless{}}\NormalTok{ n) \{}
        \ControlFlowTok{if}\NormalTok{ (tries }\SpecialCharTok{\textgreater{}}\NormalTok{ max\_tries) }\FunctionTok{stop}\NormalTok{(}\StringTok{"acceptance rate too low for this alpha: envelope is too loose"}\NormalTok{)}
\NormalTok{        m }\OtherTok{\textless{}{-}}\NormalTok{ n }\SpecialCharTok{{-}} \FunctionTok{length}\NormalTok{(out)}
\NormalTok{        x }\OtherTok{\textless{}{-}} \FunctionTok{rcauchy}\NormalTok{(m) }\SpecialCharTok{*} \FunctionTok{sqrt}\NormalTok{(}\DecValTok{2} \SpecialCharTok{*}\NormalTok{ alpha }\SpecialCharTok{{-}} \DecValTok{1}\NormalTok{) }\SpecialCharTok{+}\NormalTok{ (alpha }\SpecialCharTok{{-}} \DecValTok{1}\NormalTok{)}
\NormalTok{        tries }\OtherTok{\textless{}{-}}\NormalTok{ tries }\SpecialCharTok{+}\NormalTok{ m}
\NormalTok{        keep }\OtherTok{\textless{}{-}} \FunctionTok{log}\NormalTok{(}\FunctionTok{runif}\NormalTok{(m)) }\SpecialCharTok{\textless{}} \FunctionTok{dgamma}\NormalTok{(x, }\AttributeTok{shape =}\NormalTok{ alpha, }\AttributeTok{log =} \ConstantTok{TRUE}\NormalTok{) }\SpecialCharTok{{-}} \FunctionTok{log\_g\_gamma}\NormalTok{(x, alpha)}
\NormalTok{        out }\OtherTok{\textless{}{-}} \FunctionTok{c}\NormalTok{(out, x[keep])}
\NormalTok{    \}}
    \FunctionTok{structure}\NormalTok{(out[}\FunctionTok{seq\_len}\NormalTok{(n)], }\AttributeTok{accept.rate =}\NormalTok{ n }\SpecialCharTok{/}\NormalTok{ tries)}
\NormalTok{\}}
\end{Highlighting}
\end{Shaded}

\subsection{Test cases and efficiency}\label{test-cases-and-efficiency}

\begin{Shaded}
\begin{Highlighting}[]
\FunctionTok{set.seed}\NormalTok{(}\DecValTok{3}\NormalTok{)}
\NormalTok{alphas\_test }\OtherTok{\textless{}{-}} \FunctionTok{c}\NormalTok{(}\FloatTok{2.1}\NormalTok{, }\FloatTok{2.5}\NormalTok{, }\FloatTok{4.5}\NormalTok{)}
\NormalTok{gamma\_tests }\OtherTok{\textless{}{-}} \FunctionTok{lapply}\NormalTok{(alphas\_test, }\ControlFlowTok{function}\NormalTok{(a) }\FunctionTok{sample\_gamma\_rej}\NormalTok{(}\DecValTok{2000}\NormalTok{, a))}
\NormalTok{knitr}\SpecialCharTok{::}\FunctionTok{kable}\NormalTok{(}\FunctionTok{data.frame}\NormalTok{(}
    \AttributeTok{alpha =}\NormalTok{ alphas\_test,}
    \AttributeTok{accept\_rate =} \FunctionTok{vapply}\NormalTok{(gamma\_tests, attr, }\FunctionTok{numeric}\NormalTok{(}\DecValTok{1}\NormalTok{), }\StringTok{"accept.rate"}\NormalTok{),}
    \AttributeTok{ks\_stat =} \FunctionTok{mapply}\NormalTok{(}\ControlFlowTok{function}\NormalTok{(x, a) }\FunctionTok{suppressWarnings}\NormalTok{(}\FunctionTok{ks.test}\NormalTok{(x, }\StringTok{"pgamma"}\NormalTok{, a))}\SpecialCharTok{$}\NormalTok{statistic,}
\NormalTok{                      gamma\_tests, alphas\_test)}
\NormalTok{), }\AttributeTok{digits =} \DecValTok{4}\NormalTok{, }\AttributeTok{col.names =} \FunctionTok{c}\NormalTok{(}\StringTok{"alpha"}\NormalTok{, }\StringTok{"acceptance rate"}\NormalTok{, }\StringTok{"KS statistic vs true Gamma"}\NormalTok{))}
\end{Highlighting}
\end{Shaded}

{

\begin{longtable}[]{@{}rrr@{}}

\caption{\label{tbl-gamma-tests}Acceptance rate falls sharply as alpha
grows: the Cauchy envelope is tuned to match curvature at alpha = 2 and
gets looser as the Gamma density becomes more Gaussian-shaped further
from that point.}

\tabularnewline

\toprule\noalign{}
alpha & acceptance rate & KS statistic vs true Gamma \\
\midrule\noalign{}
\endhead
\bottomrule\noalign{}
\endlastfoot
2.1 & 0.4927 & 0.0185 \\
2.5 & 0.3967 & 0.0145 \\
4.5 & 0.0488 & 0.0223 \\

\end{longtable}

}

\begin{Shaded}
\begin{Highlighting}[]
\FunctionTok{par}\NormalTok{(}\AttributeTok{mfrow =} \FunctionTok{c}\NormalTok{(}\DecValTok{1}\NormalTok{, }\DecValTok{3}\NormalTok{), }\AttributeTok{mar =} \FunctionTok{c}\NormalTok{(}\DecValTok{4}\NormalTok{, }\DecValTok{4}\NormalTok{, }\DecValTok{2}\NormalTok{, }\DecValTok{1}\NormalTok{))}
\ControlFlowTok{for}\NormalTok{ (i }\ControlFlowTok{in} \DecValTok{1}\SpecialCharTok{:}\DecValTok{3}\NormalTok{) \{}
\NormalTok{    a }\OtherTok{\textless{}{-}}\NormalTok{ alphas\_test[i]}
    \FunctionTok{qqplot}\NormalTok{(gamma\_tests[[i]], }\FunctionTok{rgamma}\NormalTok{(}\DecValTok{2000}\NormalTok{, a), }\AttributeTok{xlab =} \StringTok{"rejection sample"}\NormalTok{, }\AttributeTok{ylab =} \StringTok{"rgamma sample"}\NormalTok{,}
           \AttributeTok{main =} \FunctionTok{sprintf}\NormalTok{(}\StringTok{"alpha = \%g (accept \%.1f\%\%)"}\NormalTok{, a, }\DecValTok{100} \SpecialCharTok{*} \FunctionTok{attr}\NormalTok{(gamma\_tests[[i]], }\StringTok{"accept.rate"}\NormalTok{)))}
    \FunctionTok{abline}\NormalTok{(}\DecValTok{0}\NormalTok{, }\DecValTok{1}\NormalTok{, }\AttributeTok{col =} \StringTok{"firebrick"}\NormalTok{, }\AttributeTok{lwd =} \DecValTok{2}\NormalTok{)}
\NormalTok{\}}
\end{Highlighting}
\end{Shaded}

\begin{figure}[H]

\centering{

\pandocbounded{\includegraphics[keepaspectratio]{chapters/unit10_rejection_sampling_files/figure-pdf/fig-gamma-tests-1.pdf}}

}

\caption{\label{fig-gamma-tests}Q-Q plots of rejection-sampled draws
against the true Gamma(alpha,1) quantiles, for three of the test cases:
the envelope produces exact draws regardless of alpha, only the cost per
draw changes.}

\end{figure}%

The Q-Q plots confirm exactness at every \(\alpha\): rejection sampling
never trades accuracy for speed, only speed changes. Efficient
random-number generators for the Gamma distribution (as used inside
\texttt{rgamma} itself) use envelopes that stay tight across a much
wider range of \(\alpha\); see Marsaglia \& Tsang (2000, \emph{ACM
TOMS}) for the algorithm R actually uses.

\section{Adaptive Rejection Sampling for Log-Concave
Targets}\label{adaptive-rejection-sampling-for-log-concave-targets}

Designing an envelope by hand, as above, is the hard part of rejection
sampling. For a broad and important class of targets it can be
automated.

\subsection{Log-concavity and the tangent
hull}\label{log-concavity-and-the-tangent-hull}

A density \(f\) is \textbf{log-concave} if \(\log f(x)\) is concave.
Many posteriors are: the normal, the gamma with shape \(\ge 1\), the
beta with both parameters \(\ge 1\), logistic-regression posteriors
under normal priors, and any product of log-concave factors (so most
conjugate and many non-conjugate exponential- family posteriors
qualify).

A concave function lies below every one of its tangent lines. So for
tangent lines \(\ell_1,\ldots,\ell_k\) to \(\log f\) at points
\(x_1 < \cdots < x_k\), the piecewise-linear \textbf{upper hull} \[
\log g(x) = \min_j \ell_j(x)
\] satisfies \(\log g \ge \log f\) everywhere, hence
\(g = e^{\log g} \ge f\): a valid envelope, built automatically from a
handful of tangent points and their slopes (the log-density's
derivative), with no hand-crafted formula. Because \(g\) is
piecewise-exponential, both its normalizing integral and its inverse CDF
are closed-form on each piece, so sampling from \(g\) is cheap.

\subsection{The adaptive algorithm}\label{the-adaptive-algorithm}

\textbf{Adaptive rejection sampling} (Gilks \& Wild, 1992) refines the
hull as it goes:

\begin{enumerate}
\def\labelenumi{\arabic{enumi}.}
\tightlist
\item
  Start with a few tangent points, at least one on each side of the
  mode.
\item
  Draw \(x \sim g\), \(U \sim \text{Unif}(0,1)\); accept if
  \(U < f(x)/g(x)\).
\item
  Whenever a point is \textbf{rejected}, add it to the tangent set and
  rebuild the hull.
\end{enumerate}

Each rejection adds a tangent exactly where the envelope was loose, so
the acceptance rate rises toward \(1\) as sampling proceeds --- the
envelope adapts itself to the target instead of being fixed in advance.
Only \(\log f\) and its derivative are required (a derivative-free
variant uses secants in place of tangents). The \texttt{ars} package
implements this for any user-supplied log-concave \texttt{f}.

\subsection{Example: Truncated Normal}\label{example-truncated-normal}

The normal density is log-concave, so \texttt{ars} applies directly to a
normal truncated to \([l, u]\) --- useful whenever a parameter has a
hard bound (a variance, a probability, a physically constrained rate).

\begin{Shaded}
\begin{Highlighting}[]
\DocumentationTok{\#\#\# direct rejection sampling for a truncated normal: draw from N(0,1),}
\DocumentationTok{\#\#\# keep values inside [lb, ub]; vectorized batch{-}and{-}refill}
\NormalTok{sample\_tnorm\_drs }\OtherTok{\textless{}{-}} \ControlFlowTok{function}\NormalTok{(n, }\AttributeTok{lb =} \SpecialCharTok{{-}}\ConstantTok{Inf}\NormalTok{, }\AttributeTok{ub =} \ConstantTok{Inf}\NormalTok{) \{}
\NormalTok{    out }\OtherTok{\textless{}{-}} \FunctionTok{numeric}\NormalTok{(}\DecValTok{0}\NormalTok{)}
    \ControlFlowTok{while}\NormalTok{ (}\FunctionTok{length}\NormalTok{(out) }\SpecialCharTok{\textless{}}\NormalTok{ n) \{}
\NormalTok{        x }\OtherTok{\textless{}{-}} \FunctionTok{rnorm}\NormalTok{(n }\SpecialCharTok{{-}} \FunctionTok{length}\NormalTok{(out))}
\NormalTok{        out }\OtherTok{\textless{}{-}} \FunctionTok{c}\NormalTok{(out, x[x }\SpecialCharTok{\textgreater{}=}\NormalTok{ lb }\SpecialCharTok{\&}\NormalTok{ x }\SpecialCharTok{\textless{}=}\NormalTok{ ub])}
\NormalTok{    \}}
\NormalTok{    out}
\NormalTok{\}}

\DocumentationTok{\#\#\# adaptive rejection sampling via the ars package}
\NormalTok{sample\_tnorm\_ars }\OtherTok{\textless{}{-}} \ControlFlowTok{function}\NormalTok{(n, lb, ub) \{}
\NormalTok{    logf }\OtherTok{\textless{}{-}} \ControlFlowTok{function}\NormalTok{(x) }\FunctionTok{dnorm}\NormalTok{(x, }\AttributeTok{log =} \ConstantTok{TRUE}\NormalTok{)}
\NormalTok{    fprima }\OtherTok{\textless{}{-}} \ControlFlowTok{function}\NormalTok{(x) }\SpecialCharTok{{-}}\NormalTok{x}
    \FunctionTok{ars}\NormalTok{(n, }\AttributeTok{f =}\NormalTok{ logf, }\AttributeTok{fprima =}\NormalTok{ fprima, }\AttributeTok{x =} \FunctionTok{c}\NormalTok{(lb, (lb }\SpecialCharTok{+}\NormalTok{ ub) }\SpecialCharTok{/} \DecValTok{2}\NormalTok{, ub),}
        \AttributeTok{lb =} \ConstantTok{TRUE}\NormalTok{, }\AttributeTok{ub =} \ConstantTok{TRUE}\NormalTok{, }\AttributeTok{xlb =}\NormalTok{ lb, }\AttributeTok{xub =}\NormalTok{ ub)}
\NormalTok{\}}
\end{Highlighting}
\end{Shaded}

\begin{Shaded}
\begin{Highlighting}[]
\FunctionTok{set.seed}\NormalTok{(}\DecValTok{4}\NormalTok{)}
\NormalTok{n }\OtherTok{\textless{}{-}} \DecValTok{1000}
\NormalTok{windows }\OtherTok{\textless{}{-}} \FunctionTok{list}\NormalTok{(}\FunctionTok{c}\NormalTok{(}\SpecialCharTok{{-}}\DecValTok{5}\NormalTok{, }\SpecialCharTok{{-}}\DecValTok{4}\NormalTok{), }\FunctionTok{c}\NormalTok{(}\SpecialCharTok{{-}}\DecValTok{50}\NormalTok{, }\SpecialCharTok{{-}}\DecValTok{40}\NormalTok{), }\FunctionTok{c}\NormalTok{(}\DecValTok{100}\NormalTok{, }\DecValTok{110}\NormalTok{))}
\NormalTok{timing }\OtherTok{\textless{}{-}} \FunctionTok{lapply}\NormalTok{(windows, }\ControlFlowTok{function}\NormalTok{(w) \{}
\NormalTok{    t\_ars }\OtherTok{\textless{}{-}} \FunctionTok{system.time}\NormalTok{(}\FunctionTok{sample\_tnorm\_ars}\NormalTok{(n, w[}\DecValTok{1}\NormalTok{], w[}\DecValTok{2}\NormalTok{]))[}\StringTok{"elapsed"}\NormalTok{]}
\NormalTok{    t\_naive }\OtherTok{\textless{}{-}} \ControlFlowTok{if}\NormalTok{ (w[}\DecValTok{1}\NormalTok{] }\SpecialCharTok{\textgreater{}} \SpecialCharTok{{-}}\DecValTok{10} \SpecialCharTok{\&\&}\NormalTok{ w[}\DecValTok{2}\NormalTok{] }\SpecialCharTok{\textless{}} \DecValTok{10}\NormalTok{) }\FunctionTok{system.time}\NormalTok{(}\FunctionTok{sample\_tnorm\_drs}\NormalTok{(n, w[}\DecValTok{1}\NormalTok{], w[}\DecValTok{2}\NormalTok{]))[}\StringTok{"elapsed"}\NormalTok{] }\ControlFlowTok{else} \ConstantTok{NA}
    \FunctionTok{c}\NormalTok{(}\AttributeTok{lb =}\NormalTok{ w[}\DecValTok{1}\NormalTok{], }\AttributeTok{ub =}\NormalTok{ w[}\DecValTok{2}\NormalTok{], }\AttributeTok{ars\_seconds =}\NormalTok{ t\_ars, }\AttributeTok{naive\_seconds =}\NormalTok{ t\_naive)}
\NormalTok{\})}
\NormalTok{knitr}\SpecialCharTok{::}\FunctionTok{kable}\NormalTok{(}\FunctionTok{do.call}\NormalTok{(rbind, timing), }\AttributeTok{digits =} \DecValTok{4}\NormalTok{,}
             \AttributeTok{col.names =} \FunctionTok{c}\NormalTok{(}\StringTok{"lower bound"}\NormalTok{, }\StringTok{"upper bound"}\NormalTok{, }\StringTok{"ars (s)"}\NormalTok{, }\StringTok{"naive rejection (s)"}\NormalTok{))}
\end{Highlighting}
\end{Shaded}

{

\begin{longtable}[]{@{}rrrr@{}}

\caption{\label{tbl-ars-timing}Adaptive rejection sampling (ars) vs
naive rejection sampling from an untruncated N(0,1): naive sampling
degrades catastrophically as the truncation window moves into the tail,
while ars does not, because it never wastes draws outside {[}lb, ub{]}
in the first place.}

\tabularnewline

\toprule\noalign{}
lower bound & upper bound & ars (s) & naive rejection (s) \\
\midrule\noalign{}
\endhead
\bottomrule\noalign{}
\endlastfoot
-5 & -4 & 0.003 & 1.677 \\
-50 & -40 & 0.006 & NA \\
100 & 110 & 0.003 & NA \\

\end{longtable}

}

For \([-50,-40]\) or \([100,110]\) the naive sampler's acceptance
probability is astronomically small
(\(P(-50<Z<-40) \approx 10^{-221}\)), so it is left as \texttt{NA}
rather than run to timeout; \texttt{ars} samples these windows in a
fraction of a second because its envelope is built from the truncated
density itself, never from the untruncated normal.

\begin{Shaded}
\begin{Highlighting}[]
\FunctionTok{par}\NormalTok{(}\AttributeTok{mfrow =} \FunctionTok{c}\NormalTok{(}\DecValTok{1}\NormalTok{, }\DecValTok{3}\NormalTok{), }\AttributeTok{mar =} \FunctionTok{c}\NormalTok{(}\DecValTok{4}\NormalTok{, }\DecValTok{4}\NormalTok{, }\DecValTok{2}\NormalTok{, }\DecValTok{1}\NormalTok{))}
\FunctionTok{set.seed}\NormalTok{(}\DecValTok{5}\NormalTok{)}
\ControlFlowTok{for}\NormalTok{ (w }\ControlFlowTok{in}\NormalTok{ windows) \{}
\NormalTok{    x }\OtherTok{\textless{}{-}} \FunctionTok{sample\_tnorm\_ars}\NormalTok{(n, w[}\DecValTok{1}\NormalTok{], w[}\DecValTok{2}\NormalTok{])}
    \FunctionTok{hist}\NormalTok{(x, }\AttributeTok{breaks =} \DecValTok{25}\NormalTok{, }\AttributeTok{col =} \StringTok{"steelblue"}\NormalTok{, }\AttributeTok{border =} \StringTok{"white"}\NormalTok{,}
         \AttributeTok{main =} \FunctionTok{sprintf}\NormalTok{(}\StringTok{"[\%g, \%g]"}\NormalTok{, w[}\DecValTok{1}\NormalTok{], w[}\DecValTok{2}\NormalTok{]), }\AttributeTok{xlab =} \StringTok{"x"}\NormalTok{)}
\NormalTok{\}}
\end{Highlighting}
\end{Shaded}

\begin{figure}[H]

\centering{

\pandocbounded{\includegraphics[keepaspectratio]{chapters/unit10_rejection_sampling_files/figure-pdf/fig-ars-1.pdf}}

}

\caption{\label{fig-ars}Histograms of adaptive-rejection-sampled draws
in three increasingly extreme truncation windows: ars has no trouble
with any of them.}

\end{figure}%

\section{Summary}\label{summary-4}

\begin{itemize}
\tightlist
\item
  Rejection sampling turns draws from an envelope \(g \ge f\) into
  \emph{exact} iid draws from \(f\), using only the ratio \(f/g\); no
  normalizing constants are needed. The geometric view --- uniform
  points under \(g\), keep those under \(f\) --- explains why this works
  and why efficiency is exactly the ratio of the two areas.
\item
  A good envelope is tight (high acceptance) and easy to sample; finding
  one by hand, as in the Cauchy envelope for the Gamma distribution
  above, requires matching the target's tails and curvature and can be
  delicate.
\item
  Adaptive rejection sampling automates envelope construction for any
  log-concave target via tangent lines to \(\log f\), tightening the
  envelope on every rejection.
\item
  Limitations: an envelope with a bounded ratio \(g/f\) may not exist
  when \(f\) has heavier tails than any convenient \(g\), and in high
  dimensions the acceptance rate of a product envelope typically decays
  exponentially with dimension, for the same curse-of-dimensionality
  reason quadrature grids fail (Unit 8). These limitations motivate
  \textbf{importance sampling}, which reweights instead of rejecting,
  and \textbf{Markov chain Monte Carlo}, which builds a chain that
  adapts to the target as it runs --- the subjects of the next two
  chapters.
\end{itemize}

\bookmarksetup{startatroot}

\chapter{Importance Sampling}\label{importance-sampling}

\begin{Shaded}
\begin{Highlighting}[]
\FunctionTok{library}\NormalTok{(ars)}
\end{Highlighting}
\end{Shaded}

Rejection sampling turns draws from an envelope \(g \ge f\) into exact
draws from \(f\), but every rejected proposal is wasted, and in the
tails of a poor envelope almost every proposal is rejected.
\textbf{Importance sampling} keeps every proposal instead of discarding
most of them: draw from an easier function \(g\) and attach a
\emph{weight} to each draw that corrects for sampling from the wrong
distribution. The price is that the resulting sample is no longer an
exact iid sample from \(f\) --- it is a \emph{weighted} sample --- but
nothing is ever thrown away, and no envelope condition \(g \ge f\) is
required.

\section{The Reweighting Idea}\label{the-reweighting-idea}

\subsection{The goal}\label{the-goal}

As in the last two chapters, \(f(\theta)\) is an \textbf{unnormalized}
function of interest --- typically \(f(\theta) = L(\theta)\pi(\theta)\)
for a likelihood \(L\) and prior \(\pi\) --- and we want two kinds of
integral: \[
A = E_f\big[a(\theta)\big] = \frac{\int a(\theta)\,f(\theta)\,d\theta}{\int f(\theta)\,d\theta},
\qquad\qquad
C_f = \int f(\theta)\,d\theta,
\] a posterior expectation and a normalizing constant (the marginal
likelihood when \(f = L\pi\)). Plain Monte Carlo needs draws from
\(f/C_f\), which we generally cannot produce.

Suppose instead we can draw from another function \(g(\theta)\) (also
possibly unnormalized, with \(C_g = \int g(\theta)\,d\theta\)). Wherever
\(f\) puts mass, \(g\) must too: \[
\{\theta : g(\theta) > 0\} \ \supseteq\ \{\theta : f(\theta) > 0\}.
\] Unlike rejection sampling, \(g\) need \textbf{not} dominate \(f\)
everywhere --- there is no requirement \(g \ge f\). This is what makes
importance sampling applicable in situations where no usable envelope
exists.

\subsection{Two estimators}\label{two-estimators}

Define the \textbf{importance weight}
\(w(\theta) = f(\theta)/g(\theta)\). For any \(a\), \[
E_f\big[a(\theta)\big]
= \int a(\theta)\,\frac{f(\theta)}{C_f}\,d\theta
= \frac{C_g}{C_f}\, E_g\big[a(\theta)\,w(\theta)\big],
\] and taking \(a \equiv 1\) gives \(C_f/C_g = E_g[w(\theta)]\).
Combining the two, \[
E_f\big[a(\theta)\big] = \frac{E_g\big[a(\theta)\,w(\theta)\big]}{E_g\big[w(\theta)\big]}.
\] Draw \(\theta_1,\ldots,\theta_S \overset{iid}{\sim} g/C_g\) and
compute \(w_s = w(\theta_s)\). Two estimators follow directly:

\begin{itemize}
\tightlist
\item
  \textbf{Ratio of normalizing constants}:
  \(\widehat{C_f/C_g} = \frac1S\sum_s w_s\), unbiased for \(C_f/C_g\);
  if \(g\) is already a normalized density (\(C_g=1\)) this estimates
  \(C_f\) itself.
\item
  \textbf{Self-normalized expectation}:
  \(\hat A = \sum_{s=1}^S a(\theta_s)\,\tilde w_s\), with normalized
  weights \(\tilde w_s = w_s / \sum_j w_j\) summing to \(1\). This turns
  the proposal sample into a \emph{weighted sample from \(f\)}: any
  unknown constants in \(f\) and \(g\) cancel in the ratio, so
  \(\hat A\) needs only \(f\) and \(g\) up to a proportionality
  constant. It is a ratio of two averages and so is biased with bias
  \(O(1/S)\), but consistent.
\end{itemize}

\subsection{Effective sample size}\label{effective-sample-size}

\(\hat A\) is a weighted average of \(a(\theta_s)\): draws where
\(f \gg g\) count more, draws where \(f \ll g\) count less. How many
\emph{equally weighted} draws is the weighted sample worth? \[
S_{\text{eff}} = \frac{1}{\displaystyle\sum_{s=1}^S \tilde w_s^2}.
\] If the weights are all equal, \(\tilde w_s = 1/S\) and
\(S_{\text{eff}} =
1/(S \cdot 1/S^2) = S\): the proposal is shaped just like the target,
and every draw counts fully. If one weight dominates all the others,
\(\tilde w_s
\approx 1\) for that draw and \(S_{\text{eff}} \approx 1\): only one
draw is doing any work. A variance heuristic makes this precise:
\(\text{Var}_g(\hat A) \approx \text{Var}_f(a)/S_{\text{eff}}\), the
same form as for an iid sample of size \(S_{\text{eff}}\) from \(f\)
itself.

The catch, central to this chapter, is that \(S_{\text{eff}}\) is
computed \textbf{only from the weights actually observed}. A proposal
that never visits the part of \(\theta\)-space where \(f\) and \(g\)
disagree most will report a reassuring \(S_{\text{eff}}\) while being
badly wrong --- this is exactly the subject of the ``bad stories''
below.

\begin{Shaded}
\begin{Highlighting}[]
\DocumentationTok{\#\#\# effective sample size from raw (possibly unnormalized) weights}
\NormalTok{ess }\OtherTok{\textless{}{-}} \ControlFlowTok{function}\NormalTok{(w) }\FunctionTok{sum}\NormalTok{(w)}\SpecialCharTok{\^{}}\DecValTok{2} \SpecialCharTok{/} \FunctionTok{sum}\NormalTok{(w}\SpecialCharTok{\^{}}\DecValTok{2}\NormalTok{)}
\DocumentationTok{\#\#\# numerically stable log(sum(exp(lx)))}
\NormalTok{log\_sum\_exp }\OtherTok{\textless{}{-}} \ControlFlowTok{function}\NormalTok{(lx) \{ m }\OtherTok{\textless{}{-}} \FunctionTok{max}\NormalTok{(lx); m }\SpecialCharTok{+} \FunctionTok{log}\NormalTok{(}\FunctionTok{sum}\NormalTok{(}\FunctionTok{exp}\NormalTok{(lx }\SpecialCharTok{{-}}\NormalTok{ m))) \}}
\DocumentationTok{\#\#\# effective sample size computed from LOG weights (avoids under/overflow)}
\NormalTok{log\_ess }\OtherTok{\textless{}{-}} \ControlFlowTok{function}\NormalTok{(lw) }\FunctionTok{exp}\NormalTok{(}\DecValTok{2} \SpecialCharTok{*} \FunctionTok{log\_sum\_exp}\NormalTok{(lw) }\SpecialCharTok{{-}} \FunctionTok{log\_sum\_exp}\NormalTok{(}\DecValTok{2} \SpecialCharTok{*}\NormalTok{ lw))}
\end{Highlighting}
\end{Shaded}

\section{A Good Story: Estimating a Normal Tail
Probability}\label{a-good-story-estimating-a-normal-tail-probability}

\subsection{From plain Monte Carlo to importance
sampling}\label{from-plain-monte-carlo-to-importance-sampling}

For \(X \sim N(0,1)\), take \(f(x) = \varphi(x)\,I(a<x<b)\), the
standard normal density truncated to \([a,b]\), so that
\(C_f = P(a<X<b)\). Plain Monte Carlo draws \(X_i \sim N(0,1)\) and
averages the indicator \(I(a<X_i<b)\); this wastes every draw that lands
outside \([a,b]\); for a proposal \(g\) supported on \([a,b]\) itself,
every draw contributes.

The simplest such \(g\) is uniform on \([a,b]\):
\(g(x) = I(a\le x\le b)/(b-a)\), so \(w(x) = (b-a)\varphi(x)\) and
\(\widehat{C_f} = (b-a)\cdot\frac1S\sum_s
\varphi(x_s)\) --- the \emph{Monte Carlo} version of the midpoint rule
from the quadrature chapter: a random grid instead of an equally spaced
one, evaluating the same weighted sum.

\begin{Shaded}
\begin{Highlighting}[]
\NormalTok{est\_normprob\_mc   }\OtherTok{\textless{}{-}} \ControlFlowTok{function}\NormalTok{(A, S) \{ X }\OtherTok{\textless{}{-}} \FunctionTok{rnorm}\NormalTok{(S); }\FunctionTok{mean}\NormalTok{((X }\SpecialCharTok{\textgreater{}}\NormalTok{ A[}\DecValTok{1}\NormalTok{]) }\SpecialCharTok{\&}\NormalTok{ (X }\SpecialCharTok{\textless{}}\NormalTok{ A[}\DecValTok{2}\NormalTok{])) \}}
\NormalTok{est\_normprob\_imps }\OtherTok{\textless{}{-}} \ControlFlowTok{function}\NormalTok{(A, S) \{ X }\OtherTok{\textless{}{-}} \FunctionTok{runif}\NormalTok{(S, A[}\DecValTok{1}\NormalTok{], A[}\DecValTok{2}\NormalTok{]); }\FunctionTok{mean}\NormalTok{(}\FunctionTok{dnorm}\NormalTok{(X)) }\SpecialCharTok{*}\NormalTok{ (A[}\DecValTok{2}\NormalTok{] }\SpecialCharTok{{-}}\NormalTok{ A[}\DecValTok{1}\NormalTok{]) \}}

\FunctionTok{set.seed}\NormalTok{(}\DecValTok{1}\NormalTok{)}
\NormalTok{cases }\OtherTok{\textless{}{-}} \FunctionTok{list}\NormalTok{(}\FunctionTok{c}\NormalTok{(}\SpecialCharTok{{-}}\DecValTok{2}\NormalTok{, }\DecValTok{2}\NormalTok{), }\FunctionTok{c}\NormalTok{(}\DecValTok{5}\NormalTok{, }\DecValTok{6}\NormalTok{))}
\NormalTok{mse\_tbl }\OtherTok{\textless{}{-}} \FunctionTok{do.call}\NormalTok{(rbind, }\FunctionTok{lapply}\NormalTok{(cases, }\ControlFlowTok{function}\NormalTok{(A) \{}
\NormalTok{    tp }\OtherTok{\textless{}{-}} \FunctionTok{pnorm}\NormalTok{(A[}\DecValTok{2}\NormalTok{]) }\SpecialCharTok{{-}} \FunctionTok{pnorm}\NormalTok{(A[}\DecValTok{1}\NormalTok{])}
\NormalTok{    mc  }\OtherTok{\textless{}{-}} \FunctionTok{replicate}\NormalTok{(}\DecValTok{1000}\NormalTok{, }\FunctionTok{est\_normprob\_mc}\NormalTok{(A, }\DecValTok{100}\NormalTok{))}
\NormalTok{    imp }\OtherTok{\textless{}{-}} \FunctionTok{replicate}\NormalTok{(}\DecValTok{1000}\NormalTok{, }\FunctionTok{est\_normprob\_imps}\NormalTok{(A, }\DecValTok{100}\NormalTok{))}
    \FunctionTok{data.frame}\NormalTok{(}\AttributeTok{interval =} \FunctionTok{sprintf}\NormalTok{(}\StringTok{"(\%g, \%g)"}\NormalTok{, A[}\DecValTok{1}\NormalTok{], A[}\DecValTok{2}\NormalTok{]),}
               \AttributeTok{true\_prob =} \FunctionTok{formatC}\NormalTok{(tp, }\AttributeTok{format =} \StringTok{"e"}\NormalTok{, }\AttributeTok{digits =} \DecValTok{3}\NormalTok{),}
               \AttributeTok{mse\_mc =} \FunctionTok{formatC}\NormalTok{(}\FunctionTok{mean}\NormalTok{((mc }\SpecialCharTok{{-}}\NormalTok{ tp)}\SpecialCharTok{\^{}}\DecValTok{2}\NormalTok{), }\AttributeTok{format =} \StringTok{"e"}\NormalTok{, }\AttributeTok{digits =} \DecValTok{3}\NormalTok{),}
               \AttributeTok{mse\_importance =} \FunctionTok{formatC}\NormalTok{(}\FunctionTok{mean}\NormalTok{((imp }\SpecialCharTok{{-}}\NormalTok{ tp)}\SpecialCharTok{\^{}}\DecValTok{2}\NormalTok{), }\AttributeTok{format =} \StringTok{"e"}\NormalTok{, }\AttributeTok{digits =} \DecValTok{3}\NormalTok{))}
\NormalTok{\}))}
\NormalTok{knitr}\SpecialCharTok{::}\FunctionTok{kable}\NormalTok{(mse\_tbl, }\AttributeTok{align =} \FunctionTok{c}\NormalTok{(}\StringTok{"l"}\NormalTok{, }\StringTok{"r"}\NormalTok{, }\StringTok{"r"}\NormalTok{, }\StringTok{"r"}\NormalTok{))}
\end{Highlighting}
\end{Shaded}

{

\begin{longtable}[]{@{}lrrr@{}}

\caption{\label{tbl-tail-good}Mean squared error of the
truncated-probability estimate over 1000 replications, for a moderate
and a rare event, S = 100 draws per replication.}

\tabularnewline

\toprule\noalign{}
interval & true\_prob & mse\_mc & mse\_importance \\
\midrule\noalign{}
\endhead
\bottomrule\noalign{}
\endlastfoot
(-2, 2) & 9.545e-01 & 4.421e-04 & 2.002e-03 \\
(5, 6) & 2.857e-07 & 8.160e-14 & 1.295e-15 \\

\end{longtable}

}

For the moderate interval \((-2,2)\) both methods work; for the rarer
event \((5,6)\), where plain Monte Carlo almost never lands an
observation inside the interval, the importance-sampling error is many
orders of magnitude smaller --- the uniform proposal \emph{guarantees}
every draw falls in \([a,b]\).

\subsection{An underflowing tail
probability}\label{an-underflowing-tail-probability}

Push the interval further out, to \(A = (50,100)\). Now even
\(\varphi(50)\) itself underflows to exactly \(0\) in double precision
(its log is about \(-1251\), far below the smallest representable log),
so every calculation must be carried out on the log scale.

\begin{Shaded}
\begin{Highlighting}[]
\NormalTok{a }\OtherTok{\textless{}{-}} \DecValTok{50}\NormalTok{; b }\OtherTok{\textless{}{-}} \DecValTok{100}
\NormalTok{log\_minus\_exp }\OtherTok{\textless{}{-}} \ControlFlowTok{function}\NormalTok{(la, lb) la }\SpecialCharTok{+} \FunctionTok{log1p}\NormalTok{(}\SpecialCharTok{{-}}\FunctionTok{exp}\NormalTok{(lb }\SpecialCharTok{{-}}\NormalTok{ la))}
\NormalTok{log\_exact }\OtherTok{\textless{}{-}} \FunctionTok{log\_minus\_exp}\NormalTok{(}\FunctionTok{pnorm}\NormalTok{(}\SpecialCharTok{{-}}\NormalTok{a, }\AttributeTok{log =} \ConstantTok{TRUE}\NormalTok{), }\FunctionTok{pnorm}\NormalTok{(}\SpecialCharTok{{-}}\NormalTok{b, }\AttributeTok{log =} \ConstantTok{TRUE}\NormalTok{))}
\NormalTok{log\_exact}
\end{Highlighting}
\end{Shaded}

\begin{verbatim}
[1] -1254.831
\end{verbatim}

\subsection{Matching the proposal to the tail: a better
envelope}\label{matching-the-proposal-to-the-tail-a-better-envelope}

A uniform proposal on \((50,100)\) still works in principle, but
\(\varphi\) decays so fast across the interval that the weights vary
enormously: nearly all of the probability mass under \(\varphi\)
restricted to \((50,100)\) sits within a sliver near \(x=50\), so a
uniform draw almost never lands where the weight is large. A
\textbf{shifted exponential} proposal with rate \(\lambda \approx a\)
matches that decay and keeps the weights nearly constant: \[
g(x) = \frac{\lambda e^{-\lambda(x-a)}}{1-e^{-\lambda(b-a)}}\,I(a\le x\le b).
\]

\begin{Shaded}
\begin{Highlighting}[]
\NormalTok{S }\OtherTok{\textless{}{-}} \FloatTok{1e4}\NormalTok{; lam }\OtherTok{\textless{}{-}}\NormalTok{ a}
\NormalTok{x\_u }\OtherTok{\textless{}{-}} \FunctionTok{runif}\NormalTok{(S, a, b);  lw\_u }\OtherTok{\textless{}{-}} \FunctionTok{log}\NormalTok{(b }\SpecialCharTok{{-}}\NormalTok{ a) }\SpecialCharTok{+} \FunctionTok{dnorm}\NormalTok{(x\_u, }\AttributeTok{log =} \ConstantTok{TRUE}\NormalTok{)}
\NormalTok{u\_e }\OtherTok{\textless{}{-}} \FunctionTok{runif}\NormalTok{(S); x\_e }\OtherTok{\textless{}{-}}\NormalTok{ a }\SpecialCharTok{{-}} \FunctionTok{log}\NormalTok{(}\DecValTok{1} \SpecialCharTok{{-}}\NormalTok{ u\_e }\SpecialCharTok{*}\NormalTok{ (}\DecValTok{1} \SpecialCharTok{{-}} \FunctionTok{exp}\NormalTok{(}\SpecialCharTok{{-}}\NormalTok{lam }\SpecialCharTok{*}\NormalTok{ (b }\SpecialCharTok{{-}}\NormalTok{ a)))) }\SpecialCharTok{/}\NormalTok{ lam        }\CommentTok{\# inverse cdf}
\NormalTok{lg\_e }\OtherTok{\textless{}{-}} \FunctionTok{log}\NormalTok{(lam) }\SpecialCharTok{{-}}\NormalTok{ lam }\SpecialCharTok{*}\NormalTok{ (x\_e }\SpecialCharTok{{-}}\NormalTok{ a) }\SpecialCharTok{{-}} \FunctionTok{log1p}\NormalTok{(}\SpecialCharTok{{-}}\FunctionTok{exp}\NormalTok{(}\SpecialCharTok{{-}}\NormalTok{lam }\SpecialCharTok{*}\NormalTok{ (b }\SpecialCharTok{{-}}\NormalTok{ a)))}
\NormalTok{lw\_e }\OtherTok{\textless{}{-}} \FunctionTok{dnorm}\NormalTok{(x\_e, }\AttributeTok{log =} \ConstantTok{TRUE}\NormalTok{) }\SpecialCharTok{{-}}\NormalTok{ lg\_e}

\NormalTok{tail\_tbl }\OtherTok{\textless{}{-}} \FunctionTok{data.frame}\NormalTok{(}
    \AttributeTok{proposal   =} \FunctionTok{c}\NormalTok{(}\StringTok{"uniform on (a, b)"}\NormalTok{, }\StringTok{"shifted exponential"}\NormalTok{),}
    \AttributeTok{log\_C\_f    =} \FunctionTok{c}\NormalTok{(}\FunctionTok{log\_sum\_exp}\NormalTok{(lw\_u) }\SpecialCharTok{{-}} \FunctionTok{log}\NormalTok{(S), }\FunctionTok{log\_sum\_exp}\NormalTok{(lw\_e) }\SpecialCharTok{{-}} \FunctionTok{log}\NormalTok{(S)),}
    \AttributeTok{S\_eff      =} \FunctionTok{c}\NormalTok{(}\FunctionTok{log\_ess}\NormalTok{(lw\_u), }\FunctionTok{log\_ess}\NormalTok{(lw\_e))}
\NormalTok{)}
\NormalTok{knitr}\SpecialCharTok{::}\FunctionTok{kable}\NormalTok{(tail\_tbl, }\AttributeTok{digits =} \FunctionTok{c}\NormalTok{(}\DecValTok{0}\NormalTok{, }\DecValTok{4}\NormalTok{, }\DecValTok{1}\NormalTok{))}
\end{Highlighting}
\end{Shaded}

{

\begin{longtable}[]{@{}lrr@{}}

\caption{\label{tbl-tail-proposals}Two proposals for the same rare
event, S = 10,000 draws: log-scale weights avoid the underflow that
would occur if densities were multiplied directly.}

\tabularnewline

\toprule\noalign{}
proposal & log\_C\_f & S\_eff \\
\midrule\noalign{}
\endhead
\bottomrule\noalign{}
\endlastfoot
uniform on (a, b) & -1254.489 & 7.8 \\
shifted exponential & -1254.831 & 10000.0 \\

\end{longtable}

}

Both proposals give a log-probability close to -1254.831, but their
effective sample sizes are worlds apart: the exponential proposal keeps
essentially all \ensuremath{10^{4}} draws informative, while the uniform
proposal's effective sample size collapses because almost all of its
weight concentrates on the handful of draws nearest \(a\). Shaping \(g\)
like \(f\) is not just elegant --- it is what makes the computation
efficient.

\section{Two Bad Stories: When Reweighting Fails
Silently}\label{two-bad-stories-when-reweighting-fails-silently}

\(S_{\text{eff}}\) diagnoses only \emph{one} kind of failure: weights
that vary wildly among the draws actually observed. It says nothing
about mass that \(g\) never visits at all, and a badly wrong estimate
can carry a perfectly healthy \(S_{\text{eff}}\).

\subsection{A proposal that misses a
mode}\label{a-proposal-that-misses-a-mode}

Let
\(f(\theta) = \tfrac12\varphi(\theta+5) + \tfrac12\varphi(\theta-5)\) (a
50/50 mixture, so \(C_f=1\) and \(E_f(\theta)=0\)) and
\(g(\theta) = \varphi(\theta+5)\). Both are positive everywhere, so the
support requirement holds. But \(g\) places essentially all of its mass
near \(\theta=-5\) and, for \(\theta\) near \(-5\), \[
w(\theta) = \tfrac12 + \tfrac12\,\frac{\varphi(\theta-5)}{\varphi(\theta+5)}
= \tfrac12 + \tfrac12\, e^{10\theta} \ \approx\ \tfrac12,
\] so every weight the sampler actually sees looks almost identical.

\begin{Shaded}
\begin{Highlighting}[]
\FunctionTok{set.seed}\NormalTok{(}\DecValTok{2}\NormalTok{)}
\NormalTok{S }\OtherTok{\textless{}{-}} \DecValTok{5000}
\NormalTok{f\_mix }\OtherTok{\textless{}{-}} \ControlFlowTok{function}\NormalTok{(t) }\FloatTok{0.5} \SpecialCharTok{*} \FunctionTok{dnorm}\NormalTok{(t }\SpecialCharTok{+} \DecValTok{5}\NormalTok{) }\SpecialCharTok{+} \FloatTok{0.5} \SpecialCharTok{*} \FunctionTok{dnorm}\NormalTok{(t }\SpecialCharTok{{-}} \DecValTok{5}\NormalTok{)}
\NormalTok{g\_mix }\OtherTok{\textless{}{-}} \ControlFlowTok{function}\NormalTok{(t) }\FunctionTok{dnorm}\NormalTok{(t }\SpecialCharTok{+} \DecValTok{5}\NormalTok{)}
\NormalTok{th }\OtherTok{\textless{}{-}} \FunctionTok{rnorm}\NormalTok{(S, }\SpecialCharTok{{-}}\DecValTok{5}\NormalTok{)}
\NormalTok{w  }\OtherTok{\textless{}{-}} \FunctionTok{f\_mix}\NormalTok{(th) }\SpecialCharTok{/} \FunctionTok{g\_mix}\NormalTok{(th)}
\NormalTok{wt }\OtherTok{\textless{}{-}}\NormalTok{ w }\SpecialCharTok{/} \FunctionTok{sum}\NormalTok{(w)}

\FunctionTok{par}\NormalTok{(}\AttributeTok{mfrow =} \FunctionTok{c}\NormalTok{(}\DecValTok{1}\NormalTok{, }\DecValTok{2}\NormalTok{), }\AttributeTok{mar =} \FunctionTok{c}\NormalTok{(}\DecValTok{4}\NormalTok{, }\DecValTok{4}\NormalTok{, }\DecValTok{2}\NormalTok{, }\DecValTok{1}\NormalTok{))}
\FunctionTok{curve}\NormalTok{(f\_mix, }\SpecialCharTok{{-}}\DecValTok{9}\NormalTok{, }\DecValTok{9}\NormalTok{, }\AttributeTok{lwd =} \DecValTok{3}\NormalTok{, }\AttributeTok{col =} \StringTok{"firebrick"}\NormalTok{, }\AttributeTok{xlab =} \FunctionTok{expression}\NormalTok{(theta), }\AttributeTok{ylab =} \StringTok{""}\NormalTok{,}
      \AttributeTok{main =} \StringTok{"target f and proposal g"}\NormalTok{, }\AttributeTok{ylim =} \FunctionTok{c}\NormalTok{(}\DecValTok{0}\NormalTok{, }\FloatTok{0.42}\NormalTok{))}
\FunctionTok{curve}\NormalTok{(g\_mix, }\AttributeTok{add =} \ConstantTok{TRUE}\NormalTok{, }\AttributeTok{lwd =} \DecValTok{3}\NormalTok{, }\AttributeTok{col =} \StringTok{"steelblue"}\NormalTok{)}
\FunctionTok{rug}\NormalTok{(th[}\DecValTok{1}\SpecialCharTok{:}\DecValTok{300}\NormalTok{], }\AttributeTok{col =} \FunctionTok{adjustcolor}\NormalTok{(}\StringTok{"steelblue"}\NormalTok{, }\FloatTok{0.5}\NormalTok{))}
\FunctionTok{legend}\NormalTok{(}\StringTok{"topright"}\NormalTok{, }\AttributeTok{bty =} \StringTok{"n"}\NormalTok{, }\AttributeTok{lwd =} \DecValTok{3}\NormalTok{, }\AttributeTok{col =} \FunctionTok{c}\NormalTok{(}\StringTok{"firebrick"}\NormalTok{, }\StringTok{"steelblue"}\NormalTok{), }\AttributeTok{legend =} \FunctionTok{c}\NormalTok{(}\StringTok{"f"}\NormalTok{, }\StringTok{"g"}\NormalTok{))}
\FunctionTok{hist}\NormalTok{(wt }\SpecialCharTok{*}\NormalTok{ S, }\AttributeTok{breaks =} \DecValTok{40}\NormalTok{, }\AttributeTok{col =} \StringTok{"gray85"}\NormalTok{, }\AttributeTok{border =} \StringTok{"white"}\NormalTok{, }\AttributeTok{xlab =} \FunctionTok{expression}\NormalTok{(S }\SpecialCharTok{*} \FunctionTok{tilde}\NormalTok{(w)[s]),}
     \AttributeTok{main =} \FunctionTok{sprintf}\NormalTok{(}\StringTok{"weights: S\_eff = \%.0f of \%d"}\NormalTok{, }\FunctionTok{ess}\NormalTok{(w), S))}
\end{Highlighting}
\end{Shaded}

\begin{figure}[H]

\centering{

\pandocbounded{\includegraphics[keepaspectratio]{chapters/unit11_importance_sampling_files/figure-pdf/fig-missed-mode-1.pdf}}

}

\caption{\label{fig-missed-mode}Left: target f (mixture of two normals)
and proposal g (only the left component). Right: normalized weights look
almost flat --- S\_eff is close to S --- even though the estimate is
badly wrong.}

\end{figure}%

\begin{Shaded}
\begin{Highlighting}[]
\NormalTok{missed\_mode\_tbl }\OtherTok{\textless{}{-}} \FunctionTok{data.frame}\NormalTok{(}
    \AttributeTok{quantity =} \FunctionTok{c}\NormalTok{(}\StringTok{"C\_f / C\_g (true 1)"}\NormalTok{, }\StringTok{"E\_f[theta] (true 0)"}\NormalTok{, }\StringTok{"S\_eff (of S = 5000)"}\NormalTok{),}
    \AttributeTok{estimate =} \FunctionTok{c}\NormalTok{(}\FunctionTok{mean}\NormalTok{(w), }\FunctionTok{sum}\NormalTok{(th }\SpecialCharTok{*}\NormalTok{ wt), }\FunctionTok{ess}\NormalTok{(w))}
\NormalTok{)}
\NormalTok{knitr}\SpecialCharTok{::}\FunctionTok{kable}\NormalTok{(missed\_mode\_tbl, }\AttributeTok{digits =} \DecValTok{3}\NormalTok{)}
\end{Highlighting}
\end{Shaded}

{

\begin{longtable}[]{@{}lr@{}}

\caption{\label{tbl-missed-mode}Missed-mode failure: the effective
sample size looks fine, but both quantities of interest are far from the
truth.}

\tabularnewline

\toprule\noalign{}
quantity & estimate \\
\midrule\noalign{}
\endhead
\bottomrule\noalign{}
\endlastfoot
C\_f / C\_g (true 1) & 0.50 \\
E\_f{[}theta{]} (true 0) & -4.96 \\
S\_eff (of S = 5000) & 5000.00 \\

\end{longtable}

}

The estimate of \(C_f/C_g\) comes out near 0.5 instead of \(1\), and
\(E_f[\theta]\) near -4.96 instead of \(0\) --- both wrong by a wide
margin --- while \(S_{\text{eff}}\) reports essentially the full sample
size. The draws that would reveal the problem, near \(\theta=+5\), have
probability under \(g\) of order \(\Phi(-10) \approx 10^{-23}\): they
simply never appear, at any practical \(S\).

\subsection{A proposal narrower than the target's
support}\label{a-proposal-narrower-than-the-targets-support}

The same failure in its starkest form:
\(f(\theta) = I(0\le\theta\le2)\), \(g(\theta) = I(0\le\theta\le1)\).
The true ratio is \(C_f/C_g = 2\), but every draw from \(g\) lands in
\((0,1)\), where \(w(\theta)=f(\theta)/g(\theta)=1\) identically.

\begin{Shaded}
\begin{Highlighting}[]
\FunctionTok{set.seed}\NormalTok{(}\DecValTok{3}\NormalTok{)}
\NormalTok{f\_indicator }\OtherTok{\textless{}{-}} \ControlFlowTok{function}\NormalTok{(x) }\FunctionTok{as.numeric}\NormalTok{(x }\SpecialCharTok{\textgreater{}=} \DecValTok{0} \SpecialCharTok{\&}\NormalTok{ x }\SpecialCharTok{\textless{}=} \DecValTok{2}\NormalTok{)   }\CommentTok{\# f = I(0 \textless{}= theta \textless{}= 2)}
\NormalTok{g\_indicator }\OtherTok{\textless{}{-}} \ControlFlowTok{function}\NormalTok{(x) }\FunctionTok{as.numeric}\NormalTok{(x }\SpecialCharTok{\textgreater{}=} \DecValTok{0} \SpecialCharTok{\&}\NormalTok{ x }\SpecialCharTok{\textless{}=} \DecValTok{1}\NormalTok{)   }\CommentTok{\# g = I(0 \textless{}= theta \textless{}= 1)}
\NormalTok{support\_tbl }\OtherTok{\textless{}{-}} \FunctionTok{data.frame}\NormalTok{(}
    \AttributeTok{S =} \FunctionTok{c}\NormalTok{(}\FloatTok{1e2}\NormalTok{, }\FloatTok{1e3}\NormalTok{, }\FloatTok{1e4}\NormalTok{, }\FloatTok{1e5}\NormalTok{),}
    \AttributeTok{estimate =} \FunctionTok{sapply}\NormalTok{(}\FunctionTok{c}\NormalTok{(}\FloatTok{1e2}\NormalTok{, }\FloatTok{1e3}\NormalTok{, }\FloatTok{1e4}\NormalTok{, }\FloatTok{1e5}\NormalTok{), }\ControlFlowTok{function}\NormalTok{(S) \{}
\NormalTok{        th }\OtherTok{\textless{}{-}} \FunctionTok{runif}\NormalTok{(S, }\DecValTok{0}\NormalTok{, }\DecValTok{1}\NormalTok{)                              }\CommentTok{\# draws from g}
        \FunctionTok{mean}\NormalTok{(}\FunctionTok{f\_indicator}\NormalTok{(th) }\SpecialCharTok{/} \FunctionTok{g\_indicator}\NormalTok{(th))            }\CommentTok{\# weight is 1 for every draw}
\NormalTok{    \})}
\NormalTok{)}
\NormalTok{knitr}\SpecialCharTok{::}\FunctionTok{kable}\NormalTok{(support\_tbl, }\AttributeTok{digits =} \FunctionTok{c}\NormalTok{(}\DecValTok{0}\NormalTok{, }\DecValTok{3}\NormalTok{))}
\end{Highlighting}
\end{Shaded}

{

\begin{longtable}[]{@{}rr@{}}

\caption{\label{tbl-support-violation}Estimated C\_f / C\_g (true value
2) as the sample size grows: no amount of additional sampling helps,
because g never visits (1, 2{]}.}

\tabularnewline

\toprule\noalign{}
S & estimate \\
\midrule\noalign{}
\endhead
\bottomrule\noalign{}
\endlastfoot
1e+02 & 1 \\
1e+03 & 1 \\
1e+04 & 1 \\
1e+05 & 1 \\

\end{longtable}

}

The estimate sits at \(1\) regardless of \(S\): this is a bias that more
sampling cannot fix, because it comes from where \(g\) looks, not from
how much it looks there. The reverse situation --- \(g\) wider than
\(f\) --- is harmless: draws outside \(\text{supp}(f)\) simply receive
weight \(0\) and are correctly down-weighted. Together with the
missed-mode example, the lesson is that \(S_{\text{eff}}\) alone cannot
certify a good proposal; the proposal must be checked against what is
actually known about where \(f\) puts its mass.

\subsection{Explore it interactively}\label{explore-it-interactively-2}

The app below reproduces the missed-mode example with an adjustable
proposal: move its mean and spread and watch the effective sample size
stay high while the estimated mean drifts away from the true value of
\(0\) whenever the proposal fails to reach both components of the
target.

\begin{Shaded}
\begin{Highlighting}[]
\NormalTok{\#| label: missed{-}mode{-}app}
\NormalTok{\#| standalone: true}
\NormalTok{\#| viewerHeight: 620}

\NormalTok{library(shiny)}

\NormalTok{f\_mix \textless{}{-} function(t) 0.5 * dnorm(t + 5) + 0.5 * dnorm(t {-} 5)}
\NormalTok{ess   \textless{}{-} function(w) sum(w)\^{}2 / sum(w\^{}2)}

\NormalTok{ui \textless{}{-} fluidPage(}
\NormalTok{    titlePanel(NULL),}
\NormalTok{    sidebarLayout(}
\NormalTok{        sidebarPanel(}
\NormalTok{            width = 4,}
\NormalTok{            sliderInput("mean\_g", "proposal mean", min = {-}9, max = 9, value = {-}5, step = 0.5),}
\NormalTok{            sliderInput("sd\_g", "proposal sd", min = 0.5, max = 8, value = 1, step = 0.5),}
\NormalTok{            helpText("Target f is a 50/50 mixture of N({-}5, 1) and N(5, 1), so the ",}
\NormalTok{                     "truth is C\_f/C\_g = 1 and E\_f(theta) = 0 when g integrates to 1. ",}
\NormalTok{                     "A narrow proposal centred on one mode gives a deceptively healthy ",}
\NormalTok{                     "S\_eff while both estimates are wrong; only a proposal wide enough ",}
\NormalTok{                     "or centred between the modes recovers the truth."),}
\NormalTok{            tableOutput("tab")}
\NormalTok{        ),}
\NormalTok{        mainPanel(}
\NormalTok{            width = 8,}
\NormalTok{            plotOutput("curves", height = "260px"),}
\NormalTok{            plotOutput("weights", height = "260px")}
\NormalTok{        )}
\NormalTok{    )}
\NormalTok{)}

\NormalTok{server \textless{}{-} function(input, output, session) \{}

\NormalTok{    draws \textless{}{-} reactive(\{}
\NormalTok{        set.seed(1)}
\NormalTok{        th \textless{}{-} rnorm(5000, input$mean\_g, input$sd\_g)}
\NormalTok{        w  \textless{}{-} f\_mix(th) / dnorm(th, input$mean\_g, input$sd\_g)}
\NormalTok{        list(th = th, w = w)}
\NormalTok{    \})}

\NormalTok{    output$curves \textless{}{-} renderPlot(\{}
\NormalTok{        d \textless{}{-} draws()}
\NormalTok{        par(mar = c(4, 4, 2, 1))}
\NormalTok{        curve(f\_mix, {-}12, 12, lwd = 3, col = "firebrick", xlab = expression(theta), ylab = "",}
\NormalTok{              main = "target f and proposal g", ylim = c(0, 0.45))}
\NormalTok{        curve(dnorm(x, input$mean\_g, input$sd\_g), add = TRUE, lwd = 3, col = "steelblue")}
\NormalTok{        rug(d$th[1:300], col = adjustcolor("steelblue", 0.5))}
\NormalTok{    \})}

\NormalTok{    output$weights \textless{}{-} renderPlot(\{}
\NormalTok{        d \textless{}{-} draws(); wt \textless{}{-} d$w / sum(d$w)}
\NormalTok{        par(mar = c(4, 4, 2, 1))}
\NormalTok{        hist(wt * length(d$w), breaks = 40, col = "gray85", border = "white",}
\NormalTok{             xlab = expression(S * tilde(w)[s]),}
\NormalTok{             main = sprintf("S\_eff = \%.0f of \%d", ess(d$w), length(d$w)))}
\NormalTok{    \})}

\NormalTok{    output$tab \textless{}{-} renderTable(\{}
\NormalTok{        d \textless{}{-} draws(); wt \textless{}{-} d$w / sum(d$w)}
\NormalTok{        data.frame(quantity = c("C\_f / C\_g (true 1)", "E\_f[theta] (true 0)", "S\_eff (of 5000)"),}
\NormalTok{                   estimate = c(mean(d$w), sum(d$th * wt), ess(d$w)))}
\NormalTok{    \}, digits = 3, rownames = FALSE)}
\NormalTok{\}}

\NormalTok{shinyApp(ui, server)}
\end{Highlighting}
\end{Shaded}

\begin{tcolorbox}[enhanced jigsaw, arc=.35mm, bottomrule=.15mm, bottomtitle=1mm, breakable, colback=white, colbacktitle=quarto-callout-note-color!10!white, colframe=quarto-callout-note-color-frame, coltitle=black, left=2mm, leftrule=.75mm, opacityback=0, opacitybacktitle=0.6, rightrule=.15mm, title=\textcolor{quarto-callout-note-color}{\faInfo}\hspace{0.5em}{Reproducing this figure}, titlerule=0mm, toprule=.15mm, toptitle=1mm]

See the callout in \hyperref[numerical-quadrature]{Numerical Quadrature}
for how to install the \texttt{shinylive} extension used above.

\end{tcolorbox}

\section{Estimating a Truncated-Normal
Expectation}\label{estimating-a-truncated-normal-expectation}

Return to a truncated \(N(0,1)\) on \([l,u]\), and estimate
\(E\big[a(X)\mid l<X<u\big]\) for \(a(x)=x^2\). Two competing samplers:

\begin{itemize}
\tightlist
\item
  \textbf{Importance sampling} with \(g(x) = I(l\le x\le u)/(u-l)\)
  (uniform), giving \(w(x) = \varphi(x)\) up to the constant \((u-l)\)
  that cancels in the self-normalized estimator.
\item
  \textbf{Rejection sampling} via adaptive rejection sampling
  (\texttt{ars} package, applicable because \(\log\varphi\) is concave),
  producing exact iid draws.
\end{itemize}

\begin{Shaded}
\begin{Highlighting}[]
\DocumentationTok{\#\#\# importance{-}sampling estimate of E[a(X) | l \textless{} X \textless{} u], X \textasciitilde{} N(0,1)}
\NormalTok{est\_tnorm\_imps }\OtherTok{\textless{}{-}} \ControlFlowTok{function}\NormalTok{(a\_fn, l, u, S) \{}
\NormalTok{    X }\OtherTok{\textless{}{-}} \FunctionTok{runif}\NormalTok{(S, l, u)}
\NormalTok{    W }\OtherTok{\textless{}{-}} \FunctionTok{dnorm}\NormalTok{(X)}
\NormalTok{    ahat }\OtherTok{\textless{}{-}} \FunctionTok{sum}\NormalTok{(}\FunctionTok{a\_fn}\NormalTok{(X) }\SpecialCharTok{*}\NormalTok{ W) }\SpecialCharTok{/} \FunctionTok{sum}\NormalTok{(W)}
    \FunctionTok{attr}\NormalTok{(ahat, }\StringTok{"ess"}\NormalTok{) }\OtherTok{\textless{}{-}} \FunctionTok{ess}\NormalTok{(W)}
\NormalTok{    ahat}
\NormalTok{\}}

\DocumentationTok{\#\#\# exact draws via adaptive rejection sampling (log phi is concave)}
\NormalTok{sample\_tnorm\_ars }\OtherTok{\textless{}{-}} \ControlFlowTok{function}\NormalTok{(n, l, u) \{}
\NormalTok{    logf   }\OtherTok{\textless{}{-}} \ControlFlowTok{function}\NormalTok{(x) }\FunctionTok{dnorm}\NormalTok{(x, }\AttributeTok{log =} \ConstantTok{TRUE}\NormalTok{)}
\NormalTok{    fprima }\OtherTok{\textless{}{-}} \ControlFlowTok{function}\NormalTok{(x) }\SpecialCharTok{{-}}\NormalTok{x}
    \FunctionTok{ars}\NormalTok{(n, }\AttributeTok{f =}\NormalTok{ logf, }\AttributeTok{fprima =}\NormalTok{ fprima, }\AttributeTok{x =} \FunctionTok{c}\NormalTok{(l, (l }\SpecialCharTok{+}\NormalTok{ u) }\SpecialCharTok{/} \DecValTok{2}\NormalTok{, u),}
        \AttributeTok{lb =} \ConstantTok{TRUE}\NormalTok{, }\AttributeTok{ub =} \ConstantTok{TRUE}\NormalTok{, }\AttributeTok{xlb =}\NormalTok{ l, }\AttributeTok{xub =}\NormalTok{ u)}
\NormalTok{\}}
\end{Highlighting}
\end{Shaded}

The original version of \texttt{sample\_tnorm\_ars} always drew exactly
\(10{,}000\) points regardless of the requested \texttt{n}; the version
above respects \texttt{n}, which matters for the timing comparison next.

\begin{Shaded}
\begin{Highlighting}[]
\NormalTok{a\_fn }\OtherTok{\textless{}{-}} \ControlFlowTok{function}\NormalTok{(x) x}\SpecialCharTok{\^{}}\DecValTok{2}
\NormalTok{midpoint\_ref }\OtherTok{\textless{}{-}} \ControlFlowTok{function}\NormalTok{(l, u, }\AttributeTok{n =} \DecValTok{100000}\NormalTok{) \{}
\NormalTok{    h }\OtherTok{\textless{}{-}}\NormalTok{ (u }\SpecialCharTok{{-}}\NormalTok{ l) }\SpecialCharTok{/}\NormalTok{ n; xm }\OtherTok{\textless{}{-}}\NormalTok{ l }\SpecialCharTok{+}\NormalTok{ (}\FunctionTok{seq\_len}\NormalTok{(n) }\SpecialCharTok{{-}} \FloatTok{0.5}\NormalTok{) }\SpecialCharTok{*}\NormalTok{ h}
    \FunctionTok{sum}\NormalTok{(}\FunctionTok{a\_fn}\NormalTok{(xm) }\SpecialCharTok{*} \FunctionTok{dnorm}\NormalTok{(xm)) }\SpecialCharTok{*}\NormalTok{ h }\SpecialCharTok{/}\NormalTok{ (}\FunctionTok{pnorm}\NormalTok{(u) }\SpecialCharTok{{-}} \FunctionTok{pnorm}\NormalTok{(l))}
\NormalTok{\}}

\NormalTok{intervals }\OtherTok{\textless{}{-}} \FunctionTok{list}\NormalTok{(}\FunctionTok{c}\NormalTok{(}\DecValTok{1}\NormalTok{, }\DecValTok{2}\NormalTok{), }\FunctionTok{c}\NormalTok{(}\SpecialCharTok{{-}}\DecValTok{1}\NormalTok{, }\DecValTok{1}\NormalTok{))}
\NormalTok{tnorm\_tbl }\OtherTok{\textless{}{-}} \FunctionTok{do.call}\NormalTok{(rbind, }\FunctionTok{lapply}\NormalTok{(intervals, }\ControlFlowTok{function}\NormalTok{(iv) \{}
\NormalTok{    l }\OtherTok{\textless{}{-}}\NormalTok{ iv[}\DecValTok{1}\NormalTok{]; u }\OtherTok{\textless{}{-}}\NormalTok{ iv[}\DecValTok{2}\NormalTok{]; ref }\OtherTok{\textless{}{-}} \FunctionTok{midpoint\_ref}\NormalTok{(l, u)}
\NormalTok{    t\_imps }\OtherTok{\textless{}{-}} \FunctionTok{system.time}\NormalTok{(EA\_imps }\OtherTok{\textless{}{-}} \FunctionTok{replicate}\NormalTok{(}\DecValTok{100}\NormalTok{, }\FunctionTok{est\_tnorm\_imps}\NormalTok{(a\_fn, l, u, }\DecValTok{1000}\NormalTok{)))}
\NormalTok{    t\_rej  }\OtherTok{\textless{}{-}} \FunctionTok{system.time}\NormalTok{(EA\_rej  }\OtherTok{\textless{}{-}} \FunctionTok{replicate}\NormalTok{(}\DecValTok{100}\NormalTok{, }\FunctionTok{mean}\NormalTok{(}\FunctionTok{a\_fn}\NormalTok{(}\FunctionTok{sample\_tnorm\_ars}\NormalTok{(}\DecValTok{1000}\NormalTok{, l, u)))))}
    \FunctionTok{data.frame}\NormalTok{(}\AttributeTok{interval =} \FunctionTok{sprintf}\NormalTok{(}\StringTok{"(\%g, \%g)"}\NormalTok{, l, u), }\AttributeTok{reference =}\NormalTok{ ref,}
               \AttributeTok{mean\_imps =} \FunctionTok{mean}\NormalTok{(EA\_imps), }\AttributeTok{mse\_imps =} \FunctionTok{mean}\NormalTok{((EA\_imps }\SpecialCharTok{{-}}\NormalTok{ ref)}\SpecialCharTok{\^{}}\DecValTok{2}\NormalTok{),}
               \AttributeTok{ms\_imps =} \DecValTok{1000} \SpecialCharTok{*}\NormalTok{ t\_imps[[}\StringTok{"elapsed"}\NormalTok{]] }\SpecialCharTok{/} \DecValTok{100}\NormalTok{,}
               \AttributeTok{mean\_rej =} \FunctionTok{mean}\NormalTok{(EA\_rej), }\AttributeTok{mse\_rej =} \FunctionTok{mean}\NormalTok{((EA\_rej }\SpecialCharTok{{-}}\NormalTok{ ref)}\SpecialCharTok{\^{}}\DecValTok{2}\NormalTok{),}
               \AttributeTok{ms\_rej =} \DecValTok{1000} \SpecialCharTok{*}\NormalTok{ t\_rej[[}\StringTok{"elapsed"}\NormalTok{]] }\SpecialCharTok{/} \DecValTok{100}\NormalTok{)}
\NormalTok{\}))}
\NormalTok{knitr}\SpecialCharTok{::}\FunctionTok{kable}\NormalTok{(tnorm\_tbl, }\AttributeTok{digits =} \FunctionTok{c}\NormalTok{(}\DecValTok{0}\NormalTok{, }\DecValTok{4}\NormalTok{, }\DecValTok{4}\NormalTok{, }\DecValTok{6}\NormalTok{, }\DecValTok{2}\NormalTok{, }\DecValTok{4}\NormalTok{, }\DecValTok{6}\NormalTok{, }\DecValTok{2}\NormalTok{))}
\end{Highlighting}
\end{Shaded}

{

\begin{longtable}[]{@{}
  >{\raggedright\arraybackslash}p{(\linewidth - 14\tabcolsep) * \real{0.1268}}
  >{\raggedleft\arraybackslash}p{(\linewidth - 14\tabcolsep) * \real{0.1408}}
  >{\raggedleft\arraybackslash}p{(\linewidth - 14\tabcolsep) * \real{0.1408}}
  >{\raggedleft\arraybackslash}p{(\linewidth - 14\tabcolsep) * \real{0.1268}}
  >{\raggedleft\arraybackslash}p{(\linewidth - 14\tabcolsep) * \real{0.1127}}
  >{\raggedleft\arraybackslash}p{(\linewidth - 14\tabcolsep) * \real{0.1268}}
  >{\raggedleft\arraybackslash}p{(\linewidth - 14\tabcolsep) * \real{0.1268}}
  >{\raggedleft\arraybackslash}p{(\linewidth - 14\tabcolsep) * \real{0.0986}}@{}}

\caption{\label{tbl-tnorm-compare}Importance sampling versus adaptive
rejection sampling for E{[}X\^{}2 \textbar{} l \textless{} X \textless{}
u{]}, averaged over 100 replications at S = 1000 draws each. Both
recover the midpoint-rule reference to the precision shown; timings are
milliseconds per replication.}

\tabularnewline

\toprule\noalign{}
\begin{minipage}[b]{\linewidth}\raggedright
interval
\end{minipage} & \begin{minipage}[b]{\linewidth}\raggedleft
reference
\end{minipage} & \begin{minipage}[b]{\linewidth}\raggedleft
mean\_imps
\end{minipage} & \begin{minipage}[b]{\linewidth}\raggedleft
mse\_imps
\end{minipage} & \begin{minipage}[b]{\linewidth}\raggedleft
ms\_imps
\end{minipage} & \begin{minipage}[b]{\linewidth}\raggedleft
mean\_rej
\end{minipage} & \begin{minipage}[b]{\linewidth}\raggedleft
mse\_rej
\end{minipage} & \begin{minipage}[b]{\linewidth}\raggedleft
ms\_rej
\end{minipage} \\
\midrule\noalign{}
\endhead
\bottomrule\noalign{}
\endlastfoot
(1, 2) & 1.9859 & 1.9871 & 0.000649 & 0.05 & 1.9854 & 0.000687 & 3.24 \\
(-1, 1) & 0.2911 & 0.2907 & 0.000060 & 0.03 & 0.2923 & 0.000099 &
3.22 \\

\end{longtable}

}

Both estimators are unbiased for this smooth, bounded target and agree
with the midpoint-rule reference; importance sampling is markedly faster
here because it needs no accept/reject loop or envelope construction, at
the cost of returning a \emph{weighted} rather than an exact sample.

\section{Application: Marginal Likelihood via Importance
Sampling}\label{application-marginal-likelihood-via-importance-sampling}

\subsection{Setup}\label{setup-1}

The \hyperref[numerical-quadrature]{quadrature} and
Laplace-approximation chapters compute the marginal likelihood
\(P(y) = \int L(\theta)\pi(\theta)\,d\theta\) of the same conjugate-free
Gaussian model --- unknown mean \(\mu\) and \(w=\log\sigma^2\),
independent normal priors --- by a product grid and by a
mode-and-Hessian approximation, respectively. Importance sampling offers
a third route, and one that comes with a built-in accuracy diagnostic
(\(S_{\text{eff}}\)) that neither of the other two methods provides
directly.

Take \(g(\mu,w)\) to be the \textbf{Gaussian approximation to the
posterior} obtained by a Laplace fit: locate the posterior mode
\((\hat\mu,\hat w)\) by optimization, take
\(g = N_2\big((\hat\mu,\hat w),\, H^{-1}\big)\) for \(H\) the Hessian of
\(-\log[L\pi]\) at the mode, and importance-sample from \(g\) with
\(f = L\pi\): \[
\widehat{P(y)} = \frac{1}{S}\sum_{s=1}^S w_s, \qquad
w_s = \frac{L(\theta_s)\pi(\theta_s)}{g(\theta_s)}, \qquad \theta_s \sim g.
\] Because \(g\) is shaped like the posterior near its mode, this is a
``good story'' proposal by construction --- in sharp contrast to
sampling directly from the (typically much wider, uninformative) prior,
which is exactly the \texttt{log\_marlik\_mc} function used elsewhere in
this book as a naive-Monte-Carlo check: it is importance sampling with
\(g = \pi\), and its weights are just the likelihood values
\(L(\theta_s)\).

\subsection{An efficient
implementation}\label{an-efficient-implementation-1}

\begin{Shaded}
\begin{Highlighting}[]
\NormalTok{log\_lik   }\OtherTok{\textless{}{-}} \ControlFlowTok{function}\NormalTok{(x, mu, w) }\FunctionTok{sum}\NormalTok{(}\FunctionTok{dnorm}\NormalTok{(x, mu, }\FunctionTok{exp}\NormalTok{(w }\SpecialCharTok{/} \DecValTok{2}\NormalTok{), }\AttributeTok{log =} \ConstantTok{TRUE}\NormalTok{))}
\NormalTok{log\_prior }\OtherTok{\textless{}{-}} \ControlFlowTok{function}\NormalTok{(mu, w, mu\_0, sigma\_mu, w\_0, sigma\_w) \{}
    \FunctionTok{dnorm}\NormalTok{(mu, mu\_0, sigma\_mu, }\AttributeTok{log =} \ConstantTok{TRUE}\NormalTok{) }\SpecialCharTok{+} \FunctionTok{dnorm}\NormalTok{(w, w\_0, sigma\_w, }\AttributeTok{log =} \ConstantTok{TRUE}\NormalTok{)}
\NormalTok{\}}
\NormalTok{neg\_log\_post }\OtherTok{\textless{}{-}} \ControlFlowTok{function}\NormalTok{(theta, x, mu\_0, sigma\_mu, w\_0, sigma\_w) \{}
    \SpecialCharTok{{-}}\FunctionTok{log\_lik}\NormalTok{(x, theta[}\DecValTok{1}\NormalTok{], theta[}\DecValTok{2}\NormalTok{]) }\SpecialCharTok{{-}}
        \FunctionTok{log\_prior}\NormalTok{(theta[}\DecValTok{1}\NormalTok{], theta[}\DecValTok{2}\NormalTok{], mu\_0, sigma\_mu, w\_0, sigma\_w)}
\NormalTok{\}}

\DocumentationTok{\#\#\# log density of N(mu, A\^{}\{{-}1\}) at every column of a 2 x S matrix of points at once,}
\DocumentationTok{\#\#\# replacing a per{-}point apply() call with one vectorized quadratic form}
\NormalTok{log\_dmvnorm\_batch }\OtherTok{\textless{}{-}} \ControlFlowTok{function}\NormalTok{(Theta, mu, A) \{}
\NormalTok{    Tc }\OtherTok{\textless{}{-}}\NormalTok{ Theta }\SpecialCharTok{{-}}\NormalTok{ mu}
\NormalTok{    quad }\OtherTok{\textless{}{-}} \FunctionTok{colSums}\NormalTok{((A }\SpecialCharTok{\%*\%}\NormalTok{ Tc) }\SpecialCharTok{*}\NormalTok{ Tc)}
    \FloatTok{0.5} \SpecialCharTok{*}\NormalTok{ (}\SpecialCharTok{{-}}\FunctionTok{nrow}\NormalTok{(Theta) }\SpecialCharTok{*} \FunctionTok{log}\NormalTok{(}\DecValTok{2} \SpecialCharTok{*}\NormalTok{ pi) }\SpecialCharTok{+} \FunctionTok{sum}\NormalTok{(}\FunctionTok{log}\NormalTok{(}\FunctionTok{svd}\NormalTok{(A)}\SpecialCharTok{$}\NormalTok{d)) }\SpecialCharTok{{-}}\NormalTok{ quad)}
\NormalTok{\}}

\DocumentationTok{\#\#\# log marginal likelihood by importance sampling from a Laplace{-}shaped Gaussian}
\NormalTok{log\_marlik\_imps }\OtherTok{\textless{}{-}} \ControlFlowTok{function}\NormalTok{(x, mu\_0, sigma\_mu, w\_0, sigma\_w, S) \{}
\NormalTok{    fit }\OtherTok{\textless{}{-}} \FunctionTok{nlm}\NormalTok{(neg\_log\_post, }\FunctionTok{c}\NormalTok{(}\FunctionTok{mean}\NormalTok{(x), }\FunctionTok{log}\NormalTok{(stats}\SpecialCharTok{::}\FunctionTok{var}\NormalTok{(x))), }\AttributeTok{hessian =} \ConstantTok{TRUE}\NormalTok{,}
               \AttributeTok{x =}\NormalTok{ x, }\AttributeTok{mu\_0 =}\NormalTok{ mu\_0, }\AttributeTok{sigma\_mu =}\NormalTok{ sigma\_mu, }\AttributeTok{w\_0 =}\NormalTok{ w\_0, }\AttributeTok{sigma\_w =}\NormalTok{ sigma\_w)}
\NormalTok{    mu\_hat }\OtherTok{\textless{}{-}}\NormalTok{ fit}\SpecialCharTok{$}\NormalTok{estimate; H }\OtherTok{\textless{}{-}}\NormalTok{ fit}\SpecialCharTok{$}\NormalTok{hessian}
\NormalTok{    L }\OtherTok{\textless{}{-}} \FunctionTok{t}\NormalTok{(}\FunctionTok{chol}\NormalTok{(}\FunctionTok{solve}\NormalTok{(H)))                        }\CommentTok{\# Sigma = H\^{}\{{-}1\} = L L\textquotesingle{}}
\NormalTok{    Theta }\OtherTok{\textless{}{-}}\NormalTok{ L }\SpecialCharTok{\%*\%} \FunctionTok{matrix}\NormalTok{(}\FunctionTok{rnorm}\NormalTok{(}\DecValTok{2} \SpecialCharTok{*}\NormalTok{ S), }\DecValTok{2}\NormalTok{, S) }\SpecialCharTok{+}\NormalTok{ mu\_hat}

\NormalTok{    log\_g }\OtherTok{\textless{}{-}} \FunctionTok{log\_dmvnorm\_batch}\NormalTok{(Theta, mu\_hat, H)}
\NormalTok{    log\_f }\OtherTok{\textless{}{-}} \SpecialCharTok{{-}}\FunctionTok{apply}\NormalTok{(Theta, }\DecValTok{2}\NormalTok{, neg\_log\_post, }\AttributeTok{x =}\NormalTok{ x, }\AttributeTok{mu\_0 =}\NormalTok{ mu\_0, }\AttributeTok{sigma\_mu =}\NormalTok{ sigma\_mu,}
                     \AttributeTok{w\_0 =}\NormalTok{ w\_0, }\AttributeTok{sigma\_w =}\NormalTok{ sigma\_w)}
\NormalTok{    lw }\OtherTok{\textless{}{-}}\NormalTok{ log\_f }\SpecialCharTok{{-}}\NormalTok{ log\_g}
    \FunctionTok{structure}\NormalTok{(}\FunctionTok{log\_sum\_exp}\NormalTok{(lw) }\SpecialCharTok{{-}} \FunctionTok{log}\NormalTok{(S), }\AttributeTok{ess =} \FunctionTok{log\_ess}\NormalTok{(lw))}
\NormalTok{\}}

\DocumentationTok{\#\#\# log marginal likelihood by naive Monte Carlo from the prior (g = pi): a}
\DocumentationTok{\#\#\# second, independent importance{-}sampling estimator with an unfavourable g}
\NormalTok{log\_marlik\_prior }\OtherTok{\textless{}{-}} \ControlFlowTok{function}\NormalTok{(x, mu\_0, sigma\_mu, w\_0, sigma\_w, S) \{}
\NormalTok{    mus }\OtherTok{\textless{}{-}} \FunctionTok{rnorm}\NormalTok{(S, mu\_0, sigma\_mu); ws }\OtherTok{\textless{}{-}} \FunctionTok{rnorm}\NormalTok{(S, w\_0, sigma\_w)}
\NormalTok{    lw }\OtherTok{\textless{}{-}} \FunctionTok{vapply}\NormalTok{(}\FunctionTok{seq\_len}\NormalTok{(S), }\ControlFlowTok{function}\NormalTok{(i) }\FunctionTok{log\_lik}\NormalTok{(x, mus[i], ws[i]), }\FunctionTok{numeric}\NormalTok{(}\DecValTok{1}\NormalTok{))}
    \FunctionTok{structure}\NormalTok{(}\FunctionTok{log\_sum\_exp}\NormalTok{(lw) }\SpecialCharTok{{-}} \FunctionTok{log}\NormalTok{(S), }\AttributeTok{ess =} \FunctionTok{log\_ess}\NormalTok{(lw))}
\NormalTok{\}}
\end{Highlighting}
\end{Shaded}

Replacing the two \texttt{apply()} calls used in earlier drafts of this
code --- one for the proposal log-density, one for the log-posterior ---
with a single vectorized quadratic form (\texttt{log\_dmvnorm\_batch})
removes \(S\) separate R function calls for the proposal evaluation;
only the data-dependent log-likelihood still needs one call per draw,
since it sums over the \(n\) observations for each \(\theta_s\)
separately.

\subsection{Checking accuracy}\label{checking-accuracy}

\begin{Shaded}
\begin{Highlighting}[]
\FunctionTok{set.seed}\NormalTok{(}\DecValTok{4}\NormalTok{)}
\NormalTok{x\_small }\OtherTok{\textless{}{-}} \FunctionTok{rnorm}\NormalTok{(}\DecValTok{10}\NormalTok{);  x\_large }\OtherTok{\textless{}{-}} \FunctionTok{rnorm}\NormalTok{(}\DecValTok{200}\NormalTok{)}

\NormalTok{check\_tbl }\OtherTok{\textless{}{-}} \FunctionTok{do.call}\NormalTok{(rbind, }\FunctionTok{lapply}\NormalTok{(}\FunctionTok{list}\NormalTok{(}\AttributeTok{n10 =}\NormalTok{ x\_small, }\AttributeTok{n200 =}\NormalTok{ x\_large), }\ControlFlowTok{function}\NormalTok{(x) \{}
\NormalTok{    reps }\OtherTok{\textless{}{-}} \FunctionTok{replicate}\NormalTok{(}\DecValTok{500}\NormalTok{, }\FunctionTok{log\_marlik\_imps}\NormalTok{(x, }\DecValTok{0}\NormalTok{, }\DecValTok{5}\NormalTok{, }\DecValTok{0}\NormalTok{, }\DecValTok{5}\NormalTok{, }\DecValTok{200}\NormalTok{), }\AttributeTok{simplify =} \ConstantTok{FALSE}\NormalTok{)}
\NormalTok{    vals }\OtherTok{\textless{}{-}} \FunctionTok{vapply}\NormalTok{(reps, as.numeric, }\FunctionTok{numeric}\NormalTok{(}\DecValTok{1}\NormalTok{))}
\NormalTok{    esss }\OtherTok{\textless{}{-}} \FunctionTok{vapply}\NormalTok{(reps, attr, }\FunctionTok{numeric}\NormalTok{(}\DecValTok{1}\NormalTok{), }\AttributeTok{which =} \StringTok{"ess"}\NormalTok{)}
    \FunctionTok{data.frame}\NormalTok{(}\AttributeTok{mean\_log\_ml =} \FunctionTok{mean}\NormalTok{(vals), }\AttributeTok{sd\_log\_ml =} \FunctionTok{sd}\NormalTok{(vals), }\AttributeTok{mean\_ess =} \FunctionTok{mean}\NormalTok{(esss))}
\NormalTok{\}))}
\NormalTok{check\_tbl }\OtherTok{\textless{}{-}} \FunctionTok{cbind}\NormalTok{(}\AttributeTok{dataset =} \FunctionTok{c}\NormalTok{(}\StringTok{"n = 10"}\NormalTok{, }\StringTok{"n = 200"}\NormalTok{), check\_tbl)}
\NormalTok{knitr}\SpecialCharTok{::}\FunctionTok{kable}\NormalTok{(check\_tbl, }\AttributeTok{digits =} \DecValTok{4}\NormalTok{)}
\end{Highlighting}
\end{Shaded}

{

\begin{longtable}[]{@{}llrrr@{}}

\caption{\label{tbl-marlik-check}Log marginal likelihood by importance
sampling (Laplace-shaped proposal, S = 200) across 500 replications, for
two simulated data sets.}

\tabularnewline

\toprule\noalign{}
& dataset & mean\_log\_ml & sd\_log\_ml & mean\_ess \\
\midrule\noalign{}
\endhead
\bottomrule\noalign{}
\endlastfoot
n10 & n = 10 & -19.2284 & 0.1447 & 133.4968 \\
n200 & n = 200 & -285.8823 & 0.0091 & 196.4443 \\

\end{longtable}

}

The Monte Carlo standard deviation across replications is small relative
to the estimate itself, and the effective sample size stays close to the
nominal \(S=200\) in both cases --- evidence that the Laplace-shaped
proposal tracks the posterior well even though it is, after all, only a
second-order approximation to it.

\subsection{Why the proposal matters: Laplace-shaped versus prior
sampling}\label{why-the-proposal-matters-laplace-shaped-versus-prior-sampling}

The prior is a valid proposal for the same target --- every posterior
draw is also a prior draw with positive density --- so
\texttt{log\_marlik\_prior} is a legitimate importance-sampling
estimator too. But the prior is far more diffuse than the posterior,
especially once \(n\) is not tiny, so it is a much poorer match to
\(f = L\pi\).

\begin{Shaded}
\begin{Highlighting}[]
\FunctionTok{set.seed}\NormalTok{(}\DecValTok{5}\NormalTok{)}
\NormalTok{ns }\OtherTok{\textless{}{-}} \FunctionTok{c}\NormalTok{(}\DecValTok{5}\NormalTok{, }\DecValTok{20}\NormalTok{, }\DecValTok{100}\NormalTok{, }\DecValTok{400}\NormalTok{)}
\NormalTok{ess\_tbl }\OtherTok{\textless{}{-}} \FunctionTok{do.call}\NormalTok{(rbind, }\FunctionTok{lapply}\NormalTok{(ns, }\ControlFlowTok{function}\NormalTok{(n) \{}
\NormalTok{    x }\OtherTok{\textless{}{-}} \FunctionTok{rnorm}\NormalTok{(n)}
    \FunctionTok{data.frame}\NormalTok{(}\AttributeTok{n =}\NormalTok{ n,}
               \AttributeTok{ess\_laplace\_proposal =} \FunctionTok{attr}\NormalTok{(}\FunctionTok{log\_marlik\_imps}\NormalTok{(x, }\DecValTok{0}\NormalTok{, }\DecValTok{5}\NormalTok{, }\DecValTok{0}\NormalTok{, }\DecValTok{5}\NormalTok{, }\DecValTok{500}\NormalTok{), }\StringTok{"ess"}\NormalTok{),}
               \AttributeTok{ess\_prior\_sampling   =} \FunctionTok{attr}\NormalTok{(}\FunctionTok{log\_marlik\_prior}\NormalTok{(x, }\DecValTok{0}\NormalTok{, }\DecValTok{5}\NormalTok{, }\DecValTok{0}\NormalTok{, }\DecValTok{5}\NormalTok{, }\DecValTok{500}\NormalTok{), }\StringTok{"ess"}\NormalTok{))}
\NormalTok{\}))}
\NormalTok{knitr}\SpecialCharTok{::}\FunctionTok{kable}\NormalTok{(ess\_tbl, }\AttributeTok{digits =} \FunctionTok{c}\NormalTok{(}\DecValTok{0}\NormalTok{, }\DecValTok{1}\NormalTok{, }\DecValTok{1}\NormalTok{))}
\end{Highlighting}
\end{Shaded}

{

\begin{longtable}[]{@{}rrr@{}}

\caption{\label{tbl-marlik-ess-vs-n}Effective sample size (of S = 500
draws) for the Laplace-shaped proposal versus prior sampling, as the
sample size n of the simulated data grows. Both target the same
posterior for (mu, w).}

\tabularnewline

\toprule\noalign{}
n & ess\_laplace\_proposal & ess\_prior\_sampling \\
\midrule\noalign{}
\endhead
\bottomrule\noalign{}
\endlastfoot
5 & 258.8 & 19.5 \\
20 & 341.4 & 3.3 \\
100 & 476.0 & 1.0 \\
400 & 496.0 & 1.0 \\

\end{longtable}

}

As \(n\) grows the posterior sharpens around the truth while the prior
does not, so prior sampling's effective sample size deteriorates
steadily; the Laplace-shaped proposal, which re-centers and re-scales
itself around the posterior mode for every data set, keeps pace. This is
the marginal-likelihood instance of the same lesson as the missed-mode
example above: a fixed, data-independent proposal becomes an
increasingly bad match to a target that concentrates as more data
arrive, and only a proposal that adapts to the target avoids the
resulting collapse in efficiency.

\section{Summary}\label{summary-5}

\begin{itemize}
\tightlist
\item
  Importance sampling reweights draws from a proposal \(g\) by
  \(w(\theta) = f(\theta)/g(\theta)\) instead of accepting or rejecting
  them, using \(\widehat{C_f/C_g} = \frac1S\sum_s w_s\) and the
  self-normalized \(\hat A = \sum_s a(\theta_s)\tilde w_s\); only the
  support condition \(\{g>0\}\supseteq\{f>0\}\) is required, not an
  envelope \(g \ge f\).
\item
  A good proposal shaped like the target on its support gives
  near-constant weights and \(S_{\text{eff}} \approx S\) (the
  tail-probability and Laplace-proposal examples).
\item
  A proposal that misses a mode or part of the support fails
  \textbf{silently}: the observed weights can look perfectly healthy
  while both \(C_f/C_g\) and \(E_f[a]\) are badly wrong, and no amount
  of extra sampling fixes it. \(S_{\text{eff}}\) diagnoses only the
  noisy failure (wildly varying observed weights), never the silent one.
\item
  For the marginal likelihood, a proposal built from a Laplace
  approximation at the posterior mode remains efficient as the data
  accumulate and the posterior sharpens, while a fixed, data-independent
  proposal such as the prior does not --- reinforcing why the
  mode-finding methods from earlier chapters remain useful even once we
  turn to sampling-based methods.
\end{itemize}

\bookmarksetup{startatroot}

\chapter{Markov Chain Monte Carlo}\label{markov-chain-monte-carlo}

\begin{Shaded}
\begin{Highlighting}[]
\FunctionTok{library}\NormalTok{(truncnorm)}
\FunctionTok{library}\NormalTok{(rstan)}
\FunctionTok{rstan\_options}\NormalTok{(}\AttributeTok{auto\_write =} \ConstantTok{TRUE}\NormalTok{)}
\end{Highlighting}
\end{Shaded}

The last two chapters drew from a \textbf{fixed} proposal \(g\): an
envelope for rejection sampling, an importance distribution for
reweighting. Both work well when \(g\) can be made to resemble the
target \(f\) over its whole support, but in more than a handful of
dimensions no fixed \(g\) does --- the target occupies an ever-thinner
sliver of the space \(g\) covers, exactly the curse of dimensionality
met already for quadrature grids. \textbf{Markov chain Monte Carlo
(MCMC)} answers this differently: instead of one global proposal, build
a \textbf{random process}
\(\theta^{(0)}, \theta^{(1)}, \theta^{(2)}, \ldots\) that moves
\emph{locally} --- each step depends only on where the chain currently
is --- and is engineered so that, after enough steps, \(\theta^{(t)}\)
is (approximately) distributed as the target \(f\). The price is that
consecutive draws are dependent and the guarantee is only asymptotic;
the payoff is that the same local-move idea scales to any dimension.

This chapter develops the Markov chain theory that makes this work, then
builds three increasingly automatic ways to construct such a chain ---
Gibbs sampling, Metropolis--Hastings, and Hamiltonian Monte Carlo ---
and ends with Stan, which automates the last of these for essentially
arbitrary models.

\section{Markov Chains}\label{markov-chains}

\subsection{Definition and the convergence
question}\label{definition-and-the-convergence-question}

A sequence \(\theta^{(0)}, \theta^{(1)}, \theta^{(2)}, \ldots\) is a
\textbf{Markov chain} if the future depends on the past only through the
present, \[
P\big(\theta^{(i)} \mid \theta^{(0)}, \ldots, \theta^{(i-1)}\big) = T\big(\theta^{(i)} \mid \theta^{(i-1)}\big),
\] where \(T\), the \textbf{transition distribution}, does not depend on
\(i\). A chain is completely specified by a starting value (or
distribution) and \(T\); simulating it just means repeatedly drawing
\(\theta^{(i)} \sim T(\cdot \mid \theta^{(i-1)})\).

The question MCMC answers is: given a target \(\pi\) known only up to a
constant, can we \emph{design} a \(T\) such that, no matter where the
chain starts, \(\theta^{(i)}\) converges in distribution to \(\pi\) as
\(i \to \infty\)? If so, running the chain for a long time and averaging
\(a(\theta^{(i)})\) over the later iterations estimates
\(E_\pi[a(\theta)]\) --- this is what ``MCMC'' means in practice.

\subsection{A worked example: two random walks on a
cycle}\label{a-worked-example-two-random-walks-on-a-cycle}

To make ``designing \(T\)'' concrete before turning to theory, consider
a chain on the ten states \(\{0, 1, \ldots, 9\}\), arranged on a circle.
At each step it moves one position clockwise with probability \(0.45\),
one position counter-clockwise with probability \(0.45\), and stays put
with probability \(0.1\) --- a \emph{cyclic} random walk, since state
\(9\)'s clockwise neighbour wraps around to state \(0\):

\begin{Shaded}
\begin{Highlighting}[]
\NormalTok{sim\_cyclic\_walk }\OtherTok{\textless{}{-}} \ControlFlowTok{function}\NormalTok{(ini, iters, }\AttributeTok{n =} \DecValTok{10}\NormalTok{) \{}
\NormalTok{    forward  }\OtherTok{\textless{}{-}} \ControlFlowTok{function}\NormalTok{(i) }\ControlFlowTok{if}\NormalTok{ (i }\SpecialCharTok{==}\NormalTok{ n }\SpecialCharTok{{-}} \DecValTok{1}\NormalTok{) }\DecValTok{0} \ControlFlowTok{else}\NormalTok{ i }\SpecialCharTok{+} \DecValTok{1}
\NormalTok{    backward }\OtherTok{\textless{}{-}} \ControlFlowTok{function}\NormalTok{(i) }\ControlFlowTok{if}\NormalTok{ (i }\SpecialCharTok{==} \DecValTok{0}\NormalTok{) n }\SpecialCharTok{{-}} \DecValTok{1} \ControlFlowTok{else}\NormalTok{ i }\SpecialCharTok{{-}} \DecValTok{1}
\NormalTok{    mc }\OtherTok{\textless{}{-}} \FunctionTok{integer}\NormalTok{(iters }\SpecialCharTok{+} \DecValTok{1}\NormalTok{); mc[}\DecValTok{1}\NormalTok{] }\OtherTok{\textless{}{-}}\NormalTok{ ini}
    \ControlFlowTok{for}\NormalTok{ (i }\ControlFlowTok{in} \DecValTok{2}\SpecialCharTok{:}\NormalTok{(iters }\SpecialCharTok{+} \DecValTok{1}\NormalTok{)) \{}
\NormalTok{        u }\OtherTok{\textless{}{-}} \FunctionTok{runif}\NormalTok{(}\DecValTok{1}\NormalTok{)}
\NormalTok{        mc[i] }\OtherTok{\textless{}{-}} \ControlFlowTok{if}\NormalTok{ (u }\SpecialCharTok{\textless{}} \FloatTok{0.45}\NormalTok{) }\FunctionTok{forward}\NormalTok{(mc[i }\SpecialCharTok{{-}} \DecValTok{1}\NormalTok{])}
                 \ControlFlowTok{else} \ControlFlowTok{if}\NormalTok{ (u }\SpecialCharTok{\textgreater{}} \FloatTok{0.55}\NormalTok{) }\FunctionTok{backward}\NormalTok{(mc[i }\SpecialCharTok{{-}} \DecValTok{1}\NormalTok{])}
                 \ControlFlowTok{else}\NormalTok{ mc[i }\SpecialCharTok{{-}} \DecValTok{1}\NormalTok{]}
\NormalTok{    \}}
\NormalTok{    mc}
\NormalTok{\}}
\end{Highlighting}
\end{Shaded}

Now change only the two boundary states: at \(0\), ``backward'' does
nothing (stays at \(0\)) instead of wrapping to \(9\); at \(9\),
``forward'' does nothing instead of wrapping to \(0\). This is a
\emph{reflecting} walk on a line segment rather than a cycle:

\begin{Shaded}
\begin{Highlighting}[]
\NormalTok{sim\_reflecting\_walk }\OtherTok{\textless{}{-}} \ControlFlowTok{function}\NormalTok{(ini, iters, }\AttributeTok{n =} \DecValTok{10}\NormalTok{) \{}
\NormalTok{    forward  }\OtherTok{\textless{}{-}} \ControlFlowTok{function}\NormalTok{(i) }\ControlFlowTok{if}\NormalTok{ (i }\SpecialCharTok{==}\NormalTok{ n }\SpecialCharTok{{-}} \DecValTok{1}\NormalTok{) n }\SpecialCharTok{{-}} \DecValTok{1} \ControlFlowTok{else}\NormalTok{ i }\SpecialCharTok{+} \DecValTok{1}
\NormalTok{    backward }\OtherTok{\textless{}{-}} \ControlFlowTok{function}\NormalTok{(i) }\ControlFlowTok{if}\NormalTok{ (i }\SpecialCharTok{==} \DecValTok{0}\NormalTok{) }\DecValTok{0} \ControlFlowTok{else}\NormalTok{ i }\SpecialCharTok{{-}} \DecValTok{1}
\NormalTok{    mc }\OtherTok{\textless{}{-}} \FunctionTok{integer}\NormalTok{(iters }\SpecialCharTok{+} \DecValTok{1}\NormalTok{); mc[}\DecValTok{1}\NormalTok{] }\OtherTok{\textless{}{-}}\NormalTok{ ini}
    \ControlFlowTok{for}\NormalTok{ (i }\ControlFlowTok{in} \DecValTok{2}\SpecialCharTok{:}\NormalTok{(iters }\SpecialCharTok{+} \DecValTok{1}\NormalTok{)) \{}
\NormalTok{        u }\OtherTok{\textless{}{-}} \FunctionTok{runif}\NormalTok{(}\DecValTok{1}\NormalTok{)}
\NormalTok{        mc[i] }\OtherTok{\textless{}{-}} \ControlFlowTok{if}\NormalTok{ (u }\SpecialCharTok{\textless{}} \FloatTok{0.4}\NormalTok{) }\FunctionTok{forward}\NormalTok{(mc[i }\SpecialCharTok{{-}} \DecValTok{1}\NormalTok{])}
                 \ControlFlowTok{else} \ControlFlowTok{if}\NormalTok{ (u }\SpecialCharTok{\textgreater{}} \FloatTok{0.8}\NormalTok{) }\FunctionTok{backward}\NormalTok{(mc[i }\SpecialCharTok{{-}} \DecValTok{1}\NormalTok{])}
                 \ControlFlowTok{else}\NormalTok{ mc[i }\SpecialCharTok{{-}} \DecValTok{1}\NormalTok{]}
\NormalTok{    \}}
\NormalTok{    mc}
\NormalTok{\}}
\end{Highlighting}
\end{Shaded}

\begin{Shaded}
\begin{Highlighting}[]
\FunctionTok{set.seed}\NormalTok{(}\DecValTok{1}\NormalTok{)}
\NormalTok{cyc }\OtherTok{\textless{}{-}} \FunctionTok{sim\_cyclic\_walk}\NormalTok{(}\DecValTok{2}\NormalTok{, }\DecValTok{20000}\NormalTok{)}
\NormalTok{refl }\OtherTok{\textless{}{-}} \FunctionTok{sim\_reflecting\_walk}\NormalTok{(}\DecValTok{2}\NormalTok{, }\DecValTok{20000}\NormalTok{)}
\FunctionTok{par}\NormalTok{(}\AttributeTok{mfrow =} \FunctionTok{c}\NormalTok{(}\DecValTok{1}\NormalTok{, }\DecValTok{2}\NormalTok{), }\AttributeTok{mar =} \FunctionTok{c}\NormalTok{(}\DecValTok{4}\NormalTok{, }\DecValTok{4}\NormalTok{, }\DecValTok{2}\NormalTok{, }\DecValTok{1}\NormalTok{))}
\FunctionTok{barplot}\NormalTok{(}\FunctionTok{table}\NormalTok{(cyc[}\SpecialCharTok{{-}}\NormalTok{(}\DecValTok{1}\SpecialCharTok{:}\DecValTok{200}\NormalTok{)]) }\SpecialCharTok{/} \FunctionTok{length}\NormalTok{(cyc[}\SpecialCharTok{{-}}\NormalTok{(}\DecValTok{1}\SpecialCharTok{:}\DecValTok{200}\NormalTok{)]), }\AttributeTok{col =} \StringTok{"steelblue"}\NormalTok{,}
        \AttributeTok{xlab =} \StringTok{"state"}\NormalTok{, }\AttributeTok{ylab =} \StringTok{"frequency"}\NormalTok{, }\AttributeTok{main =} \StringTok{"cyclic walk (wraps around)"}\NormalTok{)}
\FunctionTok{abline}\NormalTok{(}\AttributeTok{h =} \FloatTok{0.1}\NormalTok{, }\AttributeTok{lty =} \DecValTok{2}\NormalTok{)}
\FunctionTok{barplot}\NormalTok{(}\FunctionTok{table}\NormalTok{(refl[}\SpecialCharTok{{-}}\NormalTok{(}\DecValTok{1}\SpecialCharTok{:}\DecValTok{200}\NormalTok{)]) }\SpecialCharTok{/} \FunctionTok{length}\NormalTok{(refl[}\SpecialCharTok{{-}}\NormalTok{(}\DecValTok{1}\SpecialCharTok{:}\DecValTok{200}\NormalTok{)]), }\AttributeTok{col =} \StringTok{"firebrick"}\NormalTok{,}
        \AttributeTok{xlab =} \StringTok{"state"}\NormalTok{, }\AttributeTok{ylab =} \StringTok{"frequency"}\NormalTok{, }\AttributeTok{main =} \StringTok{"reflecting walk (bounces off ends)"}\NormalTok{)}
\end{Highlighting}
\end{Shaded}

\begin{figure}[H]

\centering{

\pandocbounded{\includegraphics[keepaspectratio]{chapters/unit11_mcmc_files/figure-pdf/fig-two-walks-1.pdf}}

}

\caption{\label{fig-two-walks}Long-run state frequencies for the cyclic
walk (left, converges to uniform on all ten states) and the reflecting
walk (right, concentrates almost entirely at state 9).}

\end{figure}%

The two chains use almost the same rule and yet settle into completely
different long-run behaviour, and the theory below explains why. In the
cyclic walk every state has exactly the same \emph{local} picture ---
two neighbours, reached with equal probability \(0.45\) in each
direction --- so no state is structurally different from any other, and
the equilibrium distribution treats them all alike: uniform on
\(\{0,\ldots,9\}\). The reflecting walk changes two things at once, and
both matter. First, its forward and backward probabilities are
themselves unequal (\(0.4\) versus \(0.2\), with \(0.4\) probability of
staying put), so even on a cycle this chain would drift persistently in
one direction --- though a \emph{cyclic} chain with that same asymmetry
would still converge to uniform, because every state would still be
structurally identical to every other. Second, and decisively, states
\(0\) and \(9\) are no longer structurally identical to the interior
states: an attempted move off either end is turned into a self-loop
instead of wrapping around, so probability that would have left the
segment is returned to the same state. Combined with the rightward
drift, this reflecting boundary acts as a trap at state \(9\)
specifically --- mass drifts right, and once it arrives it accumulates
there rather than passing through --- which is exactly the extreme,
almost-all-mass-at-\(9\) pattern in the right panel above. This is the
invariance and detailed-balance machinery below made visible: a
transition rule under which every state looks alike to the chain
preserves the uniform distribution, and one that treats some states
differently does not.

\subsection{Invariance, detailed balance, and
convergence}\label{invariance-detailed-balance-and-convergence}

\(\pi\) is an \textbf{invariant} (or stationary) distribution of
transition \(T\) if \[
\int \pi(\theta)\, T(\theta \to \theta') \, d\theta = \pi(\theta') \qquad \text{for all } \theta'.
\] In words: if \(\theta \sim \pi\) and one step of the chain is
applied, the result is again distributed as \(\pi\) --- \(\pi\) is a
fixed point of \(T\). Invariance alone does not guarantee the chain
\emph{converges} to \(\pi\) from an arbitrary start (a chain that never
moves leaves every distribution invariant), so two more conditions are
needed: \textbf{aperiodicity} (the chain does not cycle
deterministically through a fixed sequence of states) and
\textbf{irreducibility} (every state can eventually be reached from
every other). Under all three, \(\theta^{(i)} \to \pi\) in distribution
from any starting point, and ergodic averages converge,
\(\frac1n\sum_i a(\theta^{(i)}) \to E_\pi[a(\theta)]\).

Invariance is a \emph{global} condition (it involves integrating over
all \(\theta\)) and is awkward to check directly for a proposed \(T\).
The standard shortcut is a \emph{local}, pairwise condition. \(T\)
satisfies \textbf{detailed balance} with respect to \(\pi\) if \[
\pi(\theta)\, T(\theta \to \theta') = \pi(\theta')\, T(\theta' \to \theta) \qquad \text{for all } \theta, \theta',
\] i.e.~the flow of probability from \(\theta\) to \(\theta'\) under
\(\pi\) exactly balances the reverse flow (the transition \(T\) itself
need not be symmetric --- only the \(\pi\)-weighted flow is). Detailed
balance implies invariance: integrating both sides over \(\theta\), \[
\int \pi(\theta)\, T(\theta \to \theta') \, d\theta = \int \pi(\theta')\, T(\theta' \to \theta) \, d\theta
= \pi(\theta') \underbrace{\int T(\theta' \to \theta)\, d\theta}_{=\,1} = \pi(\theta').
\] Every construction in this chapter --- Gibbs, Metropolis--Hastings,
Hamiltonian Monte Carlo --- is a recipe for building a \(T\) that
satisfies detailed balance (or, for Gibbs, invariance directly) with
respect to a \emph{given} target \(\pi\), so that simulating the chain
long enough and discarding an initial \textbf{burn-in} period produces
(dependent) draws from \(\pi\).

\section{Gibbs Sampling}\label{gibbs-sampling}

\subsection{The idea and why it works}\label{the-idea-and-why-it-works}

Suppose \(\theta = (\theta_1, \theta_2)\) and both full conditionals
\(\pi(\theta_1 \mid \theta_2)\) and \(\pi(\theta_2 \mid \theta_1)\) can
be sampled directly, even though the joint \(\pi(\theta_1, \theta_2)\)
has no closed form or cannot be sampled directly. The \textbf{Gibbs
sampler} alternates \[
\theta_1^{(t)} \sim \pi\big(\theta_1 \mid \theta_2^{(t-1)}\big), \qquad
\theta_2^{(t)} \sim \pi\big(\theta_2 \mid \theta_1^{(t)}\big).
\] Each half-step leaves \(\pi\) invariant: if
\((\theta_1, \theta_2) \sim \pi\), then \(\theta_2 \sim \pi(\theta_2)\)
(its marginal), and replacing \(\theta_1\) by a fresh draw from
\(\pi(\theta_1 \mid \theta_2)\) produces a pair with joint density
\(\pi(\theta_2)\,\pi(\theta_1 \mid \theta_2) = \pi(\theta_1, \theta_2)\)
--- exactly \(\pi\) again. Cycling through both coordinates therefore
also leaves \(\pi\) invariant, and as long as the full conditionals are
everywhere positive the chain is irreducible and aperiodic. No proposal,
no acceptance step, no tuning: Gibbs sampling only needs full
conditionals that can be sampled.

\subsection{Deriving the full conditionals for normal
data}\label{deriving-the-full-conditionals-for-normal-data}

Take the running example for this half of the book:
\(x_1, \ldots, x_n \mid \mu, \sigma^2 \overset{iid}{\sim} N(\mu, \sigma^2)\),
with independent priors \(\mu \sim N(\mu_0, \sigma_0^2)\) and
\(\sigma^2 \sim \mathrm{IG}(a_0, \lambda_0)\) (inverse-gamma:
\(1/\sigma^2 \sim \mathrm{Gamma}(a_0, \lambda_0)\), so \(a_0\) is a
shape and \(\lambda_0\) a rate/scale). The joint posterior is \[
P(\mu, \sigma^2 \mid x) \ \propto\ e^{-(\mu-\mu_0)^2 / (2\sigma_0^2)} \cdot (\sigma^2)^{-(a_0+1)} e^{-\lambda_0/\sigma^2} \cdot (\sigma^2)^{-n/2} \exp\Big\{-\sum_{i=1}^n \tfrac{(x_i-\mu)^2}{2\sigma^2}\Big\},
\] two-dimensional with no closed-form normalizing constant, but each
conditional is standard.

\textbf{Conditional of \(\mu\).} Collect the factors that involve
\(\mu\) and expand
\(\sum_i(x_i-\mu)^2 = n\mu^2 - 2n\mu\bar x + \sum_i x_i^2\): \[
P(\mu \mid \sigma^2, x) \ \propto\ \exp\Big\{-\tfrac12\Big[\big(\tfrac{1}{\sigma_0^2} + \tfrac{n}{\sigma^2}\big)\mu^2 - 2\big(\tfrac{\mu_0}{\sigma_0^2} + \tfrac{n\bar x}{\sigma^2}\big)\mu\Big]\Big\},
\] which is the kernel of \(N(m, v)\) with
\(1/v = 1/\sigma_0^2 + n/\sigma^2\) and
\(m/v = \mu_0/\sigma_0^2 + n\bar x/\sigma^2\) --- a precision-weighted
average of the prior mean and the sample mean, with \(n/\sigma^2\) the
weight the data contribute. As \(n \to \infty\), \(m \to \bar x\) and
\(v \approx \sigma^2/n\): the prior is overwhelmed.

\textbf{Conditional of \(\sigma^2\).} Collect the factors involving
\(\sigma^2\): \[
P(\sigma^2 \mid \mu, x) \ \propto\ (\sigma^2)^{-(a_0+1)} e^{-\lambda_0/\sigma^2} \cdot (\sigma^2)^{-n/2} e^{-\sum_i(x_i-\mu)^2/(2\sigma^2)}
= (\sigma^2)^{-\left(\frac n2 + a_0 + 1\right)} \exp\Big\{-\tfrac{\lambda_0 + \frac12\sum_i(x_i-\mu)^2}{\sigma^2}\Big\},
\] again an inverse-gamma kernel,
\(\sigma^2 \mid \mu, x \sim \mathrm{IG}\big(\tfrac n2 + a_0,\ \tfrac12\sum_i(x_i-\mu)^2 + \lambda_0\big)\):
the prior contributes \(a_0\) ``pseudo-observations'' worth of shape and
\(\lambda_0\) worth of sum-of-squares, and the data add \(n\) and
\(\sum_i(x_i-\mu)^2\) to each.

\subsection{Implementation}\label{implementation-3}

\begin{Shaded}
\begin{Highlighting}[]
\NormalTok{gibbs\_norm }\OtherTok{\textless{}{-}} \ControlFlowTok{function}\NormalTok{(iters, x, mu0, sigma0, alpha, w, }\AttributeTok{mu\_init =} \DecValTok{0}\NormalTok{, }\AttributeTok{s2\_init =} \DecValTok{1}\NormalTok{) \{}
\NormalTok{    n }\OtherTok{\textless{}{-}} \FunctionTok{length}\NormalTok{(x); sumx }\OtherTok{\textless{}{-}} \FunctionTok{sum}\NormalTok{(x)}
\NormalTok{    out }\OtherTok{\textless{}{-}} \FunctionTok{matrix}\NormalTok{(}\ConstantTok{NA\_real\_}\NormalTok{, iters, }\DecValTok{2}\NormalTok{, }\AttributeTok{dimnames =} \FunctionTok{list}\NormalTok{(}\ConstantTok{NULL}\NormalTok{, }\FunctionTok{c}\NormalTok{(}\StringTok{"mu"}\NormalTok{, }\StringTok{"s2"}\NormalTok{)))}
\NormalTok{    mu }\OtherTok{\textless{}{-}}\NormalTok{ mu\_init; s2 }\OtherTok{\textless{}{-}}\NormalTok{ s2\_init}
    \ControlFlowTok{for}\NormalTok{ (t }\ControlFlowTok{in} \FunctionTok{seq\_len}\NormalTok{(iters)) \{}
\NormalTok{        v  }\OtherTok{\textless{}{-}} \DecValTok{1} \SpecialCharTok{/}\NormalTok{ (n }\SpecialCharTok{/}\NormalTok{ s2 }\SpecialCharTok{+} \DecValTok{1} \SpecialCharTok{/}\NormalTok{ sigma0)}
\NormalTok{        mu }\OtherTok{\textless{}{-}} \FunctionTok{rnorm}\NormalTok{(}\DecValTok{1}\NormalTok{, (sumx }\SpecialCharTok{/}\NormalTok{ s2 }\SpecialCharTok{+}\NormalTok{ mu0 }\SpecialCharTok{/}\NormalTok{ sigma0) }\SpecialCharTok{*}\NormalTok{ v, }\FunctionTok{sqrt}\NormalTok{(v))}
\NormalTok{        s2 }\OtherTok{\textless{}{-}} \DecValTok{1} \SpecialCharTok{/} \FunctionTok{rgamma}\NormalTok{(}\DecValTok{1}\NormalTok{, (alpha }\SpecialCharTok{+}\NormalTok{ n) }\SpecialCharTok{/} \DecValTok{2}\NormalTok{, (alpha }\SpecialCharTok{*}\NormalTok{ w }\SpecialCharTok{+} \FunctionTok{sum}\NormalTok{((x }\SpecialCharTok{{-}}\NormalTok{ mu)}\SpecialCharTok{\^{}}\DecValTok{2}\NormalTok{)) }\SpecialCharTok{/} \DecValTok{2}\NormalTok{)}
\NormalTok{        out[t, ] }\OtherTok{\textless{}{-}} \FunctionTok{c}\NormalTok{(mu, s2)}
\NormalTok{    \}}
\NormalTok{    out}
\NormalTok{\}}
\end{Highlighting}
\end{Shaded}

A preallocated matrix filled by a plain \texttt{for} loop replaces the
original \texttt{replicate()}-with-\texttt{\textless{}\textless{}-}
pattern: it does the same work with one direct assignment per iteration
instead of a closure mutating variables in its enclosing scope, and it
makes the recurrence --- each row built from the previous one ---
visually explicit.

\begin{Shaded}
\begin{Highlighting}[]
\FunctionTok{set.seed}\NormalTok{(}\DecValTok{2}\NormalTok{)}
\NormalTok{x }\OtherTok{\textless{}{-}} \FunctionTok{rnorm}\NormalTok{(}\DecValTok{50}\NormalTok{, }\DecValTok{10}\NormalTok{, }\DecValTok{2}\NormalTok{)}
\NormalTok{draws\_gibbs }\OtherTok{\textless{}{-}} \FunctionTok{gibbs\_norm}\NormalTok{(}\DecValTok{10000}\NormalTok{, x, }\AttributeTok{mu0 =} \DecValTok{0}\NormalTok{, }\AttributeTok{sigma0 =} \FloatTok{1e10}\NormalTok{, }\AttributeTok{alpha =} \FloatTok{1e{-}5}\NormalTok{, }\AttributeTok{w =} \FloatTok{1e{-}5}\NormalTok{)}
\NormalTok{burn }\OtherTok{\textless{}{-}} \DecValTok{1}\SpecialCharTok{:}\DecValTok{200}
\NormalTok{mu\_g  }\OtherTok{\textless{}{-}}\NormalTok{ draws\_gibbs[}\SpecialCharTok{{-}}\NormalTok{burn, }\StringTok{"mu"}\NormalTok{]}
\NormalTok{sd\_g  }\OtherTok{\textless{}{-}} \FunctionTok{sqrt}\NormalTok{(draws\_gibbs[}\SpecialCharTok{{-}}\NormalTok{burn, }\StringTok{"s2"}\NormalTok{])}

\FunctionTok{par}\NormalTok{(}\AttributeTok{mfrow =} \FunctionTok{c}\NormalTok{(}\DecValTok{2}\NormalTok{, }\DecValTok{2}\NormalTok{), }\AttributeTok{mar =} \FunctionTok{c}\NormalTok{(}\DecValTok{4}\NormalTok{, }\DecValTok{4}\NormalTok{, }\DecValTok{2}\NormalTok{, }\DecValTok{1}\NormalTok{))}
\FunctionTok{plot}\NormalTok{(mu\_g, }\AttributeTok{type =} \StringTok{"l"}\NormalTok{, }\AttributeTok{col =} \StringTok{"steelblue"}\NormalTok{, }\AttributeTok{xlab =} \StringTok{"iteration"}\NormalTok{, }\AttributeTok{ylab =} \FunctionTok{expression}\NormalTok{(mu), }\AttributeTok{main =} \StringTok{"trace of mu"}\NormalTok{)}
\FunctionTok{plot}\NormalTok{(}\FunctionTok{density}\NormalTok{(mu\_g), }\AttributeTok{col =} \StringTok{"steelblue"}\NormalTok{, }\AttributeTok{lwd =} \DecValTok{2}\NormalTok{, }\AttributeTok{main =} \StringTok{"posterior density of mu"}\NormalTok{, }\AttributeTok{xlab =} \FunctionTok{expression}\NormalTok{(mu))}
\FunctionTok{plot}\NormalTok{(sd\_g, }\AttributeTok{type =} \StringTok{"l"}\NormalTok{, }\AttributeTok{col =} \StringTok{"firebrick"}\NormalTok{, }\AttributeTok{xlab =} \StringTok{"iteration"}\NormalTok{, }\AttributeTok{ylab =} \FunctionTok{expression}\NormalTok{(sigma), }\AttributeTok{main =} \StringTok{"trace of sigma"}\NormalTok{)}
\FunctionTok{plot}\NormalTok{(}\FunctionTok{density}\NormalTok{(sd\_g), }\AttributeTok{col =} \StringTok{"firebrick"}\NormalTok{, }\AttributeTok{lwd =} \DecValTok{2}\NormalTok{, }\AttributeTok{main =} \StringTok{"posterior density of sigma"}\NormalTok{, }\AttributeTok{xlab =} \FunctionTok{expression}\NormalTok{(sigma))}
\end{Highlighting}
\end{Shaded}

\begin{figure}[H]

{\centering \pandocbounded{\includegraphics[keepaspectratio]{chapters/unit11_mcmc_files/figure-pdf/gibbs-demo-1.pdf}}

}

\caption{Gibbs sampler on 50 simulated N(10, 2\^{}2) observations: trace
plots and posterior densities for mu and sigma, with a vague prior.}

\end{figure}%

\begin{Shaded}
\begin{Highlighting}[]
\NormalTok{gibbs\_summary }\OtherTok{\textless{}{-}} \FunctionTok{data.frame}\NormalTok{(}
    \AttributeTok{parameter =} \FunctionTok{c}\NormalTok{(}\StringTok{"mu"}\NormalTok{, }\StringTok{"sigma"}\NormalTok{),}
    \AttributeTok{mean =} \FunctionTok{c}\NormalTok{(}\FunctionTok{mean}\NormalTok{(mu\_g), }\FunctionTok{mean}\NormalTok{(sd\_g)),}
    \AttributeTok{sd   =} \FunctionTok{c}\NormalTok{(}\FunctionTok{sd}\NormalTok{(mu\_g), }\FunctionTok{sd}\NormalTok{(sd\_g)),}
    \AttributeTok{q025 =} \FunctionTok{c}\NormalTok{(}\FunctionTok{quantile}\NormalTok{(mu\_g, }\FloatTok{0.025}\NormalTok{), }\FunctionTok{quantile}\NormalTok{(sd\_g, }\FloatTok{0.025}\NormalTok{)),}
    \AttributeTok{q975 =} \FunctionTok{c}\NormalTok{(}\FunctionTok{quantile}\NormalTok{(mu\_g, }\FloatTok{0.975}\NormalTok{), }\FunctionTok{quantile}\NormalTok{(sd\_g, }\FloatTok{0.975}\NormalTok{))}
\NormalTok{)}
\NormalTok{knitr}\SpecialCharTok{::}\FunctionTok{kable}\NormalTok{(gibbs\_summary, }\AttributeTok{digits =} \DecValTok{3}\NormalTok{,}
             \AttributeTok{col.names =} \FunctionTok{c}\NormalTok{(}\StringTok{"parameter"}\NormalTok{, }\StringTok{"mean"}\NormalTok{, }\StringTok{"sd"}\NormalTok{, }\StringTok{"2.5\%"}\NormalTok{, }\StringTok{"97.5\%"}\NormalTok{),}
             \AttributeTok{align =} \FunctionTok{c}\NormalTok{(}\StringTok{"l"}\NormalTok{, }\StringTok{"r"}\NormalTok{, }\StringTok{"r"}\NormalTok{, }\StringTok{"r"}\NormalTok{, }\StringTok{"r"}\NormalTok{))}
\end{Highlighting}
\end{Shaded}

{

\begin{longtable}[]{@{}lrrrr@{}}

\caption{\label{tbl-gibbs-summary}Posterior summaries from the Gibbs
sampler (after discarding a 200-iteration burn-in).}

\tabularnewline

\toprule\noalign{}
parameter & mean & sd & 2.5\% & 97.5\% \\
\midrule\noalign{}
\endhead
\bottomrule\noalign{}
\endlastfoot
mu & 10.141 & 0.330 & 9.489 & 10.791 \\
sigma & 2.293 & 0.239 & 1.885 & 2.809 \\

\end{longtable}

}

\subsection{Data augmentation: sampling with censored
data}\label{data-augmentation-sampling-with-censored-data}

Gibbs sampling extends immediately to models with \textbf{latent
variables}, by treating the latent quantities as extra parameters to
sample. Suppose the \(x_i\) are not observed exactly but only rounded
down to the nearest multiple of some interval length \(l\) (ages
recorded in whole years, measurements binned to the nearest instrument
tick): we observe \(y_i\) with \(x_i \in [y_i, y_i + l)\). \textbf{Data
augmentation} adds \(x_1,\ldots,x_n\) as unknowns and Gibbs-samples all
three blocks: \[
x_i^{(t)} \sim P\big(x_i \mid \mu^{(t-1)}, \sigma^{2(t-1)}, y_i\big), \qquad
\mu^{(t)} \sim P\big(\mu \mid x^{(t)}, \sigma^{2(t-1)}\big), \qquad
\sigma^{2(t)} \sim P\big(\sigma^2 \mid x^{(t)}, \mu^{(t)}\big).
\] The last two steps are exactly the conditionals derived above,
applied to the imputed \(x\). The first is a normal density restricted
to the interval \([y_i, y_i+l)\) and renormalized --- a
\textbf{truncated normal}, sampled directly with
\texttt{truncnorm::rtruncnorm} rather than by rejection:

\begin{Shaded}
\begin{Highlighting}[]
\NormalTok{gibbs\_norm\_interval }\OtherTok{\textless{}{-}} \ControlFlowTok{function}\NormalTok{(iters, y, mu0, sigma0, alpha, w, }\AttributeTok{l =} \DecValTok{1}\NormalTok{) \{}
\NormalTok{    n }\OtherTok{\textless{}{-}} \FunctionTok{length}\NormalTok{(y)}
\NormalTok{    mu }\OtherTok{\textless{}{-}} \DecValTok{0}\NormalTok{; s2 }\OtherTok{\textless{}{-}} \DecValTok{1}\NormalTok{; x }\OtherTok{\textless{}{-}} \FunctionTok{rep}\NormalTok{(}\DecValTok{0}\NormalTok{, n)}
\NormalTok{    out }\OtherTok{\textless{}{-}} \FunctionTok{matrix}\NormalTok{(}\ConstantTok{NA\_real\_}\NormalTok{, iters, }\DecValTok{2} \SpecialCharTok{+}\NormalTok{ n)}
    \ControlFlowTok{for}\NormalTok{ (t }\ControlFlowTok{in} \FunctionTok{seq\_len}\NormalTok{(iters)) \{}
\NormalTok{        x }\OtherTok{\textless{}{-}} \FunctionTok{rtruncnorm}\NormalTok{(n, }\AttributeTok{a =}\NormalTok{ y, }\AttributeTok{b =}\NormalTok{ y }\SpecialCharTok{+}\NormalTok{ l, }\AttributeTok{mean =}\NormalTok{ mu, }\AttributeTok{sd =} \FunctionTok{sqrt}\NormalTok{(s2))}
\NormalTok{        sumx }\OtherTok{\textless{}{-}} \FunctionTok{sum}\NormalTok{(x)}
\NormalTok{        v  }\OtherTok{\textless{}{-}} \DecValTok{1} \SpecialCharTok{/}\NormalTok{ (n }\SpecialCharTok{/}\NormalTok{ s2 }\SpecialCharTok{+} \DecValTok{1} \SpecialCharTok{/}\NormalTok{ sigma0)}
\NormalTok{        mu }\OtherTok{\textless{}{-}} \FunctionTok{rnorm}\NormalTok{(}\DecValTok{1}\NormalTok{, (sumx }\SpecialCharTok{/}\NormalTok{ s2 }\SpecialCharTok{+}\NormalTok{ mu0 }\SpecialCharTok{/}\NormalTok{ sigma0) }\SpecialCharTok{*}\NormalTok{ v, }\FunctionTok{sqrt}\NormalTok{(v))}
\NormalTok{        s2 }\OtherTok{\textless{}{-}} \DecValTok{1} \SpecialCharTok{/} \FunctionTok{rgamma}\NormalTok{(}\DecValTok{1}\NormalTok{, (alpha }\SpecialCharTok{+}\NormalTok{ n) }\SpecialCharTok{/} \DecValTok{2}\NormalTok{, (alpha }\SpecialCharTok{*}\NormalTok{ w }\SpecialCharTok{+} \FunctionTok{sum}\NormalTok{((x }\SpecialCharTok{{-}}\NormalTok{ mu)}\SpecialCharTok{\^{}}\DecValTok{2}\NormalTok{)) }\SpecialCharTok{/} \DecValTok{2}\NormalTok{)}
\NormalTok{        out[t, ] }\OtherTok{\textless{}{-}} \FunctionTok{c}\NormalTok{(mu, s2, x)}
\NormalTok{    \}}
    \FunctionTok{list}\NormalTok{(}\AttributeTok{mu =}\NormalTok{ out[, }\DecValTok{1}\NormalTok{], }\AttributeTok{sd =} \FunctionTok{sqrt}\NormalTok{(out[, }\DecValTok{2}\NormalTok{]), }\AttributeTok{x =}\NormalTok{ out[, }\SpecialCharTok{{-}}\NormalTok{(}\DecValTok{1}\SpecialCharTok{:}\DecValTok{2}\NormalTok{)])}
\NormalTok{\}}
\end{Highlighting}
\end{Shaded}

\begin{Shaded}
\begin{Highlighting}[]
\FunctionTok{set.seed}\NormalTok{(}\DecValTok{3}\NormalTok{)}
\NormalTok{x\_true }\OtherTok{\textless{}{-}} \FunctionTok{rnorm}\NormalTok{(}\DecValTok{200}\NormalTok{, }\DecValTok{100}\NormalTok{, }\DecValTok{40}\NormalTok{)}
\NormalTok{y\_obs  }\OtherTok{\textless{}{-}} \FunctionTok{floor}\NormalTok{(x\_true }\SpecialCharTok{/} \DecValTok{10}\NormalTok{) }\SpecialCharTok{*} \DecValTok{10}
\NormalTok{fit\_cens }\OtherTok{\textless{}{-}} \FunctionTok{gibbs\_norm\_interval}\NormalTok{(}\DecValTok{5000}\NormalTok{, y\_obs, }\AttributeTok{mu0 =} \DecValTok{0}\NormalTok{, }\AttributeTok{sigma0 =} \FloatTok{1e10}\NormalTok{, }\AttributeTok{alpha =} \FloatTok{1e{-}5}\NormalTok{, }\AttributeTok{w =} \FloatTok{1e{-}5}\NormalTok{, }\AttributeTok{l =} \DecValTok{10}\NormalTok{)}
\NormalTok{burn2 }\OtherTok{\textless{}{-}} \DecValTok{1}\SpecialCharTok{:}\DecValTok{200}
\NormalTok{knitr}\SpecialCharTok{::}\FunctionTok{kable}\NormalTok{(}\FunctionTok{data.frame}\NormalTok{(}
    \AttributeTok{quantity =} \FunctionTok{c}\NormalTok{(}\StringTok{"true mu (unknown to the model)"}\NormalTok{, }\StringTok{"true sigma"}\NormalTok{, }\StringTok{"posterior mean of mu"}\NormalTok{, }\StringTok{"posterior mean of sigma"}\NormalTok{),}
    \AttributeTok{value =} \FunctionTok{c}\NormalTok{(}\DecValTok{100}\NormalTok{, }\DecValTok{40}\NormalTok{, }\FunctionTok{mean}\NormalTok{(fit\_cens}\SpecialCharTok{$}\NormalTok{mu[}\SpecialCharTok{{-}}\NormalTok{burn2]), }\FunctionTok{mean}\NormalTok{(fit\_cens}\SpecialCharTok{$}\NormalTok{sd[}\SpecialCharTok{{-}}\NormalTok{burn2]))}
\NormalTok{), }\AttributeTok{digits =} \DecValTok{2}\NormalTok{, }\AttributeTok{col.names =} \FunctionTok{c}\NormalTok{(}\StringTok{"quantity"}\NormalTok{, }\StringTok{"value"}\NormalTok{), }\AttributeTok{align =} \FunctionTok{c}\NormalTok{(}\StringTok{"l"}\NormalTok{, }\StringTok{"r"}\NormalTok{))}
\end{Highlighting}
\end{Shaded}

{

\begin{longtable}[]{@{}lr@{}}

\caption{\label{tbl-gibbs-censored}Recovering mu and sigma from data
rounded down to the nearest 10, n = 200.}

\tabularnewline

\toprule\noalign{}
quantity & value \\
\midrule\noalign{}
\endhead
\bottomrule\noalign{}
\endlastfoot
true mu (unknown to the model) & 100.00 \\
true sigma & 40.00 \\
posterior mean of mu & 100.63 \\
posterior mean of sigma & 38.71 \\

\end{longtable}

}

The rounding destroys a great deal of information (each \(y_i\) only
says which decade \(x_i\) fell in), yet the augmented Gibbs sampler
still recovers \(\mu\) and \(\sigma\) close to their true values,
because it correctly propagates the \emph{uncertainty} in each \(x_i\)
rather than, say, treating \(y_i\) as if it were exact.

\section{Metropolis--Hastings}\label{metropolishastings}

\subsection{The algorithm}\label{the-algorithm-2}

Gibbs sampling needs full conditionals that are both known and directly
sample-able --- a strong requirement outside conjugate models.
\textbf{Metropolis--Hastings (MH)} needs only a target \(\pi(\theta)\)
computable up to a constant and \emph{any} proposal distribution
\(\tilde T(\theta^* \mid \theta)\) we can sample from. Starting from
\(\theta^{(0)}\), repeat:

\begin{enumerate}
\def\labelenumi{\arabic{enumi}.}
\tightlist
\item
  Draw a proposal \(\theta^* \sim \tilde T(\cdot \mid \theta^{(t-1)})\).
\item
  Compute the acceptance ratio
  \(r = \min\Big\{\dfrac{\pi(\theta^*)\,\tilde T(\theta^{(t-1)} \mid \theta^*)}{\pi(\theta^{(t-1)})\,\tilde T(\theta^* \mid \theta^{(t-1)})},\ 1\Big\}\).
\item
  With probability \(r\) set \(\theta^{(t)} = \theta^*\) (accept);
  otherwise set \(\theta^{(t)} = \theta^{(t-1)}\) (reject \textbf{and
  repeat the current value}).
\end{enumerate}

Only the ratio \(\pi(\theta^*)/\pi(\theta^{(t-1)})\) is needed, so any
normalizing constant of \(\pi\) cancels --- exactly as in rejection
sampling, but here a rejected proposal still contributes a draw (the
repeated current value) rather than being discarded, which is what lets
MH work even when \(\pi\) has no available envelope. One can show (by
checking detailed balance directly, splitting into the ``propose and
accept'' and ``reject and stay'' cases and using the
\(\min\{\cdot,\cdot\}\) identity \(\min\{a,b\} = \min\{b,a\}\)) that
this \(T\) leaves \(\pi\) invariant regardless of what \(\tilde T\) is,
as long as its support is compatible with \(\pi\)'s.

\subsection{Random-walk Metropolis}\label{random-walk-metropolis}

If \(\tilde T\) is symmetric,
\(\tilde T(\theta^* \mid \theta) = \tilde T(\theta \mid \theta^*)\) ---
e.g.~\(\theta^* \sim N(\theta^{(t-1)}, \varsigma^2)\) --- the proposal
terms cancel and \(r = \min\{\pi(\theta^*)/\pi(\theta^{(t-1)}), 1\}\):
uphill proposals are always accepted, downhill ones accepted in
proportion to how much worse they are. \(\varsigma\), the \textbf{step
size}, is the only tuning parameter.

\begin{Shaded}
\begin{Highlighting}[]
\NormalTok{met\_gauss }\OtherTok{\textless{}{-}} \ControlFlowTok{function}\NormalTok{(iters, log\_f, stepsizes, ini\_value, ...) \{}
\NormalTok{    state }\OtherTok{\textless{}{-}}\NormalTok{ ini\_value; no\_var }\OtherTok{\textless{}{-}} \FunctionTok{length}\NormalTok{(state)}
\NormalTok{    logf }\OtherTok{\textless{}{-}} \FunctionTok{log\_f}\NormalTok{(ini\_value, ...)}
    \ControlFlowTok{if}\NormalTok{ (}\SpecialCharTok{!}\FunctionTok{is.finite}\NormalTok{(logf)) }\FunctionTok{stop}\NormalTok{(}\StringTok{"Initial value has zero probability"}\NormalTok{)}
\NormalTok{    out }\OtherTok{\textless{}{-}} \FunctionTok{matrix}\NormalTok{(}\ConstantTok{NA\_real\_}\NormalTok{, iters, no\_var); n\_acc }\OtherTok{\textless{}{-}} \DecValTok{0}
    \ControlFlowTok{for}\NormalTok{ (t }\ControlFlowTok{in} \FunctionTok{seq\_len}\NormalTok{(iters)) \{}
\NormalTok{        proposal }\OtherTok{\textless{}{-}} \FunctionTok{rnorm}\NormalTok{(no\_var, state, stepsizes)}
\NormalTok{        new\_logf }\OtherTok{\textless{}{-}} \FunctionTok{log\_f}\NormalTok{(proposal, ...)}
        \ControlFlowTok{if}\NormalTok{ (}\FunctionTok{log}\NormalTok{(}\FunctionTok{runif}\NormalTok{(}\DecValTok{1}\NormalTok{)) }\SpecialCharTok{\textless{}}\NormalTok{ new\_logf }\SpecialCharTok{{-}}\NormalTok{ logf) \{}
\NormalTok{            state }\OtherTok{\textless{}{-}}\NormalTok{ proposal; logf }\OtherTok{\textless{}{-}}\NormalTok{ new\_logf; n\_acc }\OtherTok{\textless{}{-}}\NormalTok{ n\_acc }\SpecialCharTok{+} \DecValTok{1}
\NormalTok{        \}}
\NormalTok{        out[t, ] }\OtherTok{\textless{}{-}}\NormalTok{ state}
\NormalTok{    \}}
    \FunctionTok{attr}\NormalTok{(out, }\StringTok{"accept\_rate"}\NormalTok{) }\OtherTok{\textless{}{-}}\NormalTok{ n\_acc }\SpecialCharTok{/}\NormalTok{ iters}
\NormalTok{    out}
\NormalTok{\}}
\end{Highlighting}
\end{Shaded}

An indicator-shaped target needs no special handling: sampling a normal
truncated to \([l, u]\) by targeting
\(\pi(\theta) \propto \varphi(\theta)\,I(l \le \theta \le u)\), any
proposal landing outside \([l,u]\) has \(\pi(\theta^*) = 0\), hence
\(r = 0\), and is automatically rejected.

\begin{Shaded}
\begin{Highlighting}[]
\NormalTok{log\_tnorm }\OtherTok{\textless{}{-}} \ControlFlowTok{function}\NormalTok{(x, mu, sigma, lb, ub) \{}
    \ControlFlowTok{if}\NormalTok{ (x }\SpecialCharTok{\textgreater{}}\NormalTok{ lb }\SpecialCharTok{\&\&}\NormalTok{ x }\SpecialCharTok{\textless{}}\NormalTok{ ub) }\FunctionTok{dnorm}\NormalTok{(x, mu, sigma, }\AttributeTok{log =} \ConstantTok{TRUE}\NormalTok{) }\ControlFlowTok{else} \SpecialCharTok{{-}}\ConstantTok{Inf}
\NormalTok{\}}
\FunctionTok{par}\NormalTok{(}\AttributeTok{mfrow =} \FunctionTok{c}\NormalTok{(}\DecValTok{1}\NormalTok{, }\DecValTok{3}\NormalTok{), }\AttributeTok{mar =} \FunctionTok{c}\NormalTok{(}\DecValTok{4}\NormalTok{, }\DecValTok{4}\NormalTok{, }\DecValTok{2}\NormalTok{, }\DecValTok{1}\NormalTok{))}
\ControlFlowTok{for}\NormalTok{ (step }\ControlFlowTok{in} \FunctionTok{c}\NormalTok{(}\FloatTok{0.1}\NormalTok{, }\DecValTok{1}\NormalTok{, }\DecValTok{2}\NormalTok{)) \{}
\NormalTok{    mc }\OtherTok{\textless{}{-}} \FunctionTok{met\_gauss}\NormalTok{(}\DecValTok{5000}\NormalTok{, log\_tnorm, }\AttributeTok{stepsizes =}\NormalTok{ step, }\AttributeTok{ini\_value =} \FloatTok{4.5}\NormalTok{, }\AttributeTok{mu =} \DecValTok{0}\NormalTok{, }\AttributeTok{sigma =} \DecValTok{2}\NormalTok{, }\AttributeTok{lb =} \DecValTok{4}\NormalTok{, }\AttributeTok{ub =} \DecValTok{5}\NormalTok{)}
    \FunctionTok{plot}\NormalTok{(mc, }\AttributeTok{type =} \StringTok{"l"}\NormalTok{, }\AttributeTok{col =} \StringTok{"steelblue"}\NormalTok{, }\AttributeTok{xlab =} \StringTok{"iteration"}\NormalTok{, }\AttributeTok{ylab =} \FunctionTok{expression}\NormalTok{(theta),}
         \AttributeTok{main =} \FunctionTok{sprintf}\NormalTok{(}\StringTok{"step = \%.1f, accept = \%.2f"}\NormalTok{, step, }\FunctionTok{attr}\NormalTok{(mc, }\StringTok{"accept\_rate"}\NormalTok{)))}
\NormalTok{\}}
\end{Highlighting}
\end{Shaded}

\begin{figure}[H]

\centering{

\pandocbounded{\includegraphics[keepaspectratio]{chapters/unit11_mcmc_files/figure-pdf/fig-tnorm-mh-1.pdf}}

}

\caption{\label{fig-tnorm-mh}Random-walk MH targeting N(0, 2\^{}2)
truncated to {[}4, 5{]}, for three step sizes.}

\end{figure}%

\subsection{Choosing the step size}\label{choosing-the-step-size}

The chain's draws are dependent, and what determines the precision of an
ergodic average is not the acceptance rate itself but the resulting
\textbf{autocorrelation}: for a stationary chain with lag-\(k\)
autocorrelation \(\rho_k\), \[
\mathrm{Var}\Big(\tfrac1T\sum_{t=1}^T \theta^{(t)}\Big) \approx \frac{\sigma^2}{T}\Big(1 + 2\sum_{k \ge 1} \rho_k\Big) = \frac{\sigma^2}{T_\text{eff}},
\] so the tuning goal is small autocorrelation, not a high acceptance
rate as such. A step too small is accepted almost every time but crawls
(highly autocorrelated); a step too large is almost always rejected, so
the chain barely moves either. Both extremes waste the same \(T\)
iterations on far fewer effective ones --- the figure below (three step
sizes on a \(N(0,1)\) target) makes this visible directly in the trace
and the ACF.

\begin{Shaded}
\begin{Highlighting}[]
\NormalTok{mh\_std\_normal }\OtherTok{\textless{}{-}} \ControlFlowTok{function}\NormalTok{(sig, }\AttributeTok{n =} \DecValTok{2000}\NormalTok{) \{}
\NormalTok{    th }\OtherTok{\textless{}{-}} \FunctionTok{numeric}\NormalTok{(n); th[}\DecValTok{1}\NormalTok{] }\OtherTok{\textless{}{-}} \DecValTok{0}\NormalTok{; acc }\OtherTok{\textless{}{-}} \DecValTok{0}
    \ControlFlowTok{for}\NormalTok{ (t }\ControlFlowTok{in} \DecValTok{2}\SpecialCharTok{:}\NormalTok{n) \{}
\NormalTok{        p }\OtherTok{\textless{}{-}}\NormalTok{ th[t }\SpecialCharTok{{-}} \DecValTok{1}\NormalTok{] }\SpecialCharTok{+} \FunctionTok{rnorm}\NormalTok{(}\DecValTok{1}\NormalTok{, }\DecValTok{0}\NormalTok{, sig)}
        \ControlFlowTok{if}\NormalTok{ (}\FunctionTok{runif}\NormalTok{(}\DecValTok{1}\NormalTok{) }\SpecialCharTok{\textless{}} \FunctionTok{dnorm}\NormalTok{(p) }\SpecialCharTok{/} \FunctionTok{dnorm}\NormalTok{(th[t }\SpecialCharTok{{-}} \DecValTok{1}\NormalTok{])) \{ th[t] }\OtherTok{\textless{}{-}}\NormalTok{ p; acc }\OtherTok{\textless{}{-}}\NormalTok{ acc }\SpecialCharTok{+} \DecValTok{1}\NormalTok{ \} }\ControlFlowTok{else}\NormalTok{ th[t] }\OtherTok{\textless{}{-}}\NormalTok{ th[t }\SpecialCharTok{{-}} \DecValTok{1}\NormalTok{]}
\NormalTok{    \}}
    \FunctionTok{list}\NormalTok{(}\AttributeTok{th =}\NormalTok{ th, }\AttributeTok{acc =}\NormalTok{ acc }\SpecialCharTok{/}\NormalTok{ (n }\SpecialCharTok{{-}} \DecValTok{1}\NormalTok{))}
\NormalTok{\}}
\NormalTok{res }\OtherTok{\textless{}{-}} \FunctionTok{lapply}\NormalTok{(}\FunctionTok{c}\NormalTok{(}\FloatTok{0.1}\NormalTok{, }\FloatTok{2.4}\NormalTok{, }\DecValTok{40}\NormalTok{), mh\_std\_normal)}
\NormalTok{labs }\OtherTok{\textless{}{-}} \FunctionTok{c}\NormalTok{(}\StringTok{"too small"}\NormalTok{, }\StringTok{"about right"}\NormalTok{, }\StringTok{"too large"}\NormalTok{)}
\FunctionTok{par}\NormalTok{(}\AttributeTok{mfrow =} \FunctionTok{c}\NormalTok{(}\DecValTok{2}\NormalTok{, }\DecValTok{3}\NormalTok{), }\AttributeTok{mar =} \FunctionTok{c}\NormalTok{(}\DecValTok{4}\NormalTok{, }\DecValTok{4}\NormalTok{, }\DecValTok{2}\NormalTok{, }\DecValTok{1}\NormalTok{))}
\ControlFlowTok{for}\NormalTok{ (k }\ControlFlowTok{in} \DecValTok{1}\SpecialCharTok{:}\DecValTok{3}\NormalTok{) }\FunctionTok{plot}\NormalTok{(res[[k]]}\SpecialCharTok{$}\NormalTok{th[}\DecValTok{1}\SpecialCharTok{:}\DecValTok{500}\NormalTok{], }\AttributeTok{type =} \StringTok{"l"}\NormalTok{, }\AttributeTok{col =} \StringTok{"steelblue"}\NormalTok{, }\AttributeTok{xlab =} \StringTok{"t"}\NormalTok{, }\AttributeTok{ylab =} \FunctionTok{expression}\NormalTok{(theta}\SpecialCharTok{\^{}}\NormalTok{(t)),}
                     \AttributeTok{ylim =} \FunctionTok{c}\NormalTok{(}\SpecialCharTok{{-}}\FloatTok{3.5}\NormalTok{, }\FloatTok{3.5}\NormalTok{), }\AttributeTok{main =} \FunctionTok{sprintf}\NormalTok{(}\StringTok{"\%s: accept \%.2f"}\NormalTok{, labs[k], res[[k]]}\SpecialCharTok{$}\NormalTok{acc))}
\ControlFlowTok{for}\NormalTok{ (k }\ControlFlowTok{in} \DecValTok{1}\SpecialCharTok{:}\DecValTok{3}\NormalTok{) }\FunctionTok{acf}\NormalTok{(res[[k]]}\SpecialCharTok{$}\NormalTok{th, }\AttributeTok{lag.max =} \DecValTok{50}\NormalTok{, }\AttributeTok{main =} \StringTok{"ACF"}\NormalTok{)}
\end{Highlighting}
\end{Shaded}

\begin{figure}[H]

\centering{

\pandocbounded{\includegraphics[keepaspectratio]{chapters/unit11_mcmc_files/figure-pdf/fig-stepsize-1.pdf}}

}

\caption{\label{fig-stepsize}Step size too small, about right, and too
large: trace (top) and autocorrelation (bottom) for a random walk
targeting N(0, 1).}

\end{figure}%

For random-walk Metropolis on a smooth target, theory recommends an
acceptance rate near \(0.234\) in high dimension (up to about \(0.44\)
in one dimension) as the rate that minimizes autocorrelation for a given
computational budget --- a useful rule of thumb, not a target to hit
exactly.

\subsection{Example: the same normal-data model, by
Metropolis--Hastings}\label{example-the-same-normal-data-model-by-metropolishastings}

\(\sigma^2 > 0\) makes a symmetric normal random walk directly on
\(\sigma^2\) wasteful (many proposals land below zero), so reparametrize
\(w = \log \sigma^2\), which ranges over all of \(\mathbb R\). By the
change-of-variables formula, the posterior in \((\mu, w)\) picks up a
Jacobian factor \(e^w\) relative to the posterior in
\((\mu, \sigma^2)\); omitting it targets the wrong distribution.

\begin{Shaded}
\begin{Highlighting}[]
\NormalTok{log\_post\_mu\_w }\OtherTok{\textless{}{-}} \ControlFlowTok{function}\NormalTok{(mu, w, x, mu0, sigma0, alpha, wprior) \{}
\NormalTok{    s2 }\OtherTok{\textless{}{-}} \FunctionTok{exp}\NormalTok{(w)}
    \FunctionTok{sum}\NormalTok{(}\FunctionTok{dnorm}\NormalTok{(x, mu, }\FunctionTok{sqrt}\NormalTok{(s2), }\AttributeTok{log =} \ConstantTok{TRUE}\NormalTok{)) }\SpecialCharTok{+} \FunctionTok{dnorm}\NormalTok{(mu, mu0, }\FunctionTok{sqrt}\NormalTok{(sigma0), }\AttributeTok{log =} \ConstantTok{TRUE}\NormalTok{) }\SpecialCharTok{+}
\NormalTok{        (}\SpecialCharTok{{-}}\NormalTok{(alpha }\SpecialCharTok{/} \DecValTok{2} \SpecialCharTok{+} \DecValTok{1}\NormalTok{) }\SpecialCharTok{*} \FunctionTok{log}\NormalTok{(s2) }\SpecialCharTok{{-}}\NormalTok{ alpha }\SpecialCharTok{*}\NormalTok{ wprior }\SpecialCharTok{/}\NormalTok{ (}\DecValTok{2} \SpecialCharTok{*}\NormalTok{ s2)) }\SpecialCharTok{+}\NormalTok{ w   }\CommentTok{\# "+ w" is the log Jacobian}
\NormalTok{\}}
\NormalTok{mh\_norm }\OtherTok{\textless{}{-}} \ControlFlowTok{function}\NormalTok{(iters, stepsizes, x, mu0, sigma0, alpha, w) \{}
\NormalTok{    log\_f }\OtherTok{\textless{}{-}} \ControlFlowTok{function}\NormalTok{(state) }\FunctionTok{log\_post\_mu\_w}\NormalTok{(state[}\DecValTok{1}\NormalTok{], state[}\DecValTok{2}\NormalTok{], x, mu0, sigma0, alpha, w)}
\NormalTok{    mc }\OtherTok{\textless{}{-}} \FunctionTok{met\_gauss}\NormalTok{(iters, }\ControlFlowTok{function}\NormalTok{(state) }\FunctionTok{log\_f}\NormalTok{(state), stepsizes, }\AttributeTok{ini\_value =} \FunctionTok{c}\NormalTok{(}\DecValTok{0}\NormalTok{, }\DecValTok{0}\NormalTok{))}
    \FunctionTok{list}\NormalTok{(}\AttributeTok{mu =}\NormalTok{ mc[, }\DecValTok{1}\NormalTok{], }\AttributeTok{sd =} \FunctionTok{exp}\NormalTok{(}\FloatTok{0.5} \SpecialCharTok{*}\NormalTok{ mc[, }\DecValTok{2}\NormalTok{]), }\AttributeTok{accept\_rate =} \FunctionTok{attr}\NormalTok{(mc, }\StringTok{"accept\_rate"}\NormalTok{))}
\NormalTok{\}}
\end{Highlighting}
\end{Shaded}

\begin{Shaded}
\begin{Highlighting}[]
\NormalTok{scales }\OtherTok{\textless{}{-}} \FunctionTok{c}\NormalTok{(}\DecValTok{20}\NormalTok{, }\DecValTok{10}\NormalTok{, }\DecValTok{4}\NormalTok{, }\DecValTok{2}\NormalTok{)}
\NormalTok{fits }\OtherTok{\textless{}{-}} \FunctionTok{lapply}\NormalTok{(scales, }\ControlFlowTok{function}\NormalTok{(s) }\FunctionTok{mh\_norm}\NormalTok{(}\DecValTok{10000}\NormalTok{, }\FunctionTok{c}\NormalTok{(s, s) }\SpecialCharTok{/} \FunctionTok{sqrt}\NormalTok{(}\FunctionTok{length}\NormalTok{(x)), x, }\DecValTok{0}\NormalTok{, }\FloatTok{1e10}\NormalTok{, }\FloatTok{1e{-}5}\NormalTok{, }\FloatTok{1e{-}5}\NormalTok{))}
\FunctionTok{par}\NormalTok{(}\AttributeTok{mfcol =} \FunctionTok{c}\NormalTok{(}\DecValTok{2}\NormalTok{, }\DecValTok{4}\NormalTok{), }\AttributeTok{mar =} \FunctionTok{c}\NormalTok{(}\DecValTok{4}\NormalTok{, }\DecValTok{4}\NormalTok{, }\DecValTok{2}\NormalTok{, }\DecValTok{1}\NormalTok{))}
\ControlFlowTok{for}\NormalTok{ (k }\ControlFlowTok{in} \FunctionTok{seq\_along}\NormalTok{(fits)) \{}
    \FunctionTok{plot}\NormalTok{(fits[[k]]}\SpecialCharTok{$}\NormalTok{mu[}\DecValTok{1}\SpecialCharTok{:}\DecValTok{1000}\NormalTok{], }\AttributeTok{type =} \StringTok{"l"}\NormalTok{, }\AttributeTok{col =} \StringTok{"steelblue"}\NormalTok{, }\AttributeTok{xlab =} \StringTok{"t"}\NormalTok{, }\AttributeTok{ylab =} \FunctionTok{expression}\NormalTok{(mu),}
         \AttributeTok{main =} \FunctionTok{sprintf}\NormalTok{(}\StringTok{"scale \%g, accept \%.2f"}\NormalTok{, scales[k], fits[[k]]}\SpecialCharTok{$}\NormalTok{accept\_rate))}
    \FunctionTok{acf}\NormalTok{(fits[[k]]}\SpecialCharTok{$}\NormalTok{mu[}\SpecialCharTok{{-}}\NormalTok{(}\DecValTok{1}\SpecialCharTok{:}\DecValTok{500}\NormalTok{)], }\AttributeTok{main =} \StringTok{"ACF"}\NormalTok{)}
\NormalTok{\}}
\end{Highlighting}
\end{Shaded}

\begin{figure}[H]

\centering{

\pandocbounded{\includegraphics[keepaspectratio]{chapters/unit11_mcmc_files/figure-pdf/fig-mh-stepsize-search-1.pdf}}

}

\caption{\label{fig-mh-stepsize-search}Preliminary MH runs at four
step-size scales, targeting mu (top: trace, bottom: ACF after discarding
500 iterations).}

\end{figure}%

Scale \(4\) gives the least autocorrelated trace among the four,
consistent with an acceptance rate closer to the \(0.2\)--\(0.4\) range
than the very low or very high rates at the extremes. Running a long
chain at that scale and comparing to the Gibbs sampler on the same data:

\begin{Shaded}
\begin{Highlighting}[]
\NormalTok{mh\_long }\OtherTok{\textless{}{-}} \FunctionTok{mh\_norm}\NormalTok{(}\DecValTok{50000}\NormalTok{, }\FunctionTok{c}\NormalTok{(}\DecValTok{4}\NormalTok{, }\DecValTok{4}\NormalTok{) }\SpecialCharTok{/} \FunctionTok{sqrt}\NormalTok{(}\FunctionTok{length}\NormalTok{(x)), x, }\DecValTok{0}\NormalTok{, }\FloatTok{1e10}\NormalTok{, }\FloatTok{1e{-}5}\NormalTok{, }\FloatTok{1e{-}5}\NormalTok{)}
\NormalTok{keep\_mh }\OtherTok{\textless{}{-}} \SpecialCharTok{{-}}\NormalTok{(}\DecValTok{1}\SpecialCharTok{:}\DecValTok{1000}\NormalTok{)}
\NormalTok{coherence }\OtherTok{\textless{}{-}} \FunctionTok{data.frame}\NormalTok{(}
    \AttributeTok{method =} \FunctionTok{c}\NormalTok{(}\StringTok{"Metropolis{-}Hastings"}\NormalTok{, }\StringTok{"Gibbs"}\NormalTok{),}
    \AttributeTok{mean\_mu =} \FunctionTok{c}\NormalTok{(}\FunctionTok{mean}\NormalTok{(mh\_long}\SpecialCharTok{$}\NormalTok{mu[keep\_mh]), }\FunctionTok{mean}\NormalTok{(mu\_g)),}
    \AttributeTok{sd\_mu   =} \FunctionTok{c}\NormalTok{(}\FunctionTok{sd}\NormalTok{(mh\_long}\SpecialCharTok{$}\NormalTok{mu[keep\_mh]), }\FunctionTok{sd}\NormalTok{(mu\_g)),}
    \AttributeTok{mean\_sigma =} \FunctionTok{c}\NormalTok{(}\FunctionTok{mean}\NormalTok{(mh\_long}\SpecialCharTok{$}\NormalTok{sd[keep\_mh]), }\FunctionTok{mean}\NormalTok{(sd\_g)),}
    \AttributeTok{sd\_sigma   =} \FunctionTok{c}\NormalTok{(}\FunctionTok{sd}\NormalTok{(mh\_long}\SpecialCharTok{$}\NormalTok{sd[keep\_mh]), }\FunctionTok{sd}\NormalTok{(sd\_g))}
\NormalTok{)}
\NormalTok{knitr}\SpecialCharTok{::}\FunctionTok{kable}\NormalTok{(coherence, }\AttributeTok{digits =} \DecValTok{3}\NormalTok{, }\AttributeTok{align =} \FunctionTok{c}\NormalTok{(}\StringTok{"l"}\NormalTok{, }\StringTok{"r"}\NormalTok{, }\StringTok{"r"}\NormalTok{, }\StringTok{"r"}\NormalTok{, }\StringTok{"r"}\NormalTok{))}
\end{Highlighting}
\end{Shaded}

{

\begin{longtable}[]{@{}lrrrr@{}}

\caption{\label{tbl-mh-vs-gibbs}Posterior mean and sd from
Metropolis-Hastings versus Gibbs sampling, same data and priors: the two
exact samplers of the same posterior agree up to Monte Carlo error.}

\tabularnewline

\toprule\noalign{}
method & mean\_mu & sd\_mu & mean\_sigma & sd\_sigma \\
\midrule\noalign{}
\endhead
\bottomrule\noalign{}
\endlastfoot
Metropolis-Hastings & 10.144 & 0.33 & 2.294 & 0.237 \\
Gibbs & 10.141 & 0.33 & 2.293 & 0.239 \\

\end{longtable}

}

\subsection{Metropolis-within-Gibbs}\label{metropolis-within-gibbs}

The two ideas combine directly. If \(\theta = (\theta_1, \theta_2)\) and
\(\pi(\theta_1 \mid \theta_2)\) has no convenient sampler while
\(\pi(\theta_2 \mid \theta_1)\) does, replace only the first Gibbs step
by a Metropolis--Hastings update that targets
\(\pi(\theta_1 \mid \theta_2)\) (evaluated at the \emph{current}
\(\theta_2\)) and leave the second step as an exact conditional draw.
Each block transition leaves \(\pi\) invariant on its own --- the Gibbs
argument applies verbatim to any
\(\pi(\theta_j \mid \theta_{-j})\)-invariant transition, not only an
exact conditional draw --- so the composite sweep leaves the full joint
\(\pi\) invariant. This ``Gibbs where conjugate, Metropolis elsewhere''
pattern, iterated one parameter or block at a time, is how most
general-purpose Bayesian software before Hamiltonian methods actually
worked.

\section{Hamiltonian Monte Carlo}\label{hamiltonian-monte-carlo}

\subsection{Why random walks struggle}\label{why-random-walks-struggle}

A random-walk proposal has no sense of direction: reaching a point one
posterior standard deviation away takes \(O(d)\) steps in \(d\)
dimensions, and far more along a narrow, correlated ridge, since most
proposed directions immediately step off the ridge and get rejected.
Gibbs suffers the same way when parameters are strongly correlated,
since each step moves along only one coordinate axis.
\textbf{Hamiltonian Monte Carlo (HMC)} uses the \emph{gradient}
\(\nabla \log \pi(\theta)\) to propose points far away with high
probability of acceptance.

Introduce an auxiliary \textbf{momentum} variable \(p \in \mathbb R^d\)
and the joint density \(\pi(\theta, p) \propto \exp\{-H(\theta,p)\}\)
with \[
H(\theta, p) = \underbrace{-\log \pi(\theta)}_{\text{potential } U(\theta)} + \underbrace{\tfrac12 p^\top M^{-1} p}_{\text{kinetic } K(p)},
\] so that the \(\theta\)-marginal is exactly the target and
\(p \sim N_d(0, M)\) independently. Physically: a particle at position
\(\theta\) with momentum \(p\) rolling frictionlessly on the potential
surface \(U(\theta)\). Hamilton's equations move \((\theta,p)\) so as to
\emph{conserve} \(H\) exactly and are volume-preserving and reversible;
discretizing them with the \textbf{leapfrog} integrator (step size
\(\epsilon\), \(L\) steps), \[
p \leftarrow p + \tfrac\epsilon2 \nabla\log\pi(\theta), \qquad \theta \leftarrow \theta + \epsilon M^{-1}p, \qquad p \leftarrow p + \tfrac\epsilon2\nabla\log\pi(\theta),
\] keeps volume-preservation and reversibility exactly (negating \(p\)
and re-running gives back the start) while introducing only an
\(O(\epsilon^2)\) error in \(H\), which a Metropolis accept/reject step
then corrects exactly.

\subsection{The algorithm}\label{the-algorithm-3}

Starting from \(\theta^{(0)}\), repeat: draw a fresh
\(p \sim N_d(0,M)\); run \(L\) leapfrog steps from
\((\theta^{(t-1)}, p)\) to reach \((\theta^*, p^*)\); accept
\(\theta^{(t)} = \theta^*\) with probability
\(\min\{\exp[H(\theta^{(t-1)},p) - H(\theta^*,p^*)], 1\}\), else repeat
\(\theta^{(t-1)}\). Drawing \(p\) is a Gibbs step from its own
conditional (it doesn't depend on \(\theta\)), and the
leapfrog-plus-accept step is a Metropolis update with a deterministic,
reversible proposal --- so each half leaves \(\pi(\theta,p)\) invariant,
and because \(H\) is nearly conserved, acceptance stays high (typically
\(0.65\)--\(0.8\)) even though \(\theta^*\) may be far from
\(\theta^{(t-1)}\).

\begin{Shaded}
\begin{Highlighting}[]
\NormalTok{hmc }\OtherTok{\textless{}{-}} \ControlFlowTok{function}\NormalTok{(log\_pi, grad\_log\_pi, theta, eps, L, iters, }\AttributeTok{M =} \FunctionTok{diag}\NormalTok{(}\FunctionTok{length}\NormalTok{(theta))) \{}
\NormalTok{    d }\OtherTok{\textless{}{-}} \FunctionTok{length}\NormalTok{(theta); Minv }\OtherTok{\textless{}{-}} \FunctionTok{solve}\NormalTok{(M); R }\OtherTok{\textless{}{-}} \FunctionTok{chol}\NormalTok{(M)}
\NormalTok{    out }\OtherTok{\textless{}{-}} \FunctionTok{matrix}\NormalTok{(}\ConstantTok{NA\_real\_}\NormalTok{, iters, d); n\_acc }\OtherTok{\textless{}{-}} \DecValTok{0}
    \ControlFlowTok{for}\NormalTok{ (t }\ControlFlowTok{in} \FunctionTok{seq\_len}\NormalTok{(iters)) \{}
\NormalTok{        p0 }\OtherTok{\textless{}{-}} \FunctionTok{as.vector}\NormalTok{(}\FunctionTok{t}\NormalTok{(R) }\SpecialCharTok{\%*\%} \FunctionTok{rnorm}\NormalTok{(d))}
\NormalTok{        th }\OtherTok{\textless{}{-}}\NormalTok{ theta; p }\OtherTok{\textless{}{-}}\NormalTok{ p0 }\SpecialCharTok{+}\NormalTok{ eps }\SpecialCharTok{/} \DecValTok{2} \SpecialCharTok{*} \FunctionTok{grad\_log\_pi}\NormalTok{(th)}
        \ControlFlowTok{for}\NormalTok{ (l }\ControlFlowTok{in} \FunctionTok{seq\_len}\NormalTok{(L)) \{}
\NormalTok{            th }\OtherTok{\textless{}{-}}\NormalTok{ th }\SpecialCharTok{+}\NormalTok{ eps }\SpecialCharTok{*} \FunctionTok{as.vector}\NormalTok{(Minv }\SpecialCharTok{\%*\%}\NormalTok{ p)}
            \ControlFlowTok{if}\NormalTok{ (l }\SpecialCharTok{\textless{}}\NormalTok{ L) p }\OtherTok{\textless{}{-}}\NormalTok{ p }\SpecialCharTok{+}\NormalTok{ eps }\SpecialCharTok{*} \FunctionTok{grad\_log\_pi}\NormalTok{(th)}
\NormalTok{        \}}
\NormalTok{        p }\OtherTok{\textless{}{-}}\NormalTok{ p }\SpecialCharTok{+}\NormalTok{ eps }\SpecialCharTok{/} \DecValTok{2} \SpecialCharTok{*} \FunctionTok{grad\_log\_pi}\NormalTok{(th)}
\NormalTok{        H0 }\OtherTok{\textless{}{-}} \SpecialCharTok{{-}}\FunctionTok{log\_pi}\NormalTok{(theta) }\SpecialCharTok{+} \FloatTok{0.5} \SpecialCharTok{*} \FunctionTok{sum}\NormalTok{(p0 }\SpecialCharTok{*}\NormalTok{ (Minv }\SpecialCharTok{\%*\%}\NormalTok{ p0))}
\NormalTok{        H1 }\OtherTok{\textless{}{-}} \SpecialCharTok{{-}}\FunctionTok{log\_pi}\NormalTok{(th)    }\SpecialCharTok{+} \FloatTok{0.5} \SpecialCharTok{*} \FunctionTok{sum}\NormalTok{(p  }\SpecialCharTok{*}\NormalTok{ (Minv }\SpecialCharTok{\%*\%}\NormalTok{ p))}
        \ControlFlowTok{if}\NormalTok{ (}\FunctionTok{log}\NormalTok{(}\FunctionTok{runif}\NormalTok{(}\DecValTok{1}\NormalTok{)) }\SpecialCharTok{\textless{}}\NormalTok{ H0 }\SpecialCharTok{{-}}\NormalTok{ H1) \{ theta }\OtherTok{\textless{}{-}}\NormalTok{ th; n\_acc }\OtherTok{\textless{}{-}}\NormalTok{ n\_acc }\SpecialCharTok{+} \DecValTok{1}\NormalTok{ \}}
\NormalTok{        out[t, ] }\OtherTok{\textless{}{-}}\NormalTok{ theta}
\NormalTok{    \}}
    \FunctionTok{attr}\NormalTok{(out, }\StringTok{"accept\_rate"}\NormalTok{) }\OtherTok{\textless{}{-}}\NormalTok{ n\_acc }\SpecialCharTok{/}\NormalTok{ iters}
\NormalTok{    out}
\NormalTok{\}}
\end{Highlighting}
\end{Shaded}

\subsection{HMC versus random-walk Metropolis on a correlated
target}\label{hmc-versus-random-walk-metropolis-on-a-correlated-target}

\begin{Shaded}
\begin{Highlighting}[]
\FunctionTok{set.seed}\NormalTok{(}\DecValTok{4}\NormalTok{)}
\NormalTok{Sigma }\OtherTok{\textless{}{-}} \FunctionTok{matrix}\NormalTok{(}\FunctionTok{c}\NormalTok{(}\DecValTok{1}\NormalTok{, }\FloatTok{0.98}\NormalTok{, }\FloatTok{0.98}\NormalTok{, }\DecValTok{1}\NormalTok{), }\DecValTok{2}\NormalTok{); Sinv }\OtherTok{\textless{}{-}} \FunctionTok{solve}\NormalTok{(Sigma)}
\NormalTok{log\_pi\_corr  }\OtherTok{\textless{}{-}} \ControlFlowTok{function}\NormalTok{(z) }\SpecialCharTok{{-}}\FloatTok{0.5} \SpecialCharTok{*} \FunctionTok{sum}\NormalTok{(z }\SpecialCharTok{*}\NormalTok{ (Sinv }\SpecialCharTok{\%*\%}\NormalTok{ z))}
\NormalTok{grad\_pi\_corr }\OtherTok{\textless{}{-}} \ControlFlowTok{function}\NormalTok{(z) }\SpecialCharTok{{-}}\FunctionTok{as.vector}\NormalTok{(Sinv }\SpecialCharTok{\%*\%}\NormalTok{ z)}

\NormalTok{rw }\OtherTok{\textless{}{-}} \FunctionTok{matrix}\NormalTok{(}\ConstantTok{NA\_real\_}\NormalTok{, }\DecValTok{2000}\NormalTok{, }\DecValTok{2}\NormalTok{); cur }\OtherTok{\textless{}{-}} \FunctionTok{c}\NormalTok{(}\SpecialCharTok{{-}}\DecValTok{2}\NormalTok{, }\SpecialCharTok{{-}}\DecValTok{2}\NormalTok{); lp }\OtherTok{\textless{}{-}} \FunctionTok{log\_pi\_corr}\NormalTok{(cur); acc }\OtherTok{\textless{}{-}} \DecValTok{0}
\ControlFlowTok{for}\NormalTok{ (t }\ControlFlowTok{in} \DecValTok{1}\SpecialCharTok{:}\DecValTok{2000}\NormalTok{) \{}
\NormalTok{    pr }\OtherTok{\textless{}{-}}\NormalTok{ cur }\SpecialCharTok{+} \FloatTok{0.3} \SpecialCharTok{*} \FunctionTok{rnorm}\NormalTok{(}\DecValTok{2}\NormalTok{); lpp }\OtherTok{\textless{}{-}} \FunctionTok{log\_pi\_corr}\NormalTok{(pr)}
    \ControlFlowTok{if}\NormalTok{ (}\FunctionTok{log}\NormalTok{(}\FunctionTok{runif}\NormalTok{(}\DecValTok{1}\NormalTok{)) }\SpecialCharTok{\textless{}}\NormalTok{ lpp }\SpecialCharTok{{-}}\NormalTok{ lp) \{ cur }\OtherTok{\textless{}{-}}\NormalTok{ pr; lp }\OtherTok{\textless{}{-}}\NormalTok{ lpp; acc }\OtherTok{\textless{}{-}}\NormalTok{ acc }\SpecialCharTok{+} \DecValTok{1}\NormalTok{ \}}
\NormalTok{    rw[t, ] }\OtherTok{\textless{}{-}}\NormalTok{ cur}
\NormalTok{\}}
\NormalTok{draws\_hmc }\OtherTok{\textless{}{-}} \FunctionTok{hmc}\NormalTok{(log\_pi\_corr, grad\_pi\_corr, }\AttributeTok{theta =} \FunctionTok{c}\NormalTok{(}\SpecialCharTok{{-}}\DecValTok{2}\NormalTok{, }\SpecialCharTok{{-}}\DecValTok{2}\NormalTok{), }\AttributeTok{eps =} \FloatTok{0.05}\NormalTok{, }\AttributeTok{L =} \DecValTok{30}\NormalTok{, }\AttributeTok{iters =} \DecValTok{2000}\NormalTok{)}

\NormalTok{g }\OtherTok{\textless{}{-}} \FunctionTok{seq}\NormalTok{(}\SpecialCharTok{{-}}\FloatTok{3.5}\NormalTok{, }\FloatTok{3.5}\NormalTok{, }\AttributeTok{length.out =} \DecValTok{80}\NormalTok{)}
\NormalTok{Z }\OtherTok{\textless{}{-}} \FunctionTok{outer}\NormalTok{(g, g, }\FunctionTok{Vectorize}\NormalTok{(}\ControlFlowTok{function}\NormalTok{(a, b) }\FunctionTok{exp}\NormalTok{(}\SpecialCharTok{{-}}\FloatTok{0.5} \SpecialCharTok{*} \FunctionTok{c}\NormalTok{(a, b) }\SpecialCharTok{\%*\%}\NormalTok{ Sinv }\SpecialCharTok{\%*\%} \FunctionTok{c}\NormalTok{(a, b))))}
\FunctionTok{par}\NormalTok{(}\AttributeTok{mfrow =} \FunctionTok{c}\NormalTok{(}\DecValTok{1}\NormalTok{, }\DecValTok{3}\NormalTok{), }\AttributeTok{mar =} \FunctionTok{c}\NormalTok{(}\DecValTok{4}\NormalTok{, }\DecValTok{4}\NormalTok{, }\DecValTok{2}\NormalTok{, }\DecValTok{1}\NormalTok{))}
\ControlFlowTok{for}\NormalTok{ (nm }\ControlFlowTok{in} \FunctionTok{c}\NormalTok{(}\StringTok{"random walk"}\NormalTok{, }\StringTok{"HMC"}\NormalTok{)) \{}
\NormalTok{    d }\OtherTok{\textless{}{-}} \ControlFlowTok{if}\NormalTok{ (nm }\SpecialCharTok{==} \StringTok{"random walk"}\NormalTok{) rw }\ControlFlowTok{else}\NormalTok{ draws\_hmc}
\NormalTok{    ar }\OtherTok{\textless{}{-}} \ControlFlowTok{if}\NormalTok{ (nm }\SpecialCharTok{==} \StringTok{"random walk"}\NormalTok{) acc }\SpecialCharTok{/} \DecValTok{2000} \ControlFlowTok{else} \FunctionTok{attr}\NormalTok{(draws\_hmc, }\StringTok{"accept\_rate"}\NormalTok{)}
    \FunctionTok{contour}\NormalTok{(g, g, Z, }\AttributeTok{nlevels =} \DecValTok{6}\NormalTok{, }\AttributeTok{drawlabels =} \ConstantTok{FALSE}\NormalTok{, }\AttributeTok{col =} \StringTok{"gray60"}\NormalTok{,}
            \AttributeTok{xlab =} \FunctionTok{expression}\NormalTok{(theta[}\DecValTok{1}\NormalTok{]), }\AttributeTok{ylab =} \FunctionTok{expression}\NormalTok{(theta[}\DecValTok{2}\NormalTok{]),}
            \AttributeTok{main =} \FunctionTok{sprintf}\NormalTok{(}\StringTok{"\%s: accept \%.2f"}\NormalTok{, nm, ar))}
    \FunctionTok{lines}\NormalTok{(d[}\DecValTok{1}\SpecialCharTok{:}\DecValTok{100}\NormalTok{, ], }\AttributeTok{col =} \ControlFlowTok{if}\NormalTok{ (nm }\SpecialCharTok{==} \StringTok{"random walk"}\NormalTok{) }\StringTok{"firebrick"} \ControlFlowTok{else} \StringTok{"steelblue"}\NormalTok{, }\AttributeTok{lwd =} \FloatTok{1.5}\NormalTok{)}
    \FunctionTok{points}\NormalTok{(d[}\DecValTok{1}\SpecialCharTok{:}\DecValTok{100}\NormalTok{, ], }\AttributeTok{pch =} \DecValTok{19}\NormalTok{, }\AttributeTok{cex =} \FloatTok{0.5}\NormalTok{, }\AttributeTok{col =} \ControlFlowTok{if}\NormalTok{ (nm }\SpecialCharTok{==} \StringTok{"random walk"}\NormalTok{) }\StringTok{"firebrick"} \ControlFlowTok{else} \StringTok{"steelblue"}\NormalTok{)}
\NormalTok{\}}
\NormalTok{a1 }\OtherTok{\textless{}{-}} \FunctionTok{acf}\NormalTok{(rw[, }\DecValTok{1}\NormalTok{], }\AttributeTok{lag.max =} \DecValTok{60}\NormalTok{, }\AttributeTok{plot =} \ConstantTok{FALSE}\NormalTok{)}\SpecialCharTok{$}\NormalTok{acf}
\NormalTok{a2 }\OtherTok{\textless{}{-}} \FunctionTok{acf}\NormalTok{(draws\_hmc[, }\DecValTok{1}\NormalTok{], }\AttributeTok{lag.max =} \DecValTok{60}\NormalTok{, }\AttributeTok{plot =} \ConstantTok{FALSE}\NormalTok{)}\SpecialCharTok{$}\NormalTok{acf}
\FunctionTok{plot}\NormalTok{(}\DecValTok{0}\SpecialCharTok{:}\DecValTok{60}\NormalTok{, a1, }\AttributeTok{type =} \StringTok{"h"}\NormalTok{, }\AttributeTok{lwd =} \DecValTok{2}\NormalTok{, }\AttributeTok{col =} \StringTok{"firebrick"}\NormalTok{, }\AttributeTok{xlab =} \StringTok{"lag"}\NormalTok{, }\AttributeTok{ylab =} \StringTok{"ACF"}\NormalTok{, }\AttributeTok{ylim =} \FunctionTok{c}\NormalTok{(}\SpecialCharTok{{-}}\FloatTok{0.2}\NormalTok{, }\DecValTok{1}\NormalTok{), }\AttributeTok{main =} \StringTok{"autocorrelation"}\NormalTok{)}
\FunctionTok{lines}\NormalTok{((}\DecValTok{0}\SpecialCharTok{:}\DecValTok{60}\NormalTok{) }\SpecialCharTok{+} \FloatTok{0.3}\NormalTok{, a2, }\AttributeTok{type =} \StringTok{"h"}\NormalTok{, }\AttributeTok{lwd =} \DecValTok{2}\NormalTok{, }\AttributeTok{col =} \StringTok{"steelblue"}\NormalTok{)}
\FunctionTok{legend}\NormalTok{(}\StringTok{"topright"}\NormalTok{, }\AttributeTok{bty =} \StringTok{"n"}\NormalTok{, }\AttributeTok{lwd =} \DecValTok{2}\NormalTok{, }\AttributeTok{col =} \FunctionTok{c}\NormalTok{(}\StringTok{"firebrick"}\NormalTok{, }\StringTok{"steelblue"}\NormalTok{), }\AttributeTok{legend =} \FunctionTok{c}\NormalTok{(}\StringTok{"random walk"}\NormalTok{, }\StringTok{"HMC"}\NormalTok{))}
\end{Highlighting}
\end{Shaded}

\begin{figure}[H]

\centering{

\pandocbounded{\includegraphics[keepaspectratio]{chapters/unit11_mcmc_files/figure-pdf/fig-hmc-vs-rw-1.pdf}}

}

\caption{\label{fig-hmc-vs-rw}Target: bivariate normal with correlation
0.98. Left/middle: first 100 draws over the posterior contours for
random-walk Metropolis and HMC, same iteration budget. Right:
autocorrelation of the first coordinate -- HMC crosses the ridge in one
trajectory and its draws decorrelate almost immediately.}

\end{figure}%

\subsection{Explore it interactively}\label{explore-it-interactively-3}

Leapfrog step size \(\epsilon\) and step count \(L\) both need tuning:
too large an \(\epsilon\) makes the discretized trajectory unstable
(rejections rise sharply); too small an \(L\) reduces HMC to a random
walk; too large an \(L\) lets the trajectory turn back on itself and
waste computation. The app below runs the \texttt{hmc()} function above
on the same correlated target and lets you vary \(\epsilon\) and \(L\)
directly.

\begin{Shaded}
\begin{Highlighting}[]
\NormalTok{\#| label: hmc{-}tuning{-}app}
\NormalTok{\#| standalone: true}
\NormalTok{\#| viewerHeight: 620}

\NormalTok{library(shiny)}

\NormalTok{Sigma \textless{}{-} matrix(c(1, 0.98, 0.98, 1), 2)}
\NormalTok{Sinv  \textless{}{-} solve(Sigma)}
\NormalTok{log\_pi\_corr  \textless{}{-} function(z) {-}0.5 * sum(z * (Sinv \%*\% z))}
\NormalTok{grad\_pi\_corr \textless{}{-} function(z) {-}as.vector(Sinv \%*\% z)}

\NormalTok{hmc \textless{}{-} function(log\_pi, grad\_log\_pi, theta, eps, L, iters) \{}
\NormalTok{    d \textless{}{-} length(theta)}
\NormalTok{    out \textless{}{-} matrix(NA\_real\_, iters, d); n\_acc \textless{}{-} 0}
\NormalTok{    for (t in seq\_len(iters)) \{}
\NormalTok{        p0 \textless{}{-} rnorm(d)}
\NormalTok{        th \textless{}{-} theta; p \textless{}{-} p0 + eps / 2 * grad\_log\_pi(th)}
\NormalTok{        for (l in seq\_len(L)) \{}
\NormalTok{            th \textless{}{-} th + eps * p}
\NormalTok{            if (l \textless{} L) p \textless{}{-} p + eps * grad\_log\_pi(th)}
\NormalTok{        \}}
\NormalTok{        p \textless{}{-} p + eps / 2 * grad\_log\_pi(th)}
\NormalTok{        H0 \textless{}{-} {-}log\_pi(theta) + 0.5 * sum(p0\^{}2)}
\NormalTok{        H1 \textless{}{-} {-}log\_pi(th)    + 0.5 * sum(p\^{}2)}
\NormalTok{        if (log(runif(1)) \textless{} H0 {-} H1) \{ theta \textless{}{-} th; n\_acc \textless{}{-} n\_acc + 1 \}}
\NormalTok{        out[t, ] \textless{}{-} theta}
\NormalTok{    \}}
\NormalTok{    attr(out, "accept\_rate") \textless{}{-} n\_acc / iters}
\NormalTok{    out}
\NormalTok{\}}

\NormalTok{ui \textless{}{-} fluidPage(}
\NormalTok{    titlePanel(NULL),}
\NormalTok{    sidebarLayout(}
\NormalTok{        sidebarPanel(}
\NormalTok{            width = 4,}
\NormalTok{            sliderInput("eps", "leapfrog step size (epsilon)", min = 0.01, max = 1, value = 0.05, step = 0.01),}
\NormalTok{            sliderInput("L", "number of leapfrog steps (L)", min = 1, max = 60, value = 30, step = 1),}
\NormalTok{            sliderInput("iters", "iterations", min = 200, max = 3000, value = 1000, step = 200),}
\NormalTok{            helpText("Target: bivariate normal with correlation 0.98. Try eps = 0.05, ",}
\NormalTok{                     "L = 30 (efficient), then eps = 0.5 (unstable, low acceptance), ",}
\NormalTok{                     "then L = 2 (random{-}walk{-}like, high autocorrelation)."),}
\NormalTok{            tableOutput("tab")}
\NormalTok{        ),}
\NormalTok{        mainPanel(}
\NormalTok{            width = 8,}
\NormalTok{            plotOutput("traj", height = "360px"),}
\NormalTok{            plotOutput("acfplot", height = "220px")}
\NormalTok{        )}
\NormalTok{    )}
\NormalTok{)}

\NormalTok{server \textless{}{-} function(input, output, session) \{}

\NormalTok{    fit \textless{}{-} reactive(\{}
\NormalTok{        set.seed(1)}
\NormalTok{        hmc(log\_pi\_corr, grad\_pi\_corr, theta = c({-}2, {-}2), eps = input$eps, L = input$L, iters = input$iters)}
\NormalTok{    \})}

\NormalTok{    output$traj \textless{}{-} renderPlot(\{}
\NormalTok{        d \textless{}{-} fit()}
\NormalTok{        g \textless{}{-} seq({-}3.5, 3.5, length.out = 60)}
\NormalTok{        Z \textless{}{-} outer(g, g, Vectorize(function(a, b) exp({-}0.5 * c(a, b) \%*\% Sinv \%*\% c(a, b))))}
\NormalTok{        par(mar = c(4, 4, 2, 1))}
\NormalTok{        contour(g, g, Z, nlevels = 6, drawlabels = FALSE, col = "gray60",}
\NormalTok{                xlab = expression(theta[1]), ylab = expression(theta[2]),}
\NormalTok{                main = sprintf("accept rate: \%.2f", attr(d, "accept\_rate")))}
\NormalTok{        k \textless{}{-} min(100, nrow(d))}
\NormalTok{        lines(d[1:k, ], col = "steelblue", lwd = 1.2)}
\NormalTok{        points(d[1:k, ], pch = 19, cex = 0.5, col = "steelblue")}
\NormalTok{    \})}

\NormalTok{    output$acfplot \textless{}{-} renderPlot(\{}
\NormalTok{        d \textless{}{-} fit()}
\NormalTok{        par(mar = c(4, 4, 1, 1))}
\NormalTok{        acf(d[, 1], lag.max = 60, col = "steelblue", main = "autocorrelation of theta\_1")}
\NormalTok{    \})}

\NormalTok{    output$tab \textless{}{-} renderTable(\{}
\NormalTok{        d \textless{}{-} fit()}
\NormalTok{        data.frame(quantity = c("acceptance rate", "mean theta\_1", "sd theta\_1"),}
\NormalTok{                   value = c(attr(d, "accept\_rate"), mean(d[, 1]), sd(d[, 1])))}
\NormalTok{    \}, digits = 3, rownames = FALSE)}
\NormalTok{\}}

\NormalTok{shinyApp(ui, server)}
\end{Highlighting}
\end{Shaded}

\begin{tcolorbox}[enhanced jigsaw, arc=.35mm, bottomrule=.15mm, bottomtitle=1mm, breakable, colback=white, colbacktitle=quarto-callout-note-color!10!white, colframe=quarto-callout-note-color-frame, coltitle=black, left=2mm, leftrule=.75mm, opacityback=0, opacitybacktitle=0.6, rightrule=.15mm, title=\textcolor{quarto-callout-note-color}{\faInfo}\hspace{0.5em}{Reproducing this figure}, titlerule=0mm, toprule=.15mm, toptitle=1mm]

The chunk above uses the
\href{https://github.com/quarto-ext/shinylive}{shinylive} extension,
which compiles the app to WebAssembly so it runs in the reader's browser
with no Shiny server. Install it once per project with

\begin{Shaded}
\begin{Highlighting}[]
\ExtensionTok{quarto}\NormalTok{ add quarto{-}ext/shinylive}
\end{Highlighting}
\end{Shaded}

and add \texttt{filters:\ {[}shinylive{]}} to the document or project
YAML. Only packages available in webR may be used inside the app, so the
code above is restricted to base R and \texttt{shiny}.

\end{tcolorbox}

\section{Stan and Probabilistic
Programming}\label{stan-and-probabilistic-programming}

\subsection{Automating HMC}\label{automating-hmc}

Every sampler above needed something derived by hand: Gibbs needed full
conditionals worked out algebraically; Metropolis--Hastings needed only
a target density but still needed step sizes tuned by trial and error;
HMC needed an analytic gradient and needed \(\epsilon\) and \(L\) tuned
by hand. \textbf{Stan} is a probabilistic programming language that
automates all of this: given a model written declaratively (data,
parameters, and a log-density), it computes gradients by automatic
differentiation and runs an adaptive variant of HMC --- the
\textbf{No-U-Turn Sampler (NUTS)}, which picks a good \(L\) on the fly
by extending the leapfrog trajectory until it starts to double back ---
with step size tuned during a warm-up phase. The user never derives a
conditional or a gradient.

\subsection{The running example, in
Stan}\label{the-running-example-in-stan}

To confirm Stan lands on the same posterior as the Gibbs and MH samplers
above, fit the identical normal-mean/variance model with the identical
(vague) priors and data:

\begin{Shaded}
\begin{Highlighting}[]
\NormalTok{stan\_normal\_code }\OtherTok{\textless{}{-}} \StringTok{"}
\StringTok{data \{}
\StringTok{  int\textless{}lower=0\textgreater{} n;}
\StringTok{  vector[n] x;}
\StringTok{  real mu0;}
\StringTok{  real\textless{}lower=0\textgreater{} sigma0;}
\StringTok{  real\textless{}lower=0\textgreater{} alpha;}
\StringTok{  real\textless{}lower=0\textgreater{} w;}
\StringTok{\}}
\StringTok{parameters \{}
\StringTok{  real mu;}
\StringTok{  real\textless{}lower=0\textgreater{} sigma2;}
\StringTok{\}}
\StringTok{model \{}
\StringTok{  mu \textasciitilde{} normal(mu0, sqrt(sigma0));}
\StringTok{  sigma2 \textasciitilde{} inv\_gamma(alpha / 2, alpha * w / 2);}
\StringTok{  x \textasciitilde{} normal(mu, sqrt(sigma2));}
\StringTok{\}}
\StringTok{"}
\end{Highlighting}
\end{Shaded}

\begin{Shaded}
\begin{Highlighting}[]
\NormalTok{fit\_stan }\OtherTok{\textless{}{-}} \FunctionTok{stan}\NormalTok{(}
    \AttributeTok{model\_code =}\NormalTok{ stan\_normal\_code,}
    \AttributeTok{data =} \FunctionTok{list}\NormalTok{(}\AttributeTok{n =} \FunctionTok{length}\NormalTok{(x), }\AttributeTok{x =}\NormalTok{ x, }\AttributeTok{mu0 =} \DecValTok{0}\NormalTok{, }\AttributeTok{sigma0 =} \FloatTok{1e10}\NormalTok{, }\AttributeTok{alpha =} \FloatTok{1e{-}5}\NormalTok{, }\AttributeTok{w =} \FloatTok{1e{-}5}\NormalTok{),}
    \AttributeTok{chains =} \DecValTok{4}\NormalTok{, }\AttributeTok{iter =} \DecValTok{2000}\NormalTok{, }\AttributeTok{warmup =} \DecValTok{1000}\NormalTok{, }\AttributeTok{seed =} \DecValTok{5}\NormalTok{, }\AttributeTok{refresh =} \DecValTok{0}
\NormalTok{)}
\end{Highlighting}
\end{Shaded}

\begin{Shaded}
\begin{Highlighting}[]
\NormalTok{stan\_summary }\OtherTok{\textless{}{-}} \FunctionTok{summary}\NormalTok{(fit\_stan, }\AttributeTok{pars =} \FunctionTok{c}\NormalTok{(}\StringTok{"mu"}\NormalTok{, }\StringTok{"sigma2"}\NormalTok{))}\SpecialCharTok{$}\NormalTok{summary}
\NormalTok{stan\_sigma\_draws }\OtherTok{\textless{}{-}} \FunctionTok{sqrt}\NormalTok{(}\FunctionTok{as.matrix}\NormalTok{(fit\_stan, }\AttributeTok{pars =} \StringTok{"sigma2"}\NormalTok{)[, }\DecValTok{1}\NormalTok{])}
\NormalTok{three\_way }\OtherTok{\textless{}{-}} \FunctionTok{data.frame}\NormalTok{(}
    \AttributeTok{method =} \FunctionTok{c}\NormalTok{(}\StringTok{"Gibbs"}\NormalTok{, }\StringTok{"Metropolis{-}Hastings"}\NormalTok{, }\StringTok{"Stan (NUTS)"}\NormalTok{),}
    \AttributeTok{mean\_mu =} \FunctionTok{c}\NormalTok{(}\FunctionTok{mean}\NormalTok{(mu\_g), }\FunctionTok{mean}\NormalTok{(mh\_long}\SpecialCharTok{$}\NormalTok{mu[keep\_mh]), stan\_summary[}\StringTok{"mu"}\NormalTok{, }\StringTok{"mean"}\NormalTok{]),}
    \AttributeTok{sd\_mu   =} \FunctionTok{c}\NormalTok{(}\FunctionTok{sd}\NormalTok{(mu\_g), }\FunctionTok{sd}\NormalTok{(mh\_long}\SpecialCharTok{$}\NormalTok{mu[keep\_mh]), stan\_summary[}\StringTok{"mu"}\NormalTok{, }\StringTok{"sd"}\NormalTok{]),}
    \AttributeTok{mean\_sigma =} \FunctionTok{c}\NormalTok{(}\FunctionTok{mean}\NormalTok{(sd\_g), }\FunctionTok{mean}\NormalTok{(mh\_long}\SpecialCharTok{$}\NormalTok{sd[keep\_mh]), }\FunctionTok{mean}\NormalTok{(stan\_sigma\_draws)),}
    \AttributeTok{sd\_sigma   =} \FunctionTok{c}\NormalTok{(}\FunctionTok{sd}\NormalTok{(sd\_g), }\FunctionTok{sd}\NormalTok{(mh\_long}\SpecialCharTok{$}\NormalTok{sd[keep\_mh]), }\FunctionTok{sd}\NormalTok{(stan\_sigma\_draws))}
\NormalTok{)}
\NormalTok{knitr}\SpecialCharTok{::}\FunctionTok{kable}\NormalTok{(three\_way, }\AttributeTok{digits =} \DecValTok{3}\NormalTok{, }\AttributeTok{align =} \FunctionTok{c}\NormalTok{(}\StringTok{"l"}\NormalTok{, }\StringTok{"r"}\NormalTok{, }\StringTok{"r"}\NormalTok{, }\StringTok{"r"}\NormalTok{, }\StringTok{"r"}\NormalTok{))}
\end{Highlighting}
\end{Shaded}

{

\begin{longtable}[]{@{}lrrrr@{}}

\caption{\label{tbl-stan-vs-hand}Posterior mean and sd: hand-written
Gibbs and Metropolis-Hastings versus Stan's NUTS sampler, same data,
priors, and model. All three samplers target the identical posterior and
agree up to Monte Carlo error.}

\tabularnewline

\toprule\noalign{}
method & mean\_mu & sd\_mu & mean\_sigma & sd\_sigma \\
\midrule\noalign{}
\endhead
\bottomrule\noalign{}
\endlastfoot
Gibbs & 10.141 & 0.330 & 2.293 & 0.239 \\
Metropolis-Hastings & 10.144 & 0.330 & 2.294 & 0.237 \\
Stan (NUTS) & 10.137 & 0.329 & 2.301 & 0.245 \\

\end{longtable}

}

The three completely independent implementations --- hand-derived
conditionals, a hand-tuned random walk, and Stan's automatic NUTS ---
agree to within Monte Carlo noise, which is exactly the coherence check
this book has used throughout: independent methods with no shared code
path landing on the same answer is evidence that none of them is wrong.

\subsection{Where Stan earns its keep}\label{where-stan-earns-its-keep}

The normal-mean/variance model above is simple enough that Gibbs
sampling is not just possible but easy, so Stan's advantage is not
visible yet. Its real value shows up on models where hand-deriving
conditionals or gradients would be a serious undertaking: hierarchical
models with many levels of latent structure, time-series models with an
unobserved latent state at every time point, or anything with a
non-conjugate likelihood/prior combination. A stochastic-volatility
model --- where a latent log-variance \(h_t\) follows its own
autoregressive process and each observation's variance is \(e^{h_t}\)
--- is a standard example of this kind: writing a Gibbs sampler for it
requires deriving a conditional for every latent \(h_t\) given its
neighbours, which is neither conjugate nor pleasant, while in Stan it
only needs the model written down.

A first, direct transcription of the model --- sampling each \(h_t\)
centred on \(\mu + \phi(h_{t-1}-\mu)\) exactly as the model
specification reads --- compiles and runs, but Stan reports \(\hat R\)
up to \(1.24\), effective sample sizes in the tens out of a thousand
draws, and a low-BFMI warning: the chains have not mixed. This is not a
bug in Stan but a known geometric pathology (a ``funnel'': when \(\tau\)
is small the conditional distribution of each \(h_t\) given its
neighbours becomes extremely narrow, and a constant leapfrog step size
cannot be simultaneously right for the narrow and wide parts of the
posterior). The standard fix is a \textbf{non-centred
reparameterization}: sample standard-normal innovations
\(h_t^{\text{raw}}\) and reconstruct \(h_t\) deterministically, \[
h_t^{\text{raw}} \sim N(0,1), \qquad h_t = \mu + \phi(h_{t-1}-\mu) + \tau\, h_t^{\text{raw}},
\] which moves the awkward \(\tau\)-dependent scaling out of the sampled
parameters and into a deterministic transform --- the same idea as the
\(\log\sigma\) and logit reparameterizations used earlier in this book
to keep a sampler's geometry well behaved, applied here to a latent time
series instead of a single scale parameter.

\begin{Shaded}
\begin{Highlighting}[]
\NormalTok{stan\_sv\_code }\OtherTok{\textless{}{-}} \StringTok{"}
\StringTok{data \{}
\StringTok{  int\textless{}lower=0\textgreater{} T;}
\StringTok{  vector[T] y;}
\StringTok{\}}
\StringTok{parameters \{}
\StringTok{  real mu;}
\StringTok{  real\textless{}lower={-}1,upper=1\textgreater{} phi;}
\StringTok{  real\textless{}lower=0.01\textgreater{} tausq;}
\StringTok{  real h0\_raw;}
\StringTok{  vector[T] h\_raw;}
\StringTok{\}}
\StringTok{transformed parameters \{}
\StringTok{  real\textless{}lower=0\textgreater{} tau = sqrt(tausq);}
\StringTok{  real h0 = mu + tau * h0\_raw;}
\StringTok{  vector[T] h;}
\StringTok{  h[1] = mu + phi * (h0 {-} mu) + tau * h\_raw[1];}
\StringTok{  for (t in 2:T)}
\StringTok{    h[t] = mu + phi * (h[t {-} 1] {-} mu) + tau * h\_raw[t];}
\StringTok{\}}
\StringTok{model \{}
\StringTok{  mu \textasciitilde{} normal({-}10, 5);}
\StringTok{  phi \textasciitilde{} uniform({-}1, 1);}
\StringTok{  tausq \textasciitilde{} inv\_gamma(2.5, 0.025);}
\StringTok{  h0\_raw \textasciitilde{} std\_normal();}
\StringTok{  h\_raw \textasciitilde{} std\_normal();}
\StringTok{  y \textasciitilde{} normal(0, exp(h / 2));}
\StringTok{\}}
\StringTok{"}
\end{Highlighting}
\end{Shaded}

\begin{Shaded}
\begin{Highlighting}[]
\FunctionTok{set.seed}\NormalTok{(}\DecValTok{6}\NormalTok{)}
\NormalTok{N }\OtherTok{\textless{}{-}} \DecValTok{200}\NormalTok{; mu\_sv }\OtherTok{\textless{}{-}} \DecValTok{3}\NormalTok{; phi\_sv }\OtherTok{\textless{}{-}} \FloatTok{0.3}\NormalTok{; tau\_sv }\OtherTok{\textless{}{-}} \DecValTok{1}\NormalTok{; h0\_sv }\OtherTok{\textless{}{-}} \DecValTok{2}
\NormalTok{h }\OtherTok{\textless{}{-}} \FunctionTok{numeric}\NormalTok{(N); y }\OtherTok{\textless{}{-}} \FunctionTok{numeric}\NormalTok{(N)}
\NormalTok{h[}\DecValTok{1}\NormalTok{] }\OtherTok{\textless{}{-}} \FunctionTok{rnorm}\NormalTok{(}\DecValTok{1}\NormalTok{, mu\_sv }\SpecialCharTok{+}\NormalTok{ phi\_sv }\SpecialCharTok{*}\NormalTok{ (h0\_sv }\SpecialCharTok{{-}}\NormalTok{ mu\_sv), tau\_sv)}
\ControlFlowTok{for}\NormalTok{ (i }\ControlFlowTok{in} \DecValTok{2}\SpecialCharTok{:}\NormalTok{N) h[i] }\OtherTok{\textless{}{-}} \FunctionTok{rnorm}\NormalTok{(}\DecValTok{1}\NormalTok{, mu\_sv }\SpecialCharTok{+}\NormalTok{ phi\_sv }\SpecialCharTok{*}\NormalTok{ (h[i }\SpecialCharTok{{-}} \DecValTok{1}\NormalTok{] }\SpecialCharTok{{-}}\NormalTok{ mu\_sv), tau\_sv)}
\ControlFlowTok{for}\NormalTok{ (t }\ControlFlowTok{in} \DecValTok{1}\SpecialCharTok{:}\NormalTok{N) y[t] }\OtherTok{\textless{}{-}} \FunctionTok{rnorm}\NormalTok{(}\DecValTok{1}\NormalTok{, }\DecValTok{0}\NormalTok{, }\FunctionTok{exp}\NormalTok{(h[t] }\SpecialCharTok{/} \DecValTok{2}\NormalTok{))}

\NormalTok{fit\_sv }\OtherTok{\textless{}{-}} \FunctionTok{stan}\NormalTok{(}
    \AttributeTok{model\_code =}\NormalTok{ stan\_sv\_code, }\AttributeTok{data =} \FunctionTok{list}\NormalTok{(}\AttributeTok{y =}\NormalTok{ y, }\AttributeTok{T =}\NormalTok{ N),}
    \AttributeTok{chains =} \DecValTok{2}\NormalTok{, }\AttributeTok{iter =} \DecValTok{1000}\NormalTok{, }\AttributeTok{warmup =} \DecValTok{500}\NormalTok{, }\AttributeTok{seed =} \DecValTok{6}\NormalTok{, }\AttributeTok{refresh =} \DecValTok{0}
\NormalTok{)}
\end{Highlighting}
\end{Shaded}

\begin{verbatim}
Warning: Bulk Effective Samples Size (ESS) is too low, indicating posterior means and medians may be unreliable.
Running the chains for more iterations may help. See
https://mc-stan.org/misc/warnings.html#bulk-ess
\end{verbatim}

\begin{Shaded}
\begin{Highlighting}[]
\FunctionTok{par}\NormalTok{(}\AttributeTok{mfrow =} \FunctionTok{c}\NormalTok{(}\DecValTok{1}\NormalTok{, }\DecValTok{3}\NormalTok{), }\AttributeTok{mar =} \FunctionTok{c}\NormalTok{(}\DecValTok{4}\NormalTok{, }\DecValTok{4}\NormalTok{, }\DecValTok{2}\NormalTok{, }\DecValTok{1}\NormalTok{))}
\FunctionTok{plot}\NormalTok{(y, }\AttributeTok{type =} \StringTok{"l"}\NormalTok{, }\AttributeTok{col =} \StringTok{"steelblue"}\NormalTok{, }\AttributeTok{xlab =} \StringTok{"t"}\NormalTok{, }\AttributeTok{ylab =} \StringTok{"y"}\NormalTok{, }\AttributeTok{main =} \StringTok{"simulated returns"}\NormalTok{)}
\NormalTok{sv\_summary }\OtherTok{\textless{}{-}} \FunctionTok{summary}\NormalTok{(fit\_sv, }\AttributeTok{pars =} \FunctionTok{c}\NormalTok{(}\StringTok{"mu"}\NormalTok{, }\StringTok{"phi"}\NormalTok{, }\StringTok{"tau"}\NormalTok{))}\SpecialCharTok{$}\NormalTok{summary}
\FunctionTok{hist}\NormalTok{(}\FunctionTok{as.matrix}\NormalTok{(fit\_sv, }\AttributeTok{pars =} \StringTok{"phi"}\NormalTok{), }\AttributeTok{breaks =} \DecValTok{30}\NormalTok{, }\AttributeTok{col =} \StringTok{"gray85"}\NormalTok{, }\AttributeTok{border =} \StringTok{"white"}\NormalTok{,}
     \AttributeTok{main =} \StringTok{"posterior of phi (true 0.3)"}\NormalTok{, }\AttributeTok{xlab =} \FunctionTok{expression}\NormalTok{(phi))}
\FunctionTok{abline}\NormalTok{(}\AttributeTok{v =}\NormalTok{ phi\_sv, }\AttributeTok{col =} \StringTok{"firebrick"}\NormalTok{, }\AttributeTok{lwd =} \DecValTok{2}\NormalTok{)}
\FunctionTok{hist}\NormalTok{(}\FunctionTok{as.matrix}\NormalTok{(fit\_sv, }\AttributeTok{pars =} \StringTok{"tau"}\NormalTok{), }\AttributeTok{breaks =} \DecValTok{30}\NormalTok{, }\AttributeTok{col =} \StringTok{"gray85"}\NormalTok{, }\AttributeTok{border =} \StringTok{"white"}\NormalTok{,}
     \AttributeTok{main =} \StringTok{"posterior of tau (true 1)"}\NormalTok{, }\AttributeTok{xlab =} \FunctionTok{expression}\NormalTok{(tau))}
\FunctionTok{abline}\NormalTok{(}\AttributeTok{v =}\NormalTok{ tau\_sv, }\AttributeTok{col =} \StringTok{"firebrick"}\NormalTok{, }\AttributeTok{lwd =} \DecValTok{2}\NormalTok{)}
\end{Highlighting}
\end{Shaded}

\begin{figure}[H]

\centering{

\pandocbounded{\includegraphics[keepaspectratio]{chapters/unit11_mcmc_files/figure-pdf/fig-stan-sv-1.pdf}}

}

\caption{\label{fig-stan-sv}Simulated stochastic-volatility data (left)
and posteriors of the persistence phi and innovation scale tau under the
non-centred parameterization (true values 0.3 and 1, red lines): both
intervals are wide, reflecting genuine weak identification from a single
200-point series, but both cover the truth.}

\end{figure}%

\begin{Shaded}
\begin{Highlighting}[]
\NormalTok{sv\_tbl }\OtherTok{\textless{}{-}} \FunctionTok{data.frame}\NormalTok{(}\AttributeTok{variable =} \FunctionTok{rownames}\NormalTok{(sv\_summary), sv\_summary[, }\FunctionTok{c}\NormalTok{(}\StringTok{"mean"}\NormalTok{, }\StringTok{"sd"}\NormalTok{, }\StringTok{"Rhat"}\NormalTok{, }\StringTok{"n\_eff"}\NormalTok{)],}
                      \AttributeTok{row.names =} \ConstantTok{NULL}\NormalTok{, }\AttributeTok{check.names =} \ConstantTok{FALSE}\NormalTok{)}
\NormalTok{knitr}\SpecialCharTok{::}\FunctionTok{kable}\NormalTok{(sv\_tbl, }\AttributeTok{digits =} \DecValTok{3}\NormalTok{, }\AttributeTok{align =} \FunctionTok{c}\NormalTok{(}\StringTok{"l"}\NormalTok{, }\StringTok{"r"}\NormalTok{, }\StringTok{"r"}\NormalTok{, }\StringTok{"r"}\NormalTok{, }\StringTok{"r"}\NormalTok{))}
\end{Highlighting}
\end{Shaded}

{

\begin{longtable}[]{@{}lrrrr@{}}

\caption{\label{tbl-stan-sv}Posterior summaries for the
stochastic-volatility model's top-level parameters, non-centred
parameterization: Rhat near 1 and effective sample sizes in the hundreds
indicate the chains have converged and mixed, in contrast to the centred
version.}

\tabularnewline

\toprule\noalign{}
variable & mean & sd & Rhat & n\_eff \\
\midrule\noalign{}
\endhead
\bottomrule\noalign{}
\endlastfoot
mu & 3.039 & 0.216 & 1.005 & 394.177 \\
phi & 0.603 & 0.232 & 1.003 & 166.911 \\
tau & 0.618 & 0.203 & 1.002 & 202.156 \\

\end{longtable}

}

The syntax above is modern Stan (\texttt{=} for deterministic
assignment, \texttt{\textasciitilde{}} sampling statements, no
deprecated \texttt{\textless{}-} or \texttt{\_log} density functions,
which older Stan code --- including earlier drafts of this example ---
used and which current Stan releases no longer accept). After the
reparameterization, \(\hat R \approx 1\) and effective sample sizes in
the hundreds confirm the chains have converged; the \(95\%\) intervals
for \(\phi\) and \(\tau\) are wide but do contain their true simulated
values, which is the honest outcome for a weakly identified latent-state
model fit to a single moderate-length series, and a materially better
diagnosis than the centred version's confidently-wrong-looking but
actually-unconverged chains. That a model no one in this chapter
hand-derived a sampler for can be fit --- and, just as importantly,
checked for convergence --- with a few lines of declarative model code
is the point of the whole section.

\section{Summary}\label{summary-6}

\begin{itemize}
\tightlist
\item
  A Markov chain converges to a target \(\pi\) from any starting point
  if its transition \(T\) leaves \(\pi\) invariant and is aperiodic and
  irreducible; \textbf{detailed balance}, a pairwise condition, is the
  standard route to invariance.
\item
  \textbf{Gibbs sampling} cycles through exact draws from each full
  conditional; free of tuning, but needs conditionals that are both
  known and directly sample-able, and data augmentation extends it to
  models with latent variables.
\item
  \textbf{Metropolis--Hastings} accepts or rejects a proposal from
  \emph{any} distribution with probability depending on the
  target-density ratio (times a proposal-symmetry correction);
  random-walk MH needs a tuned step size, aimed at low autocorrelation
  rather than high acceptance.
\item
  \textbf{Hamiltonian Monte Carlo} adds momentum and follows
  gradient-driven trajectories, discretized by the leapfrog integrator
  and corrected by a Metropolis step, giving distant, nearly independent
  proposals even on strongly correlated targets --- at the cost of
  needing a gradient and two tuning parameters (\(\epsilon\), \(L\)).
\item
  \textbf{Stan} automates gradient computation and adaptive HMC (NUTS)
  for essentially any model expressible in its declarative syntax,
  turning what would otherwise be a substantial hand-derivation (as the
  stochastic-volatility example shows) into a direct transcription of
  the model.
\end{itemize}

This is the last of the sampling-methods chapters in the book:
quadrature (Chapter 8) and Laplace's method (Chapter 9) work in low
dimension by grid or by a local Gaussian approximation; rejection and
importance sampling (Chapters 10--11) draw independently from a fixed
proposal; Gibbs, Metropolis--Hastings, HMC, and Stan (this chapter) let
a Markov chain do the work instead, and are, in practice, how most
modern Bayesian computation is actually done.




\end{document}
